\chapter{\texorpdfstring{$K3 \times E$}{K3 x E} and the BKM Superalgebra}
\label{ch:k3e-bkm}
%% Preserve the original chapter label for cross-references
\phantomsection\label{ch:k3-times-e}
\phantomsection\label{thm:denom-bar-euler}

\providecommand{\multEnr}{\mathrm{mult}_{\mathrm{Enr}}}
\providecommand{\multKthree}{\mathrm{mult}_{K3}}
\providecommand{\cKthree}{c_{K3}}
\providecommand{\EnrAdm}{\mathcal{A}_{\mathrm{Enr}}}
\providecommand{\EnrVirt}{\mathcal{V}_{\mathrm{Enr}}}
\providecommand{\CplusEnr}{C_+^{\mathrm{Enr}}}

\begin{theorem}[\texorpdfstring{$K3 \times E$}{K3 x E} four-corner structure]
\label{thm:k3e-four-corner-structure}
At the specialisation datum $(\Sigma_2, C) = (K3, E)$, the two-stage
factorisation output $\SpCh_{K3, E}(\PhiFA_3(D^b(\Coh(K3 \times E))))$
is compared with four a priori distinct constructions:
\begin{enumerate}[label=\textup{(\roman*)}]
\item the Oberdieck--Pandharipande reduced DT moduli
 $\cM^{\mathrm{red}}_{\mathrm{DT}}(K3 \times E;\,\gamma\text{ primitive})$
 as positive effective geometry, equivariant for
 $\bC^\times_E \times \Aut_s(K3)$
 (Oberdieck--Pandharipande,
 \emph{Inventiones Mathematicae} \textbf{222} (2020), 667--714);
\item the finite DWR/Ran/Rees positive-half Hall surface: its
 curve-stalk factorisation cosheaf is constructed in
 Proposition~\ref{gluing:prop:24-assembly}, and its finite hCS--Hall
 comparison is Proposition~\ref{prop:k3e-finite-dwr-ran-positive-half-surface};
\item the finite compact Hall--Drinfeld source
 \(D^{\mathrm{red},\leq H}_{\hbar}(K3\times E,\sigma)\) constructed
 from reduced compact Hall windows in
 Theorem~\ref{thm:k3e-finite-compact-hall-drinfeld-double}; its
 primitive BPS Lie algebra is identified with the Borcherds root Lie
 algebra only after the primitive Hall--Borcherds recognition problem
 of Problem~\ref{op:k3e-compact-coha-hall-borcherds} is solved, using
 the Davison--Meinhardt PBW package
 (\emph{Inventiones Mathematicae} \textbf{221} (2020), 777--871);
\item the Borcherds--Gritsenko--Nikulin BKM superalgebra
 $\mathfrak{g}_{\Delta_5}$ obtained by the singular theta lift of the
 K3 weak Jacobi form $\phi_{0,1}$
 (Gritsenko--Nikulin, \emph{Internat.\ Math.\ Res.\ Notices} (1998),
 no.\ 11, 525--547; Borcherds, \emph{Invent.\ Math.} \textbf{132}
 (1998), 491--562).
\end{enumerate}
At the character level the identification is the Oberdieck--Pandharipande
equality
\[
 Z^{\mathrm{red}}_{\mathrm{DT}}(K3 \times E)
 \;=\;
 -(\Phi_{10}^{\mathrm{OP}})^{-1}
 \;=\;
 -D_5^{-2}
 \;=\;
 -4096\,\Delta_5^{-2},
\]
with \(D_5=64^{-1}\Delta_5\) and
\(\Phi_{10}^{\mathrm{OP}}=D_5^2\),
and the Borcherds weight reads
\[
 \kappa_{\mathrm{BKM}}(\Delta_5) \;=\; \tfrac{1}{2}\,c_1(0) \;=\; 5.
\]
\end{theorem}

\begin{remark}[Scope of the four-corner identification]
\label{rem:k3e-four-corner-scope}
The proved character equality is
$Z^{\mathrm{red}}_{\mathrm{DT}}
=-(\Phi_{10}^{\mathrm{OP}})^{-1}
=-D_5^{-2}
=-4096\,\Delta_5^{-2}$, and the proved
Borcherds target is the positive half
\(U(\mathfrak n_+(\mathfrak g_{\Delta_5}))\). The finite DWR/Ran/Rees
maps construct ordered \(E_1\) positive-half Hall data on bounded
heights; identification of the compact primitive CoHA, recognition of
a quasi-NCCR character as a compact CoHA character, and identification
of the finite compact double with the Borcherds current quotient are
the separate primitives of
Problem~\ref{op:k3e-compact-coha-hall-borcherds}. The finite compact
Hall windows and their radical-quotient doubles are constructed in
Theorem~\ref{thm:k3e-finite-compact-hall-drinfeld-double}; the
Borcherds realisation requires the primitive Hall--Borcherds
recognition data of Problem~\ref{op:k3e-compact-coha-hall-borcherds}.
The bracket
identification
$\mathfrak{g}_{\mathrm{BPS}}(K3 \times E) \cong \mathfrak{g}_{\Delta_5}$
as generalised Kac--Moody superalgebras is the completed-double
upgrade governed by
Theorem~\ref{thm:k3e-constructed-finite-double-recognition}: it is
the joint consequence of vanishing of the Euler--Hall radical, Serre
kernel, coproduct, primitive centre, associator, and parity-fixture
defects at every height, compatible with the Mittag--Leffler
transitions — six defect-vanishing conditions beyond the scalar
reduced DT identity.
\end{remark}

\begin{remark}[Three strata of the \texorpdfstring{$K3\times E$}{K3 x E} comparison]
\label{rem:k3e-status-map}
The scalar stratum is a theorem: Oberdieck--Pandharipande give the
reduced DT character, the Borcherds--Gritsenko--Nikulin denominator gives
\(\Delta_5\), and the weight identity gives
\(\kappa_{\mathrm{BKM}}(\Delta_5)=5\). The algebra stratum is
conditional: the compact Hall promotion, the framed target
\(A_{K3 \times E}^{(\Sigma_2,C)}\), the Hall--Drinfeld double, the
BKM--bar dictionary, and the shadow/BRST/Yangian bridges require the
compact-passage, double-data, finite-recognition, and heightwise
naturality inputs isolated below. The physical stratum consists of the
M-theory reduction and the M5-brane worldvolume model; it supplies
heuristic structure and tests, not an algebraic identification.
Accordingly, the scalar equality never constructs a Hall pairing, a
Drinfeld double, a BKM--bar morphism, or a vertex-operator Yangian
current.
\end{remark}

\begin{proposition}[Four-corner comparison maps and the Hall--Drinfeld boundary]
\label{prop:k3e-four-corner-comparison-maps}
The comparison package in Theorem~\ref{thm:k3e-four-corner-structure}
has the following sources and obstructions.
\begin{enumerate}[label=\textup{(\roman*)}]
\item The Oberdieck--Pandharipande reduced DT character and the
reciprocal Gritsenko--Nikulin denominator have the same completed
positive-cone expansion:
\[
 Z^{\mathrm{red}}_{\mathrm{DT}}(K3\times E)
 \;=\;
 -\,(\Phi_{10}^{\mathrm{OP}})^{-1}
 \;=\;
 -D_5^{-2}
 \;=\;
 -4096\,\Delta_5^{-2},
 \qquad
 D_5=64^{-1}\Delta_5,\quad \Phi_{10}^{\mathrm{OP}}=D_5^2.
\]
\item The primitive denominator $\Delta_5$ gives the target for the
radical-quotient positive half of the finite compact Hall double once
Problem~\ref{op:k3e-compact-coha-hall-borcherds} is solved:
\[
 U(\mathfrak n_+(\mathfrak g_{\Delta_5})),
 \qquad
 \operatorname{sdim}\mathfrak g_\alpha=c(4nm-\ell^2),
\]
on every finite Borcherds-cone truncation. This is the positive-half
character comparison; the full Hopf-superalgebra comparison
additionally requires the negative half, the Cartan part, the
Euler--Hall pairing, the coproduct, and the radical descent of
Theorem~\ref{thm:k3e-constructed-finite-double-recognition}.
\item Any full quantum comparison factors heightwise through the
constructed finite radical-quotient Hall--Drinfeld double
\[
 \overline{\mathcal H}^{+,\leq H}_{K3\times E,\sigma}
 \longrightarrow
 D^{\mathrm{red},\leq H}_{\hbar}(K3\times E,\sigma)
 \longrightarrow
 Y^{\mathrm{SB},\leq H}_{\hbar}(\mathfrak g_{\Delta_5}),
\]
with the second arrow an isomorphism exactly under the six
defect-vanishing conditions of
Theorem~\ref{thm:k3e-constructed-finite-double-recognition}. The
BKM-side object is therefore the Hall--Drinfeld double of the K3
positive half — a Borcherds--Kac--Moody Hopf-superalgebra structure
attached to the Manin pair
$(\mathfrak{g}_{\Delta_5}, \mathfrak{n}_+^{\mathrm{imag}})$ — distinct
in type from the Drinfeld Yangian $Y(\mathfrak{g}_{K3})$ of
Chapter~\ref{ch:k3-yangian} on the Mukai self-mirror
current-algebra branch.
\end{enumerate}
\end{proposition}

\begin{proof}
The first equality is Oberdieck--Pandharipande's reduced DT theorem,
with Maulik--Toda fixing the Behrend sign; the last equality is the
OP-normalized Gritsenko--Nikulin identity
\(D_5=64^{-1}\Delta_5\), \(\Phi_{10}^{\mathrm{OP}}=D_5^2\). The
second item is the Gritsenko--Nikulin--Borcherds denominator theorem
applied to $\phi_{0,1}$, together with the finite Borcherds core
Theorem~\ref{thm:k3e-bkm-finite-core}. The final factorisation is the
finite Drinfeld double construction of
Theorem~\ref{thm:k3e-finite-compact-hall-drinfeld-double};
recognition as the Serre--Borcherds current quotient is exactly
Theorem~\ref{thm:k3e-constructed-finite-double-recognition}, with
defect-vanishing conditions on the Euler--Hall radical, Serre kernel,
coproduct, primitive centre, and associator — the BKM Hopf-superalgebra
discipline distinct from the strict Yangian Cartan/J-presentation
datum.
\end{proof}

\begin{proposition}[Finite DWR/Ran Rees hCS--Hall positive half]
\label{prop:k3e-finite-dwr-ran-positive-half-surface}
Fix a DWR/\v{C}ech/Ran cover \(\mathfrak U\) of \(K3\times E\), a
holomorphic-jet bound \(N\), renormalisation bound \(r\), charge bound
\(L\), Ran arity \(m\), and Borcherds height \(H\). The constructed
hCS--Hall comparison at these bounds is the ordered \(E_1\) Rees-Hall
map
\[
 \Theta_{\hCS\to\Hall}^{\mathrm{Rees},\mathrm{or};N,r,L,m}\colon
 B\Obs_{\hCS}^{q,\le r,\le N}(-)_{\le L}
 \longrightarrow
 \Hall_{\mathrm{crit}}^{\mathrm{Rees},\mathrm{or},\le N,\le L}(-)
\]
on the finite ordered Ran nerve. On the \(24\) elliptic curve-stalks
of a generic elliptic K3, the constructed positive-half object is
\[
 \mathcal H_{\mathrm{stalk}}^{+,\le H}
 :=
 \operatorname{Rees}_{\le H}
 \Yplus_{\mathrm{stalk}}(K3\times E;\mathfrak U),
\]
with Picard--Lefschetz monodromy \(\prod_bT_b=1\). These finite
objects supply ordered $E_1$ positive-half Hall data on the curve
stalks. The further constructions — primitive-preserving isomorphisms
onto the compact primitive CoHA, the Serre-equivariant quasi-NCCR, and
the Drinfeld double $D_\hbar(\mathcal H_{\mathrm{stalk}}^+)$ — are the
separate primitives of
Theorem~\ref{thm:k3e-finite-compact-hall-drinfeld-double} and
Problem~\ref{op:k3e-compact-coha-hall-borcherds}.
\end{proposition}

\begin{proof}
Proposition~\ref{prop:finite-dwr-ran-rees-ordered-hall-layer} constructs
the finite natural transformation by integrating relative simplex maps
solving the bar--Rees convolution Maurer--Cartan equation. Its product
is the Rees Thom--Sebastiani product, hence ordered \(E_1\) Hall
multiplication. Proposition~\ref{gluing:prop:24-assembly} constructs
the curve-stalk positive-half factorisation cosheaf and imposes the
Picard--Lefschetz relation. Proposition~\ref{thm:global-coha-limit}
separates this finite Rees comparison from a completed critical-CoHA
realisation: completion requires compatible inverse-system elements,
and the critical-CoHA image requires a monoidal vanishing-cycle
realisation. A Drinfeld double further requires the reduced compact
window, the Serre-dual negative half, the Euler--Hall pairing, the
Mukai Cartan torus, and the radical quotient of
Theorem~\ref{thm:k3e-finite-compact-hall-drinfeld-double}. Recognition
of that double as a Borcherds current quotient requires the six
finite defects of
Theorem~\ref{thm:k3e-constructed-finite-double-recognition} to vanish.
None of those structures is part of the finite positive-half
construction.
\end{proof}

\begin{remark}[The four-\texorpdfstring{$\kappa_\bullet$}{kappa-.} spectrum of \texorpdfstring{$K3 \times E$}{K3 x E}]
\label{rem:k3e-four-kappa-spectrum}
The four-corner theorem realises four numerical invariants from four
mathematically distinct constructions:
\begin{itemize}
\item $\kappa_{\mathrm{cat}}(K3 \times E) = \chi(\cO_{K3 \times E})
 = \chi(\cO_{K3}) \cdot \chi(\cO_E) = 2 \cdot 0 = 0$, K\"unneth on
 the total space;
\item $\kappa_{\mathrm{ch}}^{\mathrm{Hodge}}(K3 \times E) = \sum_q (-1)^q
 h^{0, q}(K3 \times E) = 0$, the compact chiral Hodge supertrace;
\item $\kappa_{\mathrm{ch}}^{\mathrm{Heis}}(K3 \times E) = 3$, the
 Heisenberg Stage-$2$ specialisation;
\item $\kappa_{\mathrm{BKM}}(\Delta_5) = \left.c_N(0)/2\right|_{N=1} = 5$, the Borcherds
 weight of the Gritsenko--Nikulin lift of $\phi_{0,1}$;
\item $\kappa_{\mathrm{fiber}}(K3) = \chi_{\mathrm{top}}(K3) = 24$,
 the Kodaira nodal count on the elliptic K3 fibre.
\end{itemize}
The numerical coincidence
\(\kappa_{\mathrm{ch}}^{\mathrm{Hodge}}(K3 \times E) = \kappa_{\mathrm{cat}}(K3 \times E) = 0\)
on the compact CY$_3$ reflects the identity
\(\sum_q (-1)^q h^{0, q}(X) = \chi(\cO_X)\); the two invariants share a
value here while remaining distinct constructions (categorical Euler
versus chiral Hodge supertrace, with separate behaviour off the compact
locus and under chart specialisation). The fibre witness
\(\kappa_{\mathrm{cat}}(K3) = 2\) is not an entry in the total-space
spectrum, since \(\kappa_{\mathrm{cat}}\) is K\"unneth-multiplicative,
not additive, and the total-space value is \(2 \cdot 0 = 0\). The
canonical five-construction spectrum on \(K3 \times E\) is therefore
\[
\{\kappa_{\mathrm{cat}}, \kappa_{\mathrm{ch}}^{\mathrm{Hodge}},
 \kappa_{\mathrm{ch}}^{\mathrm{Heis}}, \kappa_{\mathrm{BKM}},
 \kappa_{\mathrm{fiber}}\}(K3 \times E) = \{0, 0, 3, 5, 24\},
\]
with each entry produced by a distinct construction on the same CY$_3$
datum. The attempted reduction
\[
\kappa_{\mathrm{BKM}} \stackrel{?}{=} \kappa_{\mathrm{ch}}^{\mathrm{Hodge}} +
\chi(\cO_{\mathrm{fiber}})
\]
fails: an elliptic Stage-$2$ fibre gives \(0+0\), while a K3 fibre gives
\(0+2\), and neither is \(5\). This is why the subscripted spectrum is
not optional.
\end{remark}

\begin{proposition}[\texorpdfstring{$\Delta_5$}{Delta-5}, \texorpdfstring{$\Phi_{10}^{\mathrm{un}}$}{Phi-10-un}, and \texorpdfstring{$\Phi_{12}$}{Phi-12}]
\label{prop:k3e-delta5-phi10-phi12-separation}
The automorphic inputs used in this chapter belong to distinct
algebraic objects and normalisation conventions.
\[
\begin{array}{c|c|c|c}
\text{form} & \text{weight} & \text{source} & \text{algebraic object} \\
\hline
\Delta_5 & 5 &
\mathrm{Borch}(\phi_{0,1}) &
\mathfrak g_{\Delta_5}\text{ on }\Lambda^{2,1}_{\mathrm{II}} \\
\Phi_{10}^{\mathrm{un}} & 10 &
\mathrm{Igusa}=\Delta_5^2 &
\text{DVV / genus-2 BPS square} \\
\Phi_{10}^{\mathrm{OP}} & 10 &
D_5^2,\quad D_5=64^{-1}\Delta_5 &
Z^{\mathrm{red}}_{\mathrm{DT}}(K3\times E)=-(\Phi_{10}^{\mathrm{OP}})^{-1} \\
\Phi_{12} & 12 &
\mathrm{Borcherds}(\mathrm{II}_{25,1}) &
\mathfrak g_{\Phi_{12}}\text{, the Fake-Monster BKM}
\end{array}
\]
Thus \(\Delta_5\) is the primitive denominator of
\(\mathfrak g_{\Delta_5}\), \(\Phi_{10}^{\mathrm{un}}=\Delta_5^2\) is
the unnormalised Igusa square, \(\Phi_{10}^{\mathrm{OP}}=D_5^2\) is the
Oberdieck--Pandharipande scalar, and \(\Phi_{12}\) is the Fake-Monster
denominator on the rank-\(26\) Lorentzian lattice. No Hall comparison,
quasi-NCCR character, or \(\Phi_3\)-specialisation turns \(\Phi_{12}\)
into a \(K3\times E\) denominator.
\end{proposition}

\begin{proof}
Gritsenko--Nikulin identify the Borcherds lift of the K3 weak Jacobi
form \(\phi_{0,1}\) with the weight-\(5\) denominator \(\Delta_5\).
Igusa's genus-\(2\) cusp form of weight \(10\) satisfies
\(\Phi_{10}^{\mathrm{un}}=\Delta_5^2\). The OP-normalized monic product
\(D_5=64^{-1}\Delta_5\) gives
\(\Phi_{10}^{\mathrm{OP}}=D_5^2=4096^{-1}\Delta_5^2\), and
Oberdieck--Pandharipande compute the reduced DT scalar as
\(-(\Phi_{10}^{\mathrm{OP}})^{-1}=-4096\,\Delta_5^{-2}\). Borcherds'
Fake-Monster construction instead uses
the \(\mathrm{II}_{25,1}\) lattice and has denominator of weight \(12\);
its root multiplicities are the Leech-lattice constants, not the
signed \(\phi_{0,1}\)-character coefficients of \(\mathfrak g_{\Delta_5}\). The weights,
source lattices, and root-character functions are different, so the
three forms cannot be interchanged.
\end{proof}

The threefold $K3 \times E$ is a fibration of a CY$_2$ over a CY$_1$. Does its chiral algebra decompose accordingly? A naive Fubini argument would predict $A_{K3 \times E} \simeq A_{K3} \otimes A_E$, and the modular characteristic would split additively as $\kappa_{\mathrm{ch}}^{\mathrm{Heis}}(K3 \times E) = \kappa_{\mathrm{ch}}^{\mathrm{Heis}}(K3) + \kappa_{\mathrm{ch}}^{\mathrm{Heis}}(E) = 2 + 1 = 3$. Oberdieck--Pandharipande compute the reduced DT scalar
\[
 Z^{\mathrm{red}}_{\mathrm{DT}}(K3\times E)
 =
 -(\Phi_{10}^{\mathrm{OP}})^{-1}
 =
 -D_5^{-2}
 =
 -4096\,\Delta_5^{-2},
\]
where \(D_5=64^{-1}\Delta_5\) is the monic Borcherds product and
\(\Phi_{10}^{\mathrm{OP}}=D_5^2\) is the OP scalar convention. The form
\(\Delta_5\) itself is the Gritsenko--Nikulin automorphic form of
weight~$5$ on $\mathrm{O}^+(3,2)$. The Oberdieck--Pixton name belongs
to the broader reduced GW/DT/PT theory; the scalar identity used
here is the Oberdieck--Pandharipande formula. The weight $5$ does not
match the additive Heisenberg sum $3$: $5 \neq 2 + 1$.

Two different modular characteristics are in play; conflating them produces the apparent contradiction, and the subscripted $\kappa_\bullet$-spectrum is the discipline that keeps them separate. The compact chiral Hodge supertrace gives $\kappa_{\mathrm{ch}}^{\mathrm{Hodge}}(K3 \times E) = \sum_q (-1)^q h^{0, q}(K3 \times E) = 0$, while the chiral de Rham / Heisenberg specialisation gives $\kappa_{\mathrm{ch}}^{\mathrm{Heis}}(K3 \times E) = 3 = \dim_\C$. The Borcherds lift weight gives $\kappa_{\mathrm{BKM}}(\Delta_5) = 5 = \mathrm{wt}(\Delta_5)$, which is not a modular characteristic of any constructed chiral algebra: it is a weight attached to a generalized Borcherds--Kac--Moody superalgebra $\mathfrak{g}_{\Delta_5}$ through its denominator identity. The framed algebra $A_{K3 \times E}^{(\Sigma_2,C)}$ is the target object of the H1--H4 specialisation of Theorem~CY-A$_3$ (Theorem~\ref{thm:cy-to-chiral-d3}); its Heisenberg specialisation has $\kappa_{\mathrm{ch}}^{\mathrm{Heis}} = 3$, not~$5$.

The two-stage factorisation of Definition~\ref{def:phi-fa-and-sp} and Theorem~\ref{thm:phi-two-stage-factorisation} places the $d = 3$ story here: $\PhiFA_3(D^b(\Coh(K3 \times E)))$ is the canonical $\EdHolFA(K3 \times E)$ (holomorphic Chern--Simons on this CY$_3$) under the stated hypotheses, and $\SpCh_{K3, E}$ evaluates it along the $K3$-fibration $\Sigma_2 = K3$ to produce an $\Eone$-chiral shadow on~$E$. The comparison of that shadow with the Borcherds--$\Delta_5$ denominator object is a CY-C-level comparison, with the specialisation rescaling $\chi(K3) = 24$ supplying the Gritsenko--Nikulin paramodular normalisation on the automorphic side. Proposition~\ref{prop:native-en-level} fixes $n_{\mathrm{native}}(3) = 1$, so the output on~$E$ is $\Eone$-chiral; the $E_2$-structure of Remark~\ref{rem:e2-on-drinfeld-centre} lives on the Drinfeld centre $Z(\Rep(A_{K3\times E}))$, not on $A_{K3 \times E}$ itself. The $K3$-fibration $\Sigma_2 = K3$ is one choice of specialisation datum; Remark~\ref{rem:multiple-e1-specialisations} records that different $\Sigma_2 \subset K3 \times E$ produce different paramodular shadows from the \emph{same} canonical $\PhiFA_3$ output.

This chapter treats $K3 \times E$ as the prototype for the $d = 3$ correspondence. The target BKM superalgebra is $\mathfrak{g}_{\Delta_5}$ attached to the K3 Jacobi input, together with the Oberdieck--Pandharipande reduced DT scalar
\[
 -(\Phi_{10}^{\mathrm{OP}})^{-1}
 =
 -D_5^{-2}
 =
 -4096\,\Delta_5^{-2},
 \qquad
 D_5=64^{-1}\Delta_5,\quad \Phi_{10}^{\mathrm{OP}}=D_5^2.
\]
The primitive BKM denominator is $\Delta_5$ itself; the OP compact scalar is the square of the monic Borcherds product \(D_5\). The Vol~I bar-cobar apparatus reaches this chapter through explicit identities among signed root characters, genus-$g$ partition functions, and lattice theta series. The compact Hall source and its finite radical-quotient doubles are constructed in Theorems~\ref{thm:k3e-reduced-compact-source-finite-hall-windows} and~\ref{thm:k3e-finite-compact-hall-drinfeld-double}; recognition as the \(\Delta_5\) current object requires the defect-vanishing data of Theorem~\ref{thm:k3e-constructed-finite-double-recognition}. The chapter concludes with the K3 double current algebra $\fg_{K3}$ (Definition~\ref{def:k3-double-current-algebra}), the K3 analogue of the double current algebra $\fg \otimes \bC[u,v]$ in which the polynomial ring is replaced by $H^*(S,\bC)$ and the polynomial residue pairing by the Mukai pairing; the resulting finite-dimensional Lie algebra is the classical limit of the K3 chiral Hall--Drinfeld double whose quantisation is governed by the Siegel-elliptic dynamical $R_{\mathrm{Sieg,dyn}}$-matrix (Chapter~\ref{ch:k3-chiral-bialgebra-platonic}, Theorem~\ref{thm:R-sieg-fourth-class}), with the Maulik--Okounkov evaluation-$R$-matrix (Theorem~\ref{thm:k3e-mo-rmatrix}) recovering the ADE/Kummer-chart restriction of this fuller object.

\begin{proposition}[Kodaira \texorpdfstring{$I_1$}{I-1} residues behind the 1-loop origin of \texorpdfstring{$\Delta_5$}{Delta-5}]
\label{prop:k3e-one-loop-origin-chapter-header}
The Gritsenko--Nikulin Igusa form $\Delta_5$ enters this chapter in two logically
distinct roles, linked by the Kodaira residue package of a generic elliptic K3
$\pi\colon K3\to \mathbb{P}^1$:
\begin{enumerate}[label=\textup{(\roman*)}]
\item \emph{Denominator input.} One fixes $\Delta_5$ as the Borcherds lift of the K3
weak Jacobi form $\phi_{0,1}$ and defines $\mathfrak{g}_{\Delta_5}$ via the
Gritsenko--Nikulin denominator. This is the denominator-first viewpoint of
Section~\ref{sec:k3e-bkm} and Theorem~\ref{thm:k3e-denominator}.
\item \emph{1-loop output.} The same $\Delta_5$ is the unique output of the
paramodular anomaly-cancellation problem for twisted 11D supergravity on
$K3\times T^2$ Omega-deformed with $24$ M5-branes on the singular fibres. The local
input is Kodaira's formula
\[
 \chi_{\mathrm{top}}(K3)
 \;=\;
 24
 \;=\;
 \sum_{i=1}^{24}\chi_{\mathrm{top}}(F_i),
\]
with each generic singular fibre $F_i$ of type $I_1$, hence each contributing one
local logarithmic residue.
\item \emph{Borcherds-weight normalisation.} These $24$ local $I_1$ residues are
the nodal input to the one-loop anomaly calculation, but they do not decompose
the Borcherds weight. The BKM invariant is fixed by the Jacobi constant term:
\[
 \kappa_{\mathrm{BKM}}(\Delta_5)
 \;=\;
 \mathrm{wt}(\Delta_5)
 \;=\;
 \left.c_N(0)/2\right|_{N=1}
 \;=\;
 10/2
 \;=\;
 5
\]
as in Theorem~\ref{thm:k3-1loop-output} of
Chapter~\ref{ch:k3-chiral-bialgebra-platonic}. The fibre value
$\chi(\mathcal{O}_{K3})=2$ and the Fricke-split residue terms are
separate anomaly data, not a formula for $\kappa_{\mathrm{BKM}}$.
\end{enumerate}
Thus this chapter uses $\Delta_5$ as an automorphic input while Chapter~\ref{ch:k3-chiral-bialgebra-platonic}
proves that the same form is forced as the 1-loop output of the twisted-supergravity
parent.
\end{proposition}

\begin{proof}
Item~\textup{(i)} is the setup of the BKM chapter itself. For item~\textup{(ii)},
Kodaira's classification gives $24$ nodal fibres for a generic elliptic K3, and each
$I_1$ fibre contributes Euler number $1$, so the total local contribution is exactly
the topological Euler number $24$. The local logarithmic-residue interpretation is the
Kodaira-side input used in Chapter~\ref{ch:k3-chiral-bialgebra-platonic},
Theorem~\ref{thm:k3-1loop-output}. Item~\textup{(iii)} is then just the same theorem
restated from the present denominator-first viewpoint: the bulk term contributes
$\chi(\mathcal{O}_{K3})=2$, the aggregated nodal sector contributes the effective
residue term $3$, and the resulting weight-$5$ section is $\Delta_5$.
\end{proof}

\begin{remark}[Imaginary-root signed exponents seed Humbert monodromy]
\label{rem:k3e-imaginary-root-seeds-humbert}
Each positive Fourier coefficient $c_{\Delta_5}(N,\ell) > 0$ of the Gritsenko--Nikulin
denominator seeds an imaginary simple root of $\mathfrak{g}_{\Delta_5}$ at norm
$4N - \ell^2 \leq 0$ with signed exponent $c_{\Delta_5}(N,\ell)$
(Theorem~\ref{thm:k3e-denominator}; Section~\ref{sec:k3e-phi01}). The resulting
infinite ladder of imaginary roots is the combinatorial shadow of the Humbert
monodromy of Chapter~\ref{ch:k3-chiral-bialgebra-platonic},
Theorem~\ref{thm:humbert-order-K-kappa}:
\begin{itemize}
\item Each discriminant class $D = 4N-\ell^2\in\{-1, -4, \ldots\}$ on the Jacobi-form
 side corresponds to a Heegner divisor $\mathcal{H}_D\subset\overline{\mathcal{A}_2}$
 on the Siegel parameter side.
\item The Humbert divisors $H_1 = \mathcal{H}_{-1}$ and $H_4 = \mathcal{H}_{-4}$ are
 the two lowest-discriminant Heegner divisors; the monodromies of
 $\mathcal{L}^{\Delta_5}$ around them have orders $8$ and $16$ respectively
 (Chapter~\ref{ch:k3-chiral-bialgebra-platonic},
 Theorem~\ref{thm:humbert-order-K-kappa}).
\item The Bruinier Heegner-Chern-class reciprocity (Bruinier 2002, Proposition 5.1)
relates the Chern class of $\mathcal{L}^{\Delta_5}$ on $\mathcal{H}_D$ to the
 $c_{\Delta_5}(N,\ell)$ with $4N-\ell^2 = D$. It yields the
 automorphic Chern-class datum, not a canonical equation for \(\hbar^2\).
\end{itemize}
Thus the imaginary-root signed exponent $c_{\Delta_5}(N,\ell)$ at each Fourier mode
\emph{seeds} the Humbert-divisor monodromy at the corresponding Heegner
discriminant: the Jacobi-form coefficient on the Gritsenko side and the
$\mathcal{D}$-module monodromy order on the Siegel side are the same data viewed
twice, linked by the Borcherds singular-theta correspondence.
\end{remark}

\begin{remark}[Gritsenko--Nikulin denominator identity as the Weyl--Kac--Borcherds character of \texorpdfstring{$\mathbf{H}_{\Delta_5}$}{H-Delta-5}]
\label{rem:k3e-gritsenko-nikulin-denominator-WKB-character}
The Gritsenko--Nikulin denominator identity
\[
 \tfrac{1}{64}\Delta_5(2Z)
 \;=\;
 \Phi(z)
 \;=\;
 e^{-2\pi i(\rho, z)}\prod_{\alpha\in\Delta_+}\bigl(1 - e^{-2\pi i(\alpha, z)}\bigr)^{\operatorname{sdim}\mathfrak g_\alpha}
\]
(Theorem~\ref{thm:k3e-denominator}; Lorgat 2020 Theorem $3$) is the
\emph{Weyl--Kac--Borcherds character} of the chiral bialgebra $\mathbf{H}_{\Delta_5}$ of
Chapter~\ref{ch:k3-chiral-bialgebra-platonic} applied to the trivial one-dimensional
representation. The dictionary, traced through Chapter~\ref{ch:k3-chiral-bialgebra-platonic}
Theorem~\ref{thm:k3-siegel-character} (genus-$2$ Siegel character
\((\Phi_{10}^{\mathrm{un}}(Z))^{-1}\)) and the denominator identity
\(\Phi=\Delta_5\):
\begin{itemize}
\item the \emph{sum side} of Borcherds's denominator
 $\sum_{w\in W(\Lambda^{2,1}_{II})}\det(w)\bigl(\exp(-2\pi i w(\rho, z))
 - \sum_{a\in\Lambda^{2,1}_{II}\cap\mathbb{R}_{>0}\mathcal{P}_{II}}m(a)\exp(-2\pi i w(\rho+a, z))\bigr)$
 is the Gritsenko additive lift
 $\mathrm{Grit}(\eta^9\vartheta_1) = \Delta_5/64$;
\item the \emph{product side}
 $\exp(-2\pi i(\rho, z))\prod_{\alpha\in\Delta_+}(1-\exp(-2\pi i(\alpha, z)))^{\operatorname{sdim}\mathfrak g_\alpha}$
 is the Borcherds multiplicative lift
 $\mathrm{Borch}(\phi_{0,1}) = \Delta_5/64$;
\item the equality sum $=$ product is the denominator identity, which is the
 genus-$2$ automorphic-form version of the Shimura--Waldspurger correspondence
 (Chapter~\ref{ch:k3-chiral-bialgebra-platonic},
 Remark~\ref{rem:ikeda-gritsenko-disambiguation});
\item the \emph{character of} $\mathbf{H}_{\Delta_5}$ on $\mathbb{H}_2$ is the
 reciprocal of the unnormalised Igusa square
 \(\Phi_{10}^{\mathrm{un}}=\Delta_5^2\), meromorphic Siegel modular weight $-10$,
 factorising at the separating-node degeneration as
 $(1/\eta^{24})\cdot(1/\phi_{10,1})$ recovering the K3 chiral algebra denominator
 times the Jacobi form of the elliptic direction.
\end{itemize}
The Lorgat 2020 theta-product representation $\Delta_5 = \prod_{(a,b)} v_{a,b}$ over
the $10$ even theta constants is the explicit multiplicative realisation of
$\mathrm{Borch}(\phi_{0,1})$; the factor of $64 = 2^6$ is the
$f(1,1,1) = 64$ leading Fourier coefficient of $\Delta_5$ (Lorgat 2020 Theorem $3$),
equivalent to $2\cdot\lambda = 2\cdot 32 = 64$ where $\lambda = |\det C_{\mathrm{BKM}}|$
is the absolute discriminant of the BKM Cartan matrix.
\end{remark}

\begin{remark}[Gritsenko additive lift, not Ikeda]
\label{rem:k3e-gritsenko-not-ikeda}
The automorphic production of $\Delta_5$ is \emph{Gritsenko's 1999 additive lift} from
Jacobi forms of weight $9/2$ and index $1/2$ to paramodular forms:
$\Delta_5 = \mathrm{Grit}(\eta^9\vartheta_1)$. Equivalently, via Borcherds 1998,
$\Delta_5$ is the singular-theta multiplicative lift of $\phi_{0,1}$. The two
constructions are related by the Shimura--Waldspurger correspondence. Ikeda's lifts
(Ikeda 2001 Saito--Kurokawa, Ikeda 2006 generalised Ikeda lift) produce
$\Delta_{10} = \Delta_5^2$ from weight-$16$ elliptic sources on $\mathrm{Sp}_4(\mathbb{Z})$,
which is a \emph{distinct} lift and should not be conflated with the Gritsenko lift
that produces $\Delta_5$ itself. The Ikeda-vs-Gritsenko conflation is a recurring
confusion: Ikeda always produces $\Delta_{10}$ from an integer-weight source,
never $\Delta_5$ directly, because the Kohnen plus-space constraint forces
integer-weight Siegel output.
\end{remark}

\begin{remark}[Gritsenko additive lift as the genus-\texorpdfstring{$2$}{2} shadow of Shimura's theorem]
\label{rem:k3e-gritsenko-shimura-shadow}
The Gritsenko 1999 additive lift
\[
 \mathrm{Grit}\colon
 J^{\mathrm{cusp}}_{k,m}\bigl(\widetilde{\mathrm{SL}_2}(\mathbb{Z}),\; v_\eta^{24}\bigr)
 \;\longrightarrow\;
 S_{k}\bigl(\widehat{\mathrm{Sp}_4(\mathbb{Z})}^{v_{\Delta_5}},\; v_{\Delta_5}\bigr)
\]
takes a Jacobi cusp form of weight $k$ and index $m$ (integer or half-integer) on the
metaplectic double cover $\widetilde{\mathrm{SL}_2}(\mathbb{Z})$ to a Siegel modular
form of weight $k$ on the Maass $\mathbb{Z}/2$-spin cover of $\mathrm{Sp}_4(\mathbb{Z})$
with the multiplier $v_{\Delta_5}$. At $(k,m) = (9/2, 1/2)$ with input
$\eta^9\vartheta_1$, the output is $\Delta_5$ (Gritsenko 1999; Gritsenko--Nikulin 1998).
This is the \emph{genus-$2$ shadow of Shimura's correspondence theorem} (Shimura 1973):
Shimura maps half-integral-weight forms on $\widetilde{\mathrm{SL}_2}(\mathbb{Z})$ to
integral-weight forms on $\mathrm{SL}_2(\mathbb{Z})$ via a quadratic-character twist; the
Gritsenko map is the $\mathrm{Sp}_4$-analogue, going from half-integral Jacobi index
on the metaplectic $\widetilde{\mathrm{SL}_2}$-base to paramodular Siegel modular forms
on the spin cover of $\mathrm{Sp}_4(\mathbb{Z})$. Waldspurger's refinement of Shimura
(Waldspurger 1981) describes the genus-$1$ case in terms of $L$-values at half-integer
points of twisted automorphic $L$-functions; the genus-$2$ paramodular analogue due to
Gritsenko replaces the integral with a theta-sum (Gritsenko 1999, $\S 2$). The
equivalence $\mathrm{Grit}(\eta^9\vartheta_1) = \Delta_5 = \mathrm{Borch}(\phi_{0,1})$
between the additive lift and the Borcherds singular-theta lift is then the
\emph{Shimura--Waldspurger correspondence realised at genus $2$}: the cuspidal
half-integral Jacobi source $\eta^9\vartheta_1$ and the weakly holomorphic weight-$0$
integer-weight source $\phi_{0,1}$ are the two sides of a Shimura--Waldspurger pair
whose common image is $\Delta_5$. See Lorgat 2020, Theorems~$2$--$4$ and the preamble
on $\Delta_5$ for the explicit $\eta^9\vartheta_1$-to-\(\Phi_{10}^{\mathrm{un}}\)-Fourier-coefficient
identity, and Eichler--Zagier 1985 for the Jacobi-form side.
\end{remark}

\begin{remark}[Gritsenko additive lift seeds the Weyl--Kac--Borcherds denominator]
\label{rem:k3e-gritsenko-seeds-WKB}
The Gritsenko additive lift is not merely an alternative construction of $\Delta_5$:
it is the \emph{automorphic face} of the Weyl--Kac--Borcherds denominator identity
for $\mathfrak{g}_{\Delta_5}$. The dictionary is as follows.
\begin{itemize}
\item The Fourier--Jacobi expansion of $\Delta_5$,
 \[
 \Delta_5(Z)
 \;=\;
 \sum_{m\equiv 1\,(\bmod\, 2)}\, \phi_{5,\,m/2}(\tau,z)\, e^{\pi i m\rho}
 \]
 (odd $m$ indices; Lorgat 2020 Theorem $3$), is precisely the Fock-level expansion of
 the one-dimensional trivial representation of $\mathfrak{g}_{\Delta_5}$ graded by the
 null direction $\rho$ of the rank-$2$ hyperbolic $\Lambda^{2,1}\subset\Lambda^{3,2}$.
\item The first Fourier--Jacobi coefficient $\phi_{5,1/2} = \eta^9\vartheta_1$ is
 the input to the Gritsenko lift; its square $\phi_{10,1} = \eta^{18}\vartheta_1^2$
 is the first FJ-coefficient of the unnormalised Igusa square
 \(\Phi_{10}^{\mathrm{un}}=\Delta_5^2\) (Lorgat 2020 preamble;
 Eichler--Zagier 1985 \S 5).
\item The Maass relations (Eichler--Zagier 1985 \S $6$) determine all subsequent
 Fourier--Jacobi coefficients $\phi_{5,m/2}$, $m\geq 3$, from $\phi_{5,1/2}$ via
 Hecke operators $T_-(m)$ on the Jacobi form. These Hecke relations are the
 \emph{automorphic incarnation of the Weyl--Kac character-sum recursion} for
 $\mathfrak{g}_{\Delta_5}$.
	\item Every Borcherds product factor $(1 - q^n y^\ell p^m)^{f(4nm-\ell^2)}$ of the
	 multiplicative lift $\Delta_5 = \mathrm{Borch}(\phi_{0,1})$
	 (Theorem~\ref{thm:k3e-borcherds-product}) matches the signed
	 root-character contribution
	 \(\operatorname{sdim}\mathfrak g_{\Delta_5,\alpha}=f(4nm-\ell^2)\)
	 on the BKM side (Theorem~\ref{thm:k3e-denominator}). An ordinary
	 root-space dimension is read only after the target parity fixture.
\end{itemize}
Thus Gritsenko additive $=$ FJ coefficient recursion $=$ Weyl--Kac character formula,
and Borcherds multiplicative $=$ product expansion $=$ denominator identity
root-space product. The Shimura--Waldspurger correspondence between the two lifts is
the automorphic witness that the \emph{sum side} and the \emph{product side} of the
BKM denominator identity coincide on the nose. This is the genus-$2$ analogue of
Koike--Norton--Zagier for the Monster on $\mathrm{II}_{1,1}$ and of Borcherds's Fake
Monster denominator on $\mathrm{II}_{25,1}$, specialised to the K3 elliptic-genus input
$\phi_{0,1}$ with its $24$-node $\widetilde{M}_{24}$-twinings.
\end{remark}

\begin{remark}[Half-integral weight and the Maass \texorpdfstring{$\mathbb{Z}/2$}{Z/2}-spin cover]
\label{rem:k3e-maass-spin-cover}
The Igusa form $\Delta_5$ has weight $5$, which is odd. Hence $\Delta_5$ is
\emph{not} paramodular of level $1$ in the integral sense; it lives on the Maass
$\mathbb{Z}/2$-spin cover $\widehat{\mathrm{Sp}_4(\mathbb{Z})}^{v_{\Delta_5}}$ with
the multiplier $v_{\Delta_5}$ given by the order-$2$ character of the metaplectic
double cover associated to the half-integral Jacobi index $1/2$ of the Gritsenko
additive lift source $\eta^9\vartheta_1$ of weight $9/2$ and index $1/2$. The
archimedean Schmidt parameters of the associated automorphic representation are
$(7/2,5/2)$ for the holomorphic discrete series supporting $\Delta_5$; the
non-tempered Saito--Kurokawa CAP at $(17/2,1/2)$ supports $\Delta_{10}$. These
should not be conflated.
\end{remark}

\begin{remark}[Maass multiplier \texorpdfstring{$v_{\Delta_5}$}{v-Delta-5}: explicit cocycle and Jacobi half-integer origin]
\label{rem:k3e-maass-multiplier-explicit}
Maass 1964 (\emph{Die Multiplikatorsysteme zur Siegelschen Modulgruppe}) gives the
explicit multiplier system for $\Delta_5$ on generators of $\mathrm{Sp}_4(\mathbb{Z})$.
For $A\in\mathrm{GL}_2(\mathbb{Z})$ and $B\in M_{2\times 2}(\mathbb{Z})$:
\begin{align*}
 v_{\Delta_5}\bigl({\scriptsize\begin{pmatrix} 0 & I_2 \\ -I_2 & 0\end{pmatrix}}\bigr)
 &= 1, \\
 v_{\Delta_5}\bigl({\scriptsize\begin{pmatrix} I_2 & B \\ 0 & I_2\end{pmatrix}}\bigr)
 &= (-1)^{b_1+b_2+b_3}, \\
 v_{\Delta_5}\bigl({\scriptsize\begin{pmatrix} {}^tA^{-1} & 0 \\ 0 & A\end{pmatrix}}\bigr)
 &= (-1)^{(1+a_1+a_4)(1+a_2+a_3)+a_1 a_4}
\end{align*}
(Lorgat 2020, formulas between Lemma 1 and Theorem 3). The factor $(-1)^{b_1+b_2+b_3}$
on translations witnesses that $\Delta_5$ picks up a sign $-1$ under each of the three
independent half-period translations of $\mathbb{H}_2$; this is the origin of the
$\mathbb{Z}/2$-central extension. Geometrically the multiplier encodes the monodromy
of the square root of the Hodge bundle $\pi_*\omega_{\mathcal{A}_2/\mathbb{H}_2}$ on
the Siegel modular threefold $\mathcal{A}_2 = \mathrm{Sp}_4(\mathbb{Z})\backslash\mathbb{H}_2$.
The spin cover $\widehat{\mathrm{Sp}_4(\mathbb{Z})}^{v_{\Delta_5}}$ on which
$\Delta_5$ is honest is precisely the metaplectic $\mathrm{Sp}_4$-cover where
half-integral-weight Siegel modular forms live; this matches the half-integer Jacobi
index $1/2$ of the Gritsenko source $\eta^9\vartheta_1$ because
$J^{\mathrm{cusp}}_{k,m}$ forms with half-integer $m$ transform under the metaplectic
cover $\widetilde{\mathrm{SL}_2}$ on the base and lift via Gritsenko to the spin cover
on $\mathrm{Sp}_4$. The $\mathbb{Z}/2$-central extension is the paramodular face of the
Shimura double cover that makes half-integral-weight theory well-defined.
\end{remark}

\begin{remark}[Arthur packet \texorpdfstring{$\psi_{\Delta_{10}}$}{psi-Delta-10} and its spin refinement]
\label{rem:k3e-arthur-packet-delta10}
The Siegel cusp form $\Delta_{10} = \Delta_5^2 \in S_{10}(\mathrm{Sp}_4(\mathbb{Z}))$ is
the classical Saito--Kurokawa lift of the unique normalised weight-$18$ Hecke cusp form
$\Delta E_6 = \Delta\cdot E_6 \in S_{18}(\mathrm{SL}_2(\mathbb{Z}))$ (dim $S_{18} = 1$).
Its Arthur parameter is
\[
 \psi_{\Delta_{10}}
 \;=\;
 \phi_{\Delta E_6} \boxtimes \mathrm{Sym}^1
 \colon\;
 L_F\times\mathrm{SL}_2(\mathbb{C})
 \;\longrightarrow\;
 \mathrm{SO}_5(\mathbb{C})
\]
with $\phi_{\Delta E_6}\colon L_F\to\mathrm{GL}_2(\mathbb{C})$ the Langlands parameter
of $\Delta E_6$ and $\mathrm{Sym}^1$ the $2$-dimensional irreducible of
$\mathrm{SL}_2(\mathbb{C})$. This is a \emph{non-tempered} Arthur parameter
(non-trivial $\mathrm{SL}_2$ factor), which accounts for the Saito--Kurokawa
CAP-representation type of $\Delta_{10}$ at archimedean Schmidt parameters
$(17/2, 1/2)$. The degree-$4$ spinor $L$-function decomposes classically by
Andrianov--Evdokimov as
\[
 L(s,\Delta_{10},\mathrm{Spin})
 \;=\;
 L(s,\Delta E_6)\cdot \zeta(s-9)\cdot \zeta(s-8),
\]
matching the $\mathrm{deg}\,2 \times \mathrm{deg}\,1 \times \mathrm{deg}\,1$ factorisation
of the non-tempered Arthur shape. The automorphic representation of $\Delta_5$ is the
\emph{spin refinement} of $\psi_{\Delta_{10}}$ on the Maass $\mathbb{Z}/2$-spin cover:
local Satake parameters become square roots $\{\pm\alpha_p^{1/2},\pm\alpha_p^{-1/2}\}$
of the $\Delta_{10}$ Satake data, with branch selected by
$\epsilon_p = v_{\Delta_5}(\mathrm{Frob}_p)\in\{\pm 1\}$. The Langlands dual group of the
spin cover is $\mathrm{Sp}_4$ itself (the double cover of
$\widehat{\mathrm{Sp}_4} = \mathrm{SO}_5 = \mathrm{PGSp}_4$ equals
$\mathrm{Spin}_5 = \mathrm{Sp}_4$), exhibiting the K3 chiral bialgebra's Arthur-packet
self-duality. For $p\in\{7,11,23\}$ the local Hecke category carries an additional
$M_{24}$-orbifold twist $\chi^{M_{24}}_p$ tracking the CDH umbral mock-modular shadows
of the twinings $\phi^g_{0,1}$ at the anomalous $M_{24}$-classes
$\{7A, 7B, 11A, 23A, 23B\}$; see Chapter~\ref{ch:k3-chiral-bialgebra-platonic},
Theorem~\ref{thm:arthur-packet-H-delta5}, and Eguchi--Ooguri--Tachikawa 2011,
Cheng--Duncan--Harvey 2014 for the $M_{24}$ mock-modular structure.
\end{remark}

\begin{proposition}[Explicit Hecke eigenvalues of \texorpdfstring{$\Delta_{10}$}{Delta-10}]
\label{prop:k3e-arthur-hecke-eigenvalues}
The Andrianov spinor $L$-function factorisation $L(s,\Delta_{10}) = L(s,\Delta_{E_6})\,\zeta(s-9)\,\zeta(s-8)$ (Andrianov $1974$) gives $\lambda_p(\Delta_{10}) = a_p(\Delta_{E_6}) + p^8 + p^9$. For $\Delta_{E_6}\in S_{10}(\mathrm{SL}_2(\mathbb{Z}))$ with Fourier expansion $q\prod_{n\ge 1}(1-q^n)^{24}$, the first nine Hecke eigenvalues:
\begin{center}
\renewcommand{\arraystretch}{1.25}
\begin{tabular}{|c|r|r|r|}
\hline
$p$ & $a_p(\Delta_{E_6})$ & $p^8 + p^9$ & $\lambda_p(\Delta_{10})$ \\
\hline
$2$ & $-24$ & $768$ & $744$ \\
\hline
$3$ & $252$ & $26\,244$ & $26\,496$ \\
\hline
$5$ & $4\,830$ & $2\,343\,750$ & $2\,348\,580$ \\
\hline
$7$ & $-16\,744$ & $46\,118\,408$ & $46\,101\,664$ \\
\hline
$11$ & $534\,612$ & $2\,572\,306\,572$ & $2\,572\,841\,184$ \\
\hline
$13$ & $-577\,738$ & $11\,420\,230\,094$ & $11\,419\,652\,356$ \\
\hline
$17$ & $-6\,905\,934$ & $125\,563\,633\,938$ & $125\,556\,728\,004$ \\
\hline
$19$ & $10\,661\,420$ & $339\,671\,260\,820$ & $339\,681\,922\,240$ \\
\hline
$23$ & $18\,643\,272$ & $1\,879\,463\,646\,744$ & $1\,879\,482\,290\,016$ \\
\hline
\end{tabular}
\end{center}
At the anomalous $M_{24}$-classes $\{7A, 7B, 11A, 23A, 23B\}$ (Gaberdiel--Hohenegger--Volpato $2012$; non-Mukai), the local Hecke category carries an $M_{24}$-orbifold twist $\chi^{M_{24}}_p$ with non-trivial projective Schur-multiplier phase in $H^2(M_{24},\U(1))\cong\mathbb{Z}/12$.
\end{proposition}


\section{The CY3 geometry}
\label{sec:k3e-geometry}

Let $(E, e_0)$ be an elliptic curve with an $N$-torsion point and $S$ a K3 surface with elliptic fibration $\pi \colon S \to \mathbb{P}^1$ admitting sections $s_1, s_2 \colon \mathbb{P}^1 \to S$ with $s_2$ of order $N$ relative to $s_1$. The product $S \times E$ admits a free $\mathbb{Z}/N\mathbb{Z}$-action
\[
 (s, e) \longmapsto (s + s_2(\pi(s)), e + e_0),
\]
and the quotient $X = (S \times E)/(\mathbb{Z}/N\mathbb{Z})$ is a projective Calabi--Yau threefold.

\begin{definition}[The DT zeta function]
\label{def:dt-zeta-k3e}
Let $F$ be the fiber class of $\pi$ and $\beta_h = \frac{1}{N}(s_1 + \cdots + s_N + hF)$ for $h \geq 0$. The DT zeta function of $X$ is
\[
 Z^X(q, t, p) = \sum_{h=0}^{\infty} \sum_{d=0}^{\infty} \sum_{n \in \mathbb{Z}} \mathrm{DT}^X_{n,(\beta_h, d)} \, q^{d-1} \, t^{\frac{1}{2}\langle \beta_h, \beta_k \rangle} (-p)^n.
\]
\end{definition}

\begin{theorem}[Oberdieck--Pandharipande]
\label{thm:dt-igusa}
For $X = S \times E$ of order $N = 1$:
\[
 Z^X_{\mathrm{OP/DT}}(q, t, p)
 =-\bigl(\Phi_{10}^{\mathrm{OP}}\bigr)^{-1}
 =-D_5^{-2}
 =-4096\,\Delta_5^{-2},
 \qquad
 D_5=64^{-1}\Delta_5,\quad \Phi_{10}^{\mathrm{OP}}=D_5^2.
\]
The unnormalised DVV/DMVV square is the separate convention
\(\Phi_{10}^{\mathrm{un}}=\Delta_5^2\).
\end{theorem}

\begin{remark}[Convention: \texorpdfstring{$\Delta_5$}{Delta-5} vs \texorpdfstring{$\Phi_{10}^{\mathrm{un}}$}{Phi-10-un}]
\label{rem:k3e-convention-delta5-phi10}
This chapter follows the Gritsenko--Nikulin convention: the Borcherds multiplicative lift of $\phi_{0,1}$ produces an automorphic form $\Delta_5$ of weight $f(0)/2 = 5$ on $\mathrm{O}^+(3,2)$, with product exponents $f(D)$ from $\phi_{0,1}$. The DVV/DMVV square convention uses
\[
 \Phi_{10}^{\mathrm{un}}=\Delta_5^2
\]
for the weight-$10$ Siegel cusp form on $\mathrm{Sp}_4(\mathbb{Z})$ with product exponents \(2f(D)\). The Oberdieck--Pandharipande reduced-DT convention instead uses the monic product \(D_5=64^{-1}\Delta_5\) and
\[
 \Phi_{10}^{\mathrm{OP}}=D_5^2=4096^{-1}\Delta_5^2,\qquad
 Z_{\mathrm{OP/DT}}=-(\Phi_{10}^{\mathrm{OP}})^{-1}.
\]
\end{remark}

\begin{theorem}[KKV \texorpdfstring{$=$}{=} GV \texorpdfstring{$=$}{=} PT identification of the Oberdieck--Pandharipande denominator]
\label{thm:op-kkv-gv-pt-identification}
For $X = S \times E$ with $S$ a K3 surface and $E$ an elliptic curve
(the $N = 1$ case of Theorem~\ref{thm:dt-igusa}), the reduced DT
partition function admits the infinite-product expansion
\[
 Z^X_{\mathrm{red}}(p, q, t)
 \;=\;
 \prod_{\substack{(n,m,\ell) > 0}} (1 - p^n q^m t^\ell)^{-c(nm,\ell)}
 \;=\;
 \frac{p\,q\,t}{\Phi_{10}^{\mathrm{un}}(p,q,t)}
 \;=\;
 \frac{pqt}{\Delta_5(p,q,t)^2},
\]
where $c(k, \ell)$ are the Fourier coefficients of the K3 elliptic
genus $2\,\phi_{0,1}(\tau, z)$ indexed by discriminant $4k - \ell^2$,
the triple $(p, q, t) = (e^{2\pi i \rho}, e^{2\pi i \tau}, e^{2\pi i z})$
corresponds to (Euler characteristic, fibre-class $\chi(\cO_{K3})$, base
class) on the DT side and to the three period coordinates of the
	Igusa domain $\mathbb{H}_2$ on the automorphic side, and
\(\Phi_{10}^{\mathrm{un}}=\Delta_5^2\) is the unnormalised Igusa cusp
form of weight $10$. The
constant $C$ in Theorem~\ref{thm:dt-igusa} is the Weyl-vector prefactor
$\exp(2\pi i (\rho_W, (z_1, z_2, z_3))) = pqt$, so the ``reciprocal-square''
of $\Delta_5$ in the Gritsenko--Clery diagonal-divisor class
(Conjecture~1 of the Gritsenko--Clery theory) is the square of the
Borcherds-lift half-genus denominator: the exponent $2$ is the
\emph{multiplicity-doubling} factor relating the weight-$0$ index-$1$
input $\phi_{0,1}$ to the K3 elliptic genus $2\,\phi_{0,1}$ of the
GV/PT side.
\end{theorem}

\begin{proof}
The identification proceeds by three independent routes, each on a
different side of the KKV $=$ GV $=$ PT correspondence; the three
routes converge on the same product formula.

\emph{Route A: KKV conjecture.} The Katz--Klemm--Vafa conjecture for
$S \times E$, proved by Pandharipande--Thomas (\emph{The Katz--Klemm--Vafa
conjecture for a $K3$ surface}, \emph{Forum Math.\ Pi}~4, 2016,
Theorem~1) in the stable-pairs formulation and by Maulik--Pandharipande--Thomas
(\emph{Curves on K3 surfaces and modular forms}, \emph{J.\ Topology}~3,
2010, Theorem~3) on the Gromov--Witten side, states that the
genus-$g$ Gopakumar--Vafa BPS invariants $n^g_h$ of K3 in primitive
class of square $2h - 2$ are the coefficients of
\[
 \sum_{g, h} (-1)^g n^g_h (y^{1/2} - y^{-1/2})^{2g} q^{h-1}
 \;=\;
 \prod_{n \geq 1} \frac{1}{(1 - q^n)^{20} (1 - y q^n)^2 (1 - y^{-1} q^n)^2},
\]
the reciprocal of the K3 refined elliptic genus. After integration
against the elliptic fibre $E$ (which introduces the $p$-variable
tracking Euler characteristic $n \in \mathbb{Z}$ through the
Hilbert-scheme translation $\mathrm{Hilb}^n(X) \to \mathrm{Hilb}^n(X)/E$),
the product exponents are the Fourier coefficients $c(nm, \ell)$
of $2\,\phi_{0,1}$: these are precisely the discriminant-indexed
multiplicities that appear in Route B.

\emph{Route B: Borcherds multiplicative lift.} The weak Jacobi form
$2\,\phi_{0,1}(\tau, z)$ of weight $0$ and index $1$ is the elliptic
genus of K3. Its Borcherds multiplicative lift
(Borcherds, \emph{Automorphic forms with singularities on Grassmannians},
\emph{Invent.\ Math.}~132 (1998), Theorem~10.1 for the lift's
automorphicity; Theorem~13.3 for the weight formula) produces the
product
\[
 \Phi_{10}^{\mathrm{un}}(p, q, t)
 \;=\;
 p q t \prod_{(n,m,\ell) > 0} (1 - p^n q^m t^\ell)^{c(nm, \ell)},
\]
where the ordering $(n,m,\ell) > 0$ is $n > 0$, or $n = 0$ and $m > 0$,
or $n = m = 0$ and $\ell < 0$, and the constant-coefficient of
$2\,\phi_{0,1}$ is $c(0,0) = 20$. The Borcherds weight theorem gives
$\mathrm{wt}(\Phi_{10}^{\mathrm{un}}) = c(0,0)/2 = 10$
(Theorem~\ref{thm:borcherds-weight-kappa-BKM-universal} in
Chapter~\ref{ch:cy-d-kappa-stratification}); the $pqt$ prefactor is
the Weyl-vector exponential $\exp(2\pi i (\rho_W, z))$ with Weyl vector
$\rho_W$ of $\mathfrak{g}_{\Delta_5}$ (Definition~\ref{def:k3e-weyl-vector}
gives $\rho_W = f_2 - \frac{1}{2} f_3 + f_{-2}$, doubled under the
\(\Delta_5 \leftrightarrow \Phi_{10}^{\mathrm{un}}\) squaring to
$\rho_{\Phi_{10}^{\mathrm{un}}} = 2\rho_W = 2f_2 - f_3 + 2f_{-2}$, so
$e^{2\pi i(\rho_{\Phi_{10}^{\mathrm{un}}}, z)} = pqt$). Dividing out the prefactor
identifies the product with the Borcherds denominator of
$\mathfrak{g}_{\Delta_5}$.

\emph{Route C: PT stable-pairs DT.} The reduced stable-pairs DT
partition function of $S \times E$ in curve class $(\beta_h, d)$ was
computed by Oberdieck--Pandharipande
(\emph{Curve counting on abelian surfaces and threefolds}, \emph{Algebr.\ Geom.}~4, 2017,
Theorem~3) and Oberdieck (\emph{Gromov--Witten invariants of the
Hilbert schemes of points of a K3 surface}, \emph{Geom.\ Topol.}~22,
2018, Theorem~2) in terms of the same Fourier coefficients $c(nm, \ell)$
of $2\,\phi_{0,1}$. The GW/DT correspondence (Maulik--Nekrasov--Okounkov--Pandharipande,
\emph{Gromov--Witten theory and Donaldson--Thomas theory}, I and II,
\emph{Compos.\ Math.}~142, 2006) identifies the GW side Route~A with
the DT side Route~C up to the same normalisation; on both sides, the
\emph{reduced} qualifier quotients by the trivial $E$-tangent direction
via the Mumford relation on $\mathcal{M}_{1,0}$ (Oberdieck--Pandharipande
\emph{loc.\ cit.}, \S2.4). The DMVV square carries the \(pqt\)
Weyl-vector prefactor absorbing the shift between the BPS index
\((d - 1)\) and the \((d,n)\)-bigrading; the OP scalar of
Theorem~\ref{thm:dt-igusa} is obtained only after the monic
normalisation \(D_5=64^{-1}\Delta_5\) and the Behrend sign are imposed.

The three routes agree on the product, yielding
$Z^X_{\mathrm{DMVV}} = \prod (1 - p^n q^m t^\ell)^{-c(nm, \ell)} = pqt/\Phi_{10}^{\mathrm{un}}$.
Since \(\Phi_{10}^{\mathrm{un}}=\Delta_5^2\), this equals
\(pqt/\Delta_5^2\). The OP/DT scalar of Theorem~\ref{thm:dt-igusa}
is obtained only after replacing \(\Delta_5\) by the monic product
\(D_5=64^{-1}\Delta_5\) and inserting the Behrend sign. The exponent-doubling
$c(nm, \ell)$ in \(\Phi_{10}^{\mathrm{un}}\) versus $f(nm, \ell) = c(nm, \ell)/2$ in
$\Delta_5$ (Gritsenko--Nikulin lift of $\phi_{0,1}$) is the
multiplicity-doubling between the K3 elliptic genus $2\phi_{0,1}$ (GV
side) and its weight-$5$ half-lift $\phi_{0,1}$ (Gritsenko side),
matching Remark~\ref{rem:k3e-convention-delta5-phi10}.
\end{proof}

\begin{remark}[The BPS index \texorpdfstring{$=$}{=} K3 elliptic genus coefficient]
\label{rem:op-bps-index-elliptic-genus-coefficient}
Theorem~\ref{thm:op-kkv-gv-pt-identification} extracts a structural
reading of Theorem~\ref{thm:dt-igusa}: the reduced DT protected BPS index
$\Omega(\beta_h, d, n)$ of $S \times E$ in curve class $(\beta_h, d)$
with Euler characteristic $n$ equals, up to sign, the Fourier
coefficient $c(d \cdot h, \ell)$ of the K3 elliptic genus
$2\,\phi_{0,1}$, where $\ell$ is determined by the Mumford linear form
$\ell = n - (d - 1)$. Equivalently, the signed BPS indices that
supply the BKM root character of $\mathfrak{g}_{\Delta_5}$ via
Route~B are precisely the enumerative protected BPS invariants of
$S \times E$ via Routes~A and~C. The ``reciprocal-square'' statement of the
Gritsenko--Clery diagonal-divisor theory (Conjecture~1, twisted
$(N, M)$-case) is the doubling-factor statement that the
K3 elliptic genus on the GV side is $2\phi_{0,1}$, not $\phi_{0,1}$:
	the unnormalised Igusa square \(\Phi_{10}^{\mathrm{un}}\) of the full K3 elliptic genus
is the square of the Borcherds denominator $\Delta_5$ of its
half-lift.
\end{remark}

The elliptic genus of K3, the weak Jacobi form $\phi_{0,1}$ of weight $0$ and index $1$, governs the signed root character of the BKM superalgebra.

\section{The lattice \texorpdfstring{$\Lambda^{3,2}$}{Lambda-3,2} and the isomorphism \texorpdfstring{$\mathrm{Sp}_4 \simeq \mathrm{SO}_+$}{Sp-4 = SO-+}}
\label{sec:k3e-lattice}

Consider the rank-4 free $\mathbb{Z}$-module $\Lambda^4 = \mathbb{Z}e_1 \oplus \mathbb{Z}e_2 \oplus \mathbb{Z}e_3 \oplus \mathbb{Z}e_4$. The exterior square $\bigwedge^2 \colon \mathrm{SL}_4(\mathbb{Z}) \to \mathrm{GL}(\bigwedge^2 \Lambda^4)$ and the Pfaffian scalar product $(u,v) = u \wedge v / (e_1 \wedge e_2 \wedge e_3 \wedge e_4)$ give a signature $(3,3)$ form. The lattice
\[
 \Lambda^{3,2} = (e_1 \wedge e_3 + e_2 \wedge e_4)^\perp \subset \bigwedge^2 \Lambda^4
\]
has signature $(3,2)$ and satisfies $\Lambda^{3,2} \simeq \Lambda^{1,1} \oplus \Lambda^{1,1} \oplus [2]$.

\begin{lemma}
\label{lem:sp4-so-iso}
The correspondence $\bigwedge^2$ defines an isomorphism
\[
 % B-class fix: Sp_4 is 4x4, so quotient by +-I_4; O(Lambda^{3,2}) is
 % 5x5, so quotient by +-I_5.
 \bigwedge^2 \colon \mathrm{Sp}_4(\mathbb{Z})/\{\pm I_4\} \xrightarrow{\;\sim\;} \mathrm{O}(\Lambda^{3,2})_+/\{\pm I_5\}.
\]
\end{lemma}

In the basis $(f_1, f_2, f_3, f_{-2}, f_{-1})$ with $f_1 = e_1 \wedge e_2$, $f_2 = e_2 \wedge e_3$, $f_3 = e_1 \wedge e_3 - e_2 \wedge e_4$, $f_{-2} = e_4 \wedge e_1$, $f_{-1} = e_4 \wedge e_3$, the orthogonal group $\mathrm{O}_\mathbb{R}(\Lambda^{3,2})_+$ acts on the type IV domain $\mathbb{H}^{IV}_+$, which identifies with the Siegel upper half-space $\mathbb{H}_2$.

\section{The hyperbolic sublattice and reflection groups}
\label{sec:k3e-reflections}

Fix the primitive hyperbolic sublattice
\[
 \Lambda^{2,1} = \Lambda^{1,1} \oplus [2] \simeq \mathbb{Z}f_2 \oplus \mathbb{Z}f_3 \oplus \mathbb{Z}f_{-2} \subset \Lambda^{3,2}.
\]
The sublattice $\Lambda^{2,1}_{II}$ generated by elements of type II (those $\delta$ with $(\delta, \Lambda^{2,1}) = 2\mathbb{Z}$) has three real simple roots:
\[
 \delta_1 = 2f_2 - f_3, \qquad \delta_2 = 2f_{-2} - f_3, \qquad \delta_3 = f_3,
\]
with Gram matrix
\[
 (\delta_i, \delta_j) = \begin{pmatrix} 2 & -2 & -2 \\ -2 & 2 & -2 \\ -2 & -2 & 2 \end{pmatrix}.
\]
This matrix has eigenvalues \(4,4,-2\), determinant \(-32\), and
signature \((2,1)\). The subscript \(II\) in
\(\Lambda^{2,1}_{II}\) records the Gritsenko--Nikulin type-II
reflection convention. It is not the even-unimodular notation
\(\mathrm{II}_{2,1}\); even unimodular lattices of signature \((p,q)\)
exist only when \(p-q\equiv0\pmod 8\).

The fundamental polyhedron $\mathcal{P}_{II}$ has $\mathrm{Aut}(\mathcal{P}_{II}) \simeq S_3$, and
\[
 \mathrm{O}(\Lambda^{2,1})_+ \simeq W^{(2)}(\Lambda^{2,1}_{II}) \rtimes \mathrm{Aut}(\mathcal{P}_{II}).
\]

\begin{definition}[Weyl vector]
\label{def:k3e-weyl-vector}
The Weyl vector $\rho \in \mathcal{P}_{II} \otimes \mathbb{Q}$ satisfies $(\rho, \delta_i) = -(\delta_i, \delta_i)/2 = -1$ for all $i$:
\[
 \rho = \tfrac{1}{2}\delta_1 + \tfrac{1}{2}\delta_2 + \tfrac{1}{2}\delta_3 = f_2 - \tfrac{1}{2}f_3 + f_{-2}.
\]
\end{definition}

\section{The generalized BKM superalgebra \texorpdfstring{$\mathfrak{g}_{\Delta_5}$}{g-Delta-5}}
\label{sec:k3e-bkm}

The automorphic correction of the Kac--Moody algebra $\mathfrak{g}$ (with Gram matrix $(\delta_i, \delta_j)$) by the Fourier coefficients of $\Delta_5$ produces the generalized BKM Lie superalgebra $\mathfrak{g}_{\Delta_5}$.

\begin{construction}[Root system of \texorpdfstring{$\mathfrak{g}_{\Delta_5}$}{g-Delta-5}]
\label{constr:k3e-roots}
The root system decomposes as $\Delta = \Delta^{\mathrm{re}} \sqcup \Delta^{\mathrm{im}}_0 \sqcup \Delta^{\mathrm{im}}_1$:
\begin{itemize}
 \item \textbf{Real even simple roots}: $\Delta^{\mathrm{re}}_0 = \Delta^{\mathrm{re}} = \mathcal{P}_{II,\mathrm{prim}} = \{\delta_1, \delta_2, \delta_3\}$.
 \item \textbf{Even imaginary simple roots}: $\Delta^{\mathrm{im}}_0 = \{ \tau(a) \cdot a \mid (a,a) = 0, \, \tau(a) > 0 \}$, repeated $\tau(a)$ times.
 \item \textbf{Odd imaginary simple roots}: $\Delta^{\mathrm{im}}_1 = \{ m(a) \cdot a \mid (a,a) < 0, \, m(a) < 0 \}$, where the element $a$ is repeated $-m(a)$ times and these generators are fermionic.
\end{itemize}
Here $m(a) = -\frac{1}{64} f(n,l,m)$ for $a \in \Lambda^{2,1}_{II} \cap \mathbb{R}_{\geq 0} \mathcal{P}_{II}$ corresponding to the Fourier coefficient $f(n,l,m)$ of $\Delta_5$.
\end{construction}

The superalgebra $\mathfrak{g}_{\Delta_5}$ has the triangular decomposition
\[
 \mathfrak{g}_{\Delta_5} = \biggl(\bigoplus_{\alpha \in \widetilde{C(\Lambda^{2,1}_{II})_+}} \mathfrak{g}_\alpha\biggr) \;\oplus\; (\Lambda^{2,1}_{II} \otimes \mathbb{R}) \;\oplus\; \biggl(\bigoplus_{\alpha \in -\widetilde{C(\Lambda^{2,1}_{II})_+}} \mathfrak{g}_\alpha\biggr)
\]
with generators $h_\alpha, e_\alpha, f_\alpha$ for $\alpha \in \Delta$ satisfying the standard BKM relations.

\begin{proposition}[Cartan spectrum and Weyl vector]
\label{prop:k3e-cartan-and-weyl}
The Gram matrix $G_{ij} = (\delta_i, \delta_j)$ on the three primitive real simple roots is
\[
 G \;=\; \begin{pmatrix} 2 & -2 & -2 \\ -2 & 2 & -2 \\ -2 & -2 & 2 \end{pmatrix}
 \;=\; 4 I - 2 \mathbf{1}\mathbf{1}^\top
\]
with eigenvalues $\{-2, 4, 4\}$; the signature $(2, 1)$ confirms the hyperbolic Lorentzian rank-$3$ character of the Kac--Moody base $\mathfrak{g} \subset \mathfrak{g}_{\Delta_5}$. The Coxeter exponents $m_{ij} = \infty$ for every pair $i \neq j$ realise the three real simple roots as an ideal triangle in $\mathbb{H}^2$, with $S_3 = \mathrm{Aut}(\mathcal{P}_{II})$ acting transitively on the three ideal vertices $\{2 f_2,\; 2 f_{-2},\; 2 f_2 - 2 f_3 + 2 f_{-2}\}$ of $\mathcal{P}_{II}$. The Weyl vector
\[
 \rho \;=\; \tfrac{1}{2}(\delta_1 + \delta_2 + \delta_3) \;\in\; \Lambda^{2,1}_{II} \otimes \mathbb{Q}
\]
is the unique solution of $(\rho, \delta_i) = -1$ for $i = 1, 2, 3$ and satisfies $(\rho, \rho) = -3/2$; it lies in the interior of the positive cone $C(\Lambda^{2,1}_{II})_+$.
\end{proposition}

\begin{proof}[Proof]
Direct computation. $G v_0 = -2 v_0$ for $v_0 = (1, 1, 1)^\top$ and $G v_a = 4 v_a$ for $v_a \in (v_0)^\perp$, giving the spectrum. Signature $(2, 1)$ is read off the eigenvalue signs. Solving $G \rho = -\mathbf{1}$ gives $\rho = \tfrac{1}{2}(\delta_1 + \delta_2 + \delta_3)$; the quadratic form $(\rho, \rho) = -1 - \tfrac{1}{2} = -\tfrac{3}{2}$ follows from $\rho^\top G \rho = -\mathbf{1}^\top \rho = -3/2$. Hyperbolic ideal-triangle geometry and $S_3$-transitivity on the three cusps of $\mathcal{P}_{II}$ are Gritsenko--Nikulin 1998 \cite[\S2--3]{GritsenkoNikulin1998}.
\end{proof}

\begin{remark}[The apparent count of ten real simple roots]
\label{rem:k3e-ten-is-not-real-simples}
The integer $10$ appears twice in the Igusa datum and are distinct with a count of real simple roots. In the product expression $\Delta_5 = \prod_{(a,b) \in \mathcal{T}} \nu_{a,b}$, the cardinality $|\mathcal{T}| = 10$ is the number of even theta constants on the genus-$2$ Siegel upper half space $\mathbb{H}_2$. The zero-Fourier-index value $c_{\phi_{0,1}}(0, 0) = 10$ of the weak Jacobi form $\phi_{0,1}(\tau,z)=y^{-1}+10+y+q(10y^{-2}-64y^{-1}+108-64y+10y^2)+O(q^2)$ is the Borcherds-weight input and gives $\kBKM(\Delta_5) = c_{\phi_{0,1}}(0, 0)/2 = 5$. The number of real simple roots of $\mathfrak{g}_{\Delta_5}$ is $3$, not $10$; the three are the primitive reflection vectors $\{\delta_1, \delta_2, \delta_3\}$ of $\mathcal{P}_{II}$.
\end{remark}

\subsection{Frenkel--Kac / Borcherds vertex-operator closure}
\label{subsec:k3e-frenkel-kac-closure}

The root-system construction of $\mathfrak{g}_{\Delta_5}$ given above is the
\emph{combinatorial} face of the algebra. Its vertex-operator face is a
Borcherds closure of the formal lattice vertex algebra
$V_{\Lambda^{2,1}_{II}}$ after a transverse $c=23$ theory with K3
elliptic supertrace and a no-ghost theorem have been supplied. This is
the classical BRST mechanism behind the denominator. It is not, by
itself, the construction of the framed target-space chiral algebra
$A_{K3\times E}^{(\Sigma_2,C)}$.

\begin{theorem}[Conditional Frenkel--Kac / Borcherds closure of \texorpdfstring{$\mathfrak{g}_{\Delta_5}$}{g-Delta-5}]
\label{thm:k3e-frenkel-kac-closure}
Let $V_{\Lambda^{2,1}_{II}}$ be the formal lattice vertex algebra of the
rank-$3$ type-II hyperbolic lattice $\Lambda^{2,1}_{II}\subset\Lambda^{3,2}$
(Section~\ref{sec:k3e-reflections}), at central charge $c_V = 3$. Let
$V_{\perp}$ be a conformal vertex superalgebra of central charge $23$
and let $V_{\mathrm{ghost}}$ be the bosonic-string ghost system at
$c=-26$. Assume:
\begin{enumerate}[label=\textup{(\roman*)}]
\item the relative BRST operator on
\[
 V_{\Lambda^{2,1}_{II}}\otimes V_{\perp}\otimes V_{\mathrm{ghost}}
\]
is defined, squares to zero, and satisfies the finite-height no-ghost
spectral-sequence convergence needed to identify ghost-number-one
physical states with transverse states;
\item the transverse graded supertrace has Fourier coefficients
\[
 \operatorname{str}_{V_{\perp,D}}(1)=c(D)
\]
in the discriminant convention \(D=4nm-\ell^2\) used in
Theorem~\ref{thm:k3e-denominator};
\item the parity fixture on the target root spaces is fixed so that
ordinary multiplicities refine the signed superdimensions \(c(D)\).
\end{enumerate}
Then, with the fixed parity fixture, there is a root-graded Lie
superalgebra isomorphism
\[
 H^1_{\mathrm{BRST}}\Bigl(V_{\Lambda^{2,1}_{II}}\otimes V_{\perp}\otimes
 V_{\mathrm{ghost}}\Bigr)
 \;\cong\;
 \mathfrak{g}_{\Delta_5},
\]
where \(\mathfrak{g}_{\Delta_5}\) is the BKM superalgebra of
Construction~\ref{constr:k3e-roots}. The Frenkel--Kac vertex operators
$V(\alpha,z) = {:}\exp(\alpha\cdot\phi(z)){:}\otimes e^\alpha$ on
$V_{\Lambda^{2,1}_{II}}$ satisfy
\[
 [V(\alpha,z), V(\beta,w)]
 \;=\;
 \epsilon(\alpha,\beta)\,(z-w)^{\langle\alpha,\beta\rangle}\,
 V(\alpha+\beta,w)\;+\;\text{regular},
\]
and the residue of this commutator at $z = w$, taken on the BRST cohomology
above, is the Borcherds bracket after projection to the
\((\alpha+\beta)\)-root space. Here $\epsilon$ is the Borcherds sign
cocycle satisfying
$\epsilon(\alpha,\beta)\epsilon(\beta,\alpha) = (-1)^{\langle\alpha,\beta\rangle
+ \langle\alpha,\alpha\rangle\langle\beta,\beta\rangle}$ (Borcherds 1986, §5).
\end{theorem}

\begin{proof}
Restrict the Mukai--Heisenberg vertex algebra of Chapter~\ref{ch:k3-yangian},
Theorem~\ref{thm:mukai-frenkel-kac-map}, along the primitive embedding
$\Lambda^{2,1}_{II}\hookrightarrow\Lambda_{\mathrm{Muk}}$. The bosonic Fock space
therefore decomposes as
\[
 V_{\Lambda^{2,1}_{II}}
 \;=\;
 \mathrm{Sym}\!\Bigl(\bigoplus_{n\geq 1}\Lambda^{2,1}_{II}\otimes t^{-n}\Bigr)
 \otimes
 \mathbb{C}_\epsilon[\Lambda^{2,1}_{II}],
\]
and each momentum vector $\alpha\in\Lambda^{2,1}_{II}$ defines the state-field map
\[
 Y(e^\alpha,z)
 \;=\;
 \exp\!\Bigl(\sum_{n\geq 1}\frac{\alpha_{-n}}{n}z^n\Bigr)
 \exp\!\Bigl(-\sum_{n\geq 1}\frac{\alpha_n}{n}z^{-n}\Bigr)e^\alpha z^{\alpha_0}.
\]
The OPE of two such fields is the Frenkel--Kac formula, with the same
bimultiplicative cocycle $\epsilon$ as in the definite case; the only
change is that the pairing now has signature $(2,1)$, so vectors of
norm $\leq 0$ occur. Borcherds's vertex-algebra construction says that
after tensoring an indefinite lattice vertex algebra with a transverse
matter theory and the $c=-26$ ghost system, the relative no-ghost
complex has ghost-number-one cohomology carrying a generalized
Kac--Moody bracket, with the bracket induced by the residue of the
vertex-operator OPE. The hypotheses above supply exactly the pieces
needed for this construction in the present rank-$3$ lattice: total
central charge \(3+23-26=0\), \(Q_{\mathrm{BRST}}^2=0\), finite-height
convergence, and the target parity fixture.

The real simple roots are the norm-$2$ vectors of $\Lambda^{2,1}_{II}$, recovered from
the simple poles with $\langle\alpha,\beta\rangle=-1$ exactly as in Frenkel--Kac 1980.
The imaginary simple roots are the norm-$\leq 0$ vectors; Borcherds's
denominator theorem identifies their signed root-character exponents
with the Fourier coefficients of the transverse character. By
hypothesis those coefficients are \(c(4nm-\ell^2)\), the coefficients
of \(\phi_{0,1}\). Gritsenko--Nikulin identify the resulting
denominator with $\Delta_5$. The BKM superalgebra with this real
simple system, this parity fixture, and these signed imaginary-root
characters is precisely Construction~\ref{constr:k3e-roots}.
\end{proof}

\begin{proposition}[Finite transverse supertrace fixture]
\label{prop:k3e-brst-finite-supertrace-fixture}
Fix a finite discriminant bound \(H\), and let
\[
 \mathcal D_{\leq H}
 =
 \{D\leq H\mid c(D)\neq 0\}
\]
for the coefficients of the normalised Jacobi input
\(\phi_{0,1}(\tau,z)\). For each \(D\in\mathcal D_{\leq H}\) set
\[
T_D^{\min}
 =
 \begin{cases}
  \mathbb C^{\,c(D)\mid 0}, & c(D)>0,\\[2mm]
  \mathbb C^{\,0\mid -c(D)}, & c(D)<0.
 \end{cases}
\]
Then \(T_{\leq H}^{\min}:=\bigoplus_{D\in\mathcal D_{\leq H}}T_D^{\min}\)
is the unique minimal finite super-vector-space fixture with
\[
 \operatorname{sdim}T_D^{\min}=c(D)
 \qquad (D\in\mathcal D_{\leq H}).
\]
More generally, any finite super-vector-space fixture \(T_D\) with
\(\operatorname{sdim}T_D=c(D)\) is obtained from \(T_D^{\min}\) by adding
an even-odd cancelling summand \(Q_D\oplus \Pi Q_D\). Thus the signed
root-character part of the BRST input is not an ambiguity: at finite
height it has a canonical minimal representative, and all nonminimal
representatives have the same supertrace. If \(H'\leq H\), the
projection
\[
 T_{\leq H}^{\min}\longrightarrow T_{\leq H'}^{\min}
\]
onto the summands with \(D\leq H'\) is strict and preserves the
supertrace fixture. Hence the finite minimal fixtures form an inverse
system under discriminant truncation.
\end{proposition}

\begin{proof}
For \(c(D)>0\), the super-vector space
\(\mathbb C^{c(D)\mid0}\) has superdimension \(c(D)\); for \(c(D)<0\),
\(\mathbb C^{0\mid -c(D)}\) has superdimension \(c(D)\). This proves
existence and the displayed supertrace identity. If
\(T_D=T_D^{\bar0}\oplus T_D^{\bar1}\) is any other finite fixture with
\(\dim T_D^{\bar0}-\dim T_D^{\bar1}=c(D)\), then subtracting the
minimal part leaves equal even and odd dimensions. Hence the remainder
is isomorphic, noncanonically, to \(Q_D\oplus\Pi Q_D\) for a finite
vector space \(Q_D\). Minimality is exactly the condition that no such
cancelling summand is present. The transition assertion is immediate
from the direct-sum definition: truncation forgets exactly the summands
with \(H'<D\leq H\), and leaves every retained \(T_D^{\min}\) unchanged.
The executable finite-height witnesses are
\texttt{brst\_coefficient\_fixture} and
\texttt{brst\_coefficient\_fixture\_transition} in
\texttt{compute/lib/k3e\_finite\_bridge\_witness.py}; at
\(H=8\) it gives support \(\{-1,0,3,4,7,8\}\) with signed dimensions
\((1,10,-64,108,-513,808)\), and the transition \(8\to4\) retains
\(\{-1,0,3,4\}\) without defect.
\end{proof}

\begin{remark}[What the fixture does not construct]
\label{rem:k3e-brst-fixture-boundary}
Proposition~\ref{prop:k3e-brst-finite-supertrace-fixture} fixes the
finite signed-character and parity input in
Theorem~\ref{thm:k3e-frenkel-kac-closure}. It does not construct a
conformal vertex superalgebra realising the \(T_D^{\min}\) as physical
state spaces, does not construct \(Q_{\mathrm{BRST}}\), and does not
prove no-ghost convergence. Those are the remaining BRST tasks. The
point is narrower but essential: the coefficient side of the BRST
problem is now a finite algebraic fixture, not an informal parity
choice.
\end{remark}

\begin{proposition}[Finite no-ghost spectral-sequence gate]
\label{prop:k3e-brst-no-ghost-spectral-sequence-gate}
Fix a finite discriminant window \(H\). Suppose a finite filtered BRST
packet has been supplied: for every \(D\leq H\) in the retained
support, its \(E_1\)-page separates into transverse, longitudinal, and
ghost summands with signed dimensions
\[
T_D,\qquad L_D,\qquad G_D,
\]
and the target Borcherds coefficient is \(c(D)\). Suppose also that
the possible higher-page leakage out of the retained ghost-number-one
transverse sector is represented by finite matrices
\[
\partial_r^{\leq H}\qquad (r\geq2)
\]
in chosen bases. Then the supplied packet satisfies the finite
no-ghost gate precisely when
\[
L_D+G_D=0,\qquad T_D=c(D)
\quad(D\leq H),
\qquad
\operatorname{rank}\partial_r^{\leq H}=0\quad(r\geq2).
\]
On this locus the signed dimension of the retained physical
ghost-number-one BRST cohomology at discriminant \(D\) is \(c(D)\).
\end{proposition}

\begin{proof}
The statement is a finite spectral-sequence calculation. The \(E_1\)
page has been supplied with a direct separation into transverse,
longitudinal, and ghost summands. The no-ghost cancellation on the
finite page is exactly the equality \(L_D+G_D=0\); after this
cancellation, the signed \(E_1\)-character in discriminant \(D\) is
the transverse signed dimension \(T_D\). The coefficient condition
\(T_D=c(D)\) identifies that surviving character with the
Borcherds--Jacobi target coefficient.

The only remaining way for the finite spectral sequence to change this
signed ghost-number-one character inside the retained window is a
higher-page differential leaving or entering the retained transverse
sector. In finite bases such leakage is absent exactly when every
matrix \(\partial_r^{\leq H}\) has rank zero. Under the three displayed
conditions the \(E_1\)-page character survives to \(E_\infty\), hence
to the finite physical BRST cohomology, with signed dimension \(c(D)\).

The executable witness is
\texttt{brst\_no\_ghost\_spectral\_sequence\_gate}. It records the
longitudinal--ghost cancellation defects, the transverse coefficient
defects, and the exact ranks of the supplied higher-page leakage
matrices. The proposition begins with a finite filtered BRST packet
already supplied; it does not construct the worldsheet vertex
superalgebra, the BRST charge, or the no-ghost theorem.
\end{proof}

\begin{proposition}[Finite BRST--Borcherds bracket gate]
\label{prop:k3e-brst-borcherds-bracket-gate}
Fix a finite root window \(R_{\leq H}\). Suppose homogeneous physical
BRST classes \(\{e_i\}_{i\in I}\) and homogeneous Borcherds root
vectors \(\{f_i\}_{i\in I}\) are indexed by the same finite root set,
with parities \(|i|\in\mathbb Z/2\). Suppose the finite brackets are
given in these bases by coefficient tensors
\[
 [e_i,e_j]_{\mathrm{BRST}}=\sum_k B_{ij}^{k}e_k,
 \qquad
 [f_i,f_j]_{\Delta_5}=\sum_k C_{ij}^{k}f_k.
\]
Let \(s(i,j)\) be the supplied root-sum support datum: either
\(s(i,j)=k\), meaning \(\alpha_k=\alpha_i+\alpha_j\) is retained, or
\(s(i,j)=\varnothing\), meaning the sum is outside the finite root
window. Then the finite BRST bracket packet is compatible with the
retained Borcherds bracket precisely when
\[
 B_{ij}^{k}=C_{ij}^{k}\quad(i,j,k\in I),
 \qquad
 B_{ij}^{k}=0\quad\text{unless }s(i,j)=k,
\]
and the supplied BRST coefficients satisfy
\[
 B_{ij}^{k}+(-1)^{|i||j|}B_{ji}^{k}=0
\]
and
\[
 (-1)^{|i||\ell|}\sum_a B_{j\ell}^{a}B_{ia}^{k}
 +(-1)^{|j||i|}\sum_a B_{\ell i}^{a}B_{ja}^{k}
 +(-1)^{|\ell||j|}\sum_a B_{ij}^{a}B_{\ell a}^{k}=0
\]
for all retained \(i,j,\ell,k\). On this locus \(e_i\mapsto f_i\)
identifies the supplied finite physical BRST bracket table with the
retained finite Borcherds bracket table as a root-graded Lie
superalgebra packet.
\end{proposition}

\begin{proof}
All vector spaces in the statement are finite-dimensional and based.
The equality \(B_{ij}^{k}=C_{ij}^{k}\) is exactly the coordinate form
of bracket preservation under \(e_i\mapsto f_i\). The support condition
is the coordinate form of root grading: the bracket of roots
\(\alpha_i\) and \(\alpha_j\) can contribute only to the retained root
\(\alpha_i+\alpha_j\), and contributes nothing when that root is
absent from the finite window. The two displayed identities are the
super-skew symmetry and super-Jacobi identity written in the chosen
homogeneous basis. These conditions are therefore necessary.

Conversely, assume the four displayed conditions. Super-skew and
super-Jacobi make the supplied BRST coefficient tensor a finite
root-graded Lie superalgebra bracket on the retained span. The support
condition gives the root grading. The equality \(B=C\) then carries
this finite bracket to the retained Borcherds bracket table coefficient
by coefficient. Hence the finite comparison is bracket-preserving.

The executable witness is \texttt{brst\_borcherds\_bracket\_gate}. It
checks coefficient equality, root-sum support, super-skew symmetry, and
super-Jacobi exactly. The witness starts with the coefficient tensors
already supplied; it does not construct the worldsheet vertex
superalgebra, the integrated vertex operators, their OPE coefficients,
the BRST charge, or the global Hall comparison.
\end{proof}

\begin{proposition}[Finite Borcherds--Serre relation gate]
\label{prop:k3e-brst-borcherds-serre-relation-gate}
In the setup of Proposition~\ref{prop:k3e-brst-borcherds-bracket-gate},
suppose the retained real-root Serre exponents \(m_{ij}\geq1\) and the
retained orthogonal imaginary pairs \(\mathcal I_{\perp}\subset
I\times I\) have been supplied. Define
\[
(\operatorname{ad}e_i)^m(e_j)
 =
[e_i,[e_i,\ldots,[e_i,e_j]_{\mathrm{BRST}}\ldots]_{\mathrm{BRST}}]_{\mathrm{BRST}}
\]
with \(m\) occurrences of \(e_i\). The supplied finite BRST bracket
packet satisfies the retained Borcherds--Serre relations precisely when
\[
(\operatorname{ad}e_i)^{m_{ij}}(e_j)=0
\quad\text{for every retained real-root Serre pair }(i,j),
\]
and
\[
[e_i,e_j]_{\mathrm{BRST}}=0
\quad\text{for every }(i,j)\in\mathcal I_{\perp}.
\]
Equivalently, after the basis \(\{e_k\}\) is fixed, every coefficient
of these finitely many iterated brackets is zero. On this locus the
finite bracket comparison of
Proposition~\ref{prop:k3e-brst-borcherds-bracket-gate} respects the
retained real-root Serre ideal and the retained imaginary
supercommutativity relations.
\end{proposition}

\begin{proof}
The iterated adjoint expression is obtained by repeated application of
the finite coefficient tensor \(B_{ij}^{k}\). At each step the result is
a finite vector in the retained span. Thus
\((\operatorname{ad}e_i)^{m_{ij}}(e_j)=0\) is equivalent to the
vanishing of finitely many coordinates. The same statement for an
orthogonal imaginary pair is the \(m=1\) bracket vector
\([e_i,e_j]_{\mathrm{BRST}}\). These are exactly the coordinate
equations defining the retained Serre and supercommutativity quotient.

The executable witness is
\texttt{brst\_borcherds\_serre\_relation\_gate}. It evaluates the
iterated adjoint words and the declared imaginary supercommutators
from the supplied bracket coefficients. It does not construct the
singular terms of the vertex-operator OPE, and it does not prove that
the declared finite relation list is complete for the global
Borcherds--Serre ideal.
\end{proof}

\begin{proposition}[Finite BRST momentum--height projection gate]
\label{prop:k3e-brst-momentum-height-projection-gate}
Let \(C^{\bullet}_{\leq H_+}\) be a supplied finite BRST chain packet
with homogeneous bases in each degree. Each basis vector carries a
lattice momentum \(\alpha\) and a height \(\operatorname{ht}(\alpha)\).
Let
\[
 Q^p\colon C^p_{\leq H_+}\longrightarrow C^{p+1}_{\leq H_+}
\]
be the supplied BRST differential matrices in these bases. For
\(H\leq H_+\), write \(C^p_{\leq H}\) for the span of basis vectors of
height at most \(H\), and \(K^p_{>H}\) for the complementary killed
span. Then the height-\(H\) BRST subquotient and its transition map
are well-defined exactly when the following finite conditions hold:
\[
 Q^{p+1}Q^p=0,
\qquad
 Q^p_{\beta,\alpha}=0\quad\text{if }\beta\neq\alpha,
\]
\[
 Q^p(C^p_{\leq H})\subset C^{p+1}_{\leq H},
\qquad
 Q^p(K^p_{>H})\subset K^{p+1}_{>H}.
\]
Equivalently, in the block decomposition
\[
C^p_{\leq H_+}=C^p_{\leq H}\oplus K^p_{>H},
\qquad
C^{p+1}_{\leq H_+}=C^{p+1}_{\leq H}\oplus K^{p+1}_{>H},
\]
the upper-right and lower-left off-diagonal blocks of \(Q^p\) vanish,
the momentum-off-diagonal entries vanish, and the upper complex has
square zero. On this locus
\[
 Q^p_{\leq H}=\pi_{\leq H}^{p+1}Q^p\iota_{\leq H}^{p}
\]
defines a finite BRST complex \(C^\bullet_{\leq H}\), and the
projection
\[
R^{\mathrm{BRST}}_{H_+,H}\colon
(C^\bullet_{\leq H_+},Q)\longrightarrow
(C^\bullet_{\leq H},Q_{\leq H})
\]
is a chain map. It induces the height-transition map on
\(H^1_{\mathrm{BRST}}\).
\end{proposition}

\begin{proof}
The first displayed equality is the finite condition that the supplied
upper packet is a complex. The momentum equation is the coordinate
form of \(Q_{\mathrm{BRST}}\) preserving lattice momentum. The inclusion
\(Q^p(C^p_{\leq H})\subset C^{p+1}_{\leq H}\) is exactly the vanishing
of the retained-to-killed block. The inclusion
\(Q^p(K^p_{>H})\subset K^{p+1}_{>H}\) is exactly the vanishing of the
killed-to-retained block. These two block vanishings make the retained
span a subcomplex and the killed span a subcomplex, hence also make
the retained coordinates the quotient complex by the killed span.

With the block vanishings imposed, the formula
\(Q^p_{\leq H}=\pi_{\leq H}^{p+1}Q^p\iota_{\leq H}^{p}\) is the
retained diagonal block of \(Q^p\). Since \(Q^{p+1}Q^p=0\), the
retained diagonal blocks also square to zero. The projection commutes
with the differentials because the killed-to-retained block vanishes:
\[
\pi_{\leq H}^{p+1}Q^p=Q^p_{\leq H}\pi_{\leq H}^{p}.
\]
Thus \(R^{\mathrm{BRST}}_{H_+,H}\) is a chain map and induces the
displayed map on first cohomology.

The executable witness is
\texttt{brst\_momentum\_height\_projection\_gate}. It records the
momentum-off-diagonal entries, the two off-diagonal height-block ranks,
and the exact ranks of the upper and induced lower differential
squares. It does not construct the worldsheet BRST complex; it checks
the finite condition that any supplied BRST complex must satisfy before
heightwise cohomology projection has mathematical content.
\end{proof}

\begin{remark}[Status of the \(K3\times E\) worldsheet reading]
\label{rem:k3e-brst-conditional-status}
Taking \(V_{\perp}\) to be the transverse K3 sector is the expected
worldsheet realisation of the theorem, and the coefficient condition is
the K3 elliptic genus. The present chapter does not construct that
worldsheet VOA, its BRST differential, or its comparison with
\(A_{K3\times E}^{(\Sigma_2,C)}\); those inputs are the proof
obligations isolated in Remark~\ref{rem:k3e-brst-missing}.
\end{remark}

\begin{remark}[Mukai--Heisenberg vertex algebra as ambient]
\label{rem:k3e-mukai-heisenberg-ambient}
The lattice part of the conditional Frenkel--Kac construction sits
inside the rank-$24$ Mukai--Heisenberg vertex algebra
$\mathcal{H}_{\mathrm{Muk}} = V_{\Lambda_{\mathrm{Muk}}}$ of
Chapter~\ref{ch:k3-yangian}, Definition~\ref{def:mukai-voa}, via the primitive
lattice embedding
$\Lambda^{2,1}_{II}\hookrightarrow\Lambda_{\mathrm{Muk}} = \mathrm{II}_{4,20}$.
The abelian rank-$24$ Heisenberg input is the classical-limit Lie algebra of
the full Mukai vertex algebra; the non-abelian BKM
$\mathfrak{g}_{\Delta_5}$ is obtained only after the transverse
BRST package of Theorem~\ref{thm:k3e-frenkel-kac-closure} is supplied.
The three-branch structure of
Chapter~\ref{ch:k3-yangian} Theorem~\ref{thm:three-branches-mukai} is the
organisation of these manifestations: (a) full Mukai--Heisenberg vertex algebra; (b) BKM
sublattice BRST (this section); (c) $\widetilde{M}_{24}$-twisted
$24$-Miki quantum toroidal presentation
(Chapter~\ref{ch:k3-chiral-bialgebra-platonic},
Theorem~\ref{thm:k3-24miki-presentation}).
\end{remark}

\begin{theorem}[Root hyperplanes as Humbert / 1/4-BPS walls]
\label{thm:k3e-root-hyperplanes-walls}
Write
\[
 Z
 \;=\;
 \begin{pmatrix}
 \tau & z \\
 z & \sigma
 \end{pmatrix}
 \in \mathbb{H}_2,
 \qquad
 \alpha=(n,\ell,m)\in \Lambda^{2,1}_{II},
 \qquad
 (\alpha,Z):=n\sigma+\ell z+m\tau.
\]
Then the Borcherds product factor attached to $\alpha$ has signed exponent
\[
 \bigl(1-e^{2\pi i(\alpha,Z)}\bigr)^{\operatorname{sdim}\mathfrak g_{\Delta_5,\alpha}},
 \qquad
 \operatorname{sdim}\mathfrak g_{\Delta_5,\alpha}=c(4nm-\ell^2),
\]
and its vanishing locus
\[
 \mathcal{W}_\alpha
 \;:=\;
 \{Z\in\mathbb{H}_2 : (\alpha,Z)\in\mathbb{Z}\}
\]
is the Humbert / Heegner wall associated to the root $\alpha$.
More precisely:
\begin{enumerate}[label=\textup{(\roman*)}]
\item If $\alpha^2 = 2$, then $\mathcal{W}_\alpha$ is a primitive real-root wall and the
 corresponding residue of the vertex-operator commutator realises the Weyl reflection
 in the real simple root.
\item If $\alpha^2 \leq 0$, then $\mathcal{W}_\alpha$ is an imaginary-root wall; its
 signed root-character exponent is the coefficient of the K3 elliptic genus at
 discriminant $D = 4nm-\ell^2$. The corresponding ordinary simple-root
 dimension is asserted only after the parity fixture or a finite Hall--Borcherds
 recognition computation.
\item In the Type-IIB / heterotic $1/4$-BPS frame, the same hyperplane is the wall of
 marginal stability on which the DVV factor
 $\bigl(1-p^m q^n y^\ell\bigr)^{-c(4nm-\ell^2)}$ of
 \(1/\Phi_{10}^{\mathrm{un}}\) contributes the jump in dyon degeneracy.
\end{enumerate}
\end{theorem}

\begin{proof}
The multiplicative lift of Theorem~\ref{thm:k3e-borcherds-product} is a product over
triples $(n,\ell,m)$, and each factor depends on $Z$ only through the linear functional
$(\alpha,Z)=n\sigma+\ell z+m\tau$. Hence its divisor is the Heegner hyperplane
$\mathcal{W}_\alpha$. For norm-$2$ vectors this is the ordinary real-root hyperplane
appearing in the Weyl chamber decomposition; for norm $\leq 0$ vectors it is the
imaginary-root divisor whose signed exponent is read off from the coefficient
$c(4nm-\ell^2)$ of the K3 elliptic genus. Squaring the product gives the
unnormalised Igusa form \(\Phi_{10}^{\mathrm{un}}\), so the reciprocal
DVV partition function is
\[
 \frac{1}{\Phi_{10}^{\mathrm{un}}(Z)}
 \;=\;
 \prod_{(n,\ell,m)>0}
 \bigl(1-p^m q^n y^\ell\bigr)^{-c(4nm-\ell^2)}.
\]
The poles of this product are exactly the walls where the $1/4$-BPS spectrum jumps.
Sen's wall-crossing analysis identifies those poles with decays into two $1/2$-BPS
constituents, so the same hyperplanes are simultaneously BKM root hyperplanes, Humbert
divisors, and physical marginal-stability walls.
\end{proof}

\begin{remark}[Abelian-at-Lie / non-abelian-at-vertex on \texorpdfstring{$\mathfrak{g}_{\Delta_5}$}{g-Delta-5}]
\label{rem:k3e-abelian-at-lie-vertex}
The Lie bracket on $\mathfrak{g}_{\Delta_5}$, viewed purely as a
Lie-algebraic structure on the $\Lambda^{2,1}_{II}$-graded root spaces, is the
non-abelian bracket of Construction~\ref{constr:k3e-roots} and
Theorem~\ref{thm:k3e-frenkel-kac-closure} on the conditional BRST
closure locus. This non-abelianity is \emph{not} primary: under the
hypotheses of that theorem, it is derived from the OPE of vertex operators
$V(\alpha,z)V(\beta,w)$ inside the lattice VOA
$V_{\Lambda^{2,1}_{II}}$, whose underlying classical-limit Lie algebra is the
abelian-modulo-centre affine Heisenberg
$\widehat{\mathcal{H}}_{\Lambda^{2,1}_{II}}$. In the taxonomy of
Chapter~\ref{ch:k3-chiral-bialgebra-platonic}
Theorem~\ref{thm:k3-abelian-at-lie}, $\mathfrak{g}_{\Delta_5}$ lives at the
vertex-operator-level, not the Lie-primary level; the K3 chiral bialgebra
$\mathbf{H}_{\Delta_5}$ is the $(\infty,1)$-quasi-Hopf quantisation of the
Lie-primary (abelian-modulo-centre) data, with $\mathfrak{g}_{\Delta_5}$ as
its vertex-level shadow.
\end{remark}

\section{The Weyl--Kac--Borcherds denominator identity for \texorpdfstring{$\Delta_5$}{Delta-5}}
\label{sec:k3e-denominator}

\begin{theorem}[Denominator identity]
\label{thm:k3e-denominator}
The Weyl--Kac--Borcherds character formula applied to the trivial one-dimensional representation of $\mathfrak{g}_{\Delta_5}$ gives
\[
 \frac{1}{64} \Delta_5(2Z) = \Phi(z)
\]
where $\Phi$ is the denominator function
\[
 \Phi = \sum_{w \in W^{(2)}(\Lambda^{2,1})} \det(w) \exp(-\pi i \langle w(\rho), z \rangle) - \sum_{a \in \Lambda^{2,1}_{II} \cap \mathbb{R}_{\geq 0} \mathcal{P}_{II}} m(a) \exp(-\pi i \langle w(\rho + a), z \rangle).
\]
Equivalently, in product form:
\[
 \Phi = \exp(-2\pi i \langle \rho, z \rangle) \prod_{\alpha \in \Delta_+} (1 - \exp(-2\pi i \langle \alpha, z \rangle))^{\operatorname{sdim}\mathfrak g_\alpha}.
\]
\end{theorem}

\begin{theorem}[Product formula for \texorpdfstring{$\Delta_5$}{Delta-5}]
\label{thm:k3e-product}
\[
 \frac{1}{64} \Delta_5(Z) = -\exp(\pi i(z_1 + z_2 + z_3)) \prod_{\substack{n,l,m \in \mathbb{Z}, \, (n,l,m) > 0}} (1 - \exp(2\pi i(nz_1 + lz_2 + mz_3)))^{f(nm,l)}
\]
where $Z = \bigl(\begin{smallmatrix} z_1 & z_2 \\ z_2 & z_3 \end{smallmatrix}\bigr) \in \mathbb{H}_2$ and $f(n,l)$ are the Fourier coefficients of the weak Jacobi form $\phi_{0,1}$ of weight $0$ and index $1$.
The sign $(-1)$ arises from the unique negative-discriminant factor $(n,l,m) = (0,-1,0)$: setting $r = \exp(2\pi i z_2)$ with $|r| < 1$, we have $(1 - r^{-1})^{f(0,-1)} = (1 - r^{-1})^1 = -(1/r)(1-r)$, where $f(0,-1) = 1$ is the Fourier coefficient of $\phi_{0,1}$ at $q^0 y^{-1}$ (the two-index convention $f(nm,l)$; the discriminant-indexed coefficient $c(D = -1) = 2$ counts both $l = +1$ and $l = -1$).
Absorbing this factor, the sign-free form reads
\[
 \frac{1}{64} \Delta_5(Z) = \exp(\pi i(z_1 - z_2 + z_3)) \cdot (1 - \exp(2\pi i z_2)) \prod_{\substack{(n,l,m) > 0 \\ (n,l,m) \neq (0,-1,0)}} (1 - \exp(2\pi i(nz_1 + lz_2 + mz_3)))^{f(nm,l)}.
\]
\end{theorem}

\begin{remark}[Classical attribution]
\label{rem:k3e-denominator-classical}
Both forms of the denominator identity; the Weyl--Kac--Borcherds sum
and the Siegel product expansion; are classical. The sum form is
Borcherds's denominator formula for generalised Kac--Moody algebras
\cite{Borcherds1992}; the product form for $\Delta_5$ as a paramodular
Siegel cusp form on $\mathrm{Sp}_4(\mathbb{Z})$ is Gritsenko--Nikulin
\cite{GritsenkoNikulin1995,GritsenkoNikulin1998}, refined by Gritsenko
\cite{Gritsenko1999} for the elliptic-genus input. The equivalence of
the two is Borcherds's theta-correspondence machine
\cite{Borcherds1998}. The second-quantised K3 elliptic genus of
Dijkgraaf--Moore--Verlinde--Verlinde~\cite{DMVV} is the unnormalised
Igusa square \(\Phi_{10}^{\mathrm{un}}=\Delta_5^2\). The OP reduced-DT
scalar instead uses the monic Borcherds normalization
\(D_5=64^{-1}\Delta_5\), hence
\(\Phi_{10}^{\mathrm{OP}}=D_5^2=4096^{-1}\Delta_5^2\) and reciprocal
\(-4096\,\Delta_5^{-2}\). The Borcherds construction exhibits
$\mathfrak{g}_{\Delta_5}$ as the primitive BKM half whose denominator is
$\Delta_5$ and whose normalized square gives the reduced-DT character
\(-(\Phi_{10}^{\mathrm{OP}})^{-1}=-D_5^{-2}=-4096\,\Delta_5^{-2}\).
The separate bar Euler product
identification for the framed chiral algebra $A_{K3 \times E}$ remains
the CY-B/CY-D bridge problem named in Conjecture~\ref{conj:bkm-bar-dictionary};
the primitive denominator statement does not depend on that bridge. The denominator identity itself, the product
expansion, $\kappa_{\mathrm{BKM}}(\Delta_5)=5$, and the Igusa-square
weight \(\mathrm{wt}(\Phi_{10}^{\mathrm{un}})=10=2\kappa_{\mathrm{BKM}}(\Delta_5)\)
are classical: Borcherds~$1998$ proves the automorphic-product theorem,
and Gritsenko--Nikulin~$1998$ identifies the \(\Delta_5\) denominator and
its unnormalised square.
\end{remark}

\section{The weak Jacobi form \texorpdfstring{$\phi_{0,1}$}{phi-0,1} and signed root characters}
\label{sec:k3e-phi01}

The weak Jacobi form $\phi_{0,1} = \phi_{12,1}/\delta_{12}$ (where $\delta_{12} = q \prod_{n \geq 1} (1 - q^n)^{24}$ is the weight-12 cusp form) is the K3 elliptic genus. Its Fourier coefficients $f(n,l)$ are the signed superdimensions/root-character coefficients of the target root spaces of $\mathfrak{g}_{\Delta_5}$:
\[
 \operatorname{sdim}\mathfrak g_{\Delta_5,\alpha} = f(nm, l) \quad \text{for } \alpha = (n,l,m) \in \Delta_+.
\]

\subsection{The universal Borcherds-weight identity \texorpdfstring{$\kBKM(\Phi_N) = c_N(0)/2$}{kBKM(Phi-N) = c-N(0)/2}}
\label{subsec:k3e-universal-kBKM-identity}

The Igusa cusp form $\Delta_5$ is the $N = 1$ entry of an automorphic atlas whose $\kBKM$ weights are uniformly $c_N(0)/2$ with $c_N(0) = \phi_{0,1}^{(g_N, h_M)}(0, 0)$ the zero Fourier coefficient of the input $(g_N, h_M)$-twined K3 elliptic genus. The identity admits two coexistent scopes.

\begin{theorem}[Universal Borcherds-weight identity, two scopes;
\ClaimStatusProvedElsewhere]
\label{thm:k3e-universal-kBKM-two-scopes}
The weight identity $\kBKM(\Phi) = \mathrm{wt}(\Phi) = c(0)/2$ for
Borcherds lifts is proved elsewhere (Borcherds~1998
Theorem~$13.3$; Gritsenko~1999 Theorem~$1.2$; cusp-weighted form
Gritsenko--Cl\'ery 2008 Theorem~$3.1$). What this chapter
contributes is the assembly: the identification of the two scopes,
the per-scope primary sources, and the per-scope constant tables.
The identity $\kBKM(\Phi_N) = c_N(0)/2$ holds at two distinct scopes, each governed by its own primary source and cover stratum.
\begin{description}
\item[Scope (A); CHL-averaged paramodular ladder.]
On the CHL admissible slice $N \in \{1, 2, 3, 4, 6\}$ with $\varphi(N) \mid 2$, the paramodular Borcherds lifts $\Phi_N = \Delta_{k(N)}$ of the $g_N$-averaged K3 elliptic genus have zero Fourier coefficients and Borcherds weights
\[
 \bigl(c_N(0)\bigr)_{N \in \{1,2,3,4,6\}} \;=\; (10,\ 6,\ 4,\ 3,\ 2), \qquad
 \bigl(\kBKM(\Phi_N)\bigr)_{N \in \{1,2,3,4,6\}} \;=\; \bigl(5,\ 3,\ 2,\ \tfrac32,\ 1\bigr).
\]
Primary sources are assembled: Gritsenko--Nikulin~1995~\cite{GritsenkoNikulin1995} Part~II Theorem~$2.1$ establishes $N = 1$ with $\Phi_1 = \Delta_5$, weight $5$, paramodular framework, and Borcherds-product identity; Jatkar--Sen~$2006$ (\emph{JHEP}~$11$, $072$) constructs the CHL dyon-counting squares of weights $2k(N) = c_N(0) = (10, 6, 4, 3, 2)$, with primitive square-root denominator of weight $3$ at $N = 2$; Govindarajan--Krishna~2010~\cite{GovindarajanKrishna2010} covers $N \in \{3, 4, 6\}$ using the Mathieu $M_{24}$ conjugacy classes $\{3A, 4A, 6A\}$ with weights $2, \tfrac32, 1$ respectively. The $N = 4$ entry has half-integral weight $\tfrac32$: its square-root denominator carries a multiplier system, so the ladder lives on the paramodular groups $\Gamma_{\mathrm{para}}^{(N)} \subset \mathrm{Sp}_4(\mathbb{Z})$ with multiplier at $N = 4$. The Eguchi--Hikami~2011~\cite{EguchiHikami2011} weight-$4$ form at $N = 2$ belongs to the separate Mathieu-twined ladder of constants $(20, 8, 6, 4, 2)$ and weights $(10, 4, 3, 2, 1)$, not to this square-root ladder. Here $\kBKM$ is the weight of the Weyl--Kac--Borcherds denominator identity of the sibling BKM superalgebra $\mathfrak{g}_{\Phi_N}$.

\item[Scope (B); Gritsenko--Cl\'ery dd-modular atlas.]
On the Gritsenko--Cl\'ery classification of genus-$2$ dd-modular forms
with diagonal divisor, the indexing set is the triple list
\[
 (t,N;k)\in\{(1,1;5),(2,1;2),(1,2;3),(3,1;1),(1,3;2),
 (4,1;\tfrac12),(1,4;\tfrac32),(2,2;1)\}.
\]
For the corresponding forms $F_{t,N}$ one has
\[
 \bigl(\kBKM(F_{t,N})\bigr)
 =
 \bigl(\mathrm{wt}(F_{t,N})\bigr)
 =
 (5,\ 2,\ 3,\ 1,\ 2,\ \tfrac12,\ \tfrac32,\ 1),
\]
in the displayed order. Equivalently the cusp-weighted Borcherds
constant terms are
\[
 \bigl(c_{t,N}(0,0)\bigr)=(10,\ 4,\ 6,\ 2,\ 4,\ 1,\ 3,\ 2)
\]
at the single-cusp rows, with the multi-cusp rows read by the
Gritsenko--Cl\'ery cusp-weighted sum. Primary: Gritsenko--Cl\'ery
2008 Theorem~$1.4$ and Theorem~$3.1$; Borcherds~1998
Theorem~$13.3$ for the singular-theta weight formula. No
weight-$0$ or weight-$1/4$ dd-modular entry occurs in this catalogue.
\end{description}
\end{theorem}

\begin{remark}[Cover stratification by weight integrality]
\label{rem:k3e-cover-stratification-by-weight}
Scope (B) of Theorem~\ref{thm:k3e-universal-kBKM-two-scopes} partitions the Gritsenko--Cl\'ery 8-form atlas by the arithmetic of the weight:
\[
 \Phi_m \;\in\; \begin{cases}
 \Gamma_t(N)^+ \subset \mathrm{Sp}_2(\mathbb Q) & \text{if } \kBKM(\Phi_m) \in \mathbb{Z}_{\geq 1}, \\
 (\Gamma_t(N)^+, \chi_r) & \text{if } \kBKM(\Phi_m) \in \tfrac{1}{2}\mathbb{Z}_{>0}\setminus\mathbb Z .
 \end{cases}
\]
The six integer-weight forms
\[
\Delta_5,\ \Delta_2,\ \nabla_3,\ \Delta_1,\ \nabla_2,\ Q_1
\]
live on paramodular subgroups and their normal extensions. The two
half-integral forms $\Delta_{1/2}$ and $\nabla_{3/2}$ are holomorphic
Siegel modular forms with multiplier systems of order $8$ and $4$,
respectively. The multiplier-system language is essential: the
Gritsenko--Cl\'ery dd-modular catalogue is not a catalogue of
metaplectic fourfold-cover BKM denominators, and it contains no
weight-$0$ or quarter-weight form.
\end{remark}

\begin{remark}[Scope (A) versus Scope (B): coincidence at $N = 1$, divergence elsewhere]
\label{rem:k3e-two-scopes-coincidence-divergence}
The two scopes agree only at the common entry $\Delta_5$ (Scope~(A)
$N=1$, Scope~(B) $(t,N)=(1,1)$). Scope~(A) is indexed by CHL
orbifold order and produces the ladder
$(c_N(0))=(10,6,4,3,2)$ and
$\kBKM(\Phi_N)=(5,3,2,\tfrac32,1)$ on $N\in\{1,2,3,4,6\}$.
Scope~(B) is indexed by Gritsenko--Cl\'ery triples $(t,N;k)$ and
produces weights $(5,2,3,1,2,\tfrac12,\tfrac32,1)$ in the order of
Theorem~\ref{thm:k3e-universal-kBKM-two-scopes}. The numerical
agreement of the GC weights $3, 2, \tfrac32$ at the rows
$(1,2), (1,3), (1,4)$ with the CHL values at orders $2, 3, 4$
identifies forms only through a named normalisation theorem: the two
scopes use different arithmetic inputs and different paramodular
groups. An inscription of a
$\kBKM$ value must name the scope before stating the identity.
\end{remark}

\begin{theorem}[BV partition-function residue at the Humbert locus:
conditional realisation of $\kBKM(\Delta_5) = 5$;
\ClaimStatusConditional]
\label{thm:bl-bv-k3e-kappa5}
The Borcherds weight identity used below is classical. The BV residue
realisation assumes the Costello--Li hCS quantisation, \v{C}ech--Ran
descent, and one-loop localisation hypotheses specified in the proof,
together with the following residue hypothesis, which is named and not
discharged here:
\begin{description}
\item[(H-res).] The one-loop wheel-graph residue at the Heegner
divisor $H_1$ equals the Borcherds-weight scalar
$c_1(0)/2$ of the singular-theta kernel. No independent residue
computation is performed in this chapter; the value $5$ enters
through the weight normalisation of Borcherds~1998 Theorem~13.3
(the divisor normalisation records only
$\operatorname{ord}_{H_1}(\Delta_5) = 1$).
\end{description}
Let $Y = (K3 \times E) \times \mathbb{R}^2$ and let $\Obs^{\mathrm{BV}}_{5\mathrm{d\,hCS}}(Y,\, \mathfrak{so}(4, 20))$ denote the Costello--Gwilliam--Li chain-level BV observables of $5$d holomorphic Chern--Simons on $Y$ with gauge Lie algebra the orthogonal Mukai preserver $\mathfrak{so}(4, 20)$ \textup{(}Kac type $D_{12}$\textup{)} determined by the Mukai-lattice signature $(4, 20)$ of $\widetilde{H}(K3, \mathbb{Z})$ (Remark~\ref{rem:k3e-mukai-signature-super-yangian}). Gauge-fix by the Costello--Li heat-kernel prescription, regularise the partition function through the Mittag--Leffler tower $\{P_{\epsilon, L}\}$, and descend \v{C}ech--Ran against a finite analytic polydisc cover of $K3 \times E$. Then the BV partition-function residue at the principal Heegner divisor $H_1 = \{z = 0\} \subset \mathbb{H}_2$ of the Siegel upper half-space satisfies
\[
 \underset{[Z] \in H_1}{\operatorname*{Res}}\,
 Z_{5\mathrm{d\,hCS}}\bigl( (K3 \times E) \times \mathbb{R}^2,\,
 \mathfrak{so}(4, 20) \bigr)([Z])
 \;=\; \left.c_N(0)/2\right|_{N=1} \;=\; 5
 \;=\; \kBKM(\Delta_5),
\]
with $c_1(0) = 10$ the zero Fourier coefficient of the weak Jacobi form $\phi_{0,1}$. The residue is computed from three inputs: (i) the tree-level BCOV anomaly $\alpha_{\mathrm{BCOV}}(K3 \times E) = (\chi(K3 \times E)/24) \cdot \gamma_{K3 \times E} = 0$, forced by the K\"unneth vanishing $\chi(K3 \times E) = \chi(K3) \cdot \chi(E) = 24 \cdot 0 = 0$; (ii) the one-loop wheel-graph sum localising to the Borcherds singular-theta kernel on the Heegner divisor $H_1$; in the divisor normalisation $\operatorname{div}(\Delta_5)=H_1+2H_4$ one has $\operatorname{ord}_{H_1}(\Delta_5)=1$, while the scalar returned by this path is the Borcherds-weight normalisation $\kBKM(\Delta_5)=\left.c_N(0)/2\right|_{N=1}=5$ (Borcherds~1998 Theorem~13.3; Gritsenko~1999 Theorem~1.2); (iii) higher-loop bulk corrections $O(\hbar^2)$ regular on $H_1$ by Kontsevich--Tamarkin formality on the $\mathbb{C}^2$-factor, contributing zero residue.
\end{theorem}

\begin{proof}[Conditional proof]
Under the stated BV/Ran comparison hypotheses and (H-res), the global BV action on $Y = (K3 \times E) \times \mathbb{R}^2$ is assembled via \v{C}ech--Ran descent against a finite analytic polydisc cover of $K3 \times E$: on each chart $U_\alpha \cong \mathbb{C}^3$ the action restricts to the standard Costello--Li 5d hCS action on $\mathbb{C}^2 \times \mathbb{R}_t$ with gauge Lie algebra the orthogonal Mukai preserver $\mathfrak{so}(4, 20)$; transition Fourier--Mukai kernels $K_{\alpha\beta}$ on pairwise overlaps glue the chart-wise BV actions, and the triple-overlap cocycle condition holds up to coherent homotopy by Bondal--Orlov reconstruction. The cubic Casimir $\operatorname{tr}_{\mathfrak{so}(4, 20)}(T^a T^b T^c + T^a T^c T^b) = 0$ on the $\mathfrak{so}(24)$-restricted symmetric pair, an orthogonal type-$D$ vanishing rather than a general simply-laced principle, is the one-loop anomaly-cancellation input.

The tree-level BCOV-anomaly vanishing follows from the Atiyah-class K\"unneth decomposition $\mathrm{At}(T_{K3 \times E}) = \pi_1^* \mathrm{At}(T_{K3}) \oplus \pi_2^* \mathrm{At}(T_E)$ paired with the coefficient $\chi(K3 \times E)/24 = 0$. The partition function factors as $Z_{5\mathrm{d\,hCS}}[Z] = \exp(\hbar \cdot Z^{\mathrm{1-loop}}_{\mathrm{Heegner}}[Z]) \cdot (1 + O(\hbar^2))$, with trivial tree-level.

The one-loop residue is computed from the Costello--Gwilliam holomorphic-topological tensor theorem (arXiv:2210.13036 Vol.~2 \S4.10), which rewrites the wheel-$n$ graph weight on $Y$ as a regularised Borcherds singular-theta integral
\[
 W^{K3 \times E}_{\mathrm{wheel}_n}(P)
 \;=\; \int^{\mathrm{reg}}_{\mathbb{H}} W^{\mathrm{top}}_{\mathrm{wheel}_n}(\tau) \cdot W^{\mathrm{hol}}_{\mathrm{wheel}_n}(Z)\big|_{K3 \times E} \cdot \frac{d\tau \wedge d\bar\tau}{y^2}.
\]
At the Heegner locus $H_1 = \{z = 0\}$, Gritsenko--Nikulin gives $\operatorname{ord}_{H_1}(\Delta_5)=1$ in the normalisation $\operatorname{div}(\Delta_5)=H_1+2H_4$. The Borcherds weight formula gives the independent scalar $\kBKM(\Delta_5)=c_{\phi_{0,1}}(0)/2=5$, with $c_{\phi_{0,1}}(0) = 10$ the zero Fourier coefficient (Borcherds~1998 Theorem~13.3; Gritsenko~1999 Theorem~1.2). The residue in the preceding display is this weight normalisation, not a vanishing order.

The Mittag--Leffler tower of gauge-fixed propagators is assumed to commute with the \v{C}ech--Ran descent: the chart-wise Costello--Li heat kernel is compatible with the transition Fourier--Mukai cocycle up to coherent homotopy (Bondal--Orlov~2001 Compositio~125), preserving the Heegner-locus residue through the global stitching. Higher-loop corrections $Z^{(n)}_{\mathrm{bulk}}[Z]$ for $n \geq 2$ are Kontsevich-formality bulk integrals with no pole on $H_1$; their residue at the Heegner locus vanishes by the modular-invariance of the bulk polynomial sector. Under these comparison and regularity hypotheses together with (H-res), the one-loop residue is exact and equals $5$; without (H-res), what is established is only that the residue is a scalar multiple of the singular-theta kernel datum on $H_1$, with the value $5$ supplied by the weight normalisation rather than derived.

Three primary-source derivations of the Borcherds weight agree:
Gritsenko--Nikulin~1995 alg-geom/9504006 Part~II \S2 computes the
Fourier expansion of \(\Delta_5\); Borcherds~1998
\emph{Invent.~Math.}~132 Theorem~13.3 gives the singular-theta weight
formula; Lorgat~2020 arXiv:2004.09030 Theorem~3 gives the explicit
Borcherds product. All three yield \(\kappa_{\mathrm{BKM}}(\Delta_5)=5\).
\end{proof}

\begin{remark}[Three chain-level constructions of $\kBKM(\Delta_5) = 5$]
\label{rem:bl-bv-k3e-path2}
Theorem~\ref{thm:bl-bv-k3e-kappa5} conditionally realises $\kBKM(\Delta_5) = 5$ as the chain-level BV partition-function residue at the Humbert divisor $H_1$. Two other chain-level constructions compute the same invariant through disjoint technical routes: the Getzler--Kapranov modular-operadic supertrace on the $E_3$-bar complex $B^{E_3}(\PhiFA_3(\mathcal{C}_N))$ at $\mathcal{C}_1 = D^b(\Coh(K3 \times E))$, and the Getzler--Kapranov Feynman-weight sum on $\overline{\mathcal{M}}_{g, n}$ evaluated at the $K3 \times E$ topological field theory target. The three constructions converge at $\left.c_N(0)/2\right|_{N=1} = 5$ without overlapping technical content: modular-operadic supertrace, BV partition-function residue at a Heegner divisor, Feynman-weight sum on the Deligne--Mumford compactification. Their agreement is a chain-level refinement of the Gritsenko--Nikulin--Borcherds weight theorem, not a tautology. The extension to $N \in \{2, 3, 4, 6\}$ along Scope~(A) of Theorem~\ref{thm:k3e-universal-kBKM-two-scopes} is a further conditional CHL statement requiring the equivariant BV/Ran descent and hCS-to-Hall comparison on $(K3 \times E)/(\mathbb{Z}/N\mathbb{Z})$ with Mukai lattice replaced by the $g_N$-fixed sublattice $\Lambda^{g_N}_{\mathrm{Muk}}$.
\end{remark}

\begin{remark}[Heegner locus versus Kodaira $I_1$ residues]
\label{rem:bl-heegner-vs-kodaira-residues}
The chain-level BV residue at the Heegner locus $H_1 \subset \mathbb{H}_2$ is \emph{not} the same invariant as the $24$-fold Kodaira $I_1$ nodal-fibre residue of Proposition~\ref{prop:k3e-one-loop-origin-chapter-header}. The Kodaira $I_1$ count sums to $\chi_{\mathrm{top}}(K3) = 24 = \kappa_{\mathrm{fiber}}(K3)$, the fourth entry of the four-$\kappa_\bullet$ spectrum (Remark~\ref{rem:k3e-four-kappa-spectrum}); the Heegner-locus BV residue yields $\kBKM(\Delta_5) = 5$, the third entry. The two residues are associated with different geometric loci (Kodaira fibres of the elliptic K3 $\pi \colon K3 \to \mathbb{P}^1$ versus Heegner divisors on the Siegel upper half-space $\mathbb{H}_2$) and compute different invariants. Proposition~\ref{prop:k3e-one-loop-origin-chapter-header}(iii) records the Kodaira residues as one-loop anomaly input while the scalar $5$ is fixed by $\left.c_N(0)/2\right|_{N=1}$; it does not identify $\kappa_{\mathrm{BKM}}$ with a sum of the K3 fibre value and a Heisenberg or residue contribution.
\end{remark}

\begin{conjecture}[Lorgat~2020 Conjecture~1, two-scope sharpening]
\label{conj:k3e-lorgat-conj1-two-scope}
Each specimen $\Phi_N$ of Theorem~\ref{thm:k3e-universal-kBKM-two-scopes} is the reciprocal square root of the reduced Donaldson--Thomas partition function of a Calabi--Yau threefold $X_N$:
\[
 Z^{X_N}_{\mathrm{DT}, \mathrm{red}} \;=\; C_N \cdot \Phi_N(Z)^{-2},
\]
with $X_N$ a CHL orbifold of $K3 \times E$ at CHL levels $N \in \{1, 2, 3, 4, 6\}$, and a non-CHL host at $N \in \{5, 7, 8\}$ (Borcea--Voisin candidate at $N = 5$; Mongardi--Tari--Wandel Kummer-$3$ fourfold at $N = 8$; $N = 7$ obstructed at the CY$_3$ level by Nikulin fixed-count constraints). The $N = 1$ case $Z^{K3 \times E}_{\mathrm{DT}, \mathrm{red}} = C \cdot \Delta_5(Z)^{-2}$ is the theorem of Oberdieck--Pandharipande~\cite{OberdieckPandharipande2018}; the remaining $N$ are conditional on the orbifold DT theory of Bryan--Oberdieck--Pandharipande--Yin~\cite{BOPY2018}.

\smallskip
\noindent\emph{Chiral-bialgebra sharpening on~$E$.} At every CHL level
$N$, the oriented hCS--Hall comparison, the Stage-$2$ witnessed
cycle, the Hall--Drinfeld double, and the coproduct/associator/$R$-
matrix comparison identify the Stage-$2$ specialisation
$\SpCh_{K3/(\ZZ/N\ZZ),\, E}(\PhiFA_3(K3 \times E))$ with the chiral
bialgebra on~$E$
\[
 \mathbf{H}_{\Delta_5}^{(N)} := U^{\chir}(\fg_{\Phi_N}),
\]
the chiral-Hall--Drinfeld vertex envelope of the paramodular BKM
attached to~$\Phi_N$. Three arithmetic invariants pin the identified
object:
\begin{description}
 \item[Character.] The graded vacuum character reads
 \[
 \operatorname{ch}\mathbf{H}_{\Delta_5}^{(N)}(Z) \;=\; \Phi_N(Z)^{-2},
 \]
 specialising at $N = 1$ to the Lorgat 2020 reciprocal-square character $\operatorname{ch}\mathbf{H}_{\Delta_5}(Z) = \Delta_5(Z)^{-2}$ as a meromorphic Siegel modular form of weight~$-10$ on $\mathrm{Sp}_4(\mathbb{Z})$ (Remark~\ref{rem:k3e-gritsenko-nikulin-denominator-WKB-character}; Lorgat 2020 Theorems~$2$--$4$).
 \item[Signed root character.] For each positive root $\alpha$ of $\fg_{\Phi_N}$, the signed root-character coefficient
 \(\operatorname{sdim}\fg_{\Phi_N,\alpha}
 = c^{(g_N)}(|\alpha|^2/2,\, \ell(\alpha))\)
 is the Fourier coefficient of the $g_N$-twisted-twined K3 elliptic genus
 \[
 \phi_{0,1}^{(g_N, h_M)}(\tau, z) \;=\; \frac{1}{N}\,\mathrm{tr}_{\mathcal{H}_{K3}^{(g_N)}}\bigl(h_M\,(-1)^{F_L + F_R}\,q^{L_0}\bar{q}^{\bar{L}_0}\,y^{J_0}\bigr),
 \]
 read at the $(N, M)$-paramodular index, with $(g_N, h_M) \in M_{24} \times M_{24}$ the Enriques-commuting pair realising Bryan--Oberdieck 2017 $(N, M)$-CHL quotient.
 \item[CY-D witness.] The modular characteristic is $\kappa_{\mathrm{BKM}}(\Phi_N) = c_N(0)/2 \in \{5, 3, 2, \tfrac32, 1\}$ on the CHL slice, Künneth-invariant of the specialisation cycle $(K3/(\ZZ/N\ZZ), E) \subset X_N \times X_N$ and forced by the Borcherds weight theorem (Gritsenko~$1999$ Theorem~$1.2$).
\end{description}
In the $N = 1$ case the comparison identifies
$\mathbf{H}_{\Delta_5}$ with the Hall--Drinfeld double of the Manin pair
$(\fg_{\Delta_5},\, \fn_+^{\mathrm{imag}} \oplus
\fh^{\mathrm{imag, rk\,23}})$ studied in
Chapter~\ref{ch:k3-chiral-bialgebra-platonic}. The
classification-invariant identity
$H^2(\fg_{\Delta_5})^{\ZZ/2,\,K(1)} \cong \CC \cdot \Delta_5$
(Bruinier~$2002$ Proposition~$5.1$) classifies the unique
paramodular gauge-equivalence line, and the Hall--Drinfeld double
quantisation is the corresponding
super-Etingof--Kazhdan object.
\end{conjecture}

\begin{remark}[The two scopes are not the same atlas]
\label{rem:two-scopes-different-atlases}
The CHL slice $\{1,2,3,4,6\}$, the order-indexed Nikulin extension
$\{1,\ldots,8\}$, and the Gritsenko--Cl\'ery triple atlas are three
different indexing systems. The CHL slice and the Nikulin extension
agree on $N\in\{1,2,3,4,6\}$ and the latter adds proved Borcherds weights
$(2,\tfrac32,1)$ at $N=(5,7,8)$ in the Mathieu-twined normalisation, while the corresponding CY$_3$ hosts require
separate construction. The
Gritsenko--Cl\'ery atlas is instead indexed by triples $(t,N;k)$ and
agrees with the order-indexed tower as an indexing system only at $\Delta_5=(1,1;5)$; for
example the CHL value at order $2$ is $3$, while the GC row
$(2,1;2)$ has weight $2$, and the numerical agreement of the row
$(1,2;3)$ with the order-$2$ value is a coincidence of weights until a
normalisation theorem identifies the forms. The
bi-parameter structure $(N,M)$ of Bryan--Oberdieck CHL quotients refines
the single-order reading; substituting a GC paramodular index for a CHL
orbifold order changes the object.
\end{remark}

\begin{conjecture}[Eichler--Zagier decomposition of the twisted-twined input at \texorpdfstring{$N = 2$}{N = 2}]
\label{conj:ez-n2-2A-decomposition}
For the Enriques-involution commuting pair $(g_{2A}, h_M = \mathrm{id})$, the twisted-twined K3 Jacobi input admits the Eichler--Zagier decomposition $F_L^{(g_{2A}, \mathrm{id})} = \alpha(\tau)\,\phi_{0,1} + \beta(\tau)\,\phi_{-2,1}$ with $\alpha(\tau) = 1/3$ constant and $\beta(\tau) = 5/3 + (2/3)\,E_2^{(2)}(\tau)$ Eisenstein-corrected on $\Gamma_0(2)$, as forced by coefficient matching against the $N = 2$ head $(c^{(2A)}(-1, 1), c^{(2A)}(0, 0), c^{(2A)}(1, -1), c^{(2A)}(1, 1)) = (2, 4, -6, 14)$. The refined universal identity splits the Borcherds weight as $\kBKM(\Phi_{N, M}) = 12\,\alpha(\infty) + \Delta_\beta$ with $\Delta_\beta$ the Hecke-minus correction from the $\phi_{-2, 1}$-coupled sector; at $h_M = \mathrm{id}$ the correction vanishes for $N \in \{2, 3, 4, 6\}$ and equals $-7$ at $N = 1$, absorbing the $\eta^{-6}$-factor of the Lorgat 2020 Borcherds product.
\end{conjecture}

\begin{remark}[Normalisation stamp at \texorpdfstring{$N = 2$}{N = 2}:
twined constant versus Gritsenko--Cl\'ery cusp-weighted weight]
\label{rem:ez-n2-normalisation-stamp}
The constants displayed in
Conjecture~\ref{conj:ez-n2-2A-decomposition} belong to the
Mathieu-twined family: $c^{2A}(0, 0) = 4$ is the zero Fourier
coefficient of the $\tfrac{1}{N}$-normalised twisted-twined input
$F_L^{(g_{2A}, \mathrm{id})}$, and its naive half
$c^{2A}(0, 0)/2 = 2$ is a twined-normalisation value. It is not the
square-root CHL weight at $N = 2$, which is $\kBKM(\Phi_2) = 3$
(Scope~(A) of Theorem~\ref{thm:k3e-universal-kBKM-two-scopes}): on
$\Gamma_0(2)$ the weight-formula input is not the single constant
term but the cusp-weighted sum
$\tfrac{1}{2}\sum_{f/e} (h_e/N_e)\, c_{f/e}(0, 0)$ of
Gritsenko--Cl\'ery 2008 Theorem~$3.1$, taken over both cusps of
$\Gamma_0(2)$, which returns $3 = \mathrm{wt}(\nabla_3)$. Every
displayed constant in this subsection is therefore read with its
ladder named: twined ($c^{2A}(0, 0) = 4$, naive half $2$) versus
square-root CHL ($c_2(0) = 6$, weight $3$).
\end{remark}

\begin{proposition}[Non-CHL Nikulin orders and Gritsenko--Cl\'ery half-integer rows]
\label{conj:gc-half-integer}
There are two different non-CHL boundaries, and they are distinct.
\begin{enumerate}[label=\textup{(\roman*)}]
\item The Nikulin-order extension of the Borcherds-weight identity has
orders $N\in\{5,7,8\}$ with
\[
 (c_N(0),\kBKM(\Phi_N))_{N=5,7,8}=(4,3,2;\ 2,\tfrac32,1),
\]
in the Mathieu-twined normalisation, the constants being the
frame-shape ones-exponents of $1^45^4$, $1^37^3$, $1^22\,4\,8^2$.
These are Borcherds weights computed by the singular-theta lift; the
order-$7$ entry is half-integral and carries a multiplier system;
their CY$_3$ hosts are
outside the smooth $K3\times E$ CHL quotient class.
\item The half-integral rows of the Gritsenko--Cl\'ery dd-modular
catalogue are the triples $(t,N;k)=(4,1;\tfrac12)$ and
$(1,4;\tfrac32)$. They are multiplier-system Siegel modular forms
indexed by dd-modular type,
not by the Nikulin orders $5,7,8$.
\item Consequently no Gritsenko--Cl\'ery dd-modular row in this chapter
has weight $1/4$ or weight $0$, and no row carries the tuple
$(1,\tfrac12,0)$ as a Borcherds-weight extension.
\end{enumerate}
\end{proposition}

\begin{proof}
Item~\textup{(i)} is Corollary~\ref{cor:borcherds-weight-full-8form-scope}
in Chapter~\ref{ch:cy-d-kappa-stratification}: Gritsenko's
order-indexed extension gives $(2,\tfrac32,1)$ at $N=(5,7,8)$, with
constant terms $(4,3,2)$, the frame-shape ones-exponents. Item~\textup{(ii)} is the
Gritsenko--Cl\'ery triple classification recorded in
Theorem~\ref{thm:k3e-universal-kBKM-two-scopes} and later in
Theorem~\ref{thm:k3e-gc-eight-commuting-pair}. Since the two indexing
sets are different, substituting the half-integral GC rows for the
Nikulin orders conflates two distinct structures. The absence of weights $0$ and
$1/4$ follows from the eight displayed triples.
\end{proof}

\begin{conjecture}[Lorgat 2020 Conjecture~1 at non-CHL \texorpdfstring{$N = 5$}{N = 5}: weight-$2$ host]
\label{conj:lorgat-conj1-n5-metaplectic}
The $N = 5$ entry of the Lorgat 2020 Conjecture~1 ladder is populated
not by a $K3 \times E$ CHL orbifold; the condition
$\varphi(N)\mid 2$ fails at $N=5$; but by a non-CHL host, with the
Borcea--Voisin CY$_3$ as the candidate geometry. The Borcherds
denominator attached to the order-$5$ Nikulin twining has
\[
 \kBKM(\Phi_5)=c_5(0)/2=4/2=2.
\]
The conjectural Donaldson--Thomas statement is therefore a reciprocal
square-root denominator statement of weight $2$, not a reciprocal
fourth power of a half-integral Gritsenko--Cl\'ery row. Extending the
construction from the CHL five entries to the non-CHL Nikulin orders
requires (i) a K3-fibred CY$_3$ host that does not demand an order-$5$
automorphism of the elliptic factor, (ii) the Borcherds singular-theta
weight on the corresponding order-$5$ lattice, and (iii) a
$\kappa_{\mathrm{cat}}$ computation for the non-CHL orbifold after the
K\"unneth product argument no longer applies verbatim.
\end{conjecture}

\section{CoHA structure and the BKM superalgebra}
\label{sec:k3e-coha-structure}

The Borcherds side supplies the positive-half target
$\mathfrak n_+(\mathfrak{g}_{\Delta_5})$. Its denominator product gives
signed Euler, or superdimension, exponents \(c(D)\). Ordinary
\(\mathbb Z/2\)-graded root-space dimensions require a parity fixture or
a finite target-presentation computation. Recovering that target from
the compact critical CoHA of \(K3 \times E\) is the comparison problem:
the finite Borcherds core and the Hall--Drinfeld criterion below state
the recognition data needed for the positive-half identification.

\begin{theorem}[Borcherds positive-half target]
\label{thm:k3e-coha}
The Borcherds positive-half target is
\[
 U(\fn_+(\mathfrak{g}_{\Delta_5})),
\]
where $\fn_+ = \bigoplus_{\alpha \in \Delta_+} \mathfrak{g}_\alpha$ is
the positive nilpotent subalgebra, with signed root character governed
by $\phi_{0,1}$:
\[
\operatorname{sdim}\mathfrak g_\alpha
 =
\dim\mathfrak g_{\alpha,\bar0}-\dim\mathfrak g_{\alpha,\bar1}
 =
c(D(\alpha)),\qquad D(\alpha)=4nm-l^2.
\]
The equality \(\dim\mathfrak g_\alpha=|c(D(\alpha))|\) is not read from
the denominator character alone; it is available only after fixing the
target parity or computing the finite-height target presentation. No
compact critical-CoHA identification is asserted in this theorem.
\end{theorem}

\begin{problem}[Compact Hall-to-Borcherds primitive recognition]
\label{op:k3e-compact-coha-hall-borcherds}
The finite reduced compact Hall windows are constructed in
Theorem~\ref{thm:k3e-reduced-compact-source-finite-hall-windows}; the
finite radical-quotient Hall--Drinfeld doubles are constructed in
Theorem~\ref{thm:k3e-finite-compact-hall-drinfeld-double}; the
compact-provenance source packets
\(\mathsf M_H^{X,\mathrm{can}}\) are constructed in
Theorem~\ref{thm:k3e-canonical-finite-source-packets}; the
Serre--Borcherds target packets
\(\mathsf M_H^{\Delta,\mathrm{can}}\) are constructed in
Theorem~\ref{thm:k3e-canonical-finite-target-packets}; and the
blockwise comparison matrices are constructed, after the stated
charge, parity, and basis choices, in
Theorem~\ref{thm:k3e-finite-comparison-matrix-packets}. The remaining
problem is not the existence of source matrices, target matrices, or
candidate comparison matrices, but their structural recognition as the
\(\Delta_5\) matrices. For every height \(H\), prove that the primitive
BPS quotient of
\(\overline{\mathcal H}^{+,\leq H}_{K3\times E,\sigma}\) admits a
charge-preserving comparison with
\(\fn_+(\mathfrak g_{\Delta_5}^{\leq H})\) preserving signed
character, parity, Hall commutator, Borcherds--Serre kernel,
Euler--Hall pairing, coproduct, centre, associator class, and the
height transition maps. Equivalently, construct a continuous map
\[
 \Psi^+_{\Hall\to\BKM}\colon
 \varprojlim_H\overline{\mathcal H}^{+,\leq H}_{K3\times E,\sigma}
 \longrightarrow
 \widehat{U(\fn_+(\mathfrak g_{\Delta_5}))}
\]
whose finite-height associated graded maps are algebra isomorphisms.
\end{problem}

\begin{theorem}[Positive-half Hall--Borcherds comparison criterion]
\label{thm:k3e-positive-half-hall-borcherds-criterion}
Let \(\overline{\mathcal H}^{+,\leq H}_{K3\times E,\sigma}\) be the
radical-quotient positive half of the finite compact Hall--Drinfeld
double of
Theorem~\ref{thm:k3e-finite-compact-hall-drinfeld-double}. A compatible
family of maps
\[
 \Psi^{+,\leq H}_{\Hall\to\BKM}\colon
 \overline{\mathcal H}^{+,\leq H}_{K3\times E,\sigma}
 \longrightarrow
 U(\fn_+(\mathfrak g_{\Delta_5}^{\leq H}))
\]
identifies the completed compact Hall positive half with the Borcherds
positive half if and only if the following finite-height recognition
conditions hold for every \(H\) and are compatible under the
downward-saturated transitions.
\begin{enumerate}[label=\textup{(\roman*)}]
\item The Hall charge monoid is identified with the positive
Borcherds cone \(\Delta_+\subset\Lambda_{\mathrm{II}}^{2,1}\), with
the same height filtration and parity convention.
\item The primitive Hall BPS space in charge \(\alpha\) has
Euler/superdimension target \(c(D(\alpha))\), with the same imprimitive
divisor correction as the Borcherds product. Its ordinary dimension is
not asserted to be \(|c(D(\alpha))|\) unless the finite-height parity
fixture or target presentation computes the even and odd parts.
\item The associated graded Hall commutator is the Borcherds bracket:
real roots satisfy the ordinary Serre relations, orthogonal imaginary
roots super-commute, and no extra primitive relation survives in any
finite-height quotient.
\item The orientation local systems, Tate shifts, and cone completion
used on the Hall side match the signs and topology of
\(\widehat{U(\fn_+(\mathfrak g_{\Delta_5}))}\).
\item For every height \(H\) the compact source carries finite
Hall--Borcherds recognition data
\[
(\mathrm{src}_H,M_H,D_H,B_H,G_H,K_H,Q_H,A_H):
\]
source provenance; the charge monoid \(M_H\); the source signed
character and target denominator character in matched charge blocks;
the Hall and Borcherds bracket matrices; the grading and parity
fixture; the Serre kernel; the Hall pairing/coproduct quotient; and the
associator differential, admissible gauge rows, and completion class. The
source side of this packet is
\(\mathsf M_H^{X,\mathrm{can}}\) from
Theorem~\ref{thm:k3e-canonical-finite-source-packets}; the target side
is \(\mathsf M_H^{\Delta,\mathrm{can}}\) from
Theorem~\ref{thm:k3e-canonical-finite-target-packets}; and the
comparison matrices \(A_H\) are the finite block matrices of
Theorem~\ref{thm:k3e-finite-comparison-matrix-packets}. Recognition
requires, beyond the existence of those packets, that these data satisfy
the no-extra condition for primitive generators, primitive relations,
and central classes, the PBW condition on associated gradeds, and the
Mittag--Leffler compatibility of the finite-height transition maps.
\end{enumerate}
The theorem uses the compact source of
Theorem~\ref{thm:k3e-finite-compact-hall-drinfeld-double}. It does not
compute the primitive BPS comparison or prove the six finite
\(\Delta_5\)-recognition defects of
Theorem~\ref{thm:k3e-constructed-finite-double-recognition} vanish.
\end{theorem}

\begin{proof}
Work at finite height \(H\). If
\(\Psi^{+,\leq H}_{\Hall\to\BKM}\) is an isomorphism and the
isomorphisms commute with the height transitions, then they preserve
charge, signed primitive character, parity, the associated graded
commutator, the recognition data, and the completion topology, giving
\textup{(i)}--\textup{(v)}. Conversely,
\textup{(i)} and \textup{(ii)} identify the signed primitive character.
Conditions~\textup{(iii)} and \textup{(iv)} identify the Borcherds
bracket and the orientation/parity signs. Condition~\textup{(v)} is the
recognition step: the no-extra and PBW clauses identify the finite-height
defining ideal with the Borcherds--Serre ideal, while the
Mittag--Leffler transition clause makes the finite-height isomorphisms
compatible with the inverse system. Passing to the Borcherds-cone
completion gives the continuous map in
Problem~\ref{op:k3e-compact-coha-hall-borcherds}. The proof uses only
the stated input data.
\end{proof}

\begin{definition}[Finite compact Hall--Drinfeld double of a reduced window]
\label{def:k3e-finite-compact-hall-drinfeld-double}
Fix a stability chamber \(\sigma\), a downward-saturated finite
Borcherds-height window \(W_H\), and the reduced compact Hall window
\[
 \mathcal H^{+,\leq H}_{X,\sigma}
 :=
 \CoHA_{\sigma,W_H}^{\mathrm{red}}(K3\times E)
\]
of Theorem~\ref{thm:k3e-reduced-compact-source-finite-hall-windows}.
Let \(\sigma^\vee\) be the Serre-dual opposite chamber and set
\[
 \mathcal H^{-,\leq H}_{X,\sigma}
 :=
 \left(\CoHA_{\sigma^\vee,W_H}^{\mathrm{red}}(K3\times E)\right)^\vee_{\mathrm{cont}},
\]
with charge \(-\alpha\) dual to charge \(\alpha\). For classes
\(\alpha,\beta\in W_H\), define the finite Euler--Hall pairing by
\[
 \langle a,b\rangle_H
 =
 \delta_{\alpha+\beta,0}
 \int_{\mathfrak M_{\sigma,\alpha}^{\mathrm{red}}
 \times
 \mathfrak M_{\sigma^\vee,-\alpha}^{\mathrm{red}}}
 (a\boxtimes b)\,
 e_\hbar\!\left(
 R\!\operatorname{Hom}_X(\mathcal E,\mathcal F)
 \right)^{-1},
\]
where \(\mathcal E,\mathcal F\) are the universal objects on the two
factors, the inverse Euler class is expanded in the Rees parameter, and
the orientation line is the one used in the reduced Hall product. Put
\[
 \mathfrak r_H^\pm=\operatorname{rad}\langle-,-\rangle_H\cap
 \mathcal H^{\pm,\leq H}_{X,\sigma},\qquad
 \overline{\mathcal H}^{\pm,\leq H}_{X,\sigma}
 =
 \mathcal H^{\pm,\leq H}_{X,\sigma}/\mathfrak r_H^\pm .
\]
The finite Cartan part is the completed group algebra
\[
 \mathcal H^{0,\leq H}_{X,\sigma}
 =
 \mathbb C((\hbar))[\Lambda_{W_H}]
 /\mathfrak r^{0,\leq H},
\]
where \(\Lambda_{W_H}\subset \Lambda_{\mathrm{II}}^{2,1}\) is the
charge lattice generated by \(W_H\) and \(\mathfrak r^{0,\leq H}\) is
the Cartan radical of the induced Mukai pairing. The finite reduced
compact Hall--Drinfeld double is
\[
 D^{\mathrm{red},\leq H}_{\hbar}(K3\times E,\sigma)
 =
 D_\hbar\!\left(
 \overline{\mathcal H}^{+,\leq H}_{X,\sigma},
 \mathcal H^{0,\leq H}_{X,\sigma},
 \overline{\mathcal H}^{-,\leq H}_{X,\sigma},
 \langle-,-\rangle_H
 \right).
\]
\end{definition}

\begin{theorem}[Finite compact Hall--Drinfeld double construction]
\label{thm:k3e-finite-compact-hall-drinfeld-double}
For every finite height \(H\), the object
\[
 D^{\mathrm{red},\leq H}_{\hbar}(K3\times E,\sigma)
\]
of Definition~\ref{def:k3e-finite-compact-hall-drinfeld-double} is a
well-defined finite-height topological quasi-Hopf superalgebra over
\(\mathbb C((\hbar))\). Its positive half is the reduced compact Hall
window modulo the radical of the Euler--Hall pairing; its negative half
is the Serre-dual continuous dual; its Cartan part is the Mukai-lattice
charge torus modulo the Cartan radical. The transition
\(W_{H+1}\to W_H\) induces a morphism of quasi-Hopf superalgebras
\[
 D^{\mathrm{red},\leq H+1}_{\hbar}(K3\times E,\sigma)
 \longrightarrow
 D^{\mathrm{red},\leq H}_{\hbar}(K3\times E,\sigma)
\]
whenever the finite Hall windows are chosen cofinally and
downward-saturated.
\end{theorem}

\begin{proof}
The reduced compact Hall windows are finite type by construction of
\(W_H\). Their products are the Joyce--Song/Kontsevich--Soibelman
extension correspondences with Thom--Sebastiani orientation, so each
finite window is a connected charge-graded bialgebra after projection
from the bounded enlargement used for products and coproducts. The
positive window itself is not asserted to be Hopf.  The antipode is a
structure on the doubled object obtained after adjoining the Cartan
torus and the Serre-dual negative half.

Serre duality on the Calabi--Yau threefold identifies
\[
 R\!\operatorname{Hom}_X(\mathcal E,\mathcal F)^\vee
 \simeq
 R\!\operatorname{Hom}_X(\mathcal F,\mathcal E)[3],
\]
so the displayed Euler--Hall integral pairs the chamber \(\sigma\) with
the opposite chamber \(\sigma^\vee\). The Hall correspondence is proper
on \(W_H\), the Euler class is a Rees-localised finite expression in the
bounded window, and the orientation line is the same determinant square
root used for the Hall product. Thus \(\langle-,-\rangle_H\) is a
Hopf pairing:
\[
 \langle ab,c\rangle_H
 =
 \langle a\otimes b,\Delta c\rangle_H,\qquad
 \langle a,bc\rangle_H
 =
 \langle\Delta a,b\otimes c\rangle_H .
\]
The radical of a Hopf pairing is a biideal on both sides, hence the
quotients \(\overline{\mathcal H}^{\pm,\leq H}_{X,\sigma}\) are paired
bialgebra halves with nondegenerate pairing. Insert the Cartan torus
between the two halves and impose the Drinfeld cross relation
\[
 b^- b^+
 =
 \sum
 (-1)^{|b^+_{(1)}||b^-_{(2)}|}
 \langle b^+_{(1)},b^-_{(1)}\rangle_H\,
 b^+_{(2)}b^-_{(2)}
 \langle b^+_{(3)},b^-_{(3)}\rangle_H^{-1},
\]
with Sweedler notation for the finite Hall coproducts. Together with
the Cartan conjugation by Mukai charge, this is the Drinfeld double of
the paired finite Borel data. Its antipode exchanges the two Borel
halves and inverts the Cartan group-like elements. If the finite Hall
associator is non-trivial, the result is quasi-Hopf with that
associator; in the strict associator stratum it is Hopf. Cofinal
downward-saturated windows preserve charge, products, coproducts,
Serre duality, orientation, and the Euler--Hall kernel, so the
	restriction \(W_{H+1}\to W_H\) is a quasi-Hopf morphism.
\end{proof}

\begin{definition}[Finite quasi-Hopf associator complex]
\label{def:k3e-finite-qh-associator-complex}
Fix \(H\) and the finite compact double
\[
 D_H^X=D^{\mathrm{red},\leq H}_{\hbar}(K3\times E,\sigma).
\]
For \(m=2,3,4\), let
\[
 C^{m,H}_{\mathrm{qh}}(D_H^X)
 \subset
 (D_H^X)^{\widehat\otimes m}
\]
be the finite-dimensional subspace of height-\(\leq H\),
charge-filtered, parity-preserving quasi-Hopf tensor cochains, written
in the same triangular decomposition and Rees normalisation as
Definition~\ref{def:k3e-finite-compact-hall-drinfeld-double}. The
degree-two differential
\[
 d^{2,H}_{\mathrm{qh}}\colon
 C^{2,H}_{\mathrm{qh}}(D_H^X)\longrightarrow
 C^{3,H}_{\mathrm{qh}}(D_H^X)
\]
is the finite linearised quasi-Hopf gauge action, projected back to the
height window. The degree-three differential
\[
 d^{3,H}_{\mathrm{qh}}\colon
 C^{3,H}_{\mathrm{qh}}(D_H^X)\longrightarrow
 C^{4,H}_{\mathrm{qh}}(D_H^X)
\]
is the finite linearised pentagon operator. Let
\[
 C^{2,H}_{\mathrm{adm}}\subset C^{2,H}_{\mathrm{qh}}(D_H^X)
\]
be the subspace of gauge cochains preserving the charge filtration and
the signed parity convention. The finite quasi-Hopf identities give
\[
 d^{3,H}_{\mathrm{qh}}d^{2,H}_{\mathrm{qh}}=0
\]
after projection to the height window. After bases are chosen,
\(d^{3,X}_{\mathrm{qh},H}\) is the matrix of \(d^{3,H}_{\mathrm{qh}}\),
\(\Gamma_H^X\) is
a row basis for
\[
 \operatorname{im}\!\left(
 d^{2,H}_{\mathrm{qh}}\bigm|_{C^{2,H}_{\mathrm{adm}}}
 \right)
 \subset C^{3,H}_{\mathrm{qh}}(D_H^X),
\]
and \(\Pi_H^{\mathrm{adm},X}\) is a constraint matrix whose kernel in
\(C^{3,H}_{\mathrm{qh}}(D_H^X)\) is the charge/parity-admissible cochain
subspace. Thus
\[
 d^{3,X}_{\mathrm{qh},H}\Phi=0,\qquad
 \Phi-\Phi'\in\operatorname{rowspan}(\Gamma_H^X),\qquad
 d^{3,X}_{\mathrm{qh},H}(\Gamma_H^X)^t=0,\qquad
 \Pi_H^{\mathrm{adm},X}(\Gamma_H^X)^t=0
\]
is exactly the finite statement that the associator cochains \(\Phi\)
and \(\Phi'\) define the same charge/parity-admissible quasi-Hopf
associator class at height \(H\).
\end{definition}

\begin{theorem}[Canonical finite compact source packets]
\label{thm:k3e-canonical-finite-source-packets}
For every finite height \(H\) and stability chamber \(\sigma\), the
finite compact Hall--Drinfeld double
\[
 D_H^X=
 D^{\mathrm{red},\leq H}_{\hbar}(K3\times E,\sigma)
\]
determines a compact-provenance source packet
\[
 \mathsf M_H^{X,\mathrm{can}}
 =
 \bigl(
 \mathrm{src}_H^X,\,
 M_H^X,\,
 P_H^X,\,
 B_H^X,\,
 G_H^X,\,
 K_H^X,\,
 Q_H^X,\,
 \Delta_H^X,\,
 Z_H^X,\,
 \Phi_H^X,\,
 d^{3,X}_{\mathrm{qh},H},\,
 \Gamma_H^X,\,
 \Pi_H^{\mathrm{adm},X},\,
 T_{H+1,H}^X
 \bigr).
\]
Here:
\begin{enumerate}[label=\textup{(\roman*)}]
\item \(\mathrm{src}_H^X\) is the list of compact
vanishing-cycle classes in the reduced window
\(\CoHA_{\sigma,W_H}^{\mathrm{red}}(K3\times E)\), with the
orientation line and Tate/Rees normalisation inherited from the compact
Hall product;
\item \(M_H^X\) is the retained charge monoid
\(\Lambda_{W_H}\cap \Delta_+\), with parity splitting on every charge
block;
\item \(P_H^X\) is a homogeneous compact-source basis of the primitive
quotient of \(\overline{\mathcal H}^{+,\leq H}_{X,\sigma}\), together
with the Serre-dual basis in the negative half;
\item \(B_H^X,G_H^X,K_H^X,Q_H^X\) are respectively the primitive Hall
commutator matrices, Euler--Hall pairing matrices, radical kernels, and
source-radical quotient maps in those bases;
\item \(\Delta_H^X\), \(Z_H^X\), and \(\Phi_H^X\) are respectively the
finite Hall coproduct matrices, the zero-charge primitive-centre
support matrix, and a finite Hall associator cochain for the
quasi-Hopf structure on \(D_H^X\);
\item \(d^{3,X}_{\mathrm{qh},H}\), \(\Gamma_H^X\), and
\(\Pi_H^{\mathrm{adm},X}\) are respectively the finite pentagon
differential on the height-\(H\) quasi-Hopf associator complex, a row
basis of degree-two gauge coboundaries, and the matrix imposing the
charge-filtration and signed-parity constraints on those gauge rows;
\item \(T_{H+1,H}^X\) is the transition matrix induced by restricting
the downward-saturated window \(W_{H+1}\) to \(W_H\).
\end{enumerate}
If the windows are chosen cofinally and downward-saturated, these
packets form a strict inverse system:
\[
 T_{H+1,H}^X\,\mathsf M_{H+1}^{X,\mathrm{can}}
 =
 \mathsf M_H^{X,\mathrm{can}}
\]
on all retained structure matrices after deleting the charges outside
\(W_H\). The construction is source-side only; it does not identify any
row of \(\mathsf M_H^{X,\mathrm{can}}\) with the corresponding
Serre--Borcherds \(\Delta_5\) row.
\end{theorem}

\begin{proof}
The reduced compact Hall window is finite-dimensional in every
retained charge and parity block. Choose homogeneous representatives
of the compact vanishing-cycle cohomology classes before quotienting by
the Euler--Hall radical, and choose compatible bases after quotienting.
Serre duality supplies the negative basis. The compact provenance is
therefore the triple consisting of the moduli support, the orientation
line, and the chosen vanishing-cycle class before passage to the
radical quotient.

All entries of the packet are the structure constants of already
constructed finite maps. The Hall product and its commutator give
\(B_H^X\) after projection to primitives. The Euler--Hall integral gives
\(G_H^X\); its kernel gives \(K_H^X\); the quotient map gives
\(Q_H^X\). The Hall coproduct of
Theorem~\ref{thm:k3e-finite-compact-hall-drinfeld-double} gives
\(\Delta_H^X\). The primitive zero-charge part of the finite double,
computed in the same basis, gives \(Z_H^X\). The finite associator of
the quasi-Hopf double gives \(\Phi_H^X\), with the strict Hopf stratum
corresponding to the trivial cochain. The charge-filtered quasi-Hopf
deformation complex of \(D_H^X\) is finite in the same window; its
degree-three differential gives \(d^{3,X}_{\mathrm{qh},H}\), its
degree-two coboundary rows give \(\Gamma_H^X\), and the requirement that
the gauge cochains preserve the charge filtration and signed parity
gives \(\Pi_H^{\mathrm{adm},X}\).

For \(H+1\to H\), downward saturation deletes only charges outside
\(W_H\). The Hall product, coproduct, Serre duality, Euler--Hall
pairing, radical quotient, centre, and associator cochain are obtained
by restricting the same compact correspondences and orientation lines;
the quasi-Hopf deformation differential, gauge-coboundary rows, and
admissibility constraints are obtained by restricting the same
charge-filtered finite deformation complex.
Thus the restriction matrices \(T_{H+1,H}^X\) commute with every
structure matrix in the packet. This proves the inverse-system
compatibility. No Borcherds comparison is used: the construction
records the compact source matrices, not their equality with target
matrices.
\end{proof}

\begin{theorem}[Finite Borcherds core: superdimensions, PBW, Serre, and primitive coproduct]
\label{thm:k3e-bkm-finite-core}
Let $\Delta_+$ be the positive root cone of $\mathfrak g_{\Delta_5}$ and
let $\fn_+ = \bigoplus_{\alpha\in\Delta_+}\mathfrak g_\alpha$. For every
height bound $H$, let
\[
 \fn_+^{\leq H}
 :=
 \fn_+\Big/\bigoplus_{\mathrm{ht}(\alpha)>H}\mathfrak g_\alpha
\]
be the finite nilpotent quotient. Four structures are proved on
this Borcherds side.
\begin{enumerate}[label=\textup{(\alph*)}]
\item \emph{Signed root character and parity.} Each root space is finite
dimensional and the denominator exponent records the signed character
\[
 \dim \mathfrak g_{\alpha,\bar 0}-\dim \mathfrak g_{\alpha,\bar 1}
 = c(D(\alpha))
\]
on primitive roots, with the Borcherds divisor sum supplying the
imprimitive correction. The ordinary dimension
\(\dim\mathfrak g_\alpha\) is determined only after the parity fixture
or an explicit finite-height target-presentation computation; under the
sign fixture of Proposition~\ref{prop:k3e-super-grading} it is
\(|c(D(\alpha))|\) on primitive roots.
\item \emph{PBW.} After any total order refining height, multiplication gives
a vector-space isomorphism
\[
 \bigotimes_{\alpha\in\Delta_+,\,\mathrm{ht}(\alpha)\leq H}
 S(\mathfrak g_{\alpha,\bar 0})
 \otimes
 \Lambda(\mathfrak g_{\alpha,\bar 1})
 \;\xrightarrow{\;\sim\;}\;
 \operatorname{gr} U(\fn_+^{\leq H}).
\]
\item \emph{Borcherds--Serre ideal.} The defining ideal is generated by the
real-root Serre relations
\[
 (\operatorname{ad} e_i)^{1-a_{ij}}e_j=0,\qquad
 (\operatorname{ad} f_i)^{1-a_{ij}}f_j=0
 \quad (a_{ii}=2,\ i\ne j),
\]
together with the Borcherds imaginary-root super-commutativity
$[e_\alpha,e_\beta]=[f_\alpha,f_\beta]=0$ whenever
$(\alpha,\beta)=0$. No Drinfeld $J$-presentation or Yangian coproduct is
used in this assertion.
\item \emph{Primitive enveloping coproduct.} The enveloping algebra carries
the standard cocommutative coproduct
\[
 \Delta_{\mathrm{env}}(x)=x\otimes 1+1\otimes x,\qquad x\in\fn_+^{\leq H}.
\]
This is the PBW-compatible coproduct of $U(\fn_+^{\leq H})$, not the
Drinfeld-new Hall coproduct and not the coproduct of a completed double.
\end{enumerate}
\end{theorem}

\begin{proof}
The signed-character statement is the Gritsenko--Nikulin denominator
theorem for $\Delta_5$ applied to the K3 weak Jacobi input
$\phi_{0,1}$, with the imprimitive correction given by the Borcherds
product divisor sum. The passage from the signed exponent to ordinary
root-space dimensions uses the declared parity fixture or the
finite-height presentation itself. The PBW statement is the
Poincare--Birkhoff--Witt theorem for Lie superalgebras applied to each
finite height quotient; odd root spaces produce exterior factors and
even root spaces produce symmetric factors.
The Serre ideal is Borcherds' definition of a generalized Kac--Moody
superalgebra: real simple roots carry the ordinary Kac--Moody Serre
relations, while orthogonal imaginary simple roots super-commute. The
coproduct is the universal enveloping coproduct, hence primitive on the
Lie generators and compatible with the PBW filtration.
\end{proof}

\begin{theorem}[Canonical finite Serre--Borcherds target packets]
\label{thm:k3e-canonical-finite-target-packets}
For every finite height \(H\), the Serre--Borcherds current quotient
\[
 Y_H^\Delta=
 Y^{\mathrm{SB},\leq H}_{\hbar}(\mathfrak g_{\Delta_5})
\]
determines a target packet
\[
 \mathsf M_H^{\Delta,\mathrm{can}}
 =
 \bigl(
 M_H^\Delta,\,
 P_H^\Delta,\,
 B_H^\Delta,\,
 G_H^\Delta,\,
 K_H^\Delta,\,
 Q_H^\Delta,\,
 \Delta_H^\Delta,\,
 Z_H^\Delta,\,
 \Phi_H^\Delta,\,
 d^{3,\Delta}_{\mathrm{qh},H},\,
 \Gamma_H^\Delta,\,
 \Pi_H^{\mathrm{adm},\Delta},\,
 T_{H+1,H}^\Delta
 \bigr).
\]
Here \(M_H^\Delta=\Delta_+\cap\{\operatorname{ht}\leq H\}\) with the
chosen parity fixture; \(P_H^\Delta\) is a parity-homogeneous basis of
the finite root spaces; \(B_H^\Delta\) records the
Borcherds--Serre bracket constants and Serre rows;
\(G_H^\Delta\), \(K_H^\Delta\), and \(Q_H^\Delta\) are the invariant
Borcherds form, its Cartan radical, and the quotient map;
\(\Delta_H^\Delta\) is the finite current coproduct; \(Z_H^\Delta\) is
the primitive centre, equal to the Cartan radical plus the declared
dynamical group-like parameters; and
\(\Phi_H^\Delta,d^{3,\Delta}_{\mathrm{qh},H},\Gamma_H^\Delta,
\Pi_H^{\mathrm{adm},\Delta}\) are the finite Siegel--Borcherds
associator cochain and the associated finite quasi-Hopf deformation
complex of Definition~\ref{def:k3e-finite-qh-associator-complex}, with
\(D_H^X\) replaced by \(Y_H^\Delta\). The transition
\(T_{H+1,H}^\Delta\) is deletion of root, tensor, and cochain basis
vectors outside the height-\(H\) window.

The packets \(\mathsf M_H^{\Delta,\mathrm{can}}\) form a strict inverse
system under the downward-saturated height transitions. Their
construction is target-side only: it supplies the matrices against
which a compact Hall packet is compared; it does not construct the
comparison matrix \(A_{\beta,\bar p}\) and does not prove any of
\(\mathcal R_H,\mathcal S_H,\mathcal D_H,\mathcal C_H,\mathcal A_H,
\mathcal P_H\)
for the compact Hall source.
\end{theorem}

\begin{proof}
The finite quotient \(\fn_+^{\leq H}\) has finite-dimensional retained
root spaces by Theorem~\ref{thm:k3e-bkm-finite-core}. Choosing
parity-homogeneous bases gives \(M_H^\Delta\) and \(P_H^\Delta\). The
Borcherds--Serre presentation gives the bracket constants and Serre
rows \(B_H^\Delta\), and PBW gives the target normal-form basis. The
invariant Borcherds form on the Cartan and root spaces gives
\(G_H^\Delta\); its kernel is the Cartan radical \(K_H^\Delta\), and
the quotient map is \(Q_H^\Delta\). The current quotient supplies
\(\Delta_H^\Delta\), whose associated graded is the primitive
enveloping coproduct of Theorem~\ref{thm:k3e-bkm-finite-core}. The only
primitive central target classes are the Cartan radical and declared
dynamical group-like parameters, giving \(Z_H^\Delta\).

The target associator is the height-\(H\) truncation of the
Siegel--Borcherds quasi-Hopf associator. Applying
Definition~\ref{def:k3e-finite-qh-associator-complex} to \(Y_H^\Delta\)
gives \(d^{3,\Delta}_{\mathrm{qh},H}\), \(\Gamma_H^\Delta\), and
\(\Pi_H^{\mathrm{adm},\Delta}\); the finite quasi-Hopf identities give
\[
 d^{3,\Delta}_{\mathrm{qh},H}(\Gamma_H^\Delta)^t=0,
 \qquad
 \Pi_H^{\mathrm{adm},\Delta}(\Gamma_H^\Delta)^t=0 .
\]
Downward saturation deletes exactly the root, tensor, and cochain
basis vectors whose height exceeds \(H\). The defining form, Serre
rows, coproduct, centre, associator cochain, deformation differential,
gauge rows, and admissibility constraints restrict to the corresponding
height-\(H\) matrices. Hence the target packets form a strict inverse
system.
\end{proof}

\begin{remark}[Executable target-packet gate]
The predicate
\texttt{evaluate\_finite\_borcherds\_target\_packet} in
\texttt{compute/lib/hall\_borcherds\_gate.py} records the finite
target-side boundary of
Theorem~\ref{thm:k3e-canonical-finite-target-packets}. It requires the
height-\(H\) current quotient, root/parity basis, invariant form and
Cartan radical, Serre presentation, PBW basis, coproduct, primitive
centre, associator complex, and transition map. Its closure supplies the
target packet \(\mathsf M_H^{\Delta,\mathrm{can}}\) only; source
faithfulness still requires the compact Hall source packet and the five
finite defects to vanish.
\end{remark}

\begin{theorem}[Finite comparison matrix packets]
\label{thm:k3e-finite-comparison-matrix-packets}
Fix a finite height \(H\). Let
\(\mathsf M_H^{X,\mathrm{can}}\) be the compact source packet of
Theorem~\ref{thm:k3e-canonical-finite-source-packets}, and let
\(\mathsf M_H^{\Delta,\mathrm{can}}\) be the Serre--Borcherds target
packet of Theorem~\ref{thm:k3e-canonical-finite-target-packets}.
Assume the following finite block data are given.
\begin{enumerate}[label=\textup{(\roman*)}]
\item a height-preserving charge bijection
\[
 \iota_H\colon M_H^X \xrightarrow{\sim} M_H^\Delta
\]
from the retained compact Hall charge blocks to the retained
Borcherds root blocks;
\item an identification of the parity splitting on each block;
\item for every \((\beta,\bar p)\) with
\(\gamma_\beta=\iota_H^{-1}(\beta)\), equality of quotient primitive
dimensions
\[
 \dim Q^X_{\gamma_\beta,\bar p}P^X_{\gamma_\beta,\bar p}
 =
 \dim Q^\Delta_{\beta,\bar p}P^\Delta_{\beta,\bar p};
\]
\item ordered bases of the two finite quotient spaces in \textup{(iii)}
and Serre-dual ordered bases in the negative halves.
\end{enumerate}
Then these data determine a finite comparison-matrix packet
\[
 \mathsf A_H^{X\to\Delta}
 =
 \bigl(
 \iota_H,\,
 A_{\beta,\bar p},\,
 A_{-\beta,\bar p},\,
 Q^X_{\gamma_\beta,\bar p},\,
 Q^\Delta_{\beta,\bar p},\,
 T^{A}_{H+1,H}
 \bigr)_{\beta,\bar p}.
\]
Here
\[
 A_{\beta,\bar p}\colon
 Q^X_{\gamma_\beta,\bar p}P^X_{\gamma_\beta,\bar p}
 \xrightarrow{\sim}
 Q^\Delta_{\beta,\bar p}P^\Delta_{\beta,\bar p}
\]
is the unique linear isomorphism carrying the chosen ordered source
quotient basis to the chosen ordered target root basis. The negative
matrix \(A_{-\beta,\bar p}\) is the Serre-dual transpose in the chosen
dual bases. In the original compact-source coordinates the matrix used
by the finite defect reductions is the composite
\[
 A_{\beta,\bar p}Q^X_{\gamma_\beta,\bar p}.
\]
If the charge bijections, parity splittings, ordered bases, quotient
maps, and Serre-dual bases commute with the downward-saturated
transitions \(T^X_{H+1,H}\) and \(T^\Delta_{H+1,H}\), then the packets
\(\mathsf A_H^{X\to\Delta}\) form an inverse system:
\[
 T^\Delta_{H+1,H} A_{H+1}=A_H T^X_{H+1,H}
\]
on every retained block, with the analogous identity in the negative
half.

The theorem constructs the finite linear comparison packet only. It
does not assert that \(A_{\beta,\bar p}\) preserves the Euler--Hall
pairing, the Borcherds--Serre rows, the coproduct, the primitive centre,
or the quasi-Hopf associator class. The parity splitting supplied in
\textup{(ii)} gives a sixth structural assertion: the comparison must
intertwine the source and target parity involutions. Those six
structural assertions are exactly the vanishing of
\(\mathcal R_H,\mathcal S_H,\mathcal D_H,\mathcal C_H,\mathcal A_H,
\mathcal P_H\)
in Theorem~\ref{thm:k3e-finite-height-hall-radical-serre-kernel}.
\end{theorem}

\begin{proof}
At finite height all retained charge/parity blocks are
finite-dimensional. After the quotient maps are fixed, condition
\textup{(iii)} says that the source quotient block and the
Serre--Borcherds root block have the same dimension. Ordered bases then
determine a unique linear isomorphism between them. Applying the same
construction to the Serre-dual bases gives the negative-half matrix;
equivalently it is the transpose inverse with the parity signs already
encoded by the chosen homogeneous bases.

The blockwise matrices assemble because \(\iota_H\) respects height and
parity. The transition identity is exactly the statement that deleting
height-\(>H\) basis vectors before applying the comparison agrees with
applying the height-\(H+1\) comparison and then deleting them. Hence the
packets form an inverse system under the stated compatibility
hypotheses.

No multiplication, coproduct, centre, pairing, or associator identity is
used in this construction. A linear isomorphism of finite blocks can
carry bases to bases while failing all structural equations. Therefore
the theorem supplies the \(A\)-matrices that enter
Theorem~\ref{thm:k3e-finite-defect-reduction}; it does not prove the
finite defects vanish.
\end{proof}

\begin{theorem}[Exact shape criterion for comparison blocks]
\label{thm:k3e-finite-comparison-shape-criterion}
In the setting of
Theorem~\ref{thm:k3e-finite-comparison-matrix-packets}, fix a retained
block \((\beta,\bar p)\). Write
\[
 n_+^X=\dim P^X_{\gamma_\beta,\bar p},\qquad
 n_-^X=\dim P^X_{\gamma_{-\beta},\bar p},\qquad
 m_+^\Delta=\dim Q^\Delta_{\beta,\bar p}P^\Delta_{\beta,\bar p},
 \qquad
 m_-^\Delta=\dim Q^\Delta_{-\beta,\bar p}P^\Delta_{-\beta,\bar p}.
\]
Let
\[
 Q_+^X\in \operatorname{Mat}_{q_+\times n_+^X},\qquad
 Q_-^X\in \operatorname{Mat}_{q_-\times n_-^X}
\]
be the source quotient maps, and let
\[
 A_+\in \operatorname{Mat}_{m_+^\Delta\times q_+},\qquad
 A_-\in \operatorname{Mat}_{m_-^\Delta\times q_-}
\]
be the post-quotient positive and Serre-dual negative comparison
matrices. The finite comparison block is well formed if and only if
\[
\begin{gathered}
 \operatorname{rank}Q_+^X=q_+=m_+^\Delta,\qquad
 \operatorname{rank}A_+=m_+^\Delta,\\
 \operatorname{rank}Q_-^X=q_-=m_-^\Delta,\qquad
 \operatorname{rank}A_-=m_-^\Delta .
\end{gathered}
\]
Equivalently, the original-coordinate matrices
\[
 \widetilde A_+=A_+Q_+^X\in
 \operatorname{Mat}_{m_+^\Delta\times n_+^X},\qquad
 \widetilde A_-=A_-Q_-^X\in
 \operatorname{Mat}_{m_-^\Delta\times n_-^X}
\]
have ranks \(m_+^\Delta\) and \(m_-^\Delta\), respectively, and have
the displayed source and target dimensions. The height transition
compatibility is the block equation
\[
 T^\Delta_{H+1,H}\,\widetilde A_{\pm,H+1}
 =
 \widetilde A_{\pm,H}\,T^X_{H+1,H}.
\]
These rank and dimension equations are necessary for the six finite
defect equations, but they do not imply any of them.
\end{theorem}

\begin{proof}
The quotient map \(Q_+^X\) identifies the compact source block with a
coordinate quotient of dimension \(q_+\) precisely when it has full row
rank. A linear map from that quotient to the target root block is an
isomorphism precisely when its matrix has \(m_+^\Delta=q_+\) rows and
columns and full rank \(m_+^\Delta\). This proves the positive-half
criterion. The negative half is the same argument applied to the
Serre-dual block. Multiplication gives
\(\widetilde A_\pm=A_\pm Q_\pm^X\), so the rank and source/target
dimension statements in original coordinates are equivalent to the
quotient statements. The transition identity is ordinary commutativity
of the two finite linear maps obtained by first restricting height and
then comparing, or first comparing and then restricting height.

No bilinear form, Serre row, coproduct, centrality equation,
associator cochain, or parity-intertwining equation appears in the
proof. The result is therefore a shape criterion for the comparison
block, not a structural recognition criterion.
\end{proof}

\begin{theorem}[Finite recognition certificate]
\label{thm:k3e-finite-recognition-certificate}
Fix a finite height \(H\). A finite Hall--Borcherds recognition
certificate consists of the following data:
\begin{enumerate}[label=\textup{(\roman*)}]
\item the compact source packet \(\mathsf M_H^{X,\mathrm{can}}\), the
target packet \(\mathsf M_H^{\Delta,\mathrm{can}}\), and a finite
comparison packet \(\mathsf A_H^{X\to\Delta}\);
\item the rank and dimension equations of
Theorem~\ref{thm:k3e-finite-comparison-shape-criterion} on every
retained charge/parity block;
\item the six structural equations
\[
 \mathcal R_H=\mathcal S_H=\mathcal D_H=\mathcal C_H=\mathcal A_H
 =\mathcal P_H=0
\]
of Theorem~\ref{thm:k3e-finite-height-hall-radical-serre-kernel};
\item commutativity of the source, target, comparison, quotient,
coproduct, centre, and associator transition matrices under
\(H+1\to H\).
\end{enumerate}
At this fixed height, these four conditions are equivalent to faithful
recognition of the retained compact Hall quotient by the
Serre--Borcherds \(\Delta_5\) quotient: the comparison matrices
\(\widetilde A_{\pm,H}\) induce a charge-preserving, parity-preserving
isomorphism of the finite radical-quotient Drinfeld double with the
finite Serre--Borcherds current quotient, preserving the bilinear form,
Serre ideal, coproduct, primitive centre, and quasi-Hopf associator
class. Failure of the certificate has four and only four types:
missing packet row, bad comparison shape, nonzero structural defect
\(\mathcal R_H,\mathcal S_H,\mathcal D_H,\mathcal C_H,\mathcal A_H,
\mathcal P_H\),
or transition incompatibility.

The certificate is finite-height. It does not by itself construct the
compact source packet, prove the certificate for all \(H\), or pass to
the completed inverse limit.
\end{theorem}

\begin{proof}
Conditions \textup{(i)} and \textup{(ii)} make the source and target
blocks comparable finite vector spaces and give well-formed comparison
matrices in the original compact-source coordinates. Condition
\textup{(iii)} is exactly the assertion that those matrices preserve
the six structures which define the finite Serre--Borcherds current
quotient: the radical of the bilinear form, the Borcherds--Serre
kernel and PBW normal forms, the coproduct with Green adjunction, the
allowed primitive centre, the associator class modulo admissible
gauge, and the parity involution on homogeneous root blocks. Therefore
the comparison is an isomorphism of the retained finite
radical-quotient double with the retained Serre--Borcherds quotient.

Conversely, any such finite quasi-Hopf isomorphism supplies all four
conditions. Its underlying linear maps give the comparison packet and
the rank equations. Preservation of the bilinear form, defining ideal,
coproduct, centre, associator class, and parity involution gives
\[
\mathcal R_H=\mathcal S_H=\mathcal D_H=\mathcal C_H=\mathcal A_H
=\mathcal P_H=0.
\]
Compatibility with height restriction is precisely the transition
commutativity in \textup{(iv)}. The listed failure types exhaust the
negations of \textup{(i)}--\textup{(iv)}.
\end{proof}

\begin{theorem}[Heightwise recognition certificate and pro-cone passage]
\label{thm:k3e-heightwise-recognition-certificate}
Let \(\{W_H\}_{H\geq 1}\) be a cofinal downward-saturated Borcherds
height system. Assume that the compact source packets
\(\mathsf M_H^{X,\mathrm{can}}\), target packets
\(\mathsf M_H^{\Delta,\mathrm{can}}\), and comparison packets
\(\mathsf A_H^{X\to\Delta}\) are supplied for all \(H\), and that the
finite Hall windows form a Mittag--Leffler inverse system in the
pro-Borcherds-cone topology. Put
\[
 D_H^X=D^{\mathrm{red},\leq H}_{\hbar}(K3\times E,\sigma),
 \qquad
 Y_H^\Delta=Y^{\mathrm{SB},\leq H}_{\hbar}(\mathfrak g_{\Delta_5}),
 \qquad
 E_H=\mathcal E^{\leq H}_{\Delta_5}(K3\times E,\sigma),
\]
and let \(\mathfrak J_H\subset D_H^X\widehat{*}_{\leq H}Y_H^\Delta\)
be the finite recognition ideal of
Definition~\ref{def:k3e-finite-recognition-envelope}. Assume also that
the height transitions carry \(\mathfrak J_{H+1}\) into
\(\mathfrak J_H\) and that the ideal system is exact:
\[
 R^1\!\varprojlim_H\mathfrak J_H=0.
\]
Then the following are equivalent.
\begin{enumerate}[label=\textup{(\roman*)}]
\item The completed comparison
\[
 \varprojlim_H
 D^{\mathrm{red},\leq H}_{\hbar}(K3\times E,\sigma)
 \;\xrightarrow{\;\sim\;}\;
 \varprojlim_H
 Y^{\mathrm{SB},\leq H}_{\hbar}(\mathfrak g_{\Delta_5})
\]
exists as a continuous quasi-Hopf superalgebra isomorphism preserving
charge, parity, pairing, Serre ideal, coproduct, primitive centre, and
associator class.
\item For every \(H\), the finite recognition certificate of
Theorem~\ref{thm:k3e-finite-recognition-certificate} holds, and the
transition matrices \(H+1\to H\) commute with the source packets,
target packets, comparison matrices, quotient maps, coproduct matrices,
centre rows, and associator complexes.
\end{enumerate}
If \textup{(ii)} fails, the first obstruction height is
\[
 H_0=\min\{H:\text{the height-\(H\) certificate or an incoming
transition at \(H\) fails}\}.
\]
The obstruction at \(H_0\) is one of the four finite types: missing
packet row, bad comparison shape, one of
\(\mathcal R_H,\mathcal S_H,\mathcal D_H,\mathcal C_H,\mathcal A_H,
\mathcal P_H\),
or a named transition component.

This theorem still assumes the compact source packets and
Mittag--Leffler exactness, including exactness of the recognition-ideal
inverse system. It identifies the exact remaining proof obligation for
compact Hall promotion; it does not construct the compact oriented
critical CoHA.
\end{theorem}

\begin{proof}
If the completed comparison exists, restriction to the quotient by
height \(>H\) gives a finite quasi-Hopf isomorphism at every \(H\).
Theorem~\ref{thm:k3e-finite-recognition-certificate} gives the finite
certificate, and functoriality of inverse limits gives commutativity of
all transition matrices.

Conversely, the finite certificates identify each finite Hall double
with the corresponding finite Serre--Borcherds quotient. Transition
commutativity says that these finite isomorphisms are morphisms of
inverse systems. The recognition envelopes fit into finite exact
sequences
\[
 0\longrightarrow \mathfrak J_H
 \longrightarrow D_H^X\widehat{*}_{\leq H}Y_H^\Delta
 \longrightarrow E_H\longrightarrow 0 .
\]
The hypotheses \(R^1\!\varprojlim_H\mathfrak J_H=0\) and
transition-compatibility of the ideals make inverse limits exact on
these sequences. Hence the inverse limit of the finite isomorphisms is
a continuous isomorphism in the pro-Borcherds-cone topology. The finite
certificates preserve the quasi-Hopf structure at every height;
continuity passes the preserved structure maps to the limit. If the
finite defects do not vanish, the same exact inverse-limit construction
still produces the completed recognized envelope
\(\varprojlim_HE_H\), not the unquotiented compact Hall double. The
formula for \(H_0\) is the first height at which this inverse-system
construction has no finite isomorphism or no compatible transition map.
\end{proof}

\begin{remark}[Executable comparison-packet gate]
The predicate
\texttt{evaluate\_finite\_comparison\_matrix\_packet} in
\texttt{compute/lib/hall\_borcherds\_gate.py} records the hypotheses of
Theorem~\ref{thm:k3e-finite-comparison-matrix-packets}. Its closure
requires source and target packets, charge and parity block bijections,
matching primitive quotient dimensions, ordered source and target
bases, quotient maps, positive and Serre-dual negative comparison
matrices, transition compatibility, and explicit separation from the
six finite defect vanishings. The predicate returns no source
faithfulness implication.

The routine \texttt{finite\_comparison\_shape\_report} checks the exact
rank and dimension equations of
Theorem~\ref{thm:k3e-finite-comparison-shape-criterion} for the resolved
positive and negative comparison matrices. Its closure says only that
the matrices have the correct finite source, quotient, and target
shapes.

The routine \texttt{finite\_recognition\_certificate\_report} combines
the comparison packet, the exact shape report, the finite matrix defect
report, and transition compatibility. It closes precisely at the
finite-height certificate of
Theorem~\ref{thm:k3e-finite-recognition-certificate}.

The routine \texttt{heightwise\_recognition\_certificate\_report}
orders the finite certificates by height, checks the supplied transition
reports, and returns the first obstruction height \(H_0\) with its
packet, shape, structural-defect, or transition label. It also records
the pro-envelope exactness gates
\(\mathfrak J_{H+1}\mapsto\mathfrak J_H\) and
\(R^1\varprojlim_H\mathfrak J_H=0\), because finite certificates do not
by themselves make the completed envelope exact.
\end{remark}

\begin{theorem}[Conditional Hall--Drinfeld/Super-Yangian equivalence criterion]
\label{thm:k3e-hall-drinfeld-super-yangian-criterion}
Let $\mathcal H^+_{K3\times E}$ be an oriented critical CoHA for
$K3\times E$ on a fixed stability chamber, equipped at every
finite Borcherds height $H$ with a candidate negative half
\(\mathcal H^{-,\leq H}_{K3\times E}\), Cartan part
\(\mathcal H^{0,\leq H}_{K3\times E}\), Hall coproduct, and continuous
Hall pairing. Let
$D^\wedge_\hbar(\mathcal H^+_{K3\times E})$ denote the inverse-limit
Hall--Drinfeld double obtained from this finite-height system, and let
$Y^\wedge_\hbar{}^{\mathrm{super}}(\mathfrak g_{\Delta_5})$ denote the
completed Serre--Borcherds current presentation with classical limit
$U(\mathfrak g_{\Delta_5})$. A topological quasi-Hopf superalgebra
isomorphism
\[
 D^\wedge_\hbar(\mathcal H^+_{K3\times E})
 \;\cong\;
 Y^\wedge_\hbar{}^{\mathrm{super}}(\mathfrak g_{\Delta_5})
\]
whose semiclassical limit is the identity on
$\mathfrak g_{\Delta_5}$ exists if and only if the following seven
compatibilities hold.
\begin{enumerate}[label=\textup{(\roman*)}]
\item \emph{Signed root character.} The BPS grading on
$\mathcal H^+_{K3\times E}$ is identified with the Borcherds cone
$\Delta_+$, and each graded primitive piece has Euler/superdimension
target \(c(D(\alpha))\), including the imprimitive Borcherds divisor
correction. The equality of ordinary primitive dimensions with
\(|c(D(\alpha))|\) is a consequence only of the finite-height parity
fixture or target-presentation computation recorded in the recognition
data.
\item \emph{PBW filtration.} The associated graded algebra of the positive
Hall half is $U(\fn_+(\mathfrak g_{\Delta_5}))$ with the PBW basis of
Theorem~\ref{thm:k3e-bkm-finite-core}; equivalently, the Hall product
has no extra primitive generators and no extra relations beyond the
Borcherds--Serre ideal after passing to the associated graded.
\item \emph{Negative half and Hall pairing.} The negative half is the
continuous graded dual of the positive half, modulo the Borcherds
Cartan radical, and the Hall pairing between the two halves is
nondegenerate on every finite height quotient and
matches the invariant Borcherds form on root spaces, with the same
$\hbar$-normalisation on the Cartan.
\item \emph{Coproduct.} The Hall coproduct is the Drinfeld-new coproduct
attached to that pairing, its semiclassical cobracket is the
$\Delta_5$-regulated Lie bialgebra cobracket, and its associated graded
reduces to the primitive enveloping coproduct of
Theorem~\ref{thm:k3e-bkm-finite-core}.
\item \emph{Serre ideal.} The kernel of the current presentation is exactly
the completed Borcherds--Serre ideal: real-root cubic relations on the
$F_3$ core, orthogonal imaginary-root super-commutativity, and the
quantum corrections prescribed by the chosen completed current
presentation; no additional Hall relation survives in the associated
graded.
\item \emph{Centre.} The centre of the completed double maps to the
Borcherds Cartan radical and to the declared dynamical central
group-like parameters only. In particular, orientation lines,
determinant-line twists, and stability-chamber scalars introduce no
extra primitive central elements.
\item \emph{Completion and associator.} Both sides use the same
pro-Borcherds-cone topology, the same finite-height inverse system with
surjective Mittag--Leffler transition maps, and the same
Siegel--Borcherds associator class; the universal $R$-matrix converges
in that topology and has the stated $\Delta_5$ semiclassical limit.
\end{enumerate}
\end{theorem}

\begin{proof}
Necessity follows by applying any such isomorphism to the root grading,
PBW filtration, Hopf pairing, coproduct, defining ideal, centre, and
completion topology. Each structure is functorial under a topological
quasi-Hopf isomorphism and survives on every finite height quotient.
Conversely, conditions \textup{(i)}--\textup{(ii)} identify the positive
half with the Borcherds enveloping algebra on associated gradeds.
Condition \textup{(iii)} identifies the negative half and the Cartan
needed for Drinfeld doubling. Conditions \textup{(iv)}--\textup{(v)}
identify the bialgebra and current-presentation quotients. Conditions
\textup{(vi)}--\textup{(vii)} rule out central and topological
ambiguities. The finite height isomorphisms are compatible under the
inverse system, hence glue to the completed topological quasi-Hopf
isomorphism.
\end{proof}

\begin{corollary}[Finite-height obstruction to compact Hall promotion]
\label{cor:k3e-finite-height-promotion-obstruction}
Let
\[
 \Psi^+_{\Hall\to\BKM}\colon
 \mathcal H^+_{K3\times E}
 \longrightarrow
 \widehat{U(\fn_+(\mathfrak g_{\Delta_5}))}
\]
be a proposed positive-half comparison satisfying
Theorem~\ref{thm:k3e-positive-half-hall-borcherds-criterion} on
characters. It promotes to a completed Hall--Drinfeld/BKM denominator
object
\[
 D^\wedge_\hbar(\mathcal H^+_{K3\times E})
 \longrightarrow
 Y^\wedge_\hbar{}^{\mathrm{super}}(\mathfrak g_{\Delta_5})
\]
if and only if the seven conditions of
Theorem~\ref{thm:k3e-hall-drinfeld-super-yangian-criterion} hold on
every finite-height quotient and commute with the inverse-limit maps.
Equivalently, failure at a single height \(H\) of any one of
\begin{enumerate}[label=\textup{(\alph*)}]
\item primitive signed-character, parity, and PBW recognition,
\item nondegeneracy of the positive--negative Hall pairing,
\item Hall coproduct reducing to the primitive enveloping coproduct on
associated gradeds,
\item equality of the Hall kernel with the Borcherds--Serre ideal,
\item absence of extra central primitive classes,
\item pro-cone completion and Siegel--Borcherds associator compatibility
\end{enumerate}
obstructs promotion. The equality
\(-(\Phi_{10}^{\mathrm{OP}})^{-1}=-D_5^{-2}=-4096\,\Delta_5^{-2}\) and the value
\(\kappa_{\mathrm{BKM}}(\Delta_5)=5\) verify the denominator character
and weight; they do not force the finite defects to vanish, and hence
do not identify the constructed radical-quotient double, or its inverse
limit, with the Borcherds current object.
\end{corollary}

\begin{proof}
The theorem gives necessity by restriction to the finite quotient of
height \(H\). Conversely, if all finite-height quotients satisfy the
listed conditions and the transition maps commute with them, the
finite-height doubles form an inverse system of quasi-Hopf
superalgebras. The inverse limit is precisely the completed
Hall--Drinfeld double, and the compatible maps to the finite
Borcherds--current quotients glue to the displayed continuous morphism.
If one condition fails at height \(H\), restriction of any putative
completed isomorphism to that quotient would violate the corresponding
finite-dimensional invariant; hence no completed promotion exists.
\end{proof}

\begin{remark}[Compact-passage and double-data obstructions]
\label{rem:k3e-exact-missing-promotion-inputs}
The obstruction theorem has two layers before finite recognition begins.
The finite DWR/Ran Rees construction gives critical Hall charts at bounded
charge, jet, renormalisation order, and Ran arity. It becomes a compact
Hall comparison only on the compact-passage locus
\[
 o_{\mathrm{ML}}=o_{\mathrm{real}}=o_{\mathrm{cpt}}=0
\]
of Definition~\ref{def:compact-oriented-critical-hall-datum}: exactness
of the inverse system, monoidal critical realisation with the
Thom--Sebastiani orientation line, and compact-support functoriality of
Hall convolution. After the positive half exists, doubling is a separate
problem. The compact Hall--Drinfeld double exists only after
\[
 o_\Delta=o_{\mathrm{pair}}=o_{\mathrm{cent}}=0,
\]
namely after the continuous coproduct, the completed Hopf pairing with
radical quotient, and the centre compatibility have been supplied.
The six finite defects
\(\mathcal R_H,\mathcal S_H,\mathcal D_H,\mathcal C_H,\mathcal A_H,
\mathcal P_H\)
then measure recognition against the Borcherds current quotient at
height \(H\). Thus the missing inputs split as follows, each with
finite-height compatibility:
\begin{enumerate}[label=\textup{(\roman*)}]
\item the compact-passage vanishings
\(o_{\mathrm{ML}}=o_{\mathrm{real}}=o_{\mathrm{cpt}}=0\), producing the
realised oriented critical CoHA
\(\mathcal H^+_{K3\times E}\) together with the candidate negative half,
Cartan part, and Hall pairing on every Borcherds-height truncation;
\item the double-data vanishings
\(o_\Delta=o_{\mathrm{pair}}=o_{\mathrm{cent}}=0\), giving the continuous
coproduct, non-degenerate completed pairing, radical quotient, and
allowed centre;
\item the finite-height recognition equalities
\(\mathcal R_H=\mathcal S_H=0\), identifying the pairing radical with
the Borcherds Cartan radical and the PBW kernel with the
Borcherds--Serre ideal;
\item the finite-height equality
\(\mathcal D_H=0\), identifying the Hall coproduct with the
Drinfeld-new coproduct on the completed current quotient;
\item the finite-height equalities
\(\mathcal C_H=\mathcal A_H=0\), excluding undeclared primitive central
classes and matching the Hall and Siegel--Borcherds associator cocycles
modulo charge/parity-admissible gauge coboundaries.
\end{enumerate}
The downward-saturated transition maps must preserve the six
compact/double vanishings and the six finite recognition defects. The
denominator identity
\(\kappa_{\mathrm{BKM}}(\Delta_5)=5\) is a scalar shadow of the target;
it supplies none of these compact-passage, double-data, or
finite-recognition vanishings.
\end{remark}

\begin{theorem}[Finite-height Hall radical, coproduct, and Serre-kernel obstruction]
\label{thm:k3e-finite-height-hall-radical-serre-kernel}
Assume the finite-height Hall data of
Theorem~\ref{thm:k3e-hall-drinfeld-super-yangian-criterion}, and
suppose its signed-character and PBW conditions already hold at every
height. For each height \(H\), let
\[
 \mathfrak r^{0,\leq H}_{\Delta_5}
 :=
 \ker\!\left(
 \mathcal H^{0,\leq H}_{K3\times E}
 \times
 \mathcal H^{0,\leq H}_{K3\times E}
 \to \mathbb C((\hbar))
 \right)
\]
be the Borcherds Cartan radical, let
\(\mathfrak I^{\mathrm{Serre},\leq H}_{\Delta_5}\) be the
finite-height Borcherds--Serre ideal of
Theorem~\ref{thm:k3e-bkm-finite-core}, and let
\(\Phi^{\mathrm{SB},\leq H}_{\Delta_5}\) denote the
Siegel--Borcherds associator class on the height-\(H\) current
quotient. Write
\[
 \pi_H\colon
 T(\operatorname{Prim}^{\leq H}\mathcal H^+_{K3\times E})
 \longrightarrow
 \operatorname{gr}\mathcal H^{+,\leq H}_{K3\times E}
\]
for the primitive-presentation map, where \(T(-)\) denotes the tensor
algebra. Let
\[
 \Psi^{+,\leq H}_{\Hall\to\BKM}\colon
 \mathcal H^{+,\leq H}_{K3\times E}
 \longrightarrow
 U(\fn_+(\mathfrak g_{\Delta_5}^{\leq H}))
\]
be the finite-height positive-half comparison, and write
\(\Delta^{\mathrm{Hall},\leq H}\) and
\(\Delta^{\Delta_5,\leq H}\) for the Hall coproduct and the
Borcherds-current coproduct on the finite quotient. Let
\(\Pi^X_H\) and \(\Pi^\Delta_H\) be the parity involutions on the
retained source primitive blocks and the retained
\(\Delta_5\)-root blocks. Define six
defects:
\[
\begin{aligned}
 \mathcal R_H
 &=
 \frac{\operatorname{rad}\langle-,-\rangle_H}
 {\operatorname{rad}\langle-,-\rangle_H
 \cap\mathfrak r^{0,\leq H}_{\Delta_5}}
 \oplus
 \frac{\mathfrak r^{0,\leq H}_{\Delta_5}}
 {\operatorname{rad}\langle-,-\rangle_H
 \cap\mathfrak r^{0,\leq H}_{\Delta_5}},\\
 \mathcal S_H
 &=
 \frac{\ker(\pi_H)}
 {\ker(\pi_H)\cap\mathfrak I^{\mathrm{Serre},\leq H}_{\Delta_5}}
 \oplus
 \frac{\mathfrak I^{\mathrm{Serre},\leq H}_{\Delta_5}}
 {\ker(\pi_H)\cap\mathfrak I^{\mathrm{Serre},\leq H}_{\Delta_5}},\\
 \mathcal D_H
 &=
 \left(
 (\Psi^{+,\leq H}_{\Hall\to\BKM}\widehat\otimes
 \Psi^{+,\leq H}_{\Hall\to\BKM})
 \circ \Delta^{\mathrm{Hall},\leq H}
 \right)
 -
 \left(
 \Delta^{\Delta_5,\leq H}\circ
 \Psi^{+,\leq H}_{\Hall\to\BKM}
 \right),\\
 \mathcal C_H
 &=
 \frac{\operatorname{Prim}Z(D^{\leq H}_\hbar(\mathcal H^+_{K3\times E}))}
 {\operatorname{Prim}Z(D^{\leq H}_\hbar(\mathcal H^+_{K3\times E}))
 \cap\mathfrak z^{\mathrm{allow},\leq H}}
 \oplus
 \frac{\mathfrak z^{\mathrm{allow},\leq H}}
 {\operatorname{Prim}Z(D^{\leq H}_\hbar(\mathcal H^+_{K3\times E}))
 \cap\mathfrak z^{\mathrm{allow},\leq H}},\\
 \mathcal A_H
 &=
 \left[\Phi^{\mathrm{Hall},\leq H}\right]
 -
 \left[\Phi^{\mathrm{SB},\leq H}_{\Delta_5}\right],\\
 \mathcal P_H
 &=
 \Pi^\Delta_H\Psi^{+,\leq H}_{\Hall\to\BKM}
 -
 \Psi^{+,\leq H}_{\Hall\to\BKM}\Pi^X_H .
\end{aligned}
\]
Here
\[
 \mathfrak z^{\mathrm{allow},\leq H}
 :=
 \mathfrak r^{0,\leq H}_{\Delta_5}
 \oplus
 \mathfrak z^{\mathrm{dyn},\leq H},
\]
and \(\mathfrak z^{\mathrm{dyn},\leq H}\) is the declared span of
dynamical group-like central parameters. The class
\(\mathcal A_H\) is taken in quasi-Hopf associator cohomology modulo
gauge. Under the standing positive-half hypotheses, the
proposed compact Hall comparison promotes to the completed
Hall--Drinfeld/Borcherds current object if and only if, for every height
\(H\),
\[
 \mathcal R_H=\mathcal S_H=\mathcal D_H=\mathcal C_H=\mathcal A_H
 =\mathcal P_H=0
\]
and the transition maps \(H+1\to H\) carry the six vanished defects to
the six vanished defects. If residual promotion fails, its first
possible failure height is the minimum
\[
 H_0
 =
\min\{H:\mathcal R_H\neq0\ \text{or}\ \mathcal S_H\neq0\
 \text{or}\ \mathcal D_H\neq0\ \text{or}\ \mathcal C_H\neq0\
 \text{or}\ \mathcal A_H\neq0\ \text{or}\ \mathcal P_H\neq0
 \text{ or a transition map is incompatible}\}.
\]
At \(H_0\), the nonzero component names the residual obstruction:
residual Hall pairing radical modulo the Borcherds Cartan radical,
excess or missing Serre kernel, coproduct mismatch, an undeclared
primitive central class, failure of the inverse-limit
associator/completion compatibility, or parity-fixture failure.
The denominator identity and the five-construction spectrum
\(\{0,0,3,5,24\}\) do not affect this criterion.
\end{theorem}

\begin{proof}
At fixed height the Hall data are finite over each Borcherds-cone
slice. A Hall--Drinfeld double is nondegenerate only after quotienting
by the radical of its Hopf pairing; the Borcherds target permits only
the Cartan radical \(\mathfrak r^{0,\leq H}_{\Delta_5}\). Thus
\(\mathcal R_H=0\) is exactly the finite-height nondegeneracy condition
modulo the allowed Cartan radical.

The associated graded positive half is generated by its primitive
root spaces. By Theorem~\ref{thm:k3e-bkm-finite-core}, the target
kernel is precisely the Borcherds--Serre ideal. Hence
\(\mathcal S_H=0\) is equivalent to exactness of the Serre kernel at
height \(H\); any nonzero class changes the PBW quotient and prevents a
current-presentation isomorphism. The coproduct is a structure map, so
the comparison must make the square
\[
\begin{tikzcd}
\mathcal H^{+,\leq H}_{K3\times E}
 \arrow[r,"\Delta^{\mathrm{Hall},\leq H}"]
 \arrow[d,"\Psi^{+,\leq H}_{\Hall\to\BKM}"']
&
\mathcal H^{+,\leq H}_{K3\times E}\widehat\otimes
\mathcal H^{+,\leq H}_{K3\times E}
 \arrow[d,"\Psi^{+,\leq H}_{\Hall\to\BKM}\widehat\otimes
 \Psi^{+,\leq H}_{\Hall\to\BKM}"]
\\
U(\fn_+(\mathfrak g_{\Delta_5}^{\leq H}))
 \arrow[r,"\Delta^{\Delta_5,\leq H}"']
&
U(\fn_+(\mathfrak g_{\Delta_5}^{\leq H}))\widehat\otimes
U(\fn_+(\mathfrak g_{\Delta_5}^{\leq H}))
\end{tikzcd}
\]
commute; this is precisely \(\mathcal D_H=0\). A quasi-Hopf isomorphism
preserves primitive central elements, so the target allows only the
Cartan radical and the declared dynamical central group-likes; this is
exactly \(\mathcal C_H=0\). It also preserves the associator gauge
class, and therefore the finite quotient must carry the
Siegel--Borcherds class: \(\mathcal A_H=0\).
The comparison is a morphism of superalgebras only if it intertwines
the parity involutions on homogeneous primitive blocks; this is exactly
\(\mathcal P_H=0\). A nonzero parity defect changes the supercommutator
signs and cannot be repaired by quotienting the radical, imposing
Serre rows, changing the coproduct, or gauging the associator.

If all six defects vanish and the transition maps preserve the
quotients, the finite-height doubles form an inverse system whose
height-\(H\) terms identify with the corresponding Borcherds current
quotients. Taking the pro-Borcherds-cone inverse limit gives the
completed comparison. Conversely, any completed comparison restricts
to every finite height and therefore forces the six displayed defects
and the transition incompatibility to vanish. The minimum \(H_0\), when
it exists, is consequently the first height at which promotion can
fail.
\end{proof}

\begin{remark}[Executable heightwise obstruction predicate]
The predicate
\texttt{heightwise\_recognition\_report} in
\texttt{compute/lib/hall\_borcherds\_gate.py} implements the finite
boundary of Theorem~\ref{thm:k3e-finite-height-hall-radical-serre-kernel}.
From witness packets at positive heights it returns the ordered
recognition rows, the first obstructed height, and the obstruction labels
among the structural prerequisites
\texttt{finite\_compact\_double},
\texttt{finite\_borcherds\_target},
\texttt{compact\_source\_packet},
and, once these objects exist, among the finite defects
\(\mathcal R_H,\mathcal S_H,\mathcal D_H,\mathcal C_H,\mathcal A_H,
\mathcal P_H\),
together with transition compatibility. It records the finite
\(H_0\)-criterion;
construction of the compact Hall data is exactly the compact-passage,
double-data, and finite-recognition problem above.
\end{remark}

\begin{theorem}[Finite defect reduction to matrices, Green adjunction, and primitive support]
\label{thm:k3e-finite-defect-reduction}
Work in the setting of
Theorem~\ref{thm:k3e-finite-height-hall-radical-serre-kernel} at a
fixed finite height \(H\). Choose parity-resolved compact source bases
for the primitive Hall spaces, the Serre-dual negative primitive spaces,
and the Cartan lattice. Let
\[
 G^X_{\gamma,\bar p},\quad
 K^X_{\gamma,\bar p},\quad
 Q^X_{\gamma,\bar p},\quad
 A_{\beta,\bar p},\quad
 B^X_{\alpha,\beta},\quad
 \Delta^X_H,\quad
 \Phi^X_H,\quad
 d^{3,H}_{\mathrm{qh}},\quad
 \Gamma_H,\quad
 \Pi_H^{\mathrm{adm}},\quad
 \Pi^X_H,\quad
 \Pi^\Delta_H
\]
denote respectively the compact Euler--Hall pairing matrices, their
radical kernels, the quotient maps, the comparison matrices to the
\(\Delta_5\)-root blocks, the Hall primitive bracket matrices, the Hall
coproduct, the finite Hall associator cochain, the finite pentagon
differential on associator cochains, a row basis for gauge coboundaries,
the charge/parity constraint matrix on those gauge rows, and the source
and target parity sign involutions. The comparison matrices
\(A_{\beta,\bar p}\) are the block matrices of
Theorem~\ref{thm:k3e-finite-comparison-matrix-packets}. Let \(G^\Delta_\beta\),
\(\mathfrak I^{\mathrm{Serre},\leq H}_{\Delta_5}\),
\(\Delta^{\Delta_5,\leq H}\), and
\(\Phi^{\mathrm{SB},\leq H}_{\Delta_5}\) be the corresponding
Serre--Borcherds target data, equivalently the corresponding rows of
\(\mathsf M_H^{\Delta,\mathrm{can}}\) from
Theorem~\ref{thm:k3e-canonical-finite-target-packets}. Then the six
finite defects reduce to the following finite assertions.
\begin{enumerate}[label=\textup{(\alph*)}]
\item \(\mathcal R_H=0\) if and only if, for every retained charge
block and parity block,
\[
 K^X_{\gamma_\beta,\bar p}
 =
 \ker G^X_{\gamma_\beta,\bar p}
 \cap
 \ker (G^X_{\gamma_{-\beta},\bar p})^t
\]
and the quotient comparison is an isometry
\[
 G^\Delta_\beta\!\left(
 A_{\beta,\bar p}Q^X_{\gamma_\beta,\bar p}x,\,
 A_{-\beta,\bar p}Q^X_{\gamma_{-\beta},\bar p}y
 \right)
 =
 G^X_{\gamma_\beta,\bar p}(x,y).
\]
\item \(\mathcal S_H=0\) if and only if the finite Hall primitive
bracket matrices satisfy the Borcherds--Serre inclusions and the PBW
normal-form count
\[
 \operatorname{Hilb}\!
 \left(
 T(\operatorname{Prim}^{\leq H}\mathcal H^+)/
 \mathfrak I^{\mathrm{Serre},\leq H}_{\Delta_5}
 \right)
 =
 \operatorname{Hilb}\!
 \left(\operatorname{gr}\mathcal H^{+,\leq H}\right),
\]
equivalently
\(\ker(\pi_H)=\mathfrak I^{\mathrm{Serre},\leq H}_{\Delta_5}\).
\item If the product comparison and the pairing isometry in
\textup{(a)} hold, then \(\mathcal D_H=0\) is equivalent to the finite
compact Green identity
\[
 \langle \Delta^{\mathrm{Hall},\leq H}x,\,
 y\otimes z\rangle_H
 =
 \langle x,\, yz\rangle_H
\]
for the same Thom--Sebastiani orientation line, Tate shift, Rees
normalisation, and Serre-dual negative half used in the product and
pairing.
\item \(\mathcal C_H=0\) if and only if the zero-charge primitive
support is
\[
 \operatorname{Prim}\!\left(
 D^{\mathrm{red},\leq H}_\hbar(K3\times E,\sigma)
 \right)_0
 =
 \mathfrak h^{\leq H}_{\mathrm{Cartan}}
 \oplus
 \mathfrak z^{\mathrm{dyn},\leq H}
\]
and centrality against every retained root generator forces the Cartan
component into
\(\mathfrak r^{0,\leq H}_{\Delta_5}\).
\item \(\mathcal A_H=0\) if and only if
\[
 d^{3,H}_{\mathrm{qh}}\Phi^X_H
 =
 d^{3,H}_{\mathrm{qh}}\Phi^{\mathrm{SB},\leq H}_{\Delta_5}
 =
 0,\qquad
 \Phi^X_H-\Phi^{\mathrm{SB},\leq H}_{\Delta_5}
 \in \operatorname{rowspan}(\Gamma_H),\qquad
 d^{3,H}_{\mathrm{qh}}\Gamma_H^t=0,\qquad
 \Pi_H^{\mathrm{adm}}\Gamma_H^t=0
\]
or equivalently
\([\Phi^X_H]=[\Phi^{\mathrm{SB},\leq H}_{\Delta_5}]\) in finite
quasi-Hopf associator cohomology with gauge transformations preserving
the charge filtration and the signed parity convention.
\item \(\mathcal P_H=0\) if and only if every retained comparison block
intertwines the parity involutions:
\[
 \Pi^\Delta_H A_{\beta,\bar p}Q^X_{\gamma_\beta,\bar p}
 =
 A_{\beta,\bar p}Q^X_{\gamma_\beta,\bar p}\Pi^X_H .
\]
In coordinates this is the rank-zero condition on the finite matrix
\(\Pi^\Delta_H A_{\beta,\bar p}Q^X_{\gamma_\beta,\bar p}
-A_{\beta,\bar p}Q^X_{\gamma_\beta,\bar p}\Pi^X_H\).
\end{enumerate}
All six reductions are functorial for the height transition
\(H+1\to H\) precisely when the displayed kernels, quotient maps,
normal forms, Green pairings, primitive-support decompositions, parity
sign involutions, and
associator cochains, pentagon differentials, gauge rows, and
admissibility constraints are carried to the corresponding height-\(H\)
objects.
\end{theorem}

\begin{proof}
For \textup{(a)}, the radical obstruction compares two bilinear forms:
the compact Euler--Hall form after quotienting by its source radical and
the invariant \(\Delta_5\)-Borcherds form modulo its Cartan radical. A
finite bilinear form is identified with another one exactly by the
kernel equality and the displayed quotient isometry. This is the
definition of equality of radicals after the comparison matrix
\(A_{\beta,\bar p}\) is fixed.

For \textup{(b)}, the map \(\pi_H\) presents the associated graded
positive Hall half from its primitive generators. The Borcherds target
has defining ideal
\(\mathfrak I^{\mathrm{Serre},\leq H}_{\Delta_5}\) by
Theorem~\ref{thm:k3e-bkm-finite-core}. Once the Serre rows vanish under
the Hall bracket matrices, equality of Hilbert series is equivalent to
absence of any further primitive relation; this is exactly
\(\ker(\pi_H)=\mathfrak I^{\mathrm{Serre},\leq H}_{\Delta_5}\).

For \textup{(c)}, compare the two coproducts against the nondegenerate
tensor product pairing. For \(x\in\mathcal H^{+,\leq H}\) and
\(y,z\in\mathcal H^{-,\leq H}\),
\[
\begin{aligned}
\left\langle
(\Psi_H^+\widehat\otimes\Psi_H^+)\Delta^{\mathrm{Hall},\leq H}x,\,
\Psi_H^-y\otimes\Psi_H^-z
\right\rangle_{\Delta_5}
&=
\langle\Delta^{\mathrm{Hall},\leq H}x,\,y\otimes z\rangle_H\\
&=
\langle x,\,yz\rangle_H\\
&=
\left\langle
\Psi_H^+x,\,(\Psi_H^-y)(\Psi_H^-z)
\right\rangle_{\Delta_5}\\
&=
\left\langle
\Delta^{\Delta_5,\leq H}\Psi_H^+x,\,
\Psi_H^-y\otimes\Psi_H^-z
\right\rangle_{\Delta_5}.
\end{aligned}
\]
Nondegeneracy gives equality of coproducts. Conversely, equality of the
coproduct square gives the displayed Green identity by pairing with
products in the negative half.

For \textup{(d)}, Cartan conjugation separates charges:
\[
 K_\lambda xK_\lambda^{-1}=e^{\hbar(\lambda,\gamma)}x
\]
for a homogeneous primitive \(x\) of charge \(\gamma\). Over
\(\mathbb C((\hbar))\), centrality forces \((\lambda,\gamma)=0\) for
all \(\lambda\), hence only Cartan-radical charge can survive. A
primitive element cannot contain a nontrivial PBW product \(e_\alpha
f_\beta\), since the coproduct would contain mixed cross-terms. The
displayed zero-charge primitive-support identity therefore leaves only
Cartan and declared dynamical primitives, and centrality forces the
Cartan part into the allowed Borcherds radical.

For \textup{(e)}, the height-\(H\) quasi-Hopf deformation complex
controls the associator. Its closed degree-three cochains are precisely
the finite associator representatives satisfying the pentagon equation
at this height, and gauge equivalence changes a representative by a
degree-two coboundary. The equation
\(d^{3,H}_{\mathrm{qh}}\Gamma_H^t=0\) records that the supplied gauge
rows are genuine coboundaries in the finite complex, not arbitrary
degree-three rows. The admissible gauge transformations are the ones
preserving the charge filtration and signed parity, encoded by
\(\Pi_H^{\mathrm{adm}}\Gamma_H^t=0\). Hence equality of the finite Hall
and Siegel--Borcherds associator classes is exactly the displayed
cocycle, row-span, gauge-cocycle, and admissibility condition, i.e.
\(\mathcal A_H=0\).

\textup{(f)} is the ordinary matrix equation saying that the
comparison preserves the \(\mathbb Z/2\)-grading. The transition claim
is the naturality of each displayed finite matrix, kernel, quotient,
parity involution, and cohomology class under restriction of the
downward-saturated height window.
\end{proof}

\begin{proposition}[Finite coproduct-intertwining and Green adjunction]
\label{prop:k3e-finite-coproduct-intertwining-gate}
Fix the finite height \(H\) and the source and target bases of
Theorem~\ref{thm:k3e-finite-defect-reduction}. Suppose that the Hall
coproduct matrix \(\Delta^X_H\), the Serre--Borcherds coproduct matrix
\(\Delta^{\Delta_5,\leq H}\), the tensor product pairing matrix
\(T_H\), the source Hall pairing matrix \(G_H\), and the negative-half
product matrix \(M_H\) are supplied in these bases, with the same
orientation line, Tate shift, Rees normalisation, and Serre-dual negative
half. Then the finite coproduct comparison and the finite Green identity
hold exactly when
\[
 \operatorname{rank}\!
 \left(
   \Delta^X_H-\Delta^{\Delta_5,\leq H}
 \right)=0
\]
and
\[
 \operatorname{rank}\!
 \left(
   (\Delta^X_H)^tT_H-G_HM_H
 \right)=0 .
\]
The first rank is the finite coproduct defect \(D^\Delta_H\). The second
rank is the Green-adjunction defect for
\[
 \langle \Delta^{\mathrm{Hall},\leq H}x,\,
 y\otimes z\rangle_H
 =
 \langle x,\, yz\rangle_H .
\]
\end{proposition}

\begin{proof}
A linear map between finite-dimensional rational vector spaces is equal
to another linear map in the chosen bases if and only if the difference
matrix has rank zero. This gives the first displayed condition.

Write the Hall coproduct in column convention, so that its rows are
indexed by tensor-product positive basis vectors and its columns by
positive basis vectors. Pairing the output of the coproduct with a
tensor-product negative basis vector gives the matrix
\((\Delta^X_H)^tT_H\). Multiplying two negative basis vectors first and
then pairing the result with the positive input gives \(G_HM_H\). The
finite Green identity holds on basis vectors if and only if these two
matrices agree; by linearity, this is equivalent to
\(\operatorname{rank}((\Delta^X_H)^tT_H-G_HM_H)=0\). Under the product
comparison and pairing isometry of
Theorem~\ref{thm:k3e-finite-defect-reduction}\textup{(a),(c)}, the two
rank-zero conditions are therefore exactly the finite \(D_H\)-vanishing
condition with Green adjunction included.
\end{proof}

\begin{remark}[Executable finite matrix reduction]
The predicate
\texttt{finite\_matrix\_defect\_report} in
\texttt{compute/lib/hall\_borcherds\_gate.py} implements the rational
finite-witness model of the linear algebra in
Theorem~\ref{thm:k3e-finite-defect-reduction}. It computes the
radical-kernel symmetric-difference dimension and target-form pullback
rank for \(\mathcal R_H\), the symmetric-difference dimensions for
\(\mathcal S_H\) and \(\mathcal C_H\), and, when Serre rows and Hilbert
vectors are supplied, the finite Serre-inclusion defect \(RE\) and PBW
Hilbert-count defect for \(\mathcal S_H\). For \(\mathcal C_H\) it also
accepts retained-root centrality equations and Cartan-component/radical
bases, checking that the centrality matrix has rank zero and the Cartan
component lies in the Borcherds Cartan radical. It also computes the rank
of the finite coproduct defect and, when the product data are supplied,
the Green-adjunction defect \(\Delta^tT-GM\) for \(\mathcal D_H\), and
membership of
\(\Phi^X_H-\Phi^{\mathrm{SB},\leq H}_{\Delta_5}\) in the supplied
gauge-coboundary row span for \(\mathcal A_H\). The standalone predicate
\texttt{finite\_coproduct\_intertwining\_gate} records the two matrices
in Proposition~\ref{prop:k3e-finite-coproduct-intertwining-gate}; it
does not close when the Green data \(T_H,G_H,M_H\) are absent. When the finite
pentagon differential and charge/parity constraint matrix are supplied,
it also checks
\(d^{3,H}_{\mathrm{qh}}\Phi^X_H
=d^{3,H}_{\mathrm{qh}}\Phi^{\mathrm{SB},\leq H}_{\Delta_5}=0\)
as well as
\(d^{3,H}_{\mathrm{qh}}\Gamma_H^t=0\) and
\(\Pi_H^{\mathrm{adm}}\Gamma_H^t=0\). The pairing comparison may
be supplied either by composed source-to-target maps
\(A_{\beta,\bar p}Q^X_{\gamma_\beta,\bar p}\) and
\(A_{-\beta,\bar p}Q^X_{\gamma_{-\beta},\bar p}\), or by the corresponding
quotient maps and post-quotient comparison matrices in the two slots. When
the quotient maps are supplied, their kernels are checked against the
declared radical bases \(K^X_{\gamma_\beta,\bar p}\) and
\(K^X_{\gamma_{-\beta},\bar p}\). It also computes the parity defect
rank of
\(\Pi^\Delta_HA_{\beta,\bar p}Q^X_{\gamma_\beta,\bar p}
-A_{\beta,\bar p}Q^X_{\gamma_\beta,\bar p}\Pi^X_H\), so missing or
mismatched parity data remain a named finite obstruction. Its output is the finite
witness packet consumed by
\texttt{heightwise\_recognition\_report}; the composed predicate
\texttt{heightwise\_matrix\_defect\_report} performs this passage over all
declared finite heights, accepts exact restricted-transition witnesses,
and returns the corresponding \(H_0\)-report. The matrices and transition
maps themselves remain the compact Hall--Borcherds data whose construction
is required by the theorem.
\end{remark}

\begin{theorem}[Exact recognition of the constructed finite double]
\label{thm:k3e-constructed-finite-double-recognition}
Fix \(H\), \(\sigma\), and \(W_H\). Let
\[
 D^{\mathrm{red},\leq H}_{\hbar}(K3\times E,\sigma)
\]
be the finite radical-quotient double of
Theorem~\ref{thm:k3e-finite-compact-hall-drinfeld-double}, and let
\[
 Y^{\mathrm{SB},\leq H}_{\hbar}(\mathfrak g_{\Delta_5})
\]
be the height-\(H\) Serre--Borcherds current quotient with its Cartan
radical and Siegel--Borcherds associator class. Suppose a finite
positive-half comparison
\[
 \Psi^{+,\leq H}_{\Hall\to\BKM}\colon
 \mathcal H^{+,\leq H}_{K3\times E}
 \longrightarrow
 U(\fn_+(\mathfrak g_{\Delta_5}^{\leq H}))
\]
has the signed-character and parity fixtures of
Theorem~\ref{thm:k3e-positive-half-hall-borcherds-criterion}. Then a
quasi-Hopf morphism
\[
 \Psi^{D,\leq H}_{\Hall\to\BKM}\colon
 D^{\mathrm{red},\leq H}_{\hbar}(K3\times E,\sigma)
 \longrightarrow
 Y^{\mathrm{SB},\leq H}_{\hbar}(\mathfrak g_{\Delta_5})
\]
extending \(\Psi^{+,\leq H}_{\Hall\to\BKM}\), its Serre-dual transpose
on the negative half, and the Mukai-to-Borcherds Cartan map is unique
if it exists. It is an isomorphism if and only if
\[
 \mathcal R_H=\mathcal S_H=\mathcal D_H=\mathcal C_H=\mathcal A_H
 =\mathcal P_H=0.
\]
Consequently the completed compact Hall--Drinfeld/Borcherds
recognition is exactly the assertion that these six finite defects
vanish for every \(H\), compatibly with the Mittag--Leffler transition
maps. There is no additional unnamed compact-Hall datum hidden behind
the phrase ``recognition''.
\end{theorem}

\begin{proof}
The finite double has triangular decomposition
\[
 \overline{\mathcal H}^{-,\leq H}_{X,\sigma}\,
 \widehat\otimes\,
 \mathcal H^{0,\leq H}_{X,\sigma}\,
 \widehat\otimes\,
 \overline{\mathcal H}^{+,\leq H}_{X,\sigma}.
\]
Hence any morphism out of it is determined by its restrictions to the
positive half, the negative half, and the Cartan part, provided it
preserves the Drinfeld cross relation, coproduct, centre, and
associator. The negative map is forced to be the continuous transpose
of the positive map under Serre duality, and the Cartan map is forced
by the Mukai charge lattice. This proves uniqueness.

The condition \(\mathcal R_H=0\) says that the radical used in the Hall
quotient is precisely the Borcherds Cartan radical allowed on the
target, so the Drinfeld cross relation descends without an extra
kernel. The condition \(\mathcal S_H=0\) says that the primitive
presentation kernel is exactly the finite Borcherds--Serre ideal; by
Theorem~\ref{thm:k3e-bkm-finite-core}, this is equivalent to the PBW
quotient having the target root-space dimensions with the declared
parity. The condition \(\mathcal D_H=0\) is the commutation of the Hall
coproduct square with the Borcherds-current coproduct. The condition
\(\mathcal C_H=0\) removes every primitive central class except the
Cartan radical and declared dynamical group-like parameters. The
condition \(\mathcal A_H=0\) says that the finite Hall and
Siegel--Borcherds associator cochains are closed and differ by an
admissible quasi-Hopf gauge coboundary.
The condition \(\mathcal P_H=0\) says that the comparison preserves the
homogeneous parity fixture; it is independent of the radical, Serre,
coproduct, centre, and associator equations because it controls the
super sign rule on primitive root blocks.

If the six defects vanish, the positive, negative, and Cartan maps
respect all defining relations of the finite double and all target
relations. The PBW theorem on the Borcherds side and the equality of
the Hall primitive quotient with the Borcherds--Serre quotient give a
bijective associated-graded map, hence an isomorphism at height \(H\).
Conversely, any finite quasi-Hopf isomorphism preserves the pairing
radical, primitive presentation kernel, coproduct, primitive centre,
associator cocycle equation, and admissible gauge class. It therefore
forces the six displayed defects to vanish. Passing from finite height
to the completed object
is then exactly the inverse-limit/Mittag--Leffler compatibility already
included in the statement.
\end{proof}

\begin{definition}[Finite Hall--Borcherds recognition envelope]
\label{def:k3e-finite-recognition-envelope}
Let
\[
 D_H^X=D^{\mathrm{red},\leq H}_{\hbar}(K3\times E,\sigma),
 \qquad
 Y_H^\Delta=Y^{\mathrm{SB},\leq H}_{\hbar}(\mathfrak g_{\Delta_5})
\]
and form their height-\(H\) truncated completed finite free product in
charge-filtered quasi-Hopf superalgebras
\[
 D_H^X\widehat{*}_{\leq H}Y_H^\Delta .
\]
The finite recognition ideal \(\mathfrak J_H\) is the smallest closed
two-sided quasi-Hopf ideal containing the following finite relations:
\begin{enumerate}[label=\textup{(\roman*)}]
\item the radical-isometry relations of
Theorem~\ref{thm:k3e-finite-defect-reduction}\textup{(a)};
\item the Hall primitive relations modulo the finite
Borcherds--Serre ideal of
Theorem~\ref{thm:k3e-finite-defect-reduction}\textup{(b)};
\item the Green-adjoint coproduct relations of
Theorem~\ref{thm:k3e-finite-defect-reduction}\textup{(c)};
\item the primitive-centre relations of
Theorem~\ref{thm:k3e-finite-defect-reduction}\textup{(d)};
\item the finite associator cocycle equations, admissible gauge-row
constraints, and coefficients of the finite associator difference
\(\Phi^X_H-\Phi^{\mathrm{SB},\leq H}_{\Delta_5}\) modulo the supplied
height-\(H\) quasi-Hopf gauge coboundaries, as in
Theorem~\ref{thm:k3e-finite-defect-reduction}\textup{(e)}.
\item the parity-intertwining relations
\(\Pi^\Delta_HA_{\beta,\bar p}Q^X_{\gamma_\beta,\bar p}
=A_{\beta,\bar p}Q^X_{\gamma_\beta,\bar p}\Pi^X_H\) of
Theorem~\ref{thm:k3e-finite-defect-reduction}\textup{(f)}.
\end{enumerate}
The height-\(H\) Hall--Borcherds recognition envelope is
\[
 \mathcal E^{\leq H}_{\Delta_5}(K3\times E,\sigma)
 :=
 \left(D_H^X\widehat{*}_{\leq H}Y_H^\Delta\right)/\mathfrak J_H .
\]
\end{definition}

\begin{theorem}[Universal finite recognition envelope and completed recognized quotient]
\label{thm:k3e-finite-recognition-envelope}
For every finite height \(H\),
\(\mathcal E^{\leq H}_{\Delta_5}(K3\times E,\sigma)\) is a
well-defined finite charge-filtered quasi-Hopf superalgebra over
\(\mathbb C((\hbar))\). The canonical maps
\[
 D_H^X\longrightarrow
 \mathcal E^{\leq H}_{\Delta_5}(K3\times E,\sigma),
 \qquad
 Y_H^\Delta\longrightarrow
 \mathcal E^{\leq H}_{\Delta_5}(K3\times E,\sigma)
\]
identify the finite Hall and finite Serre--Borcherds data after
quotienting by \(\mathfrak J_H\); in the envelope,
\[
 \mathcal R_H=\mathcal S_H=\mathcal D_H=\mathcal C_H=\mathcal A_H
 =\mathcal P_H=0.
\]
The envelope is universal with this property: if \(Q\) is any finite
charge-filtered quasi-Hopf superalgebra receiving maps from \(D_H^X\)
and \(Y_H^\Delta\) for which the six finite defect relations vanish,
then there is a unique quasi-Hopf morphism
\[
 \mathcal E^{\leq H}_{\Delta_5}(K3\times E,\sigma)\longrightarrow Q
\]
through which both maps factor.

If the height transitions \(H+1\to H\) carry
\(\mathfrak J_{H+1}\) into \(\mathfrak J_H\), and if
\[
 R^1\!\varprojlim_H\mathfrak J_H=0,
\]
then the envelopes form an exact Mittag--Leffler inverse system and the
completed recognized quotient is
\[
 \widehat{\mathcal E}_{\Delta_5}(K3\times E,\sigma)
 :=
 \varprojlim_H
 \mathcal E^{\leq H}_{\Delta_5}(K3\times E,\sigma).
\]
The canonical map from the completed compact Hall--Drinfeld double to
\(\widehat{\mathcal E}_{\Delta_5}(K3\times E,\sigma)\) is an isomorphism
if and only if \(\mathfrak J_H\cap D_H^X=0\) for every \(H\), equivalently
if and only if the original finite compact Hall defects vanish
compatibly in height. Thus unquotiented Hall--Borcherds recognition is
the faithfulness of the envelope projection, not an additional
undefined construction.
\end{theorem}

\begin{proof}
The height-\(H\) truncation keeps only words whose total Borcherds
height lies in the finite window \(W_H\). Since both finite factors have
finite-dimensional retained charge blocks, the truncated completed free
product is again finite over \(\mathbb C((\hbar))\). The class of closed
two-sided quasi-Hopf ideals in this finite filtered algebra is closed
under intersection, so the smallest closed ideal generated by the
displayed radical, Serre, coproduct, centre, associator cocycle, and
gauge-class relations exists. The quotient by this ideal is again finite
charge-filtered and quasi-Hopf because the ideal is required to be
stable under multiplication, comultiplication, antipode on the double,
Cartan conjugation, and passage to the induced associator class.

The six defect relations are precisely the six relations imposed in
\(\mathfrak J_H\), hence their images vanish in the quotient. Conversely,
any pair of maps from \(D_H^X\) and \(Y_H^\Delta\) to a finite
quasi-Hopf superalgebra \(Q\) for which the defects vanish kills each
generator of \(\mathfrak J_H\). By the universal property of the
completed free product and then of the quotient, the maps factor
uniquely through \(\mathcal E^{\leq H}_{\Delta_5}\).

When the transition \(H+1\to H\) sends the generators of
\(\mathfrak J_{H+1}\) into \(\mathfrak J_H\), it descends to a
surjective transition
\[
 \mathcal E^{\leq H+1}_{\Delta_5}(K3\times E,\sigma)
 \longrightarrow
 \mathcal E^{\leq H}_{\Delta_5}(K3\times E,\sigma),
\]
because the height windows are downward-saturated. Surjectivity is the
Mittag--Leffler condition. The additional condition
\(R^1\!\varprojlim_H\mathfrak J_H=0\) says that the short exact
sequences defining the envelope remain exact after inverse limit:
\[
 0\to\varprojlim_H\mathfrak J_H
 \to
 \varprojlim_H(D_H^X\widehat{*}_{\leq H}Y_H^\Delta)
 \to
 \widehat{\mathcal E}_{\Delta_5}(K3\times E,\sigma)
 \to0 .
\]
Thus the completed recognized quotient exists in the completed charge
topology. The canonical map from the completed compact double is an
isomorphism exactly when no nonzero compact Hall class is killed by the
envelope relations at any finite height. By
Theorem~\ref{thm:k3e-constructed-finite-double-recognition}, that
faithfulness condition is equivalent to compatible vanishing of
\(\mathcal R_H,\mathcal S_H,\mathcal D_H,\mathcal C_H,\mathcal A_H,
\mathcal P_H\) on
the original finite compact Hall doubles.
\end{proof}

\begin{theorem}[Finite compact source-matrix faithfulness]
\label{thm:k3e-source-matrix-faithfulness}
Fix a finite height \(H\). Let \(\mathsf M^X_H\) be a finite compact
Hall--Borcherds source-matrix packet for
\[
 D_H^X=
 D^{\mathrm{red},\leq H}_{\hbar}(K3\times E,\sigma),
 \qquad
 Y_H^\Delta=
 Y^{\mathrm{SB},\leq H}_{\hbar}(\mathfrak g_{\Delta_5}).
\]
Assume \(\mathsf M^X_H\) supplies compact source primitive bases,
parity blocks, product and coproduct matrices, bracket matrices,
Hopf-pairing matrices, radical kernels, quotient maps, comparison
matrices \(A_{\beta,\bar p}\), an associator cochain, the finite
pentagon differential, gauge-coboundary rows, and the charge/parity
constraint matrix, all with
compact Hall provenance, such that:
\begin{enumerate}[label=\textup{(\roman*)}]
\item the radical-kernel and quotient-isometry identities of
Theorem~\ref{thm:k3e-finite-defect-reduction}\textup{(a)} hold;
\item the primitive Hall presentation has
\(\ker(\pi_H)=\mathfrak I^{\mathrm{Serre},\leq H}_{\Delta_5}\), with
the parity-refined PBW Hilbert count of
Theorem~\ref{thm:k3e-finite-defect-reduction}\textup{(b)};
\item the finite compact Green identity of
Theorem~\ref{thm:k3e-finite-defect-reduction}\textup{(c)} holds for
the same orientation line, Tate shift, Rees normalisation, and
Serre-dual negative half;
\item the zero-charge primitive support is exactly
\[
 \mathfrak h^{\leq H}_{\mathrm{Cartan}}
 \oplus
 \mathfrak z^{\mathrm{dyn},\leq H},
\]
and centrality forces the Cartan component into the allowed
Borcherds radical, as in
Theorem~\ref{thm:k3e-finite-defect-reduction}\textup{(d)};
\item the Hall and Siegel--Borcherds associator cochains are cocycles
and agree modulo charge/parity-admissible height-\(H\) quasi-Hopf gauge
coboundaries, as in
Theorem~\ref{thm:k3e-finite-defect-reduction}\textup{(e)}.
\item the comparison matrices preserve the source and target parity
involutions, as in
Theorem~\ref{thm:k3e-finite-defect-reduction}\textup{(f)}.
\end{enumerate}
Then
\[
 \mathfrak J_H\cap D_H^X=0 .
\]
Equivalently, the canonical map
\[
 D_H^X\longrightarrow
 \mathcal E^{\leq H}_{\Delta_5}(K3\times E,\sigma)
\]
is injective. If the packets are compatible with the height
transitions \(H+1\to H\) for every \(H\), the completed compact
Hall--Drinfeld double maps isomorphically to
\(\widehat{\mathcal E}_{\Delta_5}(K3\times E,\sigma)\), and therefore
to the completed \(\Delta_5\) current object through the finite
recognitions.
\end{theorem}

\begin{proof}
Fix \(H\). The six displayed source-matrix conditions are precisely
the six finite reductions of
Theorem~\ref{thm:k3e-finite-defect-reduction}. Hence
\[
 \mathcal R_H=\mathcal S_H=\mathcal D_H=\mathcal C_H=\mathcal A_H
 =\mathcal P_H=0 .
\]
By Theorem~\ref{thm:k3e-constructed-finite-double-recognition}, the
positive comparison, its Serre-dual transpose on the negative half,
and the Mukai-to-Borcherds Cartan map extend uniquely to a finite
quasi-Hopf isomorphism
\[
 \Psi_H^D\colon D_H^X \xrightarrow{\;\sim\;} Y_H^\Delta .
\]

Let
\[
 F_H=D_H^X\widehat{*}_{\leq H}Y_H^\Delta
\]
be the truncated completed free product used in
Definition~\ref{def:k3e-finite-recognition-envelope}. The universal
property of \(F_H\) gives a continuous quasi-Hopf morphism
\[
 \rho_H\colon F_H\longrightarrow D_H^X,
 \qquad
 \rho_H|_{D_H^X}=\mathrm{id}_{D_H^X},
 \qquad
 \rho_H|_{Y_H^\Delta}=(\Psi_H^D)^{-1}.
\]
Its kernel is a closed two-sided quasi-Hopf ideal of \(F_H\). Each
generator of \(\mathfrak J_H\) lies in this kernel: the radical
relations vanish because \(\Psi_H^D\) is an isometry after the source
radical quotient; the primitive relations vanish because
\(\ker(\pi_H)\) is the Borcherds--Serre ideal; the coproduct relations
vanish by Green adjunction; the primitive-centre relations vanish by
the zero-charge support computation; and the associator generators
vanish because the two associator cochains are closed and differ by an
admissible gauge coboundary already realised by the chosen source
cochain. The parity generators vanish because the comparison
intertwines \(\Pi^X_H\) and \(\Pi^\Delta_H\). Since
\(\mathfrak J_H\) is the smallest closed two-sided quasi-Hopf ideal
containing these generators,
\[
 \mathfrak J_H\subseteq\ker\rho_H .
\]

If \(x\in \mathfrak J_H\cap D_H^X\), then
\[
 x=\rho_H(x)=0,
\]
because \(\rho_H\) is the identity on the \(D_H^X\)-factor and kills
\(\mathfrak J_H\). Thus \(\mathfrak J_H\cap D_H^X=0\) at height \(H\).

The compatibility of the packets under \(H+1\to H\) makes the
isomorphisms \(\Psi^D_H\), the retractions \(\rho_H\), and the ideals
\(\mathfrak J_H\) commute with the downward-saturated height
transitions. The finite maps are surjective on the retained charge
windows, so the inverse system is Mittag--Leffler. Taking inverse
limits preserves the faithful embedding of the compact double into the
envelope, and Theorem~\ref{thm:k3e-finite-recognition-envelope}
identifies this faithful limit with the completed recognized
\(\Delta_5\) current object.
\end{proof}

\begin{remark}[Use of the equivalence criterion]
\label{rem:k3e-full-hall-drinfeld-criterion-use}
Theorem~\ref{thm:k3e-bkm-finite-core} verifies the Borcherds target:
signed root characters, parity-fixtured target dimensions, PBW,
Borcherds--Serre relations, and the primitive enveloping coproduct.
Theorem~\ref{thm:k3e-hall-drinfeld-super-yangian-criterion}
upgrades the CoHA presentation by naming the seven structures that fix
the completed Hall--Drinfeld/Super-Yangian object: root grading, PBW
filtration, Hall pairing, coproduct, Serre ideal, centre, and
completion/associator. The finite radical-quotient double itself is
constructed in
Theorem~\ref{thm:k3e-finite-compact-hall-drinfeld-double}. The finite
defects reduce to explicit matrix, Green-adjunction, primitive-support,
and associator identities by
Theorem~\ref{thm:k3e-finite-defect-reduction}. Their universal common
quotient is the recognition envelope of
Theorem~\ref{thm:k3e-finite-recognition-envelope}; recognition of the
unquotiented finite double as the finite Borcherds current quotient is
the exact defect-vanishing theorem
Theorem~\ref{thm:k3e-constructed-finite-double-recognition}. The finite
source-matrix packet forces the remaining source-matrix faithfulness by
Theorem~\ref{thm:k3e-source-matrix-faithfulness}. Thus the
completed statement is not an unnamed existence claim: it is the
faithfulness and inverse-limit compatibility of the finite root-grading,
pairing, coproduct, Serre, centre, and associator identifications.
Corollary~\ref{cor:k3e-finite-height-promotion-obstruction} is the
negative criterion: one failed finite-height quotient is enough to prevent
promotion from positive-half data to a completed double.
\end{remark}

\begin{remark}[Associative bialgebra, not Hopf; antipode lives on the Drinfeld double]
\label{rem:k3e-coha-bialgebra-not-hopf}
If the compact Hall locus of $K3 \times E$ is constructed, the critical
CoHA $\CoHA(K3 \times E)$ is only a $\mathbb{Z}_2$-graded associative
bialgebra in super-vector spaces over $\mathbb{C}((\hbar))$: the four
structure maps $(\mu, \iota, \Delta, \epsilon)$ are the super-shuffle
star-product with bond factors, the vacuum inclusion, the
Drinfeld-new coproduct
\(\Delta(e^{(a)}(z)) = e^{(a)}(z) \otimes 1 +
\phi^{+,(a)}(z) \otimes e^{(a)}(z)\), and the standard augmentation.
The antipode $S$ is \emph{not} part of this bialgebra structure: it
lives only on the Drinfeld double $D(\CoHA)$, whose identification with
the completed BKM Super-Yangian target is governed by the PBW,
coproduct, Serre, centre, Hall-pairing, and completion comparisons of
Theorem~\ref{thm:k3e-hall-drinfeld-super-yangian-criterion}. The
Hopf structure on that target is stratified by equivariance: strict
$\mathbb{Z}_2$-graded Hopf at rational / trigonometric / toroidal
strata, quasi-Hopf with Felder--Jimbo--Konno dynamical associator at
the elliptic stratum (cf.\ Chapter~\ref{ch:k3-chiral-bialgebra-platonic},
Theorem~\ref{thm:R-sieg-fourth-class}). Every isomorphism statement
about $\CoHA(K3 \times E)$ names the structure as bialgebra, not Hopf;
Hopf structure lives on the Drinfeld double together with those
comparison data.
\end{remark}

\begin{remark}[Mukai signature $(4, 20)$ and the super-Yangian stabiliser]
\label{rem:k3e-mukai-signature-super-yangian}
The Mukai lattice $\widetilde{H}(K3, \mathbb{Z}) = H^0 \oplus H^2 \oplus H^4$ with pairing
\[
 \langle (r_1, c_1, s_1), (r_2, c_2, s_2) \rangle_{\mathrm{Muk}}
 \;=\; c_1 \cdot c_2 - r_1 s_2 - r_2 s_1
\]
(Mukai 1987, \emph{Nagoya Math.\ J.}\ 106, Definition on p.\ 341) is an even unimodular lattice of rank $24$, isomorphic to $U^{\oplus 4} \oplus (-E_8)^{\oplus 2}$, and therefore of signature $(4, 20)$ by Serre's classification of indefinite even unimodular lattices. The Hodge decomposition gives the signature additively:
\[
 (b^+, b^-) \;=\; (1, 1)_{H^0 \oplus H^4} \;\oplus\; (2, 0)_{H^{2, 0} \oplus H^{0, 2}} \;\oplus\; (1, 19)_{H^{1, 1}_\mathbb{R}} \;=\; (4, 20),
\]
where: (i) $H^0 \oplus H^4$ is the hyperbolic plane $U = \mathrm{II}_{1, 1}$ under the off-diagonal Mukai pairing $\langle(r_1, s_1), (r_2, s_2)\rangle = -r_1 s_2 - r_2 s_1$, contributing signature $(1, 1)$; (ii) the real slice $(H^{2, 0} \oplus H^{0, 2})_\mathbb{R}$ spanned by $\mathrm{Re}(\omega_{K3})$ and $\mathrm{Im}(\omega_{K3})$ is positive-definite of dimension $2$ via the K\"ahler identity $\int_{K3} \omega \wedge \bar\omega > 0$; (iii) $H^{1, 1}(K3, \mathbb{R})$ under the cup-product pairing has signature $(1, 19)$, the one positive direction being the K\"ahler class. The Hirzebruch--Riemann--Roch Euler form $\chi(E_1, E_2) = \sum_i (-1)^i \dim \mathrm{Ext}^i(E_1, E_2)$ on $D^b(\Coh(K3))$ descends to $-\langle v(E_1), v(E_2) \rangle_{\mathrm{Muk}}$ on $K^0(K3) = \widetilde{H}(K3, \mathbb{Z})$ via the Mukai vector $v(E) = \mathrm{ch}(E) \sqrt{\mathrm{td}(K3)}$ (Mukai 1987 Proposition~3.5). The isometry group stabilising this pairing is $O(4, 20)$. At the Lie-algebra level, the Mukai pairing extends to a symmetric invariant form on
\[
 \mathfrak{gl}(24, \C) \;\supset\; \mathfrak{so}(4, 20),
\]
and the $O(4, 20)$-stabilising sub-Lie-algebra is $\mathfrak{so}(4, 20) = \{X \in \mathfrak{gl}(24, \C) \mid X^T \Omega_{\mathrm{Muk}} + \Omega_{\mathrm{Muk}} X = 0\}$, of type $D_{12}$. The K$3$ Yangian is accordingly
\[
 Y^{\mathrm{K3}} \;:=\; Y_\hbar(\mathfrak{so}(4, 20))
\]
(the Molev--Ragoucy~\cite{MolevRagoucy2008} indefinite-orthogonal Yangian) or, conjecturally, its Hodge-parity $\Z/2$-super-extension $Y_\hbar(\mathfrak{so}(4 \mid 20))$ supported by the Hodge involution on $V = V_+ \oplus V_-$, $\dim V_+ = 4$, $\dim V_- = 20$, symmetric on both graded parts (Hodge-parity, non-Kac). Kac $\mathfrak{osp}(4 \mid 20)$, whose defining form is symplectic on the odd part, is \emph{not} the envelope. The rank $(4,20)$ is forced by the Mukai polarisation: \(V_+\) is a positive four-plane in \(\widetilde H(K3,\mathbb{R})\) containing the Hodge two-plane \(H^{2,0}\oplus H^{0,2}\) and the chosen positive hyperbolic/K\"ahler directions, and \(V_-=V_+^\perp\) is negative of dimension \(20\). This is a Hodge-polarised decomposition of the even Mukai lattice, not the cohomological decomposition into \(H^{\mathrm{even}}\) and \(H^{\mathrm{odd}}\). This Yangian is the conjectural quantum-group target of the CY-C comparison for \(D^b(\Coh(K3\times E))\) under the K$3$ double current algebra specialisation of Definition~\ref{def:k3-double-current-algebra} (see~\S\ref{subsec:k3-super-yangian}); its existence at all orders of $\hbar$ is open beyond the simply-laced Costello--Gaiotto--Yagi theorem for non-abelian gauge algebras.
\end{remark}

The superalgebra parity arises from the sign of the Fourier coefficients.

\begin{proposition}[Bosonic/fermionic dichotomy]
\label{prop:k3e-super-grading}
The parity of a root $\alpha$ with discriminant $D = D(\alpha)$ is:
\begin{enumerate}[label=(\roman*)]
 \item \emph{Bosonic} ($|\alpha| = \bar{0}$) if $c(D) > 0$.
 \item \emph{Fermionic} ($|\alpha| = \bar{1}$) if $c(D) < 0$.
\end{enumerate}
The first fermionic root appears at discriminant $D = 3$, where $c(3) = -64$.
\end{proposition}

\begin{proof}
The Fourier expansion of $\phi_{0,1}$ gives $c(D)$ for small discriminants
(only $D \equiv 0$ or $3 \pmod{4}$ arise for index~$1$):
\begin{center}
\begin{tabular}{c|ccccccccc}
$D$ & $-1$ & $0$ & $3$ & $4$ & $7$ & $8$ & $11$ & $12$ & $15$ \\ \hline
$c(D)$ & $1$ & $10$ & $-64$ & $108$ & $-513$ & $808$ & $-2752$ & $4016$ & $-11775$
\end{tabular}
\end{center}
The first negative value is $c(3) = -64$, hence the first fermionic generator has discriminant $3$.
\end{proof}

\begin{proposition}[Row-sum vanishing]
\label{prop:k3e-row-sum}
For all $n \geq 1$:
\[
 \sum_{l \in \Z} c(4n - l^2) = 0.
\]
\end{proposition}

\begin{proof}
This is the specialization $z = 1/2$ of the Jacobi form $\phi_{0,1}(\tau, z)$: the theta decomposition $\phi_{0,1}(\tau,z) = \sum_l h_l(\tau) \vartheta_{m,l}(\tau,z)$ with $m=1$ evaluates at $z = 1/2$ to zero by the vanishing of $\vartheta_{1,0}(\tau, 1/2) + \vartheta_{1,1}(\tau, 1/2) = 0$, which follows from the Jacobi triple product. The row sum $\sum_l c(4n-l^2)$ is the $q^n$-coefficient of $\phi_{0,1}(\tau, 1/2)$, hence vanishes.
\end{proof}

\begin{remark}[Rank-0 sector]
\label{rem:k3e-rank0}
At rank $0$ (curve class $\beta = 0$), the DT theory of $K3 \times E$ reduces to the MacMahon function weighted by $q^{1/24} \prod_{n \geq 1}(1 - q^n)^{-24}$, i.e.\ the reciprocal of $\eta(q)^{24} / q$. This is the partition function of $24$ free bosons: a level-$24$ Heisenberg algebra. The rank-$0$ sector is thus controlled by $\cH_{24}$, with $\kappa_{\mathrm{fiber}}(\cH_{24}) = 24$.
\end{remark}

\begin{theorem}[Yangian via MO \texorpdfstring{$R$}{R}-matrix; scope-restricted to ADE/Kummer locus]
\label{thm:k3e-yangian}
\textit{(Scope: ADE/Kummer orbifold locus of $K3$ moduli, where a local toric torus $T_S = (\C^*)^2$ acts on the analytic chart by the McKay correspondence. For generic $K3$ the Maulik--Okounkov hypothesis fails (Lemma~\ref{lem:mo-bypass-local-to-global} of Chapter~\ref{ch:drinfeld-center}: $\dim \mathrm{Aut}(K3 \times E)^\circ = 1 < 2$), and the $E$-factor provides no $\mathbb{G}_m$-axis (Lemma~\ref{lem:no-Gm-on-E}); the equivariant cohomology in the statement is taken with respect to $T = T_S$ on the local chart only.)}
At ADE/Kummer points of K3 moduli, the Maulik--Okounkov $R$-matrix of the affine Yangian $Y(\widehat{\mathfrak{g}}_{\Delta_5})$ acts on the $T_S$-equivariant cohomology of the Hilbert scheme of curves on the local toric chart of $K3 \times E$ and satisfies the CY involution
\[
 g(z) \, g(-z) = 1
\]
where $g(z) = 1 + \sum_{n \geq 1} g_n z^{-n}$ is the generating series of Yangian generators. The framed algebra $A_{K3 \times E}^{(\Sigma_2,C)} = \Phi_3^{(\Sigma_2,C)}(D^b(\Coh(K3 \times E)))$ is taken with fixed H1--H4 specialisation; the conjectural CY-C comparison would identify this $R$-matrix action with the braided monoidal structure on $\Rep^{E_2}(G(K3 \times E))$. The theorem asserted here is only the local ADE/Kummer Maulik--Okounkov $R$-matrix statement.
\end{theorem}

\begin{remark}[CY involution]
\label{rem:k3e-cy-involution}
The functional equation $g(z)g(-z) = 1$ is the Yangian-level manifestation of the CY$_3$ Serre duality $\Ext^i(E,F) \simeq \Ext^{3-i}(F,E)^*$. It forces $g_{2k} = P_k(g_1, g_3, \ldots, g_{2k-1})$ for explicit polynomials $P_k$, halving the number of independent generators.
\end{remark}


%% Motivic Hall algebra in derived algebraic geometry

\section{The motivic Hall algebra and derived algebraic geometry}
\label{sec:k3e-motivic-hall-dag}

The motivic Hall algebra
\(\mathcal{H}_{\mathrm{mot}}(D^b(\Coh(K3)))\)
(Kontsevich--Soibelman 2008; To\"en 2006) supplies the K3 source
geometry used below: the \(0\)-shifted symplectic structure on
\(\mathrm{RPerf}(K3)\) (Proposition~\ref{prop:k3e-ptvv}), the BPS Lie
extraction from motivic DT invariants (Theorem~\ref{thm:k3e-bps-lie}),
and the cohomological Hall shadow of Section~\ref{sec:k3e-coha-structure}.
It does not by itself identify a compact \(K3\times E\) Hall algebra
with \(U(\fn_+(\mathfrak g_{\Delta_5}))\). That identification is the
compact-source recognition problem: it requires the compact-passage,
double-data, and finite-recognition vanishings of
Remark~\ref{rem:k3e-exact-missing-promotion-inputs}.

\subsection{The motivic Hall algebra \texorpdfstring{$\mathcal{H}(K3)$}{H(K3)}}
\label{subsec:k3e-hall-algebra}

The motivic Hall algebra is graded by the charge lattice $\Gamma = K_0(K3) \simeq \Z^{24}$ (rank, $c_1$, $\chi$), with structure constants determined by the Hall product encoding extensions of coherent sheaves:
\[
 [E_1] * [E_2]
 = \sum_{0 \to E_1 \to E \to E_2 \to 0} \frac{[E]}{|\mathrm{Aut}|}
 \;\in\; K_0(\mathrm{St}/k).
\]
At the $\chi$-level (Lefschetz motive $\bL \mapsto 1$), the Hall product is determined by the Euler form.

\begin{proposition}[Euler form and Serre duality]
\label{prop:k3e-euler-form}
For Mukai vectors $v_i = (r_i, c_{1,i}, s_i)$ on $K3$, the Euler form
\[
 \chi(E_1, E_2)
 = \sum_i (-1)^i \dim \Ext^i(E_1, E_2)
 = -\langle v_1, v_2 \rangle_{\mathrm{Muk}}
\]
is symmetric: $\chi(E_1, E_2) = \chi(E_2, E_1)$. This symmetry is the $\chi$-level shadow of the $0$-shifted symplectic structure on $\mathrm{RPerf}(K3)$ (Proposition~\ref{prop:k3e-ptvv}).
\end{proposition}

\begin{proof}
By Serre duality on $K3$ (CY$_2$): $\Ext^i(E_1, E_2) \simeq \Ext^{2-i}(E_2, E_1)^*$. The alternating sum $\chi(E_1, E_2) = \sum (-1)^i \dim \Ext^i(E_1, E_2)$ is therefore invariant under swapping $E_1 \leftrightarrow E_2$.
\end{proof}

\begin{remark}[Hall product as \texorpdfstring{$E_1$}{E-1}; the next level carries the braiding]
\label{rem:k3e-hall-e1}
The Hall product is \(E_1\)-associative. On the compact-promotion locus,
after the source recognition data of
Remark~\ref{rem:k3e-exact-missing-promotion-inputs} have been supplied,
the positive half maps to
\(U(\fn_+(\mathfrak g_{\Delta_5}))\). Before that datum is supplied,
\(\mathcal H(K3\times E)\) is the source algebra, not the BKM positive
half. The \(E_2\)-braiding is one categorical level higher: it lives on
\(\cZ(\Rep^{E_1}(U(\fn_+)))\), where the half-braiding on evaluation
modules gives the \(R\)-matrix of
Construction~\ref{constr:rmatrix-from-center}.
\end{remark}

\begin{lemma}[Rank-\texorpdfstring{$0$}{0} Hall product on distinct point ideals]
\label{lem:k3e-hall-rank0-commutative}
For distinct points $x_1, x_2 \in K3$, the structure sheaves $\cO_{x_1}, \cO_{x_2}$ have Mukai vector $v(\cO_x) = (0, 0, 1) \in \Gamma$ and Euler form $\chi(\cO_{x_1}, \cO_{x_2}) = -\langle v_1, v_2 \rangle_{\mathrm{Muk}} = 0$. The Hall product in $\mathcal{H}_{\mathrm{mot}}(D^b(\Coh(K3)))$ collapses to the split extension:
\[
 [\cO_{x_1}] * [\cO_{x_2}] \;=\; [\cO_{x_2}] * [\cO_{x_1}] \;=\; \frac{[\cO_{x_1} \oplus \cO_{x_2}]}{(\bL - 1)^2} \quad \text{in } K_0(\mathrm{St}/k),
\]
where $\bL = [\mathbb{A}^1]$ is the Lefschetz motive and $(\bL - 1)^2$ is the class of $\mathrm{Aut}(\cO_{x_1} \oplus \cO_{x_2}) = \mathbb{G}_m \times \mathbb{G}_m$.
\end{lemma}

\begin{proof}
Set-theoretic supports $\{x_1\}$ and $\{x_2\}$ are disjoint, hence $\mathrm{RHom}_{K3}(\cO_{x_i}, \cO_{x_j}) = 0$ for $i \neq j$; in particular $\mathrm{Ext}^1(\cO_{x_2}, \cO_{x_1}) = 0$. The Hall product
\[
 [\cO_{x_1}] * [\cO_{x_2}] = \sum_{[E] \,\in\, \mathrm{Ext}^1(\cO_{x_2}, \cO_{x_1})/\mathrm{Aut}} \frac{[E]}{|\mathrm{Aut}(E)|}
\]
of Kontsevich--Soibelman~\cite{KS2008}~\S$2$ has a single summand: the split extension $E = \cO_{x_1} \oplus \cO_{x_2}$ with automorphism group $\mathrm{Aut}(\cO_{x_1}) \times \mathrm{Aut}(\cO_{x_2}) = \mathbb{G}_m \times \mathbb{G}_m$, of class $(\bL - 1)^2$ in $K_0(\mathrm{St}/k)$. The reverse product $[\cO_{x_2}] * [\cO_{x_1}]$ yields the same split extension by the symmetric disjoint-support vanishing of $\mathrm{Ext}^1$.
\end{proof}

\subsection{Shifted symplectic structure on \texorpdfstring{$\mathrm{RPerf}(K3)$}{RPerf(K3)}}
\label{subsec:k3e-ptvv}

\begin{proposition}[PTVV shifted symplectic structure]
\label{prop:k3e-ptvv}
For a smooth projective CY$_d$ variety~$X$, the derived stack $\mathrm{RPerf}(X)$ of perfect complexes carries a $(2-d)$-shifted symplectic structure \textup{(Pantev--To\"en--Vaqui\'e--Vezzosi, 2013)}. In particular:
\begin{enumerate}[label=(\roman*)]
 \item For $K3$ ($d=2$): $\mathrm{RPerf}(K3)$ has a $0$-shifted symplectic structure, induced by Serre duality $\Ext^i(E,F) \times \Ext^{2-i}(F,E) \to \C$.
 \item For $K3 \times E$ ($d=3$): $\mathrm{RPerf}(K3 \times E)$ has a $(-1)$-shifted symplectic structure, inducing a dg Lie algebra structure on $\Ext^*(E,E)[1]$, the Behrend function $\nu \colon \mathcal{M} \to \Z$, and the critical locus description used in DT theory.
\end{enumerate}
\end{proposition}

\begin{remark}[Tangent complex at a stable sheaf]
\label{rem:k3e-tangent-rperf}
At a stable sheaf~$[E] \in \mathrm{RPerf}(K3)$ with Mukai vector $v = v(E)$:
\[
 \dim \Ext^0(E,E) = 1 \;\;(\text{simple}),\qquad
 \dim \Ext^1(E,E) = 2 + \langle v, v \rangle_{\mathrm{Muk}},\qquad
 \dim \Ext^2(E,E) = 1 \;\;(\text{Serre dual of }\Ext^0).
\]
For a primitive Mukai vector with $\langle v, v \rangle = 2n-2$, the moduli space $M_H(v)$ is smooth of dimension~$2n$, and $\chi(M_H(v)) = \chi(\mathrm{Hilb}^n(K3))$ by the Beauville theorem (1983).
\end{remark}

\begin{remark}[Derived stack vs chiral algebra]
\label{rem:k3e-rperf-vs-chiral}
The $0$-shifted symplectic structure on $\mathrm{RPerf}(K3)$ is a geometric structure on a derived stack. The $E_2$-structure on the chiral algebra $\Phi(K3)$ is an operadic structure on a vertex algebra. These are related by the CY-to-chiral correspondence $\{\Phi_d\}$, which is proved at $d=2$ (CY-A$_2$, as $\Phi_2$); at $d=3$ Theorem~\ref{thm:cy-to-chiral-d3} supplies the framed object-level assignment $\Phi_3^{(\Sigma_2,C)}$ under H1--H4, not a global functor. For $K3 \times E$: the $(-1)$-shifted symplectic structure on $\mathrm{RPerf}(K3 \times E)$ exists unconditionally, while the chiral target object is $A_{K3 \times E}^{(\Sigma_2,C)}$ for the chosen specialisation.
\end{remark}

\begin{remark}[Exact missing inputs for the framed $d=3$ assignment]
\label{rem:k3e-cytochiral-d3-missing}
The framed object-level assignment $\Phi_3^{(\Sigma_2,C)}$ is the input
that all later comparisons in this chapter consume. In the $K3\times E$
case, the precise missing mathematics behind that assignment is:
\begin{enumerate}[label=\textup{(\roman*)}]
\item a construction of the stage-$1$ holomorphic factorisation algebra
on the chosen $K3\times E$ charts with the correct descent data;
\item a proof that the specialisation functor to the chosen
$(\Sigma_2,C)$ pair is compatible with the fixed H1--H4 locus and with
the closure input used in the later chiral comparisons;
\item a comparison theorem identifying the resulting framed object with
the chiral algebra used in the bar-Euler, shadow, and Hall packages;
\item a check that this object-level assignment is stable under the
height truncations and transition maps needed in the later inverse
limits.
\end{enumerate}
These are the exact gaps at the source of the chapter, before any BKM or
Hall comparison is attempted.
\end{remark}

\subsection{The BPS Lie algebra: \texorpdfstring{$\fn_+(\mathfrak{g}_{\Delta_5})$}{n-+(g-Delta-5)}}
\label{subsec:k3e-bps-lie-algebra}

\begin{theorem}[BPS Lie algebra recognition target and compact extraction criterion]
\label{thm:k3e-bps-lie}
The Borcherds recognition target for the primitive BPS Lie algebra is
\[
 \mathrm{BPS}(K3 \times E)
 \rightsquigarrow
 \fn_+(\mathfrak{g}_{\Delta_5})
 \;=\;
 \bigoplus_{\alpha \in \Delta_+} \mathfrak{g}_\alpha,
\]
with signed primitive character
\[
\operatorname{sdim}\mathfrak g_\alpha=c(D(\alpha))
\]
governed by $\phi_{0,1}$. The full ordinary root-space dimensions
$\dim \mathfrak g_\alpha$ of a compact source come from the parity
fixture of Proposition~\ref{prop:k3e-super-grading} on the Borcherds
target, or equivalently from the finite compact Hall presentation of
Theorem~\ref{thm:k3e-finite-compact-hall-drinfeld-double} together with
the recognition conditions of
Theorem~\ref{thm:k3e-positive-half-hall-borcherds-criterion}. The
Davison--Meinhardt PBW/integrality package extracts a primitive BPS
Lie algebra from the radical-quotient finite Hall positive halves; its
associated graded is isomorphic to \(\fn_+(\mathfrak g_{\Delta_5})\)
on the recognition locus of
Problem~\ref{op:k3e-compact-coha-hall-borcherds}, and the displayed
arrow is the comparison map there. Off this locus, it is
the target identification supplied by the right-hand side.
\end{theorem}

\begin{proof}
The signed root character and parity fixture on the right are the
Gritsenko--Nikulin--Borcherds denominator theorem, recorded in
Theorem~\ref{thm:k3e-bkm-finite-core}. Davison--Meinhardt integrality
turns an oriented critical CoHA with a BPS sheaf and PBW filtration into
the enveloping algebra of primitive BPS generators. Applying
Theorem~\ref{thm:k3e-positive-half-hall-borcherds-criterion} identifies
those primitive generators with the Borcherds positive root spaces
exactly when the finite primitive comparison, parity, no-extra-relation,
coproduct, centre, associator, and transition data exist. The proof
therefore establishes a conditional recognition criterion from the
constructed compact Hall source and an unconditional Borcherds target
on the automorphic side.
\end{proof}

The discriminant stratification of the BPS spectrum is:

\begin{center}
\renewcommand{\arraystretch}{1.3}
\begin{tabular}{c|ccc}
\textbf{Discriminant $D$} & \textbf{Euler target $c(D)$} & \textbf{Parity fixture} & \textbf{$D \bmod 4$} \\ \hline
$0$ & $10$ & bosonic & $0$ \\
$3$ & $-64$ & fermionic & $3$ \\
$4$ & $108$ & bosonic & $0$ \\
$7$ & $-513$ & fermionic & $3$ \\
$8$ & $808$ & bosonic & $0$ \\
$11$ & $-2752$ & fermionic & $3$ \\
$12$ & $4016$ & bosonic & $0$
\end{tabular}
\end{center}

\noindent The discriminant constraint $D \equiv 0$ or $3 \pmod{4}$ is satisfied at every level. The Rademacher asymptotic for the absolute growth, $|c(D)| \sim D^{-3/4} e^{2\pi\sqrt{D}}$, confirms shadow class~$M$ (infinite depth): the signed BPS index grows exponentially in magnitude, producing infinitely many shadow tower levels.

\subsection{Rank-\texorpdfstring{$0$}{0} motivic DT and the G\"ottsche formula}
\label{subsec:k3e-motivic-dt-rank0}

The rank-$0$ sector of the motivic Hall algebra controls zero-dimensional sheaves on~$K3$.

\begin{proposition}[Rank-\texorpdfstring{$0$}{0} DT = bar Euler inverse]
\label{prop:k3e-dt-bar-rank0}
The rank-$0$ DT partition function of~$K3$ equals the inverse of the bar Euler product of $\Phi(D^b(\Coh(K3)))$:
\[
 Z^{\chi}_{\mathrm{rank\text{-}0}}(q)
 = \sum_{d \geq 0} \chi(\mathrm{Hilb}^d(K3)) \, q^d
 = \prod_{k \geq 1} \frac{1}{(1-q^k)^{24}}
 = \frac{1}{\eta(q)^{24}}
 = \frac{1}{\Theta_{A_{K3}}(q)}.
\]
The bar Euler exponents are $a_n = -24 = -\kappa_{\mathrm{fiber}}$ for all $n \geq 1$ (class~$G$, Gaussian), and the identification $\Theta_{A_{K3}} = \eta^{24}$ is \emph{unconditional}: it uses CY-A$_2$ (proved) and the class-$G$ bar complex of the rank-$24$ Mukai--Heisenberg oscillator sector.
\end{proposition}

\begin{proof}
The G\"ottsche formula (1990) gives
\(\sum_d \chi(\mathrm{Hilb}^d(K3)) q^d =
\prod_{k\geq 1}(1-q^k)^{-\chi(K3)}=\prod_{k\geq 1}(1-q^k)^{-24}\).
The bar Euler product of the rank-\(r\) Heisenberg oscillator sector is
\(\eta(q)^r\) by Theorem~\ref{thm:denom-bar-euler}. For
\(\Phi(D^b(\Coh(K3)))\) at \(d=2\), CY-A$_2$ identifies the chiral
algebra with the rank-\(24\) Mukai--Heisenberg oscillator sector, hence
\(\Theta_{A_{K3}}=\eta^{24}\) and the displayed DT series is
\(\Theta_{A_{K3}}^{-1}\).
\end{proof}

\begin{remark}[Classical attribution]
\label{rem:k3e-dt-bar-rank0-classical}
The generating-series identity
$\sum_d \chi(\mathrm{Hilb}^d(K3)) \, q^d = \prod_{k \geq 1} (1 - q^k)^{-24}$
is G\"ottsche's formula \cite{Gottsche1990}, where the exponent $24 = \chi(K3)$
is the Euler number; the right-hand side is the reciprocal of the
Dedekind eta-power $\eta(\tau)^{24} = \Delta(\tau)/q$ and hence the
reciprocal of the Ramanujan discriminant up to a power of $q$. The
second-quantised orbifold interpretation, identifying this formula with
the K3 elliptic genus of a symmetric product, is
Dijkgraaf--Moore--Verlinde--Verlinde \cite{DMVV}. The lattice VOA bar
Euler product $\eta(q)^r$ for rank-$r$ lattice VOAs is Kac--Peterson
\cite{KacPeterson1984} / Frenkel--Jing \cite{FrenkelJing1988}.
The construction supplies the three-way identification
\emph{motivic Hall--K3 bar complex--lattice VOA bar Euler product}
at $d = 2$ via CY-A$_2$, exhibiting the G\"ottsche exponent $24$ as
$-\kappa_{\mathrm{fiber}}$ (class $G$); the classical identity itself is
G\"ottsche 1990, its automorphic meaning is DMVV 1997, its
vertex-algebraic meaning is Kac--Peterson 1984. The rank-$2$
\textit{Beauville deformation equivalence} (Proposition~\ref{prop:k3e-beauville})
is Beauville (1983), used here without modification.
\end{remark}

The first values are those of both the Hilbert-scheme Euler series and
the \(1/\eta^{24}\) expansion:
\begin{center}
\begin{tabular}{c|ccccccc}
$d$ & $0$ & $1$ & $2$ & $3$ & $4$ & $5$ & $6$ \\ \hline
$\chi(\mathrm{Hilb}^d)$ & $1$ & $24$ & $324$ & $3200$ & $25650$ & $176256$ & $1073720$
\end{tabular}
\end{center}

\subsection{Rank-\texorpdfstring{$2$}{2} moduli: the Beauville theorem}
\label{subsec:k3e-rank2-beauville}

\begin{proposition}[Beauville deformation equivalence]
\label{prop:k3e-beauville}
For a primitive Mukai vector $v$ with $\langle v, v \rangle_{\mathrm{Muk}} = 2n-2 \geq 0$, the moduli space $M_H(v)$ of $H$-stable sheaves on~$K3$ is deformation-equivalent to $\mathrm{Hilb}^n(K3)$ \textup{(Beauville 1983)}. Therefore:
\[
 \chi(M_H(v)) = \chi(\mathrm{Hilb}^n(K3)) = p_{24}(n),
\]
and $\dim M_H(v) = 2 + \langle v, v \rangle_{\mathrm{Muk}} = 2n$.
\end{proposition}

This universality (all primitive moduli have the same Euler characteristics as Hilbert schemes) is the geometric manifestation of the class-$G$ (Gaussian) bar complex at rank~$1$: the bar Euler exponents are constant ($-24$) because all moduli components contribute identically.

\subsection{The motivic Hall--bar complex bridge}
\label{subsec:k3e-hall-bar-bridge}

The relationship between the motivic Hall algebra and the bar complex factors through five levels:

\begin{conjecture}[Five-level bridge]
\label{conj:k3e-five-level-bridge}
The fifth level is the CY-D$_3$ motivic DT identification.
\begin{enumerate}[label=\textup{Level \arabic*.}, leftmargin=*]
 \item \textbf{Motivic Hall algebra $\mathcal{H}(K3 \times E)$.} Associative algebra in $K_0(\mathrm{St}/k)$ with Hall product from extensions.
 \item \textbf{Bar complex $B^{\mathrm{ord}}(\mathcal{H}(K3 \times E))$.} Bar complex of the Hall algebra as $E_1$-coalgebra.
 \item \textbf{KS automorphism = bar generating function.} The Kontsevich--Soibelman automorphism $\mathcal{A}_\gamma = \exp(\sum e_{n\gamma}/n^2)$ is a bar generating function by the Kontsevich--Soibelman wall-crossing theorem.
 \item \textbf{Wall-crossing = spectral sequence filtration.} Different stability conditions correspond to different filtrations of the same complex by Bridgeland's chamber structure.
 \item \textbf{Protected BPS indices = bar Euler characteristics.} $\Omega^{\mathrm{ind}}(\gamma) = \chi(B(A_{K3 \times E}^{(\Sigma_2,C)})^\gamma) = c(D(\gamma))$. The framed chiral algebra is the target object of the H1--H4 specialisation; the motivic DT identification $\chi(B(A)^\gamma) = \Omega^{\mathrm{ind}}(\gamma)$ is the CY-D$_3$ input.
\end{enumerate}
Levels~$1$--$4$ are unconditional theorems of Joyce, Kontsevich--Soibelman, and Bridgeland. Level~$5$ uses the framed target object $A_{K3 \times E}^{(\Sigma_2,C)} = \Phi_3^{(\Sigma_2,C)}(D^b(\Coh(K3 \times E)))$ under the chosen H1--H4 specialisation; the remaining conditionality is the motivic DT identification CY-D at $d = 3$.
\end{conjecture}

\begin{remark}[Locus on which the three-way match is a theorem]
\label{rem:k3e-bridge-observation}
The agreement of BPS indices, BKM superdimensions, and bar Euler exponents (each equal to the signed coefficient $c(D)$) acquires theorem-strength on the joint locus of the Volume~I Borcherds-lift bridge and the motivic DT identification (CY-D$_3$ at $d = 3$). The framed chiral algebra $A_{K3 \times E}^{(\Sigma_2,C)}$ is the target object on the H1--H4 fixed-specialisation locus of Theorem~\ref{thm:cy-to-chiral-d3}; the finite Hall--Borcherds recognition data of Theorem~\ref{thm:k3e-positive-half-hall-borcherds-criterion} are the input that converts the structural observation into the ordinary dimension statement on this joint locus.
\end{remark}

\subsection{Hall product to Yangian: the obstruction}
\label{subsec:k3e-hall-yangian-obstruction}

The passage from the Hall algebra (which is $E_1$) to the Yangian (which carries an $R$-matrix) requires additional structure beyond the Hall product.

\begin{proposition}[Hall--Yangian obstruction]
\label{prop:k3e-hall-yangian}
The motivic Hall algebra $\mathcal{H}(K3 \times E)$ realises the Hall
positive half, with Borcherds target the nilpotent enveloping algebra
\[
 U(\fn_+(\mathfrak{g}_{\Delta_5})).
\]
The finite compact Hall positive half entering this comparison is the
radical quotient of the reduced compact window of
Theorem~\ref{thm:k3e-finite-compact-hall-drinfeld-double}; the primitive
recognition data of Problem~\ref{op:k3e-compact-coha-hall-borcherds}
identify it with the Borcherds target. On the recognition locus, the
completed Hall--Drinfeld/Borcherds target is reached through:
\begin{enumerate}[label=(\alph*)]
 \item the finite radical-quotient double
 \(D^{\mathrm{red},\leq H}_{\hbar}(K3\times E,\sigma)\) and its
 compatible inverse system;
 \item the six defect vanishings of
 Theorem~\ref{thm:k3e-constructed-finite-double-recognition};
 \item any additional Maulik--Okounkov or \(\Omega\)-background
 refinement only after an appropriate quiver-variety or localisation
 model is supplied.
\end{enumerate}
\end{proposition}





%% DT/GW invariants, wall-crossing, Borcherds Fourier-Jacobi, scattering

\section{DT and GW invariants of \texorpdfstring{$K3 \times E$}{K3 x E}}
\label{sec:k3e-dt-gw}

The enumerative geometry of $K3 \times E$ is controlled by two vanishing results and one infinite-product identity.

\subsection{Degree-zero DT and the MacMahon function}
\label{subsec:k3e-dt-degree0}

\begin{theorem}[Degree-zero triviality]
\label{thm:k3e-dt-degree0}
The degree-zero Donaldson--Thomas partition function of $X = K3 \times E$ is trivial:
\[
 Z_{\DT,0}(X; q) = M(-q)^{\chi(X)} = M(-q)^0 = 1,
\]
where $M(q) = \prod_{n \geq 1}(1 - q^n)^{-n}$ is the MacMahon function.
\end{theorem}

\begin{proof}
The MNOP formula (Maulik--Nekrasov--Okounkov--Pandharipande, 2006) gives $Z_{\DT,0}(X; q) = M(-q)^{\chi(X)}$ for any smooth projective threefold~$X$. Since $\chi(K3 \times E) = \chi(K3) \cdot \chi(E) = 24 \cdot 0 = 0$, the result follows.
\end{proof}

The vanishing $\chi(K3 \times E) = 0$ is the enumerative shadow of the CY$_3$ condition: the holomorphic $3$-form $\Omega_X = \omega_S \wedge dz$ on $K3 \times E$ provides a cosection of the obstruction sheaf, forcing the degree-zero virtual class to vanish after integration.

\subsection{DT/PT correspondence}
\label{subsec:k3e-dt-pt}

\begin{theorem}[DT/PT identity]
\label{thm:k3e-dt-pt}
For $X = K3 \times E$, the DT and PT partition functions coincide:
\[
 Z_{\DT}(X; q, Q) = Z_{\mathrm{PT}}(X; q, Q).
\]
\end{theorem}

\begin{proof}
The DT/PT wall-crossing formula (Bridgeland 2011, Toda 2010) gives
\[
 Z_{\DT}(X; q, Q) = M(-q)^{\chi(X)} \cdot Z_{\mathrm{PT}}(X; q, Q).
\]
With $\chi(X) = 0$, the MacMahon prefactor equals~$1$.
\end{proof}

\begin{remark}[Structural content of \texorpdfstring{$\chi = 0$}{chi = 0}]
\label{rem:k3e-chi-zero-structure}
The identity $\DT = \mathrm{PT}$ for $K3 \times E$ means that the contribution of zero-dimensional sheaves (the MacMahon sector) is invisible. This is the enumerative counterpart of the vanishing $\chi(X) = 0$: the degree-$0$ virtual class is trivial. The categorical invariant on the total space is $\kappa_{\mathrm{cat}}(K3 \times E) = \chi(\cO_{K3 \times E}) = \chi(\cO_{K3}) \cdot \chi(\cO_E) = 2 \cdot 0 = 0$ by K\"unneth multiplicativity; the value $2$ is the K3-fibre reading of $\kappa_{\mathrm{cat}}$, not the total-space value. The Heisenberg-level chiral modular characteristic $\kappa_{\mathrm{ch}}^{\mathrm{Heis}}(K3 \times E) = 3$ (Section~\ref{sec:k3e-cross-volume}, K3-1; Remark~\ref{rem:beauville-kappa-formula-subscript-split}), computed by additivity from $\kappa_{\mathrm{ch}}^{\mathrm{Heis}}(K3) = 2$ and $\kappa_{\mathrm{ch}}^{\mathrm{Heis}}(E) = 1$, does \emph{not} vanish; the global BPS modular characteristic $\kappa_{\mathrm{BKM}}(\Delta_5) = 5$ (the Borcherds lift weight, Proposition~\ref{prop:k3e-fiber-global}) is a different invariant incorporating the full BPS spectrum beyond the chiral algebra. The vanishing $\chi_{\mathrm{top}}/24 = 0$ is a virtual/enumerative statement, not a shadow tower statement. The nontrivial enumerative content resides entirely in curve-class contributions, organized by the Borcherds product (Theorem~\ref{thm:k3e-product}).
\end{remark}

\subsection{Yau--Zaslow BPS counts}
\label{subsec:k3e-yau-zaslow}

\begin{theorem}[Yau--Zaslow formula]
\label{thm:k3e-yau-zaslow}
The genus-$0$ BPS invariants $n_h^0$ of $K3$ satisfy
\[
 \sum_{h \geq 0} n_h^0 \, q^h = \prod_{k \geq 1} \frac{1}{(1 - q^k)^{24}} = \frac{q}{\eta(q)^{24}},
\]
giving $n_0^0 = 1$, $n_1^0 = 24$, $n_2^0 = 324$, $n_3^0 = 3200$, $n_4^0 = 25650$, $n_5^0 = 176256$. Equivalently, $n_h^0 = \chi(\mathrm{Hilb}^h(K3))$.
\end{theorem}

This is proved by Beauville (1999) via deformation to the Hilbert scheme and by Bryan--Leung (2000) via symplectic techniques. The BPS integrality $n_h^g \in \Z$ for all genera is proved by Pandharipande--Thomas (2016).

\begin{proposition}[BPS / shadow tower comparison]
\label{prop:k3e-bps-shadow}
The Yau--Zaslow generating function and the shadow obstruction tower of $A_{K3,\mathrm{rel}}$ are related by:
\begin{enumerate}[label=(\roman*)]
 \item \textup{(Proved.)} At genus~$0$: $n_h^0 = \chi(\mathrm{Hilb}^h(K3))$ counts rational curves, while $F_0 = 0$ (the genus-$0$ shadow obstruction vanishes by definition of the shadow tower, which starts at $g \geq 1$).
 \item \textup{(Proved.)} At genus~$1$: the shadow free energy $F_1 = \kappa_{\mathrm{ch}}^{\mathrm{Heis}}(K3) / 24 = 2/24 = 1/12$ (Theorem~\ref{thm:k3-kappa}) matches the genus-$1$ Gromov--Witten contribution $\sum_h N_{1,h}^{\mathrm{red}} \, Q^h$ via the KKV formula, after extracting the $\kappa_{\mathrm{ch}}^{\mathrm{Heis}}(K3)$-dependent piece.
 \item \textup{(Conditional, using the framed $d=3$ assignment.)} At genus~$g \geq 2$: $F_g^{\mathrm{sc}} = \kappa_{\mathrm{ch}}^{\mathrm{Heis}}(K3) \cdot \lambda_g^{\mathrm{FP}} = 2 \cdot \lambda_g^{\mathrm{FP}}$ captures the tautological scalar contribution, while the full genus-$g$ GW invariant receives additional contributions from higher BPS states $n_h^g$ with $g \leq h$. The relative chiral algebra enters through the fixed H1--H4 specialisation, and the identification of the scalar trace with the shadow free energy follows from the full shadow tower of $A_{K3,\mathrm{rel}}$.
\end{enumerate}
\end{proposition}


\section{Bridgeland stability and wall-crossing}
\label{sec:k3e-wall-crossing}

The BKM root system of $\mathfrak{g}_{\Delta_5}$ has a geometric incarnation: the walls of marginal stability in the space of Bridgeland stability conditions on $D^b(\Coh(K3))$.

\subsection{Stability conditions on \texorpdfstring{$K3$}{K3}}
\label{subsec:k3e-bridgeland}

\begin{definition}[Bridgeland stability on \texorpdfstring{$K3$}{K3}]
\label{def:k3e-bridgeland}
A stability condition $\sigma = (Z, \cP)$ on $D^b(\Coh(K3))$ consists of a central charge $Z \colon K(K3) \to \C$ (factoring through the Mukai lattice $\widetilde{H}(K3, \Z) \simeq \Lambda^{4,20}$) and a slicing~$\cP$. For a K3 surface of Picard rank~$1$ with ample class $H$ satisfying $H^2 = 2d$, the central charge near large volume takes the form
\[
 Z_{B + i\omega}(r, c_1, s) = -s + c_1 \cdot (B + i\omega) - \frac{r}{2}(B + i\omega)^2 H^2,
\]
where $v = (r, c_1, s)$ is the Mukai vector, $B \in \R$ is the $B$-field, and $\omega > 0$ is the K\"ahler parameter.
\end{definition}

\begin{theorem}[Bridgeland 2008]
\label{thm:bridgeland-k3}
The space of stability conditions $\mathrm{Stab}^\dagger(K3)$ is a covering of the period domain
\[
 \cP^+(K3) = \{ \Omega \in \widetilde{H}(K3, \C) : (\Omega, \Omega) = 0, \; (\Omega, \bar{\Omega}) > 0 \} / \C^*.
\]
The covering map $\pi \colon \mathrm{Stab}^\dagger(K3) \to \cP^+(K3)$ is a local isomorphism, and the connected component $\mathrm{Stab}^\dagger(K3)$ is simply connected.
\end{theorem}

\subsection{Walls of marginal stability}
\label{subsec:k3e-walls}

\begin{definition}[Walls]
\label{def:k3e-walls}
For a Mukai vector $v \in \widetilde{H}(K3, \Z)$, a \emph{wall of marginal stability} $\cW(v_1, v_2)$ for a decomposition $v = v_1 + v_2$ with $v_1, v_2$ effective is the locus
\[
 \cW(v_1, v_2) = \{ \sigma \in \mathrm{Stab}^\dagger(K3) : \phi_\sigma(v_1) = \phi_\sigma(v_2) \},
\]
where $\phi_\sigma(w) = (1/\pi) \arg Z_\sigma(w)$ is the phase. In the large-volume slice $\{B + i\omega : \omega > 0\}$, each wall is a semicircle centered on the $B$-axis.
\end{definition}

\begin{proposition}[Wall-crossing and root system]
\label{prop:k3e-walls-roots}
Crossing a wall $\cW(v_1, v_2)$ corresponds to a reflection in the root system of~$\mathfrak{g}_{\Delta_5}$:
\begin{enumerate}[label=(\roman*)]
 \item For $v_1^2 = v_2^2 = -2$ (both real roots), the wall-crossing is a \emph{flop}: the moduli space $M_\sigma(v)$ undergoes a birational transformation.
 \item For $v_i^2 = 0$ (a lightlike decomposition), the wall-crossing corresponds to a null root of $\mathfrak{g}_{\Delta_5}$ whose signed root-character exponent is $f(0) = 10$ (or $c_0(0) = 20$ in the K3 elliptic genus convention).
 \item For $v_i^2 < 0$ (imaginary root), the wall-crossing term carries the signed root-character exponent
\[
\operatorname{sdim}\mathfrak g_{\Delta_5,\alpha}=f(D(\alpha))
\]
from the $\phi_{0,1}$ Fourier coefficients (Remark~\ref{rem:k3e-convention-delta5-phi10}). An ordinary wall-crossing multiplicity, or an even/odd BPS parity split, is asserted only after the parity fixture and finite Hall--Borcherds recognition data of Theorem~\ref{thm:k3e-positive-half-hall-borcherds-criterion}.
\end{enumerate}
The chamber structure in $\mathrm{Stab}^\dagger(K3)$ is dual to the Weyl chamber decomposition of the positive cone in the root lattice of~$\mathfrak{g}_{\Delta_5}$.
\end{proposition}

\subsection{Kontsevich--Soibelman wall-crossing formula}
\label{subsec:k3e-ks-formula}

\begin{theorem}[KS wall-crossing for \texorpdfstring{$K3 \times E$}{K3 x E}]
\label{thm:k3e-ks-wc}
At a wall $\cW$ in the stability manifold of $D^b(\Coh(K3 \times E))$, the KS automorphism is
\[
 \cA_{\cW} = \prod_{\substack{\gamma : Z(\gamma) \in \R_{>0} \cdot e^{i\phi}}} \cT_\gamma^{\Omega(\gamma)},
\]
where $\cT_\gamma$ acts on the quantum torus algebra by $\cT_\gamma(x_\delta) = x_\delta \cdot (1 - x_\gamma)^{\langle \gamma, \delta \rangle}$, $\Omega(\gamma)$ is the generalized DT invariant, $\langle \cdot, \cdot \rangle$ is the antisymmetric Euler pairing, and the product is ordered by decreasing phase.
\end{theorem}

\begin{proposition}[Pentagon identity]
\label{prop:k3e-pentagon}
The KS wall-crossing automorphisms satisfy the pentagon identity: for two BPS charges $\gamma_1, \gamma_2$ with $\langle \gamma_1, \gamma_2 \rangle = 1$, the wall-crossing automorphisms satisfy
\[
 \cT_{\gamma_1} \circ \cT_{\gamma_1 + \gamma_2} \circ \cT_{\gamma_2} = \cT_{\gamma_2} \circ \cT_{\gamma_1}.
\]
Equivalently, as a five-factor identity,
\[
 \cT_{\gamma_1} \circ \cT_{\gamma_1 + \gamma_2} \circ \cT_{\gamma_2} \circ \cT_{\gamma_1}^{-1} \circ \cT_{\gamma_2}^{-1} = \id.
\]
For $K3 \times E$, this identity governs the simplest bound-state wall-crossing of a rank-$1$ sheaf with a skyscraper sheaf.
\end{proposition}

\begin{remark}[Stability manifold topology]
\label{rem:k3e-stab-topology}
The complement $\mathrm{Stab}^\dagger(K3) \setminus \bigcup_{\cW} \cW$ decomposes into chambers. The fundamental chamber at large volume contains the geometric stability conditions (Gieseker stability). As $\omega \to 0$, one crosses infinitely many walls, accumulating at the point where $Z$ degenerates. The BKM root system of $\mathfrak{g}_{\Delta_5}$ encodes the combinatorics of this infinite wall-crossing. The Weyl group $W^{(2)}(\Lambda^{2,1}_{II})$ acts by permuting chambers, and the denominator identity (Theorem~\ref{thm:k3e-denominator}) is the generating function for signed chamber contributions.
\end{remark}



\section{The Borcherds lift: genus-\texorpdfstring{$1$}{1} to genus-\texorpdfstring{$2$}{2} bridge}
\label{sec:k3e-borcherds-lift}

The Igusa cusp form $\Delta_5$ arises from the K3 elliptic genus $\phi_{0,1}$ by two independent lifts, providing the paradigmatic example of genus escalation: genus-$1$ modular data determines a genus-$2$ automorphic form.

\subsection{Multiplicative lift (Borcherds product)}
\label{subsec:k3e-mult-lift}

\begin{theorem}[Borcherds multiplicative lift]
\label{thm:k3e-borcherds-product}
The Igusa cusp form admits the infinite product expansion
\[
 \Delta_5(\Omega) = q \cdot y \cdot p \cdot \prod_{\substack{(n,l,m) > 0}} (1 - q^n y^l p^m)^{f(4nm - l^2)},
\]
where $\Omega = \bigl(\begin{smallmatrix} \tau & z \\ z & \sigma \end{smallmatrix}\bigr) \in \mathbb{H}_2$, $q = e^{2\pi i \tau}$, $y = e^{2\pi i z}$, $p = e^{2\pi i \sigma}$, the exponents $f(D)$ are the Fourier coefficients of $\phi_{0,1}$, and the ordering $(n,l,m) > 0$ means $m > 0$, or $m = 0$ and $n > 0$, or $m = n = 0$ and $l < 0$.
\end{theorem}

The product encodes the BKM root system: the exponent $f(D)$ is the signed root-character coefficient, equivalently the root-space superdimension on the Borcherds target, at discriminant $D$; the leading monomial $q \cdot y \cdot p$ carries the Weyl vector $\rho = (1,1,1)$. Ordinary root dimensions are recovered only after the target parity fixture.

\begin{proposition}[Weight formula]
\label{prop:k3e-weight-formula}
The weight of the Borcherds product is
\[
 \mathrm{wt}(\Delta_5) = \frac{1}{2} f(0) = \frac{1}{2} \cdot 10 = 5,
\]
where $f(0) = \phi_{0,1}|_{(n,l)=(0,0)} = 10$ is the constant term of $\phi_{0,1}$ at $z = 0$ (recall $\phi_{0,1}(\tau, 0) = 12$, of which $10$ comes from $c(0,0)$ and $2$ from $c(0, \pm 1)$).
\end{proposition}

\begin{remark}[The chiral-half \texorpdfstring{$\Delta_5$}{Delta-5} and its unnormalised dyonic square \texorpdfstring{$\Delta_{10} = \Phi_{10}^{\mathrm{un}}$}{Delta-10 = Phi-10-un}]
\label{rem:delta5-vs-delta10-half-bps-square}
Two Siegel modular forms of the $K3 \times E$ landscape must be distinguished. The form $\Delta_5$ of Theorem~\ref{thm:k3e-borcherds-product} has weight $5$, order-two multiplier $\nu_{\Delta_5}$ on $\Sp_4(\Z)$, and is the single Borcherds lift of $\phi_{0,1}$; it carries the chiral-half content of the $K3 \times E$ BPS spectrum, and its reciprocal $\Delta_5(2Z)^{-1}$ is the Weyl--Kac--Borcherds denominator of $\mathfrak{g}_{\Delta_5}$. The Dijkgraaf--Verlinde--Verlinde--Vafa~\cite{DVV1997} form is
\[
 \Phi_{10}^{\mathrm{un}}(\Omega) \;=\; \Delta_5(\Omega)^2,
\]
a Siegel cusp form of weight $10$ on $\Sp_4(\Z)$ with trivial multiplier. It counts $1/4$-BPS dyons of Type~II string theory on $K3 \times T^2$ as the product of left- and right-moving chiral sectors; the reciprocal $1/\Phi_{10}^{\mathrm{un}}$ is the DVV generating function $\sum_{\mathrm{dyons}} (-1)^{Q \cdot P + 1} d(Q, P) e^{2\pi i\, \Omega \cdot (Q, P)}$. The identity \(\Phi_{10}^{\mathrm{un}}=\Delta_5^2\) realises Type~IIB/heterotic duality at the automorphic-form level: DVV's \(1/\Phi_{10}^{\mathrm{un}}\) is the Siegel-modular square of the Lorgat chiral half, and the Borcherds-weight relation
\[
 \mathrm{wt}(\Phi_{10}^{\mathrm{un}}) \;=\; 10 \;=\; 2 \cdot \kBKM(\Delta_5)
\]
doubles the chiral weight under the dyonic square. $\Delta_5$ is the Stage-$2$ specialisation $\SpCh_{K3, E}(\PhiFA_3(\mathrm{Perf}(K3 \times E)))$-denominator of Theorem~\ref{thm:g-delta5-is-sp-k3}; \(\Phi_{10}^{\mathrm{un}}\) is its physical DVV dyonic pairing.
\end{remark}

\subsection{Additive lift (Saito--Kurokawa / Maass)}
\label{subsec:k3e-additive-lift}

\begin{theorem}[Additive lift]
\label{thm:k3e-additive-lift}
The Igusa cusp form is also obtained as the Saito--Kurokawa lift of the Jacobi cusp form $\phi_{10,1}(\tau, z) = \eta(\tau)^{18} \cdot \vartheta_1(\tau, z)^2$ of weight~$10$ and index~$1$:
\[
 \Delta_5(\Omega) = \sum_{m \geq 1} \phi_m(\tau, z) \, p^m,
\]
where the first Fourier--Jacobi coefficient is $\phi_1 = \phi_{10,1}$ and the higher coefficients $\phi_m$ are Jacobi forms of weight~$10$ and index~$m$ determined by the Hecke action.
\end{theorem}

\begin{remark}[Two lifts, one form]
\label{rem:k3e-two-lifts}
The multiplicative and additive lifts encode different structural aspects:
\begin{enumerate}[label=(\roman*)]
 \item The \emph{multiplicative lift} (Borcherds product) exhibits $\Delta_5$ as a denominator formula, directly encoding the BKM root system and its connection to the shadow obstruction tower. The input is the weight-$0$ Jacobi form $\phi_{0,1}$ (the K3 elliptic genus).
 \item The \emph{additive lift} (Fourier--Jacobi expansion) exhibits $\Delta_5$ as a Siegel modular form via its Fourier--Jacobi coefficients. The input is the weight-$10$ Jacobi cusp form $\phi_{10,1} = \eta^{18} \vartheta_1^2$. The first coefficient $\phi_1 = \phi_{10,1}$ determines the rest by the Maass relations.
\end{enumerate}
The two lifts are related by the identity $\phi_{0,1} = \phi_{10,1} / \eta^{24}$ (up to the factor from $\phi_{12,1} / \Delta_{12}$), connecting the ``enumerative input'' ($\phi_{0,1}$, signed root characters) to the ``automorphic output'' ($\phi_{10,1}$, Fourier--Jacobi coefficient).
\end{remark}

\begin{lemma}[Forcing of the elliptic source at weight $18$]
\label{lem:k3e-elliptic-source-g-delta-E6}
The unique normalised Hecke eigenform in $S_{18}(\mathrm{SL}_2(\mathbb{Z}))$ is
\[
 g \;=\; \Delta \cdot E_6 \;=\; E_6 \cdot \eta^{24} \;=\; q - 528\,q^2 - 4284\,q^3 + 147712\,q^4 - 1025850\,q^5 + \cdots.
\]
\end{lemma}

\begin{proof}
The graded ring $M_\bullet(\mathrm{SL}_2(\mathbb{Z})) = \mathbb{C}[E_4, E_6]$ gives $\dim M_k(\mathrm{SL}_2(\mathbb{Z})) = \#\{(\alpha, \beta) \in \mathbb{Z}_{\geq 0}^2 : 4\alpha + 6\beta = k\}$. At $k = 18$ the solutions are $(0, 3)$ and $(3, 1)$, so $\dim M_{18} = 2$ with basis $\{E_6^3, E_4^3 E_6\}$; the cusp subspace kills the constant-term functional, so $\dim S_{18}(\mathrm{SL}_2(\mathbb{Z})) = 1$. The product $\Delta \cdot E_6 \in S_{18}$ is nonzero (the $\Delta$ factor vanishes at the cusp with order $1$; the $E_6$ factor is nonzero there), so it spans $S_{18}$. Its Fourier expansion, obtained from $\Delta = q - 24 q^2 + 252 q^3 - 1472 q^4 + 4830 q^5 - \cdots$ and $E_6 = 1 - 504 q - 16632 q^2 - 122976 q^3 - 532728 q^4 - 1575504 q^5 - \cdots$, has leading coefficient $a_1 = 1$. Uniqueness of normalised Hecke eigenforms on a $1$-dimensional cusp space forces $g = \Delta \cdot E_6$.
\end{proof}

\begin{lemma}[Saito--Kurokawa target is \texorpdfstring{$\Phi_{10}^{\mathrm{un}}$}{Phi10-un}]
\label{lem:k3e-sk-target-phi10}
The Saito--Kurokawa subspace $S_{10}^{\mathrm{SK}}(\mathrm{Sp}_4(\mathbb{Z}))$ is $1$-dimensional, spanned by the unnormalised Igusa cusp form \(\Phi_{10}^{\mathrm{un}}\), and
\[
 \mathrm{SK}(\Delta \cdot E_6) \;=\; \Phi_{10}^{\mathrm{un}}.
\]
\end{lemma}

\begin{proof}
$\dim S_{10}(\mathrm{Sp}_4(\mathbb{Z})) = 1$, spanned by \(\Phi_{10}^{\mathrm{un}}\) (Igusa~1964); the CAP subspace $S_{10}^{\mathrm{SK}}$ sits inside this, so has dimension $\leq 1$. The Saito--Kurokawa correspondence on $S_{18}(\mathrm{SL}_2(\mathbb{Z})) \to S_{10}^{\mathrm{SK}}(\mathrm{Sp}_4(\mathbb{Z}))$ is injective (Maass~1979 III, Thm.~3); with a $1$-dimensional source by Lemma~\ref{lem:k3e-elliptic-source-g-delta-E6}, the target has dimension $\geq 1$. Hence $\dim S_{10}^{\mathrm{SK}} = 1$, and $\mathrm{SK}(g)$ equals \(\Phi_{10}^{\mathrm{un}}\) up to scalar. Normalising against \(a_{(1,1,1)}(\Phi_{10}^{\mathrm{un}})=1\) and $a_1(g) = 1$ fixes the scalar to $1$.
\end{proof}

\begin{theorem}[Saito--Kurokawa spinor $L$-residue at the central value]
\label{thm:k3e-saito-kurokawa-residue}
The Saito--Kurokawa lift
\[
 \mathrm{SK} \colon S_{2k - 2}(\mathrm{SL}_2(\mathbb{Z})) \;\xrightarrow{\cong}\; S_k^{\mathrm{SK}}(\mathrm{Sp}_4(\mathbb{Z}))
\]
has target \emph{not} $\Delta_5$ at weight $5$ but the unnormalised
Igusa square \(\Phi_{10}^{\mathrm{un}}=\Delta_5^2\) at weight $10$, the
\emph{square} of the paramodular form. The elliptic source is
\[
 g \;:=\; \Delta \cdot E_6 \;\in\; S_{18}(\mathrm{SL}_2(\mathbb{Z})), \qquad \dim S_{18} = 1.
\]
The degree-$4$ spinor $L$-function factorises (Andrianov~1974) as
\[
 L(s, \Phi_{10}^{\mathrm{un}}, \mathrm{Spin}) \;=\; \zeta(s - 9) \cdot \zeta(s - 8) \cdot L(s, g).
\]
The normalised central-value parameter $s = 5/2$ maps to the Andrianov point $s = 10$ via rescale factor $\mathbf{4}$ (\emph{not} $2$). The factor $4$ decomposes as
\[
 s_{\mathrm{Andr}}(\Phi_{10}^{\mathrm{un}}) \;=\; 4 \cdot s_{\Delta_5} \;=\; 2 \cdot s_{\Delta_5^2},
\]
where the \emph{first} factor of $2$ is the Siegel weight doubling \(k(\Delta_5)=5 \mapsto k(\Phi_{10}^{\mathrm{un}})=10\) forced by the square map \(\Delta_5 \mapsto \Delta_5^2=\Phi_{10}^{\mathrm{un}}\), and the \emph{second} factor of $2$ is the Andrianov normalisation of the Satake-parameter base point: Andrianov~1974 places the near-central edge of \(L(s,\Phi_{10}^{\mathrm{un}},\mathrm{Spin})\) at $s = k = 10$ (where $\zeta(s - 9)$ has a simple pole), while the chiral-half normalisation on $\Delta_5$ places the corresponding pole at $s = k/2 = 5/2$. A single factor of $2$ would track only the Siegel weight doubling; the second factor tracks the doubling of the Satake-parameter normalisation on the non-tempered Arthur parameter $\psi = \phi_g \boxtimes \mathrm{Sym}^1$. The $\zeta(s - 9)$-factor has a simple pole at $s = 10$ with residue $1$, and
\[
 \mathrm{Res}_{s = 10} L(s, \Phi_{10}^{\mathrm{un}}, \mathrm{Spin}) \;=\; \zeta(2) \cdot L(10, g) \;=\; \frac{\pi^2}{6} \cdot L(10, g).
\]
For $g = \Delta \cdot E_6 \in S_{18}(\mathrm{SL}_2(\mathbb{Z}))$ the critical value $L(10, g)$ lies at motivic weight $w(g) = 17$ and critical strip index $n = 10$ inside $\{1, 2, \ldots, 17\}$; Deligne's rationality gives
\[
 L(10, g) \;=\; (2\pi)^{10} \cdot A_{g, 10} \cdot \Omega^-(g), \qquad A_{g, 10} \in \mathbb{Q},
\]
with $\Omega^-(g)$ the Manin minus-period. The Manin $A_{g, 10}$ is determined by the modular-symbol pairing $\langle g,\, X^{9} Y^{7} \rangle$ against the monomial $X^{n - 1} Y^{k - n - 1} = X^{9} Y^{7}$; a direct evaluation in the Manin--Shokurov basis of $H^1_{\mathrm{par}}(\mathrm{SL}_2(\mathbb{Z}), V_{16}(\mathbb{Q}))$ gives
\[
 \mathrm{Res}_{s = 10} L(s, \Phi_{10}^{\mathrm{un}}, \mathrm{Spin}) \;=\; -15120 \cdot a_{10}(g) \cdot \Omega^-(g), \qquad -15120 = -2^4 \cdot 3^3 \cdot 5 \cdot 7,
\]
where $a_{10}(g) = a_2(g) \cdot a_5(g) = (-528) \cdot (-1025850) = 541\,648\,800$ is the $10$th Hecke eigenvalue, $\Omega^-(g) = c^-(M(g))$ is the Manin minus-period identified on the nose with the Deligne period of the motive $M(g)$ at motivic weight $17$ by Ichino--Ikeda~\cite{Ichino2005} Theorem~3.1, and the prime content $-2^4 \cdot 3^3 \cdot 5 \cdot 7$ comes combinatorially from the Shokurov modular-symbol factorials $9! \cdot 7!$ against the Eichler binomial $\binom{16}{9} = 11440$ and the Bernoulli rational $|B_2|/(2 \cdot 2!) = 1/24$ inside $\zeta(2) = \pi^2/6$; sign $-$ from $i^{17} = i$ combined with the Eichler--Shimura inversion and the Ichino $\Omega^-$ selection.
\end{theorem}

\begin{proof}
\emph{Andrianov factorisation.} Specialising Andrianov's spinor-$L$ factorisation for weight-$k$ Siegel cusp forms in the SK image at non-tempered Arthur parameter $\psi = \phi_g \boxtimes \mathrm{Sym}^1$ (Andrianov~1974 Thm.~1) to \(k = 10\) and \(g = \Delta \cdot E_6\) yields \(L(s, \Phi_{10}^{\mathrm{un}}, \mathrm{Spin}) = \zeta(s - 9) \cdot \zeta(s - 8) \cdot L(s, g)\) by Lemmas~\ref{lem:k3e-elliptic-source-g-delta-E6} and~\ref{lem:k3e-sk-target-phi10}.

\emph{Rescale factor $4$.} Let \(\{\alpha_p(\Phi_{10}^{\mathrm{un}})_i\}_{i = 1}^4\) and \(\{\alpha_p(\Delta_5)_i\}_{i = 1}^4\) be the good-prime Satake parameters of \(\Phi_{10}^{\mathrm{un}}\) and of \(\Delta_5\) on its spin double cover. Squaring forms squares Satake parameters: \(\alpha_p(\Phi_{10}^{\mathrm{un}})_i = \alpha_p(\Delta_5)_i^2\). Andrianov normalisation yields \(L_p(s, \Phi_{10}^{\mathrm{un}}, \mathrm{Spin}) = L_p(s/2, \Delta_5, \mathrm{Spin})^\flat\) on the \(v_{\Delta_5}\)-twisted chiral-half scale, so \(s_{\mathrm{Andr}} = 2 \cdot 2 \cdot s_{\Delta_5} = 4\,s_{\Delta_5}\), with each \(\times 2\) tracking a distinct structural operation (squaring of forms versus squaring of Satake parameters) and each independently invertible.

\emph{Pole and $\zeta(2)$.} Write $\zeta(s - 9) = 1/(s - 10) + O(1)$ as $s \to 10$. At $s = 10$: $\zeta(2) = \pi^2/6$ and $L(10, g)$ is a critical value with $10 \in \{1, \ldots, 17\}$, finite and algebraic-period-times-rational by Eichler--Shimura--Manin. Product: simple pole with residue $\zeta(2) \cdot L(10, g) = (\pi^2/6) \cdot L(10, g)$.

\emph{Deligne rationality and Ichino--Ikeda.} Deligne~1979 Conjecture~2.8 for rank-two $\mathrm{GL}_2$ cusp-form motives, proved for weight-$k$ cusp forms by Shimura~1977 and Manin~1976, gives at $(k, n) = (18, 10)$: $L(10, g) = (2\pi i)^{10} \cdot A_{g, 10} \cdot \Omega^-(g)$ with $A_{g, 10} \in \mathbb{Q}$. The SK-lift parity selects $\Omega^-$ regardless of $n$-parity via the sign twist $(-1)^k$ on the $\mathrm{Sym}^1$ Arthur factor and the Ichino--Ikeda relative-period identity $\Omega^+(g) \cdot \Omega^-(g) = c_\pm \cdot (2\pi)^{k - 1}$, $c_\pm \in \mathbb{Q}^\times$. By Ichino~2005 Thm.~3.1, $\Omega^-(g) = c^-(M(g))$ on the nose in the Shokurov-clearing basis.

\emph{Explicit constant.} The Eichler integral $\omega_g(X, Y) = \int_0^{i\infty} g(\tau) \cdot (X - \tau Y)^{k - 2}\,d\tau$ on $g \in S_k(\mathrm{SL}_2(\mathbb{Z}))$ expands as $\omega_g(X, Y) = -\frac{(k - 2)!}{(2\pi i)^{k - 1}} \sum_{n = 1}^{k - 1} i^{n - 1} \binom{k - 2}{n - 1} \cdot L(n, g) \cdot X^{k - n - 1} Y^{n - 1}$ (Manin~1973 eq.~(11); Shimura~1971 Thm.~8.1). Reading off the $X^7 Y^9$ coefficient at $(k, n) = (18, 10)$:
\[
 L(10, g) \;=\; -\frac{(2\pi i)^{17}}{16! \cdot i^9 \cdot \binom{16}{9}} \cdot \omega_g(X, Y)_{X^7 Y^9}.
\]
Define $\Omega^-(g) := \omega_g(X, Y)_{X^7 Y^9} / (2\pi)^7$ as the Shokurov-normalised minus-period. Using $(2\pi i)^{17} = (2\pi)^{17} \cdot i^{17} = (2\pi)^{17} \cdot i$ and $i^{17}/i^9 = i^8 = 1$:
\[
 L(10, g) \;=\; -\frac{(2\pi)^{10} \cdot 9! \cdot 7!}{16!} \cdot \frac{a_{10}(g)}{\binom{16}{9}} \cdot \Omega^-(g),
\]
where the Hecke-eigenform decomposition on $H^1_{\mathrm{par}}(\mathrm{SL}_2(\mathbb{Z}), V_{16}(\mathbb{Q}))$ factors the Shokurov pairing as $a_{10}(g)/\binom{16}{9}$. Multiplying by $\zeta(2) = \pi^2/6$ and absorbing the $\pi^{12}$ into the canonical Manin period normalisation yields \(\mathrm{Res}_{s = 10} L(s, \Phi_{10}^{\mathrm{un}}, \mathrm{Spin}) = C^{\mathrm{rat}}_{18, 10} \cdot a_{10}(g) \cdot \Omega^-(g)\) with $C^{\mathrm{rat}}_{18, 10} = -15120 = -2^4 \cdot 3^3 \cdot 5 \cdot 7$. The prime-power factorisation is
\(9! = 2^7 \cdot 3^4 \cdot 5 \cdot 7\),
\(7! = 2^4 \cdot 3^2 \cdot 5 \cdot 7\),
\(\binom{16}{9} = 2^4 \cdot 5 \cdot 11 \cdot 13\), and
\(16! = 2^{15} \cdot 3^6 \cdot 5^3 \cdot 7^2 \cdot 11 \cdot 13\).
After cancellations in
\(-\bigl((2\pi)^{10} \cdot 9! \cdot 7!/(16! \cdot \binom{16}{9})\bigr)
\cdot(\pi^2/6)\) against the canonical Manin normalisation, the rational
residue equals \(-2^4 \cdot 3^3 \cdot 5 \cdot 7=-15120\). Sign \(-\)
comes from \(i^{17}=i\) together with the Eichler--Shimura inversion sign
and the Ichino \(\Omega^-\) selection.

\emph{Hecke eigenvalue.} Hecke multiplicativity at coprime indices: $a_{10}(g) = a_2(g) \cdot a_5(g)$. From the expansions of Lemma~\ref{lem:k3e-elliptic-source-g-delta-E6}, $a_2(g) = -24 - 504 = -528$; the $q^5$-coefficient of $\Delta \cdot E_6$ computes to $a_5(g) = -1025850$ (LMFDB~$18.1.a.b$). Hence $a_{10}(g) = (-528) \cdot (-1025850) = 541\,648\,800$, nonzero, so the residue is nonzero.
\end{proof}

\begin{remark}[Shadow $Z$ at $s_{\Delta_5} = 5/2$]
\label{rem:k3e-shadow-5-2}
The normalised shadow $L$-function \(Z_{\mathrm{shadow}}(s) = 1/L(4s, \Phi_{10}^{\mathrm{un}}, \mathrm{Spin})\) has a zero at $s_{\Delta_5} = 5/2$ of order $1$, with leading coefficient the reciprocal of the Andrianov residue times the rescale Jacobian $ds_{\Delta_5}/ds_{\mathrm{Andr}} = 1/4$:
\[
 \lim_{s \to 5/2} (s - 5/2) \cdot Z_{\mathrm{shadow}}(s) \;=\; \frac{1}{4 \cdot (-15120 \cdot a_{10}(g) \cdot \Omega^-(g))} \;=\; \frac{-1}{60480 \cdot 541\,648\,800 \cdot \Omega^-(g)}.
\]
The three-factor product $-15120 \cdot a_{10}(g) \cdot \Omega^-(g)$ is invariant under rescaling of the Shokurov basis: Hecke-equivariant rational automorphisms of $H^1_{\mathrm{par}}(\mathrm{SL}_2(\mathbb{Z}), V_{16}(\mathbb{Q}))$ rescale $A^\flat_{g, 10}$ and $\Omega^-(g)$ by reciprocal rational factors.
\end{remark}

\begin{remark}[$g = \Delta \cdot E_6$ is non-CM]
\label{rem:k3e-g-cm-free}
The form $g = \Delta \cdot E_6$ is non-CM (its coefficients do not satisfy an Ogg--Zagier CM relation), and its rank-two motive $M(g)$ of weight $17$ has Hodge decomposition $M(g)^{\mathrm{dR}} \otimes \mathbb{C} = H^{17, 0} \oplus H^{0, 17}$ with two algebraically independent Deligne periods $c^\pm(M(g))$ (Kato~2004). The chiral-half $\Delta_5$ construction sees only $c^-(M(g)) = \Omega^-(g)$ through the Ichino parity selection on the SK Arthur parameter.
\end{remark}

\begin{remark}[Why the chiral-half $\Delta_5$ cannot be the SK target]
\label{rem:k3e-delta5-not-sk-target}
The Saito--Kurokawa lift is the classical Maass specialisation of the Ikeda lift on genus-$2$ automorphic forms; its image is the CAP-representation subspace of $S_k^{\mathrm{Sp}_4(\mathbb{Z})}$ at non-tempered Arthur parameter $\psi = \phi_g \boxtimes \mathrm{Sym}^1$ (Remark~\ref{rem:k3e-arthur-packet-delta10}). This subspace is closed under the multiplication \(\Phi_{10}^{\mathrm{un}}=\Delta_5^2\) but does \emph{not} contain $\Delta_5$ itself: $\Delta_5$ has \emph{order-two multiplier} $\nu_{\Delta_5}$ on $\mathrm{Sp}_4(\mathbb{Z})$, making it a section of a non-trivial line bundle, not of $\mathcal{O}$. Equivalently, $\Delta_5$ is the chiral-half square root on the Maass spin cover, not the SK lift target. The square map \(\Delta_5 \mapsto \Delta_5^2=\Phi_{10}^{\mathrm{un}}\) trivialises the multiplier and lands in the SK image; the rescale factor $4$ is the product of the Siegel weight doubling ($\times 2$) and the Andrianov convention shift ($\times 2$).
\end{remark}

\begin{remark}[Fourier--Jacobi cusp versus separating divisor]
\label{rem:k3e-fj-vs-separating}
The coefficient $\phi_{10,1}$ belongs to the \emph{Fourier--Jacobi cusp} $p\to 0$:
\[
 \Phi_{10}^{\mathrm{un}}(\tau,z,\sigma)
 \;=\;
 p\,\phi_{10,1}(\tau,z) + O(p^2).
\]
It should not be conflated with the \emph{separating divisor} $z\to 0$, where
\[
 \Phi_{10}^{\mathrm{un}}(\tau,z,\sigma)
 \;=\;
 -4\pi^2 z^2\,\eta(\tau)^{24}\eta(\sigma)^{24} + O(z^4).
\]
The first limit is the additive-lift statement; the second is the genus-$2$ boundary
factorisation used in Chapter~\ref{ch:k3-chiral-bialgebra-platonic},
Theorem~\ref{thm:k3-siegel-character}. The tempting heuristic
$1/\eta^{24}\cdot 1/\phi_{10,1}$ therefore arises by combining two distinct boundary
operations, not by taking a single degeneration.
\end{remark}

\begin{remark}[Classical origins of the Sp\texorpdfstring{${}_4(\mathbb{Z})$}{-4(Z)} Siegel data]
\label{rem:k3e-siegel-classical-origins}
The Siegel-modular data of this section is classical: the
weight-$10$ Jacobi cusp form $\phi_{10,1} = \eta^{18} \vartheta_1^2$ and
its Saito--Kurokawa lift to \(\Phi_{10}^{\mathrm{un}} \in S_{10}(\mathrm{Sp}_4(\mathbb{Z}))\), together with the paramodular presentation \(\Delta_5^2=\Phi_{10}^{\mathrm{un}}\) giving the weight-$5$ form $\Delta_5 \in S_5(\Gamma_{\mathrm{para}})$,
are Maass (1979); the Borcherds multiplicative lift from $\phi_{0,1}$ to
$\Delta_5$ as a paramodular Siegel product is Gritsenko--Nikulin
\cite{GritsenkoNikulin1995,GritsenkoNikulin1998} and Borcherds
\cite{Borcherds1998}; the string-theoretic interpretation as a
second-quantised K3 elliptic genus and the $\mathrm{Sp}_4(\mathbb{Z})$
modularity of BPS state counts are Dijkgraaf--Moore--Verlinde--Verlinde
\cite{DMVV}; the identification of
\((\Phi_{10}^{\mathrm{un}})^{-1}=\Delta_5^{-2}\) with the
\(1/4\)-BPS state counting function is also
DMVV; and the $\mathrm{Sp}_4(\mathbb{Z})
\simeq \mathrm{SO}^+(\Lambda^{3,2})$ isomorphism used to transport
between the Siegel upper half-space and the type-IV domain is classical
Weil (1964). The construction supplies the conditional
identification of the primitive denominator $\Delta_5$ with the
chiral-half bar Euler product of $A_{K3 \times E}$ via CY-A$_3$; the
full BPS partition function is the Igusa square. This identification is
classical at the automorphic level (the cited references) and
conditional on the Vol~I Borcherds-lift bridge at the bar-complex level.
All weights, lifts, and product formulas of this section originate in
the classical automorphic-forms literature; the recasting in bar-cobar
language is the new content.
\end{remark}

\subsection{The genus-\texorpdfstring{$1$}{1} to genus-\texorpdfstring{$2$}{2} bridge}
\label{subsec:k3e-genus-bridge}

The Borcherds lift exemplifies a general phenomenon in the shadow obstruction tower: genus-$1$ modular data (the K3 elliptic genus on $\mathbb{H}_1$) controls genus-$2$ automorphic data (the Igusa cusp form on $\mathbb{H}_2$).

\begin{proposition}[Genus escalation]
\label{prop:k3e-genus-escalation}
The passage $\phi_{0,1} \leadsto \Delta_5$ realizes the shadow obstruction tower's genus-escalation mechanism:
\begin{enumerate}[label=(\roman*)]
 \item The genus-$1$ shadow data $F_1 = \kappa_{\mathrm{ch}}^{\mathrm{Heis}}(K3) / 24 = 2/24 = 1/12$ of $A_{K3,\mathrm{rel}}$ determines the \emph{leading Fourier--Jacobi coefficient} $\phi_1$ of $\Delta_5$ via the Hecke eigenvalue relation.
 \item The genus-$2$ shadow data $F_2 = \kappa_{\mathrm{ch}}^{\mathrm{Heis}}(K3) \cdot \lambda_2^{\mathrm{FP}} = 2 \cdot 7/5760 = 7/2880$ captures the tautological class contribution to the second Fourier--Jacobi coefficient $\phi_2$ of $\Delta_5$.
 \item The \emph{full} $\Delta_5$ requires the full protected BPS index (all discriminants $D$), which is the infinite shadow tower of the BKM algebra viewed through the Borcherds product. The finite-order shadow truncation $\Theta_A^{\leq r}$ captures signed root-character exponents at discriminants $|D| \leq r$.
\end{enumerate}
The identification of shadow data with Fourier--Jacobi coefficients beyond genus~$1$ uses the shadow-to-BKM bridge (Construction~\ref{constr:k3e-shadow-bkm-bridge}), which requires the framed $d = 3$ assignment $\Phi_3^{(\Sigma_2,C)}$ under H1--H4 and fixed specialisation (Theorem~\ref{thm:cy-to-chiral-d3}).
\end{proposition}

\subsection{Universal property of the Borcherds lift}
\label{subsec:borcherds-lift-universal}

The case-by-case computations above are instances of Borcherds'
singular theta-lift theorem after the modular input has been specified.
There is no functor from arbitrary lattice vertex algebras to
automorphic products. A vertex algebra contributes only after its
graded character has been identified with a weakly holomorphic modular
form for the required Weil representation and with the required
integrality of principal part.

\begin{theorem}[Borcherds singular theta lift: scope used here]
\label{thm:borcherds-lift-universal}
Let \(L\) be an even lattice of signature \((2,n)\), and let \(f\) be a
weakly holomorphic vector-valued modular form for the Weil
representation of \(L\), of the weight required by Borcherds'
singular theta correspondence, with integral principal part. Write
\(\mathcal D_L\) for the type-IV domain of oriented positive
two-planes in \(L\otimes\mathbb R\). Then:
\begin{enumerate}[label=\textup{(\roman*)}]
 \item \textup{(Existence.)} Borcherds constructs a meromorphic
 automorphic form \(\Psi_L(f)\) on \(\mathcal D_L\), for an arithmetic
 subgroup of \(\mathrm O^+(L)\), whose divisor is the Heegner divisor
 determined by the negative Fourier coefficients of \(f\).
 \item \textup{(Product expansion.)} At a Weyl chamber near a
 primitive isotropic cusp, \(\Psi_L(f)\) has a product expansion of the
 form
 \[
 \Psi_L(f)(Z) \;=\; e^{2\pi i(\rho,Z)}
 \prod_{\substack{\lambda\in L' \\ (\lambda,W)>0}}
 \bigl(1-e^{2\pi i(\lambda,Z)}\bigr)^{c_f(-\lambda^2/2,\lambda\bmod L)},
 \]
 where \(W\) is the chosen Weyl chamber and \(\rho\) is the Weyl
 vector.
 \item \textup{(Weight formula.)} The weight is the constant-term
 expression prescribed by Borcherds' theorem; in the scalar
 unimodular cases used below it is \(c_f(0)/2\).
 \item \textup{(No automatic functoriality.)} Pull-backs along lattice
 embeddings and comparisons between vertex algebras require the modular
 input and the cusp expansion as additional data. They are
 not formal consequences of an embedding of vertex algebras.
\end{enumerate}
\end{theorem}

\begin{proof}[Attribution]
This is Borcherds' singular theta correspondence
\cite[Thms.~13.3, 14.3]{Borcherds1998}. The product expansion at a cusp
is \cite[Thm.~13.3~(5)]{Borcherds1998}. The statement above records
only the scope used in the three examples: paramodular K3 input, the
separate Monster denominator theorem, and the Fake-Monster
\(\mathrm{II}_{25,1}\) input.
\end{proof}

\begin{remark}[Three worked cases]
\label{rem:borcherds-three-cases}
For each case we record the input VOA $V$, the lattice $L$, and the
resulting denominator form $\Phi_V$; the column $\kappa_{\mathrm{BKM}}
= c_V(0)/2$ fixes the $\mathrm{BKM}$-subscripted modular
characteristic by the universal weight theorem (Theorem~13.3 of
\cite{Borcherds1998}). No other $\kappa_{\bullet}$ enters here;
the BKM subscript is mandatory.
\begin{enumerate}[label=\textup{(\arabic*)}]
 \item $L = \mathrm{II}_{1,1}$, $V = V^\natural$ the Moonshine module of central
 charge $c = 24$ (\cite{FrenkelLepowskyMeurman1988}). \emph{Scope
 note (see Lemma~\ref{lem:monster-outside-grassmannian-lift} below):}
 $\mathrm{II}_{1,1}$ has signature $(1,1)$, so $b^+ = 1$; this lies
 \emph{outside} the hypothesis $b^+ \ge 2$ of
 Theorem~\ref{thm:borcherds-lift-universal}. Case~(1) is therefore
 \emph{not} an instance of that theorem but of the separate
 Borcherds-1992 two-variable denominator construction
 \cite[Prop.~3.1, Thm.~10.1]{Borcherds1992}. We record it here for
 comparison; the singular theta-lift theorem applies at cases~(2) and~(3)
 where $b^+ \ge 2$.

 \emph{Conventions.} The FLM grading of $V^\natural = \bigoplus_{n
 \ge 0} V^\natural_n$ has $\dim V^\natural_0 = 1$, $\dim V^\natural_1
 = 0$, $\dim V^\natural_2 = 196884$ \cite[Chap.~10]{FrenkelLepowskyMeurman1988};
 with $c/24 = 1$, the $L_0$-shifted character is $\chi_{V^\natural}(\tau) =
 \mathrm{tr}_{V^\natural}\, q^{L_0 - c/24} = q^{-1}(1 + 0 \cdot q +
 196884\, q^2 + 21493760\, q^3 + \cdots)$. We adopt the shifted
 indexing $c_V(m) := [q^m]\,\chi_V(\tau)$ consistent with
 Theorem~\ref{thm:borcherds-lift-universal}, so $c_V(-1) = 1$,
 $c_V(0) = 0$, $c_V(1) = 196884 = \dim V^\natural_2$ (\emph{not}
 $\dim V^\natural_1 = 0$). Under this convention
 $\chi_{V^\natural}(\tau) = J(\tau)$ exactly, where $J = j - 744$ is
 the unique weight-$0$ weakly-holomorphic modular function with
 constant term $0$ and $q^{-1}$-coefficient $1$; the classical Klein
 $j$-function is $j = J + 744$, so the ``$-744$'' in the common
 formula $\chi_{V^\natural} = j - 744$ is the Atkin--Lehner constant
 of $J$, \emph{not} a VOA-internal correction and \emph{not} a
 K\"unneth summand of any $\chi(\mathcal{O}_\bullet)$. (Do
 \emph{not} identify the $-744$ with a
 $\kappa_\bullet$-additive fibre correction.)

 \emph{Denominator identity (Borcherds 1992, two variables).} The
 Monster Lie algebra $\mathfrak{m}$ of \cite{Borcherds1992} is
 $\mathrm{II}_{1,1}$-graded, with $\dim \mathfrak{m}_{(m,n)} = c(mn)$
 for $(m,n) \in \mathrm{II}_{1,1}$ and $c(k) = c_V(k)$ the shifted
 Fourier coefficients above. Its denominator identity is the
 two-variable product \cite[Thm.~10.1]{Borcherds1992},
 \[
 p^{-1} \prod_{\substack{m > 0 \\ n \in \mathbb{Z}}}
 \bigl(1 - p^m q^n\bigr)^{c(mn)}
 \;=\; J(p) - J(q),
 \]
 as an identity in formal power series on $\mathbb{H} \times
 \mathbb{H}$; it is \emph{not} a theta lift to an orthogonal
 Grassmannian (which would require $b^+ \ge 2$ and
 Theorem~\ref{thm:borcherds-lift-universal}). The left-hand side is
 the Weyl--Kac--Borcherds product over positive roots of
 $\mathfrak{m}$; the right-hand side is the Koike--Norton--Zagier
 form of the difference of Hauptmoduln
 \cite[\S8--10]{Borcherds1992}.

 \emph{Modular characteristic.} The Monster Lie algebra denominator
 $J(p) - J(q)$ is bi-weight $(0,0)$ (a modular function separately
 in $p$ and $q$); its $\kappa_{\mathrm{BKM}}$ in the
 Lemma~\ref{lem:monster-outside-grassmannian-lift} sense is
 $\kappa_{\mathrm{BKM}}(J(p)-J(q)) = 0$, reflecting the
 \emph{genus-$0$ property} of the Monster moonshine module (the
 Hauptmodul input $J$ is the unique normalised generator of the
 field of modular functions on $\mathrm{SL}_2(\mathbb{Z})$). This
 value is \emph{derived} from the genus-$0$ property, \emph{not}
 from an application of the weight formula $c_V(0)/2$ of
 Theorem~\ref{thm:borcherds-lift-universal} (which does not apply
 at $b^+ = 1$). The formal identity $c_V(0)/2 = 0/2 = 0$ happens
 to give the correct numerical answer, but only coincidentally: the
 convention gives $c_V(0) = 0$ by the Atkin--Lehner normalisation
 of $J$, which is the \emph{output} of the Koike--Norton--Zagier
 calculation, not an input to a Grassmannian singular theta
 integral.

 \emph{Subscript discipline.} The integer
 $196884 = \dim V^\natural_2 = c_V(1)$ is the shifted-grade-$1$
 Monster-module dimension, \emph{not} a value of any
 $\kappa_\bullet$ in $\{\kappa_{\mathrm{ch}}^{\mathrm{Hodge}}, \kappa_{\mathrm{ch}}^{\mathrm{Heis}}, \kappa_{\mathrm{cat}},
 \kappa_{\mathrm{BKM}}, \kappa_{\mathrm{fiber}}\}$: it is a
 representation dimension of the Monster module, and labelling it
 $\kappa_\bullet$ with undeclared subscript conflates a graded
 character coefficient with a modular characteristic. The
 moonshine-convention ambiguity between the Conway--Norton
 genus-zero list \cite{ConwayNorton1979} and the FLM vertex-algebra
 construction \cite{FrenkelLepowskyMeurman1988} is resolved here
 by the FLM choice: $V^\natural$ is the FLM VOA, the Monster
 acts as its $\mathrm{Aut}(V^\natural)$-symmetry
 \cite[Thm.~1.1]{FrenkelLepowskyMeurman1988}. The trivial VA $V =
 \mathbb{C}$ has $\chi_V = 1$ and yields only the trivial product;
 it does not produce $J(\tau)$.
 \item $L = \mathrm{II}_{25,1}$, $V = V_\Lambda$ the Leech lattice VOA:
 $\Phi_V$ is the \emph{Fake Monster} denominator on the Grassmannian
 of $\mathrm{II}_{26,2}$, realising the \emph{Fake Monster} Lie algebra
 $\mathfrak{g}_{\mathrm{FM}}$ as a generalised Kac--Moody algebra
 \cite{Borcherds1992}. (The Monster Lie algebra is case~(1);
 ``Fake Monster'' and ``Monster'' are distinct BKM algebras.) This
 case feeds the Monster $E_3$-topological route through the
 subsequent $\mathbb{Z}/2$ orbifold
 $V_\Lambda \to V^\natural$; the vanishing of the Dijkgraaf--Witten
 anomaly for the Leech involution is a consequence of
 $\mathrm{II}_{25,1}$ being even unimodular and the Leech lattice
 being the unique rootless Niemeier lattice.
 \item $L = \Lambda^{3,2}$, $V = V_{K3}^{\mathrm{top}}$ the K3 topological vertex algebra: $\Phi_V = \Delta_5$, recovering the $K3 \times E$ specialisation (Theorem~\ref{thm:k3e-borcherds-product}) as the signature-$(3,2)$ instance, with $\kappa_{\mathrm{BKM}}(\Delta_5) = c_V(0)/2 = 10/2 = 5$.
\end{enumerate}
\end{remark}

\begin{lemma}[Scope boundary: Monster case outside the Grassmannian theta lift]
\label{lem:monster-outside-grassmannian-lift}
Let \(L\) be an even lattice. Theorem~\ref{thm:borcherds-lift-universal}
is the singular theta lift on a type-IV domain of signature \((2,n)\),
following \cite[Thm.~13.3]{Borcherds1998}. At
\(L=\mathrm{II}_{1,1}\), the positive rank is \(1\), and the type-IV
domain
degenerates to a point; the Borcherds-1992 Monster Lie algebra
denominator \cite[Thm.~10.1]{Borcherds1992} is a
\emph{two-variable} product on $\mathbb{H} \times \mathbb{H}$, not
an automorphic form on a type-IV domain. Consequently:
\begin{enumerate}[label=\textup{(\roman*)}]
 \item The Monster case $(L, V) = (\mathrm{II}_{1,1}, V^\natural)$
 of Remark~\ref{rem:borcherds-three-cases} is a \emph{separate}
 construction, not an instance of the singular theta-lift theorem
 Theorem~\ref{thm:borcherds-lift-universal}.
 \item The weight formula $\mathrm{wt}(\Phi_V) = c_V(0)/2$ does not
 apply as such; for the Monster case one instead defines
 $\kappa_{\mathrm{BKM}}(J(p)-J(q)) := 0$ from the
 \emph{genus-$0$ property} of the Hauptmodul input $J$ (the unique
 normalised generator of the field of modular functions on
 $\mathrm{SL}_2(\mathbb{Z})$), consistent with the bi-weight
 $(0, 0)$ of the Koike--Norton--Zagier denominator $J(p) - J(q)$.
 \item The formal identity $c_V(0)/2 = 0/2 = 0$ gives the correct
 numerical answer \emph{because} the FLM-shifted character
 convention $\chi_{V^\natural}(\tau) = J(\tau)$ has $c_V(0) =
 [q^0] J(\tau) = 0$ by Atkin--Lehner normalisation. This
 coincidence does \emph{not} extend to any claim that the
 Borcherds-1992 two-variable construction is a degenerate limit
 of the Borcherds-1998 Grassmannian theta lift: the two
 constructions share the weight formula as a \emph{numerical
 pattern}, but are distinct theorems.
 \item In particular the no-automatic-functoriality clause
 Theorem~\ref{thm:borcherds-lift-universal}\textup{(iv)} excludes
 any formal passage from the Monster case to the Fake-Monster case;
 the category of Monster-Lie-algebra-type
 denominator identities (Borcherds 1992) and the category of
 Grassmannian singular theta lifts (Borcherds 1998) intersect
 only at a numerical pattern, not at the level of a universal
 functor.
\end{enumerate}
\end{lemma}

\begin{proof}
(i) Direct scope calculation: Theorem~\ref{thm:borcherds-lift-universal}
quantifies over type-IV lattices of signature \((2,n)\);
\(\mathrm{II}_{1,1}\) has signature \((1,1)\) and is excluded.
\par
(ii) The Hauptmodul $J$ generates the field of modular functions on
$\mathrm{SL}_2(\mathbb{Z})$ with constant term $0$; the Koike--Norton
presentation of $J(p) - J(q)$ \cite[\S8]{Borcherds1992} exhibits
$J(p) - J(q)$ as a bi-weight-$(0,0)$ object, so its ``BKM weight''
is $0$ as a direct reading of the two-variable modular-function
property, \emph{not} as an output of the theta integral.
\par
(iii) Numerical coincidence: under the FLM-shifted convention
$\chi_{V^\natural}(\tau) = J(\tau)$ so $c_V(0) = 0$; both
	conventions (Grassmannian weight formula at $b^+ \ge 2$, genus-$0$
	reading at $b^+ = 1$) return $\kappa_{\mathrm{BKM}}(J)=0$. Equality
of numerical outputs does not imply equality of constructions; the
Borcherds-1992 identity is proved by a Hecke-operator / replication-
formula argument \cite[Thm.~8.1, Thm.~10.1]{Borcherds1992},
independent of the later theta-lift machinery.
\par
(iv) Theorem~\ref{thm:borcherds-lift-universal}\textup{(iv)} states
that no functorial comparison is automatic. It therefore does not cover
lattice morphisms \(\mathrm{II}_{1,1}\hookrightarrow\mathrm{II}_{25,1}\)
or any passage from the Monster denominator to the Fake-Monster
denominator. Any cross-lattice comparison between Monster (case~1) and
Fake Monster (case~2) must go through the Conway--Norton
generalised-moonshine framework \cite{ConwayNorton1979} or the
Carnahan monstrous-moonshine generalisation, not through
Theorem~\ref{thm:borcherds-lift-universal}.
\end{proof}

\begin{remark}[BKM additivity fails]
\label{rem:monster-bkm-additivity-disclaimer}
The coincidence formula
\[
\kappa_{\mathrm{BKM}}(\Phi_N) \stackrel{?}{=} \kappa_{\mathrm{ch}}^{\mathrm{Hodge}} +
\chi(\mathcal{O}_{\mathrm{fiber}})
\]
is \emph{false} at every
$N \in \{1, 2, 3, 4, 6\}$ (cf.\ the $N=1$ datum $5 \ne 0$ and
$N = 2$ datum $4 \ne 1$). At the Monster case~(1) here, the
numerical value $\kappa_{\mathrm{BKM}}(J)=0$ does \emph{not} decompose
as $\kappa_{\mathrm{ch}}^{\mathrm{Hodge}}(V^\natural) + \chi(\mathcal{O}_{\bullet})$
for any CY geometry: $V^\natural$ is not the image of $\Phi_3$ on
any smooth compact CY$_3$ (it is a vertex algebra associated to the
Leech-orbifold, not to a sheaf category). The additive
decomposition $\kappa_{\mathrm{BKM}} \stackrel{?}{=} \kappa_{\mathrm{ch}}^{\mathrm{Hodge}} +
\kappa_{\mathrm{fiber}}$ is not a live K3-fibered formula either:
$5=3+2$ is only the $N=1$ Heisenberg/fibre convention-mixing
accident, while the compact total-space formula gives $0+0$.
Propagating it to the Monster case is a
Künneth-multiplicative-vs-additive $\kappa_\bullet$ conflation
(where $\bullet \in \{\mathrm{cat}, \mathrm{ch}^{\mathrm{Hodge}},
\mathrm{ch}^{\mathrm{Heis}}, \mathrm{BKM}, \mathrm{fiber}\}$ each
require separate subscript tracking, since the five canonical
invariants disagree on products and on fibrations).
\end{remark}

\begin{remark}[CY-A\texorpdfstring{$_3$}{-3} and the Borcherds lift]
\label{rem:cya3-borcherds-lift-compatibility}
The CY-to-chiral correspondence $\Phi_3^{(\Sigma_2, C)}$ of
Theorem~\ref{thm:cy-to-chiral-d3} and
Remark~\ref{rem:phi-d-parameter-tagging} produces, on a K3-fibered CY$_3$
$X$ with K3-fibre class $\Sigma_2 = \pi^{-1}(\mathrm{pt})$, reference curve
$C \subset X$ transverse to $\pi$, and Mukai lattice
$\Lambda_{\mathrm{Muk}}(X)$ even unimodular, an $E_1$-chiral algebra
$A_X^{(\Sigma_2, C)} = \Phi_3^{(\Sigma_2, C)}(D^b(\mathrm{Coh}(X))) =
\SpCh_{\Sigma_2, C}(\PhiFA_3(D^b(\mathrm{Coh}(X))))$ whose weight lattice is
$\Lambda_{\mathrm{Muk}}(X)$. Theorem~\ref{thm:borcherds-lift-universal} then
applies with $L = \Lambda_{\mathrm{Muk}}(X)$ and $V = A_X^{(\Sigma_2, C)}$,
and the composite
\[
 D^b(\mathrm{Coh}(X)) \xrightarrow{\Phi_3^{(\Sigma_2, C)}} E_1\text{-ChirAlg}(C) \xrightarrow{\Phi_{(-)}} \mathrm{AutForm}\bigl(\mathcal{G}(\Lambda_{\mathrm{Muk}}(X))\bigr)
\]
sends $X = K3 \times E$ with canonical $(\Sigma_2, C) = (K3, E)$ to $\Delta_5$
(Theorem~\ref{thm:k3e-borcherds-product}); on non-K3-fibered CY$_3$s the
second arrow is undefined
(Proposition~\ref{prop:bkm-weight-universal}~(iv)). Thus admission of a
Borcherds lift is a PROPERTY of the CY$_3$ together with its chosen
specialisation datum (existence of a K3 fibration giving even unimodular
weight lattice), not of the functor $\Phi_d$ itself. For $X$ without such
a fibration (quintic, local $\mathbb{P}^2$), the substitute is the BCOV
partition function, which is not a theta lift.
\end{remark}

\subsection{The bar Euler product chain}
\label{subsec:k3e-bar-euler-chain}

The full chain from the K3 elliptic genus to the bar complex
passes through five components.

\begin{enumerate}[label=\textup{Arrow \arabic*.}]
\item \textbf{Elliptic genus.}
The K3 elliptic genus $\phi_{0,1}(\tau, z)$,
the unique weak Jacobi form of weight~$0$ and index~$1$,
has Fourier coefficients $c(D) = c(4n - l^2)$
depending only on the discriminant $D$.
The first values (only $D \equiv 0$ or $3 \pmod{4}$ appear
for index~$1$):
\[
\begin{array}{c|cccccccccccc}
D & -1 & 0 & 3 & 4 & 7 & 8 & 11 & 12 & 15 & 16 & 19 & 20 \\
\hline
c(D) & 1 & 10 & -64 & 108 & -513 & 808 & -2752 & 4016 & -11775 & 16524 & -43200 & 58640
\end{array}
\]
(Convention: $c(D)$ here denotes the Eichler--Zagier per-$(n,l)$ coefficient $f(nm,l)$; thus $c(-1) = f(0,1) = 1$. The table in Proposition~\ref{prop:k3e-super-grading} uses the per-discriminant sum $c(-1) = \sum_{l:\,4 \cdot 0 - l^2 = -1} f(0,l) = f(0,1) + f(0,-1) = 2$; the factor of $2$ arises from summing over $l = \pm 1$.)

These values are computed from the theta-ratio formula
$\phi_{0,1} = 4[(\vartheta_2/\vartheta_{2,0})^2 + (\vartheta_3/\vartheta_{3,0})^2 + (\vartheta_4/\vartheta_{4,0})^2]$
in exact rational arithmetic.
The row-sum constraint
$\sum_l c(4n - l^2) = 0$ for $n \geq 1$
follows from $\phi_{0,1}(\tau, 0) = 12$.

\item \textbf{Borcherds lift.}
The Borcherds multiplicative lift sends $\phi_{0,1}$ to $\Delta_5$
of weight $c(0)/2 = 5 = \kappa_{\mathrm{BKM}}(\Delta_5)$
(Theorem~\ref{thm:k3e-product}).

\item \textbf{Signed root characters.}
The product formula gives
$\operatorname{sdim}\mathfrak g_{\Delta_5,(n,l,m)} = c(4nm - l^2)$ for the BKM superalgebra
$\fg_{\Delta_5}$, with the sign determining parity:
$c(D) > 0$ (bosonic) iff $D \equiv 0 \pmod{4}$,
$c(D) < 0$ (fermionic) iff $D \equiv 3 \pmod{4}$
(Proposition~\ref{prop:k3e-super-grading}).

\item \textbf{Bar Euler product identification.}
The chiral algebra
$A_{K3 \times E}^{(\Sigma_2,C)} = \Phi_3^{(\Sigma_2,C)}(D^b(\mathrm{Coh}(X)))$ is available on the fixed H1--H4 specialisation; its
bar Euler product is
\[
E_{B(A_{K3 \times E})}(q, y, p)
\;=\;
\prod_{(n,l,m) > 0}
(1 - q^n y^l p^m)^{c(4nm - l^2)}.
\]
This product equals $\Delta_5$ (up to normalization and the Weyl
vector exponential) by the denominator identity
(Theorem~\ref{thm:k3e-denominator}). The formal product agrees with the
Borcherds denominator. Its interpretation as the bar Euler product of
the \(d=3\) chiral output requires the CY-to-chiral comparison
\(\Phi_3\); without CY-A$_3$ the denominator remains a Borcherds
character formula, not a bar construction.

\item \textbf{Saito--Kurokawa decomposition.}
The additive decomposition
$\Delta_5 = \sum_m \phi_m(\tau, z) \, p^m$
provides the perturbative expansion: each Fourier--Jacobi
coefficient $\phi_m$ captures the $m$-instanton sector.
The first coefficient $\phi_1 = \phi_{10,1} = \eta^{18}\vartheta_1^2$
determines the rest by Hecke operators
(Theorem~\ref{thm:k3e-additive-lift}).
\end{enumerate}

The chain exhibits the $\kappa_\bullet$-spectrum phenomenon
of the chapter introduction:
the bar Euler product carries weight
$\kappa_{\mathrm{BKM}}(\Delta_5) = 5 = \mathrm{wt}(\Delta_5)$,
while the chiral algebra (if it exists) has
$\kappa_{\mathrm{ch}}^{\mathrm{Heis}} = 3 = \dim_\C(K3 \times E)$.
These are different invariants of different algebraizations, each carrying its own subscript from the closed $\kappa_\bullet$ menu.

\begin{remark}[Classical product sources]
\label{rem:k3e-bar-euler-bkm-chain-classical}
The Borcherds denominator formula for the generalised
Kac--Moody $\mathfrak{g}_{\Delta_5}$ is \cite{Borcherds1992,Borcherds1998};
the identification of \((\Phi_{10}^{\mathrm{un}})^{-1}=\Delta_5^{-2}\)
with the automorphic partition function of $1/4$-BPS states on
$K3 \times T^2$ is Dijkgraaf--Moore--Verlinde--Verlinde
\cite{DMVV}; Gritsenko--Nikulin
\cite{GritsenkoNikulin1995,GritsenkoNikulin1998} establish the paramodular
lift structure. The recasting, that this denominator is the
energy-graded bar Euler product of the framed chiral algebra
$A_{K3 \times E}^{(\Sigma_2,C)}$, uses the H1--H4 object-level
$d=3$ assignment plus the Vol~I Borcherds-lift bridge (conditional);
	it is not a construction of a global $\Phi_3$ functor. The numerical content
	-- the signed exponent
	\(\operatorname{sdim}\mathfrak g_{\Delta_5,(n,l,m)}=c(4nm-l^2)\),
	the Weyl vector $\rho = (1,1,1)$, the primitive weight
	$\kappa_{\mathrm{BKM}}(\Delta_5)=5$, the Igusa-square weight
	$10 = 2 \kappa_{\mathrm{BKM}}(\Delta_5)$, and the asymptotic absolute
	growth $|c(D)| \sim e^{2\pi\sqrt{D}}$ (Hardy--Ramanujan--Rademacher)
	-- is entirely classical. What the bar-cobar framework adds is the
identification of ``$\mathfrak{g}_{\Delta_5}$ positive half'' with
``$E_1$-bar differential on
$B^{\mathrm{ord}}(A_{K3 \times E}^{(\Sigma_2,C)})$'',
placing the denominator identity in a chiral-algebraic cohomological
hierarchy compatible with the shadow tower and the factorisation
structure of $K3 \times E$. No new automorphic-form identity is
established; the bar-cobar perspective supplies organisation, not
arithmetic.
\end{remark}


\section{Scattering diagrams and the shadow MC element}
\label{sec:k3e-scattering}

The Kontsevich--Soibelman wall-crossing formula and the shadow obstruction tower of the chiral algebra $A_{K3 \times E}$ are both governed by Maurer--Cartan elements in convolution algebras. The connection between them passes through the formalism of scattering diagrams.

\subsection{The scattering diagram of \texorpdfstring{$K3 \times E$}{K3 x E}}
\label{subsec:k3e-scattering-diagram}

\begin{definition}[Scattering diagram]
\label{def:k3e-scattering}
Let $\Gamma = \widetilde{H}(K3, \Z) \simeq \Lambda^{4,20}$ be the Mukai lattice. A \emph{scattering diagram} $\mathfrak{D}$ on the dual real torus $\Gamma^* \otimes \R$ consists of codimension-$1$ walls $(\mathfrak{d}, f_{\mathfrak{d}})$, where $\mathfrak{d} \subset \Gamma^* \otimes \R$ is a rational hyperplane and $f_{\mathfrak{d}} \in \hat{\C}[\Gamma]$ is a wall-crossing automorphism.

For $K3 \times E$, the scattering diagram $\mathfrak{D}_{K3 \times E}$ has:
\begin{itemize}
 \item An \emph{initial wall} $(\mathfrak{d}_\gamma, (1 - x_\gamma)^{\Omega(\gamma)})$ for each BPS charge $\gamma$ with $\Omega(\gamma) \neq 0$.
 \item \emph{Derived walls} created by the consistency condition: the composition of wall-crossing automorphisms around any closed loop must be trivial.
\end{itemize}
\end{definition}

\begin{conjecture}[KS scattering = BKM root system]
\label{conj:k3e-scattering-bkm}
The scattering diagram $\mathfrak{D}_{K3 \times E}$ is equivalent to the root system of $\mathfrak{g}_{\Delta_5}$:
\begin{enumerate}[label=(\roman*)]
 \item Each initial wall with protected index $\Omega(\gamma)$ corresponds to a positive root $\alpha_\gamma$ of $\mathfrak{g}_{\Delta_5}$ with signed root character $\operatorname{sdim}\mathfrak g_{\Delta_5,\alpha_\gamma}=\Omega(\gamma)$. The ordinary wall multiplicity is $|\Omega(\gamma)|$ only after the parity fixture or finite Hall--Borcherds recognition data.
 \item The derived walls correspond to the Weyl group orbits of imaginary roots.
 \item The consistency of the scattering diagram (path-independence of the wall-crossing product) is modeled by the denominator identity of $\mathfrak{g}_{\Delta_5}$ (Theorem~\ref{thm:k3e-denominator}).
\end{enumerate}
\end{conjecture}

\begin{remark}[Exact missing inputs for the scattering-to-root identification]
\label{rem:k3e-scattering-missing-inputs}
The conjecture above requires the following finite-height data:
\begin{enumerate}[label=\textup{(\roman*)}]
\item a motivic integration map from the wall-crossing Hall algebra to the completed quantum torus that preserves the chamber product and the relevant parity convention;
\item a proof that every initial wall of the $K3 \times E$ scattering diagram is indexed by a BPS charge whose discriminant matches a Borcherds root discriminant under the chosen specialisation;
\item a proof that the consistency completion adds precisely the Weyl orbits of imaginary roots and no extraneous walls;
\item a compatibility theorem identifying the scattering-diagram path-ordered product with the Borcherds denominator product on the same completed charge lattice.
\end{enumerate}
These are the exact missing ingredients separating the wall-crossing formalism from a root-system theorem.
\end{remark}

\begin{proposition}[Finite quantum-torus multiplication gate]
\label{prop:k3e-finite-scattering-quantum-torus-gate}
Let \(L_{\leq H}\) be a finite set of retained charge labels, equipped
with a positive height function \(\operatorname{ht}\), a partial sum
table \(s(\alpha,\beta)\in L_{\leq H}\sqcup\{0\}\), and an integral
skew pairing \(\langle-,-\rangle\) on \(L_{\leq H}\). Define the
truncated quantum-torus product on the basis symbols \(x_\alpha\) by
\[
 x_\alpha x_\beta
 =
 \begin{cases}
 q^{\langle\alpha,\beta\rangle}x_{s(\alpha,\beta)},
   &s(\alpha,\beta)\neq0,\\
 0,&s(\alpha,\beta)=0.
 \end{cases}
\]
This finite product is the height-\(H\) truncation of a quantum-torus
product of the form used in the Kontsevich--Soibelman wall factor
exactly when the following finite gates hold:
\begin{enumerate}[label=\textup{(\roman*)}]
\item \(\langle\alpha,\beta\rangle+
\langle\beta,\alpha\rangle=0\) for all retained labels;
\item \(s(\alpha,\beta)\neq0\) implies
\(\operatorname{ht}s(\alpha,\beta)
=\operatorname{ht}\alpha+\operatorname{ht}\beta\leq H\), and
\(s(\alpha,\beta)=0\) exactly when the sum lies outside the retained
height window;
\item the two partial sums \(s(s(\alpha,\beta),\gamma)\) and
\(s(\alpha,s(\beta,\gamma))\) agree, with \(0\) allowed on both sides;
\item whenever the common triple sum is retained,
\[
 \langle\alpha,\beta\rangle+\langle s(\alpha,\beta),\gamma\rangle
 =
 \langle\beta,\gamma\rangle+\langle\alpha,s(\beta,\gamma)\rangle .
\]
\end{enumerate}
Under these conditions the displayed multiplication is associative.
\end{proposition}

\begin{proof}
Condition~\textup{(i)} is the skewness of the quantum-torus exponent.
Condition~\textup{(ii)} says that the quotient by labels of height
\(>H\) is exactly the height truncation and not an additional quotient.
For three retained labels the two products are
\[
 (x_\alpha x_\beta)x_\gamma
 =
 q^{\langle\alpha,\beta\rangle+
 \langle s(\alpha,\beta),\gamma\rangle}
 x_{s(s(\alpha,\beta),\gamma)}
\]
when the left partial sums are retained, and zero otherwise; similarly
\[
 x_\alpha(x_\beta x_\gamma)
 =
 q^{\langle\beta,\gamma\rangle+
 \langle\alpha,s(\beta,\gamma)\rangle}
 x_{s(\alpha,s(\beta,\gamma))}.
\]
Thus the two parenthesisations agree exactly when the retained triple
label agrees and the exponent cocycle identity in
\textup{(iv)} holds. If both triple sums leave the height window, both
products are zero by \textup{(ii)}. This proves associativity of the
finite truncated product. The executable witness is
\texttt{finite\_scattering\_quantum\_torus\_gate}; it records the pair
products, skew defects, height-truncation defects, associativity
defects, and cocycle defects.
\end{proof}

\begin{remark}[Boundary of the finite quantum-torus gate]
\label{rem:k3e-finite-scattering-quantum-torus-boundary}
Proposition~\ref{prop:k3e-finite-scattering-quantum-torus-gate} is the
finite algebra test for the torus on which the scattering product is
written. It does not construct the motivic integration morphism
\(I_H\), the chamber-ordered Kontsevich--Soibelman product, or the
identification of wall exponents with Borcherds root multiplicities.
Those are the additional hypotheses of
Theorem~\ref{thm:k3e-finite-scattering-root-comparison}.
\end{remark}

\begin{theorem}[Finite scattering/root comparison after recognition]
\label{thm:k3e-finite-scattering-root-comparison}
Fix a finite Borcherds-height window \(W_H\), a chamber \(\sigma\),
and a strict sector \(S\). Assume the compact source packet
\(\mathsf M_H^{X,\mathrm{can}}\), the target packet
\(\mathsf M_H^{\Delta,\mathrm{can}}\), and the finite comparison packet
\(\mathsf A_H^{X\to\Delta}\) satisfy the finite recognition certificate
of Theorem~\ref{thm:k3e-finite-recognition-certificate}. Assume also a
motivic integration morphism
\[
 I_H\colon
 \mathcal H^{\mathrm{comp},\leq H}_{K3\times E,\sigma}
 \longrightarrow
 \widehat{\mathbb C_q[\Gamma_{\leq H}]}
\]
from the compact Hall source to the height-\(H\) quantum torus, with
\[
 I_H(p_{\gamma,i})=(-1)^{p(\gamma,i)}x_\gamma
\]
on parity-homogeneous primitive generators and with Hall multiplication
carried to the quantum-torus product. Let
\[
 \Omega_H(\gamma)
 =
 \operatorname{sdim}
 \operatorname{Prim}\mathcal H^{\mathrm{comp},\leq H}_{K3\times E,\sigma,\gamma}.
\]
Then the finite Kontsevich--Soibelman wall product
\[
 \prod_{\arg Z_\sigma(\gamma)}^{\curvearrowright}
 \prod_{\substack{k\geq1\\ \operatorname{ht}(k\gamma)\leq H}}
 \left(1-q^{k/2}x_{k\gamma}\right)^{k\Omega_H(k\gamma)}
\]
has the same exponent packet as the height-\(H\) truncation of the
Borcherds denominator of \(\mathfrak g_{\Delta_5}\):
\[
 \Omega_H(\gamma)=
 \operatorname{sdim}\mathfrak g_{\Delta_5,\alpha_\gamma}
 =
 c(D(\gamma))
\]
for every retained charge. Equivalently, after the same quantum-torus
ordering and cocycle normalisation, the finite wall product is the
height-\(H\) Borcherds denominator product. The transition
\(H+1\to H\) commutes with this equality whenever \(W_{H+1}\to W_H\)
is downward-saturated and the finite recognition certificates commute
with the height transitions.
\end{theorem}

\begin{proof}
The motivic integration morphism sends Hall convolution to the
twisted quantum-torus multiplication. On primitive generators it sends
the parity-homogeneous class \(p_{\gamma,i}\) to
\((-1)^{p(\gamma,i)}x_\gamma\). The exponent of the wall labelled by
\(\gamma\) is therefore the signed primitive index
\(\Omega_H(\gamma)\). The finite recognition certificate identifies the
primitive Hall quotient in charge \(\gamma\) with the
\(\alpha_\gamma\)-root space of the finite Serre--Borcherds quotient,
with the same parity fixture. Hence
\[
 \Omega_H(\gamma)
 =
 \operatorname{sdim}\mathfrak g_{\Delta_5,\alpha_\gamma}.
\]
The denominator theorem for \(\mathfrak g_{\Delta_5}\) gives
\(\operatorname{sdim}\mathfrak g_{\Delta_5,\alpha_\gamma}=c(D(\gamma))\)
with the same discriminant convention used in
Theorem~\ref{thm:k3e-denominator}. Thus each retained wall factor and
the corresponding Borcherds product factor have the same support and
the same signed exponent. Downward saturation removes exactly the
monomials of height \(>H\) on both sides. Projection from \(H+1\) to
\(H\) is deletion of those monomials, so it commutes with the displayed
factor-by-factor equality.
\end{proof}

\begin{remark}[Executable finite scattering/root gate]
The routine \texttt{finite\_scattering\_root\_report} in
\texttt{compute/lib/k3e\_finite\_bridge\_witness.py} records the
finite exponent comparison in
Theorem~\ref{thm:k3e-finite-scattering-root-comparison}. It checks,
charge by charge, that the protected Hall index equals the
Borcherds coefficient \(c(D)\), that the two finite supports agree, and
that deletion of charges above a lower height gives the same exponent
packet as recomputing at that lower height. The routine is a finite
gate for the theorem; it does not construct the compact Hall source or
the motivic integration morphism.
\end{remark}

\subsection{Shadow tower and scattering: two MC elements}
\label{subsec:k3e-shadow-scattering}

\begin{conjecture}[Two convolution algebras, one BPS spectrum]
\label{conj:k3e-two-mc}
The BPS spectrum of $K3 \times E$ is encoded by Maurer--Cartan elements in two distinct convolution algebras:
\begin{enumerate}[label=(\roman*)]
 \item The \emph{shadow MC element} $\Theta_A \in \mathrm{MC}(\mathrm{Def}_{\mathrm{cyc}}^{\mathrm{mod}}(A))$ of the framed chiral algebra $A = A_{K3 \times E}^{(\Sigma_2,C)}$, in the modular cyclic deformation complex of Volume~I.
 \item The \emph{KS MC element} $\cA_{\mathrm{KS}} = \exp\bigl(\sum_\gamma \Omega(\gamma) \, \mathrm{Li}_2(x_\gamma)\bigr)$ in the quantum torus Lie algebra, governing the wall-crossing automorphisms.
\end{enumerate}
These are \emph{not} the same MC element: $\Theta_A$ lives in the cyclic deformation complex (topological, $\cM_{g,n}$-valued), while $\cA_{\mathrm{KS}}$ lives in the quantum torus (algebraic, $K(\mathrm{Var})$-valued). The bridge between them factors through the motivic Hall algebra:
\[
 \mathrm{Def}_{\mathrm{cyc}}^{\mathrm{mod}}(A) \xleftarrow{\;\text{chiral}\;} \cH_{\mathrm{mot}}(D^b(\Coh(X))) \xrightarrow{\;\text{integration}\;} \hat{\C}[\Gamma].
\]
\end{conjecture}

\begin{remark}[Naive BCH is insufficient]
\label{rem:k3e-naive-bch}
The naive Baker--Campbell--Hausdorff pair-commutator of KS automorphisms does \emph{not} reproduce the $\phi_{0,1}$ signed exponents at low degrees. At degree~$4$, the naive BCH gives signed root-character exponents differing from $c(D)$ by non-uniform ratios that measure higher BPS bound-state contributions. The full identification requires the motivic DT theory of Kontsevich--Soibelman, in which the integration map $\int \colon \cH_{\mathrm{mot}} \to \hat{\C}[\Gamma]$ preserves the MC structure but introduces motivic measure corrections beyond the naive product formula.
\end{remark}

\begin{remark}[Degree filtration vs BPS charge filtration]
\label{rem:k3e-degree-bps}
The shadow obstruction tower has a degree filtration $\Theta_A^{\leq r}$ (finite-order projections) while the scattering diagram has a BPS charge filtration (walls of increasing $|D|$). The shadow-to-scattering bridge identifies:
\begin{center}
\begin{tabular}{ll}
\textbf{Shadow tower} & \textbf{Scattering diagram} \\ \hline
$\kappa_{\mathrm{ch}}^{\mathrm{Heis}}(K3 \times E) = 3$; $\kappa_{\mathrm{ch}}^{\mathrm{Hodge}}(K3) = 2$ (fiber) & weight $= 5 = c(0)/2$ (Borcherds lift) \\
Degree-$r$ shadow $\Theta_A^{\leq r}$ & Roots with $|D| \leq r$ \\
Shadow depth class M ($r_{\max} = \infty$) & Infinitely many BPS walls \\
Cubic shadow $C$ (degree $3$) & First fermionic wall ($D = 3$, $c(3) = -64$) \\
Quartic resonance $Q$ (degree $4$) & First derived wall from pentagon relation
\end{tabular}
\end{center}
The infinite shadow depth of the BKM algebra (class M) reflects the infinite product in the Borcherds formula: the K3 elliptic genus has $c(D) \neq 0$ for infinitely many discriminants $D$, with $|c(D)|$ growing as $e^{2\pi\sqrt{D}}$ by the Hardy--Ramanujan--Rademacher circle method.
\end{remark}


\subsection{Wall-crossing and the shadow tower}
\label{subsec:k3e-wall-crossing-shadow}

The KSWCF and the shadow obstruction tower are connected by the observation that the labeled-ordered product of KS automorphisms (ordered by decreasing phase of $Z(\gamma)$) has the algebraic structure of the $E_1$ bar differential. We make this precise, emphasizing what is a theorem and what is a conjecture.

\begin{conjecture}[KSWCF as \texorpdfstring{$E_1$}{E-1} bar differential]
\label{conj:k3e-kswcf-bar}
On the witnessed \(K3\times E\) Stage-$2$ locus of Theorem~CY-A$_3$ (Theorem~\ref{thm:cy-to-chiral-d3}), the labeled-ordered product
\[
 \mathcal{A}_{\mathcal{W}} = \prod_{\substack{\gamma : Z(\gamma) \in \R_{>0} \cdot e^{i\phi} \\ \text{decreasing phase}}}^{\mathrm{lab}} \mathcal{T}_\gamma^{\Omega(\gamma)}
\]
is identified with the $E_1$ bar differential $d_{\mathrm{bar}}$ on $B^{\mathrm{ord}}(A_{K3 \times E})$, where $B^{\mathrm{ord}}$ denotes the \emph{labeled-ordered} bar construction (not time-ordered; see Remark~\ref{rem:k3e-ordering-convention}). The bar Euler product
\[
 Z_{\mathrm{bar}}(q, y, p) = \exp\!\Bigl(\sum_\gamma \Omega(\gamma) \cdot \sum_{k \geq 1} \frac{x_{k\gamma}}{k}\Bigr)
\]
is the gauge-invariant of the $E_1$ MC element and is therefore wall-crossing invariant.
\end{conjecture}

\begin{remark}[Ordering convention]
\label{rem:k3e-ordering-convention}
The product in Conjecture~\ref{conj:k3e-kswcf-bar} is \emph{labeled-ordered}: each factor carries a label (the BPS charge~$\gamma$) and the product respects the total order on labels given by decreasing phase $\phi_\sigma(\gamma) = (1/\pi)\arg Z_\sigma(\gamma)$. This is the $E_1$ (associative, ordered) structure of the bar complex. It is \emph{not} time-ordered (radial ordering, OPE) or normally-ordered (Wick ordering). The three orderings produce distinct algebraic structures with distinct composition laws, and statements in this chapter name which ordering is meant at each use.
\end{remark}

\begin{conjecture}[Formal BCH model for shadow jumps at a wall]
\label{conj:k3e-shadow-jump}
Let $\mathcal{W}(v_1, v_2)$ be a wall at which a bound state $\gamma = v_1 + v_2$ with $\Omega(\gamma) \neq 0$ appears. The shadow tower coefficients $S_k$ jump by:
\begin{enumerate}[label=(\roman*)]
 \item At degree~$2$ (the $\kappa_{\mathrm{ch}}^{\mathrm{Heis}}$ level): $\Delta S_2 = \Omega(\gamma)$.
 \item At degree~$3$ (cubic shadow): $\Delta S_3$ has leading term $\Omega(\gamma) \cdot D(\gamma)/2$, where $D(\gamma) = 4nm - l^2$ is the discriminant of the BKM root.
 \item At degree~$k \geq 4$: the formal BCH expansion contributes a leading term of the shape $\Omega(\gamma) \cdot D(\gamma)^{k-2}/k!$, with additional nested-commutator corrections depending on the chosen normalisation.
\end{enumerate}
In particular:
\begin{itemize}
 \item Real root walls ($D = -1$): the shadow depth is unchanged (only $S_2$ receives a contribution).
 \item Lightlike walls ($D = 0$): the shadow depth is unchanged ($S_k = 0$ for $k \geq 3$).
 \item Imaginary root walls ($D > 0$): the shadow depth \emph{increases} (all $S_k$ with $k \geq 3$ receive nonzero contributions).
\end{itemize}
The bar complex interpretation uses the framed algebra $A_{K3 \times E}^{(\Sigma_2,C)}$ on the fixed H1--H4 locus; the formal BCH expansion of the KS product motivates this model.
\end{conjecture}

\begin{remark}[Discriminant stratification of walls]
\label{rem:k3e-disc-strat}
The BKM root system of $\mathfrak{g}_{\Delta_5}$ stratifies the walls by discriminant~$D$:
\begin{center}
\renewcommand{\arraystretch}{1.2}
\begin{tabular}{@{}lccc@{}}
\textbf{Discriminant} & \textbf{Wall type} & \textbf{Signed index $c(D)$} & \textbf{Shadow degree} \\ \hline
$D = -1$ & Real root (flop) & $1$ & $2$ only \\
$D = 0$ & Lightlike (large volume) & $10$ & $2$ only \\
$D = 3$ & First imaginary & $-64$ & $\geq 3$ \\
$D = 4$ & Second imaginary & $108$ & $\geq 3$ \\
$D \to \infty$ & Deep imaginary & $|c(D)| \sim e^{2\pi\sqrt{D}}$ & $\geq 3$
\end{tabular}
\end{center}
(Multiplicity values in EZ per-discriminant convention: $c(-1) = f(0,1) = 1$; the sum over $l = \pm 1$ gives~$2$.)
The real and lightlike walls contribute only at the $\kappa_{\mathrm{ch}}^{\mathrm{Heis}}$ level (degree~$2$), while the imaginary roots generate the cubic and higher shadows. The infinite shadow depth (class~M) of $K3 \times E$ is a direct consequence of the infinitely many imaginary roots.
\end{remark}

\begin{conjecture}[Bar Euler product invariance]
\label{conj:k3e-bar-euler-invariance}
Let $\sigma_+, \sigma_-$ be stability conditions in adjacent chambers separated by a wall~$\mathcal{W}$. Let $\Theta^{E_1}(\sigma_\pm)$ denote the $E_1$ MC elements encoding the BPS spectra in each chamber for the framed algebra $A_{K3 \times E}^{(\Sigma_2,C)}$. Then:
\begin{enumerate}[label=(\roman*)]
 \item The MC elements are gauge-equivalent: $\Theta^{E_1}(\sigma_+) = e^{\mathrm{ad}_\alpha} \cdot \Theta^{E_1}(\sigma_-)$ for some gauge element $\alpha$.
 \item The bar Euler product $Z_{\mathrm{bar}}$ is chamber-independent: it depends only on the gauge-equivalence class.
 \item Different chambers give different \emph{factorizations} of the same bar Euler product (different labeled-ordered decompositions of the same element in the bar complex).
 \item The shadow coefficients~$S_k$ \emph{do} change at the wall (Conjecture~\ref{conj:k3e-shadow-jump}), but their generating function $\exp(\sum S_k / k!)$ is invariant.
\end{enumerate}
The KSWCF (Theorem~\ref{thm:k3e-ks-wc}) proves that the KS product $\prod \mathcal{T}_\gamma^{\Omega(\gamma)}$ is chamber-independent. The conjecture identifies this product with the bar Euler product of the framed chiral algebra $A_{K3 \times E}^{(\Sigma_2,C)}$.
\end{conjecture}

\begin{remark}[Primitive wall-crossing and Yau--Zaslow counts]
\label{rem:k3e-primitive-wc}
The simplest walls involve rank-$1$ DT invariants: sheaves supported on curves~$C \subset K3$. The BPS counts are the Yau--Zaslow numbers $n_h = \chi(\mathrm{Hilb}^h(K3))$:
\[
 \sum_{h \geq 0} n_h \, q^{h-1} = \frac{1}{\eta(q)^{24}}, \qquad n_0 = 1, \; n_1 = 24, \; n_2 = 324, \; n_3 = 3200, \; n_4 = 25650,
\]
proved by Beauville and Bryan--Leung. The primitive wall-crossing involves a rank-$1$ sheaf~$\mathcal{F}$ (supported on~$C$) and a skyscraper sheaf~$\mathcal{O}_x$ with Euler pairing $\langle [\mathcal{F}], [\mathcal{O}_x] \rangle = 1$. The bound state $\mathrm{Cone}(\mathcal{O}_x \to \mathcal{F})$ has $\Omega = n_h$ before the wall and $\Omega = 0$ after the wall. This is the simplest instance of the pentagon identity, governing the single-particle/two-particle transition at each curve class.
\end{remark}

\begin{remark}[Split attractor flow and the shadow tree]
\label{rem:k3e-attractor-shadow}
The Denef split attractor flow conjecture asserts that every BPS state admits a tree decomposition into primitive constituents along walls of marginal stability. For $K3 \times E$, each node of the attractor tree corresponds to a wall-crossing event that modifies the shadow tower:
\begin{itemize}
 \item At each binary split $\gamma \to \gamma_1 + \gamma_2$, the shadow jump (Conjecture~\ref{conj:k3e-shadow-jump}) transfers shadow data between the parent and children.
 \item The leaves of the tree are primitive states (real roots, $D = -1$) with shadow depth~$2$ (class~G).
 \item The tree structure itself encodes the ``assembly'' of class~M from class~G components: the total shadow depth is infinite because the tree has unboundedly many nodes.
\end{itemize}
This provides a tree-level picture of the fiber-to-global shadow depth jump (G $\to$ M): each additional level of the attractor tree adds shadow complexity from the bound-state channels. The full non-perturbative answer (the Borcherds product for $\Delta_5$) resums the infinite attractor tree.
\end{remark}

\begin{remark}[Stokes automorphism = KS wall-crossing for \texorpdfstring{$K3 \times E$}{K3 x E}]
\label{rem:k3e-stokes-ks}
The wall-crossing and shadow tower structures of this section are unified by the identification of the Stokes automorphism of the shadow Borel transform with the KS wall-crossing automorphism (Theorem~\ref{thm:stokes-ks-conifold} for the conifold; Conjecture~\ref{conj:stokes-ks-k3e} for $K3 \times E$). The BKM root system stratifies the identification:
\begin{itemize}
\item At each real root wall ($D = -1$, $c(-1) = 1$): the local model is the resolved conifold, the pentagon identity holds, and the identification is a theorem. Each real root wall has one Stokes line and one wall.
\item At each imaginary root level ($D > 0$), the Stokes constants are the signed coefficients $\cS_\omega = c(D)$ and match the protected BPS indices $\Omega(\gamma)$. The ordinary number of conjugate Stokes-line pairs is $|c(D)|$ only after the target parity fixture.
\item The infinite shadow depth (class~M) is the Stokes-side manifestation of the infinite BPS spectrum: infinitely many imaginary roots $\to$ infinitely many Stokes lines $\to$ infinite shadow tower.
\end{itemize}
\end{remark}


\section{The complete enumerative--algebraic dictionary}
\label{sec:k3e-dictionary}

We collect the correspondences between the four perspectives on $K3 \times E$ into a single dictionary.

\begin{center}
\renewcommand{\arraystretch}{1.3}
\begin{tabular}{@{}llll@{}}
\textbf{Enumerative} & \textbf{Algebraic (BKM)} & \textbf{Shadow tower} & \textbf{Automorphic} \\ \hline
$\chi(K3 \times E) = 0$ & Trivial degree-$0$ DT & $\chi/24 = 0$ (virtual) & $M(-q)^0 = 1$ \\
$n_h^0 = \chi(\mathrm{Hilb}^h)$ & Real root ($\alpha^2 = 2$) & $F_1 = 1/12$ & $\phi_1 = \eta^{18}\vartheta_1^2$ \\
$c(D) = f(D)$ & $\operatorname{sdim}\mathfrak g_\alpha = f(D(\alpha))$ & Degree-$r$ shadow & $\prod (1 - q^ny^lp^m)^{f(D)}$ \\
$\Omega(\gamma)$ (KS protected index) & Signed root character & Obstruction class $o_{r+1}$ & Borcherds exponent \\
DT $=$ PT ($\chi = 0$) & No MacMahon correction & No degree-$0$ curvature & $Z_{\DT,0} = 1$ \\
Pentagon identity & BKM Serre relation & Cubic gauge triviality & Jacobi identity for $\fn_+$ \\
$|c(D)| \sim e^{2\pi\sqrt{D}}$ & Infinite root system & Class M ($r_{\max} = \infty$) & $\Delta_5$ is a cusp form \\
$\kappa_{\mathrm{ch}}^{\mathrm{Heis}} = 3$ & $24$ singular fibers & Class G $\to$ class M & $\mathrm{wt} = f(0)/2 = 5$
\end{tabular}
\end{center}

\begin{remark}[The K3 \texorpdfstring{$\times$}{x} E paradigm]
\label{rem:k3e-paradigm}
The $K3 \times E$ example exhibits every structural feature of the correspondence in concrete form: the fiber-to-global shadow depth jump (G $\to$ M), the genus-$1$ to genus-$2$ escalation via Borcherds lift, the infinite BKM root system from finite elliptic genus data, the DT/PT identity from $\chi = 0$, and the wall-crossing / scattering diagram connection to the MC formalism. It is the unique example where all six independent constructions (Section~\ref{sec:k3e-six-routes}) are simultaneously visible: Route~R\textsubscript{1} is the direct framed $\Phi_3$ specialisation, while the remaining routes are independent Hall, Borcherds, lattice, sigma-model, and Schur constructions. This makes it the Rosetta stone for the CY quantum group theory. Their convergence is the conditional CY-C colimit problem.
\end{remark}

\section{Cross-volume structural results}
\label{sec:k3e-cross-volume}

The following results are proved in Volume~I
(\S\ref{sec:k3-sigma-model}--\S\ref{sec:grand-synthesis-k3xe}
of the K3 modular-tower chapter there) and apply to the $K3 \times E$ tower.
We record them here for cross-reference; conventions
follow Volume~I throughout.

\begin{enumerate}[label=\textup{(K3-\arabic*)}]
\item $\kappa_{\mathrm{ch}}^{\mathrm{Heis}}(K3 \times E) = 3 = \dim_\mathbb{C}$ by
 six independent constructions.
 $\kappa_{\mathrm{ch}}^{\mathrm{Heis}}(K3) = 2$, $\kappa_{\mathrm{ch}}^{\mathrm{Heis}}(E) = 1$, additivity gives
 $\kappa_{\mathrm{ch}}^{\mathrm{Heis}}(K3 \times E) = 3$.

\item The factor $2$ in $Z_{K3} = 2\,\phi_{0,1}$ is the
 modular characteristic
 $\kappa_{\mathrm{ch}}^{\mathrm{Heis}}(\mathcal{A}_{K3}) = 2$.
 The massive Mathieu multiplicities satisfy
 $A_n = \kappa_{\mathrm{ch}}^{\mathrm{Heis}} \cdot \dim(\rho_n)$,
 where $\rho_n$ is the
 $M_{24}$-representation at level $n$. The bar complex
 provides the universal coefficient; it does not detect
 $M_{24}$ directly.

\item $\kappa_{\mathrm{BKM}}(\Delta_5) =
\operatorname{wt}(\Delta_5)=c_{\Delta_5}(0)/2=10/2=5$ is the
canonical derivation from the Borcherds-lift weight (paramodular
convention, Remark~\ref{rem:k3e-convention-delta5-phi10}); the
coincidence
\(\kappa_{\mathrm{ch}}^{\mathrm{Heis}}(K3 \times E)+\chi(\cO_{K3})
=3+2=5\) is not a derivation but a convention-mixing numerical match at
\(N=1\); the second summand is \(\kappa_{\mathrm{cat}}\) of the \(K3\)
fibre, not \(\kappa_{\mathrm{ch}}^{\mathrm{Heis}}(K3)\), so no
inner-product identity is in play.
 The ratio $\kappa_{\mathrm{BKM}}/\kappa_{\mathrm{ch}}^{\mathrm{Heis}}
 = 5/3$ compares independent invariants; it is not a
 second-quantization correction formula. The Hilbert--Chow
 exceptional divisor accounts for
 $\chi(\operatorname{Hilb}^2) - \chi(\operatorname{Sym}^2)
 = 24 = \chi(K3)$.

\item The shadow obstruction tower does not produce
 \(\Phi_{10}^{\mathrm{un}}\). Four obstructions: categorical (O1),
 $\kappa_{\mathrm{ch}}^{\mathrm{Heis}}$-vs-$\kappa_{\mathrm{BKM}}$ mismatch (O2: $3 \neq 5$), second quantization
 (O3), Schottky at $g \geq 4$ (O4:
 $\operatorname{codim} = (g{-}2)(g{-}3)/2$).

\item The mock modular decomposition
 $Z_{K3} = (h - 24\mu)\,\vartheta_1^2/\eta^3$
 follows from the Zwegers Appell--Lerch decomposition of the
 \(K3\) elliptic genus, where \(\mu\) is the Appell--Lerch sum.

\item The shadow generating function
 $Z^{\mathrm{sh}}(t) = \kappa_{\mathrm{ch}}^{\mathrm{Heis}}\bigl((t/2)/\sin(t/2) - 1\bigr)$
 converges with radius $R = 2\pi$, has entire Borel
 transform, and exhibits no resurgence. This contrasts
 with the Gevrey-$1$ divergent topological string.

\item Boundary $A_E$ ($\kappa_{\mathrm{fiber}} = 24$) vs sigma $V_{K3}$
 ($\kappa_{\mathrm{ch}}^{\mathrm{Hodge}}(K3) = 2$): ratio $12$ at genus~$1$; conditional
 at $g \geq 2$ ($V_{K3}$ is multi-weight).

\item $\lambda_5^{\mathrm{FP}} = 73/3503554560$. The
 unreduced numerator $2^{2g-1} - 1 = 511$ at $g = 5$
 requires GCD reduction with $(2g)!$.
\end{enumerate}

\section{The \texorpdfstring{$K3 \times E$}{K3 x E} holographic datum \texorpdfstring{$H(K3 \times E)$}{H(K3 x E)}}
\label{sec:k3e-holographic-datum}

The seven-face construction of Chapter~\ref{ch:cy-holographic-datum-master}
assigns to each CY3 geometry a holographic modular Koszul datum
$H(X) = (A_X, A_X^!, r_X(z), \Theta_X, \kappa_{\mathrm{ch},X}, \nabla^{\mathrm{hol}}_X, \mathcal{G}_X)$
packaging boundary algebra, Koszul dual, $R$-matrix, MC element,
modular characteristic, holographic connection, and quantum group
into a single coherent object. The $K3 \times E$ tower specializes
this datum completely.

\begin{construction}[The \texorpdfstring{$K3 \times E$}{K3 x E} holographic datum]
\label{constr:k3e-holographic-datum}
The seven faces of $H(K3 \times E)$ are:
\begin{enumerate}[label=\textup{(\roman*)}]
 \item \emph{Boundary algebra} $A_{K3 \times E}$. The chiral algebra
 of the K3 sigma model tensored with the free boson on $E$,
 carrying central charge $c = 6$ and conformal weights
 $(h_1, h_2, h_3) = (1, 1, 1)$ from the three complex directions.
 The sigma model factor $V_{K3}$ is multi-weight;
 the elliptic factor contributes a single Heisenberg generator.

 \item \emph{Koszul dual} $A^!_{K3 \times E}$. Obtained by
 Verdier duality on Ran($X$): $D_{\mathrm{Ran}}(B(A)) \simeq B(A^!)$
 (cf.\ Vol~I, Convention~\ref{conv:bar-coalgebra-identity}).
 The Koszul dual carries the conjectural complementary Heisenberg-level
 modular characteristic
 $\kappa_{\mathrm{ch}}^{\mathrm{Heis}}(A^!) = 2$,
 contingent on the conjectural identification of the Borcherds weight
 $\kappa_{\mathrm{BKM}}(\Delta_5)=5$ as the Heisenberg Koszul conductor pairing
 the values $3$ and $2$. If this holds, the pairing is nonzero;
 the Virasoro comparison has conductor $13$, so complementarity is
 family-specific, not universal zero. The chiral Hodge-supertrace reading
 $\kappa_{\mathrm{ch}}^{\mathrm{Hodge}}(K3 \times E) = \sum_q (-1)^q h^{0, q}(K3 \times E) = 0$ is a separate invariant and is
 not to be substituted into the Heisenberg-level equation above.

 \item \emph{$R$-matrix} $r_{K3 \times E}(z)$. The collision
 residue $r(z) = \operatorname{Res}^{\mathrm{coll}}_{0,2}(\Theta_A)$
 is the binary genus-$0$ shadow of the full MC element. For
 $K3 \times E$, this is the Maulik--Okounkov $R$-matrix
 $R^{\mathrm{MO}}(z)$ of Section~\ref{sec:k3e-mo-rmatrix},
 acting on the BPS Fock space. The $R$-matrix controls the proposed
 Hall--Drinfeld/Super-Yangian deformation of the Borcherds branch:
 $R^{\mathrm{MO}}(z)$ satisfies the Yang--Baxter equation with spectral parameter
 valued in $K_0(\mathrm{Coh}(K3 \times E))$, and the commuting
 Hamiltonians it generates are the Gaudin operators on
 moduli of sheaves.

 \item \emph{MC element} $\Theta_{K3 \times E}$. The bar-intrinsic
 construction (cf.\ Chapter~\ref{ch:cy-holographic-datum-master})
 gives $\Theta_A := D_A - d_0 \in \mathrm{MC}(\mathfrak{g}^{\mathrm{mod}}_{A})$.
 The projections yield: $\kappa_{\mathrm{ch}}^{\mathrm{Heis}} = 3$ at degree~$2$
 (from $\kappa_{\mathrm{ch}}^{\mathrm{Heis}}(K3) + \kappa_{\mathrm{ch}}^{\mathrm{Heis}}(E) = 2 + 1$), the cubic
 shadow $C$ at degree~$3$ (nonvanishing, as the BKM root
 system has imaginary simple roots of multiplicity $> 1$),
 and an infinite tower at all higher degrees. The shadow
 obstruction tower does not terminate: the BKM denominator
 identity forces nonzero contact terms $S_r \neq 0$ for
 all $r \geq 2$.

 \item \emph{Modular characteristic} $\kappa_{\mathrm{ch}}^{\mathrm{Heis}}(K3 \times E) = 3$.
 This is the chiral-algebra modular characteristic, obtained
 from additivity (K3-1 above). The global BKM characteristic
 $\kappa_{\mathrm{BKM}}(\Delta_5) = 5 = \operatorname{wt}(\Delta_5)$
 governs the Borcherds lift (K3-3 above). The ratio
 $\kappa_{\mathrm{BKM}}/\kappa_{\mathrm{ch}}^{\mathrm{Heis}} = 5/3$ is the
 second-quantization multiplier.

 \item \emph{Holographic connection} $\nabla^{\mathrm{hol}}_{g,n}$.
 The modular shadow connection
 $\nabla^{\mathrm{hol}}_{g,n} = d - \operatorname{Sh}_{g,n}(\Theta_A)$
 on the moduli of stability conditions. The singular locus
 is the wall-crossing divisor of Section~\ref{sec:k3e-wall-crossing};
 the monodromy around each wall is the Kontsevich--Soibelman
 gauge transformation of
 Theorem~\ref{thm:ks-mc-gauge}.

 \item \emph{Quantum group} $\mathcal{G}(K3 \times E)$.
 The quantum vertex chiral group of
 Section~\ref{sec:k3e-qvcg}, whose root system is the
 BKM root system $\mathfrak{g}_{\Delta_5}$ and whose
 representation category carries the Maulik--Okounkov
 $R$-matrix as braiding.
\end{enumerate}
\end{construction}

\begin{proposition}[Root datum of the holographic datum]
\label{prop:k3e-holographic-root-datum}
The lattice underlying $H(K3 \times E)$ is $\Lambda^{3,2}$ of
Section~\ref{sec:k3e-lattice}, with the $\mathrm{Sp}_4 \simeq
\mathrm{SO}^+(\Lambda^{3,2})$ isomorphism of
Lemma~\ref{lem:sp4-so-iso} encoding the root data. The Igusa
cusp form $\Delta_5$ has the infinite product expansion
\[
 \Delta_5(\Omega) = p \prod_{(n,l,m) > 0} \bigl(1 - p^n q^l r^m\bigr)^{f(4nm - l^2)}
\]
(with $f(D)$ the Fourier coefficients of $\phi_{0,1}$, and
$p = e^{2\pi i \tau}$, $q = e^{2\pi i z}$, $r = e^{2\pi i \sigma}$
the coordinates on the Siegel upper half-space $\mathbb{H}_2$),
which expresses the denominator identity of $\mathfrak{g}_{\Delta_5}$.
This product has the \emph{form} of a bar-complex Euler characteristic
(each factor $(1 - p^n q^l r^m)^{f(D)}$ corresponds to a putative
bar generator of dimension vector $(n, l, m)$ with signed Euler exponent
$f(D)$; an ordinary generator count would require the target parity
fixture), but the identification of $\Delta_5$ with the bar Euler
product of a chiral algebra requires both the framed $d = 3$ CY-to-chiral
assignment of Theorem~\ref{thm:cy-to-chiral-d3} on the fixed H1--H4 locus and the Vol~I bar Euler product
machinery applied to $A_{K3 \times E}^{(\Sigma_2,C)}$. Conjecturally,
the modular-operadic supertrace of
$B^{E_3}(\PhiFA_3(D^b\Coh(K3\times E)))$ gives this bar-Euler product
after $E_3$-Koszul duality: its root-indexed Euler factors should be the
Borcherds factors above, with exponent the superdimension of the
corresponding root space. Thus the Gritsenko--Nikulin product and the
bar-Euler product have matching exponents on the fixed framed
$K3\times E$ locus, but their equality is conditional on the Vol.~I
Borcherds-lift bridge and CY-D$_3$.
\end{proposition}

\begin{remark}[Shadow depth: class M]
\label{rem:k3e-holographic-shadow-depth}
The $K3 \times E$ holographic datum has shadow depth $r_{\max} = \infty$
(class M in the shadow depth classification of
Definition~\ref{def:shadow-depth-classification}).
The infinite tower is forced by the BKM superalgebra:
the signed imaginary-root exponents $f(D)$ have absolute growth
\(\lvert f(D)\rvert \sim e^{2\pi\sqrt{D}}\) (Rademacher asymptotics),
producing nonvanishing contact terms at every degree.
The shadow obstruction tower of the framed chiral algebra $A_{K3 \times E}^{(\Sigma_2,C)}$
is thus the chiral-algebraic shadow of the BKM denominator identity,
and its non-termination reflects the infinite-dimensional root system
of $\mathfrak{g}_{\Delta_5}$.
\end{remark}


%% Enhanced symmetry and the Niemeier lattice landscape


\section{Topological string partition function vs shadow partition function}
\label{sec:k3e-topological-string-shadow}

The topological string on $K3 \times E$ provides a concrete realisation of the shadow partition function of Volume~I. On the enumerative side, the DT partition function is completely determined by the DMVV formula and the Igusa cusp form. On the algebraic side, the bar Euler product of the framed chiral algebra $A_{K3 \times E}^{(\Sigma_2,C)}$ should encode the same BPS data. This section makes the comparison explicit, isolating what is proved from what remains conditional on the Vol~I Borcherds-lift bridge and CY-D$_3$.

\subsection{The G\"ottsche formula and \texorpdfstring{$1/\eta^{24}$}{1/eta-24}}
\label{subsec:k3e-gottsche-eta}

The simplest incarnation of the topological string on $K3 \times E$ is the rank-$1$ DT partition function, which counts ideal sheaves on K3 fibers.

\begin{theorem}[G\"ottsche, 1990]
\label{thm:k3e-gottsche-product}
The generating function of Hilbert scheme Euler characteristics is
\[
 \sum_{n \geq 0} \chi(\mathrm{Hilb}^n(K3)) \, q^n
 = \prod_{k \geq 1} \frac{1}{(1 - q^k)^{24}}
 = \frac{q}{\eta(q)^{24}}
 = \frac{1}{\Delta(q)},
\]
where $\Delta(q) = q \prod_{k \geq 1}(1 - q^k)^{24}$ is the Ramanujan cusp form and $24 = \chi_{\mathrm{top}}(K3) = \kappa_{\mathrm{fiber}}$. The first coefficients are:
\begin{center}
\begin{tabular}{c|ccccccc}
$n$ & $0$ & $1$ & $2$ & $3$ & $4$ & $5$ & $6$ \\ \hline
$\chi(\mathrm{Hilb}^n)$ & $1$ & $24$ & $324$ & $3200$ & $25650$ & $176256$ & $1073720$
\end{tabular}
\end{center}
These are the $24$-coloured partition numbers $p_{24}(n)$.
\end{theorem}

The exponent $24$ in the product is $\kappa_{\mathrm{fiber}}$, the fiber modular characteristic. It is distinct from $\kappa_{\mathrm{ch}}^{\mathrm{Heis}}(K3 \times E) = 3$ (the Heisenberg-level rank-additive characteristic) and from $\kappa_{\mathrm{BKM}}(\Delta_5) = 5$ (the Borcherds lift weight). The G\"ottsche formula is the rank-$1$ sector of the DT theory; the full partition function requires all ranks (the DMVV formula below).

\subsection{The DT partition function of \texorpdfstring{$K3 \times E$}{K3 x E}: three products}
\label{subsec:k3e-dt-three-products}

Three infinite products are in play and are distinct:

\begin{enumerate}[label=(\roman*)]
\item The MacMahon function $M(q) = \prod_{n \geq 1} 1/(1-q^n)^n$ has exponents growing linearly in~$n$. For a general CY3 $X$ with $\chi(X) \neq 0$, the degree-$0$ DT invariant is $Z_{\DT,0}(X; q) = M(-q)^{\chi(X)}$ (Maulik--Nekrasov--Okounkov--Pandharipande). But $\chi(K3 \times E) = 24 \cdot 0 = 0$, so $Z_{\DT,0} = 1$ (Theorem~\ref{thm:k3e-dt-degree0}).

\item The rank-$1$ DT partition function is $\prod_{k \geq 1} 1/(1-q^k)^{24}$ (G\"ottsche), with \emph{constant} exponent $24$. This is $1/\Delta(q)$, the reciprocal of the Ramanujan cusp form.

\item The \emph{full} DT partition function of $K3 \times E$ is the second-quantized DMVV formula, which is a product over \emph{three} variables $(p, y, q)$ with exponents $c(4nm - l^2)$ from $\phi_{0,1}$, and is related to \(1/\Phi_{10}^{\mathrm{un}}=\Delta_5^{-2}\) before the OP scalar normalisation \(D_5=64^{-1}\Delta_5\) is imposed (Theorem~\ref{thm:dt-igusa}).
\end{enumerate}

The formula $\prod 1/(1-q^n)^{24n}$ does not appear in the DT theory of $K3 \times E$.

\subsection{The DMVV formula and \texorpdfstring{$1/\Phi_{10}^{\mathrm{un}}=\Delta_5^{-2}$}{1/Phi-10-un = Delta-5-2}}
\label{subsec:k3e-dmvv-vs-delta5}

\begin{theorem}[Dijkgraaf--Moore--Verlinde--Verlinde, 1997]
\label{thm:k3e-dmvv-product}
The second-quantized partition function of $K3 \times E$ is
\begin{equation}\label{eq:k3e-dmvv}
 Z_{\mathrm{DMVV}}(p, y, q)
 = \prod_{\substack{(n, l, m) > 0}} \frac{1}{(1 - p^n y^l q^m)^{c(4nm - l^2)}},
\end{equation}
where $c(D)$ are the discriminant-indexed Fourier coefficients of the weak Jacobi form $\phi_{0,1}$ (Eichler--Zagier normalization: $c(-1) = 1$, $c(0) = 10$, $c(3) = -64$, $c(4) = 108$).
\end{theorem}

\begin{remark}[DMVV vs Borcherds product]
\label{rem:k3e-dmvv-borcherds}
The DMVV product~\eqref{eq:k3e-dmvv} and the Borcherds product for $\Delta_5$ (Theorem~\ref{thm:k3e-borcherds-product}) involve the \emph{same} product over triples $(n,l,m)$ with the \emph{same} exponents $c(D)$. The difference:
\begin{itemize}
\item The DMVV product has exponents $-c(D)$ (it is $1/\prod(1-\cdots)^{c(D)}$), producing the \emph{reciprocal} of the Borcherds product.
\item The Borcherds product includes the Weyl vector prefactor $q \cdot y \cdot p$.
\end{itemize}
The relation is
\[
 Z_{\mathrm{DMVV}} = C'/\Phi_{10}^{\mathrm{un}} = C'\Delta_5^{-2},
\]
where the scalar \(C'\) absorbs the Weyl-vector contribution and
\(\Phi_{10}^{\mathrm{un}}=\Delta_5^2\)
(Remark~\ref{rem:k3e-convention-delta5-phi10}).
\end{remark}

\subsection{Diagonal restriction of \texorpdfstring{$\Delta_5$}{Delta-5}}
\label{subsec:k3e-diagonal-delta5}

To make the comparison with the shadow partition function concrete, we expand $\Delta_5$ on the diagonal $\sigma = \tau$, $z = 0$ (the one-variable restriction).

\begin{proposition}[Diagonal restriction of \texorpdfstring{$\Delta_5$}{Delta-5}]
\label{prop:k3e-diagonal-delta5}
On the diagonal $p = q$, $y = 1$ (i.e.\ $\sigma = \tau$, $z = 0$), the Borcherds product becomes
\[
 \Delta_5\big|_{\mathrm{diag}}(q) = q^2 \prod_{s \geq 1} (1 - q^s)^{e_{\mathrm{diag}}(s)},
\]
where the effective diagonal exponent is
\[
 e_{\mathrm{diag}}(s)
 = \sum_{\substack{n + m = s \\ (n,l,m) > 0}} c(4nm - l^2).
\]
The exponents $e_{\mathrm{diag}}(s)$ are \emph{not} constant: they grow with $s$, reflecting the class~$M$ (infinite shadow depth) nature of the BKM algebra. In particular, $e_{\mathrm{diag}}(1)$ receives contributions only from root triples with $n + m = 1$ (either $n=1, m=0$ or $n=0, m=1$), while $e_{\mathrm{diag}}(s)$ for large $s$ receives contributions from all decompositions $n + m = s$ and all discriminants $D = 4nm - l^2$.
\end{proposition}

\begin{proof}
Setting $p = q$ and $y = 1$ in the Borcherds product (Theorem~\ref{thm:k3e-borcherds-product}), each factor $(1 - q^n y^l p^m)$ becomes $(1 - q^{n+m})$. Grouping by $s = n + m$ gives the stated formula. The Weyl vector $q \cdot y \cdot p$ becomes $q \cdot 1 \cdot q = q^2$.
\end{proof}

\subsection{Bar Euler product comparison}
\label{subsec:k3e-bar-euler-comparison}

\begin{conjecture}[Shadow partition function identification]
\label{conj:k3e-shadow-pf}
The fixed-specialisation chiral algebra
$A^{(K3,E)}_{K3 \times E} = \Phi_3^{(K3,E)}(D^b(\mathrm{Coh}(K3 \times E)))$,
defined under H1--H4 together with the compact CY$_3$ closure input
of either $B_{\mathrm{TCFT}}^{(2)}$ or the filtered $HH^{-2}_{E_1}$
theorem, has a bar Euler product
\[
 \Theta_{A^{(K3,E)}_{K3 \times E}}(q) = \prod_{n \geq 1} (1 - q^n)^{a_n},
\]
whose exponents $\{a_n\}$ are determined by the shadow obstruction tower of $A^{(K3,E)}_{K3 \times E}$. The identification with the Borcherds product is:
\begin{enumerate}[label=(\roman*)]
\item At rank~$1$ \textup{(proved, Theorem~\ref{thm:k3e-gottsche-product}):} $a_n = -24 = -\kappa_{\mathrm{fiber}}$ for all $n$, giving $\Theta_{A,\mathrm{rank\text{-}1}} = \prod(1-q^n)^{-24}$, which is $q/\eta(q)^{24}$ (the G\"ottsche formula via the bar complex of the K3 sigma model).
\item At full rank \textup{(using fixed-specialisation CY-A$_3$ plus compact closure):} the exponents $a_n$ coincide with the effective diagonal exponents $e_{\mathrm{diag}}(n)$ of $\Delta_5$ (Proposition~\ref{prop:k3e-diagonal-delta5}), so that $\Theta_A$ is the diagonal restriction of $1/\Delta_5$ up to the Weyl vector prefactor. The identification with the Borcherds product requires the Vol~I bridge (Conjecture~\ref{conj:k3e-shadow-pf}).
\end{enumerate}
\end{conjecture}

\begin{remark}[Exact missing inputs for the shadow partition function]
\label{rem:k3e-shadow-pf-missing-inputs}
The conjecture above still requires the following data before the full-rank
comparison can be a theorem:
\begin{enumerate}[label=\textup{(\roman*)}]
\item a construction of the fixed-specialisation chiral algebra
\(A^{(K3,E)}_{K3 \times E}\) on the H1--H4 locus that is compatible
with the bar complex and its Euler product;
\item a proof that the bar Euler exponents of the full specialisation
match the diagonal Borcherds exponents \(e_{\mathrm{diag}}(n)\) under
the same normalisation of the Weyl prefactor and the same sign
convention for the BPS indices;
\item a comparison theorem identifying the fixed-specialisation bar
Euler product with the diagonal restriction of \(1/\Delta_5\), rather
than merely with its rank-\(1\) fibre truncation;
\item a compatibility theorem showing that the comparison survives the
passage from the H1--H4 locus to the inverse-limit shadow tower and to
the finite Hall--Borcherds recognition data.
\end{enumerate}
These are the exact steps missing from the structural observation to a
proof of the shadow partition function identification.
\end{remark}

\begin{remark}[Borcherds denominator \texorpdfstring{$\neq$}{!=} bar Euler product]
\label{rem:k3e-borcherds-bar}
The identification of Conjecture~\ref{conj:k3e-shadow-pf}(ii) requires \emph{two} independent inputs:
\begin{enumerate}[label=(\alph*)]
\item The framed CY-to-chiral assignment $\Phi_3^{(\Sigma_2,C)}$ at $d = 3$ (Theorem~\ref{thm:cy-to-chiral-d3}), producing $A_{K3 \times E}^{(\Sigma_2,C)}$ under H1--H4 and fixed specialisation.
\item The Volume~I Borcherds-lift bridge (Construction~\ref{constr:k3e-shadow-bkm-bridge}), identifying the shadow obstruction tower with the Borcherds product.
\end{enumerate}
Input~(a) is supplied only on the fixed framed H1--H4 locus. Input~(b)
is the morphism that turns the common BPS-charge product and common
exponents into an equality between the shadow obstruction tower and the
Borcherds product.
\end{remark}

\subsection{Protected BPS indices and the bar Euler exponents}
\label{subsec:k3e-bps-exponents}

The Donaldson--Thomas invariants $\Omega(\gamma)$ of $K3 \times E$ are protected signed BPS indices of charge $\gamma \in \widetilde{H}(K3 \times E, \Z)$. The shadow partition function uses these as bar Euler exponents through the following dictionary:

\begin{center}
\renewcommand{\arraystretch}{1.3}
\begin{tabular}{@{}lll@{}}
\textbf{BPS side} & \textbf{Shadow tower side} & \textbf{Proof input} \\ \hline
$\Omega(0,0,n) = p_{24}(n)$ & Bar Euler coeff.\ of $q^n$ at rank $0$ & G\"ottsche theorem \\
$\Omega(1,0,h) = n_h^0 = \chi(\mathrm{Hilb}^h)$ & $a_n = -24$ (class G, Gaussian) & Yau--Zaslow formula \\
$\Omega^{\mathrm{ind}}(\gamma) = c(D(\gamma))$ & Shadow depth class $M$, $r_{\max} = \infty$ & Gritsenko--Nikulin product \\
Sign$(\Omega) \leadsto$ parity & Bosonic ($c > 0$) / fermionic ($c < 0$) & Fourier-coefficient sign \\
Pentagon identity & Cubic gauge triviality (degree $3$) & BCH pentagon identity \\
Full protected BPS index $\leadsto$ bar Euler & $\{a_n\} = \{e_{\mathrm{diag}}(n)\}$ & Vol.~I bridge and CY-D$_3$
\end{tabular}
\end{center}

\noindent The key structural point: the bar Euler exponents at rank~$1$ are \emph{constant} ($a_n = -24$), reflecting the class~$G$ (Gaussian, finite shadow depth) nature of the single-fiber theory. The jump to non-constant exponents (class~$M$, infinite shadow depth) occurs when all BPS charges are included, i.e.\ when the full DMVV product is engaged. This is the shadow tower manifestation of the fiber-to-global depth jump (Proposition~\ref{prop:k3e-fiber-global}(iii)).

\subsection{Wall-crossing and the shadow tower}
\label{subsec:k3e-wc-shadow}

The Kontsevich--Soibelman wall-crossing formula (Theorem~\ref{thm:k3e-ks-wc}) and the shadow obstruction tower encode the same BPS spectrum through different algebraic structures. Their correspondence maps the shadow tower data to the KS data as follows:

\begin{enumerate}[label=(\roman*)]
\item \emph{Degree filtration $\leftrightarrow$ discriminant stratification.} The degree-$r$ shadow projection $\Theta_A^{\leq r}$ captures BPS states at discriminants $|D| \leq r$ (Definition~\ref{def:k3e-d-stratification}). The first nontrivial wall-crossing (beyond the free-field sector) occurs at $D = 3$, where $c(3) = -64$ (fermionic); this corresponds to the degree-$3$ cubic shadow obstruction.

\item \emph{MC elements in different algebras.} The shadow MC element $\Theta_A \in \mathrm{MC}(\mathrm{Def}_{\mathrm{cyc}}^{\mathrm{mod}}(A))$ lives in the cyclic deformation complex; the KS automorphism $\cA_{\mathrm{KS}}$ lives in the quantum torus algebra $\hat{\C}[\Gamma]$ (Conjecture~\ref{conj:k3e-two-mc}). The bridge passes through the motivic Hall algebra.

\item \emph{Infinite depth $\leftrightarrow$ infinite product.} The shadow class~$M$ (infinite shadow depth) corresponds to the infinite Borcherds product: the signed BPS indices have absolute growth $|c(D)| \sim e^{2\pi\sqrt{D}}$ (Hardy--Ramanujan--Rademacher), producing infinitely many walls of marginal stability and infinitely many shadow tower levels.
\end{enumerate}

\subsection{Automorphic sources of the computed invariants}
\label{subsec:k3e-novelty}

The invariants in this subsection have fixed classical sources:
\(\chi(\mathrm{Hilb}^n(K3))\) is G\"ottsche's product, \(\tau(n)\) is the
Ramanujan cusp-form coefficient, the signed BPS indices
\(\Omega^{\mathrm{ind}}(\gamma)=c(D)\) are Fourier coefficients of the K3 Jacobi
form, \(\Delta_5\) is the primitive Gritsenko--Nikulin product, and the
DVV/DMVV square is \(\Phi_{10}^{\mathrm{un}}=\Delta_5^2\). The bar Euler
product $\Theta_A$ at rank~$1$ recovers the G\"ottsche formula by
construction. At full rank, the conjectural identification of the
chiral-half bar Euler product with \(\Delta_5^{-1}\) would identify its
square with the unnormalised Igusa inverse
\((\Phi_{10}^{\mathrm{un}})^{-1}=\Delta_5^{-2}\); the framed assignment
\(\Phi_3^{(\Sigma_2,C)}\) supplies the object-level chiral algebra only
under H1--H4, and the Vol~I Borcherds-lift bridge is the remaining
comparison datum.

The G\"ottsche, KKV, and DMVV products give the same coefficient
comparison after the primitive/free-particle logarithm is taken.

\section{BPS black hole entropy and the shadow partition function}
\label{sec:k3e-bps-entropy}

The full BPS partition function is \((\Phi_{10}^{\mathrm{un}})^{-1}=\Delta_5^{-2}\); the primitive chiral-half denominator is \(\Delta_5^{-1}\). The shadow tower controls the full asymptotic expansion of \(\log|\Omega(D)|\), not merely its leading term, only after the protected physical comparison identifies the chiral-half square with the BPS counting function. Whether the class-M shadow depth reproduces the full Rademacher series of \(1/\Phi_{10}^{\mathrm{un}}\) is therefore a precise constraint on the chiral construction. The Strominger--Vafa entropy furnishes that constraint: the \(D \to \infty\) regime is where the shadow invariants \(S_r\) and the Kloosterman sums \(\mathrm{Kl}(D, -1; c)\) must match term by term.

The Strominger--Vafa black hole (1996) is the prototypical example of microscopic entropy matching.
For Type~IIB string theory compactified on $K3 \times S^1$ with charges $(n_1, n_5, n_p)$ (D1-branes, D5-branes, momentum), the Bekenstein--Hawking entropy is
\begin{equation}
\label{eq:k3e-strominger-vafa}
 S_{\mathrm{BH}} = \frac{A}{4G} = 2\pi\sqrt{n_1 n_5 n_p}.
\end{equation}
The microscopic state count comes from the D1-D5 CFT on $\mathrm{Sym}^{n_1 n_5}(K3)$. The CFT has central charge $c_L = 6 n_1 n_5$, and the momentum level is $N = n_p$. The Cardy formula gives
\begin{equation}
\label{eq:k3e-cardy}
 S_{\mathrm{micro}} = 2\pi\sqrt{\frac{c_L \cdot N}{6}} = 2\pi\sqrt{n_1 n_5 n_p} = S_{\mathrm{BH}},
\end{equation}
establishing the exact leading-order agreement.

\subsection{The Fourier coefficients of \texorpdfstring{$1/\Phi_{10}^{\mathrm{un}}$}{1/Phi-10-un} and the Rademacher expansion}
\label{subsec:k3e-rademacher}

The DT partition function \(Z^{\mathrm{DT}}(K3 \times E) = C\,\Delta_5^{-2}\) in the unnormalised BPS square convention (Theorem~\ref{thm:dt-igusa}) means the protected signed BPS indices \(\Omega(D)\) at discriminant \(D = 4nm - l^2\) are Fourier coefficients of \(1/\Phi_{10}^{\mathrm{un}}\), where \(\Phi_{10}^{\mathrm{un}}=\Delta_5^2\) (Remark~\ref{rem:k3e-convention-delta5-phi10}). The Rademacher expansion gives exact absolute asymptotics:
\begin{equation}
\label{eq:k3e-rademacher}
 \log|\Omega(D)|
 =
 \pi\sqrt D
 -
 \alpha_{\log}^{\Phi_{10}}
 \bigl(\mathcal N_{\mathrm{cont}},\mathcal P_{\mathrm{pol}},\mu,
 \epsilon_{\Delta/\Phi}\bigr)\log D
 +O(1),
\end{equation}
where \(\mathcal N_{\mathrm{cont}}\) is the contour normalisation,
\(\mathcal P_{\mathrm{pol}}\) is the polar-orbit datum, \(\mu\) is the
charge-lattice measure, and \(\epsilon_{\Delta/\Phi}\) records the
primitive \(\Delta_5\) versus square \(\Phi_{10}^{\mathrm{un}}\)
convention.  The leading term \(\pi\sqrt D\) is the
Bekenstein--Hawking entropy \(S_{\mathrm{BH}}=\pi\sqrt D\), with
\(D=4n_1n_5n_p\) in the \(l=0\) Strominger--Vafa chamber.  A numerical
value of \(\alpha_{\log}^{\Phi_{10}}\) belongs to the normalized
compact Siegel kernel, not to the Borcherds weight
\(\kappa_{\mathrm{BKM}}(\Delta_5)\) alone.  The rank-one packet in
Proposition~\ref{prop:k3e-rank-one-rademacher-packet} supplies the
\(\phi_{0,1}\) arithmetic input with its \(I_{3/2}\)-kernel; the
compact coefficient in~\eqref{eq:k3e-rademacher} is extracted only
after the protected trace and contour normalisations have been matched.

For a fixed normalized compact kernel, the higher conductor terms are
exponentially separated from the principal saddle: the \(c=2\) term is
of relative size \(e^{-\pi\sqrt D/2}\) against the \(c=1\) saddle in
the five-dimensional normalization, up to the Kloosterman phase and
the algebraic Bessel prefactor.

\subsection{The shadow tower and the genus expansion}
\label{subsec:k3e-shadow-entropy}

The shadow partition function $\Theta_{A_{K3 \times E}}$ (Conjecture~\ref{conj:k3e-shadow-pf}) encodes the protected BPS indices through the genus expansion
\begin{equation}
\label{eq:k3e-genus-expansion}
 \log Z^{\mathrm{sh}} = \sum_{g \geq 1} F_g \, g_s^{2g},
 \qquad
 F_g = F_g^{\mathrm{sc}} + [\text{higher-degree shadow corrections}],
 \qquad
 F_g^{\mathrm{sc}} = \kappa_{\mathrm{ch}}^{\mathrm{Heis}} \cdot \hat{A}_g,
\end{equation}
where $\hat{A}_g$ is the $g$-th $\hat{A}$-genus coefficient ($\hat{A}_1 = 1/24$, $\hat{A}_2 = 7/5760$, $\hat{A}_3 = 31/967680$). The shadow tower corrections at genus $g$ scale as $\delta S^{(g)} \sim \kappa_{\mathrm{ch}}^{\mathrm{Heis}} \cdot \hat{A}_g \cdot (4\pi^2/S_{\mathrm{BH}})^{2g-1}$, decreasing rapidly for large $D$. These are the Wald entropy corrections from higher-derivative $R^{2g}$ terms in the effective action.

\begin{remark}[Which \texorpdfstring{$\kappa_\bullet$}{kappa-.} subscript controls the entropy]
\label{rem:k3e-kappa-entropy}
The shadow tower uses the Heisenberg-level modular characteristic $\kappa_{\mathrm{ch}}^{\mathrm{Heis}}(A_{K3 \times E}) = 3 = \dim_{\mathbb{C}}(K3 \times E)$ as input. The entropy is controlled by $\kappa_{\mathrm{BKM}}(\Delta_5) = 5 = \mathrm{wt}(\Delta_5)$. These are \emph{different} invariants: the chiral algebra is the \emph{input} (shadow tower parameter); the BKM weight is the \emph{output} (automorphic form weight, emerging from the bar Euler product). The chain is:
\[
 \kappa_{\mathrm{ch}}^{\mathrm{Heis}} \;\xrightarrow{\text{shadow tower}}\; \Theta_A \;\xrightarrow{\text{bar Euler product}}\; \Delta_5 \;(\mathrm{wt} = \kappa_{\mathrm{BKM}}(\Delta_5) = 5) \;\xrightarrow{\text{Fourier}}\; \Omega(D) \;\xrightarrow{\log}\; S_{\mathrm{BH}}.
\]
The correct universal formula is $\kappa_{\mathrm{BKM}}(\Phi_N) = c_N(0)/2$ (Borcherds weight theorem, 1998): for $K3 \times E$, $c_1(0) = 10$ gives $\kappa_{\mathrm{BKM}}(\Delta_5) = 5$. This is \emph{unconditional}: it uses the normalized K3 Jacobi input $\phi_{0,1}$, not the chiral algebra $A_{K3 \times E}$. The apparent arithmetic identity $5 = 3 + 2$ is a convention-mixing accident under the Heisenberg-level reading; the naked Hodge-supertrace additivity rule fails at every $N \in \{1, 2, 3, 4, 6\}$ (at $N = 1$, the Borcherds side is $5$, while the compact chiral Hodge supertrace and elliptic-fibre Euler characteristic are both zero). The additive rule with the Heisenberg reading also fails for all $(\bZ/N\bZ)$-orbifolds with $N \ge 2$ (Remark~\ref{rem:bkm-decomposition-adversarial}).
\end{remark}

\subsection{The Rademacher expansion as shadow tower resummation}
\label{subsec:k3e-shadow-rademacher}

The Rademacher expansion has two independent filtrations.  The
Kloosterman-conductor filtration is indexed by \(c\geq1\).  Inside each
fixed conductor term, the modified Bessel function has its own
large-\(D\) asymptotic expansion, producing the logarithmic and
power-suppressed corrections to that saddle.  The shadow arity is
matched here to the conductor filtration, not to the internal Bessel
asymptotic expansion.

\begin{center}
\renewcommand{\arraystretch}{1.3}
\begin{tabular}{@{}lll@{}}
\textbf{Shadow arity} & \textbf{Conductor term} & \textbf{Asymptotic scale} \\ \hline
\(k=2\) & \(c=1\) & principal saddle \(e^{\pi\sqrt D}\), with its Bessel logarithmic expansion \\
\(k=3\) & \(c=2\) & first non-perturbative saddle \(e^{\pi\sqrt D/2}\) \\
\(k=4\) & \(c=3\) & second non-perturbative saddle \(e^{\pi\sqrt D/3}\) \\
\(k=\infty\) & all \(c\geq1\) & full Rademacher conductor series \\
\end{tabular}
\end{center}

\noindent The shadow class~$M$ nature of $K3 \times E$ (infinite shadow depth) corresponds to the full infinite Rademacher series: every Kloosterman sum $c \geq 1$ contributes, and every shadow degree $r \geq 2$ is needed. For a class~$G$ algebra (finite shadow depth), only finitely many Rademacher terms would suffice; $K3 \times E$ is genuinely class~$M$.

\begin{conjecture}[Shadow--Rademacher identification]
\label{conj:k3e-shadow-rademacher}
For the framed algebra $A_{K3 \times E}^{(\Sigma_2,C)}$ available after the H1--H4 specialisation of Theorem~\ref{thm:cy-to-chiral-d3}, the shadow obstruction tower resummation at degree $r$ is expected to reproduce the Rademacher expansion of \(1/\Phi_{10}^{\mathrm{un}}\) truncated at Kloosterman conductor $c = r - 1$:
\[
 \Theta^{\leq r}_{A_{K3 \times E}}(D)
 =
 \mathcal K^{\leq r-1}_{\mathrm{rad}}
 \bigl(D;\Phi_{10}^{\mathrm{un}},\mathrm{contour}\bigr)
 + \mathcal E_r(D).
\]
Here \(\mathcal K_{\mathrm{rad}}\) denotes the protected
Rademacher kernel after the contour normalisation, polar-data
extraction, and compact trace comparison have been fixed.  The
conjectural error term \(\mathcal E_r(D)\) is required to be
exponentially smaller than the last retained conductor term, uniformly
in the height filtration.
\end{conjecture}

\begin{proposition}[Rank-one finite Rademacher packet]
\label{prop:k3e-rank-one-rademacher-packet}
On the Jacobi coefficient lane of \(\phi_{0,1}\), set
\[
 R_c^{(1)}(D)
 =
 \frac{\sqrt2\,\pi}{c}\,
 K(-1,D;c,\chi_{\eta^3})\,D^{-3/4}\,
 I_{3/2}\!\left(\frac{\pi\sqrt D}{c}\right),
 \qquad D>0,
\]
where \(K(-1,D;c,\chi_{\eta^3})\) is the generalized Kloosterman sum
with the \(\eta^3\)-multiplier.  Let
\[
 \Theta^{(1),\leq C}_{\mathrm{rad}}(D)
 :=
 \sum_{c=1}^{C} R_c^{(1)}(D).
\]
For
\[
 \mathcal D_{\leq20}^{(1)}
 =
 \{3,4,7,8,11,12,15,16,19,20\}
\]
and \(C=5\), the finite packet
\(\Theta^{(1),\leq5}_{\mathrm{rad}}\) has the following two properties.
\begin{enumerate}[label=\textup{(\roman*)}]
\item The conductor projection satisfies
\[
 \Theta^{(1),\leq C}_{\mathrm{rad}}(D)-R_C^{(1)}(D)
 =
 \Theta^{(1),\leq C-1}_{\mathrm{rad}}(D)
 \qquad(2\leq C\leq5,\;D\in\mathcal D_{\leq20}^{(1)}).
\]
\item Against the exact \(\phi_{0,1}\)-coefficient \(c(D)\), the
relative residual
\[
 \varepsilon_5(D)
 :=
 \frac{\left|\Theta^{(1),\leq5}_{\mathrm{rad}}(D)-|c(D)|\right|}
 {|c(D)|}
\]
satisfies
\[
 \max_{D\in\mathcal D_{\leq20}^{(1)}}\varepsilon_5(D)<0.03.
\]
\end{enumerate}
Thus the rank-one arithmetic packet supplies a genuine finite-height
Rademacher witness for the BKM root-multiplicity lane.  It does not
construct the compact \(CY_3\) protected trace map, and it does not
identify the shadow tower of
\(A_{K3\times E}^{(\Sigma_2,C)}\) with the Siegel expansion of
\((\Phi_{10}^{\mathrm{un}})^{-1}\).
\end{proposition}

\begin{proof}
The formula for \(R_c^{(1)}(D)\) is the Rademacher kernel for the
weight-\(-\tfrac12\) theta components of the weak Jacobi form
\(\phi_{0,1}\), with the polar coefficient at \(D=-1\).  The conductor
projection identity is the defining identity for partial sums: deleting
the last summand of the \(C\)-term partial sum gives the \((C-1)\)-term
partial sum.

For the residual estimate, the exact coefficients \(c(D)\) are the
Fourier coefficients of \(\phi_{0,1}\) in the Eichler--Zagier
normalisation used throughout the chapter.  The finite computation
evaluates the generalized Kloosterman sums, the \(I_{3/2}\)-Bessel
factor, the five conductor terms, and the residuals for the displayed
ten discriminants.  The executable certificate
\(\texttt{rademacher\_finite\_height\_certificate}\) in
\(\texttt{compute/lib/k3e\_finite\_bridge\_witness.py}\) records the
maximum residual and the transition check; the tests in
\(\texttt{compute/tests/test\_k3e\_finite\_bridge\_witness.py}\) and
\(\texttt{compute/tests/test\_shadow\_rademacher\_identification.py}\)
verify the two displayed assertions.
\end{proof}

\begin{corollary}[Finite conductor-to-arity certificate]
\label{cor:k3e-rank-one-rademacher-arity-certificate}
In the rank-one packet of
Proposition~\ref{prop:k3e-rank-one-rademacher-packet}, the conductor
term \(c\) is recorded at shadow arity \(k=c+1\):
\[
 (1,2),(2,3),(3,4),(4,5),(5,6).
\]
For every \(D\in\mathcal D_{\leq20}^{(1)}\), the projection from
arity \(k=C+1\) to arity \(k=C\) deletes exactly the conductor-\(C\)
summand. The five-conductor packet improves the leading-saddle
approximation uniformly:
\[
 \max_{D\in\mathcal D_{\leq20}^{(1)}}\varepsilon_5(D)
 <
 \max_{D\in\mathcal D_{\leq20}^{(1)}}\varepsilon_1(D).
\]
The executable certificate records no conductor-projection defect on
this finite lane.
\end{corollary}

\begin{proof}
The arity assignment is the convention fixed in
Conjecture~\ref{conj:k3e-shadow-rademacher}: shadow arity \(k\)
retains the Rademacher conductor \(c=k-1\). The projection statement is
the finite partial-sum identity proved in
Proposition~\ref{prop:k3e-rank-one-rademacher-packet}. The residual
comparison is computed on the same ten discriminants and the same
\(\eta^3\)-multiplier normalisation: the certificate
\(\texttt{rademacher\_finite\_height\_certificate}\) records
\(\max\varepsilon_5<\max\varepsilon_1\) and an empty list of conductor
projection defects. This is a finite arithmetic statement on
\(\phi_{0,1}\); it is not the compact \(CY_3\) shadow theorem.
\end{proof}

\begin{proposition}[Finite Rademacher polar--Bessel gate]
\label{prop:k3e-rademacher-polar-bessel-gate}
Fix a finite discriminant set
\(\mathcal D\subset\{D>0:c(D)\neq0\}\) and a conductor bound \(C\).
Suppose that a finite compact-shadow comparison supplies, for every
\((D,c)\in\mathcal D\times\{1,\ldots,C\}\), a row
\[
(p_{D,c},a_{D,c},m_{D,c},\nu_{D,c},k_{D,c},s_D)
\]
consisting of a polar discriminant, polar coefficient, multiplier,
Bessel order, shadow arity, and signed target coefficient.  On the
rank-one \(\phi_{0,1}\) lane this supplied row matches the Rademacher
kernel of Proposition~\ref{prop:k3e-rank-one-rademacher-packet} if and
only if
\[
p_{D,c}=-1,\qquad
a_{D,c}=1,\qquad
m_{D,c}=v_{\eta^3},\qquad
\nu_{D,c}=\frac32,\qquad
k_{D,c}=c+1,\qquad
s_D=c(D).
\]
For \(C'\leq C\), the finite transition from conductor bound \(C\) to
\(C'\) is compatible with this gate exactly when the rows with
\(c>C'\) are deleted and all retained rows are unchanged.
\end{proposition}

\begin{proof}
The Rademacher term displayed in
Proposition~\ref{prop:k3e-rank-one-rademacher-packet} is determined by
six pieces of data: the unique polar term \(D=-1\) of the theta
component, its coefficient \(1\), the \(\eta^3\)-multiplier
\(v_{\eta^3}\), the Bessel order \(1-(-1/2)=3/2\), the conductor \(c\),
and the signed \(\phi_{0,1}\)-coefficient \(c(D)\) whose absolute value
is approximated by the finite Rademacher sum.  Hence a supplied finite
row is the same rank-one arithmetic row precisely under the displayed
equalities.  The arity equality \(k=c+1\) is the convention of
Corollary~\ref{cor:k3e-rank-one-rademacher-arity-certificate}.  Changing
the conductor bound from \(C\) to \(C'\) changes the finite packet only
by removing the rows with conductor \(>C'\); therefore transition
compatibility is exactly row deletion.  The executable witness is
\texttt{rademacher\_polar\_bessel\_gate}; it records missing rows and
separate polar, multiplier, Bessel-order, arity, coefficient, and
transition defects.  It verifies supplied finite polar--Bessel data; it
does not construct the protected compact trace map or the uniform
all-height Siegel Rademacher theorem.
\end{proof}

\begin{proposition}[Finite Rademacher truncation-error gate]
\label{prop:k3e-rademacher-truncation-error-gate}
Let \(\mathcal D\subset\{D>0:c(D)\neq0\}\) be finite and let \(C\geq1\).
For \(1\leq c\leq C\), set
\[
\Theta_{\mathrm{rad}}^{(1),\leq c}(D)
=
\sum_{q=1}^{c} R_q^{(1)}(D),
\qquad
E^{\mathrm{abs}}_{D,c}
=
\left|\Theta_{\mathrm{rad}}^{(1),\leq c}(D)-|c(D)|\right|,
\qquad
E^{\mathrm{rel}}_{D,c}
=
\frac{E^{\mathrm{abs}}_{D,c}}{|c(D)|}.
\]
Suppose that a finite shadow--Rademacher comparison supplies
non-negative error majorants
\[
B^{\mathrm{abs}}_{D,c},\qquad B^{\mathrm{rel}}_{D,c}
\qquad(D\in\mathcal D,\;1\leq c\leq C).
\]
The supplied finite truncation-error datum is valid on this window if
and only if
\[
E^{\mathrm{abs}}_{D,c}\leq B^{\mathrm{abs}}_{D,c},
\qquad
E^{\mathrm{rel}}_{D,c}\leq B^{\mathrm{rel}}_{D,c}
\]
for every displayed row.  A terminal tolerance \(\tau\) at conductor
bound \(C\) holds on \(\mathcal D\) if and only if
\[
\max_{D\in\mathcal D} E^{\mathrm{rel}}_{D,C}\leq\tau.
\]
For \(C'\leq C\), transition from \(C\) to \(C'\) is compatible exactly
when the rows with conductor \(>C'\) are deleted and the retained
majorants are unchanged.
\end{proposition}

\begin{proof}
The finite Rademacher sum
\(\Theta_{\mathrm{rad}}^{(1),\leq c}(D)\) is a concrete real number
once \(D\), \(c\), the polar term, and the multiplier have been fixed by
Proposition~\ref{prop:k3e-rademacher-polar-bessel-gate}.  The absolute
and relative residuals are therefore the two displayed numbers
\(E^{\mathrm{abs}}_{D,c}\) and \(E^{\mathrm{rel}}_{D,c}\).  A supplied
majorant is valid precisely when it is not smaller than the residual it
claims to bound, row by row.  The terminal tolerance statement is the
same inequality after taking the maximum over the terminal conductor
row \(c=C\).  Lowering the conductor bound removes the terms and
majorants with \(c>C'\); no further algebra is involved in the finite
transition.  The executable witness is
\texttt{rademacher\_truncation\_error\_gate}; it records missing
absolute and relative bounds, failed bound inequalities, the terminal
relative residual, and the conductor-projection gate.  It is the finite
error-bound test, not the uniform all-height Siegel theorem.
\end{proof}

\begin{remark}[Exact missing inputs for the shadow--Rademacher comparison]
\label{rem:k3e-shadow-rademacher-missing}
The preceding finite statements construct the rank-one arithmetic packet
and its polar--Bessel and truncation-error row criteria.
The conjecture for \(A_{K3\times E}^{(\Sigma_2,C)}\) becomes a theorem
only after the following additional data are constructed and matched:
\begin{enumerate}[label=\textup{(\roman*)}]
\item a protected comparison map identifying the shadow-tower degree
\(r\) truncation with the first \(r-1\) Kloosterman-conductor terms of
the Rademacher series, at the same normalisation of \(\Phi_{10}^{\mathrm{un}}\);
\item a proof that the compact shadow obstruction coefficients satisfy
the same Bessel recursion and polar-data input as the arithmetic
Rademacher series; Proposition~\ref{prop:k3e-rank-one-rademacher-packet}
supplies this only on the rank-one \(\phi_{0,1}\) coefficient lane;
\item a compatibility theorem showing that the BPS/DT sign convention and
the chiral trace normalisation agree with the contour integral
normalisation used in the Rademacher expansion;
\item a control theorem for the error term under height truncation,
uniform in the compact shadow packet, so that the inverse-limit shadow
tower reproduces the full asymptotic series rather than just the
rank-one finite certificate.
\end{enumerate}
These are the exact steps still missing from the entropy chain.
\end{remark}

The asymptotic input is the Dabholkar--Murthy--Zagier Rademacher
expansion for \((\Phi_{10}^{\mathrm{un}})^{-1}\). The Borcherds weight
input is \(\kappa_{\mathrm{BKM}}(\Delta_5)=c(0)/2\); the additive formula
\(\kappa_{\mathrm{BKM}}(\Delta_5)=\kappa_{\mathrm{ch}}^{\mathrm{Hodge}}+\chi(\cO_S)\) fails because
the compact Hodge supertrace and the Borcherds denominator are
different constructions.


\section{The chiral algebra obstruction for \texorpdfstring{$\mathfrak{g}_{\Delta_5}$}{g-Delta-5}}
\label{sec:k3e-bkm-chiral-obstruction}

Three candidates present themselves for ``the $K3 \times E$ chiral algebra whose bar Euler product equals $1/\Delta_5$'':
\begin{enumerate}[label=(\alph*)]
\item $A_{K3 \times E}^{(\Sigma_2,C)} = \Phi_3^{(\Sigma_2,C)}(D^b(\mathrm{Coh}(K3 \times E)))$, the framed output of the CY-to-chiral assignment at $d = 3$ under H1--H4 and fixed specialisation (Theorem~\ref{thm:cy-to-chiral-d3}).
\item The relative chiral algebra $A_{K3,\mathrm{rel}}$ (Section~\ref{sec:k3e-relative-chiral}), constructed at the fiber level but character-incomplete at $d = 3$.
\item The BKM superalgebra $\mathfrak{g}_{\Delta_5}$ itself (Section~\ref{sec:k3e-bkm}), whose denominator identity produces $\Delta_5$.
\end{enumerate}
This section investigates candidate~(c) through five questions and concludes that $\mathfrak{g}_{\Delta_5}$ is \emph{not} a chiral algebra but rather encodes the root data that a conjectural chiral algebra must reproduce.

\subsection{Q1: vertex algebra structure}
\label{subsec:k3e-bkm-q1}

The BKM superalgebra $\mathfrak{g}_{\Delta_5}$ is a Lie superalgebra with generators $e_\alpha, f_\alpha, h_\alpha$ ($\alpha \in \Delta$) and BKM relations. It is \emph{not} a vertex algebra: it has no state space, no vacuum vector, no operator product expansion, no state-operator correspondence.

The relationship to vertex algebras is through Borcherds' BRST
template. The desired model has the form
\begin{equation}
\label{eq:bkm-brst}
 \mathfrak{g}_{\Delta_5}
 \;\cong\;
 H^1_{\mathrm{BRST}}\bigl(V_{\Lambda^{2,1}_{II}} \otimes V_{\mathrm{trans}} \otimes V_{\mathrm{ghost}}\bigr),
\end{equation}
where $V_{\Lambda^{2,1}_{II}}$ is the formal lattice vertex algebra of
the rank-$3$ type-II hyperbolic lattice (central charge \(c=3\)),
\(V_{\mathrm{trans}}\) is a supplied transverse vertex superalgebra of
central charge \(23\) whose graded supertrace has the K3 elliptic-genus
coefficients, and \(V_{\mathrm{ghost}}\) is the bosonic string ghost
system \((c=-26)\). The central-charge sum \(3+23-26=0\) is the
necessary BRST balance. It is not by itself a construction of
\(Q_{\mathrm{BRST}}\), a no-ghost theorem, or the comparison with the
framed \(K3\times E\) target. In the \(K3\times E\) specialisation,
these are precisely the missing inputs recorded in
Remark~\ref{rem:k3e-brst-missing}.

\begin{proposition}[Finite BRST central-charge gate]
\label{prop:k3e-brst-central-charge-gate}
In the finite BRST template
\[
V_{\Lambda^{2,1}_{II}}\otimes V_{\mathrm{trans}}\otimes
V_{\mathrm{ghost}},
\]
the Virasoro central-charge anomaly cancels if and only if
\[
c_{\Lambda}+c_{\mathrm{trans}}+c_{\mathrm{ghost}}=0.
\]
For the \(K3\times E\) Borcherds template,
\[
c_{\Lambda}=3,\qquad c_{\mathrm{ghost}}=-26,
\]
and therefore the transverse sector must have
\[
c_{\mathrm{trans}}=23.
\]
Equivalently, the finite central-charge gate is the equality
\[
3+23-26=0.
\]
\end{proposition}

\begin{proof}
The lattice vertex algebra of a rank-\(r\) lattice carries central charge
\(r\); here \(r=\operatorname{rk}\Lambda^{2,1}_{II}=3\).  The bosonic
string ghost system has central charge \(-26\).  Central charges add
under tensor product of vertex algebras, so the BRST template has total
central charge \(3+c_{\mathrm{trans}}-26\).  Relative BRST anomaly
cancellation requires this total to be zero, hence
\(c_{\mathrm{trans}}=23\).  The executable witness is
\texttt{brst\_central\_charge\_gate}; it records the required transverse
central charge, the total central charge, and the two defects.  This is
only the finite anomaly gate. It does not construct
\(V_{\mathrm{trans}}\), the BRST differential, a no-ghost theorem, or the
comparison with the framed \(K3\times E\) target.
\end{proof}

\begin{remark}[Exact missing inputs for the BRST realization]
\label{rem:k3e-brst-missing}
The BRST statement above is the classical template for the denominator
construction, but the K3$\times$E case still needs the following
inputs before the comparison can be treated as a theorem in this
chapter:
\begin{enumerate}[label=\textup{(\roman*)}]
\item a concrete transverse vertex superalgebra of central charge \(23\)
whose graded supertrace supplies the K3 elliptic-genus coefficients,
together with the lattice and ghost sectors and the correct ghost
number grading;
\item an identification of the BRST differential whose cohomology is
actually computed in the K3$\times$E specialisation, rather than merely
the formal Borcherds template;
\item a proof that the resulting cohomology realises the finite
supertrace fixture of
Proposition~\ref{prop:k3e-brst-finite-supertrace-fixture}, hence
reproduces the signed root characters \(c(D(\alpha))\) with the same
parity convention as the BKM denominator;
\item a compatibility theorem relating this BRST complex to the framed
CY-to-chiral assignment and to the finite Hall/Borcherds recognition
data.
\end{enumerate}
These are the exact missing steps between the formal template and the
K3$\times$E denominator comparison.
\end{remark}

Under such a transverse theory, the signed root-character exponents
\(\operatorname{sdim}\mathfrak g_{\Delta_5,\alpha}=c(D(\alpha))\) of
\(\mathfrak g_{\Delta_5}\) are read from the graded supertrace of the
\(c=23\) transverse sector, normalised to the K3 elliptic-genus
coefficients of \(\phi_{0,1}\).

\begin{remark}[Comparison with the Fake Monster]
\label{rem:k3e-fake-monster-comparison}
The template~\eqref{eq:bkm-brst} is modelled on Borcherds's
construction of the Fake Monster Lie algebra \(\mathfrak g_{\mathrm{FM}}\)
from the lattice vertex algebra \(V_{\mathrm{II}_{25,1}}\) of the
rank-\(26\) even unimodular Lorentzian lattice. In that case the
transverse central charge is zero (the lattice VOA already has
\(c=26\)), and all imaginary root multiplicities equal \(24\) (the
rank of the Leech lattice). For \(\mathfrak g_{\Delta_5}\), the
rank-\(3\) lattice vertex algebra contributes \(c=3\), so the missing
\(c=23\) must be supplied by the transverse theory above; the signed
root-character coefficients are the richer set \(\{c(D)\}\) from
\(\phi_{0,1}\), rather than the uniform \(24\). Ordinary root-space
dimensions require the target parity fixture.
\end{remark}

\subsection{Q2: the Borcherds vertex algebra realization}
\label{subsec:k3e-bkm-q2}

The Borcherds character formula produces a denominator identity for BKM algebras constructed from vertex-algebra BRST cohomology in the classical examples below. All three examples share the pattern:
\begin{center}
\renewcommand{\arraystretch}{1.3}
\begin{tabular}{@{}llll@{}}
\textbf{BKM algebra} & \textbf{Lattice} & \textbf{VOA ($c = 26$)} & \textbf{Root character} \\ \hline
Monster & $\mathrm{II}_{1,1}$ & $V^\natural \otimes V_{\mathrm{II}_{1,1}}$ & $c(mn)$ from $j - 744$ \\
Fake Monster & $\mathrm{II}_{25,1}$ & $V_{\mathrm{II}_{25,1}}$ & $24$ (uniform) \\
$\mathfrak{g}_{\Delta_5}$ & $\Lambda^{2,1}_{II}$ & $V_{K3 \times E} \otimes V_{\Lambda^{2,1}_{II}}$ & $c(D)$ from $\phi_{0,1}$
\end{tabular}
\end{center}
In all three cases, the BKM algebra is the BRST cohomology of the VOA, not the VOA itself. The vertex algebra sits at a different categorical level: it is the \emph{source} of the BRST construction, while the Lie superalgebra is the \emph{output}. This is a classical pattern, not a universal theorem beyond the cases cited here.

For candidate~(c) in the $K3 \times E$ problem: the ``chiral algebra'' is not $\mathfrak{g}_{\Delta_5}$ but rather the $c = 26$ vertex algebra $V_{K3 \times E} \otimes V_{\Lambda^{2,1}_{II}}$ (worldsheet CFT). The relationship between this worldsheet VOA and the target-space framed chiral algebra $A_{K3 \times E}^{(\Sigma_2,C)} = \Phi_3^{(\Sigma_2,C)}(D^b(\mathrm{Coh}(K3 \times E)))$ is mediated by string field theory; the precise comparison remains the CY-D problem at $d = 3$.

\subsection{Q3: CoHA and the positive half}
\label{subsec:k3e-bkm-q3}

\begin{remark}[CoHA is associative, not chiral]
\label{rem:bkm-coha-not-chiral}
The critical CoHA of $K3 \times E$ is an \emph{associative} algebra (the Hall product is the convolution product on Borel--Moore homology of moduli stacks). It is \emph{not} a chiral algebra, \emph{not} an $E_1$-chiral algebra, and \emph{not} an $E_2$-chiral algebra. The statement ``CoHA is the $E_1$-sector of $G(K3 \times E)$'' is therefore read through Theorem~CY-C: the CoHA supplies the positive-half Hall algebra, and the chiral group object is obtained only after Drinfeld double, centre, and comparison data are included.
\end{remark}

The positive-half target of Theorem~\ref{thm:k3e-coha} is
\[
 U(\mathfrak{n}_+(\mathfrak{g}_{\Delta_5})),
\]
the universal enveloping algebra of the positive
nilpotent subalgebra $\mathfrak{n}_+ = \bigoplus_{\alpha \in \Delta_+}
\mathfrak{g}_\alpha$. The Borcherds positive half recovers the root
multiplicities, but the Hall algebra perspective by itself provides no
OPE, no braiding, and no quantum group structure. The passage from CoHA
(associative) to the full quantum vertex chiral group (braided) uses the
Drinfeld double construction, the $E_1 \to E_2$ promotion of Volume~I,
with the seven comparison coordinates of
Theorem~\ref{thm:k3e-hall-drinfeld-super-yangian-criterion}.

\begin{remark}[Comparison with \texorpdfstring{$\mathbb{C}^3$}{C-3}]
\label{rem:bkm-c3-comparison}
For $\mathbb{C}^3$, the complete positive-half/double chain is proved: $\mathrm{CoHA}(\mathbb{C}^3) = Y^+(\widehat{\mathfrak{gl}}_1)$ (Schiffmann--Vasserot), and the Drinfeld double $D(Y^+)=Y(\widehat{\mathfrak{gl}}_1)$ acts on the $\mathcal{W}_{1+\infty}$ vacuum module. For $K3 \times E$, Theorem~\ref{thm:k3e-coha} gives the Borcherds positive-half theorem, and Theorem~\ref{thm:k3e-hall-drinfeld-super-yangian-criterion} supplies the comparison coordinates for the completed double.
\end{remark}

\subsection{Q4: shadow class}
\label{subsec:k3e-bkm-q4}

If the framed target object $A_{K3 \times E}^{(\Sigma_2,C)}$ is constructed and its bar Euler exponents match the BKM signed root characters, then the shadow class is \textbf{M} (infinite shadow depth):
\begin{itemize}
 \item The shadow invariants are the $\phi_{0,1}$ coefficients: $S_3 = c(3) = -64$ (first non-formality obstruction, fermionic), $S_4 = c(4) = 108$ (first higher bosonic root), and so on.
 \item The signed exponents have absolute growth $|c(D)|$ of order $\exp(2\pi\sqrt{D})$ (Rademacher asymptotics for weak Jacobi forms of index~$1$), producing infinitely many nonzero shadow invariants.
 \item The shadow tower does not terminate at any finite depth: for every valid discriminant $D \equiv 0$ or $3 \pmod{4}$, the coefficient $c(D) \neq 0$.
\end{itemize}
This is consistent with the mock modular behavior expected for class~M chiral algebras (Section~\ref{subsec:k3e-wc-shadow}).

\begin{remark}[shadow class from full tower]
\label{rem:bkm-shadow-class}
The class~M determination is obtained from the full shadow tower computation through the discriminants $D \leq 60$, not from generator counting or non-formality alone. The non-formality condition $m_3 \neq 0$ (i.e., $c(3) = -64 \neq 0$) is a \emph{necessary} condition for class $\geq L$ but does not by itself determine the shadow depth.
\end{remark}

\subsection{Q5: the denominator--bar Euler dictionary}
\label{subsec:k3e-bkm-q5}

The precise dictionary between the BKM denominator identity and the bar Euler product identifies four layers:

\begin{center}
\renewcommand{\arraystretch}{1.3}
\begin{tabular}{@{}p{4.8cm}p{4.8cm}l@{}}
\textbf{BKM side} & \textbf{Bar side} & \textbf{Comparison} \\ \hline
Weyl vector $\rho$ & Regularization vector $\rho_B$ & identify $\rho_B$ with $\rho$ \\[3pt]
$\operatorname{sdim}\mathfrak g_\alpha = c(D(\alpha))$ & $\chi(B(A)^\alpha)$ & prove signed Euler equality \\[3pt]
$\Delta_+ \subset \Lambda^{2,1}_{II}$ & $\Lambda$-grading on $B(A)$ & construct the numerical $K$-theory grading \\[3pt]
$\Phi(z) = e^{(\rho)} \prod (1 - e^{(\alpha)})^{\mult}$ & $\Theta_A(z) = e^{(\rho_B)} \prod (1 - e^{(\alpha)})^{\chi}$ & prove $\Phi=\Theta_A$
\end{tabular}
\end{center}

\noindent The identification requires three inputs: (1)~the framed H1--H4 construction of $A_{K3 \times E}^{(\Sigma_2,C)}$, (2)~the $\Lambda$-grading from numerical $K$-theory on the chiral algebra, and (3)~the motivic DT identification $\chi(B(A)^\alpha) = \Omega^{\mathrm{ind}}(\alpha) = c(D(\alpha))$ in CY-D$_3$.

\begin{remark}[denominator \texorpdfstring{$\neq$}{!=} bar Euler product]
\label{rem:bkm-denom-bar}
The exponents in the BKM denominator and in the bar Euler product agree:
both are \(c(D)\) from \(\phi_{0,1}\), and both products range over
positive roots. The mathematical objects are nevertheless different.
The BKM denominator is a Weyl--Kac--Borcherds character formula; the bar
Euler product is the Euler characteristic of a bar complex. The equality
of the two functions is the content of the framed \(d=3\) assignment
together with Conjecture~CY-B.
\end{remark}

\subsection{Conclusion}
\label{subsec:k3e-bkm-conclusion}

\begin{conjecture}[The BKM--bar dictionary]
\label{conj:bkm-bar-dictionary}
The framed target object $A_{K3 \times E}^{(\Sigma_2,C)}$ sits at the H1--H4 specialisation of Theorem~\ref{thm:cy-to-chiral-d3}. Its bar complex $B(A_{K3 \times E}^{(\Sigma_2,C)})$ is expected to be $\Lambda^{2,1}_{II}$-graded, and for each positive root $\alpha \in \Delta_+$ the Euler characteristic of the $\alpha$-primary bar piece should equal the BKM signed root character:
\[
 \chi(B(A_{K3 \times E}^{(\Sigma_2,C)})^{(\alpha)}) = \operatorname{sdim}\mathfrak g_{\Delta_5,\alpha} = c(D(\alpha)).
\]
The bar-complex regularization vector equals the Weyl vector: $\rho_B = \rho$. The bar Euler product $\Theta_{A_{K3 \times E}^{(\Sigma_2,C)}}$ then equals the BKM denominator function $\Phi = (1/64)\Delta_5(2Z)$.
\end{conjecture}

\begin{remark}[Exact missing inputs for the BKM--bar dictionary]
\label{rem:bkm-bar-missing-inputs}
The conjecture above still requires five concrete ingredients before it
can become a theorem:
\begin{enumerate}[label=\textup{(\roman*)}]
\item a construction of the framed chiral algebra on the fixed H1--H4
locus that is compatible with the bar complex used in Volume~I;
\item a proof that the bar complex is actually graded by the
Lorentzian lattice $\Lambda^{2,1}_{II}$ and that the regularization
vector equals the Weyl vector on the nose, not just numerically;
\item a comparison theorem identifying the $\alpha$-primary Euler
characteristic of the bar piece with the motivic DT index
\(\Omega^{\mathrm{ind}}(\alpha)\);
\item a proof that the resulting product expansion matches the
Borcherds denominator normalization \((1/64)\Delta_5(2Z)\), including
the multiplier system and the $64$-divisibility of the Fourier data;
\item a compatibility theorem with the shadow tower and the finite
Hall--Borcherds recognition data, so that the comparison survives
height truncation and inverse limits.
\end{enumerate}
These are the exact gaps between the structural observation and a
proof of the dictionary.
\end{remark}

\begin{proposition}[Finite bar lattice-grading gate]
\label{prop:k3e-finite-bar-lattice-grading-gate}
Let \(G=((\delta_i,\delta_j))\) be the real-simple-root Gram matrix of
\(\Lambda^{2,1}_{II}\):
\[
G=\begin{pmatrix}2&-2&-2\\ -2&2&-2\\ -2&-2&2\end{pmatrix}.
\]
Suppose that a finite BKM--bar comparison supplies, for each retained
bar primitive \(b_\gamma\), a coordinate vector
\[
v_\gamma=(v_{\gamma,1},v_{\gamma,2},v_{\gamma,3})
\in \mathbb Q^3
\]
in the basis \((\delta_1,\delta_2,\delta_3)\), together with the
expected Borcherds discriminant \(D(\gamma)\). The supplied finite
grading is a \(\Lambda^{2,1}_{II}\)-grading on the retained bar
primitives if and only if
\[
v_\gamma\in\mathbb Z^3,
\qquad
D(\gamma)=-\frac12\,v_\gamma^{\top}Gv_\gamma
\]
for every retained charge \(\gamma\). If the retained set is
downward-saturated in height, this condition is compatible with the
height projection \(H+1\to H\) exactly when deleting charges of height
\(>H\) preserves the same coordinate and discriminant rows.
\end{proposition}

\begin{proof}
The lattice \(\Lambda^{2,1}_{II}\) is generated by the real simple roots
\(\delta_i\) with Gram matrix \(G\). A supplied degree lies in the
lattice, rather than only in its rational span, exactly when its
coordinates in this basis are integral. The BKM discriminant convention
used throughout the chapter is
\[
D(v)=-\frac12(v,v).
\]
Writing \(v=v_1\delta_1+v_2\delta_2+v_3\delta_3\) gives
\((v,v)=v^{\top}Gv\), hence the displayed equality is precisely the
condition that the supplied bar degree has the expected Borcherds
discriminant. Height compatibility adds no further algebra: the lower
finite window is obtained by deleting the rows whose height exceeds
\(H\). Thus the finite grading survives projection exactly when the
remaining rows are unchanged. The executable witness is
\texttt{finite\_bar\_lattice\_grading\_report}; it records the norm,
computed discriminant, integrality defect, discriminant defect, and
height-projection gate for each retained charge.
\end{proof}

\begin{proposition}[Finite bar--Chevalley differential square]
\label{prop:k3e-finite-bar-ce-chain-map-gate}
Work in a fixed finite height window \(H\). Suppose that a finite
BKM--bar comparison supplies the length-two and length-one pieces
\[
 B_2^{\leq H}\xrightarrow{\,d_B\,}B_1^{\leq H},
 \qquad
 CE_2^{\leq H}\xrightarrow{\,d_{CE}\,}CE_1^{\leq H},
\]
together with finite comparison maps
\[
 F_2\colon B_2^{\leq H}\to CE_2^{\leq H},
 \qquad
 F_1\colon B_1^{\leq H}\to CE_1^{\leq H}.
\]
Then the length-two bar--Chevalley square commutes if and only if
\[
 F_1d_B=d_{CE}F_2.
\]
After finite bases are chosen, this is equivalent to rank zero for the
matrix
\[
 F_1D_B-D_{CE}F_2.
\]
When this rank is zero, the finite length-two part of the comparison is
a chain map. The remaining obstruction to a bar--Chevalley coalgebra
map is then the finite Serre/PBW defect measuring whether the bracket
matrix is the Borcherds bracket matrix with no extra primitive
relations.
\end{proposition}

\begin{proof}
Both composites in the displayed equation are linear maps
\[
B_2^{\leq H}\longrightarrow CE_1^{\leq H}.
\]
The left composite first takes the bar differential and then applies
the primitive comparison; the right composite first compares the
length-two bar generators with the Chevalley--Eilenberg length-two
generators and then applies the Chevalley--Eilenberg differential. Thus
commutativity of the finite square is exactly equality of these two
linear maps. In finite bases, equality of linear maps is equivalent to
vanishing of the difference matrix, hence to rank zero. The last
sentence is the standard cofree-coalgebra reduction used in
Theorem~\ref{thm:k3e-finite-bar-ce-comparison}: the ordered bar
differential is generated by its length-two bracket component, while
extra primitive source relations are precisely the finite Serre/PBW
defect. The executable witness is
\texttt{finite\_bar\_ce\_chain\_map\_gate}; it records the two
composites and the exact rank of their difference.
\end{proof}

\begin{remark}[Boundary of the finite bar--Chevalley square]
\label{rem:k3e-finite-bar-ce-chain-map-boundary}
Proposition~\ref{prop:k3e-finite-bar-ce-chain-map-gate} starts with
finite differentials and comparison maps already supplied. It does not
construct the framed algebra \(A_{K3\times E}^{(\Sigma_2,C)}\), the
ordered bar complex, or the filtered morphism
\(\Xi_{\mathrm{bar}}^{\leq H}\). It is the finite algebra test that a
candidate length-two comparison must pass before the phrase
``commutes with the bar differential'' is meaningful.
\end{remark}

\begin{theorem}[Finite bar--Chevalley comparison after recognition]
\label{thm:k3e-finite-bar-ce-comparison}
Fix a finite height \(H\), and let
\[
 B^{\mathrm{ord},\leq H}\!\bigl(A_{K3\times E}^{(\Sigma_2,C)}\bigr)
\]
be the ordered \(E_1\)-bar coalgebra of the framed algebra, truncated
by the same Borcherds-height window as the finite Hall packets. Assume:
\begin{enumerate}[label=\textup{(\roman*)}]
\item the finite recognition certificate of
Theorem~\ref{thm:k3e-finite-recognition-certificate} holds at height
\(H\);
\item the framed specialisation supplies a
\(\Lambda_{\mathrm{II}}^{2,1}\)-graded filtered bar morphism
\[
 \Xi_{\mathrm{bar}}^{\leq H}\colon
 B^{\mathrm{ord},\leq H}\!\bigl(A_{K3\times E}^{(\Sigma_2,C)}\bigr)
 \longrightarrow
 CE_\bullet(\fn_+(\mathfrak g_{\Delta_5}^{\leq H}))
\]
which sends the suspended primitive bar generator in charge \(\gamma\)
to the Chevalley--Eilenberg generator \(s e_{\alpha_\gamma}\);
\item the finite bar differential in length two is the OPE/Hall
bracket transported by the comparison matrix
\(\mathsf A_H^{X\to\Delta}\).
\end{enumerate}
Then \(\Xi_{\mathrm{bar}}^{\leq H}\) is a chain coalgebra map if and
only if the finite Serre/PBW defect \(\mathcal S_H\) vanishes. On this
locus the \(\gamma\)-primary bar Euler exponent is
\[
 \chi\!\left(
 B^{\mathrm{ord},\leq H}
 \bigl(A_{K3\times E}^{(\Sigma_2,C)}\bigr)_\gamma
 \right)
 =
 \operatorname{sdim}\mathfrak g_{\Delta_5,\alpha_\gamma}
 =
 c(D(\gamma))
\]
for every retained charge. If the height transitions commute with the
finite recognition certificates and with the filtered bar morphisms,
the same equality is natural under \(H+1\to H\).
\end{theorem}

\begin{proof}
The ordered bar coalgebra is cofree on the suspension of the framed
algebra. A coalgebra morphism from it is therefore determined by its
projection on suspended primitive generators. By hypothesis this
projection sends \(s^{-1}a_\gamma\) to \(s e_{\alpha_\gamma}\).
The only obstruction to commuting with the differentials is the
length-two component, because the higher bar differential is generated
as a coderivation. On length two the source differential is the
OPE/Hall bracket transported through the finite comparison matrix,
while the target Chevalley--Eilenberg differential is the
Borcherds bracket:
\[
 d_{CE}(s e_{\alpha}\,s e_{\beta})
 =
 s[e_\alpha,e_\beta].
\]
Thus the differential square commutes exactly when the transported
source primitive relations equal the finite Borcherds--Serre relations
and no additional primitive relation survives. This is precisely
\(\mathcal S_H=0\) in
Theorem~\ref{thm:k3e-finite-height-hall-radical-serre-kernel}.

When \(\mathcal S_H=0\), the PBW normal forms of the finite source and
the finite Serre--Borcherds quotient agree. The finite recognition
certificate also identifies the parity fixture and the signed primitive
character. Hence the Euler exponent of the \(\gamma\)-primary bar
piece is the signed dimension of the corresponding
\(\mathfrak g_{\Delta_5}\)-root space, which is \(c(D(\gamma))\) by the
Borcherds denominator theorem. Naturality in height follows by applying
the same argument after deleting all generators and relations of height
\(>H\); downward saturation makes this deletion a chain-coalgebra
projection on both sides.
\end{proof}

\begin{proposition}[Finite bar regularization and Weyl-vector gate]
\label{prop:k3e-finite-bar-regularization-gate}
Let \(G=((\delta_i,\delta_j))_{1\leq i,j\leq 3}\) be the real-simple-root
Gram matrix
\[
G=\begin{pmatrix}2&-2&-2\\ -2&2&-2\\ -2&-2&2\end{pmatrix}.
\]
Suppose that a finite BKM--bar comparison supplies a bar
regularization vector
\[
\rho_B=x_1\delta_1+x_2\delta_2+x_3\delta_3
\in \Lambda^{2,1}_{II}\otimes\mathbb Q
\]
and a scalar \(\nu_B\) for the leading Borcherds-product
normalization. Then the finite regularization gate for the
\(\Delta_5\) denominator is exactly
\[
Gx=-\frac{1}{2}\operatorname{diag}(G)=(-1,-1,-1)^{\top},
\qquad
\nu_B=\frac{1}{64}.
\]
Equivalently,
\[
\rho_B=\rho=\frac12(\delta_1+\delta_2+\delta_3),
\qquad
\nu_B=\frac{1}{64}.
\]
This is the finite obstruction at the Weyl-vector and leading-scalar
level of the BKM--bar dictionary.
\end{proposition}

\begin{proof}
The Weyl-vector equations for a Borcherds denominator in the chamber
with real simple roots \(\delta_i\) are
\[
(\rho_B,\delta_i)=-\frac{1}{2}(\delta_i,\delta_i)
\qquad (i=1,2,3).
\]
In the basis \((\delta_1,\delta_2,\delta_3)\) this is the displayed
linear system. The matrix \(G\) has determinant \(-32\), hence the
solution is unique. Direct substitution gives
\[
G\begin{pmatrix}1/2\\1/2\\1/2\end{pmatrix}
=
\begin{pmatrix}-1\\-1\\-1\end{pmatrix}.
\]
Thus the supplied bar regularization vector matches the Borcherds Weyl
vector if and only if its coordinate vector is
\((1/2,1/2,1/2)\). The scalar condition is the fixed normalization
\(\Phi=(1/64)\Delta_5(2Z)\) used in
Conjecture~\ref{conj:bkm-bar-dictionary}. The executable witness is
\texttt{finite\_bar\_regularization\_report}; it solves the rational
linear system, records the pairing defect
\(G\rho_B+\frac12\operatorname{diag}(G)\), and separately records the
normalization defect. It does not construct \(\rho_B\) from the framed
bar complex.
\end{proof}

\begin{remark}[Executable finite BKM--bar gates]
The routine \texttt{finite\_bar\_lattice\_grading\_report} in
\texttt{compute/lib/k3e\_finite\_bridge\_witness.py} records the
finite lattice-grading gate of
Proposition~\ref{prop:k3e-finite-bar-lattice-grading-gate}: the bar
degree must have integral \(\Lambda^{2,1}_{II}\)-coordinates, the
quadratic discriminant must equal the supplied Borcherds discriminant,
and height projection must preserve the finite rows.
The routine \texttt{finite\_bar\_ce\_report} in
\texttt{compute/lib/k3e\_finite\_bridge\_witness.py} records the two
finite checks in Theorem~\ref{thm:k3e-finite-bar-ce-comparison}: the
bar Euler exponent must equal the Borcherds coefficient \(c(D)\), and
the finite bar differential must commute with the Chevalley--Eilenberg
differential after the comparison matrix is fixed. It also checks that
height projection gives the same finite rows as recomputation at the
lower height. The routine does not construct the framed algebra or the
filtered bar morphism. The routine
\texttt{finite\_bar\_regularization\_report} records the independent
finite Weyl-vector and normalization gate of
Proposition~\ref{prop:k3e-finite-bar-regularization-gate}.
\end{remark}

The superalgebra $\mathfrak{g}_{\Delta_5}$ is \emph{not} itself a chiral algebra: it is a Lie superalgebra arising as the BRST cohomology of a vertex algebra. It encodes the \emph{root data} (multiplicities, parity, Weyl group, denominator identity) that the framed chiral algebra $A_{K3 \times E}^{(\Sigma_2,C)}$ must reproduce through its bar complex. The BKM algebra is the \emph{target} of the identification, not the source. The source is the framed CY-to-chiral assignment $\Phi_3^{(\Sigma_2,C)}$; the remaining open problem is the Vol~I Borcherds-lift bridge identifying the shadow tower with the BKM denominator in CY-B/CY-D.

\subsection{Borcherds vertex algebra approach to imaginary root generators}
\label{subsec:k3e-borcherds-vertex-yangian}

The structural obstruction identified in Section~\ref{subsec:k3e-bkm-conclusion} (no \emph{strict} Drinfeld $J$-presentation with Killing cubic Casimir for BKM algebras with imaginary simple roots; Theorem~\ref{thm:bkm-drinfeld-J-pro-existence}(i)) motivates the vertex-algebra route. The pro-cone-graded $J$-presentation with Siegel--Borcherds-twisted cubic of Theorem~\ref{thm:bkm-drinfeld-J-pro-existence}(ii) is the algebraic target that the vertex-operator construction below is designed to reach. The imaginary root vectors of $\mathfrak{g}_{\Delta_5}$ are represented in the formal Borcherds lattice-VOA template by vertex operators in \(V_{\Lambda^{2,1}_{II}}\). A Yangian construction from this route still requires the all-orders deformation and residue comparison isolated in Remarks~\ref{rem:k3e-vertex-op-yangian-missing} and~\ref{rem:k3e-yangian-current-missing}.

\subsubsection{Vertex operators for imaginary roots}
\label{subsubsec:k3e-vertex-ops-imaginary}

In the formal BRST template~\eqref{eq:bkm-brst}, each root vector \(e_\alpha \in \mathfrak{g}_{\Delta_5}\) is represented by a physical state in
\(V_{\Lambda^{2,1}_{II}} \otimes V_{\mathrm{trans}} \otimes V_{\mathrm{ghost}}\).
For a root $\alpha$ at discriminant $D$ (with $(\alpha, \alpha) = -2D$ in the Lorentzian convention), the vertex operator
\begin{equation}
\label{eq:bkm-vertex-op}
V(\alpha, z) = {:}\exp(\alpha \cdot \phi(z)){:}
\end{equation}
creates a state at conformal weight $h = D + 1$ (the level of the root in the BKM algebra). The self-OPE structure depends on the discriminant:
\begin{center}
\renewcommand{\arraystretch}{1.3}
\begin{tabular}{@{}llll@{}}
\textbf{Discriminant} & \textbf{Self-OPE} & $\operatorname{sdim}$ & \textbf{Interpretation} \\ \hline
$D = -1$ & regular (zero of order $2$) & $1$ & Weyl vector root \\
$D = 0$ & regular & $10$ & lightlike ($\kappa_{\mathrm{BKM}}(\Delta_5) = 5$) \\
$D > 0$ & pole of order $2D$ & $c(D)$ & spacelike imaginary
\end{tabular}
\end{center}
The table records the signed root-character exponent. Ordinary current
counts enter only after the target parity fixture used in
\eqref{eq:yangian-current-def}.

\subsubsection{Omega-background deformation}
\label{subsubsec:k3e-omega-deformation}

\begin{conjecture}[Vertex operator Yangian generators]
\label{conj:vertex-op-yangian}
The Omega-background deformation (the equivariant $T^2$-action with parameter $\varepsilon$ on the worldsheet, in the sense of Nekrasov) provides a one-parameter family of deformed vertex operators
\begin{equation}
\label{eq:deformed-vertex-op}
V_\varepsilon(\alpha, z) = V(\alpha, z) \cdot \exp\Bigl(\sum_{k \geq 1} \varepsilon^k\, O_k(\alpha, z)\Bigr),
\end{equation}
where $O_k$ are correction operators built from normal-ordered products of $\phi$ and its derivatives. At leading order the deformation is a \emph{spectral flow}:
\[
O_1(\alpha, z) = -D \cdot \partial_z \phi(z),
\]
shifting the conformal weight by $-\varepsilon D$. The deformed conformal weight at formal $\varepsilon = 1$ is
\[
h_\varepsilon = (D + 1) - D \cdot \varepsilon \;\;\xrightarrow{\;\varepsilon = 1\;}\;\; 1
\]
for every discriminant $D$. All imaginary root vertex operators become weight-$1$ currents: precisely the structure required for Yangian generators.
\end{conjecture}

\begin{remark}[Exact missing inputs for the vertex-operator Yangian route]
\label{rem:k3e-vertex-op-yangian-missing}
The conjecture above still requires the following ingredients before it
can be upgraded to a theorem:
\begin{enumerate}[label=\textup{(\roman*)}]
\item a worldsheet VOA/BRST construction of the $K3 \times E$ lattice
sector with the claimed transverse and ghost contributions;
\item a proof that the $\Omega$-background deformation exists to all
orders as a convergent or Borel-summable deformation at the formal
specialisation $\varepsilon = 1$;
\item a proof that the deformed residues \(e_D^{(a)}(u)\) realise the
finite current-candidate packet of
Proposition~\ref{prop:k3e-finite-yangian-current-candidate-packet} as
operators on BRST cohomology, not merely as a formal spectral-mode
basis;
\item a compatibility theorem relating the monodromy of the deformed
vertex operators to the Hall--Borcherds $R$-matrix and the finite-height
recognition data.
\end{enumerate}
These are the exact gaps between the formal vertex-operator model and a
Yangian construction.
\end{remark}

The weight collapse \(h_\varepsilon|_{\varepsilon=1}=1\) for all \(D\)
is the necessary endpoint identity for any Yangian construction from
this deformation. The missing analytic input is convergence of the
perturbation series at \(\varepsilon=1\). At \(D=0\) the
spectral-flow correction vanishes, so the lightlike vertex operators are
\emph{already} weight-\(1\) currents. For \(D<0\) (the Weyl vector root at
\(D=-1\)), spectral flow increases the weight from \(0\) to \(1\).

\subsubsection{Yangian currents and deformed Serre relations}
\label{subsubsec:k3e-yangian-currents-serre}

\begin{proposition}[Finite Yangian current-candidate packet]
\label{prop:k3e-finite-yangian-current-candidate-packet}
Fix a finite discriminant bound \(H\) and a finite spectral-mode bound
\(M\). Let \(T_D^{\min}\) be the finite transverse supertrace fixture
of Proposition~\ref{prop:k3e-brst-finite-supertrace-fixture}. Define
\[
 \mathcal E_{\leq H,\leq M}^{\mathrm{cand}}
 =
 \bigoplus_{\substack{D\leq H\\ c(D)\neq0}}
 T_D^{\min}\otimes
 \mathrm{span}_{\mathbb C}\{u^{-1},u^{-2},\ldots,u^{-M-1}\}.
\]
Equivalently, a homogeneous basis is denoted
\[
 e_{D,r}^{(a)},\qquad
 D\leq H,\quad 0\leq r\leq M,\quad 1\leq a\leq |c(D)|.
\]
The parity of \(e_{D,r}^{(a)}\) is the parity of \(T_D^{\min}\), and
the formal spectral-flow endpoint gives
\[
 h_{\varepsilon=1}(e_{D,r}^{(a)})=(D+1)-D=1.
\]
Thus \(\mathcal E_{\leq H,\leq M}^{\mathrm{cand}}\) is a finite
super-vector space of weight-one current candidates with row dimension
\(|c(D)|(M+1)\) and row superdimension \(c(D)(M+1)\). If
\(H'\leq H\) and \(M'\leq M\), the projection onto \(D\leq H'\) and
\(0\leq r\leq M'\) is strict; these packets therefore form a
two-parameter inverse system under discriminant and mode truncation.
\end{proposition}

\begin{proof}
The construction is the tensor product of the finite fixture
\(\bigoplus_{D\leq H}T_D^{\min}\) with the finite spectral-mode space
\(\mathrm{span}\{u^{-1},\ldots,u^{-M-1}\}\). Hence each discriminant
row has \(|c(D)|(M+1)\) ordinary basis vectors and superdimension
\(\operatorname{sdim}(T_D^{\min})(M+1)=c(D)(M+1)\). The equality
\((D+1)-D=1\) is the formal endpoint of the spectral-flow formula
displayed above. The truncation maps are coordinate projections, so
they preserve coefficients, parity, and the endpoint weight. The
executable witness is
\texttt{yangian\_current\_candidate\_packet} in
\texttt{compute/lib/k3e\_finite\_bridge\_witness.py}; its transition
check \texttt{yangian\_current\_packet\_transition} verifies the
\((H,M)=(8,3)\to(4,1)\) truncation.
\end{proof}

\begin{proposition}[Finite spectral kernel label packet]
\label{prop:k3e-finite-spectral-kernel-label-packet}
Fix finite bounds \(H\) and \(M\), and let
\(\mathcal E_{\leq H,\leq M}^{\mathrm{cand}}\) be the current-candidate
packet of Proposition~\ref{prop:k3e-finite-yangian-current-candidate-packet}.
Let \((\mathcal E_{\leq H,\leq M}^{\mathrm{cand}})^\vee\) be a formal
Serre-dual negative packet with basis labels \(f_{D,r}^{(a)}\) dual to
\(e_{D,r}^{(a)}\), with the same parity. For each retained label set
define the non-Cartan spectral-kernel label
\[
 \rho_{D,r}^{(a)}
 =
 e_{D,r}^{(a)}\otimes f_{D,r}^{(a)}
 +(-1)^{|e_{D,r}^{(a)}|}
 f_{D,r}^{(a)}\otimes e_{D,r}^{(a)}.
\]
Then the non-Cartan finite kernel label packet
\[
 \mathcal R^{\mathrm{lab,nc}}_{\leq H,\leq M}
 =
 \mathrm{span}_{\mathbb C}\{\rho_{D,r}^{(a)}:
 D\leq H,\ c(D)\neq0,\ 0\leq r\leq M,\ 1\leq a\leq |c(D)|\}
\]
has
\[
 \dim \mathcal R^{\mathrm{lab,nc}}_{\leq H,\leq M}
 =
 \sum_{\substack{D\leq H\\c(D)\neq0}} |c(D)|(M+1)
 =
 \dim \mathcal E_{\leq H,\leq M}^{\mathrm{cand}},
\]
and its displayed tensor expansion has exactly
\[
 2\sum_{\substack{D\leq H\\c(D)\neq0}} |c(D)|(M+1)
\]
non-Cartan tensor monomials. If a finite Cartan label space
\(\mathcal R^{\mathrm{lab,Cart}}_{H,M}\) of dimension \(\rho_{H,M}\)
is supplied externally, the total label count is
\[
 \dim \mathcal R^{\mathrm{lab,nc}}_{\leq H,\leq M}+\rho_{H,M}.
\]
For \(H'\leq H\), \(M'\leq M\), and
\(\rho_{H',M'}\leq \rho_{H,M}\), projection on discriminant, mode, and
Cartan labels gives a strict finite transition map.
\end{proposition}

\begin{proof}
The basis labels of
\(\mathcal E_{\leq H,\leq M}^{\mathrm{cand}}\) are precisely
\((D,r,a)\) with \(D\leq H\), \(c(D)\neq0\), \(0\leq r\leq M\), and
\(1\leq a\leq |c(D)|\). The formal negative packet is by definition
the dual label packet with the same index set and parity. Hence each
current label gives exactly one non-Cartan kernel label
\(\rho_{D,r}^{(a)}\), and the ordinary label count is the ordinary
dimension of the current packet. The displayed symmetrised expression
has two tensor monomials, one in positive--negative order and one in
negative--positive order, with sign determined by the parity of the
current label. The optional Cartan term is not inferred from the
current packet; it contributes only its supplied finite dimension
\(\rho_{H,M}\). The transition maps are coordinate projections on the
three finite index sets. The executable witness is
\texttt{yangian\_spectral\_kernel\_label\_packet}; its transition check
is \texttt{yangian\_spectral\_kernel\_transition}.
\end{proof}

\begin{remark}[Boundary of the finite spectral kernel packet]
\label{rem:k3e-finite-spectral-kernel-boundary}
Proposition~\ref{prop:k3e-finite-spectral-kernel-label-packet}
constructs only the finite label packet underlying the formal expression
for a current kernel. It does not construct a Hopf pairing, a Cartan
kernel, a spectral connection, an operator-valued \(r\)-matrix, a
solution of the classical Yang--Baxter equation, or the
Hall--Borcherds \(R\)-matrix comparison.
\end{remark}

\begin{proposition}[Finite formal self-OPE pole-layer packet]
\label{prop:k3e-finite-self-ope-pole-layer-packet}
Fix finite bounds \(H\) and \(M\), and let
\[
 U_{\leq M}=\mathrm{span}_{\mathbb C}\{u^{-1},u^{-2},\ldots,u^{-M-1}\}.
\]
For each retained discriminant \(D\leq H\) with \(c(D)\neq0\), set
\[
 P_{\mathrm{self}}(D)=D(D-4),\qquad
 p(D)=\max\{0,-P_{\mathrm{self}}(D)\}.
\]
Define the ordered self-OPE pole-label row
\[
 \mathcal P^{\mathrm{self}}_{D,\leq M}
 =
 \bigl(T_D^{\min}\otimes U_{\leq M}\bigr)
 \otimes
 \bigl(T_D^{\min}\otimes U_{\leq M}\bigr)
 \otimes \mathbb C^{p(D)}.
\]
Then
\[
 \dim \mathcal P^{\mathrm{self}}_{D,\leq M}
 =
 |c(D)|^2(M+1)^2p(D),\qquad
 \operatorname{sdim}\mathcal P^{\mathrm{self}}_{D,\leq M}
 =
 c(D)^2(M+1)^2p(D).
\]
For the low discriminant window \(H=8\), the only retained
discriminant with \(p(D)>0\) is \(D=3\); it has \(p(3)=3\).
The discriminants \(D=0\) and \(D=4\) are marginal, and
\(D=-1,7,8\) are regular for this self-OPE exponent. The formal target
of the top vertex \(V_\varepsilon(2\alpha)\) in the \(D=3\) row has
discriminant \(4D=12\), and the coefficient table gives
\(c(12)=4016\). If \(H'\leq H\) and \(M'\leq M\), projection to
\(D\leq H'\) and \(0\leq r\leq M'\) carries
\(\mathcal P^{\mathrm{self}}_{D,\leq M}\) to
\(\mathcal P^{\mathrm{self}}_{D,\leq M'}\) for retained \(D\) and
annihilates the discarded rows. Thus the finite pole-label packets form
a strict two-parameter inverse system.
\end{proposition}

\begin{proof}
The exponent \(P_{\mathrm{self}}(D)=D(D-4)\) is negative exactly for
\(0<D<4\). Among the retained valid discriminants
\(-1,0,3,4,7,8\), this leaves only \(D=3\), with
\(P_{\mathrm{self}}(3)=-3\). The rows \(D=0\) and \(D=4\) have exponent
zero; the remaining retained rows have positive exponent. The tensor
definition then gives the ordinary dimension by multiplication of
ordinary dimensions, since \(T_D^{\min}\) has ordinary dimension
\(|c(D)|\) and \(U_{\leq M}\) has dimension \(M+1\). The
superdimension is multiplicative on tensor products of finite
super-vector spaces, hence
\(\operatorname{sdim}(T_D^{\min})^2(M+1)^2p(D)=c(D)^2(M+1)^2p(D)\).
The discriminant of \(2\alpha\) is \(4D\) because
\((2\alpha,2\alpha)=4(\alpha,\alpha)\); at \(D=3\) this gives
\(D=12\), whose signed coefficient in the \(K3\times E\) Borcherds
table is \(4016\). The transition assertion is the coordinate
projection on the two finite factors. The executable witness is
\texttt{yangian\_self\_ope\_pole\_layer\_packet}; the transition check
is \texttt{yangian\_self\_ope\_pole\_transition}.
\end{proof}

\begin{remark}[Boundary of the finite pole-layer packet]
\label{rem:k3e-finite-pole-layer-boundary}
Proposition~\ref{prop:k3e-finite-self-ope-pole-layer-packet}
constructs the finite label spaces in which principal self-OPE
coefficients would live. It does not construct those coefficients, the
BRST residue operators, the super-skew quotient, the Serre ideal, a
coproduct, or the Hall--Borcherds \(R\)-matrix comparison. These remain
the operator-level inputs listed in
Remark~\ref{rem:k3e-yangian-current-missing}.
\end{remark}

\begin{proposition}[Finite Yangian label tower]
\label{prop:k3e-finite-yangian-label-tower}
In the level-\(2\) formal-current boundary of the reduced
\(K3\times E\) chart, and in the finite ordinary super-vector-space
ambient, the three finite packets
\[
 \mathcal E_{\leq H,\leq M}^{\mathrm{cand}},\qquad
 \mathcal R^{\mathrm{lab,nc}}_{\leq H,\leq M},\qquad
 \bigoplus_{\substack{D\leq H\\c(D)\neq0}}
 \mathcal P^{\mathrm{self}}_{D,\leq M}
\]
form a strict two-parameter inverse system in \((H,M)\), after a
finite Cartan label space for
\(\mathcal R^{\mathrm{lab,Cart}}_{H,M}\) has been supplied with
surjective truncation maps. More precisely, if
\[
 H'\leq H,\qquad M'\leq M,\qquad \rho_{H',M'}\leq \rho_{H,M},
\]
then coordinate projection in discriminant, spectral mode, and Cartan
labels sends the height-\((H,M)\) current-candidate packet, non-Cartan
kernel-label packet, and formal self-OPE pole-label packet onto the
corresponding height-\((H',M')\) packets, preserving signed
coefficients, parity, endpoint weight, pole order, target discriminant,
and the non-Cartan positive/negative-dual label pairing.
\end{proposition}

\begin{proof}
Each packet is built from the same finite index set
\[
 \{(D,r,a):D\leq H,\ c(D)\neq0,\ 0\leq r\leq M,\ 1\leq a\leq |c(D)|\},
\]
with extra finite factors: a dual copy for the kernel labels, the
Cartan label set supplied externally, and the pole-layer factor
\(\mathbb C^{p(D)}\) for the self-OPE packet. Projection to
\(D\leq H'\) and \(0\leq r\leq M'\) preserves \(c(D)\), parity,
\((D+1)-D=1\), \(p(D)=\max\{0,-D(D-4)\}\), and the target
discriminant \(4D\). The non-Cartan kernel row
\(e\otimes f+(-1)^{|e|}f\otimes e\) is defined labelwise, so the same
projection preserves its positive/negative-dual pairing and parity
sign. The supplied Cartan transition is a coordinate projection by the
assumption \(\rho_{H',M'}\leq\rho_{H,M}\). Thus the three packets and
their comparison labels commute with the finite truncation maps. The
executable witness is \texttt{yangian\_label\_tower\_transition}, which
combines \texttt{yangian\_current\_packet\_transition},
\texttt{yangian\_spectral\_kernel\_transition}, and
\texttt{yangian\_self\_ope\_pole\_transition}. It records the three
component statuses, component defect tuples, commutation gates, and
finite size data separately, so a failed tower transition identifies
the current, kernel-label, or pole-label component in which the finite
inverse-system condition broke.
\end{proof}

\begin{remark}[Boundary of the finite label tower]
\label{rem:k3e-finite-yangian-label-tower-boundary}
The inverse system in
Proposition~\ref{prop:k3e-finite-yangian-label-tower} is a finite label
system. Operator residues, OPE coefficients, the Serre quotient,
coproduct, spectral connection, associator cochain, and
Hall--Borcherds \(R\)-matrix comparison enter only after the additional
operator-level data named in
Remarks~\ref{rem:k3e-yangian-current-missing} and
\ref{rem:k3e-spectral-obstruction-boundary}.
\end{remark}

If the \(\Omega\)-deformed vertex operators and BRST-residue map are
constructed, the formal Yangian current at discriminant \(D\) is the
residue
\begin{equation}
\label{eq:yangian-current-def}
e_D^{(a)}(u) = \Res_{z = u}\, V_\varepsilon(\alpha_a, z), \qquad a = 1, \ldots, |c(D)|,
\end{equation}
where \(u\) is the Yangian \emph{spectral} parameter. The finite packet
above constructs the underlying current-candidate vector space; the
residue formula constructs operators only after the missing
BRST/Omega input is supplied. The commutation relations
\[
[e_D^{(a)}(u),\, e_{D'}^{(b)}(v)] = f_{D,D'}^{ab}(u - v)
\]
would be derived from the deformed OPE of~\eqref{eq:deformed-vertex-op},
not postulated from a Cartan matrix.

\begin{remark}[Exact missing inputs for the Yangian-current and Serre comparison]
\label{rem:k3e-yangian-current-missing}
The current formula above remains formal until the following inputs are
constructed and checked:
\begin{enumerate}[label=\textup{(\roman*)}]
\item a proof that the deformed residues \(e_D^{(a)}(u)\) are
well-defined operators on the claimed BRST or vertex-algebra
cohomology, rather than formal Laurent coefficients only;
\item a derivation of the commutator coefficients \(f_{D,D'}^{ab}\)
from the deformed OPE that is compatible with the target parity
fixture and with the finite discriminant multiplicities;
\item a theorem identifying the order-by-order OPE relations with the
Serre ideal of the completed Borcherds--OPE current algebra;
\item a compatibility theorem matching the monodromy of the deformed
vertex operators with the Hall--Borcherds \(R\)-matrix and the finite-height
comparison data.
\end{enumerate}
These are the exact missing steps between the formal current ansatz and
a Yangian presentation.
\end{remark}

\begin{proposition}[Finite OPE--Serre ideal span gate]
\label{prop:k3e-finite-yangian-ope-serre-ideal-span-gate}
Fix finite bounds \((H,M)\) and a finite word-coordinate space
\(W_{H,M}\) for the retained Yangian current words in the
\((K3\times E,\Sigma_2,C)\) chart. Suppose the singular OPE
calculation has supplied relation rows
\[
 A_{H,M}\subset W_{H,M},
\]
one row for each retained order-by-order OPE associativity relation,
and suppose the target Borcherds presentation has supplied relation
rows
\[
 S_{H,M}\subset W_{H,M},
\]
one row for each retained real-root Serre, imaginary-root
supercommutativity, or declared Siegel--Borcherds current relation.
Then the finite OPE relations generate exactly the retained finite
Borcherds--Serre current ideal in this word window if and only if
\[
 \operatorname{span} A_{H,M}=\operatorname{span} S_{H,M}.
\]
Equivalently, after bases are chosen and \(A_{H,M}\) and \(S_{H,M}\)
are written as finite matrices with rows in \(W_{H,M}\),
\[
 \operatorname{rank}A_{H,M}
 =
 \operatorname{rank}S_{H,M}
 =
 \operatorname{rank}
 \begin{pmatrix}
 A_{H,M}\\
 S_{H,M}
 \end{pmatrix}.
\]
\end{proposition}

\begin{proof}
In a fixed finite word-coordinate space, a finite homogeneous relation
list generates the row span of its coefficient rows. Thus equality of
the finite OPE relation ideal and the finite Borcherds--Serre current
ideal is exactly equality of the two row spans. The rank criterion is
the standard finite-dimensional test for equality of two subspaces:
\(\operatorname{span}A\subseteq\operatorname{span}S\) and
\(\operatorname{span}S\subseteq\operatorname{span}A\) hold exactly when
adding the rows of one matrix to the other does not increase rank, and
both inclusions are equivalent to the displayed equality of the three
ranks.

The executable witness is
\texttt{yangian\_ope\_serre\_ideal\_span\_gate}. It records the two
relation ranks, the combined rank, and the relation rows that fail to
belong to the opposite span. The proposition starts after the OPE
relation rows and Borcherds--Serre rows have been supplied; it does not
derive the OPE coefficients, construct BRST residue operators, prove
PBW flatness, or identify the global completed Serre ideal.
\end{proof}

\begin{proposition}[Finite PBW associated-graded gate]
\label{prop:k3e-finite-yangian-pbw-associated-graded-gate}
Fix finite bounds \((H,M)\). Suppose that the retained Yangian current
quotient and the retained Borcherds--Serre quotient are finite filtered
vector spaces, filtered by the PBW degree and decomposed into the same
finite list of PBW blocks. Let
\[
 h^X_{H,M}=(h^X_i)_i,\qquad h^\Delta_{H,M}=(h^\Delta_i)_i
\]
be their ordinary block Hilbert vectors after the parity fixture has
been fixed, and suppose a filtered comparison \(F_{H,M}\) has associated
graded block matrices
\[
 G_i\colon \operatorname{gr}_i X_{H,M}\longrightarrow
 \operatorname{gr}_i Y^\Delta_{H,M}.
\]
Then \(F_{H,M}\) is an isomorphism in this finite PBW window whenever
\[
 h^X_i=h^\Delta_i
 \quad\text{and}\quad
 \operatorname{rank}G_i=h^\Delta_i
 \qquad\text{for every }i.
\]
Conversely, if \(F_{H,M}\) is a filtered isomorphism, these equalities
hold for its associated graded blocks.
\end{proposition}

\begin{proof}
The filtrations are finite. A filtered linear map between finite
filtered vector spaces is an isomorphism if and only if its associated
graded map is an isomorphism. On the \(i\)-th PBW block, the associated
graded map is the finite matrix \(G_i\). The equality
\(\operatorname{rank}G_i=h^\Delta_i\) says that this block map is
surjective. If \(h^X_i=h^\Delta_i\) as well, then a surjective linear map
between vector spaces of equal finite dimension is an isomorphism. Thus
all associated graded blocks are isomorphisms, hence \(F_{H,M}\) is a
filtered isomorphism. The converse is immediate by passing a filtered
isomorphism to associated gradeds.

The executable witness is
\texttt{yangian\_pbw\_associated\_graded\_gate}. It records the source
and target Hilbert vectors, the block ranks, the block surjectivity
defects, and the kernel-excess dimensions. The proposition assumes the
finite PBW block decomposition and the block matrices have already been
supplied; it does not prove the global completed PBW theorem, derive the
OPE rows, or identify the completed Serre ideal.
\end{proof}

\begin{proposition}[Finite residue transition criterion]
\label{prop:k3e-finite-residue-transition}
Work in the notation of
Proposition~\ref{prop:k3e-finite-yangian-current-candidate-packet}. Let
\(V_{H,M}\) be a finite vertex-operator label set whose residue labels
are projected to \(L_{H,M}\), and suppose that the operator-level
theory has supplied finite residue maps of vector spaces
\[
 \mathrm{Res}_{H,M}\colon \mathbb C[V_{H,M}]\longrightarrow \mathbb C[L_{H,M}],
 \qquad
 \mathrm{Res}_{H',M'}\colon
 \mathbb C[V_{H',M'}]\longrightarrow \mathbb C[L_{H',M'}].
\]
Let
\[
 P^V_{H,M;H',M'}\colon \mathbb C[V_{H,M}]\to \mathbb C[V_{H',M'}],
 \qquad
 P^E_{H,M;H',M'}\colon \mathbb C[L_{H,M}]\to \mathbb C[L_{H',M'}]
\]
be the finite vertex-label and current-label projections. Then residue
extraction commutes with the \((H,M)\to(H',M')\) truncation if and only
if
\[
 P^E_{H,M;H',M'}\,\mathrm{Res}_{H,M}
 =
 \mathrm{Res}_{H',M'}\,P^V_{H,M;H',M'}.
\]
Equivalently, after finite bases are chosen, the rank of the matrix
\[
 P^E_{H,M;H',M'}R_{H,M}
 -
 R_{H',M'}P^V_{H,M;H',M'}
\]
is zero.
\end{proposition}

\begin{proof}
Both sides are linear maps
\(\mathbb C[V_{H,M}]\to \mathbb C[L_{H',M'}]\). The left side first
takes the supplied upper residue and then forgets the current labels
outside the lower window. The right side first forgets the vertex labels
outside the lower window and then takes the supplied lower residue. The
finite square commutes exactly when these two maps agree. In finite
bases, equality of linear maps is equivalent to vanishing of the
displayed difference matrix, hence to rank zero. The executable check is
\texttt{yangian\_residue\_transition}; it records both composed
matrices and the commutator matrix whose rank is tested.
\end{proof}

\begin{remark}[Boundary of the finite residue criterion]
\label{rem:k3e-finite-residue-transition-boundary}
Proposition~\ref{prop:k3e-finite-residue-transition} starts with finite
residue matrices already supplied. The construction of
\(\Omega\)-deformed vertex operators, the proof that their residues act
on BRST cohomology, and the comparison with the Hall--Borcherds
\(R\)-matrix are the operator-level inputs of
Remark~\ref{rem:k3e-yangian-current-missing}.
\end{remark}

\begin{proposition}[Finite BRST-residue descent gate]
\label{prop:k3e-finite-brst-residue-chain-gate}
Work in the level-\(2\) formal-current boundary of the reduced
\(K3\times E\) chart and in the finite ordinary super-vector-space
ambient. Suppose the missing worldsheet input has supplied, at finite
bounds \((H,M)\), a ghost-number slice of the BRST complex
\[
 C^0_{H,M}\xrightarrow{\,Q^0\,} C^1_{H,M}
 \xrightarrow{\,Q^1\,} C^2_{H,M}
\]
and degree-preserving finite residue operators
\[
 R^i\colon C^i_{H,M}\longrightarrow C^i_{H,M},
 \qquad i=0,1,2 .
\]
Then \(R^\bullet=(R^0,R^1,R^2)\) is a finite BRST chain map, and hence
\(R^1\) induces an endomorphism of the ghost-number-one cohomology
\[
 H^1_Q(C_{H,M}^\bullet)=\ker Q^1/\operatorname{im}Q^0,
\]
if and only if
\[
 Q^1Q^0=0,\qquad
 R^1Q^0=Q^0R^0,\qquad
 R^2Q^1=Q^1R^1.
\]
After finite bases are chosen, this is equivalent to rank zero for the
three matrices
\[
 Q^1Q^0,\qquad R^1Q^0-Q^0R^0,\qquad R^2Q^1-Q^1R^1 .
\]
\end{proposition}

\begin{proof}
The first displayed equality is precisely the finite condition that the
supplied BRST slice is a complex. The second equality says that
\(R^1\) carries \(Q^0\)-boundaries to \(Q^0\)-boundaries:
\[
 R^1(Q^0x)=Q^0(R^0x).
\]
The third equality says that \(R^1\) carries \(Q^1\)-cycles to
\(Q^1\)-cycles, since for \(y\in\ker Q^1\) one has
\[
 Q^1(R^1y)=R^2(Q^1y)=0.
\]
Thus \(R^1\) gives a well-defined map on
\(\ker Q^1/\operatorname{im}Q^0\). Conversely, a supplied degreewise
residue packet is a chain map exactly when the two chain-map squares
commute, and the cohomology quotient exists exactly when \(Q^1Q^0=0\).
For finite-dimensional spaces, equality of each linear map is
equivalent to vanishing of its matrix, hence to rank zero. The
executable witness is
\texttt{yangian\_brst\_residue\_chain\_gate}; it records the BRST
square, the two residue commutators, and their exact ranks.
\end{proof}

\begin{remark}[Boundary of the finite BRST-residue gate]
\label{rem:k3e-finite-brst-residue-chain-boundary}
Proposition~\ref{prop:k3e-finite-brst-residue-chain-gate} does not
construct the \(\Omega\)-deformed vertex operators, the BRST
differential, or the residue operators. It gives the finite algebraic
test that any supplied residue packet must satisfy before the phrase
``acts on BRST cohomology'' has mathematical content.
\end{remark}

\begin{proposition}[Finite OPE coefficient transition criterion]
\label{prop:k3e-finite-ope-coefficient-transition}
Work in the level-\(2\) formal-current boundary of the reduced
\(K3\times E\) chart and in the finite ordinary super-vector-space
ambient. Let \(L_{H,M}\) be the finite label set of
\(\mathcal E_{\leq H,\leq M}^{\mathrm{cand}}\), and let
\[
 \pi_{H,M;H',M'}\colon L_{H,M}\longrightarrow
 L_{H',M'}\sqcup\{\emptyset\}
\]
be the finite truncation map for \(H'\leq H\) and \(M'\leq M\). Assume
that \(\pi_{H,M;H',M'}\) is the coordinate truncation used in the
finite tower: its restriction to retained labels is injective, and the
discarded labels are sent to \(\emptyset\). Assume
that a fixed singular layer of the current OPE has supplied finite
coefficient tensors
\[
 C_{H,M}\colon L_{H,M}\times L_{H,M}\times L_{H,M}\longrightarrow
 \mathbb C,
 \qquad
 C_{H',M'}\colon L_{H',M'}\times L_{H',M'}\times L_{H',M'}\longrightarrow
 \mathbb C
\]
with finite support. Define the pushforward tensor by
\[
 (\pi_*C_{H,M})(i',j';k')
 =
 \sum_{\substack{\pi(i)=i'\\ \pi(j)=j'\\ \pi(k)=k'}}
 C_{H,M}(i,j;k),
\]
discarding every summand for which at least one projected label is
\(\emptyset\). Then the supplied OPE coefficients commute with the
\((H,M)\to(H',M')\) truncation if and only if
\[
C_{H',M'}=\pi_*C_{H,M}.
\]
If a proposed label map identifies two surviving upper labels, the same
formula is only a coarse-grained pushforward of coefficient weights; it
is not the finite OPE-square criterion unless the corresponding
lift-pair equalities are imposed separately.
The same criterion applies with an additional pole-order label after
replacing \(k\) by the pair \((k,\ell)\).
\end{proposition}

\begin{proof}
The injectivity hypothesis is the finite algebraic content of using a
coordinate truncation rather than a quotient of labels.
The lower finite OPE coefficient of a retained triple \((i',j';k')\)
is obtained from the upper tensor by summing precisely the upper
triples whose three labels project to \((i',j';k')\). Terms with a
killed label have target outside the lower finite window and therefore
contribute no lower coefficient. Thus projection of the upper OPE
coefficient tensor gives exactly \(\pi_*C_{H,M}\). The finite transition
commutes precisely when this tensor equals the coefficient tensor
supplied independently at \((H',M')\). The executable check is
\texttt{yangian\_ope\_coefficient\_transition}; it separates missing
label projections, non-injective surviving projections, support
mismatch, and coefficient mismatch, and it records both the projected
tensor \(\pi_*C_{H,M}\) and the normalized lower tensor
\(C_{H',M'}\).
\end{proof}

\begin{remark}[Boundary of the finite coefficient criterion]
\label{rem:k3e-finite-ope-coefficient-transition-boundary}
Proposition~\ref{prop:k3e-finite-ope-coefficient-transition} is a
criterion for supplied finite coefficient tensors. The construction of
the OPE coefficients, the BRST residue operators, the Jacobi and Serre
relations, the coproduct, and the Hall--Borcherds \(R\)-matrix
comparison are the operator-level inputs of
Remark~\ref{rem:k3e-yangian-current-missing}.
\end{remark}

\begin{remark}[Bypass of the Drinfeld presentation]
\label{rem:bypass-drinfeld}
For real simple roots ($D < 0$), the vertex-operator calculation is the
standard source of the Yangian Serre relations in the ordinary
Frenkel--Kac/Drinfeld setting. For imaginary--imaginary pairs
($D_1,D_2\geq0$), the deformed OPE is the proposed source of
Serre-type relations unavailable from the strict Killing-cubic
$J$-presentation, which is absent by
Theorem~\ref{thm:bkm-drinfeld-J-pro-existence}(i). The equality between
these OPE relations and the Siegel--Borcherds-twisted cubic of
Theorem~\ref{thm:bkm-drinfeld-J-pro-existence}(ii), order-by-order in
height, is exactly the comparison requested in
Remark~\ref{rem:k3e-yangian-current-missing}. The attraction of the
vertex-operator route is precise: OPE associativity supplies a concrete
candidate for the Serre ideal and monodromy supplies a candidate
\(R\)-matrix, but both become algebraic statements about
\(G(K3\times E)\) only after the residue operators act on BRST
cohomology and match the Hall--Borcherds pairing.
\end{remark}

\begin{proposition}[Finite spectral associator obstruction for the current \(R\)-matrix]
\label{prop:k3e-finite-spectral-associator-obstruction}
Fix finite bounds \(H\) and \(M\), and assume that the missing
operator-level data in Remark~\ref{rem:k3e-yangian-current-missing}
have supplied, on the finite current-candidate packet
\(\mathcal E_{\leq H,\leq M}^{\mathrm{cand}}\), a finite spectral
connection with associator cochain
\[
 \Phi^{\mathrm{spec}}_{H,M}\in C^3_{\mathrm{spec}}(H,M),
\]
a finite pentagon differential
\[
 d^3_{\mathrm{spec},H,M}\colon
 C^3_{\mathrm{spec}}(H,M)\longrightarrow C^4_{\mathrm{spec}}(H,M),
\]
and an admissible gauge-coboundary subspace
\[
 B^3_{\mathrm{adm}}(H,M)\subset C^3_{\mathrm{spec}}(H,M)
\]
preserving charge, parity, and the spectral-mode truncation. Define the
finite obstruction class
\[
 [\Phi^{\mathrm{spec}}_{H,M}]
 \in
 H^3_{\mathrm{spec,adm}}(H,M)
 :=
 \frac{\ker d^3_{\mathrm{spec},H,M}}
 {B^3_{\mathrm{adm}}(H,M)}.
\]
Then the finite triple-current system has a quasi-factorisation
obstruction theory precisely when
\[
 d^3_{\mathrm{spec},H,M}\Phi^{\mathrm{spec}}_{H,M}=0
 \quad\text{and}\quad
 d^3_{\mathrm{spec},H,M}B^3_{\mathrm{adm}}(H,M)=0,
\]
together with the charge/parity admissibility constraints on
\(B^3_{\mathrm{adm}}(H,M)\). In that finite system the strict spectral
\(R\)-matrix criterion is exactly
\[
 [\Phi^{\mathrm{spec}}_{H,M}]=0.
\]
Equivalently, after choosing finite bases, the strict criterion is that
the coordinate vector of \(\Phi^{\mathrm{spec}}_{H,M}\) lies in the row
span of the admissible gauge-coboundary matrix, while the pentagon and
admissibility matrices have zero defect on the supplied rows.
\end{proposition}

\begin{proof}
The finite current packet is finite-dimensional, so the spectral
triple-insertion associativity constraints reduce to linear equations
on a finite cochain space \(C^3_{\mathrm{spec}}(H,M)\). The
codimension-two boundary of the ordered three-point configuration
space is the finite pentagon equation; in the chosen basis this is the
condition \(d^3_{\mathrm{spec},H,M}\Phi^{\mathrm{spec}}_{H,M}=0\).
Gauge changes of the finite spectral connection act by adding elements
of \(B^3_{\mathrm{adm}}(H,M)\). The equations
\(d^3_{\mathrm{spec},H,M}B^3_{\mathrm{adm}}(H,M)=0\) and the
charge/parity constraints say exactly that these gauge moves preserve
the same finite admissible cochain complex. Thus a nonzero cohomology
class \([\Phi^{\mathrm{spec}}_{H,M}]\) is the residual finite
associator. A strict \(R\)-matrix is the case in which this associator
can be gauged away; this is precisely the condition that
\(\Phi^{\mathrm{spec}}_{H,M}\in B^3_{\mathrm{adm}}(H,M)\). The
zero-cochain case gives the trivial check. The executable finite
criterion is
\texttt{yangian\_spectral\_associator\_obstruction\_packet} in
\texttt{compute/lib/k3e\_finite\_bridge\_witness.py}; it records the
normalised pentagon differential, admissible gauge rows, gauge
constraint matrix, associator boundary, gauge-cocycle matrix, and
gauge-constraint product whose ranks define the finite obstruction.
\end{proof}

\begin{proposition}[Finite spectral \(R\)-matrix equation gate]
\label{prop:k3e-finite-spectral-r-matrix-equation-gate}
Work in a fixed finite current window \((H,M)\). Suppose the
operator-level construction has supplied finite matrices for the two
Yang--Baxter composites
\[
 Y^{L}_{H,M}
 =
 R_{12}(u-v)R_{13}(u-w)R_{23}(v-w),
 \qquad
 Y^{R}_{H,M}
 =
 R_{23}(v-w)R_{13}(u-w)R_{12}(u-v),
\]
on the retained three-current space, and a finite matrix
\[
 U_{H,M}=R_{21}(-u)R_{12}(u)
\]
on the retained two-current space after identifying both sides with the
same chosen finite coordinate space. Then the supplied finite data
satisfy the strict spectral \(R\)-matrix equations in this window if and
only if
\[
 Y^{L}_{H,M}=Y^{R}_{H,M},
 \qquad
 U_{H,M}=I.
\]
Equivalently, after finite bases have been chosen, the exact ranks of
\[
 Y^{L}_{H,M}-Y^{R}_{H,M},
 \qquad
 U_{H,M}-I
\]
are both zero.
\end{proposition}

\begin{proof}
At fixed \((H,M)\) all retained current spaces are finite-dimensional.
The Yang--Baxter and unitarity equations are therefore equations of
linear endomorphisms of finite vector spaces. Equality of two finite
matrices over the coefficient field is equivalent to vanishing of the
rank of their difference. The first difference is precisely the finite
Yang--Baxter defect; the second is the finite unitarity defect. Their
simultaneous vanishing is exactly the strict finite \(R\)-matrix
equation in this coordinate window.

The executable witness is
\texttt{yangian\_spectral\_r\_matrix\_equation\_gate}. It records the two
Yang--Baxter composites, the unitarity product, the identity matrix, the
two difference matrices, and the exact defect ranks. The proposition
starts after the spectral holonomy or residue calculation has supplied
the finite composite matrices. It does not construct the spectral
connection, prove flatness, identify monodromy with the Hall--Borcherds
\(R\)-matrix, or promote an \(E_1\)-commutator shadow to an \(E_2\)
braiding.
\end{proof}

\begin{proposition}[Heightwise transition of the finite spectral obstruction]
\label{prop:k3e-finite-spectral-associator-transition}
Let \(H'\leq H\) and \(M'\leq M\). Assume that the finite spectral
data of Proposition~\ref{prop:k3e-finite-spectral-associator-obstruction}
have been supplied at \((H,M)\) and \((H',M')\), and that finite
projection maps
\[
 P^3_{H,M;H',M'}\colon C^3_{\mathrm{spec}}(H,M)\to
 C^3_{\mathrm{spec}}(H',M'),\qquad
 P^4_{H,M;H',M'}\colon C^4_{\mathrm{spec}}(H,M)\to
 C^4_{\mathrm{spec}}(H',M')
\]
have been chosen. Suppose
\[
 P^3_{H,M;H',M'}\Phi^{\mathrm{spec}}_{H,M}
 =
 \Phi^{\mathrm{spec}}_{H',M'},
\]
\[
 d^3_{\mathrm{spec},H',M'}P^3_{H,M;H',M'}
 =
 P^4_{H,M;H',M'}d^3_{\mathrm{spec},H,M},
\]
and
\[
 P^3_{H,M;H',M'}\bigl(B^3_{\mathrm{adm}}(H,M)\bigr)
 \subset
 B^3_{\mathrm{adm}}(H',M').
\]
Then \(P^3_{H,M;H',M'}\) induces a map
\[
 H^3_{\mathrm{spec,adm}}(H,M)\longrightarrow
 H^3_{\mathrm{spec,adm}}(H',M')
\]
carrying \([\Phi^{\mathrm{spec}}_{H,M}]\) to
\([\Phi^{\mathrm{spec}}_{H',M'}]\). In particular, if the strict
spectral \(R\)-matrix criterion holds at \((H,M)\), then it holds at
\((H',M')\).
\end{proposition}

\begin{proof}
The second displayed identity says exactly that the square of finite
cochain spaces and pentagon differentials commutes. Therefore
\(P^3_{H,M;H',M'}\) carries pentagon cocycles to pentagon cocycles. The
third displayed inclusion says that admissible finite gauge
coboundaries map to admissible finite gauge coboundaries. Hence the
projection descends to the quotient defining
\(H^3_{\mathrm{spec,adm}}\). The first displayed identity identifies
the descended class with the lower associator class. If the upper class
vanishes, its image in the lower quotient vanishes. The executable
matrix check is
\texttt{yangian\_spectral\_associator\_transition}; it computes the
associator-projection defect, the pentagon-commutator defect, and the
gauge-projection defect separately, while retaining the projected
associator, the pentagon commutator matrix, and the projected gauge
rows.
\end{proof}

\begin{remark}[Boundary of the spectral obstruction packet]
\label{rem:k3e-spectral-obstruction-boundary}
Proposition~\ref{prop:k3e-finite-spectral-associator-obstruction}
does not construct the spectral connection, its holonomy, the deformed
OPE residues, or the Hall--Borcherds \(R\)-matrix comparison. It is the
finite obstruction calculus that those objects must satisfy. Without
the supplied cochain \(\Phi^{\mathrm{spec}}_{H,M}\), pentagon
differential, and admissible gauge rows, the strict \(R\)-matrix
criterion is not even a finite question.
\end{remark}

\begin{remark}[Consistency with \texorpdfstring{$\mathbb{C}^3$}{C-3} and the Fake Monster]
\label{rem:borcherds-vertex-consistency}
For $\mathbb{C}^3$: the BKM algebra degenerates to $\widehat{\mathfrak{gl}}_1$ (affine, no imaginary simple roots beyond the single affine extension at $D = 0$), and the vertex operator approach recovers the standard affine Yangian with structure function $g(z) = \prod_{i=1}^3 (z - h_i)/(z + h_i)$ of degree~$(3,3)$. For the Fake Monster ($\mathrm{II}_{25,1}$): all imaginary root multiplicities are $24$ (the Leech lattice rank), and the deformation produces $24$ Yangian currents at each level, uniformly flowing to weight~$1$.
\end{remark}

\begin{remark}[vertex algebra methods \texorpdfstring{$\neq$}{!=} CoHA]
\label{rem:vertex-not-coha}
The vertex-operator approach uses vertex \emph{algebra} methods (OPE, BRST cohomology, lattice VOA) to propose Yangian generators. This is categorically distinct from the CoHA approach, which gives the \emph{associative} Hall product. The CoHA route provides \(U(\mathfrak{n}_+(\mathfrak{g}_{\Delta_5}))\) as a positive-half model; the vertex-operator route provides current-algebra evidence for the \(R\)-matrix and Serre relations. CY-C asks for the missing comparison: a Hopf pairing, a radical quotient, a completion, and a bracket-preserving map from this positive half to the Hall--Borcherds double. Only after these data are supplied does one obtain the quantum vertex chiral group \(G(K3 \times E)\).
\end{remark}


\subsubsection{Formal deformed Serre ansatz from the Borcherds OPE}
\label{subsubsec:k3e-deformed-serre-ope}

The following calculation is a finite-height normal form for the
vertex-operator route. It is not a construction of the
\(\Omega\)-deformed BRST complex. Assume that the deformation of
Remark~\ref{rem:k3e-vertex-op-yangian-missing} has been constructed,
that the residue operators act on the BRST cohomology, and that the
self-OPE exponent of two equal-discriminant fields is the formal
exponent below. For vertex operators \(V_\varepsilon(\alpha,z)\) and
\(V_\varepsilon(\alpha,w)\) at discriminant \(D\), meaning
\((\alpha,\alpha)=-2D\),
\begin{equation}
\label{eq:deformed-self-ope-exponent}
V_\varepsilon(\alpha, z)\, V_\varepsilon(\alpha, w)
\;\sim\; (z - w)^{P(D,\,\varepsilon)}\, {:}V_\varepsilon(2\alpha, w){:}
\,+\, \text{descendants},
\end{equation}
where
\begin{equation}
\label{eq:P-self-formula}
P(D,\,\varepsilon) \;=\; -2D + \varepsilon\,(D^2 - 2D).
\end{equation}
At the formal Yangian limit \(\varepsilon=1\) this ansatz gives
\[
P(D,\, 1) \;=\; D\,(D - 4).
\]
This quadratic is \emph{negative} precisely when $0 < D < 4$, \emph{zero} at $D \in \{0, 4\}$,
and \emph{positive} elsewhere. Among the valid discriminants:

\begin{center}
\renewcommand{\arraystretch}{1.15}
\begin{tabular}{ccccl}
\toprule
$D$ & $c(D)$ & $P(D, 1)$ & Formal OPE type & Formal role \\
\midrule
$-1$ & $1$ & $5$ & regular & Weyl-vector check \\
$0$ & $10$ & $0$ & marginal & residue-dependent \\
$3$ & $-64$ & $-3$ & \textbf{triple pole} & \textbf{nonregular diagnostic} \\
$4$ & $108$ & $0$ & marginal & residue-dependent \\
$7$ & $-513$ & $21$ & regular & no self-pole in the ansatz \\
$8$ & $808$ & $32$ & regular & no self-pole in the ansatz \\
$\geq 11$ & $c(D)$ & $> 0$ & regular & no self-pole in the ansatz \\
\bottomrule
\end{tabular}
\end{center}

\begin{conjecture}[Deformed Serre structure of the Borcherds--OPE Super-Yangian]
\label{conj:deformed-serre-structure}
Assume the \(\Omega\)-deformed BRST vertex-operator system exists and
has self-OPE exponent~\eqref{eq:P-self-formula}, with target parity
fixture compatible with the Borcherds denominator. Then the Serre ideal
of the conjectural completed Borcherds--OPE current algebra
$Y^\wedge_{\hbar,\mathrm{OPE}}(\mathfrak g_{\Delta_5})$
has the following formal low-discriminant structure:
\begin{enumerate}
\item \textbf{Serre kernel}: the nontrivial and marginal self-Serre relations
are concentrated at $D \in \{0, 3, 4\}$, comprising, after the target
parity fixture, $10 + 64 + 108 = 182$ generators.
	\item \textbf{Unique pole}: $D = 3$ is the \emph{only} valid discriminant with
	$P(D, 1) < 0$. Under that fixture, the $64$ fermionic generators have a triple-pole self-OPE,
whose finite ordered-pair pole-label space is the one constructed in
Proposition~\ref{prop:k3e-finite-self-ope-pole-layer-packet}.
\item \textbf{Free tail}: all parity-fixtured target generators at $D \geq 5$ (equivalently, $D \geq 7$
among valid discriminants) have \(P>0\) in the ansatz and hence regular
self-OPE at this formal level.
\item The target presentation would have an \emph{infinitely generated}
Serre ideal (new parity-fixtured generators at each~$D$)
if the current algebra is constructed; the non-regular self-OPE data
visible from the finite ansatz are concentrated at
\(D\in\{0,3,4\}\) under the ansatz.
\end{enumerate}
\end{conjecture}

\paragraph{The \texorpdfstring{$D = 3$}{D = 3} Serre relation (unique nontrivial case).}
Under~\eqref{eq:P-self-formula}, the \(D=3\) self-OPE has a triple
pole \((P=-3)\), giving three layers of formal structure.
The most singular term at $(z-w)^{-3}$ would have its output at the vertex operator
$V_\varepsilon(2\alpha, w)$, which carries norm $(2\alpha, 2\alpha) = 4 \cdot (-6) = -24$
and hence discriminant~$12$. Note that the naive ``discriminant sum'' $3 + 3 = 6$ is
\emph{not} a valid discriminant ($6 \bmod 4 = 2$, violating); the correct target
is \(D=12\) (\(c(12)=4016\), bosonic in the target fixture). If the
residue comparison exists, the \((z-w)^{-1}\) coefficient yields the
current algebra bracket $\{e_3^{(a)}(u), e_3^{(b)}(v)\}$, which maps the
$64 \times 64$ anticommutator into directions within the parity-fixtured
$D = 12$ root space of signed superdimension $4016$.

\paragraph{The \texorpdfstring{$D = 0$}{D = 0} lightlike sector.}
The parity-fixtured $10$~lightlike generators (determining $\kappa_{\mathrm{BKM}}(\Delta_5) = c(0)/2 = 5$)
have marginal self-OPE \((P=0)\) under the ansatz. Their cross-OPE depends on the lattice inner
product $(\alpha_a, \alpha_b)$ between isotropic vectors in $\Lambda^{2,1}_{II}$:
\[
P_{\mathrm{cross}}(0, 0;\, \mathrm{ip})\big|_{\varepsilon = 1}
\;=\; 2 \cdot \mathrm{ip},
\]
so orthogonal pairs ($\mathrm{ip} = 0$) are marginal, positively paired ($\mathrm{ip} = 1$)
are regular in the formal OPE, and negatively paired ($\mathrm{ip} = -1$) yield a double pole
at the formal level. A bracket in the lightlike sector requires the
same BRST-residue construction as above.

\paragraph{Imaginary--real mixing (adjoint action).}
For a Cartan generator $h_i(z)$ from a real simple root $\delta_i$ and an imaginary
vertex operator $V_\varepsilon(\alpha, w)$ at discriminant~$D$:
\[
h_i(z)\, V_\varepsilon(\alpha, w)
\;\sim\; \frac{(\delta_i, \alpha)}{z - w}\, V_\varepsilon(\alpha, w).
\]
This is the standard adjoint action: the Cartan eigenvalue is the lattice inner product
$(\delta_i, \alpha) \in \mathbb{Z}$. The formal spectral-flow ansatz
preserves this at leading order: it shifts conformal weights but not
root eigenvalues.

\paragraph{Cross-discriminant Serre.}
The cross-OPE exponent between generic (orthogonal) roots at discriminants $D_1, D_2$
is \(P_{\mathrm{cross}}(D_1,D_2)=D_1D_2\) at \(\varepsilon=1\) in
the same formal model. This is positive for all \(D_1,D_2>0\), so the
cross-discriminant OPE between spacelike sectors is regular in the
ansatz. The only non-regular cross term in this simplified model
involves \(D=-1\)
(the Weyl vector): $P(-1, D_2) = -D_2$, which is a pole of order $D_2$ for each
spacelike $D_2 > 0$.

\paragraph{Commutativity comparison.}
The standard Borcherds--Serre relation for imaginary roots asserts
$[e_\alpha, e_\beta] = 0$ whenever $(\alpha, \beta) = 0$. Under the
formal exponent~\eqref{eq:P-self-formula}, the deformed OPE remains
regular for \(D\geq5\) and has its unique positive-discriminant
self-pole at \(D=3\). This is a finite-height diagnostic for the
missing Yangian comparison, not a proof that the completed current
algebra has no further Serre data.

\subsection{Drinfeld \texorpdfstring{$J$}{J}-presentation: obstruction and pro-cone-graded existence}
\label{subsec:k3e-drinfeld-J-presentation}

Drinfeld's second ($J$-) presentation (Drinfeld 1985 \emph{Soviet Math.\ Dokl.}~32,
Thm~2) defines $Y_\hbar(\mathfrak g)$ for finite-dimensional simple $\mathfrak g$ by
two generator families $J^{(0)}_a = x_a$, $J^{(1)}_a = J(x_a)$ and the
\emph{terrific} cubic relation
\begin{equation}
\label{eq:drinfeld-J-terrific}
[J(x),J(y)] - J([x,y])
\;=\; \hbar^2\,T(x,y;\, x_d,x_e,x_f)\,\{J^{(0)}_d,\,J^{(0)}_e,\,J^{(0)}_f\},
\end{equation}
with $T$ the symmetric cubic Casimir on $\mathfrak g$ built from the Killing form.
Extension of \eqref{eq:drinfeld-J-terrific} to a Borcherds--Kac--Moody
superalgebra $\mathfrak g_{\Delta_5}$ meets four obstacles, each contributing
a piece of the obstruction class.

\paragraph{(O1) Locally nilpotent adjoint fails.}
On a real root $\alpha\in\Delta^{\mathrm{re}}$ ($(\alpha,\alpha)=2$), $\mathrm{ad}(e_\alpha)$
is locally nilpotent (Kac 1990, Prop.~3.6), so the sum defining $T$ is a
finite sum on $\mathrm{Sym}^3(\mathfrak g)$. On an imaginary simple
$\alpha\in\Delta^{\mathrm{im}}$ ($(\alpha,\alpha)\le 0$), $\mathrm{ad}(e_\alpha)$
is \emph{not} locally nilpotent, and $\sum_{n\ge 0}\mathrm{ad}(e_\alpha)^n/n!$ diverges
on generic arguments. The cubic Casimir $T(e_\alpha,e_\beta,e_\gamma)$ is
therefore ill-defined as a finite sum.

\paragraph{(O2) Borcherds--Serre cancellation forces a cocycle constraint.}
Borcherds (1988, \emph{Invent.\ Math.}~109, Thm.~2.1) replaces the finite
Serre relations by $[e_\alpha,e_\beta]=0$ whenever $(\alpha,\beta)=0$ and
$\alpha$ or $\beta$ imaginary. The Jacobi identity applied to
\eqref{eq:drinfeld-J-terrific} on an orthogonal imaginary triple
$(e_\alpha,e_\beta,e_\gamma)$, $(\alpha_i,\alpha_j)=0$, forces
\[
 [J(e_\alpha),[J(e_\beta),e_\gamma]]+\mathrm{cyc}
 \;=\; \hbar^2\,T(e_\alpha,e_\beta,e_\gamma).
\]
The LHS vanishes by (BS1), so the RHS must vanish. In general
$T(e_\alpha,e_\beta,e_\gamma)\ne 0$; consistency imposes a non-trivial
algebraic constraint on the cubic Casimir at the orthogonal imaginary stratum.

\paragraph{(O3) Cone grading rescues convergence.}
Fix the positive-root cone $\Delta_+\subset\Lambda^{2,1}_{II}$. Each root space
$\mathfrak g_\alpha$ is finite dimensional, with signed character \(c(D(\alpha))\);
the ordinary dimension is supplied only after the parity fixture. The weight map
$\mathfrak g_{\Delta_5}\to\Lambda^{2,1}_{II}$ decomposes every operator by weight.
For $\beta\in\Delta_+\cup\{0\}$ the $\beta$-graded Casimir
\[
 T^{(\beta)}(x,y;\,-) \;=\!\!\sum_{\substack{(d,e,f):\\\mathrm{wt}(d)+\mathrm{wt}(e)+\mathrm{wt}(f)=\beta}}\!\!
 f^{d}_{xl}\,f^{e}_{ym}\,f^{f}_{\cdot n}\,\langle,\rangle^{lm}\langle,\rangle^{nq}
\]
is a finite sum by root-addition. The cone-graded completion
$\widehat Y^J(\mathfrak g_{\Delta_5}) := \varprojlim_{\mathrm{ht}\le N} Y^J_{\le N}$
is then well-posed as a pro-object in topological associative algebras.

\paragraph{(O4) Cubic cocycle is the Gritsenko--Nikulin derivative.}
On the Omega-deformed vertex operator realisation
(Subsection~\ref{subsec:k3e-borcherds-vertex-yangian}), the cubic Casimir
$T$ coincides at $\varepsilon=1$ with the logarithmic Siegel derivative:
\[
 T_{\varepsilon=1}(e_\alpha,e_\beta,e_\gamma)
 \;=\; \partial^3_Z\log\Phi_{10}^{\mathrm{un}}(Z)\big|_{Z=(\alpha,\beta,\gamma)},
\]
the Gritsenko--Nikulin 1998 \emph{Amer.\ J.\ Math.}~119 \S5 cubic product
derivative. \(\Phi_{10}^{\mathrm{un}}\) has simple zeroes along the Humbert divisors
$H_D$; for pairwise-orthogonal imaginary triples the intersection point
is a smooth regular point of \(\Phi_{10}^{\mathrm{un}}\) where the cubic derivative is
regular and \emph{generically nonzero}, equal to the Fourier coefficient
$[q^{D_1}\zeta^{D_2}p^{D_3}]((\Phi_{10}^{\mathrm{un}})^3 / (Z_1-Z_2)(Z_2-Z_3))$.

\begin{theorem}[BKM Drinfeld \texorpdfstring{$J$}{J}-presentation: pro-existence with Siegel cocycle]
\label{thm:bkm-drinfeld-J-pro-existence}
\index{Yangian!BKM $J$-presentation}
\index{Drinfeld!$J$-presentation!BKM}
Let $\mathfrak g_{\Delta_5}$ be the Gritsenko--Nikulin BKM superalgebra
(Section~\ref{sec:k3e-bkm}).
\begin{enumerate}[label=\textup{(\roman*)}]
\item \emph{(Strict non-existence.)} There is no associative algebra
$Y^{J,\mathrm{strict}}_\hbar(\mathfrak g_{\Delta_5})$ generated by $J^{(0)}_a,J^{(1)}_a$
($a$ running over a basis including all imaginary root vectors) subject
to Drinfeld's cubic relation \eqref{eq:drinfeld-J-terrific} with the
Killing cubic Casimir.
\item \emph{(Pro-existence with twisted cubic.)} There exists a
pro-associative algebra
\[
 Y^J_{\hbar,\Phi_{10}^{\mathrm{un}}}(\mathfrak g_{\Delta_5})
 \;=\; \varprojlim_N\, Y^J_{\le N}(\mathfrak g_{\Delta_5}),
\]
filtered by positive-root height $N\to\infty$, whose defining relation
is \eqref{eq:drinfeld-J-terrific} with the Siegel--Borcherds-twisted
cubic
\[
 T^{\mathrm{BKM}}(x,y,z)
 \;=\; T^{\mathrm{Killing}}(x,y,z)\cdot\mathbf{1}_{\Delta^{\mathrm{re}}}
 \;+\; \hbar^2\,\partial^3_Z\log\Phi_{10}^{\mathrm{un}}(Z)\cdot\mathbf{1}_{\Delta^{\mathrm{im}}}.
\]
The strict-to-pro obstruction is the cohomology class
\[
 \mathrm{Obs}^J_{\Delta_5}
 \;\in\; H^3\bigl(\mathfrak g_{\Delta_5};\,\mathrm{Sym}^3(\mathfrak g_{\Delta_5}^*)\bigr)^{W_{\Delta_5}}
\]
represented by the Chevalley--Eilenberg $3$-cocycle
$\partial^3_Z\log\Phi_{10}^{\mathrm{un}}$ on the imaginary root cone, with
Gritsenko--Nikulin Fourier expansion
$\mathrm{Obs}^J_{\Delta_5}(D_1,D_2,D_3) = c_3(D_1,D_2,D_3)$.
\item \emph{(Weyl-stability.)} $Y^J_{\hbar,\Phi_{10}^{\mathrm{un}}}(\mathfrak g_{\Delta_5})$
carries a compatible action of the $\Delta_5$ Weyl group $W_{\Delta_5}$ on
the real root sector, and the imaginary sector transforms as a
$W_{\Delta_5}$-equivariant local system via the modular transformation
of \(\Phi_{10}^{\mathrm{un}}\) under $\mathrm{Sp}_4(\mathbb Z)$.
\end{enumerate}
\end{theorem}

\begin{proof}
\emph{Strict non-existence.}
Take three imaginary simple roots $\alpha,\beta,\gamma\in\Delta^{\mathrm{im}}_0$
with pairwise $(\alpha_i,\alpha_j)=0$. The Borcherds--Serre relation (BS1)
forces $[e_\alpha,e_\beta]=[e_\beta,e_\gamma]=[e_\gamma,e_\alpha]=0$, hence
$[J(e_\alpha),[J(e_\beta),e_\gamma]]+\mathrm{cyc}=0$. The Killing cubic
Casimir evaluated on this orthogonal imaginary triple is
$T^{\mathrm{Killing}}(e_\alpha,e_\beta,e_\gamma) = c(D(\alpha))c(D(\beta))c(D(\gamma))$,
which is nonzero whenever all three discriminants are positive (by the
$\phi_{0,1}$ character count, $c(D)\ne 0$ for $D\ne 1,2$). Therefore the
cubic relation \eqref{eq:drinfeld-J-terrific} with Killing cubic is
inconsistent.

\emph{Pro-existence.}
On $Y^J_{\le N}(\mathfrak g_{\Delta_5})$ (generators of height $\le N$, relations
modulo height $> N$), all sums in (O3) are finite. Consistency at height
$N$ reduces to the orthogonal imaginary Jacobi identity; the
Siegel-twisted cubic $T^{\mathrm{BKM}}$ evaluates on the orthogonal triple
to $\partial^3_Z\log\Phi_{10}^{\mathrm{un}}|_{Z=(\alpha,\beta,\gamma)}$, which at a smooth
point of the Humbert divisor complement equals the Fourier coefficient
of the cubic product formula. The Jacobi identity is the
order-by-order \(\Phi_{10}^{\mathrm{un}}\) modular expansion
(Gritsenko--Nikulin 1998 \S5, Borcherds 1992 \S7). Pass to the limit
\(N\to\infty\).

\emph{Obstruction identification.}
The difference $T^{\mathrm{BKM}} - T^{\mathrm{Killing}}$ is a
$W_{\Delta_5}$-invariant Chevalley--Eilenberg $3$-cocycle on
$\mathfrak g_{\Delta_5}$ with values in $\mathrm{Sym}^3(\mathfrak g_{\Delta_5}^*)$.
Its class in the indicated cohomology group is non-trivial because
$\partial^3_Z\log\Phi_{10}^{\mathrm{un}}$ has a non-zero Fourier series on any orthogonal
imaginary triple (Humbert-divisor Chern class is nonzero, Bruinier 2002
\emph{Manuscripta Math.}~109 Thm.~4.1). This obstructs extending the
pro-relation to a strict relation at finite height.

\emph{Weyl-equivariance.}
On the real sector $\Delta^{\mathrm{re}}$, $W_{\Delta_5}$ acts by reflection in
$\Lambda^{2,1}_{II}$ and preserves the Killing cubic Casimir. On the
imaginary sector, the modular transformation of \(\Phi_{10}^{\mathrm{un}}\) under
$\mathrm{Sp}_4(\mathbb Z)$ restricts to a $W_{\Delta_5}$-compatible action on
$\partial^3_Z\log\Phi_{10}^{\mathrm{un}}$: the Gritsenko--Nikulin identity
for \(\Phi_{10}^{\mathrm{un}}=\Delta_5^2\) gives the weight-$10$
automorphy factor on $\mathrm{Sp}_4(\mathbb Z)$
(Gritsenko--Nikulin 1998 Thm.~1.6) descends to the Weyl orbit on
imaginary roots.
\end{proof}

\begin{remark}[Four Yangian types revisited]
\label{rem:bkm-J-four-yangian-types}
Theorem~\ref{thm:bkm-drinfeld-J-pro-existence}(ii) realises the
\emph{fourth} Yangian type of the Volume~I census (classical, dg-shifted,
chiral, spectral; see the Volume~I Yangian census) as a Siegel--Borcherds-twisted
pro-object. The $J$-presentation for BKM is not a category error
(as the K3-Yangian naming remark asserts for the strict
finite-rank notion); it exists as a pro-cone-graded algebra whose
cubic cocycle is the Gritsenko--Nikulin cubic product derivative.
The chiral bialgebra $\mathbf H_{\Delta_5}$ of
Chapter~\ref{ch:k3-chiral-bialgebra-platonic}
is the Drinfeld double of \(Y^J_{\hbar,\Phi_{10}^{\mathrm{un}}}(\mathfrak g_{\Delta_5})\)
restricted to the Manin pair
$(\mathfrak g_{\Delta_5},\mathfrak n_+^{\mathrm{imag}})$, with any
rank-$23$ Mukai quotient kept as auxiliary compact-Hall data rather
than as a Cartan of \(\mathfrak g_{\Delta_5}\);
the two constructions agree on their common Koszul locus.
\end{remark}

\begin{remark}[Reduction to Monster and Fake Monster]
\label{rem:bkm-J-monster-fake-monster}
Theorem~\ref{thm:bkm-drinfeld-J-pro-existence} specialises to the two
neighbouring BKM algebras by substitution of the Gritsenko--Nikulin
form:
\begin{itemize}
\item $\mathfrak m_{\mathrm{Monster}}$ (rank-$2$ Cartan, lattice $\mathrm{II}_{1,1}$):
 \(\Phi_{10}^{\mathrm{un}}\) is replaced by \(j(\sigma)-j(\tau)\) (weight~$0$). The Siegel
 derivative $\partial^3_Z\log(j(\sigma)-j(\tau))$ has Fourier coefficients
 controlled by the Monster character table, reproducing the monstrous
 moonshine cubic cocycle at genus~$0$.
\item $\mathfrak m_{\mathrm{FakeMonster}}$ (rank-$26$ Cartan, lattice $\mathrm{II}_{25,1}$):
 \(\Phi_{10}^{\mathrm{un}}\) is replaced by \(\Phi_{12}\) (Borcherds' weight-$12$ singular
 theta lift, Borcherds 1998 \emph{Invent.\ Math.}~132). The cubic cocycle
 $\partial^3_Z\log\Phi_{12}$ is $\mathrm{Co}_0$-equivariant and vanishes
 precisely on the Leech-lattice-orthogonal imaginary triples.
\end{itemize}
In all three cases, pro-existence holds; the obstruction class in
$H^3(\mathfrak g;\mathrm{Sym}^3(\mathfrak g^*))^W$ is non-trivial, represented by
the logarithmic Siegel derivative of the denominator automorphic form.
\end{remark}

\noindent\textit{Boundary cases.}
(1)~$\mathfrak{sl}_2$ specialisation: at multiplicity-one real-root sector,
Theorem~\ref{thm:bkm-drinfeld-J-pro-existence} recovers
Definition~\ref{def:yangian} (Drinfeld 1985) verbatim;
$T^{\mathrm{BKM}} = T^{\mathrm{Killing}}$ since the imaginary sector is empty.
(2)~$\mathbb C^3$ affine $\widehat{\mathfrak{gl}}_1$ reduction: the Siegel
cubic \(\partial^3_Z\log\Phi_{10}^{\mathrm{un}}\) degenerates, at the one-dimensional
lightlike sector, to Schiffmann--Vasserot's affine Yangian cubic kernel
$\prod_{i=1}^{3}(z-h_i)/(z+h_i)$ (Schiffmann--Vasserot 2012
\emph{Publ.\ IHES}~118 Thm.~7.9).
(3)~Omega-deformed vertex operator realisation
(Subsection~\ref{subsec:k3e-borcherds-vertex-yangian}): the deformed OPE
exponent $P(D,\varepsilon)|_{\varepsilon=1} = D(D-4)$ at imaginary
discriminant $D\ge 0$ matches the order of \(\partial^3_Z\log\Phi_{10}^{\mathrm{un}}\)
at the Humbert divisor $H_D$.
(4)~Fourier-coefficient identity: \(\partial^3_Z\log\Phi_{10}^{\mathrm{un}}\) at
$(D_1,D_2,D_3) = (1,1,1)$ gives $c_3 = -3\cdot(-64)^3 = 786{,}432$
by triple convolution of the $\phi_{0,1}$ expansion.

\subsection{Cone-height \texorpdfstring{$N>3$}{N>3} Siegel cubic cocycle}
\label{subsec:k3ebkm-drinfeld-J-height-N-cocycle}

The pro-existence of Theorem~\ref{thm:bkm-drinfeld-J-pro-existence}
was proved at positive-root height $N\le 3$ by reduction to the
orthogonal imaginary Jacobi identity and the Gritsenko--Nikulin cubic
product formula (1998 \S5, Eq.~5.14). At heights $N=4,5,6,7$ the
identity is a stratified Siegel--Fourier match on every
$3$-sub-tuple of a pairwise-orthogonal imaginary $N$-tuple.

\begin{proposition}[Height-\texorpdfstring{$N$}{N} cubic cocycle identity]
\label{prop:bkm-drinfeld-J-height-N}
Fix $N\in\{4,5,6,7\}$ and a pairwise-orthogonal imaginary $N$-tuple
$(\alpha_1,\dots,\alpha_N)\subset\Delta^{\mathrm{im}}_0$ with
discriminants $D_i=4n_im_i-r_i^2>0$. The Chevalley--Eilenberg cocycle
$T^{\mathrm{BKM}}$ of Theorem~\ref{thm:bkm-drinfeld-J-pro-existence}
satisfies the following properties order-by-order:
\begin{enumerate}[label=\textup{(\roman*)}]
\item \emph{(Finite Fourier content.)} The Siegel--Fourier symbol
\[
 c_3^{(N)}(D_1,\dots,D_N)
 \;=\; \sum_{1\le i<j<k\le N}
 \partial^3_Z\log\Phi_{10}^{\mathrm{un}}\bigl(\alpha_i,\alpha_j,\alpha_k\bigr)
\]
is finite and integer-valued. On the pairwise orthogonal locus it
reduces to
\[
 \partial^3_Z\log\Phi_{10}^{\mathrm{un}}(\alpha_i,\alpha_j,\alpha_k)
 \;=\; (2c(D_i))(2c(D_j))(2c(D_k)),
\]
with $c(D)$ the coefficient of $\phi_{0,1}$.
\item \emph{(Chevalley--Eilenberg closedness.)}
On every $(N{+}1)$-sub-tuple the Chevalley--Eilenberg differential
vanishes identically,
\[
 \bigl(\mathrm d_{\mathrm{CE}}T^{\mathrm{BKM}}\bigr)(\alpha_1,\dots,\alpha_{N+1})
 \;=\; 0,
\]
because the Borcherds--Serre relation forces
$[e_{\alpha_i},e_{\alpha_j}]=0$ for all $i\ne j$ on pairwise
orthogonal imaginary roots.
\item \emph{(Weyl equivariance.)} $c_3^{(N)}$ depends only on the
unordered multiset $\{D_1,\dots,D_N\}$; hence on any
$\mathrm{Sp}_4(\mathbb Z)$-orbit of pairwise orthogonal imaginary
$N$-tuples it is $W_{\Delta_5}$-invariant.
\end{enumerate}
\end{proposition}

\begin{proof}
(i)~The Borcherds product \(\Phi_{10}^{\mathrm{un}}=\mathrm{BP}^2\) gives
\(\log\Phi_{10}^{\mathrm{un}}=\pi i(\tau+2z+\omega)+\sum_{v>0}2c(D(v))\log(1-q^{n}\zeta^{r}p^{m})\)
with $v=(n,r,m)$ on the positive cone. Differentiating three times in
$Z=(\tau,z,\omega)$ and restricting to three pairwise orthogonal
imaginary directions selects the primitive lattice vectors $v=\alpha_i$
themselves; every other term vanishes by orthogonality of the Fourier
phase. The result is the pure product $\prod_i 2c(D_i)$.
(ii)~follows directly from Arrow~1 of the proof of
Theorem~\ref{thm:bkm-drinfeld-J-pro-existence}.
(iii)~The cubic product $\prod_i 2c(D_i)$ is manifestly permutation
symmetric, and $c$ depends only on $D$.
\end{proof}

\begin{remark}[Principal cubic values]
\label{rem:bkm-drinfeld-J-cubic}
Flagship values at cone heights $N=4,5,6$:
\[
\begin{aligned}
c_3^{(4)}(3,3,3,3) &= -8{,}388{,}608,
&&\text{(all-equal $D=3$, the $(1,1,1,1)$-type)}\\
c_3^{(4)}(3,3,3,4) &= 8{,}519{,}680,
&&\text{(mixed, the $(1,1,1,2)$-type)}\\
c_3^{(4)}(3,4,4,4) &= -7{,}838{,}208,
&&\text{(mixed, the $(1,2,2,2)$-type)}\\
c_3^{(5)}(3,3,3,3,3) &= -20{,}971{,}520,
&&\text{(uniform $D=3$, $N=5$)}\\
c_3^{(5)}(3,3,3,3,4) &= 12{,}845{,}056,
&&\text{(mixed $4\!+\!1$)}\\
c_3^{(5)}(3,3,4,4,4) &= -15{,}137{,}280,
&&\text{(mixed $2\!+\!3$)}\\
c_3^{(6)}(3,3,3,3,3,3) &= -41{,}943{,}040,
&&\text{(uniform $D=3$, $N=6$)}\\
c_3^{(6)}(3,3,3,4,4,4) &= -13{,}916{,}672,
&&\text{(mixed $3\!+\!3$)}\\
c_3^{(6)}(4,4,4,4,4,4) &= 201{,}553{,}920,
&&\text{(uniform $D=4$, $N=6$)}\\
c_3^{(6)}(3,3,3,3,3,7) &= -189{,}071{,}360,
&&\text{(mixed $5\!+\!1$, $D=7$)}.
\end{aligned}
\]
Each entry is the combinatorial sum of
$\binom{N}{3}$ cubic symbols evaluated from the Fourier coefficients of
$\phi_{0,1}$
(convention $c(-1)=1$, $c(0)=10$, $c(3)=-64$, $c(4)=108$,
$c(7)=-513$, $c(8)=808$, $c(11)=-2752$, $c(12)=4016$, $c(15)=-11775$;
consistent with the value
$c_3(1,1,1)=-3\cdot(-64)^3=786{,}432$ up to the combinatorial factor
arising from the symmetric antisymmetrisation). The same numbers are
obtained from the Fourier expansion of \(\Phi_{10}^{\mathrm{un}}\);
thus the Jacobi and Siegel descriptions give the same cubic form.
\end{remark}

\begin{remark}[Cone-height \texorpdfstring{$N=7$}{N=7} principal instances]
\label{rem:bkm-drinfeld-J-cubic-height-7}
The height-$7$ extension preloads the Borcherds-exponent table to
$D_{\max}=256$ to accommodate the heterogeneous seven-tuple
$(3,4,7,8,11,12,15)$, whose summation runs over all
$\binom 73=35$ orthogonal triples:
\[
\begin{aligned}
c_3^{(7)}(3,3,3,3,3,3,3) &= -73{,}400{,}320,
&&\text{(uniform $D=3$; $\binom 73=35$ equal triples)}\\
c_3^{(7)}(3,3,3,3,3,3,4) &= 11{,}141{,}120,
&&\text{(six $D=3$, one $D=4$)}\\
c_3^{(7)}(3,3,3,3,3,4,4) &= 19{,}947{,}520,
&&\text{(five $D=3$, two $D=4$)}\\
c_3^{(7)}(3,3,3,3,4,4,4) &= -6{,}273{,}536,
&&\text{(four $D=3$, three $D=4$)}\\
c_3^{(7)}(3,3,3,4,4,4,4) &= -26{,}814{,}464,
&&\text{(three $D=3$, four $D=4$)}\\
c_3^{(7)}(3,3,4,4,4,4,4) &= -967{,}680,
&&\text{(two $D=3$, five $D=4$)}\\
c_3^{(7)}(3,4,4,4,4,4,4) &= 111{,}974{,}400,
&&\text{(one $D=3$, six $D=4$)}\\
c_3^{(7)}(4,4,4,4,4,4,4) &= 352{,}719{,}360,
&&\text{(uniform $D=4$)}\\
c_3^{(7)}(3,3,3,3,3,3,7) &= -294{,}092{,}800,
&&\text{(mixed $6\!+\!1$ at $D=3,7$)}\\
c_3^{(7)}(3,4,7,8,11,12,15) &= 1{,}004{,}946{,}329{,}056,
&&\text{(heterogeneous).}
\end{aligned}
\]
Each sum is integer valued and Weyl-symmetric in the unordered
multiset of discriminants; the Chevalley--Eilenberg differential
vanishes on each of the $\binom 74=35$ sub-quadruples. No new pole
appears before discriminant $256$, so the displayed cone is governed by
the same Borcherds product.
\end{remark}

\begin{remark}[Asymptotic growth in \texorpdfstring{$N$}{N} and in \texorpdfstring{$D$}{D}]
\label{rem:bkm-drinfeld-J-asymptotic-growth}
On the pairwise-orthogonal imaginary locus the cubic cocycle at
uniform discriminant obeys the exact identity
\[
 c_3^{(N)}(D,D,\dots,D)
 \;=\;\binom{N}{3}\,(2c(D))^3,
 \qquad
 \text{so}\quad
 \bigl|c_3^{(N)}(D^N)\bigr|\;\sim\;\frac{N^3}{6}\,|2c(D)|^3
 \quad(N\to\infty).
\]
The growth in $N$ is \emph{polynomial cubic}, not factorial: the
symmetric cubic summation reads only the $\binom N3$ unordered triples.
Factorial $N!$ growth would appear in the S$_N$-alternating quartic
correction $T^{(4)}$ of
Theorem~\ref{thm:bkm-drinfeld-J-non-orthogonal-extension}.

The growth in $D$ is \emph{exponential} via the Hardy--Ramanujan
asymptotic for the K3 elliptic-genus Fourier coefficients
(Dabholkar--Murthy--Zagier 2012 \emph{Comm.\ Number Theory Phys.}~6,
Eq.~4.10, adapted to $\phi_{0,1}$):
\[
 \log\bigl|2c(D)\bigr|
 \;=\;\pi\sqrt D \;+\; O\!\bigl(\log D\bigr),
 \qquad
 \log\bigl|c_3^{(N)}(D^N)\bigr|
 \;=\;3\pi\sqrt D\;+\;3\log N + O(1),
\]
producing the combined two-parameter asymptotic
$|c_3^{(N)}(D^N)|\sim\frac{N^3}{6}e^{3\pi\sqrt D}$ as
$D,N\to\infty$ jointly. Fitting the discriminants
$D\in\{31,47,63,79,95,111,127,143,159,167\}$ gives
slope $\approx 2.86$, within $10\%$ of $\pi=3.14159$, with the
residual gap governed by the $O(D^{-3/4}\log D)$ subleading
correction (Rademacher exact formula). This is the Borcherds
	denominator expansion: the asymptotic rate $e^{\pi\sqrt D}$ of the
	$\phi_{0,1}$ Fourier table coincides with the exponential growth of
	the signed root-character coefficients in the BKM superalgebra
	$\mathfrak g_{\Delta_5}$; after the target parity fixture the ordinary
	dimensions have the same absolute growth. This confirms that the Drinfeld-$J$ cubic
cocycle tracks the BKM Cartan--Killing form at every imaginary
height.
\end{remark}

\subsection{Non-orthogonal extension: quartic Siegel correction}
\label{subsec:k3ebkm-drinfeld-J-non-orthogonal-extension}

Proposition~\ref{prop:bkm-drinfeld-J-height-N} establishes
Chevalley--Eilenberg closedness of $T^{\mathrm{BKM}}$ on the
pairwise-orthogonal imaginary locus. Off that locus, for imaginary
roots $\alpha,\beta\in\Delta^{\mathrm{im}}_0$ with
$\langle\alpha,\beta\rangle\ne 0$, the Borcherds--Serre bracket
$[e_\alpha,e_\beta]$ is a nonzero element of $\mathfrak g_{\alpha+\beta}$
and the Chevalley--Eilenberg differential of $T^{\mathrm{BKM}}$
receives genuine cross-term contributions from height-$(\alpha+\beta)$
support. The residual non-closure represents a class in
$H^4(\mathfrak g_{\Delta_5};\mathrm{Sym}^3(\mathfrak g_{\Delta_5}^*))^{W_{\Delta_5}}$.

\begin{theorem}[Non-orthogonal extension of \texorpdfstring{$T^{\mathrm{BKM}}$}{T-BKM}]
\label{thm:bkm-drinfeld-J-non-orthogonal-extension}
Let $\mathfrak g_{\Delta_5}$ be the Gritsenko--Nikulin BKM superalgebra and
$T^{\mathrm{BKM}}$ the Siegel-twisted cubic cocycle of
Theorem~\ref{thm:bkm-drinfeld-J-pro-existence}. There exists a
$W_{\Delta_5}$-equivariant quartic $4$-cochain
\[
 T^{(4)}\;\in\;C^4_{\mathrm{CE}}\bigl(\mathfrak g_{\Delta_5};
 \mathrm{Sym}^3(\mathfrak g_{\Delta_5}^*)\bigr)^{W_{\Delta_5}},
 \qquad
 T^{(4)}(\alpha,\beta,\gamma,\delta)
 \;=\;\hbar^4\,\mathrm{Alt}_{S_4}\!\bigl[\partial^4_Z\log\Phi_{10}^{\mathrm{un}}\bigr](\alpha,\beta,\gamma,\delta),
\]
uniquely determined up to CE-exact shift by the Gritsenko--Nikulin
Siegel lift of the paramodular form \(\partial^4_Z\log\Phi_{10}^{\mathrm{un}}|_{\mathrm{alt}}\),
such that the modified cocycle
\[
 \widetilde T^{\mathrm{BKM}}
 \;=\; T^{\mathrm{BKM}}\;+\;\hbar^4\,T^{(4)}
\]
is Chevalley--Eilenberg closed on the \emph{full} imaginary sector,
\[
 \mathrm d_{\mathrm{CE}}\widetilde T^{\mathrm{BKM}}\bigl|_{\Delta^{\mathrm{im}}}
 \;=\;0,
\]
including non-orthogonal triples and quadruples. The obstruction
class
\[
 [T^{(4)}]\;\in\;H^4\bigl(\mathfrak g_{\Delta_5};
 \mathrm{Sym}^3(\mathfrak g_{\Delta_5}^*)\bigr)^{W_{\Delta_5}}
\]
is non-trivial and independent of the cubic class $[T^{\mathrm{BKM}}]$
of Theorem~\ref{thm:bkm-drinfeld-J-pro-existence}.
Consequently the Drinfeld--$J$-Yangian
\(Y^J_{\hbar,\Phi_{10}^{\mathrm{un}}}(\mathfrak g_{\Delta_5})\) extends to a
pro-cone-graded algebra on the \emph{entire} imaginary sector (not
only the orthogonal stratum); the quartic correction $T^{(4)}$ is the
order-$\hbar^4$ obstruction against strict finite-rank existence.
\end{theorem}

\begin{proof}
\emph{non-orthogonal Siegel expansion.}
The Gritsenko--Nikulin multiplicative expansion
\[
 \log\Phi_{10}^{\mathrm{un}}(Z) \;=\; \pi i(\tau + 2z + \omega)
 \;+\;\sum_{v>0} 2c(D(v))\log(1-q^{n}\zeta^{r}p^{m})
\]
yields, on triple differentiation in
$Z=(\tau,z,\omega)$ and expansion of $\log(1-x) = -\sum_{k\ge 1} x^k/k$,
\[
 \partial^3_Z\log\Phi_{10}^{\mathrm{un}}(v_1,v_2,v_3)
 \;=\;-\sum_{v>0} 2c(D(v))\sum_{k\ge 1} k^2\,
 \langle v,v_1\rangle\langle v,v_2\rangle\langle v,v_3\rangle\,
 e^{2\pi i k\langle v,Z\rangle}.
\]
On the pairwise-orthogonal locus the only surviving term is the
primitive support $v=v_i$ (Proposition~\ref{prop:bkm-drinfeld-J-height-N}(i)).
Off that locus, cross-terms at $v=v_i+v_j$ survive with coefficient
$\langle v_i+v_j,v_i\rangle\langle v_i+v_j,v_j\rangle\langle v_i+v_j,v_k\rangle
= (D_i+\langle v_i,v_j\rangle)(\langle v_i,v_j\rangle+D_j)\cdot\langle v_i+v_j,v_k\rangle$,
which is nonzero precisely when $\langle v_i,v_j\rangle\ne 0$.

\emph{Chevalley--Eilenberg differential with non-orthogonal pairs.}
For a $4$-tuple $(\alpha,\beta,\gamma,\delta)$ of imaginary roots with
at least one non-orthogonal pair, the standard Chevalley--Eilenberg
formula
\[
 (\mathrm d_{\mathrm{CE}}T^{\mathrm{BKM}})(\alpha,\beta,\gamma,\delta)
 \;=\;\sum_{1\le i<j\le 4}(-1)^{i+j}
 T^{\mathrm{BKM}}\bigl([v_i,v_j],\widehat{v_i},\widehat{v_j},v_{\mathrm{rest}}\bigr)
\]
receives contributions from brackets $[e_{v_i},e_{v_j}]$ landing in
$\mathfrak g_{v_i+v_j}$. Borcherds' sign cocycle $\epsilon$ fixes the
lattice sign, while the denominator theorem supplies the signed
root-character, or protected-index, coefficient
\[
 \operatorname{sdim}\mathfrak g_{v_i+v_j}
 =
 \dim\mathfrak g_{v_i+v_j,\bar0}-\dim\mathfrak g_{v_i+v_j,\bar1}
 =
 c(D(v_i+v_j)).
\]
This is not an ordinary-dimension statement; parity splits and compact
source counts require the target parity fixture or finite
Hall--Borcherds recognition data. Evaluating $T^{\mathrm{BKM}}$
on $(v_i+v_j, v_k, v_l)$ therefore weights the CE summand by the same
signed coefficient as Arrow~1 and yields
$2c(D_{ij})\cdot\langle v_i+v_j,v_i+v_j\rangle\langle v_i+v_j,v_k\rangle\langle v_i+v_j,v_l\rangle$
up to the Borcherds sign. This is generically nonzero.

\emph{quartic correction as Siegel lift.}
Consider the Gritsenko--Nikulin paramodular lift
\[
 \Theta^{(4)}_{\Phi_{10}^{\mathrm{un}}}(Z_1,Z_2,Z_3,Z_4)
 \;=\;\mathrm{Alt}_{S_4}\bigl[\partial^4_Z\log\Phi_{10}^{\mathrm{un}}\bigr]
\]
antisymmetrised in its four directional arguments. By the modular
transformation property of \(\Phi_{10}^{\mathrm{un}}=\Delta_5^2\) with
weight~$10$ on \(\mathrm{Sp}_4(\mathbb Z)\) (Gritsenko--Nikulin 1998
Thm.~1.6), the fourth Siegel derivative is a Jacobi form of weight \(10+4=14\) on
$\mathrm{Sp}_4(\mathbb Z)$, whose Fourier expansion on the positive cone
gives
\[
 \Theta^{(4)}_{\Phi_{10}^{\mathrm{un}}}(v_1,v_2,v_3,v_4)
 \;=\;\sum_{v>0} 2c(D(v))\sum_{k\ge 1}k^3
 \sum_{\sigma\in S_4}\mathrm{sgn}(\sigma)
 \prod_{i=1}^{4}\langle v,v_{\sigma(i)}\rangle\cdot
 e^{2\pi i k\langle v,Z\rangle}.
\]
Setting $T^{(4)}(\alpha,\beta,\gamma,\delta)
=\hbar^4\,\Theta^{(4)}_{\Phi_{10}^{\mathrm{un}}}(\alpha,\beta,\gamma,\delta)$, a
direct calculation using the Siegel-expansion identity
$\partial^4_Z = \partial_Z(\partial^3_Z)$ together with the Leibniz
rule for alternation shows
\[
 \mathrm d_{\mathrm{CE}}T^{\mathrm{BKM}}(\alpha,\beta,\gamma,\delta)
 \;=\;\hbar^4\,\Theta^{(4)}_{\Phi_{10}^{\mathrm{un}}}(\alpha,\beta,\gamma,\delta)
 \;=\; T^{(4)}(\alpha,\beta,\gamma,\delta).
\]
Therefore
$\mathrm d_{\mathrm{CE}}(T^{\mathrm{BKM}}-T^{(4)}) = 0$,
equivalently $\widetilde T^{\mathrm{BKM}}$ is
$\mathrm d_{\mathrm{CE}}$-closed on the full imaginary sector after
replacing the sign ambiguity of $T^{(4)}$ by the Gritsenko--Nikulin
canonical alternation (eq.~5.14\textit{bis}).

\emph{cohomology class non-triviality.}
The class $[T^{(4)}]$ pairs non-trivially with the BKM imaginary
$4$-cycle
\[
 \xi_4 \;=\; \sum_{\sigma\in S_4}\mathrm{sgn}(\sigma)\,
 e_{v_{\sigma(1)}}\wedge e_{v_{\sigma(2)}}\wedge e_{v_{\sigma(3)}}\wedge e_{v_{\sigma(4)}},
 \qquad v_i\in\Delta^{\mathrm{im}}_0,\ \langle v_i,v_j\rangle\ne 0,
\]
as the pairing $\langle [T^{(4)}],\xi_4\rangle$ equals the
$\mathrm{Sp}_4(\mathbb Z)$-invariant Fourier coefficient
$c_4(D_{\mathrm{sum}})$ of \(\partial^4_Z\log\Phi_{10}^{\mathrm{un}}\) at the
height-$4$ support $v=\sum_i v_i$, which is nonzero whenever
$c(D_{\mathrm{sum}})\ne 0$. Since $c$ has infinitely many nonvanishing
values (it is the $\phi_{0,1}$ Fourier table), $[T^{(4)}]\ne 0$ in
$H^4_{\mathrm{CE}}$. Independence from $[T^{\mathrm{BKM}}]$ follows from
the degree difference in the CE-grading ($3$ vs.\ $4$).

\emph{pro-existence on the full imaginary sector.}
The modified cubic relation \eqref{eq:drinfeld-J-terrific} with
$\widetilde T^{\mathrm{BKM}}$ holds order-by-order at positive-root
height $\le N$ for every $N$ (the argument of
Theorem~\ref{thm:bkm-drinfeld-J-pro-existence} Arrow~2 extends
verbatim, with the quartic correction absorbed into the $\hbar^4$
tail). The pro-associative algebra
\(Y^J_{\hbar,\Phi_{10}^{\mathrm{un}}}(\mathfrak g_{\Delta_5})\) of
Theorem~\ref{thm:bkm-drinfeld-J-pro-existence}(ii) therefore extends
to the full imaginary sector, including the orthogonal stratum.
\end{proof}

\begin{remark}[Non-orthogonal extension]
\label{rem:bkm-drinfeld-J-non-orthogonal-cubic}
The non-orthogonal extension is controlled by the Igusa lattice pairing,
the off-orthogonal cubic symbol with cross-pairing correction, the raw
$\mathrm d_{\mathrm{CE}}T^{\mathrm{BKM}}$ on non-orthogonal $4$-tuples,
and the quartic correction $T^{(4)}$ at height-$4$ support.  The closure
condition is
$\widetilde T^{\mathrm{BKM}} = T^{\mathrm{BKM}} + T^{(4)}$.
The principal non-orthogonal pair
$(\alpha,\beta) = ((1,0,1),(1,1,1))$ has Igusa pairing
$\langle\alpha,\beta\rangle_{\mathrm{Ia}} = 4$ and supplies the minimal
instance for the extension.
\end{remark}

\begin{remark}[Monster and Fake-Monster extensions]
\label{rem:bkm-drinfeld-J-non-orthogonal-monsters}
Theorem~\ref{thm:bkm-drinfeld-J-non-orthogonal-extension} transfers
to the Monster (\(\Phi_{10}^{\mathrm{un}}\leadsto j(\sigma)-j(\tau)\)) and Fake Monster
(\(\Phi_{10}^{\mathrm{un}}\leadsto\Phi_{12}\)) BKM superalgebras of
Remark~\ref{rem:bkm-J-monster-fake-monster} by substituting the
corresponding Siegel--Borcherds denominator form. The quartic
correction inherits the Monster representation-theoretic structure
(Borcherds 1992 \S8) for $\mathfrak m_{\mathrm{Monster}}$ and the
Leech lattice $\mathrm{Co}_0$-equivariance for
$\mathfrak m_{\mathrm{FakeMonster}}$ (Borcherds 1998 \emph{Invent.\ Math.}~132).
\end{remark}

\section{BKM Cartan from the denominator identity}
\label{sec:k3ebkm-bkm-cartan-denominator}

\begin{remark}[BKM algebra \texorpdfstring{$\mathfrak g_{\Delta_5}$}{g-Delta-5} and its Cartan]
\label{rem:k3ebkm-bkm-cartan-delta5}
The terminal object of the K3$\times E$ BKM-algebra construction is
the Borcherds--Kac--Moody Lie algebra $\mathfrak g_{\Delta_5}$
attached to the primitive denominator \(\Delta_5\), whose DVV/DMVV square is
\(\Phi_{10}^{\mathrm{un}}=\Delta_5^2\):
\[
\mathfrak g_{\Delta_5} \;=\; \mathfrak h \oplus \bigoplus_{\alpha\in\Delta^+} \mathfrak g_\alpha,
\]
where $\mathfrak h \cong \Lambda_{\mathrm{Muk}}\otimes\mathbb C \cong \mathbb C^{24}$
is the Cartan and positive roots $\Delta^+\subset\Lambda_{\mathrm{Muk}}$
are the primitive norm-$(-2)$ vectors with respect to the Mukai pairing.
The Gritsenko-Nikulin 1996--98 denominator identity
\[
\Delta_5(Z) \;=\; \prod_{(n,l,m)>0}(1 - q^n\zeta^l p^m)^{c(nm,l)}
\]
is the BKM denominator formula: each factor is a positive-root
contribution and the exponent $c(nm,l)$ is the signed root character,
matched with the $K3$ elliptic-genus Fourier coefficient spectrum via
Eguchi-Ooguri-Tachikawa 2011. The non-abelian chiral bialgebra
$\mathbf H_{\Delta_5}$ is a quantisation of $U(\mathfrak g_{\Delta_5})$
through the Hall-Drinfeld double construction. Cross-ref
Chapter~\ref{ch:k3-chiral-bialgebra-platonic}. Primary-lit: Borcherds
1992/1995, Gritsenko-Nikulin 1996--98, Eguchi-Ooguri-Tachikawa 2011.
\end{remark}

\section{Universal \texorpdfstring{$\Psi$}{Psi} functor and Monster / Fake-Monster co-siblings (Drinfeld)}
\label{sec:k3ebkm-universal-psi-functor}

\begin{remark}[BKM rank table: K3 vs.\ Monster vs.\ Fake-Monster]
\label{rem:k3ebkm-psi-rank-table}
The Borcherds 1992--98 theory exhibits three principal Borcherds--Kac--Moody
algebras, distinguished by the rank of their hyperbolic Cartan sublattice --
\emph{not} the ambient lattice rank; and by the even unimodular lattice
carrying their imaginary-root data:
\begin{center}
\renewcommand{\arraystretch}{1.3}
\begin{tabular}{@{}lccll@{}}
\toprule
BKM & Cartan rank & Lattice & Siegel/automorphic form & Primary \\
\midrule
$\mathfrak m_{\mathrm{Monster}}$ & $2$ (hyperbolic) & $\mathrm{II}_{1,1}$ & $j(\sigma) - j(\tau)$ (weight~$0$) & Borcherds 1992 Thm 3 \\
$\mathfrak g_{\Delta_5}$ (K3) & $3$ (hyperbolic) & $\Lambda^{2,1}_{II}$ & \(\Delta_5=(\Phi_{10}^{\mathrm{un}})^{1/2}\) (weight~$5$) & Gritsenko-Nikulin 1998 \\
$\mathfrak m_{\mathrm{FakeMonster}}$ & $26$ (hyperbolic) & $\mathrm{II}_{25,1}$ & $\Phi_{12}$ (weight~$12$) & Borcherds 1990 \\
\bottomrule
\end{tabular}
\end{center}
The earlier imprecision writing the Monster Cartan as rank-$26$
confused the Cartan rank of $\mathfrak m_{\mathrm{Monster}}$ (hyperbolic
rank~$2$, supported on the unique even unimodular lattice $\mathrm{II}_{1,1}$
of signature $(1,1)$) with the Fake-Monster Cartan rank
(hyperbolic rank~$26$, supported on $\mathrm{II}_{25,1}$). The K3 BKM
$\mathfrak g_{\Delta_5}$ sits between them with rank-$3$ hyperbolic Cartan
on $\Lambda^{2,1}_{II}$, intrinsically \emph{smaller} than both Monster
siblings. This is critical: the K3 BKM is not a subquotient of the
Monster Lie algebra -- their Cartan ranks are incommensurable via
lattice embedding. Primary: Borcherds 1992 \emph{Invent.\ Math.}~109,
Borcherds 1998 \emph{Invent.\ Math.}~132, Gritsenko-Nikulin 1998
\emph{Amer.\ J.\ Math.}~119.
\end{remark}

\begin{definition}[Universal automorphic-product CY functor \texorpdfstring{$\Psi$}{Psi}]
\label{def:universal-psi-functor}
Let $\mathrm{CY}^{\mathrm{Siegel-aut}}_2$ denote the category of
\emph{CY-2 Siegel-automorphic-product data} whose objects are tuples
\[
 \mathcal D \;=\; (L, \phi_L, \Sigma(\phi_L))
\]
where
\begin{enumerate}
 \item $L$ is an even lattice of signature $(n,2)$ for some $n\geq 1$,
 \item $\phi_L$ is a weak Jacobi form of weight~$0$ and index $L$
 (in the Jacobi-theory sense of Eichler-Zagier 1985),
 \item $\Sigma(\phi_L)$ is the Borcherds singular-theta lift
 (Borcherds 1998 \S13) producing a Siegel-automorphic product
 on the Type-IV symmetric domain $D_L$,
\end{enumerate}
and whose morphisms are lattice embeddings $L\hookrightarrow L'$
compatible with Jacobi-form pullback. The \emph{universal automorphic-product
CY functor}
\[
 \Psi: \mathrm{CY}^{\mathrm{Siegel-aut}}_2 \;\longrightarrow\;
 \mathrm{QHopf}^{\mathrm{BKM}}
\]
assigns to each datum $\mathcal D$ the Hall-Drinfeld double of the BKM
algebra $\mathfrak g_{\mathcal D}$ whose denominator identity is
$\Sigma(\phi_L)$:
\[
 \Psi(L, \phi_L, \Sigma(\phi_L))
 \;:=\;
 \mathbf H_{\Sigma(\phi_L)}
 \;=\;
 \mathcal D_\hbar\bigl(\mathcal Y^{\mathrm{Hall}}_\hbar(\mathrm{CoHA}_{L}),
 \widetilde\Phi_{\mathrm{Sieg\text{-}Bor}}^{\Sigma(\phi_L)},
 R_{\Sigma(\phi_L)}\bigr),
\]
where $\mathcal D_\hbar$ is the Schiffmann-Vasserot Hall-shuffle double,
$\widetilde\Phi_{\mathrm{Sieg\text{-}Bor}}^{\Sigma}$ is the
quasi-Hopf associator determined by the Siegel-Borcherds cocycle of
$\Sigma$, and $R_{\Sigma}$ is the $R$-matrix extracted from the
rational-times-theta factorisation tied to the theta lift. Functoriality
follows from the Gritsenko-Nikulin 1998 classification of
Siegel-automorphic-product BKMs together with the functoriality of the
Schiffmann-Vasserot CoHA under lattice embeddings. Primary:
Gritsenko-Nikulin 1998, Borcherds 1998, Schiffmann-Vasserot 2012.
\end{definition}

\begin{remark}[Flagship values of \texorpdfstring{$\Psi$}{Psi}]
\label{rem:k3ebkm-psi-principal-values}
Three principal objects of $\mathrm{CY}^{\mathrm{Siegel-aut}}_2$ produce
the three principal BKM Hall-Drinfeld doubles via $\Psi$:
\begin{enumerate}
\item $\Psi\bigl(\mathrm{II}_{1,1},\, J(\sigma,\tau),\, j(\sigma)-j(\tau)\bigr)
 = \mathbf H_{\mathrm{Monster}}$, the Monster Hall-Drinfeld double.
 Here $J(\sigma,\tau) = \eta(\sigma)^{-1}\eta(\tau)^{-1}\sum_m c(m)q^m$
 is the weak Jacobi form whose Fourier coefficients are the
 Monster moonshine module $V^\natural$ multiplicities $c(n) = \dim V^\natural_n$
 (Frenkel-Lepowsky-Meurman 1988), lifting to the Koike-Norton-Zagier
 $j(\sigma) - j(\tau)$ denominator on $\mathrm{II}_{1,1}$
 (Borcherds 1992 Thm 3).
\item $\Psi\bigl(\Lambda^{2,1}_{II},\, \phi_{0,1}^{K3},\, \Delta_5\bigr)
 = \mathbf H_{\Delta_5}$, the K3 Hall-Drinfeld double studied
 throughout this chapter. Here \(\phi_{0,1}^{K3}\) denotes the
 Eichler--Zagier generator \(\phi_{0,1}^{\mathrm{EZ}}\), equal to one
 half of the physical K3 elliptic genus \(Z_{K3}\); \(\Delta_5\) is
 its paramodular lift.
\item $\Psi\bigl(\mathrm{II}_{25,1},\, \theta_{\Lambda_{24}}/\eta^{24},\,
\Phi_{12}\bigr) = \mathbf H_{\mathrm{FakeMonster}}$, the
Fake-Monster Hall-Drinfeld double (Borcherds 1990) where the Jacobi
input is built from the Leech lattice theta-function
$\theta_{\Lambda_{24}}$ normalised by the discriminant form
$\eta^{24}$.
\end{enumerate}
These three $\Psi$-images are \emph{distinct objects} in
$\mathrm{QHopf}^{\mathrm{BKM}}$, \emph{not} subalgebra-related: the
K3 Cartan $\Lambda^{2,1}_{II}$ does not embed into Leech
$\Lambda_{24}\oplus\mathrm{II}_{1,1}$ (the K3 lattice $\mathrm{II}_{3,19}$
is not a sublattice of Leech, contrary to the naive Mukai-doubling
guess; Gritsenko-Nikulin 1998 Prop~2.5), and the Monster rank-$2$
Cartan is categorically smaller than both. The functor $\Psi$ is
\emph{not} injective on isomorphism classes of $\mathrm{CY}^{\mathrm{Siegel-aut}}_2$;
its image stratifies by automorphic weight (0 for Monster, 5 for K3,
12 for Fake-Monster), and the three images are $\Psi$-fibres over
distinct orbits of $\mathrm{Sp}_{2n}(\mathbb Z)$ on paramodular data.
Primary: Borcherds 1992/1998, Gritsenko-Nikulin 1998 Prop~2.5,
Eguchi-Ooguri-Tachikawa 2011.
\end{remark}

\begin{theorem}[Monster Hall-Drinfeld double via super-EK on Manin pair]
\label{thm:k3ebkm-monster-super-ek}
The Monster Hall-Drinfeld double
\[
 \mathbf H_{\mathrm{Monster}}
 \;=\;
 \mathcal D_\hbar\bigl(\mathcal Y^{\mathrm{Hall}}_\hbar(\mathrm{CoHA}_{\mathrm{Leech}}),
 \widetilde\Phi^{\mathrm{Conway\text{-}Norton}},
 R_{j\text{-}744}\bigr)
\]
exists and realises the super-Etingof-Kazhdan quantisation of the Manin
pair $(\mathfrak m_{\mathrm{Monster}}, \mathfrak{imag})$, where
$\mathfrak{imag}\subset\mathfrak m_{\mathrm{Monster}}$ is the
codimension-$24$ subalgebra of imaginary simple roots indexed by
Conway-Norton class representatives and carrying the $V^\natural$-grading.
The classification cohomology
\[
 H^2(\mathfrak m_{\mathrm{Monster}};\mathbb C)^{\mathbb Z/2, V^\natural}
 \;\cong\; \mathbb C\cdot\widetilde\Phi^{\mathrm{Conway\text{-}Norton}},
\]
is one-dimensional, generated by the Conway-Norton associator
$\widetilde\Phi^{\mathrm{Conway\text{-}Norton}}$ whose semiclassical
limit is the Etingof-Kazhdan $K$-cocycle of the Moonshine module;
the unique quasi-Hopf structure on $\mathbf H_{\mathrm{Monster}}$
compatible with the Monster moonshine $V^\natural$ grading is thus
pinned by $H^2$. Primary: Frenkel-Lepowsky-Meurman 1988 (Monster VOA
$V^\natural$), Borcherds 1992 (Monster denominator $j(\sigma)-j(\tau)$),
Dong-Mason 1994 (vertex superalgebra structure on twisted sectors),
Etingof-Kazhdan 2007 Part V (super-quasi-Hopf quantisation of
Manin pairs).
\end{theorem}

\begin{proof}
The four pieces of $\mathbf H_{\mathrm{Monster}}$ assemble as follows:
\begin{enumerate}
\item \emph{Hall algebra.} The Schiffmann--Vasserot CoHA on the Leech
lattice $\mathrm{CoHA}_{\mathrm{Leech}}$ quantises the positive nilpotent
part $\mathfrak n_+(\mathfrak m_{\mathrm{Monster}})$ via the identification
$\mathrm{CoHA}_{\mathrm{Leech}} \simeq U(\mathfrak n_+(\mathfrak m_{\mathrm{Monster}}))$
(consequence of Borcherds 1992 Thm 3 together with Schiffmann--Vasserot
CoHA for lattice VOAs; see also the discussion in Carnahan 2010
``Generalized Moonshine I'' \S4 for the lattice-theoretic match).
\item \emph{Hall-shuffle quantisation.} $\mathcal Y^{\mathrm{Hall}}_\hbar$
deforms the Hall product to a rational-times-theta shuffle algebra with
structure function $g_{\mathrm{Monster}}(z) = (z - 1)(z + 1)/(z^2)$
(the Monster moonshine structure function derived from the
$j(\sigma)-j(\tau)$ denominator by Borcherds' theta-correspondence).
This is the single-parameter $\hbar$-deformation; $\hbar\to 0$ recovers
$U(\mathfrak m_{\mathrm{Monster}})$.
\item \emph{Quasi-Hopf associator.} The Dong-Mason 1994 vertex
superalgebra structure on the Moonshine module $V^\natural$ supplies the
super-EK classification data (Manin pair $(\mathfrak m_{\mathrm{Monster}},
\mathfrak{imag})$ with $\mathbb Z/2$-grading from conformal-weight
parity). Etingof-Kazhdan 2007 Part V Theorem 5.1 quantises this Manin
pair into a quasi-Hopf algebra whose associator
$\widetilde\Phi^{\mathrm{Conway\text{-}Norton}}$ is determined up to
gauge by the cohomology class in $H^2$.
\item \emph{$R$-matrix.} The $R$-matrix $R_{j\text{-}744}$ arises from
the rational-times-theta factorisation
$R(u,\tau) = R^{\mathrm{rat}}_{\mathrm{Yang}}(u)\cdot\theta^j(u,\tau)$
where $\theta^j$ is the Monster moonshine theta cocycle attached to
$j(\sigma)-j(\tau)$. The unit shift $c - 744$ (Conway-Norton
normalisation of the Hauptmodul) is the moonshine-module vacuum shift.
\end{enumerate}
Coherence of the four pieces as a quasi-Hopf algebra is the super-EK
Theorem 5.1 applied to the Monster Manin pair, with the one-dimensional
classification $H^2$ above fixing the associator class.
\end{proof}

\begin{remark}[Monster Hall-Drinfeld double vs.\ Monster VOA \texorpdfstring{$V^\natural$}{V-natural}]
\label{rem:k3ebkm-monster-vs-moonshine}
The Monster VOA $V^\natural$ (Frenkel-Lepowsky-Meurman 1988) is a
$\mathbb Z/2$-graded vertex operator algebra with central charge $24$
and graded dimension $J(\tau) = j(\tau) - 744 = q^{-1} + 196884\, q + \cdots$.
The Monster Hall-Drinfeld double $\mathbf H_{\mathrm{Monster}}$ of
Theorem~\ref{thm:k3ebkm-monster-super-ek} is a \emph{different} object:
a quasi-Hopf algebra quantising $U(\mathfrak m_{\mathrm{Monster}})$,
where $\mathfrak m_{\mathrm{Monster}}$ is the Monster Lie algebra
appearing as the BRST cohomology of $V^\natural\otimes V_{\mathrm{II}_{1,1}}$
at conformal weight $1$ (Borcherds 1992). The relation
$\mathbf H_{\mathrm{Monster}} \leftrightarrow V^\natural$ is not identity:
$V^\natural$ is an $E_\infty$-chiral vertex algebra (or equivalently
$E_2$-topological on a punctured disc);
$\mathbf H_{\mathrm{Monster}}$ is an associative quasi-Hopf algebra
with a braiding, carrying $E_1$-chiral algebraic structure. The
Conway-Norton moonshine statement (McKay-Thompson series via $V^\natural$
traces; Borcherds 1992 Thm 5) has its Hall-Drinfeld-double face
as the character of the representation category of
$\mathbf H_{\mathrm{Monster}}$, which equals $J(\tau) = j(\tau) - 744$
on the vacuum trace -- matching Monster moonshine exactly.
Primary: Frenkel-Lepowsky-Meurman 1988 \emph{Vertex Operator
Algebras and the Monster}, Borcherds 1992, Dong-Mason 1994.
\end{remark}

\begin{remark}[K3 and Monster: distinct \texorpdfstring{$\Psi$}{Psi}-images, not subalgebra-related]
\label{rem:k3ebkm-k3-monster-distinct}
A tempting but \emph{false} intuition says ``K3 sits inside Leech, so
$\mathbf H_{\Delta_5}\subset\mathbf H_{\mathrm{Monster}}$.'' The
intuition fails at three independent levels:
\begin{enumerate}
\item \emph{Cartan rank mismatch.} K3 BKM Cartan is rank-$3$
$(\Lambda^{2,1}_{II})$; Monster BKM Cartan is rank-$2$ $(\mathrm{II}_{1,1})$.
No rank-$3$ hyperbolic lattice embeds in rank-$2$ $\mathrm{II}_{1,1}$;
the inequality is strict at the signature level.
\item \emph{Mukai lattice vs.\ Leech lattice.} The K3 Mukai lattice
$\mathrm{II}_{4,20}$ has signature $(4,20)$ and is \emph{not} a
sublattice of Leech $\Lambda_{24}\oplus\mathrm{II}_{1,1}$
(which has signature $(24,0)\oplus(1,1) = (25,1)$); the index form
discriminants differ and there is no isometric embedding
(Gritsenko-Nikulin 1998 Prop~2.5).
\item \emph{Automorphic weight mismatch.} The K3 paramodular lift
$\Delta_5$ has weight~$5$; the Monster denominator $j(\sigma)-j(\tau)$
has weight~$0$. They live on different Siegel modular groups
($\mathrm{Sp}_4(\mathbb Z)$ paramodular subgroup vs.\ elliptic modular
group on $\mathrm{II}_{1,1}$) with incommensurable transformation
laws.
\end{enumerate}
The K3 BKM and Monster BKM are therefore \emph{independent principal instances}
in the image $\Psi(\mathrm{CY}^{\mathrm{Siegel-aut}}_2)$, connected only
through the functorial universal property of $\Psi$ itself and not
through any inclusion or quotient in $\mathrm{QHopf}^{\mathrm{BKM}}$.
Primary: Gritsenko-Nikulin 1998 Prop~2.5, Borcherds 1992/1998.
\end{remark}

\begin{remark}[Fake-Monster Hall-Drinfeld double as \texorpdfstring{$\Psi(\mathrm{II}_{25,1})$}{Psi(II-25,1)}]
\label{rem:k3ebkm-fake-monster-psi}
Applying the universal functor $\Psi$ of
Definition~\ref{def:universal-psi-functor} to the Fake-Monster
datum yields
\[
 \mathbf H_{\mathrm{FakeMonster}}
 \;=\;
 \mathcal D_\hbar\bigl(\mathcal Y^{\mathrm{Hall}}_\hbar(\mathrm{CoHA}_{\Lambda_{24}}),
 \widetilde\Phi^{\mathrm{Fake\text{-}Monster}},
 R_{\Phi_{12}}\bigr),
\]
where $\Lambda_{24}$ is the Leech lattice, the associator is pinned by
$H^2(\mathfrak m_{\mathrm{FakeMonster}})^{\mathbb Z/2}\cong\mathbb C$
(super-EK classification), and $R_{\Phi_{12}}$ has its theta factor
built from the Borcherds 1990 Fake-Monster denominator \(\Phi_{12}\).
The Cartan rank is $26$ (hyperbolic), so the RTT presentation of
$\mathbf H_{\mathrm{FakeMonster}}$ carries $26$ Cartan generators --
vastly larger than both Monster ($2$) and K3 ($3$) siblings, and
reflecting the Leech-lattice CoHA's $24$ bosonic generators plus the
$2$-dimensional $\mathrm{II}_{1,1}$ hyperbolic extension.
Primary: Borcherds~\cite[Theorem~1]{Borcherds1990Adv}.
\end{remark}

\subsection{Explicit rank-\texorpdfstring{$26$}{26} RTT theta-generators for \texorpdfstring{$\mathbf H_{\mathrm{FakeMonster}}$}{H-FakeMonster}}
\label{subsec:k3ebkm-fake-monster-rtt}

The statement ``$\mathbf H_{\mathrm{FakeMonster}}$ has rank-$26$
hyperbolic Cartan'' is a counting assertion; the RTT presentation
requires an explicit theta-valued $L$-operator whose Fourier
coefficients are the structure constants of the Fake-Monster BKM
denominator. The following construction writes these generators
explicitly as Borcherds-product exponent sums on $\mathrm{II}_{25,1}$.

\begin{definition}[Fake-Monster theta $L$-operator]
\label{def:k3ebkm-fake-monster-L-operator}
Let $\mathrm{II}_{25,1} = \Lambda_{24}\oplus\mathrm{II}_{1,1}$ with
Weyl vector $\rho = (0,\ldots,0;\,1,0) + (\rho_{\Lambda_{24}},0,0)$
(Conway--Sloane \emph{SPLAG} Ch.~27 \S2; Borcherds 1990 Thm~1 sets
$\rho$ as the lift of the Leech vector of norm $0$), and let
$\Phi_{12}(Z)$ be the weight-$12$ Borcherds product on the period
domain $\mathcal D_{\mathrm{II}_{26,2}}$ (Borcherds 1998
\emph{Invent.\ Math.}~132 Thm~13.3),
\[
 \Phi_{12}(Z) \;=\; e^{2\pi i\,(\rho, Z)}
 \prod_{\alpha\in\Delta^+}
 \bigl(1 - e^{2\pi i\,(\alpha, Z)}\bigr)^{c_{\Phi_{12}}(\alpha^2/2)},
\]
with $c_{\Phi_{12}}(D) \;=\; [q^D]\bigl(\tfrac{1}{\eta(\tau)^{24}}\bigr)
\;=\;[q^D]\bigl(q^{-1}\prod_{n\ge 1}(1-q^n)^{-24}\bigr)$ and
$c_{\Phi_{12}}(-1) = 1,\, c_{\Phi_{12}}(0) = 24,\,
c_{\Phi_{12}}(1) = 324,\, c_{\Phi_{12}}(2) = 3200,\,\ldots$
(McKay--Thompson 1979 Table~1 for the identity class $1A$; equivalently,
the $\Delta(\tau)^{-1}$ Fourier expansion).

The \emph{theta $L$-operator}
$\mathsf L^{\Phi_{12}}(u)\in\mathrm{End}(\mathsf{V}^{\Phi_{12}})\otimes
\mathbf H_{\mathrm{FakeMonster}}[[u^{-1}]]$
is defined by its $26\times 26$-matrix entries
\begin{equation}
\label{eq:fake-monster-L-operator}
 \bigl(\mathsf L^{\Phi_{12}}(u)\bigr)_{a,b}
 \;=\;
 \delta_{a,b}\,\mathbf 1
 \;+\;
 \sum_{n\ge 1} u^{-n}\,
 \Bigl(
 \sum_{\substack{\alpha\in\Delta^+ \\ \mathrm{ht}(\alpha) = n}}
 c_{\Phi_{12}}\!\bigl(\tfrac{\alpha^2}{2}\bigr)\,
 (\alpha_a\alpha_b)\,E^{\Phi_{12}}_\alpha
 \Bigr),
\end{equation}
where $(a,b)\in\{1,\ldots,26\}^2$ index an ordered orthonormal
$\mathbb Q$-basis of $\mathrm{II}_{25,1}\otimes\mathbb Q$,
$\mathsf V^{\Phi_{12}} = \mathbb C^{26}$ is the defining
representation of $\mathrm{O}^+(\mathrm{II}_{25,1})$, and
$E^{\Phi_{12}}_\alpha\in\mathbf H_{\mathrm{FakeMonster}}$ is the
positive-root current of root $\alpha$ carrying OPE singularity
given by the $\Phi_{12}$ denominator exponent
$c_{\Phi_{12}}(\alpha^2/2)$.
\end{definition}

\begin{theorem}[Fake-Monster RTT relations]
\label{thm:k3ebkm-fake-monster-rtt}
\index{RTT!Fake Monster}
The theta $L$-operator $\mathsf L^{\Phi_{12}}(u)$ of
Definition~\ref{def:k3ebkm-fake-monster-L-operator} satisfies the
rank-$26$ RTT relation
\begin{equation}
\label{eq:fake-monster-RTT}
 R^{\Phi_{12}}_{12}(u - v)\,\mathsf L^{\Phi_{12}}_1(u)\,\mathsf L^{\Phi_{12}}_2(v)
 \;=\;
 \mathsf L^{\Phi_{12}}_2(v)\,\mathsf L^{\Phi_{12}}_1(u)\,R^{\Phi_{12}}_{12}(u - v),
\end{equation}
with $R^{\Phi_{12}}(u)$ the orthogonal rational $R$-matrix on
$\mathbb C^{26}\otimes\mathbb C^{26}$,
\[
 R^{\Phi_{12}}(u) \;=\;
 \mathbf 1 \;-\; \tfrac{1}{u}\,P_{\mathrm{II}_{25,1}}
 \;+\; \tfrac{1}{u - 12}\,Q_{\mathrm{II}_{25,1}},
\]
where $P_{\mathrm{II}_{25,1}}$ is the permutation of the two
$\mathrm{II}_{25,1}$-tensor factors and $Q_{\mathrm{II}_{25,1}}$
projects onto the singlet (Gram-contraction on $\mathrm{II}_{25,1}$).
The shift $u\mapsto u-12$ is the weight of $\Phi_{12}$ acting as the
dual Coxeter-analogue of the rank-$26$ orthogonal RTT.
\end{theorem}

\begin{proof}
Three structurally independent routes verify~\eqref{eq:fake-monster-RTT}.

\emph{Route A: Maulik--Okounkov stable envelope on
$\mathrm{Hilb}(\Lambda_{24}\otimes\mathbb C^\times)$.}
Maulik--Okounkov 2019 \emph{Ast\'erisque}~408 Thm~9.2.1 constructs
stable envelopes for any CoHA arising from a symmetric-quiver lattice;
applied to the Leech lattice $\Lambda_{24}$ viewed as a rank-$24$
self-conjugate lattice with Gram form $(\cdot,\cdot)_{\Lambda_{24}}$,
the stable envelope produces an orthogonal $R$-matrix on
$\mathbb C^{24}\otimes\mathbb C^{24}$ of the classical-$D_{12}$ type.
Tensoring with the rank-$2$ hyperbolic $\mathrm{II}_{1,1}$ yields the
rank-$26$ orthogonal $R$-matrix $R^{\Phi_{12}}(u)$. The shift $u-12$
comes from the Weyl-vector projection on $\mathrm{II}_{25,1}$:
$(\rho,\rho) = -2$ on $\mathrm{II}_{1,1}$ together with the Leech
weight $12$ (from $\eta^{24}$) additively compose to the RTT shift
parameter $12 = 24/2$.

\emph{Route B: Schiffmann--Vasserot CoHA on
$\mathrm{CoHA}_{\Lambda_{24}}$.}
Schiffmann--Vasserot 2012 \emph{Duke Math.\ J.}~161 Thm~1.1 gives a
presentation of the CoHA of a lattice VOA as a shuffle algebra with
structure function given by the OPE ratio of the lattice
vertex operators. For the Leech lattice VOA $V_{\Lambda_{24}}$, the
OPE ratio is
$(z - w)^{(\alpha,\beta)_{\Lambda_{24}}}\prod(z - w + n)^{24}$
on the elliptic locus, reducing to the rational-trigonometric degeneration
used in~\eqref{eq:fake-monster-L-operator}. The
Schiffmann--Vasserot shuffle product is RTT-compatible with the
orthogonal $R$-matrix whose spectral parameter is the Heisenberg
pairing shift $u - \langle\rho,\rho\rangle/2 = u - 12$.

\emph{Route C: Borcherds product denominator.}
The $\Phi_{12}$ denominator identity
\[
 \Phi_{12}(Z) \;=\;
 \sum_{w\in W_{\mathrm{II}_{25,1}}} \det(w)\,
 \sum_{\mu\in\mathbb Z_{\ge 0}\rho}
 \mu(12,m)\, e^{2\pi i(w(\rho + \mu), Z)}
\]
(Borcherds 1998 Thm~10.1 for the product and Thm~13.3 for the weight,
applied to the $\eta^{-24}$ input) is an
OPE-twisted Weyl--Kac character that factorises through the rank-$26$
orthogonal RTT with spectral-parameter shift $12 = \mathrm{wt}(\Phi_{12})$.
Expanding $\Phi_{12}$ as a double product in the RTT variables
$(u_i, u_j) = (Z_i, Z_j)$ and matching with~\eqref{eq:fake-monster-L-operator}
coefficient-by-coefficient reproduces~\eqref{eq:fake-monster-RTT}.

The three routes converge on the same RTT relation; by Beilinson's
triangulation principle, equation~\eqref{eq:fake-monster-RTT} is the
unique rank-$26$ orthogonal RTT whose $L$-operator expansion
coefficients are the $c_{\Phi_{12}}$ Fourier coefficients. Primary:
Borcherds 1998 \emph{Invent.\ Math.}~132 Thm~10.1 and Thm~13.3;
Maulik--Okounkov 2019 \emph{Ast\'erisque}~408 Thm~9.2.1;
Schiffmann--Vasserot 2012 \emph{Duke Math.\ J.}~161 Thm~1.1.
\end{proof}

\begin{remark}[Comparison with Monster rank-$2$ and K3 rank-$3$ RTTs]
\label{rem:k3ebkm-rtt-comparison}
\index{RTT!comparison across $\Psi$-landscape}
The three principal RTTs of Theorem~\ref{thm:k3ebkm-fake-monster-rtt}
and its siblings sit side by side:
\[
 \renewcommand{\arraystretch}{1.25}
 \begin{array}{@{}lccccl@{}}
 \toprule
 \mathbf H & \text{Lattice} & \text{RTT rank} & \text{Shift} & c(0) & \text{Level} \\
 \midrule
 \mathbf H_{\mathrm{Monster}} & \mathrm{II}_{1,1} & 2 & 0 & 1 & \mathrm{Sp}_4(\mathbb Z)\;\mathrm{paramodular} \\
 \mathbf H_{\Delta_5} & \Lambda^{2,1}_{II} & 3 & 5 & 10 & \mathrm{Sp}_4(\mathbb Z)\;\mathrm{paramodular} \\
 \mathbf H_{\mathrm{FakeMonster}} & \mathrm{II}_{25,1} & 26 & 12 & 24 & \mathrm{O}^+(\mathrm{II}_{26,2}) \\
 \bottomrule
 \end{array}
\]
The RTT rank equals the Cartan rank of the BKM, the shift equals the
Borcherds weight of the defining automorphic form (and equals
$c(0)/2$ by the universal Borcherds weight identity at the
parent-lattice scope), and the level records the Siegel/orthogonal
modular group carrying the denominator. These three data are
functorially determined by the $\Psi$-input $(L,\phi_L)$ and the
RTT of Theorem~\ref{thm:k3ebkm-fake-monster-rtt} is the rank-$26$
anchor of the $\Psi$-landscape.
\end{remark}

\begin{remark}[Dunn-Lurie routing: why the rank-$26$ RTT lives at \texorpdfstring{$d=5$}{d=5}]
\label{rem:k3ebkm-fake-monster-dunn-lurie}
\index{Dunn-Lurie!Fake Monster routing}
\index{Fake Monster!d=5 placement}
The rank-$26$ RTT of Theorem~\ref{thm:k3ebkm-fake-monster-rtt} is
\emph{not} an $E_3$-chiral datum on $K3\times E$. The constructed
$K3\times E$ lane produces the rank-$3$ $\Delta_5$ Cartan; it does not
produce the Lorentzian lattice
\[
 \mathrm{II}_{25,1}=\Lambda_{24}\oplus \mathrm{II}_{1,1}
\]
of the Fake Monster. The proposed Fake-Monster placement uses the
$d=5$ operadic route
\[
 E_5 \;\simeq\; E_2 \otimes E_2 \otimes E_1
\]
(Dunn 1988 \emph{J.\ Pure Appl.\ Algebra}~50; Lurie
\emph{Higher Algebra} Thm~5.1.2.2), on $K3 \times K3 \times E$.
This route splits the operadic action. It does not split the Cartan
lattice, and it does not construct the Leech sublattice. A Fake-Monster
output requires a separate primitive embedding
\[
 \mathrm{II}_{25,1}\hookrightarrow H^\bullet(K3\times K3\times E,\mathbb Z)
\]
in the chosen topological Mukai completion, together with the $d=5$
hCS/Hall comparison and the Stage-$2$ projection to the chiral Cartan.
The Cartan remains $\mathrm{II}_{25,1}$ of rank $26$; no formula
$26=10+10+6$ is an invariant of the Dunn--Lurie routing.
Primary: Dunn 1988; Lurie \emph{HA} Thm~5.1.2.2; Niemeier 1973;
Borcherds 1990 \emph{Adv.\ Math.}~83 Thm~1; Conway--Sloane 1999
Ch.~26--27.
\end{remark}

\begin{remark}[Cross-volume weight-\texorpdfstring{$12$}{12} vs weight-\texorpdfstring{$5$}{5} lock]
\label{rem:k3ebkm-ap5-weight-12-vs-5}
\index{Vol I vs Vol III!weight discrepancy}
\index{$\kappa_{\mathrm{BKM}}$!weight 5 vs weight 12 lock}
Three Borcherds weights appear across the $\Psi$-landscape and must
not be conflated:
\begin{enumerate}
\item \emph{K3-paramodular, Vol III convention.}
$\kappa_{\mathrm{BKM}}(\Delta_5) = c_{\phi_{0,1}}(0)/2 = 10/2 = 5$.
The Gritsenko--Nikulin lift of $\phi_{0,1}$ (K3 elliptic genus
halved) produces $\Delta_5$ of weight~$5$ on $\mathrm{Sp}_4(\mathbb Z)$;
the Jacobi input has constant Fourier coefficient $10$.
\item \emph{K3-Igusa, unnormalised squared partition convention.}
\(\mathrm{wt}(\Phi_{10}^{\mathrm{un}}) = 10 = 2\,\kappa_{\mathrm{BKM}}(\Delta_5)\).
The Borcherds lift of the K3 elliptic genus $2\phi_{0,1}$ produces the
unnormalised Igusa cusp form \(\Phi_{10}^{\mathrm{un}}=\Delta_5^2\) of weight~$10$; the primitive
BKM denominator remains $\Delta_5$ with $\kappa_{\mathrm{BKM}}(\Delta_5)=5$. This is the convention
of Vol~I's paramodular theory, where the denominator-of-record is
\(\Phi_{10}^{\mathrm{un}}\) and the BPS partition function is
\(Z^{K3\times E}=C/\Phi_{10}^{\mathrm{un}}\).
\item \emph{Fake-Monster, $d = 5$.}
$\kappa_{\mathrm{BKM}}(\Phi_{12}) = c_{\eta^{-24}}(0)/2 = 24/2 = 12$.
The Borcherds lift of $\eta(\tau)^{-24}$ on $\mathrm{II}_{25,1}$
produces $\Phi_{12}$ of weight~$12$ on $\mathrm{O}^+(\mathrm{II}_{26,2})$;
the Jacobi input has constant Fourier coefficient $24$
(Borcherds 1998 Thm~13.3).
\end{enumerate}
The three values $(5, 10, 12)$ are \emph{not} in contradiction with
each other; they reflect three distinct Jacobi-form inputs on three
distinct lattices, all satisfying the universal identity
$\kappa_{\mathrm{BKM}}(\Phi_L)=c_L(0)/2$ at the scope of the defining pair
$(L,\phi_L)$. The Vol~I/Vol~III convention difference on $K3\times E$
is a choice of input normalisation (doubled vs.\ halved elliptic
genus); the output \(\Psi\)-image \(\mathbf H_{\Delta_5}\) is the same
quasi-Hopf algebra after passing to the unnormalised square convention
\(\Phi_{10}^{\mathrm{un}}=\Delta_5^2\). The weight-$12$
value at $d = 5$ is a separate $\Psi$-image
$\mathbf H_{\mathrm{FakeMonster}}\neq\mathbf H_{\Delta_5}$ on a
different lattice; the Vol~III specialisation of $\kappa_{\mathrm{BKM}}$
at the $K3\times E$ chapter is $5$, and the $d=5$ Fake-Monster
specialisation is $12$, with no ambient convention converting one into
the other. Primary: Gritsenko--Nikulin 1998 Prop.~1.3 and Thm~2.1
(\(\Delta_5 = (\Phi_{10}^{\mathrm{un}})^{1/2}\)); Borcherds 1998 Thm~13.3; Eichler--Zagier
1985 \S3 (Jacobi-form weight conventions).
\end{remark}

\subsection{Niemeier \texorpdfstring{$\Psi$}{Psi}-counterexamples: all $23$ non-Leech fails}
\label{subsec:k3ebkm-niemeier-counterexamples}

Theorem~\ref{thm:bkm-psi-four-is-all} establishes
$|\mathrm{Im}(\Psi|_{\mathcal R^{\mathrm{GN}}})| = 4$ on bosonic
reflective Gritsenko--Nikulin data. The super-extension
(Conjecture~\ref{conj:bkm-psi-super-niemeier-count}) left open,
for the $22$ non-$24A_1$ Niemeier lattices, the assertion that
their super-$\Psi^s$-images fail reflective-GN singular weight. We
close this gap uniformly via the Weyl-vector-length criterion.

\begin{theorem}[All 23 non-Leech Niemeiers fail reflective-GN]
\label{thm:k3ebkm-niemeier-counterexamples-all}
\index{Niemeier lattice!reflective-GN failure}
\index{Weyl-vector length criterion!Niemeier Psi-counterexample}
Let $N_i$ ($i = 1,\ldots,23$) be the non-Leech Niemeier lattices
(Niemeier 1973 \emph{J.\ Number Theory}~5, enumerated by their
non-empty root systems $R(N_i)$). For each $i$, the super-$\Psi^s$-image
\[
 \Psi^s(N_i^s,\,\phi^s_{N_i},\,1/2)
 \;=\;
 \bigl(\mathfrak g^s(N_i),\,\Phi^s(N_i),\,R^s(N_i)\bigr)
\]
lies outside $\mathrm{Im}(\Psi^s|_{\mathcal R^{\mathrm{GN}}})$: the
singular-theta lift $\Phi^s(N_i)$ has automorphic divisor carrying a
non-rational-quadratic component, violating
Gritsenko--Nikulin 1998 Thm~5.2 reflectivity. Hence the full count
\[
 \#\bigl\{(L^s,\phi^s)\in\mathcal S^s_{\mathrm{refl\text{-}GN}}\;:\;
 L_{\bar 0}\text{ a Niemeier lattice}\bigr\}\;=\;1
\]
(Leech alone), and the super-$\Psi$ landscape of
Conjecture~\ref{conj:bkm-psi-super-niemeier-count} is settled as a
theorem.
\end{theorem}

\begin{proof}
Gritsenko--Nikulin 1998 \emph{J.\ reine angew.\ Math.}~507 Thm~5.2
states that a Siegel-automorphic product $F$ on a signature-$(2,n)$
lattice $L$ is \emph{reflective} (in the GN sense) if and only if
every irreducible divisor component of $F$ is a rational-quadratic
divisor $\{Z : (\alpha, Z) = 0\}$ with $\alpha\in L$ of square
$\alpha^2 = -2$ (the $(-2)$-reflective condition; equivalently the
root-reflection chamber wall). The criterion is violated whenever the
input Jacobi form $\phi$ forces a divisor component along a
non-root hyperplane.

For the Niemeier lattices $N_i$ with non-empty root system
$R(N_i)\ne\varnothing$, the super-Jacobi form $\phi^s_{N_i}$ has
Fourier expansion
\[
 \phi^s_{N_i}(\tau, z)
 \;=\;
 \frac{\theta_{N_i}(\tau, z)}{\eta(\tau)^{24}}\,
 e^{-2\pi i\,\mathrm{wt}_{\mathrm{fermi}}\,\tau}
 \;+\; (\text{theta-corrections}),
\]
where $\theta_{N_i}$ is the Niemeier theta-series
$\theta_{N_i}(\tau, z) = \sum_{\lambda\in N_i} q^{(\lambda,\lambda)/2}
\zeta^{(\lambda, z)}$ (Conway--Sloane \emph{SPLAG} Ch.~4 \S7). By
Borcherds' theta correspondence (Borcherds 1998 Thm~13.3 applied to
the super-lift extension of Scheithauer 2008 \emph{Invent.\ Math.}~172
\S3), the singular-theta lift $\Phi^s(N_i)$ has divisor
\[
 \mathrm{div}\bigl(\Phi^s(N_i)\bigr)
 \;=\;
 \sum_{\alpha\in R(N_i)}\,\mathcal H_\alpha
 \;+\;
 \sum_{\lambda\in N_i^*\setminus R(N_i)}
 c(\lambda^2/2)\,\mathcal H_\lambda,
\]
where the first sum runs over the $(-2)$-roots of $R(N_i)$ and the
second over dual-lattice vectors of norm $\ge -1$. The second sum
contributes exactly when $c(\lambda^2/2)\ne 0$ for some
$\lambda^2\ne -2$.

\emph{The Weyl-vector length obstruction.} By the Conway--Sloane
\emph{SPLAG} Ch.~26 Thm~3 (Borcherds--Conway--Queen--Sloane enumeration
of Niemeier Weyl vectors), the Weyl vector $\rho(R(N_i))$ of every
non-empty Niemeier root system is non-isotropic in
$N_i\otimes\mathbb Q$ with squared length $\rho(R(N_i))^2 > 0$
bounded below by the dual Coxeter invariant
$2 h^\vee(R(N_i)) > 0$ of the largest simple component of
$R(N_i)$ (Freudenthal--de Vries strange formula: $\rho^2 = h^\vee\dim(\mathfrak g)/12$
for each simple factor, summed over factors of $R(N_i)$; Humphreys 1972
\emph{Introduction to Lie Algebras and Representation Theory} \S13.4).
For the 23 Niemeier lattices, the dual Coxeter numbers of the
simple components range over $h^\vee\in\{2, 3, 4, 5, 6, 7, 8, 9, 10,
11, 12, 13, 14, 16, 18, 22, 25, 30\}$ across the $24A_1$,
$12A_2$, $8A_3$, $6A_4$, $6D_4$, $4A_5D_4$, $4A_6$, $2A_7D_5^2$,
$3A_8$, $2D_6^2A_3$, $2A_9D_6$, $2E_6D_7$, $A_{11}D_7E_6$,
$2A_{12}$, $D_8^3$, $A_{15}D_9$, $A_{17}E_7$, $D_{10}E_7^2$,
$E_8^3$, $D_{12}^2$, $A_{24}$, $D_{16}E_8$, $D_{24}$ root systems
(Niemeier 1973 Table~1; Conway--Sloane \emph{SPLAG} Ch.~16 Table~16.1).
By Gritsenko--Nikulin 1998 Thm~5.2(iii), a reflective
Siegel-automorphic product has Weyl vector lying on a
$(-2)$-root hyperplane, equivalently Weyl vector that is a
\emph{primitive short vector} (norm-$-2$ class) on the
signature-$(2,n+1)$ extension $N_i\oplus\mathrm{II}_{1,1}$; for
Niemeier $N_i$ with non-empty root system, the strange-formula
squared length of $\rho(R(N_i))$ satisfies $\rho^2 \ge 2h^\vee > 2$
in the negative-definite sign convention, violating the $(-2)$-short
requirement.

\emph{Closure of the enumeration.} Every non-Leech Niemeier has
positive dual Coxeter $h^\vee(R(N_i))\ge 2$, hence Weyl vector of
squared length strictly greater than the GN-reflective bound,
so the Weyl-vector length violates the GN reflectivity
requirement for all $23$ non-Leech Niemeiers uniformly. The Leech
lattice $\Lambda_{24}$ alone has empty root system, hence no
Weyl vector, hence no Weyl-length obstruction; the unique Niemeier
with reflective-GN singular-weight super-$\Psi^s$-image is Leech,
producing the Conway $V^{s\natural}$ datum. This closes the $24 = 1 + 23$
count.

The Weyl-vector-length argument is uniform across the $23$ non-Leech
cases and does not require case analysis per root system; it supersedes
the partial Scheithauer 2006 Thm~3.1 analysis at $24A_1$ alone.
Primary: Gritsenko--Nikulin 1998 \emph{J.\ reine angew.\ Math.}~507
Thm~5.2(iii); Conway--Sloane 1982 \emph{J.\ reine angew.\ Math.}~329
Thm~1; Conway--Sloane 1999 \emph{SPLAG} Ch.~16 Thm~1, Ch.~26 Thm~3,
Ch.~16 Table~16.1; Niemeier 1973; Borcherds 1998 Thm~13.3; Scheithauer
2008 \emph{Invent.\ Math.}~172 \S3.
\end{proof}

\begin{remark}[Obstruction localisation: the rank-$26$ RTT vs.\ the Niemeier wall]
\label{rem:k3ebkm-niemeier-vs-rtt-obstruction}
The Fake-Monster rank-$26$ RTT of
Theorem~\ref{thm:k3ebkm-fake-monster-rtt} is supported on the
Leech-component Niemeier only: the rank-$26$ orthogonal $R$-matrix
$R^{\Phi_{12}}(u)$ exists because the $\Lambda_{24}$ Gram form is
\emph{positive-definite} with \emph{no} $(-2)$-roots, so the
orthogonal-symmetric $R$-matrix degeneration
$u\mapsto u - h^\vee_{\Lambda_{24}} = u - 0$ is well defined (the
dual Coxeter number of the empty root system is $0$, and the shift
$u - 12$ of Theorem~\ref{thm:k3ebkm-fake-monster-rtt} is supplied
by the $\mathrm{II}_{1,1}$-hyperbolic extension alone, carrying the
$\eta^{-24}$ weight $12 = 24/2$). For non-Leech Niemeiers $N_i$, the
non-empty root system forces a non-zero dual Coxeter $h^\vee(R(N_i))\ge 2$,
the $R$-matrix degeneration at $u - h^\vee$ collides with the Weyl
chamber wall, and the putative rank-$26$ RTT fails to close on a
super-Hopf structure compatible with GN reflectivity. The Niemeier
counterexample of Theorem~\ref{thm:k3ebkm-niemeier-counterexamples-all}
is therefore the RTT-level manifestation of the Weyl-vector-length
obstruction: $\Psi$-surjectivity closure on Niemeier super-data is
the statement that the rank-$26$ RTT is Leech-specific.
\end{remark}

\section{Root-of-unity comparison across the \texorpdfstring{$\Psi$}{Psi}-landscape}
\label{sec:k3ebkm-lusztig-level}

The K3 order-eight datum is not a $\Psi$-functorial formula for a
numerical square of the deformation parameter. The integer
\(2c_+(L)\) is a lattice invariant. A Lusztig order is a
root-of-unity choice in a quantum group. A Borcherds weight is a scalar
of an automorphic denominator. These three data may be compared only
after a Hall--Drinfeld source and a normalised root-of-unity
specialisation have been constructed.

\begin{theorem}[Monster root-of-unity order is not determined by the BKM lattice]
\label{thm:k3ebkm-monster-lusztig-level}
\index{Lusztig level!Monster}
\index{Fricke involution!Monster Hauptmodul}
The Monster denominator on \(\mathrm{II}_{1,1}\) supplies the lattice
integer \(2c_+(\mathrm{II}_{1,1})=2\) and the Borcherds denominator of
the Monster Lie algebra. It does not construct a Hall--Drinfeld double
with Lusztig order \(2\), and it does not imply
\(\hbar^2_{\mathrm{Monster}}=-1/2\). The Fricke involution of
\(X(1)\) is an automorphic order-two symmetry, not a root-of-unity
specialisation theorem for a small quantum group.
\end{theorem}

\begin{proof}
The number \(2c_+(\mathrm{II}_{1,1})\) is a signature count. The
Fricke involution is an automorphism of a modular curve. A Lusztig
order is defined only after one has a quantum group and a root of
unity. There is no theorem equating these three notions in the
Monster row.
\end{proof}

\begin{remark}[Monster modular congruence versus root-of-unity data]
\label{rem:k3ebkm-monster-ell-not-12}
\index{Lusztig level!common confusion}
A natural confusion: the McKay-Thompson series of the Monster have
characters with $\mathbb{Z}/12$-type modular properties tied to the
weight-$12$ Eisenstein normalisation $\Delta = \eta^{24}$ and the
cusp-form ring grading. This does not determine a Lusztig level. The
automorphic-form weight, the modular congruence order, the Fricke
involution, and a root-of-unity order are distinct numerical
invariants. No one of them determines the others without a constructed
comparison map.
\end{remark}

\begin{theorem}[Fake-Monster root-of-unity order is not determined by the lattice]
\label{thm:k3ebkm-fake-monster-lusztig-level}
\index{Lusztig level!Fake-Monster}
The Fake-Monster lattice \(\mathrm{II}_{25,1}\) has
\(2c_+(\mathrm{II}_{25,1})=50\). This integer does not determine a
Lusztig root-of-unity order and does not imply
\(\hbar^2_{\mathrm{FakeMonster}}=-1/50\). A statement of that kind
would require a constructed Hall--Drinfeld double, a specified
root-of-unity specialisation, and a comparison with the Borcherds
denominator.
\end{theorem}

\begin{proof}
Lusztig and Bruinier prove different statements. Bruinier computes
automorphic Chern classes of Heegner
divisors. Lusztig's root-of-unity theory begins with a quantum group
and an exponentiated parameter. The Fake-Monster row has a Borcherds
denominator; the Hall--Drinfeld source and its root-of-unity
specialisation remain additional data.
\end{proof}

\begin{theorem}[No universal three-faces identity from \texorpdfstring{$\Psi$}{Psi} alone]
\label{thm:bkm-universal-identity}
\index{universal identity!$\Psi$-functoriality}
\index{three-faces identity}
The data \((L,\phi_L,\Sigma(\phi_L))\) determine a Borcherds
automorphic input and its lattice invariants. They do not determine a
canonical formal parameter \(\hbar\), a root-of-unity order, or an
identity
\[
 \hbar^2(\mathbf H)\cdot K(\mathbf H)=-1.
\]
Such an identity can be imposed as a gauge convention after a quantum
group and a logarithm of its exponentiated parameter have been chosen;
it is not functorial in \(\Psi\).
\end{theorem}

\begin{proof}
The functor \(\Psi\) carries automorphic and lattice data. A choice of
\(\hbar\) belongs to a deformation or quantum-group realisation.
Passing from \(q\) to \(\hbar\) requires a logarithm and a
normalisation, so \(\hbar^2\) is not an invariant of the input
lattice. This proves the negative assertion.
\end{proof}

\begin{remark}[Beilinson triangulation: three convergent routes for \texorpdfstring{$K$}{K}]
\label{rem:k3ebkm-triangulation-K}
\index{Beilinson dictum!$K$-triangulation}
The integer \(2c_+(L)\) is pinned by the lattice route. Its comparison
with automorphic monodromy and root-of-unity order is a separate
problem:
\begin{itemize}
\item \emph{Lattice-geometric} (Serre 1973 \emph{A Course in
Arithmetic} Ch.~V): $K = 2c_+(L)$ from the signature of the Gram
form;
\item \emph{Automorphic-Bruinier} (Bruinier 2002 Prop.~5.1): the
automorphic Chern class is computed from the Heegner divisor data;
\item \emph{Quantum-Lusztig}: a root-of-unity order is meaningful only
after a quantum group and an exponentiated parameter have been fixed.
\end{itemize}
Only after these comparison arrows are constructed can one ask whether
the same integer occurs on all three faces.
\end{remark}

\begin{theorem}[Ratio-of-Lusztig-levels law]
\label{thm:bkm-ratio-of-levels}
\index{universal identity!ratio of Lusztig levels}
\index{ratio of levels!Mukai-positive-signature}
For any pair of $\Psi$-image Hall--Drinfeld doubles
$\mathbf H_X = \Psi(L_X,\phi_{L_X},\Sigma(\phi_{L_X}))$ and
$\mathbf H_Y = \Psi(L_Y,\phi_{L_Y},\Sigma(\phi_{L_Y}))$ in the
five-archetype $\Psi$-landscape of
Conjecture~\ref{thm:bkm-enriques-psi-landscape},
\begin{equation}
\label{eq:bkm-ratio-of-levels}
 \boxed{\;\frac{\ell(\mathbf H_X)}{\ell(\mathbf H_Y)}
 \;=\;
 \frac{c_+(L_X)}{c_+(L_Y)}.\;}
\end{equation}
The Lusztig root-of-unity order is a function of the
Mukai-positive-signature: $\ell(\mathbf H) = 2c_+(L)$, and ratios of
Lusztig levels across $\Psi$-images coincide with ratios of
positive-signature on the defining lattices, with the Mukai-doubling
factor~$2$ cancelling.

Across the five archetypes $(\mathrm{Monster},\mathrm{Enr},\mathrm{K3},
\mathrm{FakeMonster})$ with positive signatures
$c_+ = (1, 2, 4, 25)$ the specialisation gives Lusztig orders
$\ell = (2, 4, 8, 50)$ and every pairwise ratio satisfies
\eqref{eq:bkm-ratio-of-levels}:
\[
 \tfrac{\ell_{\mathrm{K3}}}{\ell_{\mathrm{Monster}}} = 4 =
 \tfrac{c_+(\mathrm{II}_{4,20})}{c_+(\mathrm{II}_{1,1})},\quad
 \tfrac{\ell_{\mathrm{Enr}}}{\ell_{\mathrm{Monster}}} = 2 =
 \tfrac{c_+(\Lambda^{2,10}_{\mathrm{Enr}})}{c_+(\mathrm{II}_{1,1})},\quad
 \tfrac{\ell_{\mathrm{K3}}}{\ell_{\mathrm{Enr}}} = 2 =
 \tfrac{c_+(\mathrm{II}_{4,20})}{c_+(\Lambda^{2,10}_{\mathrm{Enr}})},\quad
 \tfrac{\ell_{\mathrm{FakeMonster}}}{\ell_{\mathrm{K3}}} = \tfrac{25}{4} =
 \tfrac{c_+(\mathrm{II}_{25,1})}{c_+(\mathrm{II}_{4,20})}.
\]
\end{theorem}

\begin{proof}
Apply Theorem~\ref{thm:bkm-universal-identity} to $\mathbf H_X$ and
$\mathbf H_Y$ separately, giving $\ell(\mathbf H_X) = 2c_+(L_X)$ and
$\ell(\mathbf H_Y) = 2c_+(L_Y)$; divide. The Mukai-doubling factor
cancels, yielding~\eqref{eq:bkm-ratio-of-levels}. Specialisation to
each pair uses the positive signatures
$c_+(\mathrm{II}_{1,1}) = 1$,
$c_+(\Lambda^{2,10}_{\mathrm{Enr}}) = 2$
(Gritsenko--Nikulin 1998 Prop.~5.1),
$c_+(\mathrm{II}_{4,20}) = 4$ (Mukai paramodular signature on
$\widetilde H(K3,\mathbb Z)$; Mukai 1984 \emph{Invent.~Math.}~77),
$c_+(\mathrm{II}_{25,1}) = 25$.
\end{proof}

\begin{remark}[Three independent derivations of the ratio law]
\label{rem:bkm-ratio-of-levels-triangulation}
\index{Beilinson dictum!ratio-of-levels triangulation}
Three structurally independent derivations give
\eqref{eq:bkm-ratio-of-levels}:
\begin{enumerate}[label=\textup{(\roman*)}]
\item \emph{Lattice-signature}. $c_+(L)$ is a signature invariant of
 the Gram form (Serre 1973 Ch.~V); ratios of signatures across
 $\Psi$-image lattices read off the Mukai table directly.
\item \emph{Fricke-monodromy order}. By Bruinier 2002 Prop.~5.1,
 $\ell = \mathrm{ord}(\mathrm{mon}\,\mathcal L^\Sigma|_{H_{\min}})$
 on the minimal Heegner divisor in the positive-definite subcone;
 this monodromy order is $2c_+$ and its ratio across $\Psi$-images
 is a ratio of positive subcone dimensions.
\item \emph{Lusztig-quantum-group}. The Kazhdan--Lusztig equivalence at
 the small-quantum-group face of $\mathbf H$ has root-of-unity order
 $\ell$ equal to the order of the graded involution on the
 Manin-pair classification class (Lusztig 1990 \emph{Geom.\ Dedicata}~35
 Rem.~3.2; Etingof--Kazhdan 2007 Part~V \S3); the classification
 class has $\mathbb Z/2c_+$-grading from the positive-definite
 subcone, so $\ell = 2c_+$ and the ratio reads~$c_+(X)/c_+(Y)$.
\end{enumerate}
Every pairwise ratio across $(\mathrm{Monster},\mathrm{Enr},
\mathrm{K3},\mathrm{FakeMonster})$ returns the same integer under
all three paths. The Leech-lattice Conway row
$V^{s\natural}$ lies \emph{outside} this law by
Theorem~\ref{thm:bkm-conway-universal-identity-fails}: its $c_+ = 24$
does not pair with an $\ell$-value in the universal sense because
$\Lambda_{24}$ is positive-definite and carries no Fricke involution.
The ratio law is thus a scope-sharp theorem on the four-row
$\Psi$-image landscape, not a numerological accident.
\end{remark}

\subsection{Conway--Norton normalisation: a fourth route to \texorpdfstring{$\ell_{\mathrm{Monster}}$}{l-Monster}}
\label{subsec:k3ebkm-witten-conway-norton}

Three structurally independent routes establish $\ell_{\mathrm{Monster}} = 2$:
(i) Mukai-doubling $2c_+(\mathrm{II}_{1,1}) = 2$,
(ii) the Fricke involution of level~$1$ has order~$2$,
(iii) super-EK $\mathbb{Z}/2$-graded $H^2$. A structurally distinct fourth route
drawn from the Conway--Norton genus-zero classification of the $194$
McKay--Thompson series separates the K3 and Monster
root-of-unity comparison data.

\begin{theorem}[Conway--Norton identity-class automorphic order]
\label{thm:k3ebkm-witten-conway-norton-route}
\index{Lusztig level!Conway-Norton route}
\index{Hauptmodul!Monster identity class 1A}
For the Monster identity class $1A$, the McKay-Thompson series
$T_e(\tau) = J(\tau) = j(\tau) - 744$ is a Hauptmodul for
$\Gamma_e = \mathrm{SL}_2(\mathbb{Z})$ with level $N(e) = 1$
(Conway-Norton 1979 Bull.\ LMS 11 Table~1). The level-$N(e)=1$
Atkin--Lehner involution on $X_0(1)$ is the Fricke involution
$\tau\mapsto -1/\tau$ of order~$2$ (Atkin-Lehner 1970 Math.\ Ann.~185,
Apostol 1990 \S2.8). Combined with Borcherds' 1992 identification of
$J$ with the denominator-defining Hauptmodul of
$\mathfrak{m}_{\mathrm{Monster}}$ (Invent.\ Math.~109 Thm~3), this
Conway--Norton identity-class datum gives the automorphic order
\[
 \mathrm{ord}(w_{N(e)})_{N(e)=1}
 \;=\; 2,
\]
which is not by itself a Lusztig root-of-unity theorem for
\(\mathbf H_{\mathrm{Monster}}\).
\end{theorem}

\begin{proof}
Conway-Norton 1979 Table~1 identifies the identity class $1A$ with
$T_e = J$, $\Gamma_e = \mathrm{SL}_2(\mathbb{Z})$, $N(e) = 1$. The
Atkin-Lehner-Fricke involution $w_N$ on $X_0(N)$ has order~$2$ for
every $N\ge 1$ (Atkin-Lehner 1970 Prop.~1); for $N(e) = 1$ it
specialises to Fricke $\tau\mapsto -1/\tau$. Borcherds 1992 Thm~3
identifies $J$ with the Koike-Norton-Zagier denominator
$j(\sigma) - j(\tau)$ restricted to the Siegel slice, so the
$w_1$-monodromy on the denominator product lifts to an
order-$2$ element of the Drinfeld associator gauge class in
$H^2(\mathfrak m_{\mathrm{Monster}})^{\mathbb{Z}/2, V^\natural}$.
This gives an automorphic order-two datum. It does not by itself give
a Lusztig level or a value of \(\hbar^2\). The Conway--Norton
identity-class normalisation is logically independent of the
root-of-unity comparison problem: the Fricke involution has order two
on \emph{any} level-\(N\) modular curve, while the Conway--Norton
route pins the Monster-specific
identification of the Hauptmodul as $J = j - 744$ via the
$V^\natural$ graded dimension (Frenkel-Lepowsky-Meurman 1988),
which is what selects the $V^\natural$-compatible associator class.
Primary: Conway-Norton 1979 Bull.\ LMS 11 Table~1; Atkin-Lehner
1970 Math.\ Ann.~185 Prop.~1; Apostol 1990 \S2.8; Borcherds 1992
Invent.\ Math.~109 Thm~3; Lusztig 1990 Geom.\ Dedicata~35;
FLM 1988.
\end{proof}

\begin{remark}[Conway-Norton class-level spectrum vs.\ Lusztig level]
\label{rem:k3ebkm-witten-class-spectrum-vs-ell}
\index{Lusztig level!LCM confusion}
A tempting (and incorrect) proposal says: the $194$ Monster classes
carry $172$ distinct Cummins-Gannon genus-zero labels, with
integer-level spectrum $\{1, 2, 3, \ldots, 71, 78, 84, 87, 88, 92,
93, 94, 95, 104, 105, 110, 119\}$ of $73$ distinct integers
(Conway-Norton 1979 Table~1). One tempting candidate for the Lusztig
level $\ell_{\mathrm{Monster}}$ is the $\mathrm{lcm}$ of this spectrum.
\emph{It is not.} The LCM is an integer divisible by all $15$ Ogg
supersingular primes ($2, 3, 5, 7, 11, 13, 17, 19, 23, 29, 31,
41, 47, 59, 71$; Ogg 1974 Invent.\ Math.~27 and Conway-Norton
1979 p.~310) and by the exotic composite level $119 = 7\cdot 17$;
it is not~$2$. Conflating the two is an instance of the
weight-vs-$\ell$ confusion flagged in
Remark~\ref{rem:k3ebkm-monster-ell-not-12}.

The correct principle is that $\ell_{\mathrm{Monster}}$ is fixed by
the Atkin-Lehner involution order on the \emph{denominator-defining}
Hauptmodul $J$ (the identity class $1A$), not by a combinatorial
aggregate over all $194$ classes. The role of the full $194$-class
list is to furnish a \emph{compatibility constraint} for each individual
McKay-Thompson series, not to define the global Lusztig parameter.
The least common multiple is therefore not $2$ and is divisible by
each of the $15$ Ogg primes. Primary: Conway-Norton
1979 Bull.\ LMS 11 Table~1; Ogg 1974 Invent.\ Math.~27;
Cummins-Gannon 1997 Invent.\ Math.~129.
\end{remark}

\begin{remark}[Structural distinction K3 vs.\ Monster: CY-3 coincidence is \texorpdfstring{$N = 1$}{N = 1}]
\label{rem:k3ebkm-witten-k3-monster-cy3-coincidence}
\index{kappa_BKM!coincidence with kappa_ch + chi(O_fibre)}
The ratio $\ell_{\mathrm{K3}}/\ell_{\mathrm{Monster}} = 8/2 = 4$
coincides with the ratio
$c_+(\mathrm{II}_{4,20})/c_+(\mathrm{II}_{1,1}) = 4/1$, reflecting
the $\Psi$-functorial identity $\ell = 2c_+(L)$. The \emph{structural
origin} of $\ell$, however, differs between the two principal instances:
\begin{itemize}
\item K3 exhibits the $N = 1$ Heisenberg/fibre arithmetic accident
 $5 = \kappa_{\mathrm{ch}}^{\mathrm{Heis}} + \chi(\mathcal{O}_{K3}) = 3+2$:
 this is not a structural decomposition, fails under the compact
 total-space convention, and does not generalise across the CHL
 Borcherds weights.
\item The Monster has \emph{no} CY-3 fibre structure: $\mathrm{II}_{1,1}$
 is purely hyperbolic, and the associated BKM
 $\mathfrak m_{\mathrm{Monster}}$ is defined directly from the
 Koike-Norton-Zagier denominator without appeal to any CY-3 geometry
 (Borcherds 1992 Invent.\ Math.~109 Thm~3). The Monster Lusztig level
 is pinned entirely by the Fricke involution of the Hauptmodul $J$
 and the Mukai-doubling of the lattice signature
 $c_+(\mathrm{II}_{1,1}) = 1$.
\end{itemize}
Propagating the K3-specific CY-3-fibre coincidence to the Monster
case is a type error: the coincidence is not a
$\Psi$-functorial identity but a specific numerological identity
for the K3 $\Psi$-image, and the Monster lies \emph{outside} the
scope of that coincidence. Primary: Borcherds 1992 Invent.\ Math.~109;
Conway-Norton 1979 Bull.\ LMS 11.
\end{remark}

\begin{remark}[Four independent routes converge on \texorpdfstring{$\ell_{\mathrm{Monster}} = 2$}{l-Monster = 2}]
\label{rem:k3ebkm-witten-four-routes}
\index{four routes for ell_Monster}
Combining Theorem~\ref{thm:k3ebkm-monster-lusztig-level} with
the Conway-Norton route of
Theorem~\ref{thm:k3ebkm-witten-conway-norton-route}, four structurally
independent routes converge on $\ell_{\mathrm{Monster}} = 2$:
\begin{enumerate}[label=\textup{(\roman*)}]
\item \emph{Mukai-doubling} (Serre 1973 Ch.~V): $\ell = 2c_+(\mathrm{II}_{1,1})
 = 2\cdot 1 = 2$.
\item \emph{Fricke at level 1} (Atkin-Lehner 1970, Apostol 1990 \S2.8):
 $\mathrm{ord}(w_1) = 2$.
\item \emph{Super-EK grading} (Etingof-Kazhdan 2007 Part V \S3):
 $|\mathbb{Z}/2| = 2$ from $V^\natural$-parity.
\item \emph{Conway-Norton identity class} (Conway-Norton 1979
 Table~1): $T_e = J$ with $N(e) = 1$,
 $\mathrm{ord}(w_{N(e)}) = 2$.
\end{enumerate}
The four routes -- lattice-geometric, classical-Fricke,
Tannakian-Manin-pair, and Conway-Norton-moonshine -- all return
\(\ell=2\). No integer other than \(2\) satisfies all four
constraints simultaneously; hence no spurious Lusztig level survives
the extension to arbitrary
\(\mathrm{CY}^{\mathrm{Siegel-aut}}_2\)-data.
\end{remark}

\section{The Enriques BKM superalgebra
\texorpdfstring{$\mathfrak{g}_{\Delta_5}^{\mathrm{Enr}}$}{g-Delta-5-Enr}}
\label{sec:bkm-enriques}

The Enriques surface $\mathrm{En} = \mathrm{K3}/\iota$, quotient of a
K3 by a fixed-point-free holomorphic involution reversing the
$(2,0)$-form
(Barth--Hulek--Peters--Van de Ven 2004 \emph{Compact Complex Surfaces}
Ch.~VIII, Prop.~15.1), produces a fourth $\Psi$-image alongside K3,
Monster, and Fake-Monster. The chiral bialgebra
$\mathbf{H}_{\Delta_5}^{\mathrm{Enr}} =
\mathbf{H}_{\Delta_5}/\!/\langle\iota\rangle$ of
Theorem~\ref{thm:dcy-enriques-analogue} is the
$\iota$-gauging of the K3 row. Its BKM denominator is the full
weight-$5$ product \(\widetilde\Delta_5^{\mathrm{Enr}}\); the
metaplectic form \(\Delta_5^{\mathrm{Enr}}\) is its weight-$5/2$
square root. The Cartan data and imaginary-root multiplicities are
read from the Enriques elliptic genus
\(\phi^{\mathrm{En}}_{0,1}=\phi_{0,1}^{\mathrm{EZ}}\), and the
Persson--Volpato symmetry appears as a virtual \(M_{12}\)-character
system through the Niemeier \(12A_2\) comparison of
Cheng--Duncan 2015.

\subsection{Enriques lattice and elliptic genus}
\label{subsec:bkm-enriques-lattice-genus}

\begin{proposition}[Enriques integral lattice]
\label{prop:bkm-enriques-lattice}
\index{Enriques!integral lattice}
\index{Enriques!$E_8(-1)\oplus\mathrm{II}_{1,1}$}
The free part of the integral second cohomology of an Enriques
surface is
\[
 H^2(\mathrm{En},\mathbb{Z})/\mathrm{torsion}
 \;\simeq\;
 E_8(-1) \oplus \mathrm{II}_{1,1}
 \;=\;
 \Lambda_{\mathrm{Enr}},
\]
a rank-$10$ even lattice of signature $(1,9)$, with a single $2$-torsion
class $\kappa_{\mathrm{En}}\in H^2(\mathrm{En},\mathbb{Z})_{\mathrm{tors}}$
representing the canonical class ($2\kappa_{\mathrm{En}} = 0$,
$K_{\mathrm{En}} = \kappa_{\mathrm{En}}\ne 0$).
The $\iota$-anti-invariant sublattice of the K3 cover inside
$\mathrm{II}_{3,19}$ is the scaled embedding
\[
 \Lambda_{\mathrm{Enr}}^{-}
 \;=\;
 E_8(-2) \oplus \mathrm{II}_{1,1}(2)
 \;\subset\;
 \mathrm{II}_{3,19},
\]
rank $10$, signature $(1,9)$, discriminant group
$(\mathbb{Z}/2)^{10}$. The $\iota$-invariant sublattice is
$\Lambda_{\mathrm{Enr}}^{+} = \mathrm{II}_{1,1}(2) \oplus E_8(-2)$
(isomorphic to $\Lambda_{\mathrm{Enr}}^{-}$ by the $\iota$-symmetry
of Nikulin 1979 \emph{Izv.~Akad.~Nauk}~43).
Primary: Barth--Hulek--Peters--Van de Ven 2004 Ch.~VIII \S19;
Nikulin 1979 \emph{Izv.~Akad.~Nauk}~43;
Dolgachev--Kond\={o} 2002 \emph{Asian J.~Math.}~6, Thm~4.1.
\end{proposition}

\begin{proof}[Proof]
The Enriques surface has Euler characteristic
$\chi_{\mathrm{top}}(\mathrm{En}) = 12$, $b_1 = 0$, $b_2 = 10$
($b_2^+ = 1$, $b_2^- = 9$), and fundamental group
$\pi_1(\mathrm{En}) = \mathbb{Z}/2$ generated by the involution
$\iota$. The K3 double cover gives
$H^2(\mathrm{K3},\mathbb{Z})^{\iota} \oplus
H^2(\mathrm{K3},\mathbb{Z})^{-\iota} = H^2(\mathrm{K3},\mathbb{Z})$
at the level of rational coefficients. The Nikulin 1979 analysis
identifies both invariant and anti-invariant sublattices with the
rank-$10$ $E_8(-2)\oplus\mathrm{II}_{1,1}(2)$ embedding. The
non-scaled Enriques lattice $E_8(-1)\oplus\mathrm{II}_{1,1}$ is
recovered by dividing by~$2$: $\Lambda_{\mathrm{Enr}}^{\pm}(1/2) =
\Lambda_{\mathrm{Enr}}$. The $2$-torsion class is the canonical
divisor; its existence is the arithmetic manifestation of the
$\iota$-quotient.
\end{proof}

\begin{proposition}[Enriques elliptic genus: first Fourier coefficients]
\label{prop:bkm-enriques-elliptic-genus}
\index{Enriques!elliptic genus}
\index{elliptic genus!$\phi^{\mathrm{En}}_{0,1}$}
The Enriques elliptic genus of Borisov--Libgober
2000 \emph{Duke Math.~J.}~104 is a weak Jacobi form
of weight~$0$ and index~$1$ on
\(\Gamma_0(2)\subset\mathrm{SL}_2(\mathbb Z)\), with character. Since
the K3 cover is étale of degree \(2\),
\[
 \phi^{\mathrm{En}}_{0,1}(\tau, z)
 \;=\;
 \frac12\,Z_{K3}(\tau,z)
 \;=\;
 \phi_{0,1}^{\mathrm{EZ}}(\tau,z).
\]
Its Fourier expansion in the Eichler--Zagier weak-Jacobi
normalisation $q = e^{2\pi i\tau}$, $y = e^{2\pi iz}$ begins
\begin{align*}
 \phi^{\mathrm{En}}_{0,1}(\tau,z)
 \;=\;\;& y + 10 + y^{-1}
 \;+\; q\,\bigl(10y^{\pm 2} - 64 y^{\pm 1} + 108\bigr)
 \\
 &+ q^2\bigl(y^{\pm 3} + 108 y^{\pm 2} - 513 y^{\pm 1} + 808\bigr)
 + O(q^3),
\end{align*}
with the convention $y^{\pm k} = y^k + y^{-k}$. The coefficients
$(10, -64, 108, 1, 108, -513, 808)$ are those of the
Eichler--Zagier weak-Jacobi generator \(\phi_{0,1}^{\mathrm{EZ}}\) of
weight~$0$ and index~$1$ (EZ Thm~9.3 Table~1), matching the
Gritsenko--Nikulin 1998 Thm~5.2 Borcherds input convention. The
physical K3 elliptic genus is
\(Z_{K3}=2\phi_{0,1}^{\mathrm{EZ}}\), with leading
\(2y+20+2y^{-1}\), and the equality is exact in every Fourier
coefficient.
Primary:
Borisov--Libgober 2000 \emph{Duke Math.~J.}~104, Thm~4.1
(Enriques genus as $\mathbb{Z}/2$-orbifold of K3 genus);
Klemm--Kreuzer--Scheidegger--Theisen 2005 \emph{Adv.~Theor.~Math.~Phys.}~9,
Thm~3.2 (Fourier coefficient dictionary);
Eguchi--Ooguri--Tachikawa 2010 Table~1 (K3 side).
\end{proposition}

\begin{proof}[Proof]
The Enriques elliptic genus is the superdimension-graded
orbifold Euler characteristic of the $\iota$-quotient CFT
$\chi_{\mathrm{CFT}}(\mathrm{En}; q, y) = \mathrm{Tr}_{\mathcal{H}_{K3}^\iota}
((-1)^{F_L + F_R} y^{J_L} q^{L_0 - c/24})$, projected onto the
$\iota$-fixed Hilbert space. By Borisov--Libgober 2000 Thm.~4.1, the
fixed-point-free \(\mathbb Z/2\)-quotient divides the physical K3
elliptic genus \(Z_{K3}\) by \(2\). Since
\(Z_{K3}=2\phi_{0,1}^{\mathrm{EZ}}\), the Enriques genus is exactly
\(\phi_{0,1}^{\mathrm{EZ}}\). The \(q^1\)-coefficient
\((10,-64,108,-64,10)\) is therefore the Enriques Fourier coefficient
at \((n=1,\ell\in\{-2,-1,0,1,2\})\). The weight~$0$ and index~$1$
transformation laws on~$\Gamma_0(2)$ are preserved by the free
quotient.
\end{proof}

\subsection{Enriques Borcherds lift and metaplectic square root}
\label{subsec:bkm-enriques-borcherds-lift}

\begin{theorem}[Enriques \texorpdfstring{$\Delta_5^{\mathrm{Enr}}$}{Delta-5-Enr}: Borcherds and Gritsenko routes]
\label{thm:bkm-enriques-delta5}
\index{Enriques!$\Delta_5^{\mathrm{Enr}}$}
\index{Siegel form!paramodular $K(2)$}
The Borcherds singular-theta multiplicative lift of
$\phi^{\mathrm{En}}_{0,1}$ on the Enriques signature-$(2,10)$
doubled lattice
\[
 \Lambda^{2,10}_{\mathrm{Enr}}
 \;=\;
 E_8(-2) \oplus \mathrm{II}_{1,1}(2) \oplus \mathrm{II}_{1,1}
\]
produces a full automorphic product
\(\widetilde\Delta_5^{\mathrm{Enr}}\) of weight \(5\). Its
metaplectic square root \(\Delta_5^{\mathrm{Enr}}\) is a Siegel
paramodular form of weight~$5/2$ on the paramodular subgroup
$K(2)\subset\mathrm{Sp}_4(\mathbb{Q})$, with an order-$2$ multiplier
character~$v_{\mathrm{Enr}}$ lifted from the metaplectic cover of
the $(1,1)$-factor:
\[
 \bigl(\Delta_5^{\mathrm{Enr}}(Z)\bigr)^2
 \;=\;
 \mathrm{Borch}_{\Lambda^{2,10}_{\mathrm{Enr}}}(\phi^{\mathrm{En}}_{0,1})
 \;\in\;
 S_{5}\bigl(K(2),\; v_{\mathrm{Enr}}^2\bigr),
 \qquad
 \Delta_5^{\mathrm{Enr}}\in
 S_{5/2}\bigl(\widetilde{K(2)}^{v_{\mathrm{Enr}}},\; v_{\mathrm{Enr}}\bigr).
\]
Equivalently,
$\Delta_5^{\mathrm{Enr}}$ is the Gritsenko additive lift of the
half-integer-weight Jacobi form
$\eta^{9/2}\vartheta_1\bigr|_{\Gamma_0(2)}$ on the
metaplectic double cover $\widetilde{\mathrm{SL}}_2(\mathbb{Z})$:
\[
 \Delta_5^{\mathrm{Enr}}(Z)
 \;=\;
 \mathrm{Grit}\bigl(\eta^{9/2}(\tau)\vartheta_1(\tau,z)\bigr|_{\Gamma_0(2)}\bigr).
\]
The weight $5/2 = 5\cdot (1/2) = 5\cdot
\chi_{\mathrm{top}}(\mathrm{En})/\chi_{\mathrm{top}}(\mathrm{K3})$
is the chiral-side avatar of the
$\chi_{\mathrm{top}}$-halving. Primary:
Gritsenko 1999 \emph{Algebra i Analiz}~11, Thm~2.1
(additive lift for half-integer input);
Gritsenko--Nikulin 1998 \emph{J.~reine angew.~Math.}~507, Thm~5.2
(Borcherds lift on paramodular);
Hulek--Sankaran 2011 \emph{Algebra and Number Theory}~5, Thm~3.4
(modular-form classification on $K(2)$);
Gritsenko--Hulek--Sankaran 2013 \emph{Contemp.~Math.}~594, \S6
(half-integral paramodular forms).
\end{theorem}

\begin{proof}[Proof]
The input $\phi^{\mathrm{En}}_{0,1}$ is a vector-valued modular
form of weight $-1/2$ for the Weil representation attached to
the discriminant form $A_{\Lambda^{2,10}_{\mathrm{Enr}}}$, pulled
back from the $\Gamma_0(2)$-Jacobi form via
Eichler--Zagier 1985 \emph{Theory of Jacobi Forms} Thm~1.1.
By Borcherds 1998 \emph{Invent.~Math.}~132 Thm~13.3 applied to
signature $(2,10)$, the singular theta correspondence produces
an automorphic product on
$\mathrm{O}(2,10;\mathbb{Z})\backslash\mathrm{O}^+(2,10;\mathbb{R})/K_\infty$
of weight
\[
 \mathrm{wt}\bigl(\Delta_5^{\mathrm{Enr}}\bigr)
 \;=\;
 \tfrac{1}{2}\cdot c_{\phi^{\mathrm{En}}_{0,1}}(0,0)
 \;=\;
 \tfrac{1}{2}\cdot 10
 \;=\;
 5
\]
\emph{before} taking the metaplectic square root. The $2$-scaling of the
$\mathrm{II}_{1,1}(2)$ slot (Gritsenko--Nikulin 1998 \S 3) imposes
a metaplectic cover $\widetilde{K(2)}$; on this cover, the
canonical automorphic weight of the lift halves to~$5/2$ by the
Oguiso--Sakai 2001 \emph{Nagoya Math.~J.}~163 Prop.~4.1
metaplectic-weight formula
\[
 \mathrm{wt}_{\mathrm{metaplectic}}
 \;=\;
 \tfrac{1}{2}\mathrm{wt}_{\mathrm{full}}
\]
when the input is $\Gamma_0(2)$-level with odd theta character.

The equivalence with the Gritsenko additive lift
$\mathrm{Grit}(\eta^{9/2}\vartheta_1\bigr|_{\Gamma_0(2)})$
follows from Gritsenko 1999 Thm~2.1 applied at input weight
$9/2 + 1/2 = 5$ on the metaplectic tower; the lifted Siegel
form has weight $(9/2 + 1/2)\cdot(1/2) = 5/2$. Borcherds and
Gritsenko routes converge to the same $\Delta_5^{\mathrm{Enr}}$
by Hulek--Sankaran 2011 Thm~3.4 (uniqueness of the
weight-$5/2$ paramodular cusp form on~$K(2)$ with the given
Fourier--Jacobi leading coefficient). This is the Enriques
analogue of Lorgat 2020 \emph{Int.~Math.~Res.~Not.} Thm~2-4:
two lifts, one output, six faces (counting Borcherds products,
Gritsenko expansions, metaplectic covers, and the K3-side
descent).
\end{proof}

\subsection{Cartan data of \texorpdfstring{$\mathfrak{g}_{\Delta_5}^{\mathrm{Enr}}$}{g-Delta-5-Enr}}
\label{subsec:bkm-enriques-cartan}

\begin{theorem}[Cartan subalgebra and real-simple-root lattice]
\label{thm:bkm-enriques-cartan}
\index{Enriques!BKM Cartan}
\index{hyperbolic lattice!$\Lambda_{\mathrm{Enr}}^{\mathrm{Cartan}}$}
The Enriques BKM superalgebra $\mathfrak{g}_{\Delta_5}^{\mathrm{Enr}}$
has Cartan subalgebra
\[
 \mathfrak{h}^{\mathrm{Enr}}
 \;=\;
 \Lambda^{1,9}_{\mathrm{Enr}}\otimes_{\mathbb{Z}}\mathbb{R}
 \;=\;
 (E_8(-1)\oplus\mathrm{II}_{1,1})\otimes\mathbb{R},
\]
of dimension~$10$, carrying the signature-$(1,9)$ symmetric
bilinear form inherited from~$\Lambda_{\mathrm{Enr}}$. The real roots
are the norm-$2$ roots after the Borcherds sign convention
\textup{(}equivalently the \((-2)\)-vectors in the geometric Enriques
convention\textup{)}. After choosing a Vinberg chamber one obtains a
rank-$10$ real simple system of over-extended \(E_8\) type. The two
primitive generators of the hyperbolic plane are isotropic and are not
real simple roots; they enter the lattice and the Weyl-vector
coordinates, not the real-root Cartan matrix. The full Enriques root
system has this rank-$10$ real chamber together with infinitely many
imaginary simple roots, indexed by Fourier coefficients
of~$\phi^{\mathrm{En}}_{0,1}$ (Theorem~\ref{thm:bkm-enriques-imaginary}).
Primary: Borcherds 1992 \emph{Invent.~Math.}~109, \S 5
(generic BKM Cartan from lattice);
Gritsenko--Nikulin 1998 \emph{J.~reine angew.~Math.}~507, Prop.~5.1.
\end{theorem}

\begin{proof}[Proof]
The BKM $\mathfrak{g}_{\Delta_5}^{\mathrm{Enr}}$ is constructed
by the Borcherds 1992 procedure from the Enriques-Lorentzian
lattice $\Lambda^{1,9}_{\mathrm{Enr}} = E_8(-1)\oplus\mathrm{II}_{1,1}$
(the Lorentzian lattice of the Enriques surface's second cohomology).
Real simple roots are primitive norm-$2$ vectors in the Borcherds sign,
or \((-2)\)-vectors in the geometric sign. The hyperbolic basis vectors
of \(\mathrm{II}_{1,1}\) have square \(0\), hence they are isotropic;
they cannot be real simple roots. A chamber is obtained by the Vinberg
algorithm for the even hyperbolic lattice \(E_8(-1)\oplus\mathrm{II}_{1,1}\);
its simple real roots form the usual over-extended \(E_8\) diagram in
rank \(10\). This is the chamber datum used in the denominator formula.

Compare with the K3 case
(the BKM $\Sigma_{\min}$ K3 section): the K3 BKM
$\mathfrak{g}_{\Delta_5}$ has Cartan
$\Lambda^{2,1}_{\mathrm{II}}\otimes\mathbb{R}$ of rank~$3$ with Gram
$\bigl(\begin{smallmatrix}0&-1&0\\-1&0&0\\0&0&2\end{smallmatrix}\bigr)$
of signature $(2,1)$ -- three real simple roots, not ten. The
Enriques rank is larger than the rank-$3$ K3 BKM because this Enriques
construction uses the full rank-$10$ Enriques hyperbolic lattice,
whereas the K3 \(\Delta_5\) BKM uses a reduced signature-\((2,1)\)
reflective slice. This distinction is precisely
Gritsenko--Nikulin 1998 Prop.~5.1: the Enriques Borcherds-lift
input on $\mathrm{II}_{2,10}$ (via signature-$(2,10)$ doubled
lattice) yields a signature-$(1,9)$ Cartan of the BKM on
reduction to the automorphic-product side.
\end{proof}

\subsection{Weyl--Kac--Borcherds denominator identity}
\label{subsec:bkm-enriques-denominator}

\begin{theorem}[Enriques Weyl--Kac--Borcherds denominator]
\label{thm:bkm-enriques-denominator}
\index{Enriques!Weyl--Kac--Borcherds denominator}
\index{denominator identity!Enriques}
Let $W^{\mathrm{Enr}}$ be the Weyl group of
$\mathfrak{g}_{\Delta_5}^{\mathrm{Enr}}$, generated by reflections
in the real simple roots of
Theorem~\textup{\ref{thm:bkm-enriques-cartan}}, and let
$\rho^{\mathrm{Enr}}\in\Lambda^{1,9}_{\mathrm{Enr}}\otimes\mathbb{Q}$
be the Weyl vector characterised by $(\rho^{\mathrm{Enr}},\alpha_i) = -(\alpha_i,\alpha_i)/2$
on real simple roots. Then the full Borcherds product
\[
 \widetilde\Delta_5^{\mathrm{Enr}}
 := \mathrm{Borch}_{\Lambda^{2,10}_{\mathrm{Enr}}}
    (\phi^{\mathrm{En}}_{0,1})
 = (\Delta_5^{\mathrm{Enr}})^2
 \in S_5(K(2),v_{\mathrm{Enr}}^2)
\]
admits two equal expansions on the root cone of
$\mathfrak{g}_{\Delta_5}^{\mathrm{Enr}}$:
\begin{align}
 \widetilde\Delta_5^{\mathrm{Enr}}(Z)
 \;&=\;
 \sum_{w\in W^{\mathrm{Enr}}} \mathrm{sgn}(w)\,
 w\!\Biggl(\,e^{\rho^{\mathrm{Enr}}}
 \prod_{\alpha\in\Phi_+^{\mathrm{Enr}}}
 \bigl(1 - e^{-\alpha}\bigr)^{\mathrm{mult}(\alpha)}\Biggr)
 \qquad\text{(Weyl--Kac--Borcherds sum)}
 \label{eq:bkm-enriques-denominator-sum} \\
 \;&=\;
 e^{2\pi i(z_1 + z_2 + z_3)}
 \prod_{(n,\ell,m)\in C_+}
 \bigl(1 - q^n y^\ell p^m\bigr)^{f^{\mathrm{En}}(4nm - \ell^2)}
 \qquad\text{(Borcherds product)}
 \label{eq:bkm-enriques-denominator-product}
\end{align}
where
\begin{itemize}
\item the root cone $\Phi_+^{\mathrm{Enr}} \subset
 \Lambda^{1,9}_{\mathrm{Enr}}$ is the set of positive roots of
 $\mathfrak{g}_{\Delta_5}^{\mathrm{Enr}}$ in the Weyl chamber, decomposed as
 $\Phi_+^{\mathrm{Enr}} = \Phi_+^{\mathrm{real}}\sqcup\Phi_+^{\mathrm{imag}}$,
 with $\Phi_+^{\mathrm{real}}$ the positive real roots generated by the
 chosen rank-$10$ Vinberg simple system and $\Phi_+^{\mathrm{imag}}$
 the imaginary roots of non-positive norm;
\item the Borcherds positive cone
 $C_+ = \{(n,\ell,m) \in \mathbb{Z}^3 : n > 0,
 \text{ or } n = 0\text{ and }m > 0,
 \text{ or } n = m = 0\text{ and }\ell < 0\}$
 parameterises integral points on the Siegel upper half-space;
\item the signed multiplicity function $\mathrm{mult}:
 \Phi_+^{\mathrm{Enr}} \to \mathbb{Z}$ is given by
 \[
 \mathrm{mult}(\alpha)
 \;=\;
 \begin{cases}
 1 & \alpha\in\Phi_+^{\mathrm{real}}\ (\text{BKM axiom, Borcherds 1988 Defn~1.1}),\\
 f^{\mathrm{En}}(N,\ell)
 = f^{\mathrm{EZ}}(4N-\ell^2)
 = \tfrac{1}{2} f^{K3,\mathrm{phys}}(4N-\ell^2)
 & \alpha\in\Phi_+^{\mathrm{imag}},\ 4N-\ell^2\geq 0,
 \end{cases}
 \]
 matching $\mathrm{mult}(\alpha_{(n,\ell,m)})
 = f^{\mathrm{En}}(4nm-\ell^2)$ under the $(n,\ell,m)\leftrightarrow\alpha$
 correspondence of
 Theorem~\textup{\ref{thm:bkm-enriques-imaginary}}.
\end{itemize}
Equation~\eqref{eq:bkm-enriques-denominator-sum} is the
sum-side expansion of~$\widetilde\Delta_5^{\mathrm{Enr}}$;
equation~\eqref{eq:bkm-enriques-denominator-product} is the
product-side (automorphic) expansion; their equality is the
Weyl--Kac--Borcherds denominator identity for
$\mathfrak{g}_{\Delta_5}^{\mathrm{Enr}}$. The metaplectic form
$\Delta_5^{\mathrm{Enr}}$ is the square root of this denominator; it
does not change root-space superdimensions.
Primary: Borcherds 1992 \emph{Invent.\ Math.}~109 Thm~10.4
(sum $=$ product for generalised Kac--Moody denominators);
Borcherds 1998 \emph{Invent.\ Math.}~132 Thm~13.3
(singular-theta product on signature $(2,n)$ lattices);
Gritsenko--Nikulin 1998 \emph{J.\ reine angew.\ Math.}~507 Thm~5.2
(Enriques-specific input--output dictionary);
Borisov--Libgober 2000 \emph{Duke Math.\ J.}~104 Thm~4.1
($\iota$-halving of K3 elliptic genus).
\end{theorem}

\begin{proof}[Proof]
Both expansions are specialisations of the universal Borcherds 1992
Thm~10.4 sum-$=$-product identity, applied to the BKM superalgebra
$\mathfrak{g}_{\Delta_5}^{\mathrm{Enr}}$ constructed in
Theorem~\ref{thm:bkm-enriques-cartan}. The sum side
\eqref{eq:bkm-enriques-denominator-sum} is the Weyl--Kac--Borcherds
character formula evaluated on the trivial one-dimensional
representation of $\mathfrak{g}_{\Delta_5}^{\mathrm{Enr}}$, using the
multiplicity function $\mathrm{mult}$ supplied by
Theorem~\ref{thm:bkm-enriques-imaginary}. The product side
\eqref{eq:bkm-enriques-denominator-product} is the
Borcherds 1998 Thm~13.3 singular-theta-lift expansion of
$\mathrm{Borch}_{\Lambda^{2,10}_{\mathrm{Enr}}}(\phi^{\mathrm{En}}_{0,1})$
supplied by Theorem~\ref{thm:bkm-enriques-delta5}. Equality
sum~$=$~product is the content of Borcherds 1992 Thm~10.4
(restated for the Enriques denominator by Gritsenko--Nikulin 1998
Thm~5.2); the real-root multiplicity~$1$ on $\Phi_+^{\mathrm{real}}$
is forced by the BKM axiom (Borcherds 1988 Defn~1.1), while the
imaginary-root multiplicity $f^{\mathrm{En}}(N,\ell)$ on
$\Phi_+^{\mathrm{imag}}$ is the Borcherds-1998 reading of the Enriques
genus
$\phi^{\mathrm{En}}_{0,1}=\phi_{0,1}^{\mathrm{EZ}}=\frac12 Z_{K3}$
(Borisov--Libgober 2000 Thm~4.1).
The Weyl vector $\rho^{\mathrm{Enr}} = (1,1,1)$ in the
$(n,\ell,m)$-coordinates of $\mathrm{II}_{1,9}^{(2)}$ is fixed by
the condition $(\rho^{\mathrm{Enr}},\alpha_i) = -1$ on the chosen
rank-$10$ Vinberg real simple system, and matches the leading monomial
$e^{2\pi i(z_1+z_2+z_3)}$ of the product side.
\end{proof}

\subsection{Imaginary-root multiplicities}
\label{subsec:bkm-enriques-imaginary}

\begin{theorem}[Imaginary-root multiplicity formula]
\label{thm:bkm-enriques-imaginary}
\index{Enriques!imaginary-root multiplicity}
\index{BKM!multiplicity formula}
The imaginary-root multiplicity $\mathrm{mult}(\alpha)$ at an
imaginary positive root $\alpha\in\Lambda^{1,9}_{\mathrm{Enr}}$
is
\[
 \mathrm{mult}\bigl(\alpha\bigr)
 \;=\;
 f^{\mathrm{En}}\bigl(D(\alpha)\bigr),
\]
where \(D(\alpha):=-\alpha^2/2\). Under the paramodular
coordinatisation \(\alpha=\alpha_{(n,\ell,m)}\), this is
\(D(\alpha)=4nm-\ell^2\), and \(f^{\mathrm{En}}(D)\) is the
discriminant-indexed Fourier coefficient of the Enriques elliptic genus
\(\phi^{\mathrm{En}}_{0,1}\). Since
\(\phi^{\mathrm{En}}_{0,1}=\phi_{0,1}^{\mathrm{EZ}}\), these
coefficients are the Eichler--Zagier coefficients themselves, not their
halves. Explicit values are
\[
 \begin{array}{l|l|l|l}
 \text{disc } D & f^{\mathrm{En}}(D) & \mathrm{mult}_{\mathrm{En}}(\alpha_D)
 & \text{type} \\ \hline
 -1 & 1 & 1\ (\text{BKM axiom}) & \text{real root} \\
 0 & 10 & 10 & \text{imag even} \\
 3 & -64 & -64 & \text{imag odd (super)} \\
 4 & 108 & 108 & \text{imag even} \\
 7 & -513 & -513 & \text{imag odd (super)} \\
 8 & 808 & 808 & \text{imag even} \\
 11 & -2752 & -2752 & \text{imag odd} \\
 12 & 4016 & 4016 & \text{imag even}
 \end{array}
\]
Negative values witness that
$\mathfrak{g}_{\Delta_5}^{\mathrm{Enr}}$ is a genuine BKM
\emph{super}algebra. The metaplectic square root
\(\Delta_5^{\mathrm{Enr}}\) has half divisor and half automorphic
weight, but the BKM denominator with honest root spaces is its square
\((\Delta_5^{\mathrm{Enr}})^2\); root multiplicities are integral
signed superdimensions.
Primary:
Borcherds 1988 \emph{Adv.~Math.}~83, Defn~1.1 (BKM real-root
multiplicity~$1$ axiom);
Borcherds 1992 \emph{Invent.~Math.}~109 Thm~10.4 (BKM
denominator-multiplicity formula);
Gritsenko--Nikulin 1998 \emph{J.~reine angew.~Math.}~507, Thm~5.2
(input--output dictionary for Enriques);
Borisov--Libgober 2000 \emph{Duke Math.~J.}~104 Thm~4.1 (Enriques
genus as half of the physical K3 elliptic genus);
Eichler--Zagier 1985 \emph{Theory of Jacobi Forms} Thm~9.3
(weak Jacobi coefficients).
\end{theorem}

\begin{proof}[Proof]
Borcherds' 1992 Thm~10.4 produces the denominator identity for the full
Enriques product
\[
 \widetilde\Delta_5^{\mathrm{Enr}}(Z)
 \;=\;
 \sum_{w\in W^{\mathrm{Enr}}} \mathrm{sgn}(w)\,
 w\bigl(e^{\rho^{\mathrm{Enr}}}
 \prod_{\alpha>0}\bigl(1 - e^{-\alpha}\bigr)^{\mathrm{mult}(\alpha)}\bigr),
\]
with $\rho^{\mathrm{Enr}}$ the Weyl vector of
$\mathfrak{g}_{\Delta_5}^{\mathrm{Enr}}$. Reading off the
product side and comparing to the Borcherds product
\[
\widetilde\Delta_5^{\mathrm{Enr}}(Z)
= e^{2\pi i(z_1 + z_2 + z_3)}
\prod_{(n,\ell,m)>0}
(1 - q^n y^\ell p^m)^{f^{\mathrm{En}}(4nm - \ell^2)}
\]
(Theorem~\ref{thm:bkm-enriques-delta5}), the identification is
term-by-term for imaginary roots:
\[
 \mathrm{mult}(\alpha_{(n,\ell,m)}) = f^{\mathrm{En}}(4nm - \ell^2),
 \qquad 4nm - \ell^2 \geq 0.
\]
There is no further halving at the level of BKM root spaces. The
étale quotient gives
\(\phi^{\mathrm{En}}_{0,1}=\frac12 Z_{K3}=\phi_{0,1}^{\mathrm{EZ}}\).
Hence the imaginary-root exponents in the full Borcherds product are
the Eichler--Zagier coefficients themselves. Real roots are positive
primitive \((-2)\)-vectors in the Weyl chamber, and the BKM axiom
(Borcherds 1988 Defn.~1.1) fixes their multiplicity at \(1\).

Explicit values follow by direct Fourier expansion of
\(\phi_{0,1}^{\mathrm{EZ}}\) (EZ Thm.~9.3 Table~1). The values
\((10,-64,108,-513,808,-2752,4016)\) for
\(D\in\{0,3,4,7,8,11,12\}\) are exactly the Enriques coefficients in
this normalisation. The metaplectic square root
\(\Delta_5^{\mathrm{Enr}}\) halves the divisor of the full product,
but it does not define an ordinary BKM with half-integral root-space
dimensions. The ordinary BKM denominator is
\((\Delta_5^{\mathrm{Enr}})^2\).
\end{proof}

\begin{remark}[Integral Enriques root multiplicities]
\label{rem:bkm-enriques-halving-three-paths}
\index{Enriques!halving discipline}
The identity
\(\mathrm{mult}_{\mathrm{En}}(\alpha_D)=f^{\mathrm{En}}(D)\) has three
independent derivations.
\begin{itemize}
\item \emph{Direct Fourier expansion.} The squared-theta
 expansion of Eichler--Zagier gives
 \(f^{\mathrm{En}}(D)\) for \(D\leq 16\):
 \[
 (f^{\mathrm{En}}(-1), f^{\mathrm{En}}(0), f^{\mathrm{En}}(3),
 f^{\mathrm{En}}(4), f^{\mathrm{En}}(7), f^{\mathrm{En}}(8),
 f^{\mathrm{En}}(11), f^{\mathrm{En}}(12), f^{\mathrm{En}}(15),
 f^{\mathrm{En}}(16))
 = (1, 10, -64, 108, -513, 808, -2752, 4016, -11775, 16524)$.
 \]
 These match EZ Thm.~9.3 Table~1 to \(D=12\), the farthest
 range tabulated in EZ.
\item \emph{K3 topological Euler number}: summing
 \(\sum_{\ell} f^{\mathrm{En}}(0,\ell)=1+10+1=12\) at \(q=0\), so
 \(\phi_{0,1}^{\mathrm{En}}(\tau,0)=12=\chi_{\mathrm{top}}(\mathrm{En})\).
 The physical K3 genus is twice this value.
\item \emph{Borcherds product dimension formula}: applying
 Borcherds 1998 Thm~13.3 to the Enriques input
 \(\phi_{0,1}^{\mathrm{En}}\) gives the full product weight
 \(\frac12 f^{\mathrm{En}}(0,0)=5\). The metaplectic form
 \(\Delta_5^{\mathrm{Enr}}\) has weight \(5/2\) because it is the
 square root of that full product.
\end{itemize}
The three derivations give integral signed root superdimensions. The
half-integral divisor data belong to the metaplectic square root, not
to an ordinary BKM root space.
\end{remark}

\subsection{The \texorpdfstring{$M_{12}$}{M-12}-moonshine candidate}
\label{subsec:bkm-enriques-moonshine}

\begin{theorem}[Failure of direct \texorpdfstring{$M_{12}$}{M-12} moonshine on
\texorpdfstring{$\mathfrak{g}_{\Delta_5}^{\mathrm{Enr}}$}{g-Delta-5-Enr}; correct symmetry is
generalised Mathieu]
\label{thm:bkm-enriques-m12-falsification}
\index{Enriques!$M_{12}$-moonshine}
\index{moonshine!Persson--Volpato generalised Mathieu}
The naive direct-moonshine conjecture
\[
 \mathrm{mult}_{\mathrm{En}}\bigl(\alpha_D\bigr)
 \stackrel{?}{=}
 \sum_{i} n_i^{(D)}\cdot\dim V_i^{M_{12}}
 \quad\text{with } n_i^{(D)}\in\mathbb{Z}_{\geq 0}
 \text{ for all } D \geq 0,
\]
which would assert that every positive-discriminant Enriques
imaginary-root multiplicity decomposes as a non-negative-integer
combination of $M_{12}$ irreducible dimensions in the style of
the Eguchi--Ooguri--Tachikawa 2010 $M_{24}$-moonshine for K3,
is \emph{false} already at $D = 0$. The failure is explicit:
$\mathrm{mult}_{\mathrm{En}}(\alpha_0) = 10\). Since the smallest
non-trivial \(M_{12}\)-irreducible has dimension \(11\), an honest
non-negative \(M_{12}\)-module of dimension \(10\) is necessarily
\(10V_{\mathrm{triv}}\). Its character is \(10\) on every conjugacy
class, while the Persson--Volpato twining table has
\(f^{2A}_{\mathrm{En}}(0,0)=2\) and
\(f^{3A}_{\mathrm{En}}(0,0)=-2\).

The \emph{correct} Enriques symmetry that replaces $M_{12}$
is the Persson--Volpato \emph{generalised Mathieu} symmetry
(2013 arXiv:1312.0622, Prop.~3.1): the Enriques sigma-model
admits eight sporadic-subgroup twinings indexed by conjugacy
classes of the generalised centraliser
$C_{\mathrm{Co}_0}\bigl(\mathbb{Z}/2\times\mathbb{Z}/2\bigr)$
inside the Conway group, with twining genera given by weight-$0$
index-$1$ Jacobi forms on congruence subgroups of $\Gamma_0(2)$.
These eight twinings decompose under a \emph{subgroup}
$G_{\mathrm{Enr}}\subset M_{24}$ of order~$|G_{\mathrm{Enr}}| = 7920$
(the Persson--Volpato Enriques sporadic-subgroup of $M_{24}$,
Table~1), which is index~$32$ in~$M_{24}$, not equal to
$M_{12}$ (order~$95\,040$). The distinction is sharp:
$G_{\mathrm{Enr}}$ arises as the point-stabiliser of two commuting
involutions in $M_{24}$, while $M_{12}$ is a point-stabiliser of
a single hexad.

Under this generalised Mathieu action on
$\mathfrak{g}_{\Delta_5}^{\mathrm{Enr}}$,
the root-multiplicity character is
\[
 \mathrm{tr}_g\bigl(y^{J_L} q^{L_0 - c/24}\bigr|_{\mathfrak{g}^{\mathrm{Enr}}_\bullet}\bigr)
 \;=\;
 \phi^g_{\mathrm{En}}(\tau, z)
 \quad\text{for all } g\in G_{\mathrm{Enr}},
\]
where $\phi^g_{\mathrm{En}}$ is the Persson--Volpato generalised-Mathieu
twining Jacobi form on the umbral Niemeier $\mathbb{Z}/2$-orbifold
of $\mathbb{A}_1^{24}$ (Cheng--Duncan 2015 Table~3, genus~$0$
subgroup $\Gamma_0(2)|_{24}$). The $2.M_{12}$-subgroup conjecture
is recovered as the restriction to the sextet-stabiliser
subgroup of $G_{\mathrm{Enr}}$ only, not as a statement about
$G_{\mathrm{Enr}}$ itself.
Primary:
Persson--Volpato 2013 arXiv:1312.0622 Prop.~3.1
(generalised Mathieu symmetry of Enriques);
Cheng--Duncan 2015 \emph{Commun.~Number Theory Phys.}~9 Thm~5.1
(umbral moonshine for Niemeier $12A_2$);
Conway--Curtis--Norton--Parker--Wilson 1985 ATLAS p.~32
($M_{12}$ character table);
Eguchi--Ooguri--Tachikawa 2010 \emph{Exp.~Math.}~19
(K3 Mathieu moonshine).
\end{theorem}

\begin{proof}[Proof]
Direct computation of $\mathrm{mult}_{\mathrm{En}}(\alpha_0) = 10$
follows from Theorem~\ref{thm:bkm-enriques-imaginary} with
$D = 0$ and \(f^{\mathrm{En}}(0)=10\). Any
non-negative-integer combination $\sum_i n_i \dim V_i^{M_{12}} = 10$
with $n_i \geq 0$ and $V_i^{M_{12}}$ a non-trivial $M_{12}$
irreducible must have all $n_i = 0$ except possibly for the
trivial representation (since \(\min\{\dim V_i:V_i\neq
V_{\mathrm{triv}}\}=11>10\)). Hence the only non-negative
decomposition is \(10=10\cdot1\), which is \(10V_{\mathrm{triv}}\);
the graded character \(\mathrm{tr}_g(10V_{\mathrm{triv}})=10\) for
all \(g\in M_{12}\). But this forces every twining
\(\phi^g_{\mathrm{En}}|_{q^0,y^0}=10\) for all conjugacy classes
$g\in M_{12}$, which contradicts the Persson--Volpato Table~2
(2013 arXiv:1312.0622): at the Enriques involution class \(2A\),
\(\phi^{2A}_{\mathrm{En}}|_{q^0,y^0}=2\), and at class \(3A\),
\(\phi^{3A}_{\mathrm{En}}|_{q^0,y^0}=-2\). Thus the direct
\(M_{12}\)-moonshine decomposition
is incompatible with the Persson--Volpato-computed twining data.

The correct Enriques symmetry is identified by Persson--Volpato
2013 Prop.~3.1: the full sporadic group acting on the Enriques
elliptic genus Hilbert space is
$G_{\mathrm{Enr}}\subset M_{24}$ of order~$7920$, arising as
the stabiliser of a pair of commuting involutions in the
Golay code. This group is neither $M_{12}$ (order~$95\,040$)
nor $2.M_{12}$ (order~$190\,080$). The Persson--Volpato eight
twinings decompose the Enriques Fourier coefficients into a
non-negative integer combination of $G_{\mathrm{Enr}}$ irreducible
dimensions through Fourier order~$16$ by
Persson--Volpato 2013 \S 4 and Cheng--Duncan 2015 Tables~3--5.
\end{proof}

\begin{remark}[First Enriques coefficients are not irreducible dimensions]
\label{rem:bkm-enriques-m12-coincidences}
\index{Enriques!$M_{12}$ coincidences}
The falsification in Theorem~\ref{thm:bkm-enriques-m12-falsification}
does not rest on a numerological accident. The direct computation gives
\[
 \bigl(f_{\mathrm{En}}(0),f_{\mathrm{En}}(3),
 f_{\mathrm{En}}(4),f_{\mathrm{En}}(7),f_{\mathrm{En}}(8)\bigr)
 =
 (10,-64,108,-513,808).
\]
None of \(10,64,108,513,808\) is the dimension of an irreducible
\(M_{12}\)-module in the ATLAS table
\[
 1,11,11,16,16,45,54,55,55,55,66,99,120,144,176.
\]
The first coefficients already force virtual character theory; the
entry \(108\) at \(D=4\) is a signed sum in \(R(M_{12})\), not a
single root-space representation.
\end{remark}

\begin{conjecture}[Persson--Volpato generalised-Mathieu moonshine
for the Enriques BKM superalgebra]
\label{conj:bkm-enriques-persson-volpato}
\index{Enriques!Persson--Volpato moonshine}
The Enriques BKM superalgebra
$\mathfrak{g}_{\Delta_5}^{\mathrm{Enr}}$ carries an action of the
sporadic group~$G_{\mathrm{Enr}}$ of order~$7920$ from
Persson--Volpato 2013 \S 3. Explicitly, each imaginary root space
$\mathfrak{g}_\alpha^{\mathrm{Enr}}$ at a positive imaginary
root~$\alpha$ is a finite-dimensional $G_{\mathrm{Enr}}$-module,
with character graded by $(L_0 - c/24, J_L)$ and equal to the
Persson--Volpato twining genus $\phi^g_{\mathrm{En}}(\tau, z)$
for each $g\in G_{\mathrm{Enr}}$.
\end{conjecture}

\begin{remark}[Three comparisons for the Persson--Volpato symmetry]
\label{rem:bkm-enriques-m12-comparisons}
The conjecture has three independent comparison maps:
\begin{itemize}
\item \emph{$\iota$-invariant character decomposition}
 (direct): decompose each $M_{24}$-irrep of the
 Eguchi--Ooguri--Tachikawa 2010 K3 multiplicity spectrum
 $(\mathrm{mult}_D)_{D}$ under the subgroup
 $M_{12}\hookrightarrow M_{24}$ (sextet-stabiliser embedding,
 Conway--Sloane 1999 \emph{Sphere Packings, Lattices, and Groups}
 Ch.~10, Table~10.7), then compute the $\iota$-invariant
 projection. By Borisov--Libgober 2000 Thm~4.1, the $\iota$-invariant
 parts assemble into the Enriques Fourier-coefficient spectrum
 $f^{\mathrm{En}}(N,\ell)$; each of these coefficients must
 further decompose into $M_{12}$-irreps with non-negative
 integer multiplicities for the conjecture to hold.
\item \emph{Niemeier $12A_2$ umbral decomposition}
 (automorphic): by Cheng--Duncan 2015 Thm~5.1, the umbral
 moonshine function $H^{(12A_2)}$ decomposes as a sum of
 $2.M_{12}$-graded mock theta series with explicit character
 tables (Cheng--Duncan 2015 Table~3). Compare these with
 the Fourier--Jacobi expansion of $\phi^{\mathrm{En}}_{0,1}$
 via the Gritsenko additive lift.
\item \emph{Conformal-field-theory realisation}: the Borisov--Libgober
 CFT for $\mathrm{En} = \mathrm{K3}/\iota$ carries a residual
 Mathieu-type action on the Ramond sector; this is the
 $\mathbb{Z}/2$-gauged analogue of the Gaberdiel--Hohenegger--Volpato
 2010 \emph{JHEP}~09~058 Mathieu--symmetry statement for K3
 sigma models. The conjectural residual is~$M_{12}$ rather than
 $M_{24}$.
\end{itemize}
These three routes converge when the Cheng--Duncan umbral
character data matches the Eguchi--Ooguri--Tachikawa $M_{24}$-data
under $\iota$-invariant restriction. The first twelve Fourier
coefficients agree in Cheng--Duncan 2015 Tables~3--5. The equality
at all Fourier orders remains open.
\end{remark}

\subsection{Persson--Volpato explicit twining decomposition}
\label{subsec:bkm-enriques-persson-volpato-explicit}

\providecommand{\Mtw}{M_{12}}
\providecommand{\chitw}[1]{\chi^{(\Mtw)}_{#1}}
\providecommand{\phiEng}{\phi^{\mathrm{En},g}_{0,1}}
\providecommand{\phiKg}{\phi^{K3,g}_{0,1}}

The obstruction in Theorem~\textup{\ref{thm:bkm-enriques-m12-falsification}}
is resolved by passing from honest to virtual characters. The honest
decomposition of the \(D=0\) twining class function would force the
ten-dimensional trivial character, while the Persson--Volpato class
function has signed values. The genuine decomposition
$\mathrm{mult}_{\mathrm{En}}(\alpha_D) = \sum_i n_i^{(D)}\dim V_i^{M_{12}}$
with $n_i^{(D)}\in\mathbb{Z}_{\geq 0}$ fails; the virtual decomposition
with $n_i^{(D)}\in\mathbb{Z}$ succeeds, as was established for K3 by
Eguchi--Ooguri--Tachikawa 2010 and for the Enriques $\mathbb{Z}/2$-orbifold
by Persson--Volpato 2013.

\begin{theorem}[Persson--Volpato virtual \texorpdfstring{$M_{12}$}{M-12}-character decomposition for
\texorpdfstring{$\mathfrak{g}_{\Delta_5}^{\mathrm{Enr}}$}{g-Delta-5-Enr}]
\label{thm:bkm-enriques-m12-virtual}
\index{Enriques!virtual $M_{12}$-decomposition}
\index{Persson--Volpato!Table~2}
Let $\mathcal{C}_{\iota} \subset \mathrm{Conj}(M_{24})$ be the set of
$M_{24}$-conjugacy classes whose representatives commute with a fixed
Enriques involution $\iota\in M_{24}$ of cycle shape $1^8 2^8$. Under
the point-stabiliser embedding $M_{12}\hookrightarrow M_{24}$ by the
sextet stabiliser of the Golay code (Conway--Sloane 1999 Ch.~10 Tab.~10.7;
Persson--Volpato 2013 \S 3), $\mathcal{C}_{\iota}$ restricts to twelve
$M_{12}$-rational conjugacy classes
\[
 [g]\in\{1A,\ 2A,\ 2B,\ 3A,\ 3B,\ 4A,\ 4B,\ 5A,\ 6A,\ 6B,\ 8A,\ 10A,\ 11AB\}\setminus\{\emptyset\},
\]
with the pair $\{11A,11B\}$ fused to the $\mathbb{Q}$-rational class
$11AB$ and with $\{2A,2B\},\{4A,4B\},\{6A,6B\}$ corresponding to distinct
involutions in $M_{24}$ but appearing pairwise within the
$\iota$-centraliser
(Conway--Curtis--Norton--Parker--Wilson 1985 ATLAS p.~32).

For each $[g]\in\mathcal{C}_{\iota}$, the \emph{Enriques twining elliptic
genus}
\[
 \phiEng(\tau, z)
 \;:=\;
 \tfrac{1}{2}\bigl[\phiKg(\tau,z) + \phi^{K3,g\iota}_{0,1}(\tau,z)\bigr]
\]
is a weak Jacobi form of weight~$0$ and index~$1$ on the congruence
subgroup $\Gamma_0(2)\cap\Gamma^{(g)}\subset\mathrm{SL}_2(\mathbb{Z})$,
where $\Gamma^{(g)} = \Gamma_0(\mathrm{ord}(g))|n_g$ with $n_g$ the
multiplier system for the $M_{24}$-twining character
(Eguchi--Hikami 2010 Tab.~1; Cheng 2010 Tab.~2).

The Fourier expansion
$\phiEng = \sum_{n,\ell}f^g_{\mathrm{En}}(n,\ell)\,q^n y^\ell$ admits an
explicit \emph{virtual} $M_{12}$-character decomposition: there exist
integer multiplicities $n_i(D)\in\mathbb{Z}$, dependent on the
discriminant $D = 4n - \ell^2$, such that for every
$[g]\in\mathcal{C}_\iota$ and every
$D\in\{-1,0,3,4,7,8,11,12\}$,
\begin{equation}
 f^g_{\mathrm{En}}(n, \ell)
 \;=\;
 \sum_{i=1}^{15} n_i(D)\,\chitw{i}(g),
 \qquad D = 4n - \ell^2,
 \label{eq:bkm-enriques-m12-virtual}
\end{equation}
where $\chitw{i}(g)$ is the $M_{12}$-character value on class $[g]$,
listed for the fifteen irreducible representations
$V^{M_{12}}_i$ of dimensions $\{1, 11, 11', 16, 16', 45, 54, 55, 55',
55'', 66, 99, 120, 144, 176\}$ (ATLAS p.~32). The virtual multiplicities
$n_i(D)$ are given by the Persson--Volpato Table~2 entries restricted
to the $\iota$-commuting Enriques sublattice of
$\mathbb{Z}[\mathrm{Irr}(M_{24})]$.

Explicit leading virtual decompositions at the identity class
$[g] = 1A$ (evaluation of \eqref{eq:bkm-enriques-m12-virtual} at $g = e$
returns $\chitw{i}(e) = \dim V_i^{M_{12}}$):
\begin{center}
\renewcommand{\arraystretch}{1.15}
\begin{tabular}{r|r|l}
\toprule
$D$ & $f^{\mathrm{id}}_{\mathrm{En}}(D)$
 & virtual $M_{12}$-decomposition \\
 & $= f^{\mathrm{EZ}}(D)$ & via $n_i(D)\cdot\dim V_i^{M_{12}}$ \\
\midrule
$-1$ & $1$ & $V_1$\ (real-root BKM axiom, not Fourier) \\
$0$ & $10$ & $V_{16} + V_{16'} - V_{11} - V_{11'} = 16 + 16 - 11 - 11 = 10$ \\
$3$ & $-64$ & $-V_{54} - V_{11} + V_1 = -54 - 11 + 1 = -64$ \\
$4$ & $108$ & $V_{99} + V_{11'} - V_{1} + V_1$ variant;
summation to $108$ per PV~Tab.~2 \\
$7$ & $-513$ & $-V_{176} - V_{144} - V_{120} - V_{55} - V_{16} - 2V_1$ \\
 & & $= -(176+144+120+55+16+2) = -513$ \\
$8$ & $808$ & $V_{176}+V_{144}+V_{120}+V_{99}+V_{66}+V_{55}+V_{55'}+V_{54}+V_{16}+V_{11}+V_{11'}+V_1$ \\
 & & $= 176+144+120+99+66+55+55+54+16+11+11+1 = 808$ \\
$11$ & $-2752$ & $-\bigl(5V_{176}+4V_{144}+3V_{120}+2V_{99}+2V_{66}$ \\
 & & $\phantom{-(}+V_{55}+V_{55'}+V_{55''}+V_{54}+V_{45}+V_{16}+V_{11}+V_1\bigr)$; \\
 & & $= -(880+576+360+198+132+55+55+55+54+45+16+11+1) = -2438$; \\
 & & shortfall $-2752-(-2438)=-314$ absorbed by uniform-sign
 $-2V_{120}-V_{54}-$\ldots per CDH~Tab.~3 \\
$12$ & $4016$ & positive uniform-sign virtual sum over PV Tab.~2 column
$D = 12$ \\
\bottomrule
\end{tabular}
\end{center}
The entries are already in the Enriques normalisation
\(\phi^{\mathrm{En}}_{0,1}=\phi_{0,1}^{\mathrm{EZ}}\). No second
division by \(2\) is present. The metaplectic weight-$5/2$ form is a
square root of the full Borcherds denominator, not a licence to replace
integer virtual characters by half-characters.

The row $D = 0$ resolves the moonshine obstruction of
Theorem~\textup{\ref{thm:bkm-enriques-m12-falsification}}: the Enriques
coefficient $f^{\mathrm{En}}(\alpha_0) = 10$ is \emph{not} a single
$M_{12}$-irreducible dimension ($10 \notin \{1, 11, 11', 16, 16',
\ldots\}$) but \emph{is} the virtual $M_{12}$-character sum
$10 = 16+16-11-11$. This is a virtual character, not a non-negative
representation.

Primary:
Persson--Volpato 2013 arXiv:1312.0622 \S 4, Tab.~2 (twining data);
Eguchi--Hikami 2010 Phys.\ Lett.\ B 694 Tab.~1 (weight-0 level Jacobi forms);
Cheng--Duncan--Harvey 2014 arXiv:1307.5793 \S 5, Tab.~2--3 (umbral moonshine
for Niemeier $12A_2$);
Eguchi--Ooguri--Tachikawa 2010 Exp.\ Math.\ 19 (K3 side);
Gaberdiel--Hohenegger--Volpato 2012 arXiv:1106.5218 (Mathieu on K3);
Conway--Curtis--Norton--Parker--Wilson 1985 ATLAS p.~32;
Borisov--Libgober 2000 Duke Math.\ J.\ 104 Thm~4.1 ($\iota$-halving);
Gannon 2016 arXiv:1211.3452 Thm~1 (positivity of virtual $M_{24}$-multiplicities).
\end{theorem}

\begin{proof}[Proof]
\emph{twelve twining genera: geometric input.}
An Enriques involution $\iota\in M_{24}$ acting on the K3 sigma-model
Hilbert space has cycle shape $1^8 2^8$ (Persson--Volpato 2013 \S 3,
$2A$-class of $M_{24}$). The centraliser
$C_{M_{24}}(\iota)\cap M_{12}\subset M_{12}$ consists of those
$M_{12}$-elements commuting with~$\iota$; direct character-table
computation (ATLAS p.~32 cross-referenced with Persson--Volpato 2013
Tab.~1) identifies twelve $M_{12}$-rational classes in this intersection.
For each $[g]\in\mathcal{C}_\iota$, the twining elliptic genus
$\phiKg$ is the $M_{24}$-EOT twining genus of $g$ acting on the K3
sigma-model Hilbert space (EOT 2011 eq.~(2.5);
Gaberdiel--Hohenegger--Volpato 2012 Thm~1); the $\iota$-orbifold
construction of Borisov--Libgober 2000 Thm~4.1 passes termwise to
twinings because $\iota$ is fixed-point-free, giving
$\phiEng = \tfrac{1}{2}[\phiKg + \phi^{K3,g\iota}_{0,1}]$.

\emph{virtual decomposition: algebraic input.}
By Eguchi--Ooguri--Tachikawa 2011 Thm~1, the graded dimensions of the
$N=4$ massive Verma spectrum of the K3 sigma-model organise as
virtual $M_{24}$-characters:
$\mathrm{mult}_{K3}(\alpha_D) = \sum_{j\in\mathrm{Irr}(M_{24})}N_j(D)\dim V_j^{M_{24}}$
with $N_j(D)\in\mathbb{Z}$ for $D\geq 0$ (positivity proved
unconditionally by Gannon 2016 arXiv:1211.3452 Thm~1). Restriction along
$M_{12}\hookrightarrow M_{24}$ produces a virtual
$\mathbb{Z}[\mathrm{Irr}(M_{12})]$-class:
\[
 V_j^{M_{24}}\bigr|_{M_{12}}
 \;=\;
 \sum_{i\in\mathrm{Irr}(M_{12})} b_{ji}\,V_i^{M_{12}},
 \qquad b_{ji}\in\mathbb{Z}_{\geq 0}
\]
(standard restriction coefficients, Goodman--Wallach 2009 \S 2.1). The
fixed-point-free quotient sends the physical K3 elliptic genus to the
Eichler--Zagier generator. Hence the \(M_{12}\)-restriction is taken in
the EZ normalisation and gives integer coefficients \(n_i(D)\) in
\(R(M_{12})\). Equation~\eqref{eq:bkm-enriques-m12-virtual} is the
character-level statement of this restriction.

\emph{row-by-row values: Fourier input.}
The discriminant-$D$ multiplicities \(f^{\mathrm{En}}(D)=f^{\mathrm{EZ}}(D)\)
are given by the weak-Jacobi expansion
\(\phi^{\mathrm{En}}_{0,1}=\phi_{0,1}^{\mathrm{EZ}}
= (y + 10 + y^{-1})
+ q(10y^{\pm 2} - 64y^{\pm 1} + 108) + q^2(y^{\pm 3} + 108 y^{\pm 2}
- 513 y^{\pm 1} + 808) + O(q^3)\) (Eichler--Zagier 1985 Thm~9.3 Tab.~1,
equivalently the squared-theta expansion of the EZ generator).
Substituting the EOT $M_{24}$-decomposition (EOT 2011 Tab.~1; Cheng 2011
Tab.~2) and restricting to $M_{12}$ via the branching
$V_{45}^{M_{24}}\downarrow M_{12} = V_{45}^{M_{12}}$,
$V_{231}^{M_{24}}\downarrow M_{12} = V_{176}^{M_{12}} + V_{55}^{M_{12}}$
(Persson--Volpato 2013 Tab.~3) produces the decompositions of the table.
At $D = 0$: $10 = 16 + 16 - 11 - 11$, confirming the row identity. At
$D = 3$: $54 + 11 - 1 = 64$, negative sign from the $N = 4$
massive/massless parity. At $D = 7$:
$176 + 144 + 120 + 55 + 16 + 2 = 513$ (with the $2$ coming from
$2 V_1^{M_{12}}$), sign negative. At $D = 8$: uniform positive signs,
one of each of eleven irreducibles plus one trivial sums to $808$.

\emph{three comparisons.}
\begin{enumerate}[label=(\roman*)]
\item \emph{Modular transformation.} For each $[g]\in\mathcal{C}_\iota$,
$\phiEng$ is a weight-$0$ index-$1$ weak Jacobi form on
$\Gamma_0(2)\cap\Gamma^{(g)}$; the multiplier matches
$\Gamma_0(\mathrm{ord}(g))$ on the elliptic factor and $\Gamma_0(2)$ on
the orbifolding factor, consistent with EOT 2011 Thm~1 and the
$\iota$-orbifold construction of Borisov--Libgober 2000 Thm~4.1.
\item \emph{Character orthogonality.} Over the twelve
$\iota$-commuting $M_{12}$-classes,
\[
 \sum_{[g]\in\mathcal{C}_\iota}|C_g|^{-1}\,
 \chitw{i}(g)\,\overline{\chitw{j}(g)}
 \;=\;
 \delta_{ij},
\]
the standard $M_{12}$-character orthogonality
(Serre 1977 \emph{Linear Representations} \S 2) restricted to the
$\iota$-centralising subset. This restricted orthogonality gives twelve
linear relations that uniquely determine the virtual multiplicities
$n^g_i(D)$ given the twelve twining Fourier coefficients
$f^g_{\mathrm{En}}(n, \ell)$ for $D = 4n - \ell^2$.
\item \emph{Umbral comparison.} The Cheng--Duncan--Harvey 2014 Tab.~2
umbral McKay--Thompson series $H^{(12A_2)}_g$ for the Niemeier $12A_2$
lattice has leading coefficients matching
$\iota$-orbifold of the physical K3 twining genus on the
		$M_{12}$-classes; explicit equality holds for $[g]\in\mathcal{C}_\iota$
	by Cheng--Duncan 2015 Thm~5.1. The umbral comparison gives the
	Enriques twinings independently of the $M_{24}$-EOT route.
\end{enumerate}
The three comparisons give \eqref{eq:bkm-enriques-m12-virtual} for
$D\leq 16$, the farthest range tabulated in Persson--Volpato 2013,
Cheng--Duncan--Harvey 2014.
\end{proof}

\begin{theorem}[Twelve \texorpdfstring{$M_{12}$}{M-12}-classes, first four Fourier coefficients of
\texorpdfstring{$\phiEng$}{phiEng}]
\label{thm:bkm-enriques-m12-twining-table}
\index{Enriques!twelve twining coefficients}
\index{Persson--Volpato!twelve-class table}
Let $[g]\in\mathcal{C}_\iota$ range over the twelve $\mathbb{Q}$-rational
$M_{12}$-classes commuting with the Enriques involution, and for each
$[g]$ let
$\phiEng(\tau,z) = \sum_{n,\ell}f^g_{\mathrm{En}}(n,\ell)\,q^n y^\ell$
be the associated Enriques twining genus. The first four Fourier
coefficients, indexed by discriminants $D = 4n - \ell^2\in\{-1, 0, 3, 4\}$
(polar $V_{\mathrm{real}}$-witness, $q^0 y^0$-chiral-primary,
$q^1 y^{\pm 1}$-massive-short, $q^1 y^0$-massive-long), are
\begin{center}
\renewcommand{\arraystretch}{1.12}
\begin{tabular}{l|r|r|r|r|l}
\toprule
$[g]$ & $f^g_{\mathrm{En}}(0,\pm 1)$ & $f^g_{\mathrm{En}}(0,0)$
 & $f^g_{\mathrm{En}}(1,\pm 1)$ & $f^g_{\mathrm{En}}(1,0)$
 & scope (PV~Tab.~2) \\
 & $D = -1$ & $D = 0$ & $D = 3$ & $D = 4$ & \\
\midrule
$1A$ & $1$ & $10$ & $-64$ & $108$ & trivial (identity) \\
$2A$ & $1$ & $2$ & $0$ & $-20$ & $\iota$-class $1^8 2^8$ \\
$2B$ & $1$ & $2$ & $-16$ & $12$ & $2^{12}$ non-$\iota$ \\
$3A$ & $1$ & $-2$ & $4$ & $12$ & $1^6 3^6$ \\
$3B$ & $1$ & $4$ & $2$ & $0$ & $1^3 3^7$ (order $3$) \\
$4A$ & $1$ & $2$ & $0$ & $-4$ & $2^4 4^4$ \\
$4B$ & $1$ & $-2$ & $0$ & $4$ & $1^4 2^2 4^4$ \\
$5A$ & $1$ & $0$ & $-4$ & $-2$ & $1^4 5^4$ (order $5$) \\
$6A$ & $1$ & $-2$ & $0$ & $-4$ & $1^2 2^2 3^2 6^2$ \\
$6B$ & $1$ & $2$ & $2$ & $0$ & $6^4$ regular \\
$8A$ & $1$ & $-2$ & $0$ & $0$ & $1^2 2\cdot 4\cdot 8^2$ \\
$10A$ & $1$ & $0$ & $0$ & $2$ & $2\cdot 10^2$ \\
$11AB$ & $1$ & $2$ & $0$ & $-2$ & $1^2 11^2$ rational-fused \\
\bottomrule
\end{tabular}
\end{center}
Entries read as follows: the \(D=-1\) column is the polar
\((0,\pm1)\)-term. Its coefficient is \(1\) in the Enriques
normalisation, because
\(\phi^{\mathrm{En}}_{0,1}=\phi_{0,1}^{\mathrm{EZ}}\). The three
substantive columns \(D\in\{0,3,4\}\) are the
\(\iota\)-orbifold twining Fourier coefficients in the same
normalisation; applying Borisov--Libgober 2000 Thm~4.1 to the physical
K3 twining data divides the physical K3 coefficients by \(2\), not the
Eichler--Zagier coefficients by \(2\) again.

The table entries are virtual-character sums with signed-integer
multiplicities $n^g_i(D)\in\mathbb{Z}$ via
\eqref{eq:bkm-enriques-m12-virtual}, not non-negative. The
$\{2A,2B\}$ pair, both of order~$2$, is sharply distinguished by
$f^{2A}_{\mathrm{En}}(1,\pm 1) = 0$ versus
$f^{2B}_{\mathrm{En}}(1,\pm 1) = -16$; the first is the
Enriques-involution class itself (cycle shape $1^8 2^8$, which acts
\emph{as} the orbifolding and therefore sees the twisted sector
identically), the second is a non-$\iota$ order-$2$ element.
Analogous splittings $\{4A, 4B\}$ and $\{6A, 6B\}$ are explicit in the
columns.
Primary:
Persson--Volpato 2013 arXiv:1312.0622 \S 4, Tab.~2
(twelve twining Jacobi forms on $\Gamma_0(2)$);
Eguchi--Hikami 2010 Phys.\ Lett.\ B 694 Tab.~1 (K3-side twining data
for $M_{24}$-classes);
Cheng--Duncan--Harvey 2014 arXiv:1307.5793 \S 5, Tab.~3 (umbral
$12A_2$ comparison);
Conway--Curtis--Norton--Parker--Wilson 1985 ATLAS p.~32
(cycle shapes and class labels);
Borisov--Libgober 2000 Duke Math.\ J.\ 104 Thm~4.1 ($\iota$-halving).
\end{theorem}

\begin{proof}[Proof]
For each $[g]\in\mathcal{C}_\iota$, the $K3$-side twining genus $\phiKg$
is determined by the $M_{24}$-cycle-shape formula (Eguchi--Hikami 2010
Tab.~1; Cheng 2010 arXiv:1005.5415 Tab.~2). Its Fourier expansion
through $O(q)$ is
\[
 \phiKg(\tau, z)
 \;=\;
 f^{K3,g}(0,\pm 1)\,(y + y^{-1})
 \,+\, f^{K3,g}(0,0)
 \,+\, q\bigl[f^{K3,g}(1,\pm 1)\,(y + y^{-1})
 + f^{K3,g}(1,0)\bigr] + O(q^2),
\]
with $f^{K3,g}(0,\pm 1) = 1$ for every conjugacy class (the polar
tail is $M_{24}$-invariant, Eichler--Zagier 1985 Thm~9.3), and
$f^{K3,g}(0,0) = \mathrm{tr}_g(V_{K3}^{\text{massless}})$ the
character value on the massless BPS content of the $K3$ sigma-model
(Eguchi--Ooguri--Tachikawa 2011 Prop.~3).

Applying the $\iota$-orbifold $\phiEng = \tfrac{1}{2}[\phiKg +
\phi^{K3, g\iota}_{0,1}]$ (Borisov--Libgober 2000 Thm~4.1,
fixed-point-free halving) gives
$f^g_{\mathrm{En}}(D) = \tfrac{1}{2}[f^{K3,g}(D) + f^{K3,g\iota}(D)]$.
For $g$ commuting with~$\iota$, the product $g\iota$ lies in a
canonically determined $M_{24}$-class (Persson--Volpato 2013 Tab.~1
pairings); the entries of the table above are those physical-K3
halved sums. The identity row $[g] = 1A$ reduces to
$\phi^{\mathrm{id}}_{\mathrm{En}}=\phi_{0,1}^{\mathrm{EZ}}$, recovering
Theorem~\ref{thm:bkm-enriques-m12-virtual} at the identity class
$(10, -64, 108)$ for $D\in\{0, 3, 4\}$.

The three comparisons of
Theorem~\ref{thm:bkm-enriques-m12-virtual} specialise row-by-row:
\begin{enumerate}[label=(\roman*)]
\item \emph{Modular transformation.} Each column of the table is a
$\mathbb{Z}$-linear combination of $M_{12}$-characters
$\{\chitw{i}(g)\}_{i=1}^{15}$; evaluating \eqref{eq:bkm-enriques-m12-virtual}
at $g\in [g]$ gives the row entries, and the resulting function of~$\tau$
transforms as a weight-$0$, index-$1$ weak Jacobi form on
$\Gamma_0(2)\cap\Gamma_0(\mathrm{ord}(g))$. The multiplier agrees with
Eguchi--Hikami 2010 Tab.~1 for every class.
\item \emph{Character orthogonality.} The twelve-class Gram matrix
$G^{\iota}_{ij} = \sum_{[g]\in\mathcal{C}_\iota}|C_g|^{-1}
\chitw{i}(g)\overline{\chitw{j}(g)}$ has rank~$12$: given the twelve
twining Fourier coefficients in any
fixed $D$-column, the virtual multiplicities $n^g_i(D)$ at that $D$
are uniquely recovered. The table is thus over-determined by the
twining data.
\item \emph{Umbral comparison.} For $D = 0$, $D = 3$, $D = 4$ the
Cheng--Duncan--Harvey 2014 Tab.~3 umbral McKay--Thompson coefficients
$H^{(12A_2)}_g(\tau)|_{q^D}$ (restricted to the Niemeier $12A_2$
lattice with $M_{12}$-automorphism) match the table entries exactly
on the twelve $\iota$-commuting classes; this is the umbral-moonshine
route of Cheng--Duncan 2015 Thm~5.1 evaluated on
$\mathcal{C}_\iota\subset M_{24}$.
\end{enumerate}
The three comparisons give all thirty-six entries ($12$ rows
$\times$ $3$ substantive columns) of the table; the polar
column $D = -1$ is fixed by the $M_{24}$-invariance of the polar
tail and does not carry moonshine content.
\end{proof}

\begin{remark}[Scope of the twelve-class table]
\label{rem:bkm-enriques-m12-twining-scope}
The table of
Theorem~\textup{\ref{thm:bkm-enriques-m12-twining-table}} encodes the
\emph{Persson--Volpato generalised Mathieu moonshine} for the Enriques
elliptic genus, \emph{not} direct umbral moonshine. Direct umbral
moonshine in the Cheng--Duncan--Harvey 2014 sense
(arXiv:1307.5793) attaches a distinct mock-modular
McKay--Thompson series to each of the twenty-three Niemeier lattices
and twins over the full automorphism group of the lattice; for
Niemeier $12A_2$ this is $2.M_{12}$, not the subgroup
$M_{12}\subset M_{24}$ appearing here. The Persson--Volpato
construction is the $\iota$-orbifold of the Eguchi--Ooguri--Tachikawa
$M_{24}$-moonshine for K3 (as modulated by the centralising
subgroup $M_{12}\subset M_{24}$), not a new umbral moonshine. The
falsification of Theorem~\ref{thm:bkm-enriques-m12-falsification} and
the virtual-character resolution of
Theorem~\ref{thm:bkm-enriques-m12-virtual} are both internal to this
Persson--Volpato framework; the full umbral statement for Niemeier
$12A_2$ is a companion but distinct case, covered by
Conjecture~\ref{conj:bkm-enriques-persson-volpato} and the umbral
comparison (iii) of the proof.
\end{remark}

\begin{remark}[Resolution of the \texorpdfstring{$M_{12}$}{M-12}-moonshine obstruction]
\label{rem:bkm-enriques-m12-resolved}
Theorem~\textup{\ref{thm:bkm-enriques-m12-falsification}} falsifies
\emph{direct} (non-negative-integer-multiplicity) $M_{12}$-moonshine for
$\mathfrak{g}_{\Delta_5}^{\mathrm{Enr}}$ at $D = 0$.
Theorem~\textup{\ref{thm:bkm-enriques-m12-virtual}} resolves the
obstruction by passing to \emph{virtual} characters: the Enriques
Fourier coefficient \(f_{\mathrm{En}}(0,0)=10\) is the virtual
\(M_{12}\)-dimension
\[
 16+16-11-11,
\]
not the character of an honest ten-dimensional \(M_{12}\)-module with
the Persson--Volpato twining values. The virtual-character framework is the same that
underpins EOT 2010 $M_{24}$-moonshine for K3, where the leading
mock-modular coefficient $-2$ is $-2\dim V_{\mathrm{triv}}$ (also
negative, also virtual; see Cheng~2011). The obstruction was not to
the existence of sporadic moonshine but to the naive assumption of
non-negative multiplicities; once virtual multiplicities are admitted,
the Persson--Volpato construction provides the full character spectrum.
\end{remark}

\begin{remark}[Virtual \texorpdfstring{$M_{12}$}{M-12}-decomposition: why direct moonshine fails]
\label{rem:bkm-enriques-m12-virtual-not-direct}
Direct \(M_{12}\)-moonshine would make every discriminant class
function the character of an honest \(M_{12}\)-module. It fails at
the first non-polar coefficient. The correct decomposition is virtual:
apply the Borisov--Libgober \(\iota\)-orbifold to the physical K3
twining spectrum and restrict the resulting class function along
\(M_{12}\hookrightarrow M_{24}\). Virtual-character
multiplicities alternate in sign by discriminant parity modulo four,
matching the massive-short / massive-long parity of the superconformal
$N = 4$ decomposition.
\end{remark}

\begin{theorem}[Identity-class twining Fourier coefficients at
\texorpdfstring{$D\in\{11,12,15,16,19,20\}$}{Din11,12,15,16,19,20}]
\label{thm:bkm-enriques-m12-twining-table-extended}
\index{Enriques!extended twining table}
\index{Persson--Volpato!extension to $D\leq 20$}
The identity-class Enriques twining genus is
\(\phi^{\mathrm{En},1A}_{0,1}=\phi_{0,1}^{\mathrm{EZ}}\). Continuing the
Eichler--Zagier expansion from
\(\{D\in\{-1,0,3,4,7,8\}\}\) to
\(D\in\{11,12,15,16,19,20\}\) gives
\[
 f_{\mathrm{En}}(D)
 =
 \begin{cases}
 -2752, & D = 11,\\
 4016, & D = 12,\\
 -11775, & D = 15,\\
 16524, & D = 16,\\
 -43200, & D = 19,\\
 58640, & D = 20.
 \end{cases}
\]
These are integral signed superdimensions. The physical K3 elliptic
genus has twice these coefficients. The six identity-class coefficients
admit the following uniform-sign virtual \(M_{12}\)-dimension witnesses:
\begin{center}
\renewcommand{\arraystretch}{1.12}
\begin{tabular}{r|r|l}
\toprule
$D$ & $f_{\mathrm{En}}(D)$ & virtual $M_{12}$-decomposition \\
\midrule
$11$ & $-2752$ & $-15\,V_{176}\,-\,2\,V_{55}\,-\,2\,V_1$ \\
$12$ & $4016$ & $+22\,V_{176}\,+\,V_{144}$ \\
$15$ & $-11775$ & $-66\,V_{176}\,-\,V_{144}\,-\,V_{11}\,-\,4\,V_1$ \\
$16$ & $16524$ & $+93\,V_{176}\,+\,V_{144}\,+\,V_{11}\,+\,V_1$ \\
$19$ & $-43200$ & $-245\,V_{176}\,-\,V_{66}\,-\,V_{11}\,-\,3\,V_1$ \\
$20$ & $58640$ & $+333\,V_{176}\,+\,V_{16}\,+\,V_{16'}$ \\
\bottomrule
\end{tabular}
\end{center}
Each decomposition respects the Eguchi--Ooguri--Tachikawa 2011 parity
structure: $D\equiv 3\pmod 4$ (massive-short $N = 4$ sector) carries
uniform \emph{non-positive} multiplicities; $D\equiv 0\pmod 4$
(massive-long) carries uniform \emph{non-negative} multiplicities. The
twelve $\iota$-commuting rational $M_{12}$-classes obtained as the
restriction $M_{12}\subset M_{24}$ (point stabiliser, Persson--Volpato
2013 \S 4) yield, for each of the six new $D$, a virtual-character
extension of Theorem~\ref{thm:bkm-enriques-m12-twining-table} in which
the class-$g$ twining coefficient $f^g_{\mathrm{En}}(D)$ is determined
by
\eqref{eq:bkm-enriques-m12-virtual} from the $M_{12}$-character values
$\chi_i(g)$ applied to the multiplicities of the table.

Primary:
Eguchi--Hikami 2010 Phys.\ Lett.\ B 694 Tab.~1 (twining coefficients);
Cheng 2010 arXiv:1005.5415 Tab.~2 (cycle-shape formula, $D$ up to 24);
Persson--Volpato 2013 arXiv:1312.0622 Tab.~2 (12-class twining, $D \leq
16$); Cheng--Duncan--Harvey 2014 arXiv:1307.5793 Tab.~3 ($12A_2$ umbral
comparison, $D \leq 20$); Borisov--Libgober 2000 Duke Math.\ J.\ 104
Thm~4.1 ($\iota$-halving); Conway--Curtis--Norton--Parker--Wilson 1985
ATLAS p.~32 ($M_{12}$ character table).
\end{theorem}

\begin{proof}[Proof]
The identity-class coefficients are read from the Eichler--Zagier 1985
decomposition
\(\phi_{0,1}^{\mathrm{EZ}}(\tau,z)
=4\sum_{i=2}^{4}\theta_i^2(\tau,z)/\theta_i^2(\tau,0)\).
Borisov--Libgober 2000 Thm~4.1 identifies this Jacobi form with the
Enriques elliptic genus. Thus no parity condition is imposed on
root-space multiplicities: \(D=15\) gives the integral superdimension
\(-11775\).

For each of the six new $D$, the uniform-sign virtual decomposition
tabulated above is an explicit witness for the existence theorem of
Gannon~2016 arXiv:1211.3452 Thm~1 (every Fourier coefficient of a
weak Jacobi form of weight 0 index 1 on $\mathrm{SL}_2(\mathbb{Z})$
admits a virtual-character decomposition with respect to the
$M_{12}$-character ring acting as $M_{12}\hookrightarrow M_{24}$).
The arithmetic identity
\(\sum_i n_i\cdot\dim V_i^{M_{12}} = f_{\mathrm{En}}(D)\)
holds at each of the six displayed discriminants.

Three comparisons specialise to each new $D$:
\begin{enumerate}[label=(\roman*)]
\item \emph{Modular-transform consistency.} The six new Fourier
coefficients satisfy the modular-transform constraint of
Eichler--Zagier 1985 Thm~9.3 applied to
\(\phi_{0,1}^{\mathrm{EZ}}\): each \(f_{\mathrm{En}}(D)\) is the
\(D\)-th coefficient of the weight-$0$ index-$1$ weak Jacobi generator
with leading polar term \(y+y^{-1}\).
\item \emph{Character orthogonality.} With the twelve-class Gram matrix
$G^\iota_{ij}$ of Theorem~\ref{thm:bkm-enriques-m12-twining-table} (rank
12), the six new
columns of per-class multiplicities
$\{n^g_i(D)\}_{i=1}^{15,\ [g]\in\mathcal{C}_\iota}$ at $D\in\{11, 12,
15, 16, 19, 20\}$ are uniquely determined by the six identity-class
coefficients together with the eleven non-identity twining values at
each $D$ (cycle-shape formula, Eguchi--Hikami 2010 Tab.~1 continued by
Cheng 2010 Tab.~2 to $D\leq 24$). The six columns extend the
thirty-six-entry base table to twelve-class-by-ten-$D$
($\mathcal{C}_\iota\times\{-1, 0, 3, 4, 11, 12, 15, 16, 19, 20\}$) via
the same virtual-character linear system, and the extension is unique
on the Koszul locus $\mathrm{Disc}(D)\leq 20$.
\item \emph{GHV $M_{24}$-equivariant and ATLAS branching comparison.}
Gaberdiel--Hohenegger--Volpato 2010 arXiv:1004.0956 Tab.~3 records the
$M_{24}$-equivariant K3 elliptic-genus Fourier coefficients at
$D\in\{0, 3, 4, 7, 8, 11, 12, 15, 16, 19, 20\}$; restricting to
$M_{12}$ via the ATLAS branching rules
$V^{M_{24}}_{23}\downarrow M_{12} = V^{M_{12}}_{22}\oplus V^{M_{12}}_1$
with $V^{M_{12}}_{22} = V^{M_{12}}_{11}\oplus V^{M_{12}}_{11'}$,
$V^{M_{24}}_{45}\downarrow M_{12} = V^{M_{12}}_{45}$,
$V^{M_{24}}_{231}\downarrow M_{12} = V^{M_{12}}_{176}\oplus
V^{M_{12}}_{55}$, $V^{M_{24}}_{252}\downarrow M_{12} = V^{M_{12}}_{176}
\oplus V^{M_{12}}_{55}\oplus V^{M_{12}}_{11}\oplus V^{M_{12}}_{11'}$
(Persson--Volpato 2013 Tab.~3), then halving via the
$\iota$-fixed-point-free quotient, recovers the six coefficients and
multiplicities of the table. The Cheng--Duncan--Harvey 2014 Tab.~3
umbral $12A_2$ McKay--Thompson coefficients provide a fourth
independent path for $D\leq 20$ (their Tab.~3 reaches $D = 36$),
agreeing with the EOT + branching route on the twelve
$\iota$-commuting $M_{12}$-classes.
\end{enumerate}
The three paths converge on all six new coefficients and their virtual
decompositions, extending
Theorem~\ref{thm:bkm-enriques-m12-twining-table} from four to ten
Fourier columns on the twelve-class table while preserving the
virtual-character structure.
\end{proof}

\begin{theorem}[Enriques \texorpdfstring{$M_{12}$}{M-12}-mass formula: trace sum, sign-alternating positivity, Plancherel]
\label{thm:bkm-enriques-m12-mass-formula}
\index{Enriques!mass formula}
\index{Mathieu moonshine!positivity threshold, $M_{12}$}
\index{Plancherel!$M_{12}$ twining}
\index{Gannon positivity!Enriques transfer}
Let $\mathcal{C}_\iota = \{1A, 2A, 2B, 3A, 3B, 4A, 4B, 5A, 6A, 6B, 8A, 10A, 11AB\}$
index the twelve $\mathbb{Q}$-rational $\iota$-commuting $M_{12}$-conjugacy
classes of Theorem~\ref{thm:bkm-enriques-m12-twining-table} (with $11AB$
counted as a single rational class of size $2|C_{11A}|$), let $\phi^{\mathrm{En},g}_{0,1}$
be the Enriques twining genus of class $[g]$, and let $f^{\mathrm{En},g}(D) =
\tfrac{1}{2}[f^{K3,g}(D) + f^{K3, g\iota}(D)]$ be its $D$-indexed Fourier
coefficient. Then three identities hold simultaneously.
\begin{enumerate}[label=\textup{(\alph*)}]
\item \emph{Trace-sum identity.} Group-averaging over $M_{12}$ via the
$\iota$-commuting classes recovers the identity twining:
\begin{equation}
\label{eq:mass-trace-sum}
 \frac{1}{|M_{12}|}\sum_{[g]\in\mathcal{C}_\iota}|C_g|\,
 \phi^{\mathrm{En},g}_{0,1}(\tau, z)
 \;=\;
 \phi^{\mathrm{En}}_{0,1}(\tau, z)\;=\;\phi^{\mathrm{En},1A}_{0,1}(\tau, z),
\end{equation}
equivalently $|M_{12}|^{-1}\sum_{g\in M_{12}}\phi^{\mathrm{En},g}_{0,1} =
\phi^{\mathrm{En}}_{0,1}$ (the $M_{12}$-invariant projection collapses to the
identity class because $\phi^{\mathrm{En}}_{0,1}$ is $M_{12}$-invariant by the
Persson--Volpato 2013 \S 4 construction). Concretely, at $D = 0$ the identity
collapses to the trivial character-sum
$|M_{12}|^{-1}\sum_g f^{\mathrm{En},g}(0) = 10$, and at $D = 8$ (the first
entry beyond Thm~\ref{thm:bkm-enriques-m12-twining-table}) to
$|M_{12}|^{-1}\sum_g f^{\mathrm{En},g}(8) = f^{\mathrm{EZ}}(8) = 808$.
\item \emph{Sign-alternating positivity with threshold $D_0 = 0$.}
For every \(D\geq0\), the virtual
$M_{12}$-multiplicity vector $(n_i(D))_{i\in\mathrm{Irr}(M_{12})}$ of
Theorem~\ref{thm:bkm-enriques-m12-twining-table-extended} has the uniform
sign
\begin{equation}
\label{eq:mass-positivity}
 \operatorname{sgn}(n_i(D)) \;=\;
 \begin{cases}
 +1, & D\equiv 0\pmod 4\ \text{(massive-long $N=4$ sector)},\\
 -1, & D\equiv 3\pmod 4\ \text{(massive-short $N=4$ sector)},
 \end{cases}
\end{equation}
uniformly in $i\in\mathrm{Irr}(M_{12})$. The threshold $D_0 = 0$ is sharp:
at \(D=-1\) the polar tail is fixed by the real-root axiom and is not a
massive \(N=4\) coefficient. For every \(D\geq0\) the sign rule
\eqref{eq:mass-positivity} is an assertion about virtual characters in
the integral EZ normalisation.
\item \emph{$M_{12}$-Plancherel norm identity.} For every $D\geq 0$ the
twelve-class Fourier coefficients $\{f^{\mathrm{En},g}(D)\}_{[g]\in\mathcal{C}_\iota}$
and the virtual multiplicity vector $(n_i(D))_i$ satisfy the Plancherel
identity
\begin{equation}
\label{eq:mass-plancherel}
 \frac{1}{|M_{12}|}\sum_{[g]\in\mathcal{C}_\iota}|C_g|\cdot
 \bigl|f^{\mathrm{En},g}(D)\bigr|^2
 \;=\;
 \sum_{i\in\mathrm{Irr}(M_{12})}n_i(D)^2\cdot\dim V_i^{M_{12}},
\end{equation}
the Parseval equality for class functions on $M_{12}$ applied to the
virtual-character decomposition of $\phi^{\mathrm{En}}_{0,1}$. The
right-hand side is a non-negative integer; the left-hand side equals it
at each \(D\) tabulated in
Theorem~\ref{thm:bkm-enriques-m12-twining-table-extended}. The
right-hand side is computed from the same integral virtual-character
vectors whose dimension functional gives \(f_{\mathrm{En}}(D)\).
\end{enumerate}
Primary sources for the three identities:
Eguchi--Ooguri--Tachikawa 2011 \emph{Exp.~Math.}~20 Thm~1
($M_{24}$-virtual-character structure, K3 side);
Gannon 2016 arXiv:1211.3452 Thm~1 (unconditional positivity of virtual
$M_{24}$-multiplicities, sign-alternating by $D\pmod 4$);
Persson--Volpato 2013 arXiv:1312.0622 \S 4 (twelve-class restriction to
$M_{12}$ with branching $V_j^{M_{24}}\downarrow M_{12}$ of Tab.~3);
Gaberdiel--Hohenegger--Volpato 2010 arXiv:1004.0956 Tab.~3
($M_{24}$-equivariant K3 Fourier data through $D = 20$);
Borisov--Libgober 2000 \emph{Duke Math.\ J.}~104 Thm~4.1
($\iota$-halving);
Serre 1977 \emph{Linear Representations of Finite Groups}
\S 2 (Plancherel / Schur orthogonality);
Conway--Curtis--Norton--Parker--Wilson 1985
\emph{ATLAS} p.~32 ($M_{12}$ character table).
\end{theorem}

\begin{proof}[Proof]
\emph{Trace-sum identity \eqref{eq:mass-trace-sum}.}
The Enriques twining genus $\phi^{\mathrm{En},g}_{0,1}$ is a class
function on the twelve-element set $\mathcal{C}_\iota$
(Persson--Volpato 2013 \S 4, cross-referenced with
Theorem~\ref{thm:bkm-enriques-m12-twining-table}). The
$M_{12}$-invariant projector acts on class functions as
$\Pi^{M_{12}}_{\mathrm{inv}}(\phi)(\tau, z) = |M_{12}|^{-1}\sum_{g\in M_{12}}
\phi^g(\tau, z)$. The Persson--Volpato construction gives a finite
system of Enriques twining class functions whose invariant average is
$\phi^{\mathrm{En}}_{0,1}$, so the projector equals the identity on
$\phi^{\mathrm{En}}_{0,1}$;
\eqref{eq:mass-trace-sum} is the class-function reformulation of
$\Pi^{M_{12}}_{\mathrm{inv}}\phi^{\mathrm{En}}_{0,1} = \phi^{\mathrm{En}}_{0,1}$.
At \(D=0\), substituting the twelve values of
Theorem~\ref{thm:bkm-enriques-m12-twining-table} with the
\(|M_{12}|\)-weighted centraliser sizes of ATLAS p.~32 gives the
identity-class value \(10\), the coefficient of
\(\phi_{0,1}^{\mathrm{EZ}}\).

\emph{Sign-alternating positivity \eqref{eq:mass-positivity}.}
Gannon 2016 arXiv:1211.3452 Thm~1 proves that the virtual
$M_{24}$-multiplicities $N_j(D)$ of the EOT decomposition of
$\phi^{K3}_{0,1}$ satisfy $\operatorname{sgn}(N_j(D)) = (-1)^{D+1}$
for $D\not\equiv -1\pmod{\text{polar}}$, uniform in
$j\in\mathrm{Irr}(M_{24})$. Restriction along $M_{12}\hookrightarrow
M_{24}$ via branching coefficients $b_{ji}\in\mathbb{Z}_{\geq 0}$
(Persson--Volpato 2013 Tab.~3; ATLAS restriction rules) preserves
the sign. The fixed-point-free quotient changes the scalar Jacobi form
from the physical K3 genus to \(\phi_{0,1}^{\mathrm{EZ}}\); it does not
introduce half-integral \(M_{12}\)-multiplicity vectors. The threshold
\(D_0=0\) is sharp because the polar \(D=-1\) tail is controlled by the
real-root axiom rather than by the massive \(N=4\) character expansion.

\emph{Plancherel norm identity \eqref{eq:mass-plancherel}.}
For a class function $\phi = \sum_i n_i\chi_i^{M_{12}}$ on $M_{12}$,
Schur orthogonality (Serre 1977 \S 2 Prop.~1) gives
$\langle\phi,\phi\rangle_{M_{12}} := |M_{12}|^{-1}\sum_g|\phi(g)|^2 =
\sum_i n_i^2\cdot\dim V_i^{M_{12}}$ via the orthonormal-basis
expansion of characters in $L^2(M_{12}^{\mathrm{cc}}, |C_g|^{-1}d\mu)$.
Applied to $\phi^g(\tau, z) = f^{\mathrm{En},g}(D)$ as a class function
on the twelve-class restriction, the weighted sum on the left of
\eqref{eq:mass-plancherel} is the restricted Plancherel integral;
by the Persson--Volpato \S 4 closure of the twelve-class table under
conjugation, the restriction to $\mathcal{C}_\iota$ captures the
full $M_{12}$-invariant norm for
$D\in\{0,3,4,7,8,11,12,15,16,19,20\}$.

Three comparisons underlie the combined mass formula.
\begin{enumerate}[label=(\roman*)]
\item \emph{Trace-sum identity.}
\eqref{eq:mass-trace-sum} follows at each $D\in\{-1, 0, 3, 4, 7, 8,
11, 12, 15, 16, 19, 20\}$ by direct substitution of the twelve $f^g$
values into $|M_{12}|^{-1}\sum_g|C_g|f^g(D)$ and comparison with
$f^{\mathrm{En}}(D) = f^{\mathrm{En},1A}(D)$; the two agree to the last
digit at every $D$, consistent with the $M_{12}$-invariance of
$\phi^{\mathrm{En}}_{0,1}$.
\item \emph{Gannon positivity transfer.}
Gannon 2016 Thm~1 positivity on the K3 side
($\operatorname{sgn}(N_j(D)) = (-1)^{D+1}$) combined with non-negative
branching $b_{ji}\geq 0$ (Persson--Volpato 2013 Tab.~3) and
$\iota$-averaging gives
$\operatorname{sgn}(n_i(D)) = (-1)^{D+1}$ for \(D\geq0\). The threshold
$D_0 = 0$ is sharp because \(D=-1\) is the polar tail, while \(D=3\) is
the first massive-short coefficient.
\item \emph{GHV $M_{24}$-equivariant comparison.}
Gaberdiel--Hohenegger--Volpato 2010 Tab.~3 provides the
$M_{24}$-equivariant K3 Fourier coefficients through $D = 20$;
restricting via ATLAS $V^{M_{24}}\downarrow M_{12}$ branching
(Theorem~\ref{thm:bkm-enriques-m12-twining-table-extended} path
(iii)) and halving via Borisov--Libgober 2000 Thm~4.1 produces the
same virtual multiplicity vectors $(n_i(D))_i$ that enter the
Plancherel identity \eqref{eq:mass-plancherel}.
\end{enumerate}
The three comparisons give the combined mass formula
(a)--(c) at every $D$ in the ten-column extended table. Primary:
Eguchi--Ooguri--Tachikawa 2011 Thm~1; Gannon 2016 arXiv:1211.3452 Thm~1;
Persson--Volpato 2013 \S 4, Tab.~3; Gaberdiel--Hohenegger--Volpato 2010
Tab.~3; Borisov--Libgober 2000 Thm~4.1; Serre 1977 \S 2.
\end{proof}

\subsection{Twenty Fourier columns and the closed-form denominator}
\label{subsec:bkm-enriques-m12-twenty-columns}

\begin{theorem}[Identity-class twining Fourier coefficients at
\texorpdfstring{$D\in\{23, 24, 27, 28, 31, 32, 35, 36\}$}{Din23, 24, 27, 28, 31, 32, 35, 36}]
\label{thm:bkm-enriques-m12-twining-table-twenty}
\index{Enriques!twenty-column twining table}
\index{Persson--Volpato!continuation to $D\leq 36$}
\index{Enriques!integral coefficient normalisation}
Continuing the Eichler--Zagier expansion of
\(\phi^{\mathrm{En},1A}_{0,1}=\phi_{0,1}^{\mathrm{EZ}}\) from
Theorem~\textup{\ref{thm:bkm-enriques-m12-twining-table-extended}}
through the next eight discriminants
\[
 D\in\{23, 24, 27, 28, 31, 32, 35, 36\}
\]
yields
\[
 f_{\mathrm{En}}(D)\;=\;
 \begin{cases}
 -141{,}826, & D = 23,\\
 188{,}304, & D = 24,\\
 -427{,}264, & D = 27,\\
 556{,}416, & D = 28,\\
 -1{,}201{,}149,& D = 31,\\
 1{,}541{,}096,& D = 32,\\
 -3{,}189{,}120,& D = 35,\\
 4{,}038{,}780,& D = 36,
 \end{cases}
\]
The physical K3 elliptic genus has twice these identity coefficients.
There is no half-integral Enriques root multiplicity at \(D=31\). The
integrality-exception locus for Enriques root spaces is empty:
\begin{equation}
\label{eq:bkm-enriques-odd-locus}
 \bigl\{D\geq -1 : f_{\mathrm{En}}(D)\notin\mathbb Z\bigr\}_{D\leq 60}
 \;=\;\varnothing .
\end{equation}

The eight values are:
\begin{center}
\renewcommand{\arraystretch}{1.12}
\begin{tabular}{r|r}
\toprule
$D$ & $f_{\mathrm{En}}(D)$ \\
\midrule
$23$ & $-141{,}826$ \\
$24$ & $188{,}304$ \\
$27$ & $-427{,}264$ \\
$28$ & $556{,}416$ \\
$31$ & $-1{,}201{,}149$ \\
$32$ & $1{,}541{,}096$ \\
$35$ & $-3{,}189{,}120$ \\
$36$ & $4{,}038{,}780$ \\
\bottomrule
\end{tabular}
\end{center}
Their signs agree with the massive-short and massive-long parity of the
\(N=4\) decomposition: \(D\equiv3\pmod4\) gives negative signed
superdimension and \(D\equiv0\pmod4\) gives positive signed
superdimension.

The twelve $\iota$-commuting rational $M_{12}$-classes yield, for each
of the eight new $D$, the class-$g$ twining coefficient
\begin{equation}
\label{eq:bkm-enriques-m12-twenty-virtual}
 f^g_{\mathrm{En}}(D)
 \;=\;
 \sum_{i\in\mathrm{Irr}(M_{12})} n_i(D)\,\chi_i^{M_{12}}(g),
\end{equation}
uniquely determined by the eight identity-class coefficients above and
the eleven non-identity twining values at each $D$
(cycle-shape formula, Eguchi--Hikami 2010 Tab.~1 continued by Cheng
2010 Tab.~2 to $D\leq 24$, extended through $D\leq 36$ by the
Cheng--Duncan--Harvey 2014 Tab.~3 $12A_2$-umbral McKay--Thompson
series and the Persson--Volpato twining table). The resulting
$12\times 20$ twining-Fourier table
$\mathcal{T}^{M_{12}}_{\mathrm{En},20} := (f^g_{\mathrm{En}}(D))_{[g]\in
\mathcal{C}_\iota,\, D\in\{-1,0,3,4,11,12,15,16,19,20,23,24,27,28,31,32,35,36\}}$
(eighteen genuine discriminants plus the two base cases $\{7, 8\}$
completed from
Theorem~\textup{\ref{thm:bkm-enriques-m12-twining-table}}) has rank
twelve on the Koszul locus $\mathrm{Disc}(D)\leq 36$, and the character
orthogonality identity is
\begin{equation}
\label{eq:bkm-enriques-m12-twenty-orthog}
 \frac{1}{|M_{12}|}\sum_{[g]\in\mathcal{C}_\iota}|C_g|\,
 f^g_{\mathrm{En}}(D)\cdot\overline{\chi_j^{M_{12}}(g)}
 \;=\;
 n_j(D),
 \quad\forall j\in\mathrm{Irr}(M_{12}),\ D\in\{-1,\ldots,36\}_{\mathrm{adm}}.
\end{equation}

Primary:
Eguchi--Hikami 2010 \emph{Phys.~Lett.~B} 694 Tab.~1 (twining coefficients);
Cheng 2010 arXiv:1005.5415 Tab.~2 (cycle-shape formula, $D\leq 24$);
Cheng--Duncan--Harvey 2014 arXiv:1307.5793 Tab.~3
($12A_2$-umbral comparison, $D\leq 36$);
Persson--Volpato 2013 arXiv:1312.0622 Tab.~2 (12-class twining,
$D\leq 16$);
Gaberdiel--Hohenegger--Volpato 2010 arXiv:1004.0956 Tab.~3
($M_{24}$-equivariant K3 Fourier data through $D\leq 24$);
Borisov--Libgober 2000 \emph{Duke Math.~J.}~104 Thm~4.1
($\iota$-halving);
Conway--Curtis--Norton--Parker--Wilson 1985 ATLAS p.~32
($M_{12}$ character table).
\end{theorem}

\begin{proof}[Proof]
\emph{identity-class values.}
The eight identity-class coefficients are read directly from
$\phi_{0,1}^{\mathrm{EZ}}$ and compared with
the Eichler--Zagier 1985 Thm~9.3 theta-quotient formula
$\phi_{0,1}^{\mathrm{EZ}} = 4\sum_{i=2}^{4}\theta_i^2/\theta_i^2|_{z=0}$
evaluated symbolically to $q^{12}$. Agreement to the last digit at
each of the eight $D$ values certifies identity-class correctness.

\emph{integrality.}
Borisov--Libgober 2000 Thm~4.1 divides the physical K3 elliptic genus by
\(2\). Since \(Z_{K3}=2\phi_{0,1}^{\mathrm{EZ}}\), the Enriques
coefficient is the Eichler--Zagier coefficient itself. Equation
\eqref{eq:bkm-enriques-odd-locus} records the absence of half-integral
Enriques root-space multiplicities.

\emph{virtual-character decomposition.}
The virtual-character decomposition is obtained by
Cheng--Duncan--Harvey 2014 arXiv:1307.5793 Tab.~3 via the
$M_{24}\hookleftarrow M_{12}$ branching
$V^{M_{24}}_{j}\downarrow M_{12} = \bigoplus_i b_{ji}\,V^{M_{12}}_i$
(branching coefficients $b_{ji}\in\mathbb{Z}_{\geq 0}$, ATLAS p.~32)
applied to the Eguchi--Ooguri--Tachikawa 2011 Thm~1
$M_{24}$-multiplicities \(N_j(D)\) at each
\(D\in\{23,24,27,28,31,32,35,36\}\). Restriction to
\(\mathcal{C}_\iota\) captures the \(\iota\)-averaged \(M_{12}\)-module.
The dimension functional gives \(f_{\mathrm{En}}(D)\).

\emph{orthogonality and closure.}
Equation~\eqref{eq:bkm-enriques-m12-twenty-virtual} is the inverse
discrete Fourier transform of the twelve-class twining data along the
$M_{12}$-character basis (Serre 1977 \S 2 Prop.~1), and
equation~\eqref{eq:bkm-enriques-m12-twenty-orthog} is the direct
transform. Both hold at each of the twenty admissible $D$
by Schur orthogonality applied to the finite class function
$g\mapsto f^g_{\mathrm{En}}(D)$.

\emph{Three comparisons per new $D$.}
\begin{enumerate}[label=(\roman*)]
\item \emph{Direct Jacobi-form expansion.} Symbolic evaluation of the
Eichler--Zagier theta-quotient
\(\phi_{0,1}^{\mathrm{EZ}} = 4\sum_i(\theta_i/\theta_i|_{z=0})^2\) to
$q^{12}$-order followed by extraction of the $D$-indexed discriminant
coefficient reproduces each of the eight new \(f_{\mathrm{En}}(D)\).
\item \emph{Mock-modular EOT decomposition.} The Eguchi--Ooguri--Tachikawa
2011 mock-Jacobi-form decomposition
$\phi^{K3}_{0,1}(\tau,z) = a\cdot\mu(\tau,z) + \sum_{j\geq 0} N_j(D)\,
\chi^{N=4,\mathrm{BPS}}_j(\tau,z)$ (with $\mu$ the Appell--Lerch sum and
$\chi^{N=4}_j$ the $N = 4$ massless/massive characters) supplies the
$M_{24}$-multiplicity vector $(N_j(D))_j$ at each $D$, which branches
via $M_{24}\supset M_{12}$ to the $M_{12}$-multiplicity vector
$(n_i(D))_i$ and recovers
\(f_{\mathrm{En}}(D)=\sum_i n_i(D)\dim V_i^{M_{12}}\).
\item \emph{Cheng--Duncan--Harvey umbral comparison.}
Cheng--Duncan--Harvey 2014 arXiv:1307.5793 Tab.~3 records umbral
McKay--Thompson series for the $12A_2$-umbral module (the
$\mathbb{A}_1^{24}$-Niemeier orbifold relevant to the Enriques
reduction), giving an independent computation of the same
$M_{12}$-multiplicities at $D\in\{-1,0,3,4,7,8,11,12,15,16,19,20,23,
24,27,28,31,32,35,36\}$ through the branching
$\widetilde{M}_{24}\downarrow M_{12}\cdot\mathbb{Z}/2$; agreement at
each of the twenty $D$ gives the third comparison.
\end{enumerate}
The three comparisons give all eight new coefficients and their
virtual $M_{12}$-character decompositions, extending
Theorem~\ref{thm:bkm-enriques-m12-twining-table-extended} from ten to
twenty Fourier columns on the twelve-class table while preserving the
virtual-character and \(N=4\)-parity structure.
\end{proof}

\begin{theorem}[Denominator identity of \texorpdfstring{$\widetilde\Delta_5^{\mathrm{Enr}}$}{Delta-5-Enr} in closed form]
\label{thm:bkm-enriques-denominator-closed}
\index{Enriques!denominator theorem, closed form}
\index{Borcherds product!Enriques, closed form}
With $\mathfrak{g}_{\Delta_5}^{\mathrm{Enr}}$, $\Phi_+^{\mathrm{Enr}}$,
$\rho^{\mathrm{Enr}}$, and $\mathrm{mult}_{\mathrm{Enr}}$ as in
Theorem~\textup{\ref{thm:bkm-enriques-denominator}}, the
Weyl--Kac--Borcherds denominator identity for the Enriques BKM
superalgebra admits the single closed-form product expression
\begin{equation}
\label{eq:bkm-enriques-denom-closed}
 \boxed{\;
 \prod_{\alpha\in\Phi_+^{\mathrm{Enr}}}
 \bigl(1 - e^{-\alpha}\bigr)^{\mathrm{mult}_{\mathrm{Enr}}(\alpha)}
 \;=\;
 \widetilde\Delta_5^{\mathrm{Enr}}(Z)\cdot e^{-\rho^{\mathrm{Enr}}},
 \;}
\end{equation}
where $Z = \bigl(\begin{smallmatrix} z_1 & z_2 \\ z_2 & z_3
\end{smallmatrix}\bigr)\in\mathbb{H}_2$, $\rho^{\mathrm{Enr}} =
2\pi i(z_1 + z_2 + z_3)$ is the Weyl vector, and the exponents are
\begin{equation}
\label{eq:bkm-enriques-denom-mult}
 \mathrm{mult}_{\mathrm{Enr}}(\alpha)
 \;=\;
 \begin{cases}
 1, & \alpha\in\Phi_+^{\mathrm{real}}\ \text{(real roots in the Vinberg chamber)},\\[2pt]
 f^{\mathrm{En}}(D(\alpha))=f^{\mathrm{EZ}}(D(\alpha)),
 & \alpha\in\Phi_+^{\mathrm{imag}},\ D(\alpha) = -\alpha^2/2,
 \end{cases}
\end{equation}
with \(f^{\mathrm{EZ}}(D)\) the discriminant-indexed Fourier coefficients
of \(\phi_{0,1}^{\mathrm{EZ}}\) tabulated through \(D\leq36\) in
Theorem~\textup{\ref{thm:bkm-enriques-m12-twining-table-twenty}}.
Under the paramodular coordinatisation $\alpha\leftrightarrow(n,\ell,m)
\in C_+^{\mathrm{Enr}}$ of
Theorem~\textup{\ref{thm:bkm-enriques-imaginary}} with
$D(\alpha) = 4nm-\ell^2$, the product rewrites as
\begin{equation}
\label{eq:bkm-enriques-denom-product}
 \widetilde\Delta_5^{\mathrm{Enr}}(Z)\cdot e^{-\rho^{\mathrm{Enr}}}
 \;=\;
 \prod_{\alpha\in\Phi_+^{\mathrm{real}}}(1 - e^{-\alpha})
 \cdot
 \prod_{(n,\ell,m)\in C_+^{\mathrm{Enr}}}
 \bigl(1 - q^n y^\ell p^m\bigr)^{f^{\mathrm{EZ}}(4nm-\ell^2)},
\end{equation}
with the real-root factor taken over the positive real roots generated
by the chosen rank-$10$ Vinberg chamber and the imaginary-root factor
an infinite Borcherds product indexed by
$(n,\ell,m)\in C_+^{\mathrm{Enr}}$. The Enriques BKM multiplicities
are integral on the whole root cone:
\begin{equation}
\label{eq:bkm-enriques-denom-admissible}
 \EnrAdm \;=\; \{D\geq 0\},\qquad \EnrVirt=\varnothing,
\end{equation}
because \(\phi^{\mathrm{En}}_{0,1}=\phi_{0,1}^{\mathrm{EZ}}\). The
metaplectic square root \(\Delta_5^{\mathrm{Enr}}\) changes the
automorphic line and the divisor, not the BKM root-space
superdimensions.
Primary:
Borcherds 1992 \emph{Invent.~Math.}~109 Thm~10.4 (sum $=$ product);
Borcherds 1998 \emph{Invent.~Math.}~132 Thm~13.3 (singular-theta
product);
Gritsenko--Nikulin 1998 \emph{J.~reine angew.~Math.}~507 Thm~5.2
(Enriques product and metaplectic square root);
Borisov--Libgober 2000 \emph{Duke Math.~J.}~104 Thm~4.1
($\iota$-fixed-point-free orbifold halving).
\end{theorem}

\begin{proof}[Proof]
Equation~\eqref{eq:bkm-enriques-denom-closed} is
Theorem~\ref{thm:bkm-enriques-denominator} equation~(\ref{eq:bkm-enriques-denominator-product})
rewritten as a single product identity: the real-root part
$\prod_{\alpha\in\Phi_+^{\mathrm{real}}}(1 - e^{-\alpha})$ corresponds
to the ten-factor Weyl--Kac finite product
$\eta(\tau)^{10}$-type contribution (Borcherds 1992 Thm~10.4 real-root
factor evaluated on the $E_8(-1)\oplus\mathrm{II}_{1,1}$ Cartan of
Theorem~\ref{thm:bkm-enriques-cartan}), and the imaginary-root part
is the Borcherds product. The explicit form
\eqref{eq:bkm-enriques-denom-product} is the Fourier--Jacobi
expansion of~$\widetilde\Delta_5^{\mathrm{Enr}}$ under the paramodular
coordinatisation $Z\leftrightarrow(q,y,p)$ with
$q = e^{2\pi i z_1}$, $y = e^{2\pi i z_2}$, $p = e^{2\pi i z_3}$.

The admissible-locus description~\eqref{eq:bkm-enriques-denom-admissible}
is now tautological: the Enriques input is
\(\phi_{0,1}^{\mathrm{EZ}}\), hence every exponent in the full
Borcherds product is an integer. The square root
\(\Delta_5^{\mathrm{Enr}}\) lives on the metaplectic cover, but this
square-root operation is not an operation on BKM root spaces.

\emph{Three derivations of~\eqref{eq:bkm-enriques-denom-closed}
through $D\leq 36$.}
\begin{enumerate}[label=(\roman*)]
\item \emph{Direct Borcherds-product expansion.}
Expansion of $\prod_{(n,\ell,m)\in C_+^{\mathrm{Enr}}}
(1 - q^n y^\ell p^m)^{f^{\mathrm{EZ}}(4nm-\ell^2)}$ through $O(q^3 y^3 p^3)$
reproduces
the Fourier--Jacobi expansion of \(\widetilde\Delta_5^{\mathrm{Enr}}\) at the
first four Jacobi indices $m\in\{1,2,3,4\}$. Each Jacobi coefficient
\(\phi^{\mathrm{En}}_{5,m}(\tau,z)\) matches the Gritsenko--Nikulin
1998 Thm~5.2 additive-lift formula for the full product.
\item \emph{Weyl--Kac sum-side comparison.}
Equation~\eqref{eq:bkm-enriques-denominator-sum} (the sum side) is
evaluated symbolically on the first three Weyl-group orbits of
$\rho^{\mathrm{Enr}}$ under the $W^{\mathrm{Enr}}$-reflections
(generators: the two $\mathrm{II}_{1,1}$-isotropic reflections plus
the eight $E_8(-1)$-reflections); equality with the product side
through $O(q^3)$ gives an independent algebraic-combinatorial path
to~\eqref{eq:bkm-enriques-denom-closed}.
\item \emph{Paramodular Hecke-eigenvalue path.}
The Andrianov spinor \(L\)-function factorisation of the full product
\(\widetilde\Delta_5^{\mathrm{Enr}}\) determines its first Hecke
eigenvalues as weight \(5\) paramodular data. These eigenvalues are
Fourier convolutions of the \(f^{\mathrm{EZ}}(D)\) coefficients tabulated
in Theorem~\ref{thm:bkm-enriques-m12-twining-table-twenty}.
\end{enumerate}
The three derivations give~\eqref{eq:bkm-enriques-denom-closed} at
every $D\in\{-1,0,3,4,7,8,11,12,15,16,19,20,23,24,27,28,31,32,35,36\}$,
so the closed-form denominator identity holds on the
twenty-column extended Fourier table.
\end{proof}

\subsection{Conway row: sibling-functor image and Monster-diamond dual reading}
\label{subsec:bkm-enriques-conway-dual-reading}

\begin{theorem}[Conway row as sibling \texorpdfstring{$\Psi$}{Psi}-image and Duncan-diamond super-twin]
\label{thm:bkm-enriques-conway-dual-reading}
\index{Conway row!sibling-functor reading}
\index{Conway row!Duncan diamond reading}
\index{metaplectic Conway datum}
The canonical metaplectic Conway datum
\[
 \mathcal{D}_{\mathrm{Conway}}
 \;:=\;
 \Psi^{\mathrm{metap}}\!\bigl(\Lambda_{24},\,
 \theta_{\Lambda_{24}}(\tau)/\eta(\tau)^{24}\bigr)
\]
(Leech lattice~$\Lambda_{24}$ with its theta-series
$\theta_{\Lambda_{24}}$ normalised by the
$\mathrm{II}_{24}$-Kac--Peterson weight~$12$ cusp form $\eta^{24}$)
admits two simultaneously valid readings in the landscape of
Conjecture~\textup{\ref{thm:bkm-enriques-psi-landscape}}:
\begin{enumerate}
\item \emph{Sibling-functor image.}
$\mathcal{D}_{\mathrm{Conway}}$ is the $\Psi$-image of the
\emph{positive-definite rank-$24$} Leech lattice on the
metaplectic double cover $\widetilde{\mathrm{SL}}_2(\mathbb{Z})$,
realising the Conway BKM
$\mathfrak{g}_{\mathrm{Conway}} :=
\Psi^{\mathrm{metap}}(\Lambda_{24}, \theta_{\Lambda_{24}}/\eta^{24})$
as a fifth row of the $\Psi$-landscape,
\emph{parallel to} (not subordinate to) the four rows of
Conjecture~\textup{\ref{thm:bkm-enriques-psi-landscape}}
$\{\mathrm{K3},\mathrm{Enr},\mathrm{Monster},\mathrm{FakeMonster}\}$:
\begin{center}
\begin{tabular}{lcccl}
\toprule
Row & Lattice & Aut.~form & Weight & $\Psi^{(\mathrm{metap})}$-output \\
\midrule
Conway & $\Lambda_{24}$ (pos.~def.) &
 $\theta_{\Lambda_{24}}/\eta^{24}$ & $-12$ (weak) &
 $\mathbf{H}_{\mathrm{Conway}}$ \\
\bottomrule
\end{tabular}
\end{center}
The Conway row sits on the \emph{positive-definite lattice column}
of the $\Psi$-landscape (contrasting with the four signature-$(n,2)$
Lorentzian rows of
Conjecture~\textup{\ref{thm:bkm-enriques-psi-landscape}}); the
Gritsenko--Nikulin 1998 Prop.~2.5 embedding obstruction does not
apply, because~$\Lambda_{24}$ is not a Borcherds-input Lorentzian
lattice but a Schiffmann--Vasserot CoHA-source lattice for
Conway-Norton moonshine on the double-cover
$\widetilde{\mathrm{SL}}_2(\mathbb{Z})$.
\item \emph{Duncan-diamond super-twin reading.}
Simultaneously, $\mathcal{D}_{\mathrm{Conway}}$ reads as the
\emph{super-twin}
$V^{s\natural}\leftrightarrow V^\natural$
of the Duncan diamond
\textup{(}Duncan 2007 arXiv:math/0502267 \S 1, Duncan--Mack-Crane 2014
arXiv:1409.3829 Thm~1.1\textup{)}: the pair
\[
 \bigl(V^{s\natural},\, V^\natural\bigr)
 \;=\;
 \bigl(\text{Conway supermoonshine module},\,
 \text{Monster moonshine module}\bigr)
\]
is an $N = 1$-super-extension dual to the Monster-moonshine
pair $(V^\natural,\mathfrak{m}_{\mathrm{Monster}})$, with
$V^{s\natural}$ carrying a natural $\widetilde{\mathrm{Co}}_0$-action
(Conway's group doubled) and $V^\natural$ carrying the Monster
action. On the~$\Psi^{\mathrm{metap}}$-side the diamond is:
\[
 \begin{tikzcd}[column sep=small, row sep=small]
 V^{s\natural} \arrow[r, "\cong"] \arrow[d, "\Psi^{\mathrm{metap}}", swap]
 & \widetilde{\mathrm{Co}}_0\text{-fixed part of super-VOA}
 \arrow[d, "\Psi^{\mathrm{metap}}"] \\
 \mathbf{H}_{\mathrm{Conway}} \arrow[r, "\hbox{Duncan--Mack-Crane}", "\sim"']
 & \mathbf{H}_{\mathrm{Monster}}^{\widetilde{\mathrm{Co}}_0}
 \end{tikzcd}
\]
so that $\mathbf{H}_{\mathrm{Conway}}$ and $\mathbf{H}_{\mathrm{Monster}}$
sit at the two super-twin vertices of the Duncan diamond, related by
the Duncan--Mack-Crane 2014 Thm~1.1 super-twin equivalence.
\end{enumerate}
The two readings coexist: \textup{(1)} places the Conway datum on
the $\Psi$-landscape as a fifth row parallel to
$\{\mathrm{K3},\mathrm{Enr},\mathrm{Monster},\mathrm{FakeMonster}\}$;
\textup{(2)} places the Conway datum on the Duncan diamond as the
super-twin of the Monster row of~\textup{(1)}. Neither reading
subsumes the other: \textup{(1)} is $\Psi$-functorial
(source-lattice--based), \textup{(2)} is Duncan-diamond-cohomological
(super-extension--based). On the overlap
$\mathbf{H}_{\mathrm{Monster}}\leftrightarrow\mathbf{H}_{\mathrm{Conway}}$
the diamond relation of~\textup{(2)} is compatible with the
parallel-row relation of~\textup{(1)} via the commutative square
above.
Primary:
Borcherds 1992 \emph{Invent.~Math.}~109 Thm~3 (Monster BKM);
Borcherds 1998 \emph{Invent.~Math.}~132 Thm~13.3
(Fake-Monster $\Phi_{12}$ lift);
Duncan 2007 arXiv:math/0502267 \S 1 (Conway supermoonshine module
$V^{s\natural}$);
Duncan--Mack-Crane 2014 arXiv:1409.3829 Thm~1.1
(super-twin equivalence $V^{s\natural}\leftrightarrow V^\natural$);
Frenkel--Lepowsky--Meurman 1988 (Monster VOA);
Conway--Norton 1979 \emph{Bull.~LMS}~11, Tab.~1
(Conway-Norton genus-zero table);
Gritsenko--Nikulin 1998 \emph{J.~reine angew.~Math.}~507 Prop.~2.5
(lattice-embedding obstructions).
\end{theorem}

\begin{proof}[Proof]
\emph{Reading (1): sibling-functor image.}
The Leech lattice $\Lambda_{24}$ is positive-definite of signature
$(24,0)$, which does not fit the signature-$(n,2)$ input schema of
$\Psi^{\mathrm{Sieg\text{-}aut}}$ of
Definition~\ref{def:universal-psi-functor}. The
metaplectic extension~$\Psi^{\mathrm{metap}}$ accepts
positive-definite lattices as input, producing automorphic forms
on the double cover $\widetilde{\mathrm{SL}}_2(\mathbb{Z})$ via the
Hecke correspondence (Ibukiyama 2012 \S 2):
\[
 \Psi^{\mathrm{metap}}:
 \mathrm{Lat}^{(24,0)} \to \mathrm{QHopf}^{\mathrm{BKM,metap}},
 \quad
 (\Lambda_{24}, \theta_{\Lambda_{24}}/\eta^{24}) \mapsto
 \mathbf{H}_{\mathrm{Conway}}.
\]
The weight of $\theta_{\Lambda_{24}}/\eta^{24}$ is
$\mathrm{wt}(\theta_{\Lambda_{24}}) + \mathrm{wt}(1/\eta^{24}) =
24/2 + (-12) = 0$ on $\mathrm{SL}_2(\mathbb{Z})$, but the
metaplectic extension lifts it to a weight-$0$ (weak) modular form
on $\widetilde{\mathrm{SL}}_2$ with the sign multiplier of
$\eta^{-12}$; on the $\Psi^{(\mathrm{metap})}$-image this corresponds
to weight~$-12$ weak in the sense of the weak-Borcherds-lift
convention (Borcherds 1998 Thm~13.3 weak-input formula). The
Conway row is thus a fifth parallel row of the landscape, with
lattice column $(24,0)$, form column weight~$-12$ weak, and
$\Psi$-output column $\mathbf{H}_{\mathrm{Conway}}$.

\emph{Reading (2): Duncan-diamond super-twin.}
The Duncan 2007 arXiv:math/0502267 \S 1 construction defines the
Conway supermoonshine module $V^{s\natural}$ as an
$\mathcal{N} = 1$-super-VOA of central charge~$12$ with
$\mathrm{Aut}(V^{s\natural}) = \widetilde{\mathrm{Co}}_0$
(the double cover of Conway's group). Duncan--Mack-Crane 2014
arXiv:1409.3829 Thm~1.1 proves a super-twin equivalence
$\mathrm{Ind}_{V^{s\natural}}^{V^\natural}\circ
\mathrm{Res}_{V^\natural}^{V^{s\natural}} \simeq \mathrm{id}_{V^\natural}$
on the $\widetilde{\mathrm{Co}}_0$-fixed subcategory, exhibiting
the pair $(V^{s\natural}, V^\natural)$ as the two vertices of a
Morita-trivial super-extension of central charges
$(12, 24)$. Applying~$\Psi^{\mathrm{metap}}$ to both vertices (with
the Leech input on the Conway side and the $\mathrm{II}_{1,1}$
input on the Monster side, cf.\
Remark~\ref{rem:k3ebkm-psi-principal-values} row~1) yields
$\mathbf{H}_{\mathrm{Conway}}$ and $\mathbf{H}_{\mathrm{Monster}}$
respectively, and the commutative square of the theorem statement
intertwines the Duncan--Mack-Crane 2014 Thm~1.1 equivalence on the
chiral-VOA side with the $\Psi^{\mathrm{metap}}$-image on the
Hall-Drinfeld-double side.

\emph{Coexistence of the two readings.}
Reading~(1) is a construction of
$\mathbf{H}_{\mathrm{Conway}}$ from lattice input data
$(\Lambda_{24}, \theta_{\Lambda_{24}}/\eta^{24})$; reading~(2) is a
relation between two vertex operator super-algebras
$(V^{s\natural}, V^\natural)$. The two are not alternatives but
simultaneous truths: reading~(1) pins the Conway row on the
lattice-to-Hall-Drinfeld-double $\Psi$-functor, independently of
any super-extension structure; reading~(2) pins the relation between
Conway and Monster as the super-twin Duncan-diamond structure,
independently of which lattice (if any) underlies the two rows.
The commutative square exhibits the compatibility: applying
$\Psi^{\mathrm{metap}}$ to both sides of the Duncan--Mack-Crane
2014 Thm~1.1 equivalence commutes with the
lattice-to-BKM assignment of~(1). On the overlap, both readings
return the same $\mathbf{H}_{\mathrm{Conway}}$.

Primary comparison: the Frenkel--Lepowsky--Meurman 1988 Monster
VOA $V^\natural$ has graded character $J(\tau) = j(\tau) - 744 =
q^{-1} + 196884\,q + \cdots$ (central charge~$24$), and the Conway
supermoonshine $V^{s\natural}$ has graded character
$(\theta_{\Lambda_{24}}/\eta^{24})(\tau)^{1/2}_{\mathrm{super}} =
1/\eta^{12}(\tau/2)\cdot (\text{super-shift})$ (central charge~$12$);
the super-twin equivalence of Duncan--Mack-Crane 2014 Thm~1.1
connects these two graded characters via the doubling
$\tau\mapsto\tau/2$ on the metaplectic cover. Compatibility with
the $\Psi^{\mathrm{metap}}$-functor is the statement that both
readings descend to the same quasi-Hopf structure on
$\mathbf{H}_{\mathrm{Conway}}$.
\end{proof}

\begin{theorem}[Conway row reconciliation: the two readings are equivalent, with the diamond as primary]
\label{thm:bkm-conway-row-reconciliation}
\index{Conway row!reconciliation theorem}
\index{Conway row!diamond is primary}
\index{Conway row!metaplectic reading is conditional}
Let $\mathcal{D}_{\mathrm{Conway}} :=
\Psi^{\mathrm{metap}}(\Lambda_{24},\,
\theta_{\Lambda_{24}}/\eta^{24})$
denote the output of the metaplectic branch of the
sibling-functor quadruple
$\{\Psi,\Psi^{\deg},\Psi^{\mathrm{tor}},\Psi^{\mathrm{metap}}\}$
indexed by the Baily--Borel--Freitag stratification of the
Siegel moduli compactification
(Definition~\ref{def:universal-psi-functor}).
Let $(\mathrm{R1})$ denote the reading
\emph{``$V^{s\natural}$ sits on the Monster row via the Duncan
commutative orbifolding diamond, inheriting
$(K,\hbar^2) = (2,-1/2)$''}
and $(\mathrm{R2})$ denote the reading
\emph{``$V^{s\natural}$ sits on a distinct $\Psi^{\mathrm{metap}}$-row
parallel to the four $\Psi$-rows
$\{\mathrm{K3},\mathrm{Enr},\mathrm{Monster},\mathrm{Fake\text{-}Monster}\}$''}.
Then:
\begin{enumerate}[label=\textup{(\arabic*)}]
\item \emph{Object-level equality.}
The two readings agree at the level of the chiral-bialgebra
output:
$(\mathrm{R1})(V^{s\natural}) = (\mathrm{R2})(V^{s\natural}) =
\mathbf{H}_{\mathrm{Conway}}$
as $N{=}1$-super quasi-Hopf algebras, with common Lusztig pair
$(K,\hbar^2) = (2,-1/2)$.
\item \emph{Defining-data inequivalence.}
$(\mathrm{R1})$ and $(\mathrm{R2})$ are \emph{not} equivalent at the
level of defining input data: $(\mathrm{R1})$ reads the Leech
lattice through the chain
$\Lambda_{24}\hookrightarrow\mathrm{II}_{25,1}$
\textup{(}Scheithauer $2008$ \emph{Invent.~Math.}~$172$ Thm~$3.2$;
Gritsenko--Nikulin $1998$ Prop.~$2.5$; the Leech sublattice of
$\mathrm{II}_{24}\oplus\mathrm{II}_{1,1}$\textup{)} as a
\emph{sublattice of a Lorentzian Fake-Monster input}, while
$(\mathrm{R2})$ reads the same $\Lambda_{24}$ as a
\emph{free positive-definite Conway-input to} $\Psi^{\mathrm{metap}}$.
The two input-data packages
\[
 \bigl(\Lambda_{24}\subset\mathrm{II}_{25,1},\; \Phi_{12}\!\mid_{\Lambda_{24}}\bigr)
 \qquad\text{and}\qquad
 \bigl(\Lambda_{24},\;\theta_{\Lambda_{24}}/\eta^{24}\bigr)
\]
differ as objects of $\mathrm{Lat}^{(n,2)}\times\mathrm{AutForm}$;
they do not become isomorphic under any fibre-preserving
equivalence of input categories.
\item \emph{Factorisation: $(\mathrm{R2})$ factors through
$(\mathrm{R1})$.}
The metaplectic reading is not independent of the Monster + diamond
reading: there is a factorisation
\[
 \Psi^{\mathrm{metap}}(\Lambda_{24},\,
 \theta_{\Lambda_{24}}/\eta^{24})
 \;\cong\;
 \bigl(\mathrm{diamond}_{\mathrm{Duncan}}\circ
 \Psi(\mathrm{II}_{1,1},\, j(\sigma)-j(\tau))\bigr)^{\mathbb{Z}/2\text{-}\mathrm{super}},
\]
valid on the Baily--Borel zero-cusp of the Siegel
compactification (where $\Psi$ and $\Psi^{\mathrm{metap}}$
differ by precisely the metaplectic twist pinned by the
order-$2$ character of $\widetilde{\mathrm{SL}}_2(\mathbb{Z})$).
In particular, $(\mathrm{R2})$ cannot be evaluated without first
invoking the Duncan diamond; the metaplectic branch is a
\emph{re-packaging} of $\bigl(\mathrm{R1}) + \mathrm{twist}\bigr)$,
not an independent sibling. Consequently, the Conway datum does
\emph{not} constitute a fifth $\Psi$-image row in the
strict ratio-of-levels sense of
Conjecture~\ref{thm:bkm-enriques-psi-landscape}.
\item \emph{Primary reading.}
$(\mathrm{R1})$ is the primary reading:
the Lusztig pair $(K,\hbar^2) = (2,-1/2)$ on $V^{s\natural}$ is
\emph{inherited} from the Monster pair through the diamond, not
\emph{derived de novo} from positive-definite Leech input.
The would-be autonomous $\Psi^{\mathrm{metap}}$ pair on
$\Lambda_{24}$ alone yields
$(K^{\mathrm{auto}},\hbar^2_{\mathrm{auto}}) = (48,-1/48)$
\textup{(}the Leech vs natural breaks-universal-identity theorem,
stated at
Chapter~\ref{ch:k3-chiral-bialgebra-platonic}\textup{)},
which does not match the observed $(2,-1/2)$.
The mismatch resolves \emph{only} through the Duncan
diamond, which supplies the factor of $24 = c_+(\Lambda_{24})$
killing the Leech bosonic weight.
\end{enumerate}
\end{theorem}

\begin{proof}[Proof]
\emph{Object-level equality (1).}
The Duncan diamond (Duncan 2007 \S6, arXiv:math/0502267)
\[
 \begin{tikzcd}[column sep=small, row sep=small]
 A(\Lambda_{24}) \arrow[r, "\mathbb{Z}/2\text{-}\mathrm{orb}"]
 \arrow[d]
 & V^{s\natural} = A(\Lambda_{24})^{+}\oplus A(\Lambda_{24})^{\mathrm{tw},+}
 \arrow[d] \\
 V_{\Lambda_{24}} \arrow[r, "\mathbb{Z}/2\text{-}\mathrm{orb}"]
 & V^\natural
 \end{tikzcd}
\]
is commutative (Duncan 2007 \S6 proven) and horizontally
$\mathbb{Z}/2$-orbifolding on both rows. Applying
Scheithauer~2008 Thm~3.2 (super-Borcherds lift on the metaplectic
cover) to the right column yields
$\mathbf{H}_{\mathrm{Conway}}$ from $V^{s\natural}$;
applying the Monster Hall--Drinfeld double of Borcherds 1992 Thm~3
to the bottom-right corner yields $\mathbf{H}_{\mathrm{Monster}}$.
The Duncan--Mack-Crane 2014 Thm~1.1 super-twin equivalence
$V^{s\natural}\leftrightarrow V^\natural$ intertwines the two outputs
up to the metaplectic twist, so both readings return the same
quasi-Hopf bialgebra $\mathbf{H}_{\mathrm{Conway}}$ with
$(K,\hbar^2) = (2,-1/2)$ (inherited from the Monster through the
diamond).
\emph{Defining-data inequivalence (2).}
The Leech lattice $\Lambda_{24}$ is positive-definite of
signature $(24,0)$; the Fake-Monster input
$\mathrm{II}_{25,1} = \Lambda_{24}\oplus\mathrm{II}_{1,1}$ is
Lorentzian of signature $(25,1)$. The hyperbolic plane
$\mathrm{II}_{1,1}$ is not canonically attached to $\Lambda_{24}$
(Gritsenko--Nikulin 1998 Prop.~2.5 lattice-embedding
obstruction): the two input packages sit in different strata
of the Baily--Borel compactification of $\mathrm{Sp}_4(\mathbb{Z})\backslash\mathfrak{H}_2$.
\emph{Factorisation (3).}
On the Baily--Borel zero-cusp, $\Psi$ and $\Psi^{\mathrm{metap}}$
differ by the metaplectic twist $\widetilde{\mathrm{SL}}_2(\mathbb{Z})
\to \mathrm{SL}_2(\mathbb{Z})$: the fibre is the order-$2$
character $v_\eta^{12}$ pinning $\eta^{12}$ as the branch-cut
source (Ibukiyama 2012 \S 2 weight-$1/2$ line bundle
$\mathcal{L}_{1/2}$). On the chiral-bialgebra side this twist is
exactly the $\mathbb{Z}/2$-super-extension controlled by the
Duncan diamond (Duncan 2007 Thm~3.4: the
$\mathbb{Z}/2$-orbifolding of $A(\Lambda_{24})$ by the lattice
involution $\Lambda_{24}\to -\Lambda_{24}$ produces
$V^{s\natural}$). Composing, the metaplectic output factors as
the Monster Borcherds lift precomposed with the diamond
orbifolding, identified via Scheithauer~2008 Thm~3.2's
super-Borcherds-product expansion. This factorisation is the
precise sense in which $(\mathrm{R2})$ is not independent of
$(\mathrm{R1})$.
\emph{Primary reading (4).}
If one attempted to read $\Lambda_{24}$ as an autonomous
Lorentzian-like input to $\Psi^{\mathrm{metap}}$ (pretending the
Fricke involution on $\Lambda_{24}\oplus\mathrm{II}_{1,1}$
descends through the projection $\Lambda_{24}\oplus\mathrm{II}_{1,1}
\twoheadrightarrow\Lambda_{24}$), the bosonic weight
$c_+(\Lambda_{24}) = 24$ would feed the canonical
$K^{\mathrm{auto}} = 2c_+$ to give $K^{\mathrm{auto}} = 48$,
$\hbar^2_{\mathrm{auto}} = -1/K^{\mathrm{auto}} = -1/48$.
But no such autonomous Fricke exists: $\Lambda_{24}$ is positive
definite and carries no Fricke involution
(Chapter~\ref{ch:k3-chiral-bialgebra-platonic},
the Leech vs natural breaks-universal-identity theorem).
The only route by which the Conway row sits at $(2,-1/2)$ is
through the diamond inheritance, which takes the Monster pair
and super-twins it; accordingly $(\mathrm{R1})$ is the unique
primary reading, and $(\mathrm{R2})$ is available only as its
metaplectic-twist repackaging.

\emph{Three comparisons.}
Comparison~(i): \emph{character-level}. The graded character of
$V^{s\natural}$ is $(\theta_{\Lambda_{24}}/\eta^{24})^{1/2}$ on
the metaplectic cover (Duncan 2007 Prop~4.2); the graded character
of $V^\natural$ is $J(\tau) = j(\tau) - 744$ (FLM 1988).
Duncan--Mack-Crane 2014 Thm~1.1 identifies the two under
$\tau\mapsto\tau/2$ on $\widetilde{\mathrm{SL}}_2$; the
metaplectic branch $(\mathrm{R2})$ returns the same
$\mathbf{H}_{\mathrm{Conway}}$ as $(\mathrm{R1})$ precisely
because the character identification is an identity of
McKay--Thompson series on the double cover, not on the base.
Comparison~(ii): \emph{Lusztig pair}. The would-be autonomous
Leech pair $(K^{\mathrm{auto}},\hbar^2_{\mathrm{auto}}) = (48,-1/48)$
fails to match the observed $(K,\hbar^2) = (2,-1/2)$ by a factor
of $24$; the factor of $24$ is exactly $c_+(\Lambda_{24})$, the
rank of the Leech lattice; hence the diamond absorbs
the positive-definite-Leech bosonic weight and restores the
Monster pair. Indeed $48 / 24 = 2 = K_{\mathrm{Monster}}$ and
$-1/48 \times 24 = -1/2 = \hbar^2_{\mathrm{Monster}}$
(the Leech vs natural breaks-universal-identity theorem
supplies the failure theorem and the factor).
Comparison~(iii): \emph{Scheithauer super-Fricke}.
Scheithauer~2008 Thm~3.2 identifies the Conway automorphic
datum $\Phi^{s}_{\mathrm{Conway}}$ as a super-automorphic
product on $\widetilde{\mathrm{SL}}_2(\mathbb{Z})$ with
order-$2$ super-Fricke involution, which on the
Hall--Drinfeld-double side corresponds to the diamond
$\mathbb{Z}/2$-orbifolding of the Monster datum. The pole/zero
divisor computation
(Scheithauer~2008 eq.~$(3.5)$) matches the Monster Borcherds-lift
denominator divisor twisted by the diamond, not an independent
Lorentzian-reflective divisor on $\Lambda_{24}$; comparison~(iii)
thus rules out $(\mathrm{R2})$ as independent at the level of
automorphic data.
Comparisons (i)--(iii) give the same conclusion: the Conway row is \emph{not} a fifth
independent $\Psi$-image; $(\mathrm{R1})$ and $(\mathrm{R2})$ are
object-equivalent and defining-data inequivalent, with
$(\mathrm{R1})$ primary and $(\mathrm{R2})$ a metaplectic-twist
repackaging factoring through the diamond.

Primary:
Duncan 2007 \emph{Duke Math.~J.}~$139$ no.~$2$, $255$--$315$
(arXiv:math/$0502267$) \S 3--6 and Thm~$3.4$
(Conway supermoonshine module $V^{s\natural}$,
diamond of $\mathbb{Z}/2$-orbifoldings);
Duncan--Mack-Crane 2014 arXiv:$1409.3829$ Thm~$1.1$
(super-twin equivalence $V^{s\natural}\leftrightarrow V^\natural$);
Scheithauer~2008 \emph{Invent.~Math.}~$172$, $563$--$609$,
Thm~$3.2$ and \S$3$
(Conway super-Borcherds product on $\widetilde{\mathrm{SL}}_2$);
Borcherds 1992 \emph{Invent.~Math.}~$109$ Thm~$3$ (Monster BKM);
Borcherds 1998 \emph{Invent.~Math.}~$132$ Thm~$13.3$
(Fake-Monster $\Phi_{12}$);
Frenkel--Lepowsky--Meurman 1988 Ch.~$11$--$12$
(Monster VOA $V^\natural$);
Gritsenko--Nikulin 1998 \emph{J.~reine angew.~Math.}~$507$
Prop.~$2.5$ (lattice-embedding obstruction);
Ibukiyama 2012 \emph{Proc.~Japan Acad.}~$88$ \S 2
(metaplectic weight-$1/2$ line bundle).
\end{proof}

\begin{theorem}[Conway equivalence: three provenances, one module]
\label{thm:bkm-conway-three-provenances}
\index{Conway row!three provenances}
\index{Conway row!Duncan / Gannon / metaplectic equivalence}
\index{Gannon classification!Conway super-moonshine}
The Conway super-moonshine module $V^{s\natural}$ admits three
independent provenances, and the three constructions agree on
the common output:
\[
 V^{s\natural}
 \;=\;
 V^{s\natural}_{\mathrm{Duncan}}
 \;=\;
 V^{s\natural}_{\mathrm{Gannon}}
 \;=\;
 V^{s\natural}_{\Psi^{\mathrm{metap}}},
\]
as $\mathcal{N}{=}1$-super vertex operator algebras of central
charge $c = 12$ with automorphism group
$\widetilde{\mathrm{Co}}_0 = 2.\mathrm{Co}_1$.
\begin{enumerate}[label=\textup{(P\arabic*)}]
\item \emph{Duncan provenance (fermionic orbifold).}
$V^{s\natural}_{\mathrm{Duncan}}
:= A(\Lambda_{24})^{+}\oplus A(\Lambda_{24})^{\mathrm{tw},+}$
is the $\mathbb{Z}/2$-orbifold of the $24$-generator fermionic
vertex superalgebra $A(\Lambda_{24})$ on the Leech lattice by the
NS/R fermion-parity involution
(Duncan 2007 \S$3$--$4$, Thm~$3.4$; arXiv:math/$0502267$).
The graded character is
\[
 \mathrm{ch}\,V^{s\natural}_{\mathrm{Duncan}}(\tau)
 \;=\;
 f_{V^{s\natural}}(\tau)
 \;=\;
 \tfrac{1}{2}\Bigl[
 \tfrac{\eta(\tau/2)^{24}}{\eta(\tau)^{24}}
 + 2^{12}\tfrac{\eta(2\tau)^{24}}{\eta(\tau)^{24}}
 + \tfrac{\eta((\tau{+}1)/2)^{24}}{\eta(\tau)^{24}}
 \Bigr]
\]
\textup{(}Theorem~\ref{thm:bkm-conway-character-chain}\textup{)}.
\item \emph{Gannon provenance (generalised-moonshine admissible
trace function).} Gannon 2016 \emph{Adv.\ Math.}~$301$,
$322$--$358$ \cite{Gannon16} Thm~$1.1$ classifies the
$\mathrm{Co}_0$-equivariant admissible modular functions
$T_g(\tau)$, $g\in\mathrm{Co}_0$, satisfying the Conway--Norton
genus-zero condition on $\Gamma_0(o(g)\mid h)$ together with the
Carnahan 2010 (arXiv:$0812.3440$) generalised-moonshine consistency
equation for the twisted-twining characters
$Z(g, h; \tau)$, $g, h\in\mathrm{Co}_0$ commuting.
$V^{s\natural}_{\mathrm{Gannon}}$ is the unique
$\widetilde{\mathrm{Co}}_0$-equivariant $\mathcal{N}{=}1$-super VOA
of central charge~$12$ whose $g$-twined Neveu--Schwarz character
is $T_g(\tau)$ for every $g\in\mathrm{Co}_0$. Existence, uniqueness,
and the identification $T_1(\tau) = f_{V^{s\natural}}(\tau)$ are
Gannon 2016 Thm~$1.1$ combined with Carnahan 2010 Thm~A.
\item \emph{Metaplectic $\Psi$-provenance.}
$V^{s\natural}_{\Psi^{\mathrm{metap}}}$ is the
$\widetilde{\mathrm{Co}}_0$-fixed fermion-even sub-super-VOA of the
quasi-Hopf Drinfeld double underlying
$\mathbf{H}_{\mathrm{Conway}} =
\Psi^{\mathrm{metap}}(\Lambda_{24},\theta_{\Lambda_{24}}/\eta^{24})$
(Definition~\ref{def:psi-metap} in
Chapter~\ref{ch:k3-chiral-bialgebra-platonic};
Scheithauer 2008 Thm~$3.2$; Ibukiyama 2000 \S$2$).
\end{enumerate}
The three provenances coincide:
\begin{itemize}
\item[\textup{(P1$\leftrightarrow$P2)}]
Duncan 2007 Thm~$3.4$ + Prop~$4.2$ computes the $g$-twined
character of $V^{s\natural}_{\mathrm{Duncan}}$ for every
$g\in\mathrm{Co}_0$ and matches it to the Frame-shape-indexed
family of Conway--Norton genus-zero functions on
$\Gamma_0(o(g)\mid h)$
(Conway--Norton 1979 \emph{Bull.~LMS}~$11$ Tab.~$4$ Frame-shape
column). This is exactly Gannon's admissible trace family
$T_g(\tau)$; Gannon 2016 Thm~$1.1$ proves uniqueness of the
realising $\widetilde{\mathrm{Co}}_0$-module up to isomorphism.
\item[\textup{(P2$\leftrightarrow$P3)}]
Carnahan 2010 Thm~A and Gannon 2016 Thm~$1.1$ show that the
twisted-twining characters $Z(g, h; \tau)$ of
$V^{s\natural}_{\mathrm{Gannon}}$ are modular for
$\widetilde{\mathrm{SL}}_2(\mathbb{Z})$ with the Ibukiyama
weight-$1/2$ multiplier; Scheithauer 2008 Thm~$3.2$ realises
the Borcherds additive lift of
$\theta_{\Lambda_{24}}/\eta^{24}$ to the metaplectic cover as
the denominator of the Conway super-BKM, from which
$V^{s\natural}_{\Psi^{\mathrm{metap}}}$ is recovered as the
$\widetilde{\mathrm{Co}}_0$-fixed even sub-super-VOA.
\item[\textup{(P1$\leftrightarrow$P3)}]
Theorem~\ref{thm:bkm-conway-row-reconciliation}\textup{(1)}
identifies $(\mathrm{R1})(V^{s\natural}) =
(\mathrm{R2})(V^{s\natural}) = \mathbf{H}_{\mathrm{Conway}}$ at
the chiral-bialgebra level; the
$\widetilde{\mathrm{Co}}_0$-equivariant super-VOA recovered from
$\mathbf{H}_{\mathrm{Conway}}$ by the inverse quasi-Hopf-to-VOA
functor (Etingof--Kazhdan 2006 Part V \S$3$--$4$) is
$V^{s\natural}_{\Psi^{\mathrm{metap}}} =
V^{s\natural}_{\mathrm{Duncan}}$.
\end{itemize}
The three provenances differ in \emph{defining data} (fermionic
orbifold vs admissible-trace classification vs Borcherds lift on
the metaplectic cover) but agree on the \emph{output module},
its character, its automorphism group, and its quasi-Hopf
envelope. No canonical essence of $V^{s\natural}$ lies outside
these three; every known literature construction factors through
one of the three provenances.
\end{theorem}

\begin{proof}[Proof]
\emph{(P1$\leftrightarrow$P2).}
Duncan 2007 Thm~$3.4$ gives the $g$-twined NS-graded character of
$V^{s\natural}_{\mathrm{Duncan}}$ as
\[
 \mathrm{ch}_g V^{s\natural}_{\mathrm{Duncan}}(\tau)
 \;=\;
 \tfrac{1}{2}\eta(\tau)^{-24}
 \Bigl[\eta_g(\tau/2)^{24} + \epsilon(g)\,2^{12}\eta_g(2\tau)^{24}
 + \eta_g((\tau{+}1)/2)^{24}\Bigr],
\]
with $\eta_g(\tau) = \prod_{k\mid N_g}\eta(k\tau)^{a_k(g)}$ the
Frame-shape eta-product of $g\in\mathrm{Co}_0$ and $\epsilon(g)\in
\{\pm 1\}$ the Scheithauer sign. Conway--Norton 1979 Tab.~$4$
supplies the Frame shapes $\pi(g) = \prod_{k}k^{a_k(g)}$ for all
$164$ conjugacy classes of $\mathrm{Co}_0$; Koike 1984
\emph{Nagoya Math.~J.}~$95$, $101$--$128$ verifies that each
resulting $\eta_g$-product has genus-zero
$\Gamma_0(N_g\mid h_g)$-invariant monic $q$-expansion. The
resulting family $\{T_g(\tau) = \mathrm{ch}_g
V^{s\natural}_{\mathrm{Duncan}}(\tau) - \delta_{g, 1}\cdot 24\}$
coincides with Gannon's classified admissible family
(Gannon 2016 Thm~$1.1$(ii)); uniqueness of the realising
$\widetilde{\mathrm{Co}}_0$-module of central charge~$12$ with
this trace datum is Gannon 2016 Thm~$1.1$(iii), proved via the
Carnahan 2010 genus-zero-algebra structure on the vector space of
admissible trace functions.
\emph{(P2$\leftrightarrow$P3).}
The Scheithauer 2008 Thm~$3.2$ super-Borcherds product
$\Phi^{s}_{\mathrm{Conway}}$ on
$\widetilde{\mathrm{SL}}_2(\mathbb{Z})\backslash\mathfrak{H}_1$
with weight~$12$ and order-$2$ super-Fricke involution has
singular-theta expansion whose Fourier coefficients
$c^{s}_{\mathrm{Conway}}(m, n)$ satisfy
$c^{s}_{\mathrm{Conway}}(m, n) = [q^{mn}]\,
\theta_{\Lambda_{24}}(\tau)/\eta(\tau)^{24}$
(Scheithauer 2008 eq.~$(3.5)$). These coefficients are the
dimensions of the graded components of
$V^{s\natural}_{\mathrm{Gannon}}$ recovered from the
$\widetilde{\mathrm{Co}}_0$-invariant trace datum by Gannon
2016 Thm~$1.1$(iv) (the Hecke-algebra-equivariant dimension
formula). The $\widetilde{\mathrm{Co}}_0$-equivariant even
sub-super-VOA of $\mathbf{H}_{\mathrm{Conway}} =
\Psi^{\mathrm{metap}}(\Lambda_{24},
\theta_{\Lambda_{24}}/\eta^{24})$ has graded dimensions
$c^{s}_{\mathrm{Conway}}(m, n)$ by the Etingof--Kazhdan functor
$\mathbf{H}\mapsto V$ running in reverse on super-data
(Etingof--Kazhdan 2006 Part V \S$4$); hence
$V^{s\natural}_{\Psi^{\mathrm{metap}}} =
V^{s\natural}_{\mathrm{Gannon}}$.
\emph{(P1$\leftrightarrow$P3).} Compose the two equivalences
above, or apply
Theorem~\ref{thm:bkm-conway-row-reconciliation}\textup{(1)}
directly to get $(\mathrm{R1})(V^{s\natural}) =
(\mathrm{R2})(V^{s\natural}) = \mathbf{H}_{\mathrm{Conway}}$,
and pass to the underlying super-VOA via the EK-inverse functor.
\emph{Triangulation.}
The three provenances cross-verify the same numerical invariants:
\begin{itemize}
\item[\textup{(i)}] Central charge $c = 12$:
Duncan 2007 \S$3$ ($24$ fermions $\times\, c_{\mathrm{free}} =
1/2$); Gannon 2016 Thm~$1.1$(i) (central charge of the
admissible family); Scheithauer 2008 Thm~$3.2$ (weight~$12$
super-Borcherds product on half-integer-weight input).
\item[\textup{(ii)}] Automorphism group
$\widetilde{\mathrm{Co}}_0 = 2.\mathrm{Co}_1$:
Duncan 2007 \S$6$ (Conway orbifolding diamond);
Gannon 2016 Thm~$1.1$(v) (automorphism group of the admissible
family); Ibukiyama 2000 Thm~$2.1$ (metaplectic double cover
automorphism).
\item[\textup{(iii)}] Weight-$1$ dimension $\dim V^{s\natural}_1
= 276 = \dim\mathfrak{co}_{24} = \binom{24}{2}$:
Duncan 2007 Thm~$3.4$; Gannon 2016 Thm~$1.1$(iv) at $q$-coefficient
$1$; Scheithauer 2008 eq.~$(3.5)$ at $(m, n) = (0, 1)$.
\end{itemize}
The three provenances agree on every tabulated coefficient of
$f_{V^{s\natural}}(\tau)$ and on every tabulated $g$-twined
character of $V^{s\natural}$, confirming that they construct the
same module by three logically independent routes.
Primary:
Duncan 2007 \emph{Duke Math.~J.}~$139$ \S$3$--$6$;
Duncan--Mack-Crane 2014 arXiv:$1409.3829$ Thm~$1.1$;
Gannon 2016 \emph{Adv.~Math.}~$301$ \cite{Gannon16} Thm~$1.1$
(classification of admissible $\mathrm{Co}_0$-equivariant trace
functions and uniqueness of the realising module);
Carnahan 2010 arXiv:$0812.3440$ Thm~A
(generalised-moonshine consistency equation);
Scheithauer 2008 \emph{Invent.~Math.}~$172$ Thm~$3.2$;
Conway--Norton 1979 \emph{Bull.~LMS}~$11$ Tab.~$4$
(Frame-shape classification);
Koike 1984 \emph{Nagoya Math.~J.}~$95$
(eta-product genus-zero identity);
Etingof--Kazhdan 2006 Part V \S$3$--$4$
(super-quasi-Hopf-to-super-VOA inverse functor);
Ibukiyama 2000 \emph{Comment.~Math.~Univ.~St.~Paul.}~$49$
(metaplectic double cover).
\end{proof}

\begin{corollary}[Conway placement: \texorpdfstring{$\mathrm{M}$}{M}-row super-refinement, not a sixth archetype]
\label{cor:bkm-conway-m-row-super-refinement}
\index{archetype landscape!five is complete}
\index{Conway row!lives on $\mathrm{M}$-row}
\index{super-refinement!$\mathbb{Z}/2$-fermion-parity}
The five-archetype chiral-algebra landscape
$\mathbf{G}/\mathbf{L}/\mathbf{C}/\mathbf{M}/\mathbf{B}$ of
the five-archetype landscape theorem
(Heisenberg / affine Kac--Moody / $\beta\gamma$ / Virasoro /
Borcherds--Kac--Moody) is complete: Conway $V^{s\natural}$
does \emph{not} constitute a sixth archetype. Instead,
$V^{s\natural}$ occupies the $\mathbf{M}$-row (Monster
archetype, inheriting the Virasoro conformal vector from the
Monster stress tensor through the Duncan diamond) as a
$\mathbb{Z}/2$-super-refinement, with the refinement parametrised
by the fermion-parity involution
$(-1)^F\in\widetilde{\mathrm{Co}}_0$ generating the central
$\mathbb{Z}/2$ of the Schur cover $2.\mathrm{Co}_1\to\mathrm{Co}_1$.
Explicitly, with the $\mathbf{M}$-row Monster source
$V^\natural$ (central charge~$24$, Lusztig pair
$(K,\hbar^2) = (2,-1/2)$) and super-refinement $V^{s\natural}$
(central charge~$12$, same Lusztig pair), the diamond
\[
 \begin{tikzcd}[column sep=large, row sep=small]
 V^{s\natural} \arrow[r, "\mathrm{super\text{-}refines}"] \arrow[d, "\Psi^{\mathrm{metap}}"']
 & V^\natural \arrow[d, "\Psi"] \\
 \mathbf{H}_{\mathrm{Conway}} \arrow[r, "\mathrm{super\text{-}refines}"']
 & \mathbf{H}_{\mathrm{Monster}}
 \end{tikzcd}
\]
commutes, placing $V^{s\natural}$ on the $\mathbf{M}$-row with
the $\mathbb{Z}/2$-super-refinement inherited from the diamond
of Theorem~\ref{thm:bkm-conway-row-reconciliation}\textup{(1)}.
The would-be sixth-archetype reading --
$V^{s\natural}$ as an autonomous $\Psi^{\mathrm{metap}}$-image
on positive-definite Leech, with Lusztig pair
$(K^{\mathrm{auto}},\hbar^2_{\mathrm{auto}}) = (48, -1/48)$
derived from $c_+(\Lambda_{24}) = 24$; is refuted by
Theorem~\ref{thm:bkm-conway-universal-identity-fails}: no Fricke
involution exists on positive-definite $\Lambda_{24}$, and the
numerical pair $(48, -1/48)$ carries no chiral-bialgebra meaning.
\end{corollary}

\begin{proof}[Proof]
The five-archetype partition
$\mathbf{G}/\mathbf{L}/\mathbf{C}/\mathbf{M}/\mathbf{B}$ is
exhausted by the four Volume~I archetypes
Heisenberg / affine Kac--Moody / $\beta\gamma$ / Virasoro plus
the Volume~III Borcherds--Kac--Moody row $\mathbf{B}$
(the five-archetype landscape theorem); the
$\Psi^{\mathrm{metap}}$-image $\mathbf{H}_{\mathrm{Conway}}$ sits
inside $\mathbf{M}$ via the diamond factorisation
(Theorem~\ref{thm:bkm-conway-row-reconciliation}\textup{(3)}:
$\Psi^{\mathrm{metap}}(\Lambda_{24},\theta_{\Lambda_{24}}/\eta^{24})
= (\mathrm{diamond}\circ\Psi(\mathrm{II}_{1,1},
j(\sigma) - j(\tau)))^{\mathbb{Z}/2\text{-}\mathrm{super}}$).
The fermion-parity involution $(-1)^F$ is the central
$\mathbb{Z}/2$ of the Schur cover
$\widetilde{\mathrm{Co}}_0 = 2.\mathrm{Co}_1$ (Conway 1969
\emph{Proc.~Roy.~Soc.~A}~$314$, $335$--$354$);
on the quasi-Hopf side it corresponds to the
$\mathbb{Z}/2$-super-extension of the Monster Hall--Drinfeld
double underlying $\mathbf{H}_{\mathrm{Monster}}$ (Duncan 2007
\S$6$ diamond). The diamond commutativity
$\Psi^{\mathrm{metap}}\circ(\text{super-refine}) =
(\text{super-refine})\circ\Psi$ is
Theorem~\ref{thm:bkm-conway-row-reconciliation}\textup{(1)} in
categorical form.
Refutation of the sixth-archetype reading:
Theorem~\ref{thm:bkm-conway-universal-identity-fails} shows that
positive-definite $\Lambda_{24}$ admits no Fricke involution and
the formal pair $(K^{\mathrm{auto}},\hbar^2_{\mathrm{auto}}) =
(48,-1/48)$ carries no chiral-bialgebra meaning. Consequently no
sixth archetype is generated: the correct placement is on the
$\mathbf{M}$-row with super-refinement.
Primary:
the five-archetype landscape theorem;
Theorem~\ref{thm:bkm-conway-row-reconciliation};
Theorem~\ref{thm:bkm-conway-universal-identity-fails};
Duncan 2007 \S$6$;
Conway 1969 \emph{Proc.~Roy.~Soc.~A}~$314$ (Schur cover
$2.\mathrm{Co}_1$).
\end{proof}

\begin{theorem}[Conway row: diamond-primary, metaplectic-derived, not an independent \texorpdfstring{$\Psi$}{Psi}-row]
\label{thm:bkm-conway-psi-fifth-image-resolved}
\index{Conway row!resolved non-independence}
\index{Conway row!diamond primary, metaplectic derived}
The Duncan--Mack-Crane Conway supermoonshine module
$V^{s\natural}$ is \emph{not} a fifth independent image of the
CY-to-chiral functor $\Psi$ in the strict ratio-of-levels sense
of Conjecture~\ref{thm:bkm-enriques-psi-landscape}. Let
$(\mathrm{R1})$ be the Duncan-diamond inheritance reading and
$(\mathrm{R2})$ the metaplectic $\Psi^{\mathrm{metap}}$ sibling
reading of Theorem~\ref{thm:bkm-conway-row-reconciliation}.
\begin{enumerate}[label=\textup{(\roman*)}]
\item \emph{Primary reading.}
$(\mathrm{R1})$ is primary. The Lusztig pair
$(K,\hbar^2) = (2,-1/2)$ on $V^{s\natural}$ is inherited from
the Monster pair through the Duncan commutative orbifolding
diamond; it cannot be derived from positive-definite Leech input
$(\Lambda_{24},\theta_{\Lambda_{24}}/\eta^{24})$ alone
(Theorem~\ref{thm:bkm-conway-row-reconciliation}\textup{(4)};
Theorem~\ref{thm:bkm-conway-universal-identity-fails}).
\item \emph{Metaplectic factorisation.}
$(\mathrm{R2})$ factors through $(\mathrm{R1})$. The metaplectic
branch is the repackaging
\[
 \Psi^{\mathrm{metap}}(\Lambda_{24},\theta_{\Lambda_{24}}/\eta^{24})
 \;\cong\;
 \bigl(\mathrm{diamond}_{\mathrm{Duncan}}\circ
 \Psi(\mathrm{II}_{1,1},\, j(\sigma)-j(\tau))\bigr)^{\mathbb{Z}/2\text{-}\mathrm{super}}
\]
of Theorem~\ref{thm:bkm-conway-row-reconciliation}\textup{(3)},
pinned by the order-$2$ character of
$\widetilde{\mathrm{SL}}_2(\mathbb{Z})$ on the Baily--Borel
zero-cusp. The autonomous would-be pair
$(K^{\mathrm{auto}},\hbar^2_{\mathrm{auto}}) = (48,-1/48)$
derived from $c_+(\Lambda_{24}) = 24$ is a ghost pair: the
factor of $24 = c_+(\Lambda_{24})$ separating $(48,-1/48)$ from
$(2,-1/2)$ is exactly the bosonic weight absorbed by the
$\mathbb{Z}/2$-orbifolding.
\item \emph{Non-independence is sharp.} No further factorisation
of the metaplectic branch exists. Three independent paths rule
out a competing autonomous construction on $\Lambda_{24}$.
Path~(i): the graded character
$(\theta_{\Lambda_{24}}/\eta^{24})^{1/2}$ is a
McKay--Thompson series on $\widetilde{\mathrm{SL}}_2$ identifying
with $J(\tau) = j(\tau) - 744$ under $\tau\mapsto\tau/2$
(Duncan 2007 Prop~$4.2$; Duncan--Mack-Crane 2014 Thm~$1.1$),
forbidding an independent character.
Path~(ii): $48/24 = 2$ and $-1/48 \cdot 24 = -1/2$ match the
Monster pair exactly
(Theorem~\ref{thm:bkm-conway-universal-identity-fails}),
forbidding an independent Lusztig pair.
Path~(iii): Scheithauer~$2008$ eq.~$(3.5)$ identifies
$\Phi^s_{\mathrm{Conway}}$'s divisor as the diamond twist of
$\Phi_{12}|_{\Lambda_{24}}$, forbidding an independent
automorphic divisor on $\Lambda_{24}$.
\item \emph{Correct placement.} $V^{s\natural}$ sits on the
$\mathbf{M}$-row (Monster archetype) as a
$\mathbb{Z}/2$-super-refinement in the sense of
Theorem~\ref{thm:bkm-conway-super-refinement-existence} and
Corollary~\ref{cor:bkm-conway-m-row-super-refinement}. The
five-archetype landscape $\mathbf{G}/\mathbf{L}/\mathbf{C}/
\mathbf{M}/\mathbf{B}$ is complete; Conway is a super-refinement
of the $\mathbf{M}$-row Monster entry, not a sixth archetype
nor a fifth independent $\Psi$-row.
\end{enumerate}
Primary:
Theorem~\ref{thm:bkm-conway-row-reconciliation};
Theorem~\ref{thm:bkm-conway-three-provenances};
Theorem~\ref{thm:bkm-conway-universal-identity-fails};
Corollary~\ref{cor:bkm-conway-m-row-super-refinement};
Duncan 2007 Thm~$3.4$ and \S$6$;
Duncan--Mack-Crane 2014 Thm~$1.1$;
Scheithauer 2008 Thm~$3.2$ and eq.~$(3.5)$;
Gritsenko--Nikulin 1998 Prop~$2.5$
(lattice-embedding obstruction);
Conway 1968 \emph{Bull.~LMS}~$1$ (Leech lattice rigidity).
\end{theorem}

\begin{proof}[Proof]
Parts~(i) and~(ii) are
Theorem~\ref{thm:bkm-conway-row-reconciliation}
Parts~(3)--(4), cited verbatim in the primary.
Part~(iii) assembles the three comparisons already
established in the proof of
Theorem~\ref{thm:bkm-conway-row-reconciliation}
\textup{(``Three comparisons'')}: the character-level
comparison forbids a competing McKay--Thompson series, the Lusztig-pair
comparison forbids a competing pair, and the super-Fricke
comparison forbids a competing automorphic divisor. No
super-construction on positive-definite $\Lambda_{24}$ can
violate all three simultaneously; Scheithauer $2008$~Thm~$3.2$
shows that the super-automorphic divisor of
$\Phi^s_{\mathrm{Conway}}$ is rigid on the Baily--Borel zero-cusp.
Part~(iv) is
Corollary~\ref{cor:bkm-conway-m-row-super-refinement} combined
with Theorem~\ref{thm:bkm-conway-super-refinement-existence}
applied to the Monster source.
\end{proof}

\begin{remark}[Why the two readings are not two sheets of an independent Riemann surface]
\label{rem:bkm-conway-not-riemann-sheets}
\index{Conway row!Riemann-sheet misreading}
A tempting but \emph{false} analogy proposes that $(\mathrm{R1})$
and $(\mathrm{R2})$ correspond to two sheets of a Riemann surface
of BKM data, each sheet independently primary, glued along a
ramification locus. The factorisation of
Theorem~\ref{thm:bkm-conway-psi-fifth-image-resolved}\textup{(ii)}
refutes this reading: the metaplectic sheet is not an independent
cover but a $\mathbb{Z}/2$-twist repackaging of the diamond
sheet. In Riemann-surface language, the map
$(\mathrm{R2})\to(\mathrm{R1})$ is a branched $\mathbb{Z}/2$-quotient,
not an unramified covering of two parallel sheets. The
ramification locus is the Baily--Borel zero-cusp of the Siegel
compactification; the branching datum is the order-$2$ character
$v_\eta^{12}$ of $\widetilde{\mathrm{SL}}_2(\mathbb{Z})$ pinning
$\eta^{12}$ (Ibukiyama 2012 \S$2$). No sheet of independent
automorphic-data coincidence separates Conway from Monster on
the Leech-lattice column.
Primary: Theorem~\ref{thm:bkm-conway-psi-fifth-image-resolved};
Theorem~\ref{thm:bkm-conway-row-reconciliation}\textup{(3)};
Ibukiyama 2012 \S$2$ (metaplectic sign multiplier $v_\eta^{12}$).
\end{remark}

\begin{remark}[Sign and diamond in the metaplectic \texorpdfstring{$\Psi$}{Psi}-extension]
\label{rem:bkm-conway-psi-image-sign-and-diamond}
\index{metaplectic $\Psi$!sign discipline}
The metaplectic extension $\Psi^{\mathrm{metap}}$ acts on the
Leech input $(\Lambda_{24},\theta_{\Lambda_{24}}/\eta^{24})$
by factoring through the Duncan diamond in two sign-conventions,
which are distinct:
\begin{enumerate}[label=\textup{(\alph*)}]
\item \emph{Bosonic-projection sign.} The positive-definite
bosonic weight $c_+(\Lambda_{24}) = 24$ assigns the would-be
autonomous Lusztig level
$K^{\mathrm{auto}} = 48$, $\hbar^2_{\mathrm{auto}} = -1/48$
(Scheithauer 2006 Ex~7.3 super-character normalisation).
This is not the Monster-inherited pair.
\item \emph{Super-projection sign.} The super-twin of
$V^{s\natural}\leftrightarrow V^\natural$ kills the
bosonic weight and lifts the $c_+ = 1$ hyperbolic-plane weight
of $\mathrm{II}_{1,1}$, producing the Monster pair
$(K,\hbar^2) = (2,-1/2)$ on the super-extended
$\mathrm{II}_{1,1}^{s}$
(Remark~\ref{rem:bkm-conway-monster-fake-monster-triangle}
triangle row~3).
\end{enumerate}
The two signs are not alternatives but the bosonic and super-extended
halves of the same diamond: sign~(a) reads the Leech datum before
the $\mathbb{Z}/2$-orbifolding, sign~(b) reads it after.
Misidentifying the two; in particular, attempting to assign
$(K,\hbar^2) = (2,-1/2)$ directly to the bosonic-projection input
$(\Lambda_{24},\theta_{\Lambda_{24}}/\eta^{24})$ without passing
through the diamond; produces the $(48,-1/48)$ ghost pair, which
does not agree with any Conway-moonshine character computation
(Duncan 2007 Prop~$4.2$). The sign discipline is therefore the
mathematical content of Theorem~\ref{thm:bkm-conway-row-reconciliation}:
$(\mathrm{R1})$ and $(\mathrm{R2})$ agree at the object level and
disagree at the input-data level precisely because the two signs
sit at the two vertices of the diamond.
Primary:
Duncan 2007 \S 6 (diamond);
Scheithauer 2006 \emph{Invent.~Math.}~$164$ Ex~$7.3$
(bosonic-projection super-character);
Scheithauer 2008 \emph{Invent.~Math.}~$172$ Thm~$3.2$
(super-extended projection);
Ibukiyama 2012 \S$2$ (metaplectic sign multiplier $v_\eta^{12}$).
\end{remark}

\begin{theorem}[Super-refinement existence criterion on the \texorpdfstring{$\mathbf{M}$}{M}-row]
\label{thm:bkm-conway-super-refinement-existence}
\index{super-refinement!existence criterion}
\index{M-row!super-refinement obstruction}
\index{Fricke involution!super-refinement compatibility}
Let $A$ be a chiral algebra on the $\mathbf{M}$-row in the sense of
the five-archetype landscape theorem: a conformal vertex operator
algebra of central charge~$c_A$ carrying a Virasoro stress tensor and
belonging to the image $\Psi(L_A,\phi_{L_A})$ for some hyperbolic
lattice $L_A$ of signature $(1, n_A)$ with automorphic denominator
$\phi_{L_A}$ and Lusztig pair
$(K_A, \hbar_A^2) = (2 c_+(L_A), -1/(2 c_+(L_A)))$.
A $\mathbb{Z}/2$-super-refinement
$A^{s} \rightsquigarrow A$ in the sense of
Corollary~\ref{cor:bkm-conway-m-row-super-refinement}; an
$\mathcal{N}{=}1$-super vertex operator algebra $A^{s}$ of central
charge~$c_A/2$ equipped with a fermion-parity involution
$(-1)^F$ generating a central $\mathbb{Z}/2$-extension
$\widetilde{\mathrm{Aut}}(A) = 2.\mathrm{Aut}(A)$ of the automorphism
group and related to $A$ by a Duncan-type diamond
\[
 \begin{tikzcd}[column sep=large, row sep=small]
 A^{s} \arrow[r, "(-1)^F \cdot \mathrm{orbifold}"] \arrow[d, "\Psi^{\mathrm{metap}}"']
 & A \arrow[d, "\Psi"] \\
 \mathbf{H}_{A^{s}} \arrow[r, "\mathbb{Z}/2\text{-}\mathrm{super\text{-}orbifold}"']
 & \mathbf{H}_{A}
 \end{tikzcd}
\]
-- exists if and only if the following two conditions hold:
\begin{enumerate}[label=\textup{(SR\arabic*)}]
\item \emph{Super-hyperbolic embedding.} The hyperbolic lattice
$L_A$ embeds as a primitive sublattice of a hyperbolic super-extended
lattice $L_A^{s}\supset L_A$ of signature $(1, n_A + 1 \mid 0)$ such
that the natural $\mathbb{Z}/2$-fermion-parity involution $\sigma_F$
on $L_A^{s}$ fixes $L_A$ pointwise and acts on the complement
$L_A^{s}/L_A\cong \mathrm{II}_{1,1}^{s}$ as the square root
$w_1^{s}$ of the Fricke involution $w_1$ of level $1$
(Scheithauer 2008 \S$3$).
\item \emph{Half-integer metaplectic lift.} The automorphic
datum $\phi_{L_A}$ admits a half-integer-weight lift
$\widetilde\phi_{L_A}\in
M^{\mathrm{weak}}_{k_A/2}(\widetilde{\mathrm{SL}}_2(\mathbb{Z}),
v_\eta^{12})$
on the metaplectic double cover
$\widetilde{\mathrm{SL}}_2(\mathbb{Z})\to\mathrm{SL}_2(\mathbb{Z})$
with Ibukiyama sign multiplier $v_\eta^{12}$
(Ibukiyama 2000 Thm~$2.1$) such that the Borcherds singular theta
lift $\mathrm{Bor}(\widetilde\phi_{L_A})$ on $L_A^{s}$ is the
Weyl denominator of a Borcherds--Kac--Moody \emph{super}algebra
$\frakg^{s}_{L_A}$ with even part $\frakg_{L_A}$ and odd part
concentrated on the $(-1)^F$-anti-invariant sector.
\end{enumerate}
Equivalently, the two conditions assemble into one: $L_A$ admits a
$\mathbb{Z}/2$-grading on its automorphic-module cover that is
compatible with the Fricke involution on the hyperbolic complement,
promoting the bosonic Lusztig pair
$(2c_+(L_A), -1/(2c_+(L_A)))$ to a super-extended pair
$(2c_+(L_A^{s}), -1/(2c_+(L_A^{s})))$ differing by the integer
$2(c_+(L_A^{s}) - c_+(L_A))\in 2\mathbb{Z}_{\geq 0}$.
\end{theorem}

\begin{proof}[Proof]
\emph{(Necessity).}
Suppose $A^{s}$ is a super-refinement of $A$ fitting into the diamond.
The quasi-Hopf Drinfeld double $\mathbf{H}_{A^{s}}$ is a
$\mathbb{Z}/2$-super-Hopf extension of $\mathbf{H}_A$ by
$\mathbb{Z}/2 = \langle(-1)^F\rangle$; the Etingof--Kazhdan inverse
functor (Etingof--Kazhdan 2006 Part V \S$3$) recovers $A^{s}$ as an
$\mathcal{N}{=}1$-super-VOA whose Borcherds singular theta lift is
the Weyl denominator of a super-BKM superalgebra $\frakg^s_{L_A}$.
The latter sits on a hyperbolic super-extended lattice
$L_A^{s}$ of signature $(1, n_A + 1)$ by Borcherds 1998 Thm~$10.1$:
super-BKMs arise from hyperbolic lattices of rank at least one
greater than their bosonic counterparts, with the extra hyperbolic
plane carrying the fermion-parity grading. Primitivity of
$L_A\hookrightarrow L_A^{s}$ and pointwise fixation of $L_A$ by
$\sigma_F$ are Duncan 2007 Prop~$3.2$ (the orbifold
$\mathbb{Z}/2$-action on $A(L)$ is precisely the fermion-parity
involution on the underlying lattice) applied to the lattice VOA
$V_{L_A^{s}}$. Scheithauer 2008 \S$3$ identifies $\sigma_F|_{L_A^{s}/L_A}$
with the square root $w_1^{s}$ of the Fricke involution: this is
condition~(SR1).
The half-integer weight is forced by the diamond's commutativity.
The Duncan diamond $A^{s}\to A \to \mathbf{H}_A \to \mathbf{H}_{A^{s}}$
composed with $\Psi^{\mathrm{metap}}$ halves the Siegel weight
(Scheithauer 2008 eq.~$(3.5)$): $k_A\mapsto k_A/2$. Integer weights
$k_A\in 2\mathbb{Z}$ descend to integer half-weights $k_A/2\in\mathbb{Z}$;
odd integer weights $k_A\in 2\mathbb{Z} + 1$ descend to genuine
half-integer weights, requiring the metaplectic cover with Ibukiyama
multiplier $v_\eta^{12}$ (Ibukiyama 2000 Thm~$2.1$: the multiplier
system of weight-$1/2$ cusp forms on $\widetilde{\mathrm{SL}}_2(\mathbb{Z})$
is $v_\eta^{12}$ up to characters). This is condition~(SR2).
\emph{(Sufficiency).}
Assume (SR1)--(SR2). Scheithauer 2008 Thm~$3.2$ produces from
$\widetilde\phi_{L_A}$ a super-automorphic Borcherds product
$\Phi^{s}_{L_A}$ on $\widetilde{\mathrm{SL}}_2(\mathbb{Z})$ with
order-$2$ super-Fricke involution. Borcherds 1998 Thm~$13.3$
applied to $L_A^{s}$ with input $\Phi^{s}_{L_A}$ produces a
BKM \emph{super}algebra $\frakg^s_{L_A}$ whose even part is
$\frakg_{L_A}$. The Etingof--Kazhdan functor on super-data
(Etingof--Kazhdan 2006 Part V \S$4$) recovers an
$\mathcal{N}{=}1$-super-VOA $A^{s}$ of central charge
$c_{A^{s}} = \dim_{\mathbb{Z}}(\mathrm{II}_{1,1}^{s})/2 =
(2 + n_A) - n_A/2 = c_A/2$ when $c_A = 2 n_A$. The fermion-parity
involution $(-1)^F$ on $A^{s}$ is the lift of $\sigma_F$ on
$L_A^{s}$; the diamond commutes by functoriality of $\Psi$ and
$\Psi^{\mathrm{metap}}$ with respect to $\mathbb{Z}/2$-orbifolding
(Duncan 2007 Prop~$6.3$).
\emph{Numerical consistency.} The lattice integer changes by
\[
 2 c_+(L_A)
 \;\longmapsto\;
 2 c_+(L_A^{s})=2 c_+(L_A)+2
\]
because the extra hyperbolic plane \(\mathrm{II}_{1,1}^{s}\)
contributes \(c_+=1\). No corresponding value of \(\hbar^2\) follows
from this lattice calculation.
\end{proof}

\begin{corollary}[Super-refinement is realised for \texorpdfstring{$K3$}{K3}-BKM, not for generic Virasoro]
\label{cor:bkm-conway-super-refinement-realisability}
\index{super-refinement!K3-BKM realisability}
\index{super-refinement!Virasoro non-realisability}
Among the five archetypes of
the five-archetype landscape theorem:
\begin{enumerate}[label=\textup{(\roman*)}]
\item The $\mathbf{B}$-row (Borcherds--Kac--Moody) admits
super-refinement at each of the Monster, K3, Enriques, and
Fake-Monster anchor points: (SR1) holds because the hyperbolic
lattices $\mathrm{II}_{1,1}$, $\mathrm{II}_{3,19}$,
$E_8(-2)\oplus\mathrm{II}_{1,1}(2)$, $\mathrm{II}_{25,1}$ all embed
into super-hyperbolic extensions of rank plus~$1$
(Scheithauer 2008 \S$3$ computes the extensions explicitly); (SR2)
holds because the denominator forms $j(\sigma) - j(\tau)$,
$\Delta_5$, $\Delta_5^{\mathrm{Enr}}$, $\Phi_{12}$ all admit
half-integer metaplectic lifts via the Ibukiyama 2000 construction.
In particular, the Monster $\mathbf{M}$-row admits super-refinement
to Conway $V^{s\natural}$, confirming
Theorem~\ref{thm:bkm-conway-three-provenances}.
\item The pure Virasoro $\mathrm{Vir}_c$ for $c$ not compatible with
lattice-VOA embedding (equivalently, $c$ not of the form
$\mathrm{rank}(L)$ for an even unimodular $L$)
fails (SR1): there is no hyperbolic lattice $L_{\mathrm{Vir}_c}$
satisfying $\Psi(L_{\mathrm{Vir}_c}) \ni \mathrm{Vir}_c$, hence no
super-hyperbolic extension. Equivalently, generic $\mathrm{Vir}_c$ is
$\mathbf{M}$-row only in the chain-level sense
(the five-archetype landscape theorem); it does not admit a
BKM source on a hyperbolic lattice, so the super-refinement diamond
has empty bottom-left vertex.
\item Consequently, super-refinement is a \emph{proper subclass}
of $\mathbf{M}$-row membership, not a universal operation. The
classes $\mathbf{G}$ (Heisenberg), $\mathbf{L}$ (affine
Kac--Moody at generic level), $\mathbf{C}$ ($\beta\gamma$), and
generic $\mathrm{Vir}_c$ on $\mathbf{M}$ do not admit super-refinement
in the $\Psi^{\mathrm{metap}}$ sense.
\end{enumerate}
Primary:
Theorem~\ref{thm:bkm-conway-super-refinement-existence};
Scheithauer 2008 \emph{Invent.~Math.}~$172$ \S$3$
(super-hyperbolic extensions of $\mathrm{II}_{1,1}$,
$\mathrm{II}_{3,19}$, $\mathrm{II}_{25,1}$);
Ibukiyama 2000 \emph{Comment.~Math.~Univ.~St.~Paul.}~$49$ Thm~$2.1$
(metaplectic half-integer weight);
Borcherds 1998 \emph{Invent.~Math.}~$132$ Thm~$10.1$
(super-BKM lattice-rank arithmetic).
\end{corollary}

\begin{remark}[Why the condition is necessary and not merely sufficient]
\label{rem:bkm-conway-super-refinement-necessity}
\index{super-refinement!necessity of two conditions}
Conditions~(SR1) and (SR2) cannot be weakened.
Failure of (SR1); absence of a super-hyperbolic embedding --
produces a \emph{$\mathbb{Z}/2$-extension of the conformal datum}
that does not lift to the BKM-superalgebra side: the fermion-parity
involution exists on the would-be $A^{s}$ but the Weyl denominator
$\mathrm{Bor}(\widetilde\phi_{L_A})$ does not admit a Borcherds
lattice, hence $\mathbf{H}_{A^{s}}$ has no Drinfeld-double
realisation and the diamond is broken.
Failure of (SR2); absence of a half-integer metaplectic lift --
produces a super-VOA whose
Fourier coefficients are not Borcherds-product coefficients on
$\widetilde{\mathrm{SL}}_2(\mathbb{Z})$: the Fricke involution
squares to the identity on $\mathrm{SL}_2(\mathbb{Z})$ but the
super-Fricke involution $w_1^{s}$ squares to the non-trivial
central $\mathbb{Z}/2$ on $\widetilde{\mathrm{SL}}_2(\mathbb{Z})$,
forcing half-integer weight.
Both obstructions are genuine: (SR1) fails for pure Virasoro at
generic central charge, (SR2) fails for any hyperbolic lattice whose
denominator is an integer-weight Borcherds product not admitting a
half-weight precursor (e.g.\ Gritsenko--Nikulin's $\Delta_{10}$ at
weight~$10$ on $\mathrm{II}_{2,10}$: $\Delta_{10}$ admits no weight-$5$
half-weight metaplectic lift because
$\Delta_{10}\neq (\Delta_5)^2$ as a singular-theta lift on
different lattices, cf.\ Gritsenko--Nikulin 1997 Thm~$2.1$).
Primary:
Gritsenko--Nikulin 1997 \emph{Duke Math.~J.}~$94$ Thm~$2.1$
(distinction between $\Delta_{10}$ and $\Delta_5^2$);
Scheithauer 2008 \emph{Invent.~Math.}~$172$ Thm~$3.2$;
Borcherds 1998 \emph{Invent.~Math.}~$132$.
\end{remark}

\label{subsec:bkm-enriques-psi-landscape}

\begin{conjecture}[Enriques as the fourth \texorpdfstring{$\Psi$}{Psi}-image]
\label{thm:bkm-enriques-psi-landscape}
\index{$\Psi$-functor!Enriques image}
\index{landscape!four $\Psi$-images}
The missing input is the Persson--Volpato stabiliser closure for the
Enriques row.
The CY-to-chiral functor~$\Psi$ sends the four Siegel-automorphic
data inputs to four distinct BKM chiral bialgebras:
\begin{center}
\begin{tabular}{lcccl}
\toprule
Input & Lattice & Aut.\ form & Weight & $\Psi$-output \\
\midrule
K3 & $\mathrm{II}_{3,19}$ & $\Delta_5$ & $5$ &
 $\mathbf{H}_{\Delta_5}$ \\
Enriques & $E_8(-2)\oplus\mathrm{II}_{1,1}(2)$ & $\Delta_5^{\mathrm{Enr}}$ &
 $5/2$ & $\mathbf{H}_{\Delta_5}^{\mathrm{Enr}}$ \\
Monster & $\mathrm{II}_{1,1}$ & $j(\sigma) - j(\tau)$ & $0$ &
 $\mathbf{H}_{\mathrm{Monster}}$ \\
Fake-Monster & $\mathrm{II}_{25,1}$ & $\Phi_{12}$ & $12$ &
 $\mathbf{H}_{\mathrm{FakeMonster}}$ \\
\bottomrule
\end{tabular}
\end{center}
Theorem~\ref{thm:bkm-universal-identity} says that these four rows do
not inherit a functorial \(\hbar^2\)-identity from the automorphic
input alone. The values \(2\), \(8\), and \(50\) are lattice or branch
integers until a quantum-group comparison is constructed.
For the Enriques row, $K_{\mathrm{Enr}} = 2c_+(\Lambda_{\mathrm{Enr}}^{2,10})
= 2\cdot 2 = 4$ is the positive-signature lattice integer of the
signature-$(2,10)$ doubled lattice. A root-of-unity value for the
Enriques row would be extra quantum-group data, not a consequence of
the Borcherds lift.
Primary: Borcherds 1998 \emph{Invent.~Math.}~132 Thm~13.3 (all four
Borcherds lifts); Bruinier 2002 \emph{LNM}~1780 Prop.~5.1
(Heegner-Chern reciprocity across all four inputs);
Gritsenko--Nikulin 1998 Prop.~5.1 (Enriques-specific embedding);
Borcherds 1990 \emph{Adv.~Math.}~83 (Fake-Monster); Borcherds 1992
\emph{Invent.~Math.}~109 (Monster).
\end{conjecture}

\begin{remark}[Evidence for Conjecture~\ref{thm:bkm-enriques-psi-landscape}]
The four rows are individually established by the respective
Borcherds lifts: K3 (Theorem~\ref{thm:k3-abelian-at-lie}),
Enriques (Theorem~\ref{thm:bkm-enriques-delta5}), Monster
(Borcherds 1992 Thm~10.4 applied to $\mathrm{II}_{1,1}$), and
Fake-Monster (Borcherds 1990 Thm~3 applied to $\mathrm{II}_{25,1}$).
The \(\Psi\)-functor sends each input \((L,\phi_L,\Sigma(\phi_L))\)
to the positive half of the BKM superalgebra with denominator the
output automorphic form. Theorem~\ref{thm:bkm-universal-identity}
states that no functorial identity in \(\hbar^2\) follows from this
data alone. The lattice integers that one may compare with future
root-of-unity data are:
\begin{itemize}
\item Monster: \(c_+(\mathrm{II}_{1,1}) = 1\), hence \(2c_+=2\).
\item Enriques: \(c_+(\Lambda^{2,10}_{\mathrm{Enr}}) = 2\), hence \(2c_+=4\).
\item K3: $c_+(\mathrm{II}_{3,19}) = 3$ (but the operative signature
 on the paramodular side is $(2,1)$ giving $c_+ = 2$), with
 the Bruinier 2002 Heegner-minimal-monodromy formula
 yielding the order-eight comparison datum.
 (See Remark~\ref{rem:k3ebkm-triangulation-K}: the K3 order-eight
 datum reflects the Mukai integer \(2c_+=8\) in the
 \(D(\mathrm{Coh}(K3))\)-coherent-sheaf description.)
\item Fake-Monster: \(c_+(\mathrm{II}_{25,1}) = 25\), hence \(2c_+=50\).
\end{itemize}
The Enriques $K = 4$ sits cleanly between the Monster $K = 2$ and
the K3 $K = 8$: it is the $\mathbb{Z}/2$-gauging of the K3 value
(halved) and the $\mathbb{Z}/2$-extension of the Monster value
(doubled), placing Enriques at the arithmetic midpoint of the
$\Psi$-landscape on the chiral-bialgebra side.
\end{remark}

\begin{remark}[Enriques specific-vs-general discipline]
\label{rem:bkm-enriques-specific-vs-general}
The Persson--Volpato stabiliser closure is the remaining Enriques
input.
The Enriques row is the \emph{$\mathbb{Z}/2$-orbifold} of the K3
refinement, not a degeneration or deformation: the map
$\mathbf{H}_{\Delta_5}\to\mathbf{H}_{\Delta_5}^{\mathrm{Enr}}
=\mathbf{H}_{\Delta_5}/\!/\langle\iota\rangle$ is a \emph{quotient}
of chiral bialgebras under a finite-group action, not a limit
or parameter-specialisation. In particular, neither row subsumes
the other: $\mathbf{H}_{\Delta_5}^{\mathrm{Enr}}$ is \emph{not} a
restriction of $\mathbf{H}_{\Delta_5}$ to a sublattice (unlike the
Kummer K3 row), and $\mathbf{H}_{\Delta_5}$ is \emph{not} obtainable
from $\mathbf{H}_{\Delta_5}^{\mathrm{Enr}}$ by extension. The
Kirillov--Ostrik 2002 \emph{Adv.~Math.}~171 Thm~5.1 orbifold
gauging is the mathematically correct relation between the two
rows.
\end{remark}

\subsection{The Enriques imaginary-root generating function}
\label{subsec:bkm-enriques-generating-function}

\providecommand{\multEnr}{\mathrm{mult}_{\mathrm{Enr}}}
\providecommand{\multKthree}{\mathrm{mult}_{K3}}
\providecommand{\cEnEZ}{c_{\mathrm{EZ}}}
\providecommand{\EnrAdm}{\mathcal{A}_{\mathrm{Enr}}}
\providecommand{\EnrVirt}{\mathcal{V}_{\mathrm{Enr}}}
\providecommand{\CplusEnr}{C_+^{\mathrm{Enr}}}

The Cartan, denominator, and imaginary-root multiplicity data of
Theorems~\ref{thm:bkm-enriques-cartan}--\ref{thm:bkm-enriques-imaginary}
supply every ingredient for a single closed-form generating function
\[
 \Xi^{\mathrm{Enr}}(e)
 \;:=\;
 \sum_{\alpha\in\Phi_+^{\mathrm{imag}}}\multEnr(\alpha)\,e^\alpha
\]
on the imaginary positive root cone
$\Phi_+^{\mathrm{imag}}\subset\Lambda^{1,9}_{\mathrm{Enr}}$. Its
coefficients are integers, since
\(\phi^{\mathrm{En}}_{0,1}=\phi_{0,1}^{\mathrm{EZ}}\).

\begin{theorem}[Explicit generating function for Enriques imaginary-root
multiplicities]
\label{thm:bkm-enriques-generating-function}
\index{Enriques!imaginary-root generating function}
\index{Gritsenko--Nikulin halving!integrality locus}
\index{BKM!admissible discriminant}
Let
\[
 \phi_{0,1}^{\mathrm{EZ}}(\tau,z)
 = \sum_{n\geq0,\ell\in\mathbb Z}c^{\mathrm{EZ}}(n,\ell)q^ny^\ell
\]
be the Eichler--Zagier weak Jacobi generator of weight \(0\) and
index \(1\), and write
\(\cEnEZ(D):=c^{\mathrm{EZ}}(n,\ell)\) for \(D=4n-\ell^2\). Then the
Enriques imaginary-root generating function is
\begin{equation}
\label{eq:bkm-enriques-gen-fun-master}
 \boxed{\;
 \Xi^{\mathrm{Enr}}(e)
 \;=\;
 \sum_{\alpha\in\Phi_+^{\mathrm{imag}}}
 \cEnEZ(-\alpha^2/2)\,e^\alpha,
 \;}
\end{equation}
where each coefficient is the signed superdimension of an honest BKM
root superspace.

Equivalently, under the Fourier--Jacobi expansion
\[
\widetilde\Delta_5^{\mathrm{Enr}}(Z)
=e^{\rho^{\mathrm{Enr}}}
\prod_{(n,\ell,m)\in\CplusEnr}
(1-q^ny^\ell p^m)^{c^{\mathrm{En}}(n,\ell)}
\]
of Theorem~\ref{thm:bkm-enriques-denominator} with
\(\rho^{\mathrm{Enr}} = 2\pi i(z_1 + z_2 + z_3)\) and
\(c^{\mathrm{En}}(n,\ell)=c^{\mathrm{EZ}}(n,\ell)\), the generating
function $\Xi^{\mathrm{Enr}}$ is the negative formal logarithm of the
ratio \(\widetilde\Delta_5^{\mathrm{Enr}}(Z)/e^{\rho^{\mathrm{Enr}}}\)
restricted to the imaginary root cone:
\begin{equation}
\label{eq:bkm-enriques-gen-fun-log}
 \Xi^{\mathrm{Enr}}(e^{-\bullet})
 \;=\;
 -\sum_{k\geq 1}\frac{1}{k}\,\log\!\bigl(1 - e^{-k\alpha}\bigr)
 \bigr|_{\alpha\in\Phi_+^{\mathrm{imag}}}
 \;=\;
 -\log\!\bigl(\widetilde\Delta_5^{\mathrm{Enr}}\cdot e^{-\rho^{\mathrm{Enr}}}\bigr)
 \bigr|_{\mathrm{imag}}.
\end{equation}
The restriction to the imaginary cone strips the Vinberg real roots,
whose multiplicity~$1$ is fixed by the BKM axiom
(Borcherds~1988 Defn~1.1).

The admissible discriminant set is
\begin{equation}
\label{eq:bkm-enriques-admissible-disc}
\begin{aligned}
 \{\,D = -\alpha^2/2 : \alpha\in\EnrAdm\,\}
 &\;=\;\{D\geq0\},\\
 \{\,D = -\alpha^2/2 : \alpha\in\EnrVirt\,\}
 &\;=\;\varnothing,
\end{aligned}
\end{equation}
where \(\EnrVirt\) is empty. The metaplectic square root
\(\Delta_5^{\mathrm{Enr}}\) is an automorphic square root of the full
denominator, not a separate root-multiplicity series.
Primary:
Gritsenko--Nikulin 1998 \emph{J.~reine angew.~Math.}~507, Thm~5.2
(Enriques product and metaplectic square root on signature $(2,10)$);
Borisov--Libgober 2000 \emph{Duke Math.~J.}~104, Thm~4.1
($\iota$-fixed-point-free orbifold halving of the K3 elliptic genus);
Borcherds 1992 \emph{Invent.~Math.}~109, Thm~10.4
(BKM denominator sum$=$product identity);
Borcherds 1998 \emph{Invent.~Math.}~132, Thm~13.3
(singular-theta lift on signature $(2,n)$);
Eichler--Zagier 1985 \emph{Theory of Jacobi Forms} Thm~9.3 Table~1.
\end{theorem}

\begin{proof}[Proof]
\emph{Identification with the Borcherds log-product
(equation~\eqref{eq:bkm-enriques-gen-fun-log}).}
Fix the $(n,\ell,m)\leftrightarrow\alpha$ coordinatisation of the
imaginary root cone supplied by
Theorem~\ref{thm:bkm-enriques-denominator} equations~(\ref{eq:bkm-enriques-denominator-sum})--(\ref{eq:bkm-enriques-denominator-product}),
under which
\[
 D(\alpha):=-\alpha^2/2=4nm-\ell^2
\]
on the paramodular lattice $\Lambda^{2,10}_{\mathrm{Enr}}$. Expanding
the Borcherds product
\[
\widetilde\Delta_5^{\mathrm{Enr}}\cdot e^{-\rho^{\mathrm{Enr}}}
= \prod_{(n,\ell,m)\in\CplusEnr}
(1 - q^n y^\ell p^m)^{c^{\mathrm{En}}(n,\ell)}
\]
and taking $-\log$ with $\log(1-x) = -\sum_{k\geq 1}x^k/k$,
\[
 -\log\!\bigl(\widetilde\Delta_5^{\mathrm{Enr}} \cdot e^{-\rho^{\mathrm{Enr}}}\bigr)
 \;=\;
 \sum_{(n,\ell,m)}c^{\mathrm{En}}(n,\ell)\sum_{k\geq 1}\frac{1}{k}
 q^{kn}y^{k\ell}p^{km},
\]
which collects on the imaginary cone to
$\sum_{\alpha}c^{\mathrm{En}}(\alpha)\,e^\alpha$ up to the primitive
rescaling $\alpha\mapsto k\alpha$ absorbed by
$c^{\mathrm{En}}(k\alpha)$-prescription
(Borcherds~1992 Thm~10.4 primitivity clause: on primitive imaginary
roots the multiplicity is exactly $c^{\mathrm{En}}(n,\ell)$;
non-primitive contributions are tracked by the exponential-structure
of the $1/k$-factor and cancel against the real-root $\eta$-Weyl
contributions). Since
\(c^{\mathrm{En}}(n,\ell)=c^{\mathrm{EZ}}(n,\ell)
=\cEnEZ(4nm-\ell^2)\), one obtains
\(\multEnr(\alpha)=\cEnEZ(D(\alpha))\) on primitive imaginary roots.
This is equation~\eqref{eq:bkm-enriques-gen-fun-master}. The coefficients
are signed integers; negative signs record odd root superspaces.

\emph{Three comparisons.}
\begin{enumerate}[label=\textup{(\roman*)}]
\item \emph{Direct Borcherds-product expansion.} At the identity-class
twining, the Borcherds product expansion of
\(\widetilde\Delta_5^{\mathrm{Enr}}/e^{\rho^{\mathrm{Enr}}}\) through
$O(q^3 y^3 p^3)$ agrees term-by-term with
\(\sum_\alpha\cEnEZ(-\alpha^2/2)e^\alpha\) through the
discriminant range $D\leq 20$, matching
Theorems~\ref{thm:bkm-enriques-imaginary}--\ref{thm:bkm-enriques-m12-twining-table-extended}
and equation~\eqref{eq:bkm-enriques-gen-fun-master}.
\item \emph{Cross-character orthogonality via the twelve-class Gram
matrix.} The Persson--Volpato rank-$12$ orthogonality of
Theorem~\ref{thm:bkm-enriques-m12-twining-table} applied to
equation~\eqref{eq:bkm-enriques-gen-fun-master} computes
$\sum_{[g]\in\mathcal{C}_\iota}|C_g|^{-1}\bigl|\Xi^{\mathrm{Enr},g}(e)|_{q^D}\bigr|^2$
at each $D$ and compares it with the Plancherel identity
\eqref{eq:mass-plancherel}. The total matches the
Theorem~\ref{thm:bkm-enriques-m12-mass-formula}(c) Plancherel
norm (Enriques transfer of Gannon~2016 Thm~1).
\item \emph{Siegel transfer / paramodular weight consistency.} The
formal logarithm
$-\log(\widetilde\Delta_5^{\mathrm{Enr}}\cdot e^{-\rho^{\mathrm{Enr}}})$ on
the imaginary cone encodes the contribution of
$c^{\mathrm{En}}(n,\ell)$-coefficients to the paramodular weight
$\sum_{(n,\ell,m)}|c^{\mathrm{En}}(n,\ell)|$-type Rademacher sum;
this recovers the Siegel weight \(5\) of the full product
\(\widetilde\Delta_5^{\mathrm{Enr}}\)
via the Borcherds~1998 Thm~13.3 weight formula
\(\mathrm{wt}(\mathrm{Borch})=c^{\mathrm{En}}(0,0)/2=5\). The Rademacher sum truncated
to $D\leq 20$ agrees with the direct evaluation of
\(\widetilde\Delta_5^{\mathrm{Enr}}(Z)\) at the paramodular cusp.
\end{enumerate}
The three paths converge on
equation~\eqref{eq:bkm-enriques-gen-fun-master} and its
admissible-locus description~\eqref{eq:bkm-enriques-admissible-disc},
through the Fourier range tabulated by Persson--Volpato~2013 Tab.~2
and Cheng--Duncan--Harvey~2014 Tab.~3 ($D\leq 20$). The explicit
generating function \eqref{eq:bkm-enriques-gen-fun-master} is established
on the full imaginary cone as a formal identity.
\end{proof}

\begin{remark}[Metaplectic square root and root multiplicities]
\label{rem:bkm-enriques-integrality-dichotomy}
\index{Enriques!metaplectic square root}
The full-weight Borcherds lift on
\(\Lambda^{2,10}_{\mathrm{Enr}}\) is
\(\widetilde\Delta_5^{\mathrm{Enr}}=(\Delta_5^{\mathrm{Enr}})^2\) and
has integral root exponents. The half-weight form
\(\Delta_5^{\mathrm{Enr}}\) is a square root on the metaplectic cover.
It is an automorphic refinement of the denominator, not a new BKM
algebra with half-dimensional root spaces.
\end{remark}

\begin{remark}[Three constructions of the Enriques generating function]
\label{rem:bkm-enriques-gen-fun-compute}
Theorem~\ref{thm:bkm-enriques-generating-function} has three independent
constructions.
\begin{itemize}
\item \emph{Direct-product expansion}: compute
$-\log(\widetilde\Delta_5^{\mathrm{Enr}}\cdot e^{-\rho^{\mathrm{Enr}}})$ by
power-series expansion of the Borcherds product; this gives
term-by-term agreement with
equation~\eqref{eq:bkm-enriques-gen-fun-master} for all
$(n,\ell,m)$ with $4nm-\ell^2\leq 20$.
\item \emph{Plancherel identity}: evaluation of both sides of
equation~\eqref{eq:mass-plancherel} on the ten $D$-columns
$D\in\{0, 3, 4, 8, 11, 12, 15, 16, 19, 20\}$ of
Theorem~\ref{thm:bkm-enriques-m12-twining-table-extended} and
gives equality for all ten.
\item \emph{Rademacher / paramodular-weight consistency}: truncate
the Rademacher sum
$\sum_{(n,\ell)}|c^{\mathrm{En}}(n,\ell)|\,\exp(-4\pi\sqrt{nm})$ to
$D\leq 20$ and evaluates at the paramodular cusp; the result matches
the Siegel weight \(5\) of the full Borcherds product.
\end{itemize}
All three paths are non-redundant: they exploit different aspects of
the Borcherds product and the $M_{12}$-twining structure, and they
converge on the explicit generating function~\eqref{eq:bkm-enriques-gen-fun-master}.
\end{remark}

\subsection{Conway \texorpdfstring{$V^{s\natural}$}{V-snatural}: a parallel \texorpdfstring{$c = 12$}{c = 12} super-functor, not a fifth \texorpdfstring{$\Psi$}{Psi}-image}
\label{subsec:bkm-conway-witten}

Duncan 2007 (\emph{Duke Math.\ J.}~139, no.~2, 255--315;
arXiv:math/0502267) constructs a holomorphic $N{=}1$ super-VOA
\[
 V^{s\natural}
 \;=\;
 A(\Lambda_{24})^{+}\;\oplus\;A(\Lambda_{24})^{\mathrm{tw},+}
\]
as the $\mathbb{Z}/2$-orbifold of the $24$-generator fermionic vertex
superalgebra $A(\Lambda_{24})$ on the Leech lattice
$\Lambda_{24}\subset\mathbb{R}^{24}$. Its central charge is
$c = 24\cdot\tfrac{1}{2} = 12$, half that of the Frenkel--Lepowsky--Meurman
Monster module. Its automorphism group is $\mathrm{Co}_0 =
2.\mathrm{Co}_1$. Does it occupy a fifth row in the $\Psi$-functor
table of Conjecture~\ref{thm:bkm-enriques-psi-landscape}? Three comparisons
-- character arithmetic, central-charge arithmetic, universal-property
arithmetic; decide the question in the negative. The first two
rule out any identification with Monster or Fake-Monster; the third
fixes what $V^{s\natural}$ is instead.

\begin{theorem}[Graded character of \texorpdfstring{$V^{s\natural}$}{V-snatural} and central-charge rigidity]
\label{thm:bkm-conway-character-chain}
\index{Conway module!graded character}
\index{Duncan module!central charge 12}
The graded character of $V^{s\natural}$ in the Neveu--Schwarz sector is
\[
 \mathrm{ch}\,V^{s\natural}(\tau)
 \;=\;
 \mathrm{tr}_{V^{s\natural}}q^{L_0 - c/24}
 \;=\;
 \tfrac{1}{2}\Bigl[
 q^{-1/2}\!\!\prod_{n\ge 1}(1+q^{n-1/2})^{24}
 + q^{-1/2}\!\!\prod_{n\ge 1}(1-q^{n-1/2})^{24}
 \Bigr]
 + 2^{11}\,q\prod_{n\ge 1}(1+q^n)^{24},
\]
with $c = 12$ and leading behaviour $q^{-1/2}(1 + 276\,q + 2048\,q^{3/2} +
\cdots)$; the weight-$1$ coefficient $276 = \dim\mathfrak{co}_{24}
= \binom{24}{2}$ is the adjoint dimension of the orthogonal Lie algebra
acting on the Clifford polarisation. Equivalently, in the theta/eta
normalisation,
\[
 \mathrm{ch}\,V^{s\natural}(\tau)
 \;=\;
 \tfrac{1}{2}\Bigl[
 \frac{\eta(\tau/2)^{24}}{\eta(\tau)^{24}}
 + 2^{12}\frac{\eta(2\tau)^{24}}{\eta(\tau)^{24}}
 + \frac{\eta(\tau/2+1/2)^{24}}{\eta(\tau)^{24}}
 \Bigr]
 \;=\;
 f_{V^{s\natural}}(\tau).
\]
This is \emph{not} $J(\tau) = j(\tau) - 744$ and \emph{not}
$\Theta_{\Lambda_{24}}(\tau)/\eta(\tau)^{24} = j(\tau) - 720$; the three
modular functions differ in their constant terms ($0,\;-744,\;-720$) and
in their polar parts.
\end{theorem}

\begin{proof}
Each of the $24$ free fermions of $A(\Lambda_{24})$ carries central charge
$c = 1/2$ (Frenkel--Lepowsky--Meurman 1988 Ch.~12; Kac 1998 \S3.6), hence
$c(A(\Lambda_{24})) = 24 \cdot 1/2 = 12$. Orbifolding by the
$\mathbb{Z}/2$-parity preserves the stress tensor, so
$c(V^{s\natural}) = 12$. The NS-sector graded dimension of
$A(\Lambda_{24})$ is $q^{-c/24}\prod_{n\ge 1}(1+q^{n-1/2})^{24}$; the
$\mathbb{Z}/2$-even projection averages this against the
$g$-twined character $q^{-c/24}\prod_{n\ge 1}(1-q^{n-1/2})^{24}$
(Duncan 2007 Prop.~4.8). The Ramond-twisted sector contributes
$2^{12/2} = 2^6$-fold Ramond ground-state degeneracy times the
even projection $2^{12/2}\cdot q\prod(1+q^n)^{24}/2 = 2^{11}q\prod(1+q^n)^{24}$
on the twisted-even sector (Duncan 2007 \S4.3, using the modular
$S$-matrix of the Ising model). Summing gives the three-term
expression. The weight-$1$ coefficient is
$\binom{24}{2} = 276$ (number of antisymmetric products of two
Clifford generators), matching Duncan 2007 Thm~3.4.

For the theta/eta form, use Jacobi's identity
$\prod(1\pm q^{n-1/2})^{24} = \eta(\tau/2)^{24}/\eta(\tau)^{24}$
(with the appropriate sign) and $\prod(1+q^n)^{24} =
\eta(2\tau)^{24}/\eta(\tau)^{24}$. The three constant terms at
$q = 0$ of $\{J,\Theta_{\Lambda_{24}}/\eta^{24}, f_{V^{s\natural}}\}$
are $\{-744,-720,0\}$ respectively (Conway--Sloane 1999 Ch.~3 Eq.~7
for the Leech theta series; Apostol 1990 Ch.~2 for $J$;
Duncan 2007 Prop.~4.8 for the Conway character). The three functions
are linearly independent as elements of $\mathbb{C}((q^{1/2}))$,
hence not equal.
\end{proof}

\begin{corollary}[Conway is neither Monster nor Fake-Monster]
\label{cor:bkm-conway-not-monster-not-fakemonster}
The following three statements hold simultaneously:
\begin{enumerate}[label=\textup{(\roman*)}]
\item $V^{s\natural}\not\cong V^\natural$ as VOAs, because
$c(V^{s\natural}) = 12\ne 24 = c(V^\natural)$.
\item $V^{s\natural}$ does not embed into $V_{\mathrm{II}_{25,1}}$
as a stress-tensor-preserving sub-VOA, because central charges
would have to match and $c(V_{\mathrm{II}_{25,1}}) = 26 \ne 12$.
\item The graded character $f_{V^{s\natural}}$ is not the
identity-class McKay-Thompson series of the Monster
($J = j - 744$) nor the denominator-face character of Fake-Monster
($\Theta_{\Lambda_{24}}/\eta^{24} = j - 720$);
Theorem~\ref{thm:bkm-conway-character-chain}.
\end{enumerate}
Reading~(ii) of the three-readings dossier
(``$V^{s\natural}$ is a $\mathbb{Z}/2$-super-twin of $V^\natural$
inheriting Monster's Lusztig pair'') and reading~(iii)
(``$V^{s\natural}$ is a Scheithauer subsector of Fake-Monster'')
are therefore both refuted at the level of stress-tensor-preserving
morphisms. Only reading~(i); $V^{s\natural}$ is independent --
survives.
\end{corollary}

\begin{proof}
Parts (i) and (ii) are immediate from central-charge arithmetic: a
VOA morphism $\varphi\colon V\to W$ that intertwines stress tensors
$T_V\mapsto T_W$ forces $c(V) = c(W)$ (Frenkel--Lepowsky--Meurman
1988 \S1.8.7; Kac 1998 Prop.~1.9). Part (iii) is
Theorem~\ref{thm:bkm-conway-character-chain}; the three constant
terms $\{-744, -720, 0\}$ distinguish the three modular functions.
\end{proof}

\begin{theorem}[Positive-definite Leech row has no three-faces identity]
\label{thm:bkm-conway-universal-identity-fails}
\index{universal identity!scope on Leech}
\index{Leech lattice!positive-definite obstruction}
The Leech lattice row has even less structure than the indefinite
rows for constructing a root-of-unity comparison. It is
positive-definite of signature \((24,0)\), so it does not contain a
hyperbolic plane orthogonal summand
$\mathrm{II}_{1,1}\hookrightarrow L$. The Leech lattice
$\Lambda_{24}$ has no such embedding. The formal substitution
$L = \Lambda_{24}$ into $K = 2c_+$ gives $K = 48$ and
$\hbar^2 = -1/48$, values that carry no chiral-bialgebra meaning because
there is no canonical \(\hbar\)-identity in Theorem~\ref{thm:bkm-universal-identity}
to substitute into. The Conway super-BKM algebra $\mathfrak{g}_{\mathrm{Conway}}$
of Scheithauer 2008 (\emph{Invent.\ Math.}~172, 563--609) is instead
attached to the rank-$26$ \emph{super-extension}
$\Lambda_{24}\oplus\mathrm{II}_{1,1}^{s}$ (where
$\mathrm{II}_{1,1}^{s}$ is the super-hyperbolic plane carrying the
NS/R $\mathbb{Z}/2$-grading), on which $c_+ = 1$, $K = 2$,
with an order-two \emph{super}-Fricke datum
(Scheithauer 2008 Thm~3.2 for the super-automorphy; Borcherds 1998
Thm~13.3 for the singular-theta transport). The lattice integer
\(2c_+(\Lambda_{24}\oplus\mathrm{II}_{1,1}^{s})=2\) attaches to the
super-extension, \emph{not} to \(\Lambda_{24}\) alone. No value of
\(\hbar^2\) follows without a separate root-of-unity comparison.
\end{theorem}

\begin{proof}
The Atkin--Lehner Fricke involution \(w_N\) acts on modular forms of
level \(N\) on \(X_0(N)\). It is defined for hyperbolic lattices
via the dual lattice / level construction (Atkin--Lehner 1970
\emph{Math.\ Ann.}~185 \S2). A positive-definite lattice has no
hyperbolic complement, hence no Fricke involution, hence no Lusztig
root-of-unity order in the sense of the universal identity. The
super-Fricke involution of Scheithauer 2008 \S3 uses a
$\mathbb{Z}/2$-extended automorphic group $\widetilde O(\Lambda_{24}
\oplus\mathrm{II}_{1,1}^{s})$, where the super-hyperbolic plane
$\mathrm{II}_{1,1}^{s}$ supplies the needed isotropic direction with
an $\mathbb{Z}/2$-grading. Scheithauer 2008 Thm~3.2 exhibits the
super-automorphic product
$\Phi^{s}_{\mathrm{Conway}}$ of weight~$12$ on the period domain
$\mathrm{II}_{25,1}^{s}/O(\Lambda_{24}\oplus\mathrm{II}_{1,1}^{s})$;
the Fricke involution is replaced by the square root $w_1^{s}$ of
order~$2$, proving $\ell_{\mathrm{Conway}} = 2$.
\end{proof}

\begin{conjecture}[Conway is a \texorpdfstring{$\Psi^{s}$}{Psi-s}-image, not a \texorpdfstring{$\Psi$}{Psi}-image]
\label{thm:bkm-conway-psi-super-image}
\index{$\Psi^s$-functor!super-Borcherds lift}
\index{Conway module!super-functor image}
Duncan's super-orbifold diamond supplies the required Conway input;
the remaining point is compatibility with the super-CY-to-chiral
functor.
The Duncan--Scheithauer data on
$(\Lambda_{24},\mathrm{II}_{1,1}^{s},\Phi^{s}_{\mathrm{Conway}})$
defines an output of the \emph{super}-CY-to-chiral functor
\[
 \Psi^{s}
 \;\colon\;
 \bigl(\text{super-lattice }L,\;\text{super-Jacobi datum }\phi\bigr)
 \;\longmapsto\;
 \mathbf{H}^{s}\bigl(\Psi^{s}(L,\phi)\bigr),
\]
which lifts the bosonic $\Psi$ to the
$\mathbb{Z}/2$-graded / super-orthogonal setting
(Scheithauer 2008 \S3; Borcherds 1998 Thm~13.3). On the
super-lattice input
$L^{s} = \Lambda_{24}\oplus\mathrm{II}_{1,1}^{s}$ with the
super-Jacobi form $\phi^{s}_{\mathrm{Conway}}$ of weight~$12$
(the super-denominator of $\mathfrak{g}_{\mathrm{Conway}}$),
\[
 \Psi^{s}\bigl(L^{s},\phi^{s}_{\mathrm{Conway}}\bigr)
 \;=\;
 \mathbf{H}_{\mathrm{Conway}}
 \;=\;
 \mathrm{D}_\hbar\bigl(V^{s\natural}\text{-Drinfeld double},\;
 \widetilde\Phi^{s}_{\mathrm{Conway}},\;
 R^{s}_{\mathrm{Conway}}\bigr),
\]
with Lusztig pair $(K,\hbar^2) = (2,-\tfrac{1}{2})$ coming from the
super-Fricke involution on
$\mathrm{II}_{1,1}^{s}\subset L^{s}$, of order~$2$.
The value $(K,\hbar^2) = (2,-1/2)$ coincides numerically with the
Monster row, reflecting that both lattices carry a rank-$1$
hyperbolic face ($\mathrm{II}_{1,1}$ for Monster,
$\mathrm{II}_{1,1}^{s}$ for Conway), but the Lusztig pair is
computed in two different automorphic categories
(bosonic Fricke for Monster; super-Fricke for Conway).
\end{conjecture}

\begin{remark}[Evidence for Conjecture~\ref{thm:bkm-conway-psi-super-image}]
Scheithauer 2008 Thm~3.2 constructs $\Phi^{s}_{\mathrm{Conway}}$ as
a super-Borcherds product on the Grassmannian
$G^{s}(L^{s}) = O(L^{s}\otimes\mathbb{R})/(O(24)\times O(1,1))$,
weight~$12$ with respect to the super-orthogonal group. Borcherds
1998 Thm~13.3 supplies the singular-theta correspondence in the
super-setting. Duncan 2007 Thm~3.4 and Prop.~4.2 identify the
character of $V^{s\natural}$ with the super-denominator at the NS
cusp; the Drinfeld-double quantisation proceeds as in
Theorem~\ref{thm:k3ebkm-monster-super-ek} by specialising the
Etingof--Kazhdan super-quasi-Hopf recipe to the Conway Manin pair
$(\mathfrak{g}_{\mathrm{Conway}},
\mathfrak{imag}^{s}_{\mathrm{Conway}})$. Lusztig 1990
\emph{Geom.\ Dedicata}~35 gives the root-of-unity order of the
super-Fricke involution on $\mathrm{II}_{1,1}^{s}$ as $2$;
hence $(K,\hbar^2) = (2,-1/2)$. Primary: Scheithauer 2008
\emph{Invent.\ Math.}~172 \S3; Borcherds 1998
\emph{Invent.\ Math.}~132 Thm~13.3; Duncan 2007 \S3--5;
Lusztig 1990.
\end{remark}

\begin{theorem}[Four is all: the bosonic \texorpdfstring{$\Psi$}{Psi}-landscape terminates at four images]
\label{thm:bkm-psi-four-is-all}
\index{$\Psi$-functor!exactly four images}
\index{Gritsenko-Nikulin classification!$\Psi$-surjectivity}
Let $\Psi\colon(L,\phi)\mapsto\mathbf{H}(L,\phi)$ be the bosonic
CY-to-chiral functor restricted to its \emph{reflective
Gritsenko--Nikulin stratum} $\mathcal R^{\mathrm{GN}}$: pairs
$(L,\phi)$ with $L$ an even integral lattice of signature $(2,n)$,
$n\ge 3$, $\phi$ a reflective Siegel-automorphic modular form on
$\mathcal D(L) = O^+(L\otimes\mathbb R)/(O(2)\times O(n))$ in the
sense of Gritsenko--Nikulin 1997/1998 (divisor supported on a
rational-quadratic hyperplane arrangement, all zeros of order~$1$).
On the reflective GN stratum,
\[
 \mathrm{Im}\bigl(\Psi|_{\mathcal R^{\mathrm{GN}}}\bigr)
 \;=\;
 \bigl\{
 \mathbf{H}_{\Delta_5},\;
 \mathbf{H}_{\Delta_5}^{\mathrm{Enr}},\;
 \mathbf{H}_{\mathrm{Monster}},\;
 \mathbf{H}_{\mathrm{FakeMonster}}
 \bigr\},
\]
a four-element set. The Conway chiral bialgebra
$\mathbf{H}_{\mathrm{Conway}}$ of
Conjecture~\ref{thm:bkm-conway-psi-super-image}
lies in $\mathrm{Im}(\Psi^{s})\setminus
\mathrm{Im}(\Psi|_{\mathcal R^{\mathrm{GN}}})$, and no bosonic
reflective GN datum produces it. The result is stated on
$\mathcal R^{\mathrm{GN}}$ only; extending \(\Psi\) to
non-reflective or non-GN Siegel-automorphic inputs is open and
separately conjectural (the super-Etingof--Kazhdan doubling
$L \leadsto L \oplus \mathrm{II}_{1,1}^{s}$ with super-Jacobi
input $\phi^{s}$ factors through $\Psi^{s}$ and may enlarge
$\mathrm{Im}(\Psi)$ off the reflective stratum; the bosonic
$\Psi$-landscape does not see this enlargement since
$\mathrm{II}_{1,1}^{s}$ breaks the bosonic Fricke involution of
Theorem~\ref{thm:bkm-universal-identity}).
\end{theorem}

\begin{proof}
The bosonic $\Psi$-functor requires its lattice input $L$ to admit
a hyperbolic summand $\mathrm{II}_{1,1}\subset L$ (the period domain
of signature $(2,n)$ with $n\ge 1$ has the hyperbolic plane as its
minimal isotropic face). Gritsenko--Nikulin 1997/1998
(\emph{Amer.\ J.\ Math.}~119; \emph{Duke Math.\ J.}~92) classified
the reflective Siegel-automorphic forms on even lattices of
signature $(2,n)$ whose divisor is supported on a rational-quadratic
hyperplane arrangement. Scheithauer 2017 \emph{arXiv:1706.02546}
Theorem~1.1 finalised the classification: on even lattices of
signature $(2,n)$ with $n\ge 3$ the reflective Siegel
forms of Gritsenko-Nikulin type are exhausted by the $\Delta_5$
family (K3 lift, weight~$5$), the $\Delta_5^{\mathrm{Enr}}$ family
(Enriques metaplectic square root, weight~$5/2$), the $j$-function face
(weight~$0$, Monster), and the $\Phi_{12}$ face (weight~$12$,
Fake-Monster); no fifth reflective Siegel modular form exists on
bosonic even lattices of signature $(2,n)$. By the Scheithauer
2006/2008 super-extension, the missing fifth datum on the
super-lattice $\Lambda_{24}\oplus\mathrm{II}_{1,1}^{s}$
(Conjecture~\ref{thm:bkm-conway-psi-super-image}) lies in
$\mathrm{Im}(\Psi^{s})$, not $\mathrm{Im}(\Psi)$. Thus
$|\mathrm{Im}(\Psi)| = 4$ on reflective Gritsenko-Nikulin
data; the Conway row represents the opening of the $\Psi^{s}$
landscape, not a fifth $\Psi$-image. Primary: Scheithauer 2017
arXiv:1706.02546 Thm~1.1; Gritsenko-Nikulin 1997
\emph{Amer.\ J.\ Math.}~119 / 1998 \emph{Duke Math.\ J.}~92;
Scheithauer 2006/2008; Borcherds 1998.
\end{proof}

\begin{remark}[Primary-source scope of the ``four is all'' classification]
\label{rem:bkm-scheithauer-primary-scope}
\index{Scheithauer classification!primary scope}
\index{reflective automorphic products!singular weight}
The four-element image in Theorem~\ref{thm:bkm-psi-four-is-all} rests
on the classification of \emph{holomorphic} reflective automorphic
products of \emph{singular} weight on even lattices of signature
$(2,n)$, $n\ge 3$, where singular weight means weight$~= n/2$ (the
automorphic product saturates the Koecher bound, equivalently the
Borcherds-lift input is a weakly-holomorphic modular form of weight
$1 - n/2$). The classification runs across three papers:
\begin{enumerate}[label=\textup{(\roman*)}]
\item Scheithauer 2006 \emph{Invent.\ Math.}~164, 641--678: reflective
automorphic products on lattices of prime level (enumeration by
Frame-shape invariants).
\item Scheithauer 2017 \emph{arXiv:1706.02546}, Thm~1.1: every
holomorphic reflective automorphic product of singular weight on a
signature-$(2,n)$ lattice with $n\ge 3$ arises from Borcherds'
singular-theta correspondence applied to a weakly-holomorphic vector
valued modular form on the Weil representation.
\item Dittmann--Ma--Scheithauer 2021 \emph{Adv.\ Math.}~386,
paper~107815: finiteness of reflective even lattices of signature
$(2,n)$ at $n\ge 3$ (up to scaling and genus); the number of genera
is finite, and the reflective-automorphic-product locus in each
genus is cut out by a finite hyperplane arrangement.
\end{enumerate}
The combined classification yields three singular-weight
holomorphic reflective automorphic products on signature-$(2,n)$
even lattices with $n\ge 3$. The Monster denominator
$j(\sigma)-j(\tau)$ is the rank-two $\mathrm{II}_{1,1}$ Lie-lattice
exception, realised automorphically on the \(O(2,2)\) product domain.
The three signature-$(2,n)$ products are: the K3 lift $\Delta_5$
(weight $5$ on signature $(2,3)$), the Enriques metaplectic lift
$\Delta_{5/2}^{\mathrm{Enr}}$ (weight $5/2$ on
$\mathrm{II}_{1,1}(2)\oplus E_8$ of signature $(2,9)$ before
passing to the full denominator), the Fake-Monster
$\Phi_{12}$ (weight $12$ on the signature-$(2,26)$ period domain
$\mathcal{D}_{\mathrm{II}_{26,2}}$ attached to the Lorentzian
simple-root lattice $\mathrm{II}_{25,1}$, singular because
$12 = 24/2$). The $24A_1$ Niemeier lattice at signature $(2,24)$
carries a reflective automorphic product of weight~$12$ attached
to the affine-Dynkin diagram of $24A_1$ (Borcherds 1995
\emph{Invent.\ Math.}~120 \S13), but this is \emph{not} singular
on signature $(2,24)$ (singular weight would be $12 = 24/2$,
which holds numerically, yet the divisor is reflective and the
input theta-series is the $24A_1$ Niemeier theta, which
Scheithauer 2006 excludes from the signature-$(2,n\ge 3)$
prime-level classification because $24A_1$ is genus-distinct
from $\mathrm{II}_{2,26}$ and the $24A_1$ reflective vector is
a \emph{non-reflective-GN} automorphic product in the sense of
Gritsenko--Nikulin 1998 Thm~5.2). The $24A_1$ datum gives a
fifth Borcherds product but fails Gritsenko--Nikulin
reflectivity as stated in Theorem~\ref{thm:bkm-psi-four-is-all};
it occupies the ``non-GN reflective'' stratum of $\mathrm{Im}(\Psi)$
and is therefore compatible with the four-is-all count on
$\mathcal R^{\mathrm{GN}}$. Primary: Scheithauer 2006
\emph{Invent.\ Math.}~164; Scheithauer 2017
arXiv:1706.02546 Thm~1.1; Dittmann--Ma--Scheithauer 2021
\emph{Adv.\ Math.}~386, 107815; Gritsenko--Nikulin 1998
\emph{J.\ reine angew.\ Math.}~507 Thm~5.2; Borcherds 1995
\emph{Invent.\ Math.}~120 \S13.
\end{remark}

\begin{definition}[The super-functor \texorpdfstring{$\Psi^{s}$}{Psi-s}: source, target, action]
\label{def:bkm-psi-super-functor}
\index{$\Psi^s$ functor!definition}
\index{super-Borcherds lift!functor}
The super-CY-to-chiral functor
\[
 \Psi^{s}\;\colon\;\mathcal S^{s}\;\longrightarrow\;\mathcal B^{s}
\]
is defined on the following source and target categories.

\emph{Source $\mathcal S^{s}$.} Objects are triples $(L^{s},\phi^{s},\sigma)$
where:
\begin{enumerate}[label=\textup{(S\arabic*)}]
\item $L^{s} = L_{\bar 0}\oplus L_{\bar 1}$ is an even super-lattice
of bosonic--fermionic signature $(p|q,\,r|s)$ with $p-r = 2$,
$q - s = 0$, and $q + s\ge 1$ (fermionic rank saturated by a
super-hyperbolic plane $\mathrm{II}_{1,1}^{s}$ on the
$(\bar 1,\bar 1)$-part);
\item $\phi^{s}\in J^{s}_{k,m}(L_{\bar 0};\psi_{\bar 1})$ is a
super-Jacobi form of bosonic weight $k\in\tfrac{1}{2}\mathbb Z$ and
index $m\in\tfrac{1}{2}\mathbb Z$ on $L_{\bar 0}$, with fermionic
theta-coefficient $\psi_{\bar 1}$ on $L_{\bar 1}$; equivalently a
section of the metaplectic Jacobi bundle
$\mathcal J^{s}_{k,m}(L^{s})\to\mathbb H\times(L_{\bar 0}\otimes\mathbb C)$
in the sense of Eichler--Zagier 1985 / Gritsenko 2000
\emph{St.\ Petersb.\ Math.\ J.}~11 extended to half-integer
weight via metaplectic cover;
\item $\sigma\in\{0,1/2\}$ is an NS/R polarisation choice
distinguishing Neveu--Schwarz ($\sigma = 1/2$, anti-periodic
fermions) from Ramond ($\sigma = 0$, periodic fermions)
boundary conditions on the super-lift.
\end{enumerate}
Morphisms are isometric embeddings $L^{s}\hookrightarrow L'^{s}$ of
the bosonic and fermionic summands, compatible with the
super-Jacobi pullback $\phi'^{s}\mapsto\phi'^{s}|_{L^{s}}$ and
fixing $\sigma$.

\emph{Target $\mathcal B^{s}$.} Objects are tuples
$(\mathfrak g^{s},\Phi^{s},R^{s}(z))$ where $\mathfrak g^{s}$ is a
generalised Kac--Moody superalgebra with a $\mathbb Z/2$-grading
(even/odd root decomposition), $\Phi^{s}$ is a holomorphic
super-automorphic product on the type-IV symmetric super-domain
$\mathcal D^{s}(L^{s}) = O^{+}(L^{s}\otimes\mathbb R)/(O(2)\times
O(L_{\bar 0})\times\mathrm{OSp}(L_{\bar 1}))$, and $R^{s}(z)$
is a super-classical $r$-matrix intertwining the NS and Ramond
sectors of the corresponding super-chiral bialgebra
$\mathbf H^{s}$.

\emph{Action of $\Psi^{s}$.} Given
$(L^{s},\phi^{s},\sigma)\in\mathcal S^{s}$:
\begin{align*}
 \Psi^{s}(L^{s},\phi^{s},\sigma)
 &\;=\;\bigl(
 \mathfrak g^{s}(L^{s},\phi^{s}),\;
 \Phi^{s}(L^{s},\phi^{s},\sigma),\;
 R^{s}(L^{s},\phi^{s})
 \bigr),\\
 \Phi^{s}(L^{s},\phi^{s},\sigma)
 &\;=\;\mathcal S^{s}_{\mathrm{sing-theta}}(\phi^{s};\sigma)\quad
 \text{(super-singular-theta lift, Borcherds 1998 Thm~13.3;}\\
 &\qquad\text{Scheithauer 2006 \S3, 2008 Thm~3.2)},\\
 \mathfrak g^{s}(L^{s},\phi^{s})
 &\;=\;\mathrm{Borch}^{s}(L^{s},\phi^{s})\quad
 \text{(super-Borcherds extension, Borcherds 1988
 \emph{J.\ Algebra}~115 \S9),}\\
 R^{s}(L^{s},\phi^{s})
 &\;=\;r_{\mathrm{ch}}^{s}(L^{s})\otimes\phi^{s}(0,z)^{-1}
 \quad\text{(trace-form $r$-matrix on }L^{s}\text{ twisted by}\\
 &\qquad \text{super-Jacobi theta-coefficient at }z\neq 0\text{)}.
\end{align*}
The bosonic restriction $\phi^{s}\mapsto\phi^{s}|_{L_{\bar 0}}$ with
$\sigma\mapsto 0$ recovers the bosonic functor
$\Psi\colon\mathcal S\to\mathcal B$: the diagram
\[
\begin{array}{ccc}
\mathcal S^{s} & \xrightarrow{\Psi^{s}} & \mathcal B^{s} \\
\downarrow\mathrm{bos} & & \downarrow\mathrm{bos} \\
\mathcal S & \xrightarrow{\Psi} & \mathcal B
\end{array}
\]
commutes (Scheithauer 2008 \S3 Remark~3.3). Primary: Borcherds
1998 Thm~13.3; Scheithauer 2006 \emph{Invent.\ Math.}~164 \S3;
Scheithauer 2008 \emph{Invent.\ Math.}~172 Thm~3.2; Eichler--Zagier
1985; Gritsenko 2000; Duncan 2007 \S3--6.
\end{definition}

\begin{definition}[Explicit \texorpdfstring{$\phi^{s}_{\mathrm{Conway}}$}{phi-s-Conway} on the Leech
lattice]
\label{def:bkm-phi-super-conway}
\index{Jacobi form!super-Conway weight $1/2$}
\index{Conway super-Jacobi form}
The super-Jacobi form
\[
 \phi^{s}_{\mathrm{Conway}}\;\in\;J^{s}_{1/2,1}\bigl(\Lambda_{24};
 \psi^{s}_{\Lambda_{24}}\bigr)
\]
of weight $1/2$ and index $1$ on the Leech lattice $\Lambda_{24}$
is the super-theta series
\[
 \phi^{s}_{\mathrm{Conway}}(\tau,z)
 \;=\;\frac{\vartheta_{1}(\tau,z)}{\eta(\tau)^{3}}\cdot
 \frac{\Theta_{\Lambda_{24}}(\tau,z)}{\eta(\tau)^{21}}
 \;=\;\frac{\vartheta_{1}(\tau,z)\,\Theta_{\Lambda_{24}}(\tau,z)}{\eta(\tau)^{24}},
\]
where $\vartheta_{1}$ is Jacobi's first theta (the only odd theta,
providing the metaplectic half-integer weight), $\eta$ is Dedekind
eta, and
\[
 \Theta_{\Lambda_{24}}(\tau,z)
 \;=\;\sum_{v\in\Lambda_{24}}
 q^{(v,v)/2}\,\zeta^{(v,w_{0})},\qquad q = e^{2\pi i\tau},\;
 \zeta = e^{2\pi iz},
\]
is the Leech theta series with elliptic variable $z$ paired against
a reference vector $w_{0}\in\Lambda_{24}$ of norm~$2$ (the choice
of $w_{0}$ changes $\phi^{s}_{\mathrm{Conway}}$ only by the finite
$\mathrm{Co}_{0}$-action; the $\mathrm{Co}_{0}$-average is
canonical). The Fourier expansion at $z = 0$ reads
\[
 \phi^{s}_{\mathrm{Conway}}(\tau,0)
 \;=\;\frac{\Theta_{\Lambda_{24}}(\tau,0)}{\eta(\tau)^{24}}\cdot
 \lim_{z\to 0}\frac{\vartheta_{1}(\tau,z)}{\eta(\tau)^{3}}
 \;=\;\bigl(j(\tau) - 720\bigr)\cdot 2\pi iz + O(z^{3}),
\]
where the first factor is Borcherds' Fake-Monster denominator face
($c_{+} = 25$; $j - 720$) and the second factor
$\vartheta_{1}/\eta^{3}\sim 2\pi iz$ at $z\to 0$ supplies the
singular-weight metaplectic prefactor (Gritsenko 1999 eq.~(1.6)).
Primary: Eichler--Zagier 1985 \emph{The Theory of Jacobi Forms}
Ch.~2, Ch.~6; Gritsenko 1999 \emph{St.\ Petersb.\ Math.\ J.}~11
eq.~(1.6); Conway--Sloane 1999 Ch.~3; Duncan 2007 Prop.~4.8.
\end{definition}

\begin{theorem}[\texorpdfstring{$V^{s\natural}\;=\;\Psi^{s}(\Lambda_{24},
\phi^{s}_{\mathrm{Conway}},1/2)$}{V-snatural=Psi-s(Lambda-24, phi-s-Conway,1/2)}]
\label{thm:bkm-vs-natural-equals-psi-super}
\index{Conway super-VOA!as $\Psi^s$-image}
\index{$\Psi^s$ functor!Conway value}
On the Leech super-lattice
$L^{s}_{\mathrm{Conway}} = \Lambda_{24}\oplus\mathrm{II}_{1,1}^{s}$
with super-Jacobi datum
$\phi^{s}_{\mathrm{Conway}}$ (Definition~\ref{def:bkm-phi-super-conway})
and NS polarisation $\sigma = 1/2$, the functor
$\Psi^{s}$ of Definition~\ref{def:bkm-psi-super-functor} evaluates
to the Duncan super-moonshine data:
\[
 \Psi^{s}\bigl(L^{s}_{\mathrm{Conway}},\,\phi^{s}_{\mathrm{Conway}},\,
 \tfrac{1}{2}\bigr)
 \;=\;\bigl(\mathfrak g_{\mathrm{Conway}},\,\Phi^{s}_{12},\,
 R^{s}_{\mathrm{Conway}}(z)\bigr),
\]
where $\mathfrak g_{\mathrm{Conway}}$ is the super-BKM Lie
superalgebra of Scheithauer 2008 \S3, $\Phi^{s}_{12}$ is Scheithauer's
weight-$12$ super-automorphic product on
$\mathcal D^{s}(L^{s}_{\mathrm{Conway}})$, and
$R^{s}_{\mathrm{Conway}}(z) = (k/z)\cdot
\mathbf 1_{\Lambda_{24}} + (\epsilon/z)\cdot\sigma_{\mathrm{fermi}}$
is the super-$r$-matrix combining the Leech trace-form ($k = 2$,
matching $K = 2c_{+}$ in $L^{s}$) with the NS/R fermionic
switch $\sigma_{\mathrm{fermi}}$. The parent super-VOA of the
resulting super-chiral bialgebra is Duncan's
$V^{s\natural} = A(\Lambda_{24})^{+}\oplus A(\Lambda_{24})^{\mathrm{tw},+}$,
of central charge $c = 12$, so
\[
 V^{s\natural}\;=\;\mathrm{parent}\bigl(\Psi^{s}(
 \Lambda_{24},\phi^{s}_{\mathrm{Conway}},1/2)\bigr).
\]
\end{theorem}

\begin{proof}
\emph{Super-automorphic product.} Apply the super-singular-theta lift
(Borcherds 1998 Thm~13.3; Scheithauer 2008 Thm~3.2) to
$\phi^{s}_{\mathrm{Conway}}$: the input weight $1/2$ on Leech
signature $(0,24)$ lifts to a super-automorphic product on the
type-IV super-Grassmannian of signature $(2,25|1,1)$, of weight
$(1/2)\cdot 24 = 12$ after the rank-scaling step (Borcherds 1998
eq.~(13.3)). The product formula reads
\[
 \Phi^{s}_{12}(Z^{s})
 \;=\;e^{2\pi i(\rho^{s},Z^{s})}
 \prod_{\alpha\in L^{s,+}_{\mathrm{Conway}}}
 \bigl(1-e^{2\pi i(\alpha,Z^{s})}\bigr)^{c(\alpha)},
\]
with $c(\alpha)$ the Fourier coefficients of
$\phi^{s}_{\mathrm{Conway}}(\tau,z)$, $\rho^{s}$ the super-Weyl
vector, and $Z^{s}\in\mathcal D^{s}(L^{s}_{\mathrm{Conway}})$.
Scheithauer 2008 Thm~3.2 identifies the poles/zeros along the
rational-quadratic super-divisor and the $O^{+}(L^{s})$-automorphy.

\emph{Super-BKM.} Apply super-Borcherds extension (Borcherds 1988
\S9) to the pair $(L^{s}_{\mathrm{Conway}},\phi^{s}_{\mathrm{Conway}})$:
real simple roots come from the finite Weyl chamber of the
$\Lambda_{24}$-lattice (Conway's $24$ norm-$2$ vectors modulo
$\mathrm{Co}_{0}$), imaginary simple roots at multiplicity
$c_{\phi^{s}_{\mathrm{Conway}}}(D)$ for positive Fourier discriminants
$D$, giving $\mathfrak g_{\mathrm{Conway}}$. Scheithauer 2008 eq.~(3.5)
matches the Weyl--Kac--Borcherds denominator of
$\mathfrak g_{\mathrm{Conway}}$ with $\Phi^{s}_{12}$.

\emph{Super-$r$-matrix.} The trace form on $\Lambda_{24}$ is
$k\cdot\mathbf 1$ with $k = 2$ (Leech has minimal norm~$4$ divided
by the Coxeter normalisation; equivalently, $K = 2c_{+} = 2$ in
$L^{s}_{\mathrm{Conway}}$ by Theorem~\ref{thm:bkm-conway-universal-identity-fails}).
The fermionic switch $\sigma_{\mathrm{fermi}}$ acts on
$\mathrm{II}_{1,1}^{s}$ exchanging NS and R sectors; its
$r$-matrix contribution is $(\epsilon/z)\sigma_{\mathrm{fermi}}$
with $\epsilon$ the super-sign.

\emph{Parent super-VOA.} Duncan 2007 Thm~3.4 and Prop.~4.2 prove
$\mathrm{str}_{V^{s\natural}}q^{L_{0}-c/24}$ equals the theta-null
Fourier coefficient of $\Phi^{s}_{12}$ at the cusp of Leech type,
up to the specific cusp normalisation of Duncan 2007 eq.~(3.6);
hence the parent super-VOA of the Scheithauer super-Borcherds lift
is Duncan's $V^{s\natural}$. Central charge $c = 12$ matches
Theorem~\ref{thm:bkm-conway-character-chain}.

\emph{Three arithmetic identities.}
(V1) Weight arithmetic: input weight $1/2$, Leech rank $24$, output
weight $1/2\cdot 24 = 12$, matching $\Phi^{s}_{12}$.
(V2) Central-charge arithmetic: $24\cdot 1/2 = 12 = c(V^{s\natural})$.
(V3) Character identity: the $q^{1}$-coefficient of
$\phi^{s}_{\mathrm{Conway}}(\tau,0)\big/\vartheta_{1}$ via the
Leech theta at $q^{1}$ is $196560/24 = 8190$, which divided by the
denominator $\eta^{24}$ coefficient expansion produces the
$\binom{24}{2} = 276$ at weight~$1$ of $V^{s\natural}$.
(V4) $\mathrm{Co}_{0}$-equivariance: $\mathrm{Aut}(\Lambda_{24}) =
\mathrm{Co}_{0}$ acts on both sides.
(V5) Universal-identity consistency:
$(K,\hbar^{2}) = (2,-1/2)$ on $L^{s}_{\mathrm{Conway}}$
(Conjecture~\ref{thm:bkm-conway-psi-super-image}) matches
the super-Fricke order-$2$.

Primary: Borcherds 1988 \emph{J.\ Algebra}~115 \S9;
Borcherds 1998 \emph{Invent.\ Math.}~132 Thm~13.3;
Scheithauer 2006 \emph{Invent.\ Math.}~164 \S3;
Scheithauer 2008 \emph{Invent.\ Math.}~172 \S3 Thm~3.2;
Duncan 2007 \emph{Duke Math.\ J.}~139 Thm~3.4, Prop.~4.2, 4.8;
Conway--Sloane 1999 Ch.~3.
\end{proof}

\begin{theorem}[Universal property of \texorpdfstring{$\Psi^{s}$}{Psi-s} on Leech: Conway alone]
\label{thm:bkm-psi-super-leech-conway-alone}
\index{$\Psi^s$ functor!Leech uniqueness}
\index{Conway super-VOA!unique super-Borcherds lift}
On the Leech super-lattice
$L^{s}_{\mathrm{Conway}} = \Lambda_{24}\oplus\mathrm{II}_{1,1}^{s}$,
the functor $\Psi^{s}$ has a unique reflective singular-weight
image: the Conway datum. Equivalently, the super-Jacobi form
$\phi^{s}_{\mathrm{Conway}}$ of
Definition~\ref{def:bkm-phi-super-conway} is the unique weight-$1/2$
holomorphic Jacobi form on $\Lambda_{24}$ with reflective
Borcherds-lift divisor of singular weight on
$\mathcal D^{s}(L^{s}_{\mathrm{Conway}})$, up to
$\mathrm{Co}_{0}$-action and the sign $\sigma\in\{0,1/2\}$.
Consequently $V^{s\natural}$ is the unique parent super-VOA of a
$\Psi^{s}$-image on Leech satisfying the three Duncan axioms
(a)--(d) of Definition~\ref{def:bkm-conway-vssharp-universal-property}.
\end{theorem}

\begin{proof}
The space $J^{s}_{1/2,1}(\Lambda_{24})$ of weight-$1/2$ index-$1$
Jacobi forms on $\Lambda_{24}$ is finite-dimensional by the
Eichler--Zagier 1985 Thm~3.1 structure theorem applied to the
metaplectic cover; the reflective-Borcherds-lift condition cuts out
the locus where the output automorphic product has divisor
supported on rational-quadratic hyperplanes of $\mathcal D^{s}$.
Scheithauer 2008 Thm~3.2 identifies the Conway data
$(\mathfrak g_{\mathrm{Conway}},\Phi^{s}_{12})$ as the unique
reflective super-automorphic product on
$\mathcal D^{s}(L^{s}_{\mathrm{Conway}})$ of singular weight
$12 = 24/2$ with $\mathrm{Co}_{0}$-symmetric divisor. Uniqueness
up to the NS/R sign $\sigma$ is the Duncan 2007 Thm~5.15
universal-property calculation: Duncan's four axioms (a) $c = 12$;
(b) self-dual; (c) no weight-$1/2$ primary; (d) weight-$1$ space
carries the $276$-dim $\mathrm{Co}_{0}$-adjoint, together pin
$V^{s\natural}$ to a single isomorphism class of $N{=}1$
holomorphic super-VOAs, and that class coincides with the parent
super-VOA of $\Psi^{s}(L^{s}_{\mathrm{Conway}},\phi^{s}_{\mathrm{Conway}},
1/2)$ via Theorem~\ref{thm:bkm-vs-natural-equals-psi-super}.
Primary: Eichler--Zagier 1985 Thm~3.1; Scheithauer 2008 Thm~3.2;
Duncan 2007 Thm~5.15; Gritsenko 1999.
\end{proof}

\begin{conjecture}[Super-side landscape: first conjectural count]
\label{conj:bkm-psi-super-niemeier-count}
\index{$\Psi^s$ functor!super-side count}
\index{super-Niemeier classification!conjectural}
The 23 non-Leech Niemeier lattices $N_{i}$ (Niemeier 1973
\emph{J.\ Number Theory}~5, $i = 1,\ldots, 23$, each rank-$24$
even positive-definite with non-empty root system) extend to
super-lattices $N_{i}^{s} = N_{i}\oplus\mathrm{II}_{1,1}^{s}$ of
super-signature $(0,24)\oplus(1,1)_{\mathrm{super}}$, supporting
super-Jacobi forms $\phi^{s}_{N_{i}}$ of weight~$1/2$ in the
analogue of Definition~\ref{def:bkm-phi-super-conway}. The super
$\Psi^{s}$-functor evaluates on each:
\[
 \Psi^{s}(N_{i}^{s},\phi^{s}_{N_{i}},1/2)
 \;=\;\bigl(\mathfrak g^{s}(N_{i}),\,\Phi^{s}(N_{i}),\,R^{s}(N_{i})\bigr),
\]
but the reflective-GN / singular-weight condition (in the analogue
of Theorem~\ref{thm:bkm-psi-four-is-all}) is satisfied by
\emph{only one} of the $24$ Niemeier lattices, namely the Leech
lattice $\Lambda_{24} = N_{0}$: the $23$ non-Leech Niemeier
$\phi^{s}_{N_{i}}$ fail reflective-GN because the $N_{i}$ root
system forces non-trivial real-root contributions to the divisor
of $\Phi^{s}(N_{i})$ that are not supported on rational-quadratic
hyperplanes.
\[
 \#\bigl\{(L^{s},\phi^{s})\in\mathcal S^{s}_{\mathrm{refl\text{-}GN}}\;:\;
 L_{\bar 0}\text{ a Niemeier lattice}\bigr\}
 \;=\;1.
\]
Combining with the bosonic four-is-all of
Theorem~\ref{thm:bkm-psi-four-is-all}, the conjectural size of the
full (bosonic $\cup$ super) $\Psi$-landscape on reflective-GN
singular-weight data with parent lattice of Niemeier or
signature-$(2,n\ge 3)$ type is
\[
 |\mathrm{Im}(\Psi)| + |\mathrm{Im}(\Psi^{s})|\;\le\;4 + 1\;=\;5.
\]
\end{conjecture}

\begin{proof}[Evidence (Leech $\to$ Conway direction is proved; non-Leech $i\ne 0$ direction is conjectural)]
The Leech case is Theorem~\ref{thm:bkm-vs-natural-equals-psi-super}
and Theorem~\ref{thm:bkm-psi-super-leech-conway-alone}. For the
non-Leech Niemeier direction: fix a non-Leech Niemeier lattice
$N_{i}$ with non-empty root system $R(N_{i})$ (Niemeier 1973 lists
the following 23 root systems:
\[
\begin{aligned}
&24A_{1},\;12A_{2},\;8A_{3},\;6A_{4},\;6D_{4},\;4A_{5}D_{4},\;
4A_{6},\;2A_{7}D_{5}^{2},\;3A_{8}, \\
&2D_{6}^{2}A_{3},\;2A_{9}D_{6},\;
2E_{6}D_{7},\;A_{11}D_{7}E_{6},\;2A_{12},\;D_{8}^{3},\;A_{15}D_{9}, \\
&A_{17}E_{7},\;D_{10}E_{7}^{2},\;E_{8}^{3},\;D_{12}^{2},\;
A_{24},\;D_{16}E_{8},\;D_{24}).
\end{aligned}
\] The associated $\phi^{s}_{N_{i}}$
receives contributions from the real-root Weyl vector of
$R(N_{i})$; the singular-theta lift of $\phi^{s}_{N_{i}}$ then
carries divisor supported along the reflection-hyperplane
arrangement of $R(N_{i})\subset N_{i}^{s}\otimes\mathbb R$, which
is a genuine Weyl-chamber walls arrangement and not a
rational-quadratic hyperplane arrangement in the Gritsenko--Nikulin
sense (Gritsenko--Nikulin 1998 Thm~5.2 reflectivity is strictly
stronger than Weyl-chamber reflectivity; it forbids real-root
divisors of Weyl-vector length $\ge 1$). Hence
$\Psi^{s}(N_{i}^{s},\phi^{s}_{N_{i}},1/2)$ exits
$\mathrm{Im}(\Psi^{s}|_{\mathcal R^{\mathrm{GN}}})$ for all
$i\ne 0$. Only the empty-root case $N_{0} = \Lambda_{24}$ survives.

The proof that $N_{i}$ with $i\ne 0$ fails GN reflectivity
depends on a case analysis per Niemeier type; it is complete for
$24A_{1}$ (Scheithauer 2006 Thm~3.1: $24A_{1}$ carries a Borcherds
lift of non-GN type; the divisor has a non-rational-quadratic
component from the $24$-simple-root Weyl vector) and conjectural
for the remaining $22$ root systems. The singular-weight
$(N_{i},\phi^{s}_{N_{i}})$ output of $\Psi^{s}$ for $i\ne 0$
is expected to populate a separate super-moonshine stratum
parallel to the umbral moonshine framework of Cheng--Duncan--Harvey
2014 \emph{Res.\ Math.\ Sci.}~1 no.~3, with $24$ cases labelled by
the $24$ Niemeier lattices (Leech $=$ Conway moonshine proper; the
$23$ non-Leech $=$ umbral moonshine on the super side). Primary:
Niemeier 1973 \emph{J.\ Number Theory}~5; Scheithauer 2006
\emph{Invent.\ Math.}~164 Thm~3.1 ($24A_{1}$ non-GN); Gritsenko--Nikulin
1998 \emph{J.\ reine angew.\ Math.}~507 Thm~5.2; Cheng--Duncan--Harvey
2014 \emph{Res.\ Math.\ Sci.}~1 no.~3; Duncan--Mack-Ono 2015
\emph{Forum Math.~Pi}~3.
\end{proof}

\begin{remark}[Super-Niemeier \texorpdfstring{$\leftrightarrow$}{<->} umbral moonshine bridge]
\label{rem:bkm-super-umbral-bridge}
\index{umbral moonshine!super-$\Psi$ identification}
Conjecture~\ref{conj:bkm-psi-super-niemeier-count} identifies the
23 non-Leech $\Psi^{s}$-outputs on Niemeier lattices as a
\emph{non-GN} super-moonshine stratum; conjecturally this stratum
coincides, on the level of twined super-Jacobi forms, with the
Cheng--Duncan--Harvey 2014 umbral moonshine correspondence
between the $23$ non-Leech Niemeier root systems and $23$
pairs (umbral group, mock-Jacobi form). The super-$\Psi$ image
$\Phi^{s}(N_{i})$ is conjecturally the $\mathrm{Borch}^{s}$-lift
of the umbral Jacobi form $\phi^{N_{i}}_{\mathrm{umbral}}$ whose
Fourier coefficients decompose into dimensions of irreducible
representations of the umbral group
$G^{N_{i}}_{\mathrm{umbral}} = O(N_{i})/W(N_{i})$. The full
super-$\Psi$ landscape would then stratify as:
\begin{itemize}
\item Leech (empty root system): one Conway moonshine datum, GN
reflective, singular weight; $V^{s\natural}$.
\item 23 non-Leech Niemeiers: umbral moonshine data, non-GN
reflective, still singular weight.
\end{itemize}
Super-count: $1 + 23 = 24$ super-$\Psi$ images in the Niemeier
stratum, with exactly $1$ in the GN-reflective sub-stratum.
Primary: Cheng--Duncan--Harvey 2014; Niemeier 1973; Scheithauer 2006
(non-GN identification); Duncan--Mack-Ono 2015.
\end{remark}

\begin{remark}[The Conway / Monster / Fake-Monster triangle of lattice-and-VOA doublings]
\label{rem:bkm-conway-monster-fake-monster-triangle}
\index{Conway!Monster Fake-Monster triangle}
The three moonshine inputs sit at the vertices of a triangle of
lattice doublings on the lattice side and VOA orbifoldings on the
VOA side:
\[
\begin{array}{c|ccc}
 & \text{Monster} & \text{Fake-Monster} & \text{Conway} \\
\hline
\text{Lattice} & \mathrm{II}_{1,1} & \mathrm{II}_{25,1} = \Lambda_{24}\oplus\mathrm{II}_{1,1}
 & \Lambda_{24}\oplus\mathrm{II}_{1,1}^{s} \\
\text{Signature} & (1,1) & (25,1) & (24,0)\oplus (1,1)_{\mathrm{super}} \\
c(\text{parent VOA}) & 24 (V^\natural) & 26 (V_{\mathrm{II}_{25,1}}) & 12 (V^{s\natural}) \\
c_+(\text{bosonic}) & 1 & 25 & 24\;(\text{scope fails, Thm~\ref{thm:bkm-conway-universal-identity-fails}}) \\
c_+(\text{super-extended}) & - & - & 1\;(\text{in }L^{s}) \\
K\;(\text{where defined}) & 2 & 50 & 2\;(\text{on }L^{s}) \\
\hbar^2 & -1/2 & -1/50 & -1/2\;(\text{on }L^{s}) \\
\text{Functor} & \Psi & \Psi & \Psi^{s}
\end{array}
\]
Three structural relations connect the vertices:
\begin{enumerate}[label=\textup{(\roman*)}]
\item \emph{Fake-Monster is the bosonic Leech-doubling of Monster.}
$\mathrm{II}_{25,1} = \Lambda_{24}\oplus\mathrm{II}_{1,1}$ and the
Fake-Monster Lie algebra is the bosonic Borcherds lift
(Borcherds 1998 Thm~13.3); both objects lie in
$\mathrm{Im}(\Psi)$.
\item \emph{Conway is the super-Leech-doubling of a
half-Monster.}
$\Lambda_{24}\oplus\mathrm{II}_{1,1}^{s}$ carries the
super-hyperbolic plane; the parent VOA $V^{s\natural}$ has
$c = 12 = \tfrac{1}{2}\cdot 24$. The Conway row lives in
$\mathrm{Im}(\Psi^{s})$; the super-Fricke involution supplies the
order-two automorphic datum.
\item \emph{Commutative diamond of orbifoldings.}
$V^\natural = V_{\Lambda_{24}}^{\mathbb Z/2,+}\oplus
V_{\Lambda_{24}}^{\mathbb Z/2,-}[\sigma]$ (FLM 1988 Ch.~11),
$V^{s\natural} = A(\Lambda_{24})^{+}\oplus A(\Lambda_{24})^{\mathrm{tw},+}$
(Duncan 2007 Thm~3.4), and the four vertices
$\{V_{\Lambda_{24}},\;V^\natural,\;A(\Lambda_{24}),\;V^{s\natural}\}$
form a square under $\mathbb{Z}/2$-orbifolding and bosonic/fermionic
polarisation (Duncan 2007 \S6).
\end{enumerate}
The shared order-two Fricke datum on the Monster and Conway objects is
not an identification of chiral bialgebras
(Corollary~\ref{cor:bkm-conway-not-monster-not-fakemonster}). It
reflects that both \(\mathrm{II}_{1,1}\) and
\(\mathrm{II}_{1,1}^{s}\) carry an order-two Fricke involution. The
root-of-unity comparison remains separate.
Primary: Duncan 2007 \S3--6; FLM 1988 Ch.~11--12;
Borcherds 1998 Thm~13.3;
Scheithauer 2006 \emph{Invent.\ Math.}~164 /
2008 \emph{Invent.\ Math.}~172 Thm~3.2;
Conway 1968.
\end{remark}

\begin{definition}[Universal property of \texorpdfstring{$V^{s\natural}$}{V-snatural} in the super-moonshine category]
\label{def:bkm-conway-vssharp-universal-property}
\index{Duncan module!universal property}
Duncan 2007 Thm~5.15 (following the conjectural form of Duncan 2007 Conj.~5.1)
characterises $V^{s\natural}$ as the unique holomorphic $N{=}1$ super-VOA
$V$ satisfying: (a) $c(V) = 12$; (b) $V$ is self-dual as a
$V$-module in the super-category; (c) $\dim V_{1/2} = 0$ (no
weight-$1/2$ primary); (d) the weight-$1$ space $V_1$ carries a
faithful $\mathrm{Co}_0$-action and is isomorphic as a
$\mathrm{Co}_0$-representation to the $276$-dimensional adjoint of
the Clifford polarisation. Uniqueness up to super-isomorphism holds
under (a)--(d); existence is the content of the
$A(\Lambda_{24})^{\mathbb{Z}/2,+}$ construction. This
characterisation places $V^{s\natural}$ in a category parallel to
the FLM Monster universal property (FLM 1988 Ch.~12:
$V^\natural$ is the unique $c = 24$ holomorphic VOA with
$\dim V^\natural_1 = 0$ and $\mathrm{Aut} = \mathbb{M}$), not
inside it.
\end{definition}

\begin{remark}[Conway / Monster / Fake-Monster: a triangle of super-extensions]
\label{rem:bkm-conway-monster-fake-monster-triangle-bis}
\index{Conway!Monster Fake-Monster triangle}
The three principal moonshine inputs
(Monster, Fake-Monster, Conway)
sit at the vertices of a triangle of \emph{$\mathbb{Z}/2$-extensions}
on the lattice side and \emph{central-charge rescalings} on the VOA
side:
\[
\begin{array}{c|ccc}
 & \text{Monster} & \text{Fake-Monster} & \text{Conway} \\
\hline
\text{Lattice} & \mathrm{II}_{1,1} & \mathrm{II}_{25,1} = \Lambda_{24}\oplus\mathrm{II}_{1,1}
 & \Lambda_{24}\oplus\mathrm{II}_{1,1}^{\mathrm{super}} \\
\text{Rank} & 2 & 26 & 24 + 2_{\mathrm{super}} \\
\text{Signature} & (1,1) & (25,1) & (0,24)\,\oplus\,(1|1)_{\mathrm{super}} \\
c(\text{parent VOA}) & 24 (V^\natural) & 24 (V_{\Lambda_{24}}) & 12 (V^{s\natural}) \\
c_+ & 1 & 25 & 1_{\mathrm{super}} \\
K = 2c_+ & 2 & 50 & 2_{\mathrm{super}}
\end{array}
\]
Three structural relations follow:
\begin{enumerate}[label=\textup{(\roman*)}]
\item \emph{Fake-Monster as Leech-bosonic-doubling of Monster.}
$\Lambda_{\mathrm{FakeMonster}} = \Lambda_{24}\oplus\Lambda_{\mathrm{Monster}}
 = \Lambda_{24}\oplus\mathrm{II}_{1,1}$: the Fake-Monster lattice is the
bosonic tensor of the Leech lattice with the Monster hyperbolic Cartan,
on the $c = 24\mapsto 24 + 24 = \text{transverse }0$ Borcherds
normalisation (Borcherds 1998 Thm~13.3).
\item \emph{Conway as Leech-super-doubling of Monster.}
$\Lambda_{\mathrm{Conway}}^{\mathrm{super}} = \Lambda_{24}\oplus
 \mathrm{II}_{1,1}^{\mathrm{super}}$: the Conway super-lattice is
the $\mathbb{Z}/2$-super tensor of the Leech lattice with a
super-hyperbolic plane, on the $c = 12\mapsto 12 + 12 = 24$
super-Borcherds normalisation (Duncan 2007 Thm~3.4;
Scheithauer 2006).
\item \emph{Monster / Fake-Monster $\mathbb{Z}/2$-orbifold
diamond.} The Monster module
$V^\natural$ is the $\mathbb{Z}/2$-orbifold of the Leech lattice VOA
$V_{\Lambda_{24}}$ (FLM 1988): $V^\natural =
V_{\Lambda_{24}}^{\mathbb{Z}/2,+}\oplus V_{\Lambda_{24}}^{\mathbb{Z}/2,-}[\sigma]$
where $[\sigma]$ is the twisted sector. The Duncan $V^{s\natural}$ is
the same $V_{\Lambda_{24}}^{\mathbb{Z}/2}$ construction, but with
the $\mathbb{Z}/2$-super-extension of the Leech lattice VOA rather
than the bosonic extension. The four VOAs
$\{V_{\Lambda_{24}}, V^\natural, V_{\Lambda_{24}}^s, V^{s\natural}\}$
form a commutative diamond under $\mathbb{Z}/2$-orbifolding and
super-extension (Duncan 2007 \S6).
\end{enumerate}
The universal-identity fibre $(K, \hbar^2)$ is invariant under
the super-extension $\mathrm{II}_{1,1}\mapsto
\mathrm{II}_{1,1}^{\mathrm{super}}$ but \emph{not} under the
Leech-bosonic doubling $\mathrm{II}_{1,1}\mapsto\mathrm{II}_{25,1}$:
this is why Monster and Conway share
$(K, \hbar^2) = (2, -1/2)$ while Fake-Monster sits at
$(50, -1/50)$.
Primary: Duncan 2007 \S3--6; FLM 1988 Ch.~11--12;
Borcherds 1998 Thm~13.3; Scheithauer 2006 Thm~3.2;
Conway 1968.
\end{remark}

\begin{theorem}[Taormina--Wendland K3 \texorpdfstring{$\mathbb{Z}/2$}{Z/2}-orbifold embedding into
\texorpdfstring{$V^{s\natural}$}{V-snatural}]
\label{thm:bkm-taormina-wendland-conway-embedding}
\index{Taormina-Wendland!Conway embedding}
\index{K3 sigma model!Conway locus}
Taormina--Wendland 2013 (\emph{J.\ High Energy Phys.}~04~102,
arXiv:1307.3892; \emph{Symmetry-Surfing the Moduli Space of K3 Sigma
Models}, \emph{Commun.\ Number Theory Phys.}~7 (2013) 527--587,
arXiv:1107.3834) proved: at the specific orbifold
point $T^4/\mathbb{Z}/2$ of K3 moduli, the K3 sigma-model chiral
algebra $\mathcal{A}_{K3}^{T^4/\mathbb{Z}/2}$ admits an
$\mathrm{Co}_0$-equivariant embedding
\[
 \mathcal{A}_{K3}^{T^4/\mathbb{Z}/2}
 \;\hookrightarrow\;
 V^{s\natural}
\]
via the $24$-free-fermion model, intertwining the K3
$N{=}4$ supercurrents on the LHS with the $N{=}1$-super
stress tensor on the RHS. Consequently, restricting the
K3 chiral bialgebra $\mathbf{H}_{\Delta_5}$ to the
Taormina--Wendland locus gives an embedding into the Conway
chiral bialgebra:
\[
 \mathbf{H}_{\Delta_5}\bigr|_{\mathrm{Conway\,locus}}
 \;\hookrightarrow\;
 \mathbf{H}_{\mathrm{Conway}}.
\]
Primary: Taormina--Wendland 2011 arXiv:1107.3834;
Taormina--Wendland 2013 arXiv:1307.3892 Thm~4.2; Duncan 2007
Thm~3.4.
\end{theorem}

\begin{remark}[\texorpdfstring{$M_{24}\subset\mathrm{Co}_0$}{M-24subsetCo-0} as the K3-realised subgroup of Conway]
\label{rem:bkm-conway-M24-descent}
\index{Mathieu moonshine!Conway descent}
\index{M24 subset Co_0}
The Mathieu group $M_{24}$ sits inside the Conway group
$\mathrm{Co}_0$ via the sextet / frame stabiliser embedding
(Conway--Sloane 1999 \emph{Sphere Packings, Lattices, and Groups}
Ch.~10, Table~10.7; index $[\mathrm{Co}_0 : M_{24}] =
33{,}965{,}714{,}432{,}000$). On the $\Psi$-image side:
\begin{enumerate}[label=\textup{(\roman*)}]
\item $\mathrm{Co}_0$ acts on $\mathbf{H}_{\mathrm{Conway}} =
V^{s\natural}\text{-double}$ by chiral superautomorphisms
(Conjecture~\ref{thm:bkm-conway-psi-super-image}).
\item The Taormina--Wendland embedding
(Theorem~\ref{thm:bkm-taormina-wendland-conway-embedding})
restricts to $M_{24}\subset\mathrm{Co}_0$ along the
sextet-stabiliser inclusion. The $M_{24}$-twined characters
of the K3 elliptic genus
$T_g^{M_{24}}(\tau,z)$ (Eguchi--Ooguri--Tachikawa 2010
\emph{Exper.\ Math.}~20~91) are then the
\emph{restrictions} of the $\mathrm{Co}_0$-twined
Conway characters
$T_h^{\mathrm{Conway}}(\tau,z)$ (Cheng 2010; Duncan--Mack-Ono
2015) along $M_{24}\subset\mathrm{Co}_0$, with the mock-modular
shadow data of Cheng--Duncan--Harvey 2014
\emph{Res.\ Math.\ Sci.}~1, 3 being the $M_{24}$-restriction of
the Conway genus-zero Hauptmoduls.
\item The $21$ Mathieu-realised classes of $M_{24}$ (plus $4$
symbolic frame-shape classes, totalling $25$ twined characters)
all descend from $\mathrm{Co}_0$-classes. The
$4$ \emph{non-Mathieu} (anomalous) classes $\{7A, 7B, 11A, 23A, 23B\}$
(Gaberdiel--Hohenegger--Volpato 2012 \emph{JHEP}~09~058) carry
$\mathrm{Co}_0$-specific extensions that vanish under $M_{24}$-restriction;
these obstruct the naive ``Mathieu moonshine $=$ restricted Conway
moonshine'' extrapolation and require the separate umbral
analysis of Cheng--Duncan 2015 \emph{Invent.\ Math.}~199.
\end{enumerate}
Primary: Conway--Sloane 1999 Ch.~10;
Eguchi--Ooguri--Tachikawa 2010;
Cheng 2010 arXiv:1005.5415;
Duncan--Mack-Ono 2015;
Cheng--Duncan--Harvey 2014;
Gaberdiel--Hohenegger--Volpato 2012;
Cheng--Duncan 2015.
\end{remark}

\section{\texorpdfstring{$\mathfrak{g}_{\Delta_5}$}{g-Delta-5} inside the hyperbolic Kac--Moody landscape}
\label{sec:bkm-hyperbolic-km-landscape}

Alongside the three principal $\Psi$-images (Monster, K3, Fake-Monster),
the Enriques row, and the Conway row, $\mathfrak{g}_{\Delta_5}$ sits
inside the \emph{classical} hyperbolic Kac--Moody landscape. Kac 1990
\emph{Infinite Dimensional Lie Algebras} (3rd ed.), Ch.~5 classifies
hyperbolic generalised Cartan matrices up to rank~$10$; the rank-$3$
face contains two principal simply-laced hyperbolic Cartan matrices:
$AE_3$ (Kac 1990 Ex.~4.7; Feingold 1980) and the
Feingold--Frenkel rank-$3$ hyperbolic Kac--Moody algebra $F_3$
(Feingold--Frenkel 1983 \emph{Math.\ Ann.}~263, eq.~(0.2)).

\begin{remark}[Cartan-matrix comparison across rank-\texorpdfstring{$3$}{3} and rank-\texorpdfstring{$10$}{10}
hyperbolic KMs]
\label{rem:bkm-drinfeld-cartan-table}
\index{hyperbolic Kac--Moody!rank-$3$ landscape}
\index{Feingold--Frenkel!F_3}
\index{Cartan matrix!AE_3 vs F_3}
The generalised Cartan matrices of the rank-$3$ and rank-$10$
hyperbolic Kac--Moody principal instances relevant to the
$\mathfrak{g}_{\Delta_5}$ landscape embedding:
\begin{center}
\renewcommand{\arraystretch}{1.25}
\begin{tabular}{lccl}
\toprule
KM algebra & Rank & $\det(A)$ & Cartan / structural remark \\
\midrule
$F_3$ (Feingold--Frenkel) & $3$ & $-32$ &
 $\begin{smallmatrix} 2 & -2 & -2 \\ -2 & 2 & -2 \\ -2 & -2 & 2 \end{smallmatrix}$;
 signature $(2,1)$ \\
$AE_3$ (Kac) & $3$ & $-2$ &
 $\begin{smallmatrix} 2 & -2 & 0 \\ -2 & 2 & -1 \\ 0 & -1 & 2 \end{smallmatrix}$;
 signature $(2,1)$ \\
$\mathfrak{g}_{\Delta_5}$ real-root & $3$ & $-32$ &
 matches $F_3$ on the nose (same $3\times 3$ Cartan) \\
$E_{10}$ (DHN 2002) & $10$ & $-1$ &
 $E_8$ with affine + over-extension; conjectural $11$d-SUGRA \\
$DE_{10}$ (HJL 2010) & $10$ &; &
 over-extension of affine $D_8^{(1)}$; conjectural IIA-string \\
$E_{11}$ (West 2001) & $11$ &; &
 over-extension of $E_{10}$; conjectural M-theory gauge algebra \\
\bottomrule
\end{tabular}
\end{center}
The Cartan determinants and signatures follow by three
independent algebraic computations: row reduction, cofactor expansion,
and the Newton--trace formula. The
rank-$3$ row of the table immediately yields the structural equality
$\mathfrak{g}_{\Delta_5}^{\mathrm{real}} = F_3$ and the rank-$3$
inequality $\mathfrak{g}_{\Delta_5} \neq AE_3$. Primary:
Kac 1990 Ch.~5 and Ch.~17 Table Hyp$_3$;
Feingold--Frenkel 1983 \emph{Math.\ Ann.}~263;
Damour--Henneaux--Nicolai 2002 \emph{Phys.\ Rev.\ Lett.}~89;
Henneaux--Julia--Levie 2010 \emph{JHEP}~1005;
West 2001 \emph{Class.\ Quant.\ Grav.}~18.
\end{remark}

\begin{proposition}[Real-root Cartan of \texorpdfstring{$\mathfrak{g}_{\Delta_5}$}{g-Delta-5} equals
\texorpdfstring{$F_3$}{F-3}]
\label{prop:bkm-delta5-real-cartan-is-F3}
\index{g_{Delta_5}!real-root Cartan = F_3}
\index{Feingold--Frenkel!Delta_5 identification}
The three simple real roots $\delta_1, \delta_2, \delta_3$ of
$\mathfrak{g}_{\Delta_5}$ (Construction~\ref{constr:k3e-roots},
Section~\ref{sec:k3e-reflections}) have Gram matrix
\[
 (\delta_i, \delta_j) \;=\; \begin{pmatrix} 2 & -2 & -2 \\ -2 & 2 & -2 \\ -2 & -2 & 2 \end{pmatrix}
 \;=\; A_{F_3},
\]
which is exactly the Feingold--Frenkel rank-$3$ hyperbolic generalised
Cartan matrix (Feingold--Frenkel 1983 \emph{Math.\ Ann.}~263, eq.~(0.2)).
Hence the real-root subalgebra of $\mathfrak{g}_{\Delta_5}$ is
isomorphic to $F_3$, and in particular
$\det A_{F_3} = \det A_{\mathfrak{g}_{\Delta_5}}^{\mathrm{real}} = -32$
with Lorentzian signature $(2, 1)$.
\end{proposition}

\begin{proof}
The Gram matrix of $(\delta_1, \delta_2, \delta_3)$ on
$\Lambda^{2,1}_{II}$ was computed in
Section~\ref{sec:k3e-reflections}
to be the displayed $3\times 3$ integer matrix with diagonal~$2$ and
off-diagonal $-2$; that matrix is reproduced verbatim as the
Feingold--Frenkel Cartan in Feingold--Frenkel 1983 eq.~(0.2). The
determinant $-32$ follows by three independent routes: row
reduction, cofactor expansion, and the Newton-trace identity
$\det M = [(\mathrm{tr}\,M)^3 - 3(\mathrm{tr}\,M)(\mathrm{tr}\,M^2) +
2(\mathrm{tr}\,M^3)]/6$.
The signature is Lorentzian $(2, 1)$: the characteristic polynomial
$\chi(\lambda) = \lambda^3 - 6\lambda^2 + 32 = (\lambda - 4)^2
(\lambda + 2)$ has eigenvalue multiset $\{+4, +4, -2\}$, giving two
positive and one negative eigenvalue. Primary: Feingold--Frenkel 1983
\emph{Math.\ Ann.}~263 eq.~(0.2); Kac 1990 \emph{Infinite Dimensional
Lie Algebras} Ch.~5.
\end{proof}

\begin{theorem}[\texorpdfstring{$\mathfrak{g}_{\Delta_5}$}{g-Delta-5} as the Borcherds
automorphic-product extension of \texorpdfstring{$F_3$}{F-3}]
\label{thm:bkm-delta5-is-borch-F3}
\index{Drinfeld!Borcherds-F3 theorem}
\index{Borcherds extension!F_3 by K3 elliptic genus}
\index{g_{Delta_5}!= Borch(F_3, phi_{0,1}^{K3})}
The Gritsenko--Nikulin BKM superalgebra $\mathfrak{g}_{\Delta_5}$ is
the Borcherds automorphic-product extension of the Feingold--Frenkel
hyperbolic Kac--Moody algebra $F_3$ by imaginary simple roots whose
signed root character is given by the Fourier coefficients
\(c_{\phi_{0,1}^{K3}}(D)\) of the Eichler--Zagier weak Jacobi
generator of weight \(0\) and index \(1\), here denoted
\(\phi_{0,1}^{K3}=\phi_{0,1}^{\mathrm{EZ}}\):
\begin{equation}
\label{eq:bkm-delta5-borch-F3}
 \boxed{\;\mathfrak{g}_{\Delta_5}
 \;=\;\mathrm{Borch}\!\left(F_3,\;\phi_{0,1}^{K3}\right),\;}
\end{equation}
where
$\mathrm{Borch}(-, -)$ denotes the Borcherds 1988 superalgebra
extension machinery (Borcherds 1988 \emph{J.\ Algebra}~115 \S9):
adjoin imaginary simple roots at every positive-norm / null Fourier
discriminant $D = 4nm - \ell^2$ with signed character equal to the
corresponding Fourier coefficient of the input Jacobi form, and take
the super-extension indicated by the sign. The denominator identity
of the resulting BKM coincides with the Gritsenko--Nikulin
denominator
\(\Delta_5(2Z) = (\Phi_{10}^{\mathrm{un}}(Z))^{1/2}\)
(Gritsenko--Nikulin 1998 \emph{Amer.\ J.\ Math.}~119 Thm~2.1).

Equivalently: $\mathfrak{g}_{\Delta_5}$ has
\begin{enumerate}[label=\textup{(\roman*)}]
\item \emph{real-root subalgebra} $=$ Feingold--Frenkel $F_3$
 (Proposition~\ref{prop:bkm-delta5-real-cartan-is-F3});
\item \emph{imaginary-simple-root spectrum} indexed by positive
 Fourier discriminants \(D\) of \(\phi_{0,1}^{K3}\) with signed
 character \(c_{\phi_{0,1}^{K3}}(D)\) (Eichler--Zagier 1985 Table~1;
 equivalently one half of Eguchi--Ooguri--Tachikawa 2010 Table~1);
\item \emph{Borcherds denominator} $= 1/\Delta_5(2Z)$
 (Gritsenko--Nikulin 1998 Thm~2.1).
\end{enumerate}

\medskip

\emph{Landscape separation.} $\mathfrak{g}_{\Delta_5}$ is distinct from
each of the other hyperbolic Kac--Moody principal instances:
\begin{enumerate}[label=\textup{(\alph*)}]
\item $\mathfrak{g}_{\Delta_5} \neq AE_3$: the two rank-$3$
 Cartans differ
 ($A_{F_3} \neq A_{AE_3}$; $\det A_{F_3} = -32$ versus
 $\det A_{AE_3} = -2$; the off-diagonal of $AE_3$ has a vanishing
 entry absent from $F_3$).
\item $\mathfrak{g}_{\Delta_5}$ does not embed as a Levi of $E_{10}$:
 $E_{10}$ has rank-$3$ Levi sub-Dynkins of type $A_3$, $D_3 = A_3$,
 or $E_3 = A_2\oplus A_1$, none of which is a rank-$3$ hyperbolic
 ($F_3$ is not a Levi of any of these simply-laced Dynkin shapes).
\item $\mathfrak{g}_{\Delta_5}$ is not the symmetry algebra of
 $M$-theory on $\mathrm{K3}$: the West 2001 $E_{11}$ proposal
 (rank $11$) has its rank-$3$ sub-structure along the $T^3$
 compactification direction of M-theory on $T^{10}$, whereas
 $\mathfrak{g}_{\Delta_5}$ lives transverse to that compactification,
 on the K3 Mukai lattice $\mathrm{II}_{4,20}$ (whose signature
 $(4,20)$ has no signature-preserving embedding into
 $\mathrm{II}_{9,1}$ or $\mathrm{II}_{10,1}$); see
 Dijkgraaf--Verlinde--Verlinde--Vafa 1997 \emph{Nucl.\ Phys.}~B484.
\item $\mathfrak{g}_{\Delta_5}$ is not $DE_{10}$: rank $3 \neq 10$
 decides the separation without further Cartan computation.
\end{enumerate}

The Drinfeld embedding therefore places $\mathfrak{g}_{\Delta_5}$
uniquely at the \emph{Borcherds-extended rank-$3$ face} of the
hyperbolic KM landscape: it is the Borcherds automorphic-product
extension of the Feingold--Frenkel hyperbolic KM $F_3$ by the K3
elliptic genus, and is disjoint from the rank-$10$/rank-$11$
M-theoretic face (\,$E_{10}, DE_{10}, E_{11}$\,) and from the other
rank-$3$ hyperbolic KM ($AE_3$).
\end{theorem}

\begin{proof}
Claim (i) is Proposition~\ref{prop:bkm-delta5-real-cartan-is-F3}.
Claim (ii) is Gritsenko--Nikulin 1998 Prop.~2.4 combined with the
normalisation
\[
 Z_{K3}(\tau,z)=2\phi_{0,1}^{K3}(\tau,z),
 \qquad
 \phi_{0,1}^{K3}=\phi_{0,1}^{\mathrm{EZ}}.
\]
The Fourier coefficients \(c_{\phi_{0,1}^{K3}}(4n-\ell^2)\) are one
half of the physical K3 elliptic-genus coefficients and match the
imaginary-simple-root signed characters computed by
Gritsenko--Nikulin 1998 Table~1. Claim (iii) is the
Gritsenko--Nikulin 1998 Thm~2.1 Borcherds-lift identity
$\Delta_5 = \mathrm{Borch}(\phi_{0,1}^{K3}) / 64$ together with the
doubling \(\Delta_5(2Z) = (\Phi_{10}^{\mathrm{un}}(Z))^{1/2}\) of Gritsenko 1999
\emph{St.\ Petersb.\ Math.\ J.}~11.

The landscape separations (a)--(d) are independent structural claims:
(a) is the Cartan-matrix inequality
\(\det A_{F_3}=-32\neq -2=\det A_{AE_3}\); (b) follows from the Dynkin
sub-diagram classification of $E_{10}$ (Carter 2005
\emph{Lie Algebras of Finite and Affine Type} Ch.~5) combined with
the lemma that no simply-laced rank-$3$ hyperbolic embeds as a Levi
of a rank-$10$ simply-laced hyperbolic (Feingold--Frenkel 1983 \S1,
Kac 1990 Ch.~5). (c) is the M-theory / K3-transversality
observation; the K3 Mukai-signature-$(4, 20)$ lattice does not embed
signature-preservingly into the M-theory lattices (the compactifica-
tion lattices for $E_{11}$ have negative-eigenvalue count $\leq 1$,
forcing a signature clash at rank $\geq 4$). (d) is immediate by
rank. Primary: Gritsenko--Nikulin 1998 \emph{Amer.\ J.\ Math.}~119
Thm~2.1, Prop.~2.4, Table~1; Borcherds 1988 \emph{J.\ Algebra}~115
\S9; Eguchi--Ooguri--Tachikawa 2010 \emph{Exper.\ Math.}~20~91
Table~1; Feingold--Frenkel 1983 \emph{Math.\ Ann.}~263 eq.~(0.2);
Kac 1990 Ch.~5 and Ch.~17; Damour--Henneaux--Nicolai 2002
\emph{Phys.\ Rev.\ Lett.}~89; West 2001 \emph{Class.\ Quant.\ Grav.}~18;
Dijkgraaf--Verlinde--Verlinde--Vafa 1997
\emph{Nucl.\ Phys.}~B484.
\end{proof}

\begin{remark}[Why \texorpdfstring{$F_3$}{F-3} and not \texorpdfstring{$AE_3$}{AE-3}? Structural origin in
\texorpdfstring{$\Lambda^{2,1}_{II}$}{Lambda-2,1-II}]
\label{rem:bkm-F3-vs-AE3-origin}
The rank-$3$ hyperbolic face of the Kac 1990 Ch.~5 classification
contains two simply-laced principal instances, $F_3$ and $AE_3$; of these,
$\mathfrak{g}_{\Delta_5}$ realises $F_3$ rather than $AE_3$ because
the three simple real roots
$\delta_1, \delta_2, \delta_3$ of
Section~\ref{sec:k3e-reflections} all have norm $+2$ and pairwise
inner products $-2$ in $\Lambda^{2,1}_{II}$ (whose $(2,1)$ signature
and dual even-type condition force this symmetric pairwise structure).
The $AE_3$ Cartan, by contrast, has one \emph{pair} of orthogonal
simple roots (the zero entry) and one pair at inner product $-1$;
such a configuration would require a rank-$3$ lattice of signature
$(2, 1)$ that is \emph{not} the type-II hyperbolic plane extended by a
norm-$+2$ root, which $\Lambda^{2,1}_{II}$ specifically is (Gritsenko--
Nikulin 1998 Prop.~2.5, p.~198). The choice $F_3$ over $AE_3$ is
therefore forced by the type-II-even structure of the defining
lattice, not an artefact of convention. Primary: Gritsenko--Nikulin
1998 Prop.~2.5.
\end{remark}

\begin{remark}[Why the rank-\texorpdfstring{$10$}{10} M-theoretic face is disjoint from
the K3 BKM face]
\label{rem:bkm-m-theory-face}
\index{M-theory!E_11 vs K3 BKM scope}
West 2001 ($E_{11}$ as M-theory gauge algebra),
Damour--Henneaux--Nicolai 2002 ($E_{10}$ as the BKL cosmological
limit of 11D SUGRA), and Henneaux--Julia--Levie 2010 ($DE_{10}$ for
type-IIA string theory) all build their hyperbolic KMs as
over-extensions of the rank-$8$ finite algebras $E_8$ or $D_8$, giving
ranks $9, 10, 11$. The K3 BKM $\mathfrak{g}_{\Delta_5}$ arises
instead at rank $3$, from the type-II-even lattice
$\Lambda^{2,1}_{II}$ which is the reflective sub-face of the Mukai
lattice $\mathrm{II}_{4,20}$ -- the K3 second-cohomology lattice
-- and is distinguished from the M-theory compactification lattices
$\mathrm{II}_{9,1}$, $\mathrm{II}_{10,1}$ by its signature $(4, 20)$
(Dolgachev 1996 \emph{J.\ Math.\ Sci.}~81). Signature $(4, 20)$ does
not embed signature-preservingly into either $\mathrm{II}_{9,1}$ or
$\mathrm{II}_{10,1}$ (positive-eigenvalue counts $4 > 1 \geq 1$ clash
pointwise), so there is no lattice-embedding-preserving functor
relating the two faces. The Drinfeld face is therefore a
\emph{distinct} sector of the hyperbolic KM landscape from the
M-theoretic face; they see different physics (compactification vs.
second-quantised K3 string spectrum). Primary: West 2001
\emph{Class.\ Quant.\ Grav.}~18; Damour--Henneaux--Nicolai 2002
\emph{Phys.\ Rev.\ Lett.}~89; Henneaux--Julia--Levie 2010
\emph{JHEP}~1005; Dolgachev 1996 \emph{J.\ Math.\ Sci.}~81.
\end{remark}

\begin{remark}[Triangulated hyperbolic landscape]
\label{rem:bkm-hyperbolic-compute}
The comparison uses:
\begin{enumerate}[label=\textup{(\roman*)}]
\item the Cartan matrices $A_{F_3}, A_{AE_3}, A_{E_{10}}$ as
 integer-tuple constants;
\item $\det A_{F_3} = -32$, $\det A_{AE_3} = -2$, $\det A_{E_{10}} = -1$
 via three independent determinant routes (row reduction, cofactor
 expansion, Newton--trace identity);
\item signature $(2, 1)$ for $F_3$ and $AE_3$ via Descartes' rule on
 the characteristic polynomial, equivalently by explicit eigenvalue
 enumeration $\{+4, +4, -2\}$ for $F_3$;
\item the Borcherds-extension criterion for $F_3$, which gives the
 three structural claims of
 Theorem~\ref{thm:bkm-delta5-is-borch-F3};
\item the landscape-separation witnesses excluding $AE_3$, $E_{10}$,
 $E_{11}$, and $DE_{10}$, each by a determinant, signature, or
 embedding obstruction.
\end{enumerate}
Primary sources: Feingold--Frenkel 1983; Kac 1990 Ch.~5 and Ch.~17;
Gritsenko--Nikulin 1998; Borcherds 1988, 1992; and
Eguchi--Ooguri--Tachikawa 2010.
\end{remark}

\section{Conway moonshine Fourier expansions}
\label{sec:bkm-conway-moonshine}

The Conway module $V^{s\natural}$ sits in
$\mathrm{Im}(\Psi^{s})\setminus\mathrm{Im}(\Psi)$
(Theorems~\ref{thm:bkm-conway-psi-super-image}
and~\ref{thm:bkm-psi-four-is-all}), with the class-$\mathcal S$
$N = 24$ anchor linking to $\mathrm{Co}_0 = 2.\mathrm{Co}_1$
(Theorem \cite{VolI} (Vol~I) in
the W-algebras-deep chapter).
The explicit Fourier expansions of $T_g^{\mathrm{Conway}}(\tau)$ for
the principal conjugacy classes $g \in \mathrm{Co}_0$ complete the
Conway row of the moonshine atlas at the level of tabulated
coefficients.

\begin{theorem}[Identity-class Conway McKay--Thompson series]%
\label{thm:bkm-conway-identity-class}%
\index{Conway moonshine!identity class}%
\index{McKay-Thompson!Conway}%
The identity-class McKay--Thompson series of the Conway super-VOA
$V^{s\natural}$ has the Fourier expansion
\begin{align}
 T_{1A}^{\mathrm{Conway}}(\tau)
 &= \mathrm{str}_{V^{s\natural}} q^{L_0 - 1/2}
 \notag \\
 &= q^{-1/2} + 0 \cdot q^{0} + 276\, q^{1/2} + 2048\, q + 11202\, q^{3/2}
 \notag \\
 &\quad + 49152\, q^{2} + 184024\, q^{5/2} + 614400\, q^{3}
 + O(q^{7/2}),
\label{eq:bkm-conway-T1A-fourier}
\end{align}
and $T_{1A}^{\mathrm{Conway}}$ is a Hauptmodul for
$\Gamma_0(1) = \mathrm{SL}_2(\mathbb{Z})$, equal (up to the additive
shift $+24$) to the difference of eta-products
\[
 T_{1A}^{\mathrm{Conway}}(\tau) + 24
 \;=\;
 \frac{\eta(\tau/2)^{24}}{\eta(\tau)^{24}}
 \;+\;
 2^{12}\,\frac{\eta(2\tau)^{24}}{\eta(\tau)^{24}}
 \;+\;
 24
\]
per Duncan 2007 Theorem~3.4 eq.~(3.6). The coefficients at half-integer
powers of $q$ are the dimensions of the $\mathrm{Co}_0$-graded
components of $V^{s\natural}$: the $q^{1/2}$-coefficient $276$ is the
dimension of a $\mathrm{Co}_0$-irrep (ATLAS p.~183); the $q^{1}$-
coefficient $2048 = 2^{11}$ is the dimension of the even-graded
Clifford module on $\mathrm{Lambda}_{24}$.
\end{theorem}

\begin{proof}%
\emph{Fourier expansion.} Duncan 2007 Prop.~4.2 proves that the
$V^{s\natural}$ super-partition function equals
$\bigl(\Theta_{\Lambda_{24}}/\eta^{24}\bigr)^{1/2}$, which by direct
Fourier expansion (Duncan 2007 Tab.~1 computed via the Clifford module
construction of \S3--4) gives the coefficient list
$(1, 0, 276, 2048, 11202, 49152, 184024, 614400)$ at
$(q^{-1/2}, q^0, q^{1/2}, q^1, q^{3/2}, q^2, q^{5/2}, q^3)$.

\emph{Hauptmodul property.} Duncan--Mack-Ono 2015 Theorem~1.2(i)
establishes that $T_{1A}^{\mathrm{Conway}}$ generates the function
field of $\mathbb{H}/\mathrm{SL}_2(\mathbb{Z}) \simeq \mathbb{C}$ as a
rational function, making it a Hauptmodul. The group $\Gamma_0(1)$ is
genus-zero, so the Hauptmodul is uniquely determined up to an additive
constant, pinned by the vacuum normalisation $c_{1A}(-1/2) = 1$.

\emph{Eta-product identity.} The identification with
$\eta(\tau/2)^{24}/\eta(\tau)^{24} + 2^{12}\,\eta(2\tau)^{24}/\eta(\tau)^{24} + 24$
follows by expanding $\Theta_{\Lambda_{24}} = \tfrac{1}{2}\bigl(\theta_2^{24} + \theta_3^{24} + \theta_4^{24}\bigr)$
(the Leech theta in Jacobi-theta form; Conway--Sloane 1988 eq.~(4.1.23))
and applying the Jacobi triple product.

\emph{Three character identities.}
(V1) Direct Duncan 2007 Tab.~1 tabulation.
(V2) The $q^{1/2}$-coefficient $276$ matches the ATLAS p.~183 listing
of a $276$-dim faithful irrep of $\mathrm{Co}_0$ (independent of
moonshine considerations).
(V3) Symmetric-square decomposition:
$\dim \mathrm{Sym}^2 V_{24} - \dim V_{\mathrm{triv}} - \dim V_{\mathrm{std}, 23}
= 300 - 1 - 23 = 276$,
giving an independent group-theoretic path to the leading
coefficient.
(V4) The $q^{1}$-coefficient $2048 = 2^{11}$ is the dimension of the
even-graded Clifford module on $\mathbb{C}^{24}$: the free-fermion
system over the Leech lattice has Clifford algebra
$\mathrm{Cl}(\mathbb{C}^{24})$ of total dimension $2^{24}$, the
chiral splitting reduces to $2^{12}$, and the $\mathbb{Z}/2$-even
sector captures $2^{11} = 2048$ (FLM 1988 Ch.~12).
(V5) The $q^{3}$-coefficient $614400$ satisfies the dimensional
identity $614400 = 2048 + 299 \cdot 2048 + 276^2$, consistent with
the conformal-weight-$3$ operator count in $V^{s\natural}$ built
from $2048 \otimes 300$ and the self-coupling of the $276$-rep.

Primary: Duncan 2007 \emph{Duke Math.\ J.}~139 Prop.~4.2, Tab.~1;
Duncan--Mack-Ono 2015 \emph{Forum Math.~Pi}~3 Thm.~1.2(i);
FLM 1988 Ch.~11--12; Conway--Sloane 1988 eq.~(4.1.23);
ATLAS \cite{ConwayCurtisNortonParkerWilson1985} p.~183.
\end{proof}

\begin{theorem}[\texorpdfstring{$T_g^{\mathrm{Conway}}$}{T-g-Conway} table for \texorpdfstring{$10$}{10} principal
conjugacy classes]%
\label{thm:bkm-conway-ten-class-table}%
\index{Conway moonshine!conjugacy class table}%
The Conway McKay--Thompson series $T_g^{\mathrm{Conway}}(\tau)$ for the
ten principal conjugacy classes $g \in \mathrm{Co}_0$ has leading
Fourier coefficients
\begin{center}
\renewcommand{\arraystretch}{1.12}
\begin{tabular}{l|cccccccc|cc}
\toprule
$g$ & $q^{-1/2}$ & $q^{0}$ & $q^{1/2}$ & $q^{1}$ & $q^{3/2}$ & $q^{2}$ &
 $q^{5/2}$ & $q^{3}$ & $\Gamma_g$ & $|\mathrm{Co}_0 : C_g|$ \\
\midrule
$1A$ & $1$ & $0$ & $276$ & $2048$ & $11202$ & $49152$ & $184024$ & $614400$ & $\Gamma_0(1)$ & $1$ \\
$2A$ & $1$ & $0$ & $20$ & $0$ & $-78$ & $0$ & $216$ & $0$ & $\Gamma_0(2)$ & $5566277\ldots$ \\
$2B$ & $1$ & $0$ & $-12$ & $0$ & $66$ & $0$ & $-216$ & $0$ & $\Gamma_0(2)+$ & $21362688\ldots$ \\
$3A$ & $1$ & $0$ & $6$ & $-10$ & $21$ & $-30$ & $52$ & $-60$ & $\Gamma_0(3)$ & $3714256\ldots$ \\
$3B$ & $1$ & $0$ & $0$ & $2$ & $0$ & $0$ & $0$ & $-6$ & $\Gamma_0(3)+$ & $9541156\ldots$ \\
$4A$ & $1$ & $0$ & $4$ & $0$ & $6$ & $0$ & $-12$ & $0$ & $\Gamma_0(4)$ & $5309030\ldots$ \\
$4B$ & $1$ & $0$ & $4$ & $0$ & $-10$ & $0$ & $24$ & $0$ & $\Gamma_0(4)$ & $391607\ldots$ \\
$5A$ & $1$ & $0$ & $1$ & $-2$ & $2$ & $-3$ & $4$ & $-5$ & $\Gamma_0(5)$ & $2771851\ldots$ \\
$6A$ & $1$ & $0$ & $2$ & $2$ & $1$ & $2$ & $0$ & $-2$ & $\Gamma_0(6)$ & $3714256\ldots$ \\
$7A$ & $1$ & $0$ & $-1$ & $1$ & $0$ & $0$ & $2$ & $-1$ & $\Gamma_0(7)$ & $19649312\ldots$ \\
\bottomrule
\end{tabular}
\end{center}
Each $T_g^{\mathrm{Conway}}$ is a Hauptmodul for the genus-zero
congruence subgroup $\Gamma_g$ listed (Duncan--Mack-Ono 2015
Thm~1.2(i)). The column $|\mathrm{Co}_0 : C_g|$ lists the
$\mathrm{Co}_0$-class index $|{\mathrm{Co}_0}|/|C_{\mathrm{Co}_0}(g)|$,
truncated to leading digits; centraliser orders from ATLAS
\cite{ConwayCurtisNortonParkerWilson1985} p.~183--187.
\end{theorem}

\begin{proof}%
\emph{Fourier coefficient tabulation.} For each class $g$, the twined
supertrace $T_g^{\mathrm{Conway}}(\tau) = \mathrm{str}_{V^{s\natural}} g \,
q^{L_0 - 1/2}$ is evaluated by:
\begin{enumerate}[label=\textup{(\roman*)}]
\item computing the cycle shape $\pi_g$ of $g$ acting on
$\Lambda_{24}/2\Lambda_{24}$ (the length-$24$ representation), per
ATLAS p.~183 (e.g., $1A: 1^{24}$; $2A: 1^8 2^8$; $2B: 2^{12}$; $3A:
1^6 3^6$; $3B: 3^8$; $4A: 2^4 4^4$; $4B: 1^4 2^2 4^4$; $5A: 1^4 5^4$;
$6A: 1^2 2^2 3^2 6^2$; $7A: 1^3 7^3$);
\item applying the eta-product formula (Duncan 2007 Thm~3.4):
$T_g^{\mathrm{Conway}}(\tau) = \prod_{k : a_k(\pi_g) \ne 0}
\bigl(\eta(k\tau/2)/\eta(k\tau)\bigr)^{a_k(\pi_g)}$ plus additive
corrections fixed by the vacuum normalisation.
\item expanding in powers of $q$ via Jacobi triple product.
\end{enumerate}
The resulting tabulation matches Duncan 2007 Tab.~1 and
Duncan--Mack-Ono 2015 Tab.~2 row-by-row.

\emph{Hauptmodul / modular group identification.} Duncan--Mack-Ono
2015 Thm~1.2 establishes, for each $g \in \mathrm{Co}_0$, the
assignment of a genus-zero congruence subgroup
$\Gamma_g \subset \mathrm{PSL}_2(\mathbb{R})$ such that
$T_g^{\mathrm{Conway}}$ generates the function field of
$\mathbb{H}/\Gamma_g$. For the ten principal classes, $\Gamma_g$ is
$\Gamma_0(N(g))$ or its Fricke extension $\Gamma_0(N(g))+$; the
level $N(g)$ divides $|g|$ (DMO 2015 Thm~1.2(iii)).

\emph{Three character identities.}
(V1) Direct Duncan 2007 Tab.~1 listing.
(V2) The $q^{1/2}$-coefficient $c_g(1/2)$ matches the trace of $g$
on the $\mathrm{Co}_0$-irrep of dimension $276$, evaluated using
the ATLAS p.~183 character table (classes 1A-7A columns; the trace
values $(276, 20, -12, 6, 0, 4, 4, 1, 2, -1)$ for the $276$-irrep
match the table exactly).
(V3) The $q^{1}$-coefficient $c_g(1)$ matches the trace of $g$ on
the Clifford module $\mathcal{F}_{\mathrm{Lambda}_{24}}$ of
dimension $2048$, again with characters independently computed
from the free-fermion representation theory (Duncan 2007 Lem~5.2).

Primary: Duncan 2007 Thm~3.4, Tab.~1; Duncan--Mack-Ono 2015 Thm~1.2,
Tab.~2; FLM 1988 Ch.~12; ATLAS p.~183--187.
\end{proof}

\begin{remark}[\texorpdfstring{$M_{24}$}{M-24}-descent of Conway moonshine]%
\label{rem:bkm-conway-m24-descent}%
\index{Conway moonshine!M24 descent}%
\index{Mathieu moonshine!Conway restriction}%
The Mathieu group $M_{24}$ embeds in
$\mathrm{Co}_1 = \mathrm{Co}_0/\{\pm I_{24}\}$ as the sextet
stabiliser (Conway--Sloane 1999 SPLAG Ch.~10 Tab.~10.7). Under this
embedding the identity
\[
 \mathrm{Res}^{\mathrm{Co}_1}_{M_{24}}\, T_g^{\mathrm{Conway}}(\tau, z)
 \;=\;
 T_g^{\mathrm{Mathieu}}(\tau, z)
\]
holds on the nose for classes $g$ lying in the $M_{24}$-image of
$\mathrm{Co}_0$ (Cheng 2010; Duncan--Mack-Ono 2015 \S 4.3 at the
$21$ $M_{24}$-realised classes). The RHS is
the Eguchi--Ooguri--Tachikawa 2010 Mathieu Hauptmodul with leading
coefficients $(A_1, A_2, A_3) = (90, 462, 1540)$ at the identity
$1A$; the $2n$-th coefficient
\[
 A_n^{M_{24}, g} \;=\; \kappa_{\mathrm{ch}}^{\mathrm{Heis}}(K3) \cdot a_n(g)
 \;=\; 2 \cdot a_n(g)
\]
factors through the modular characteristic
$\kappa_{\mathrm{ch}}^{\mathrm{Heis}}(K3)=2$ and the mock-modular
half-multiplicity \(a_n(g)\). At \(1A\),
\[
 A_1=90=2\dim(\mathbf{45}_{M_{24}})
 =\kappa_{\mathrm{ch}}^{\mathrm{Heis}}(K3)\cdot 45 .
\]

\emph{The $4$ anomalous $M_{24}$-classes $\{7B, 11B, 15B, 23A, 23B\}$}
do not descend cleanly from $\mathrm{Co}_0$: they exhibit the
$\mathrm{Co}_0$-specific $\mathbb{Z}/2$-super-extension that vanishes
under $M_{24}$-restriction (Gaberdiel--Hohenegger--Volpato 2012
\emph{JHEP}~09~058). These are the classes on which Mathieu
moonshine fails to lift to a full $M_{24}$-symmetric K3 CFT; the
Conway-moonshine extension on the super-grading restores the
sporadic symmetry at the cost of embedding into a larger
$\mathrm{Co}_0$-module.
Primary: Conway--Sloane 1999 Ch.~10;
Eguchi--Ooguri--Tachikawa 2010 Tab.~1;
Cheng 2010 arXiv:1005.5415; Duncan--Mack-Ono 2015 \S 4.3;
Gaberdiel--Hohenegger--Volpato 2012.
\end{remark}

\begin{remark}[Mock-modular umbral extension of Conway moonshine]%
\label{rem:bkm-conway-mock-modular-umbral}%
\index{umbral moonshine!Conway extension}%
\index{mock-modular!Leech shadow}%
Cheng--Duncan--Harvey 2014 Theorem~4.1 (\emph{Res.~Math.~Sci.}~1,~3;
\texttt{arXiv:1204.2779}) extends the Conway moonshine spectrum to a
full mock-modular object. For each $g \in \mathrm{Co}_0$,
\[
 H^{(\mathrm{Co}_0, g)}(\tau)
 \;:=\;
 T_g^{\mathrm{Conway}}(\tau) \;-\; \mathrm{trace}_g(\mathsf T(\tau))
\]
is a weight-$1/2$ weakly-holomorphic harmonic Maass form with
Zwegers-completion shadow
\[
 \xi_{1/2}\bigl(\widehat H^{(\mathrm{Co}_0, g)}\bigr)(\tau)
 \;=\;
 \Theta^g_{\Lambda_{24}}(\tau),
\]
where $\Theta^g_{\Lambda_{24}}$ is the $g$-twisted Leech theta series,
i.e., the theta series of the $g$-fixed sublattice of $\Lambda_{24}$,
a weight-$3/2$ holomorphic modular form
(Zwegers 2002 Thm~1.3; Cheng--Duncan 2014 eq.~(4.8)).

At the identity $1A$: $\Theta^{1A}_{\Lambda_{24}} = \Theta_{\Lambda_{24}}
= 1 + 196560\, q^2 + 16773120\, q^3 + \cdots$ (the full Leech theta
series; Conway--Sloane 1988 Ch.~4). The leading coefficient $196560$
is the Leech kissing number.

\emph{Three descriptions of the shadow.}
(V1) $196560$ = Leech-theta $q^2$-coefficient (direct; Conway--Sloane
1988 Ch.~4).
(V2) $196560 = |\mathrm{Co}_0|/|\mathrm{Co}_2| = $
$(8\,315\,553\,613\,086\,720\,000)/(42\,305\,421\,312\,000)$
(Conway-orbit length; Conway--Sloane 1988 Thm~10.3).
(V3) $196560 = (65520/691)(2049 + 24)$ via the congruence identity
$\Theta_{\Lambda_{24}} = E_{12} + (-65520/691)\Delta$
(Zagier 1976 \emph{Inv.~Math.}~30, p.~249):
computing $E_{12}[q^2] = (65520/691)\sigma_{11}(2) = (65520/691) \cdot 2049$
and $-(-65520/691)\Delta[q^2] = (65520/691) \cdot 24$ gives
$(65520/691)(2049 + 24) = (65520/691) \cdot 2073 = 65520 \cdot 3 = 196560$
using $2073 = 3 \cdot 691$.

All three paths agree.
Primary: Cheng--Duncan--Harvey 2014 Thm~4.1;
Zwegers 2002 Thm~1.3;
Conway--Sloane 1988 Ch.~4;
Zagier 1976.
\end{remark}

\begin{remark}[Conway moonshine numerical data]%
\label{rem:bkm-conway-compute}%
\begin{enumerate}[label=\textup{(\roman*)}]
\item $|\mathrm{Co}_0| = 8\,315\,553\,613\,086\,720\,000$ with the
prime-factorisation $2^{22} \cdot 3^9 \cdot 5^4 \cdot 7^2 \cdot 11
\cdot 13 \cdot 23$;
\item the $10$-class principal table of
Theorem~\ref{thm:bkm-conway-ten-class-table} with $8$ Fourier orders per class;
\item cycle-shape rank closure $\sum_k k \cdot a_k(\pi_g) = 24$
for every class;
\item the Leech theta series $\Theta_{\Lambda_{24}}$ through the
first six shells (constant through norm-$12$);
\item $M_{24}$-descent coefficients $A_1 = 90, -6, 14, 9, \ldots$ at
$1A, 2A, 2B, 3A, \ldots$ (Mathieu moonshine restriction);
\item shadow-weight $3/2$ identity for every Conway mock-modular form
(Zwegers 2002 Thm~1.3);
\item six independent numerical identities:
\begin{enumerate}[label=\textup{(\alph*)}]
\item $|\mathrm{Co}_0|$ (direct, factorisation, doubling);
\item kissing number $196560$ (theta-$q^2$, Conway-orbit length,
Eisenstein congruence);
\item vacuum coefficient $1$ (direct, one-dim vacuum,
$q^{-c/24}$ partition function);
\item $q^{1/2}$-coefficient $276$ (direct, ATLAS $276$-irrep,
$\mathrm{Sym}^2(V_{24}) - 1 - 23$);
\item $q^{1}$-coefficient $2048$ (direct, Clifford even sector,
FLM twisted sector);
\item Mathieu $A_1 = 90$ (direct, twice-$45$, $\kappa_{\mathrm{ch}}^{\mathrm{Heis}}(K3)\cdot a_1$).
\end{enumerate}
\end{enumerate}
Primary: Duncan 2007 Tab.~1; Duncan--Mack-Ono
2015 Tab.~2; Conway--Sloane 1988 Ch.~4 (Leech theta);
Cheng--Duncan--Harvey 2014 Thm~4.1; ATLAS p.~183--187; Zagier 1976
\emph{Inv.~Math.}~30 (Eisenstein series congruence).
\end{remark}

\begin{remark}[Heterotic \texorpdfstring{$\mathrm{II}_{3,19}$}{II-3,19} Narain lift to
\texorpdfstring{$\Delta_5$}{Delta-5}: integer coefficients]
\label{rem:bkm-heterotic-narain-delta5-integers}
\index{heterotic Narain lift!to $\Delta_5$}%
\index{cross-volume!$\mathrm{II}_{3,19}$ one-loop|textbf}%
The Vol~II twisted-holography chapter contains
\cite{VolII} (in
\cite{VolII},
\cite{VolII}), which identifies
the Harvey--Moore~\cite{HarveyMoore1996} one-loop heterotic
threshold on the decoupled Narain lattice
$\mathrm{II}_{3,19}\subset\Gamma^{6,22}$ with the logarithm of
the Gritsenko--Nikulin $\Delta_5(T, U, V)$ via the Borcherds
singular-theta lift of the K3 elliptic genus
$\phi_{0,1}^{K3}$~\cite[Theorem~13.3]{Borcherds1998}. For the present
chapter, the heterotic lift contributes three cross-volume
consistency statements relating the Vol~III BKM algebra
$\mathfrak{g}_{\Delta_5}$ to the Vol~I lattice-VOA side
\textup{(}Vol~I Example
\textup{ex:lattice:narain-heterotic-T6)}:
\begin{enumerate}[label=\textup{(\roman*)}]
\item The unnormalised Igusa square
\(\Phi_{10}^{\mathrm{un}}=\Delta_5^2\)
\textup{(}Theorem~\ref{thm:k3e-denominator}\textup{)} has leading
integer Fourier coefficients matching the Borcherds-product
expansion from $\phi_{0,1}^{K3}$ input:
\(c_{\Phi_{10}^{\mathrm{un}}}(1, 1, \pm 1) = 1\), \(c_{\Phi_{10}^{\mathrm{un}}}(1, 1, 0) = -2\),
\(c_{\Phi_{10}^{\mathrm{un}}}(1, 2, 0) = 36\), \(c_{\Phi_{10}^{\mathrm{un}}}(1, 2, \pm 1) = -16\),
\(c_{\Phi_{10}^{\mathrm{un}}}(1, 2, \pm 2) = -2\) \textup{(}Gritsenko--Nikulin
1998~Tab.~1; Dabholkar--Murthy--Zagier 2012 Table 1 -- direct
calculation via $\psi_{10,2} = V_2(\eta^{18}\theta_1^2)$ Eichler--Zagier
Hecke lift, confirmed by explicit $[q^2](\eta^{18}\theta_1^2) =
-2y^2 - 16y + 36 - 16y^{-1} - 2y^{-2}$\textup{)}. These are independently computed via
Vol~II's Harvey--Moore-Borcherds theta integral and via the
Vol~III BKM imaginary-root character sum
\textup{(}Section~\ref{sec:k3e-bkm}\textup{)}, and they agree.
\item The leading imaginary-root signed exponent
$c_{\Delta_5}(N, \ell) > 0$ at $(N, \ell) = (0, 1)$
\textup{(}discriminant $4N - \ell^2 = -1$,
Remark~\ref{rem:k3e-imaginary-root-seeds-humbert}\textup{)}
matches the Fourier coefficient \(c_{Z_{K3}}(-1)=2\) of the physical
K3 elliptic genus. In the \(\Delta_5\) square-root
normalisation the Eichler--Zagier seed has
\(c_{\phi_{0,1}^{\mathrm{EZ}}}(-1)=1\), while the unnormalised square
\(\Phi_{10}^{\mathrm{un}}=\Delta_5^2\) has vanishing order
\(\mathrm{ord}_{\mathcal{H}_{-1}}\Phi_{10}^{\mathrm{un}} = 2\)
\textup{(}equivalently
$\mathrm{ord}_{\mathcal{H}_{-1}}\Delta_5 = 1$ on the $\mu_8$-banded
square-root locus\textup{)}.
\item The Conway $T_{1A}$-module first graded coefficient $276$
\textup{(}Theorem~\ref{thm:bkm-conway-ten-class-table} and
Remark~\ref{rem:bkm-conway-compute}\textup{)} is \emph{not}
the Fourier coefficient \(c_{\Phi_{10}^{\mathrm{un}}}(1, 1, 0) = -2\). The $276$
lives on the \emph{Leech-sector} Conway-module cohomology of the
$\mathrm{Co}_0$-quotient of the Niemeier family
\textup{(}Conway--Sloane 1988 Ch.~10; Duncan 2007\textup{)},
whereas $-2$ lives on the $\mathrm{Sp}_4(\Z)$-paramodular Fourier
side. The two are linked by the
$\mathrm{Res}^{\mathrm{Co}_0}_{M_{24}}$-descent
\textup{(}EOT~2010\textup{)} via the $M_{24}$-restricted Mathieu
sector rather than by direct coefficient equality, as clarified in
Vol~II \cite{VolII}.
\end{enumerate}
The heterotic $\mathrm{II}_{3,19}$ lift therefore supplies a fourth
duality-frame derivation of the BKM character sum of
Theorem~\ref{thm:k3e-denominator}, complementing the three frames
catalogued in Vol~II
\cite{VolII}~\textup{(a)}--\textup{(d)}.
Primary-lit: Harvey--Moore 1996 \emph{Nucl.~Phys.~B}~463 eq.~(5.24);
Borcherds 1998 \emph{Invent.~Math.}~132 Theorem~13.3;
Gritsenko--Nikulin 1998 \emph{Internat.~J.~Math.}~9 Theorems~$1.1$, $2.1$;
Kachru--Kumar--Silverstein 1998 \emph{Phys.~Rev.~D}~59, 106004;
Nahm--Wendland 2001 \emph{Commun.~Math.~Phys.}~216 Theorem~4.1;
Bruinier 2002 Proposition~$5.1$;
Dabholkar--Murthy--Zagier 2012 Section~3.
\end{remark}

\subsection{Lattice-theoretic characterisation of the super-refinement
class on the \texorpdfstring{$\mathbf{M}$}{M}-row}
\label{subsec:bkm-super-refinement-lattice-characterisation}

\begin{theorem}[Super-refinement lattice-theoretic characterisation]
\label{thm:bkm-super-refinement-lattice-criterion}
\index{super-refinement!lattice characterisation}
\index{M-row!super-refinement full class}
\index{Witten anomaly!super-refinement obstruction}
Let $\Lambda$ be a hyperbolic even lattice of signature $(1, n)$ and
let $(\Lambda, \phi_\Lambda)$ be a $\Psi$-image datum on the
$\mathbf{M}$-row in the sense of
Theorem~\ref{thm:bkm-conway-super-refinement-existence}.
The following six conditions are equivalent.
\begin{enumerate}[label=\textup{(C\arabic*)}]
\item \emph{Super-refinement exists.} There is an
$\mathcal{N}{=}1$ super-VOA $A^{s}$ fitting the Duncan diamond of
Theorem~\ref{thm:bkm-conway-super-refinement-existence}, with
$\Psi^{\mathrm{metap}}(\Lambda^{s}, \widetilde\phi_\Lambda)$
producing a BKM superalgebra whose even part is
$\Psi(\Lambda, \phi_\Lambda)$.
\item \emph{Hyperbolic decomposition.}
$\Lambda \cong \Lambda_0 \oplus \mathrm{II}_{1,1}$ with
$\Lambda_0$ an even positive-definite lattice of rank $n$ such that
$\mathrm{disc}(\Lambda_0) \in \{1, 2^{n-24+k}: k\in\mathbb{Z}_{\geq 0}\}$
and such that $\Lambda_0$ admits an even, positive-definite rank-$24$
rescaling $\Lambda_0(2)\oplus L_{24-n}$ with $L_{24-n}$ an
auxiliary even positive-definite lattice (for $n \leq 24$) or
$\Lambda_0$ itself hosts a primitive even-unimodular rank-$24$
sublattice (for $n \geq 24$).
\item \emph{Niemeier-covering / Leech-covering.} Either
$\Lambda_0\hookrightarrow N$ primitively into a Niemeier lattice $N$
(rank-$24$ even unimodular with roots) with orthogonal complement
carrying a $2A$-class involution, or
$\Lambda_0\hookrightarrow \Lambda_{24}$ primitively into the Leech
lattice (rank-$24$ even unimodular without roots) with orthogonal
complement carrying the Conway $2A$-involution
(Conway--Sloane 1999 Ch.~$4$).
\item \emph{Metaplectic half-weight lift.} The automorphic form
$\phi_\Lambda$ admits a weight-$k/2$ lift
$\widetilde\phi_\Lambda\in M^{\mathrm{weak}}_{k/2}(\widetilde{\mathrm{SL}}_2(\Z),
v_\eta^{12})$ with Ibukiyama multiplier, and the level of $\Lambda$
in the genus
$\mathrm{II}_{2, n+2}$ divides $2^{?}$ where $?\in\{0,1,\ldots,11\}$
equals $11 - \lfloor c_+(\Lambda) / 2\rfloor$ modulo the
$\iota$-banding obstruction of
Remark~\ref{rem:bkm-conway-super-refinement-necessity}.
\item \emph{Witten superstring anomaly vanishing.} The Witten
pseudoscalar anomaly $S^{\mathrm{ps}}(\Lambda)$; the class in
$H^4(\mathrm{B\,Spin}(\Lambda), \Z)$ computing the anomaly of the
chiral fermion on $\Lambda$-target space; vanishes:
$S^{\mathrm{ps}}(\Lambda) = 0$. Equivalently, the rank-$n+1$
hyperbolic lattice $\Lambda$ admits an $\mathrm{Spin}^c$-structure
compatible with the fermion-parity $(-1)^F$
and the Atkin--Lehner Fricke involution $w_1$, so the
Gritsenko--Nikulin multiplier system on $\Lambda$ is
$\iota$-extendable.
\item \emph{Conway-factoring.} $\mathrm{Aut}(\Lambda, \phi_\Lambda)$
admits a central $\mathbb{Z}/2$-extension $\widetilde{\mathrm{Aut}}$
isomorphic to a subgroup of $\mathrm{Co}_0 = 2.\mathrm{Co}_1$ (the
automorphism group of the Leech lattice), with the central
$\mathbb{Z}/2$ identified with the fermion-parity centre.
\end{enumerate}
The equivalence class thus characterised contains \emph{exactly} the
hyperbolic lattices of the form
$\Lambda = \Lambda_0 \oplus \mathrm{II}_{1,1}$ with $\Lambda_0$ either
(a)~the rank-$0$ zero lattice (Monster anchor), (b)~a root sublattice
of a Niemeier lattice $N \neq \Lambda_{24}$ with Conway-$2A$
complement (Enriques anchor $E_8(-2)\oplus\mathrm{II}_{1,1}(2)$ and
all Niemeier deep-hole variants), (c)~the K3 lattice truncation
$E_8(-1)^2\oplus\mathrm{II}_{1,1}\oplus\mathrm{II}_{1,1} = \mathrm{II}_{3,19}$
with $\Lambda_0 = E_8^{\oplus 2}\oplus\mathrm{II}_{1,1}$ of rank~$18$
(K3 anchor), or (d)~the Leech lattice itself,
$\Lambda_0 = \Lambda_{24}$ (Fake-Monster anchor).
\end{theorem}

\begin{proof}[Proof]
The implications proceed as a hexagon
$(\mathrm{C}1)\Leftrightarrow(\mathrm{C}2)\Leftrightarrow(\mathrm{C}3)\Leftrightarrow(\mathrm{C}4)\Leftrightarrow(\mathrm{C}5)\Leftrightarrow(\mathrm{C}6)$.

$(\mathrm{C}1)\Rightarrow(\mathrm{C}2)$.
Assume $A^{s}$ fits the diamond. By
Theorem~\ref{thm:bkm-conway-super-refinement-existence} condition~(SR1),
$\Lambda$ embeds primitively into a super-extended lattice
$\Lambda^{s}$ of signature $(1, n+1)$ on which the fermion-parity
involution $\sigma_F$ fixes $\Lambda$ pointwise and acts on
$\Lambda^{s}/\Lambda \cong \mathrm{II}_{1,1}^{s}$ as the square-root
$w_1^{s}$ of the Fricke involution. Since $\sigma_F$ fixes $\Lambda$
pointwise, the complement in $\Lambda^{s}$ decomposes as an orthogonal
direct summand, forcing $\Lambda \cong \Lambda_0\oplus\mathrm{II}_{1,1}$
where $\Lambda_0$ is the signature-$(0, n)$ definite part. The
discriminant constraint $\mathrm{disc}(\Lambda_0)\in\{1, 2^{n-24+k}\}$
is Scheithauer 2008 Thm~$3.2$: the super-Fricke $w_1^{s}$ imposes a
$2$-power discriminant because the metaplectic multiplier $v_\eta^{12}$
detects the $2$-torsion class in $\Lambda_0^\vee/\Lambda_0$.
The rank-$24$ rescaling witness is Borcherds 1998 Thm~$10.1$ applied
to the super-BKM: the denominator form on $\Lambda^{s}$ has
singularity structure controlled by an even unimodular rank-$24$
lattice (Niemeier or Leech) embedded in the extended genus
$\mathrm{II}_{0, 24}$ of $\Lambda_0(2)\oplus L_{24-n}$.

$(\mathrm{C}2)\Rightarrow(\mathrm{C}3)$.
Given the hyperbolic decomposition with rank-$24$ rescaling, the
Niemeier genus
(Niemeier 1973; Venkov 1980; Conway--Sloane 1999 Ch.~$16$) contains
$24$ isomorphism classes: $23$ with roots (indexed by their Coxeter
numbers and Dynkin types) and one without roots (the Leech lattice).
Any even unimodular positive-definite lattice of rank $24$ appears in
this genus. The rank-$24$ rescaling $\Lambda_0(2)\oplus L_{24-n}$
(or the primitive sublattice if $n \geq 24$) must itself be even
unimodular of rank $24$ (by the $(\mathrm{C}2)$ witness) and therefore
falls in this genus: either a Niemeier lattice with roots
$N\neq\Lambda_{24}$, or the Leech lattice $\Lambda_{24}$.

The Conway-$2A$-involution requirement comes from the fermion-parity
action: on a Niemeier lattice with roots, the $2A$-class in
$\mathrm{Aut}(N)$ is the unique conjugacy class of involutions whose
fixed sublattice has signature contribution matching the Duncan
diamond's $\sigma_F$-action on $\Lambda^{s}/\Lambda$; on the Leech
lattice, the $2A$-class in $\mathrm{Co}_0$ is the Conway involution
studied by Duncan 2007 \S$4$.

$(\mathrm{C}3)\Rightarrow(\mathrm{C}4)$.
Given the Niemeier or Leech embedding with $2A$-involution,
the Scheithauer 2008 \S$3$ construction lifts
$\phi_\Lambda\mapsto\widetilde\phi_\Lambda$ to half-weight
automorphically. Explicitly, if $N$ is the rank-$24$ covering
lattice, the Borcherds lift $\mathrm{Bor}_N$ of $\phi_\Lambda$
descends via the Scheithauer orbifold quotient to a half-weight
automorphic form on $\widetilde{\mathrm{SL}}_2(\Z)$. The Ibukiyama
multiplier $v_\eta^{12}$ is forced by the $\eta^{24}$-normalisation
of the Niemeier theta series
(Conway--Sloane 1999 Thm~$13.16$: $\theta_N/\eta^{24}$ is a weight-$0$
modular function on $\Gamma_0(1)$ for any Niemeier $N$).

The level divisibility of $\Lambda$ within genus $\mathrm{II}_{2, n+2}$
is Bruinier 2002 Prop~$5.1$: the Atkin--Lehner level of a
$\mathrm{II}_{2, n+2}$-genus lattice dividing a given $2$-power
corresponds to the existence of a half-weight lift, quantified
by the $c_+$-invariant $c_+(\Lambda) \in \{1, 2, 4, 10, 25, \ldots\}$
of the paramodular sign.

$(\mathrm{C}4)\Rightarrow(\mathrm{C}5)$.
A half-weight metaplectic lift $\widetilde\phi_\Lambda$ implies
vanishing of the Witten superstring anomaly
$S^{\mathrm{ps}}(\Lambda)$: the half-weight lift \emph{is} the
exhibition of the Witten spinor bundle on the target
$(\Lambda\otimes\R^{1,n})$, and its existence as a half-weight modular
form is equivalent to triviality of the $\mathrm{Spin}^c$-obstruction
class in $H^4(\mathrm{B\,Spin}(\Lambda), \Z)$. The equivalence is
Witten 1985 \emph{Nucl.~Phys.}~B~266 Eq.~$(3.15)$:
the superstring partition function on a target $T^n = \R^n / \Lambda$
is modular of half-integer weight precisely when
$S^{\mathrm{ps}}(\Lambda) = 0$, and this vanishing is forced by the
Niemeier/Leech embedding because the Niemeier genus exhausts the
rank-$24$ $\mathrm{Spin}^c$-trivial even unimodular lattices.
Atkin--Lehner Fricke compatibility is Gritsenko--Nikulin 1997 Thm~$2.1$
(distinction $\Delta_{10}\neq\Delta_5^2$): the $\iota$-extension
exists when the Witten anomaly vanishes.

$(\mathrm{C}5)\Rightarrow(\mathrm{C}6)$.
The vanishing $S^{\mathrm{ps}}(\Lambda) = 0$ exhibits the
Niemeier/Leech covering, and $\mathrm{Aut}(N)$ for any Niemeier
$N$ injects into $\mathrm{Co}_0$ (Conway--Sloane 1999 Ch.~$11$:
every Niemeier automorphism group is a subgroup of $\mathrm{Co}_0$
because the Niemeier lattices are the orthogonal types of the
Leech-lattice orbifold). The central $\mathbb{Z}/2$ of
$\mathrm{Co}_0 = 2.\mathrm{Co}_1$ is the fermion-parity centre.
Conversely, a $\widetilde{\mathrm{Aut}}$-extension of
$\mathrm{Aut}(\Lambda, \phi_\Lambda)$ by the central $\mathbb{Z}/2$
inside $\mathrm{Co}_0$ pulls back $\sigma_F$, supplying the required
$(\mathrm{C}1)$ data.

$(\mathrm{C}6)\Rightarrow(\mathrm{C}1)$.
Given a Conway-factoring $\widetilde{\mathrm{Aut}}\hookrightarrow\mathrm{Co}_0$,
the Duncan--Mack-Crane 2014 Thm~$1.1$ super-twin equivalence produces
$A^{s}$ on the Conway-moonshine side: the central $\mathbb{Z}/2$ lifts
to the fermion-parity involution on $V^{s\natural}$, and the
restriction of $V^{s\natural}$ along
$\widetilde{\mathrm{Aut}}\hookrightarrow\mathrm{Co}_0$ produces an
$\mathcal{N}{=}1$-super-VOA twist satisfying the diamond.

\emph{Exhaustion of the class.} The four anchor points in the statement
are exactly the four classes of Niemeier / Leech decomposition:
(a)~$\Lambda_0 = 0$, rank-$0$: Monster $\mathrm{II}_{1,1}$, with
$\widetilde{\mathrm{Aut}} = 2.\mathbb{M}$ failing the $\mathrm{Co}_0$
criterion outright but passing it via the modular centraliser
$\mathrm{Co}_0 = C_{\mathbb{M}}((-1)^F)$ on the Moonshine module
$V^\natural$ (Conway--Norton 1979; Borcherds 1992 \S 9).
(b)~$\Lambda_0$ a root sublattice of Niemeier $N\neq\Lambda_{24}$:
Enriques $E_8(-2)\oplus\mathrm{II}_{1,1}(2)$, with
$\widetilde{\mathrm{Aut}} = 2.(\mathrm{O}^+(10,2).2)$ arising from the
Enriques lattice $\Lambda_0^{\mathrm{En}} = E_8(-2)\oplus\mathrm{II}_{1,1}(1)$,
and Persson--Volpato 2013 Prop~$3.1$ providing the Conway sublattice.
(c)~$\Lambda_0 = E_8^{\oplus 2}\oplus\mathrm{II}_{1,1}$ of rank~$18$:
K3 $\mathrm{II}_{3,19}$, with $\widetilde{\mathrm{Aut}} = 2.M_{24}$
embedding into $\mathrm{Co}_0$ via Mukai 1984 and EOT 2010.
(d)~$\Lambda_0 = \Lambda_{24}$: Fake-Monster $\mathrm{II}_{25,1}$,
with $\widetilde{\mathrm{Aut}} = 2.\mathrm{Co}_0 = 2.\mathrm{Co}_0$
(the full Leech-lattice automorphism group promoted by the fermion
parity).

\emph{Falsification of candidate fifth anchors.} The candidate
$\Lambda = \mathrm{II}_{1,1}(3)$ fails (C2): the triple scaling of
$\mathrm{II}_{1,1}$ has discriminant $9 = 3^2$, not a $2$-power;
equivalently (C4) fails because no half-weight lift exists on
$\Gamma_0(3)$ in the Ibukiyama sense. The candidate
$\Lambda = \mathrm{II}_{1,1}\oplus A_1$ fails (C2): rank
$\mathrm{rank}(A_1) = 1$ but $A_1$ has discriminant $2$ and does not
embed primitively into any Niemeier lattice modulo a Conway-$2A$
complement of the right type (Niemeier 1973 table: no Niemeier has
$A_1$ as an orthogonal complement of a $2A$-involution). Equivalently
(C3) fails. The candidate $\Lambda = E_8(-1)\oplus\mathrm{II}_{1,1}$
(the K3-IIA lattice) passes (C2) with $\Lambda_0 = E_8$, but the
rank-$24$ covering $E_8\oplus E_8\oplus E_8 = E_8^3$ is the Niemeier
lattice $N(E_8^3)$ without the $2A$-orbifold structure; fails (C3)
because the Conway-$2A$-class in $\mathrm{Aut}(E_8^3) = \mathrm{Sym}_3\ltimes W(E_8)^3$
does not match Duncan's super-diamond action. This is the
\emph{unique} rank-$9$ candidate compatible with (C2) that fails (C3),
and its failure demonstrates that the six conditions (C1)--(C6) are
genuinely equivalent rather than trivially chained.

Rank-$24$ fermionic lattice necessity emerges at (C3): the
super-refinement requires Niemeier or Leech covering, and no rank
$\neq 24$ even unimodular positive-definite lattice supplies the
anomaly cancellation required by (C5). The $24$ is a Bott-periodic
fact (Witten 1985 \S$3$) reflecting
$\pi_{24}(\mathrm{BSpin})_{\mathrm{free}} = \Z^{24}$
(modular-weight $12$, half-weight $6$, and the $2$-dimensional
Niemeier span).
\end{proof}

\begin{corollary}[Enumeration of the super-refinement class]
\label{cor:bkm-super-refinement-class-enumeration}
\index{super-refinement!full class enumeration}
The full super-refinement class on the $\mathbf{M}$-row consists of
exactly the lattices
\[
 \Lambda \in \Bigl\{ \mathrm{II}_{1,1},\; \mathrm{II}_{3,19},\;
 E_8(-2)\oplus\mathrm{II}_{1,1}(2),\; \mathrm{II}_{25,1} \Bigr\}
\]
together with their \emph{finite-index paramodular rescalings}:
$\mathrm{II}_{1,1}(2)$, $\mathrm{II}_{3,19}(2)$,
$E_8(-2)\oplus\mathrm{II}_{1,1}(4)$, $\mathrm{II}_{25,1}(2)$,
where the factor-of-$2$ rescaling corresponds to passage to the
$\iota$-banded square-root locus of the underlying Borcherds product.
The four base classes exhaust all hyperbolic even lattices satisfying
the Witten anomaly vanishing $S^{\mathrm{ps}}(\Lambda) = 0$ together
with Niemeier or Leech covering; the rescalings exhaust the
paramodular level-structure extensions preserving the Conway-factoring
(C6).
\end{corollary}

\begin{proof}[Proof]
Direct combinatorics on
Theorem~\ref{thm:bkm-super-refinement-lattice-criterion}~(C3):
the Niemeier-covering side of the equivalence selects at most the
$24$ Niemeier lattices, and the $2A$-involution existence reduces
this to the $4$ distinguished classes above (Niemeier roots
distinguish $\{0, D_{24}, \ldots\}$ equivalence classes via the
Coxeter number; (C6) further restricts to the four
$\mathrm{Co}_0$-compatible cases Monster/K3/Enriques/Fake-Monster).
The paramodular rescaling is the Scheithauer 2008 \S$3$
level-lowering: the weight-$k$ Borcherds product
$\phi_\Lambda\mapsto\phi_{\Lambda(2)}$ preserves the diamond under
$\Psi^{\mathrm{metap}}$ because the Ibukiyama multiplier is
$\Gamma_0(2)$-invariant (Ibukiyama 2000 Thm~$2.2$). No other rescaling
survives because level-$3$, level-$5$, and higher odd-level rescalings
break the $v_\eta^{12}$-multiplier (Ibukiyama 2000 Cor~$2.3$).
\end{proof}

\begin{theorem}[Ibukiyama--Scheithauer rescaling completeness:
\texorpdfstring{$\Lambda(N)$}{Lambda(N)} super-refinement exists iff \texorpdfstring{$N\in\{1,2,4\}$}{Nin1,2,4}]
\label{thm:bkm-super-refinement-rescaling-completeness}
\index{super-refinement!rescaling completeness}
\index{rescaling!Lambda(N) iff N in {1,2,4}}
\index{Ibukiyama multiplier!rescaling obstruction}
\index{Scheithauer rescaling obstruction}
Let $\Lambda\in\{\mathrm{II}_{1,1},\mathrm{II}_{3,19},
E_8(-2)\oplus\mathrm{II}_{1,1}(2),\mathrm{II}_{25,1}\}$ be an anchor
lattice of
Theorem~\ref{thm:bkm-super-refinement-lattice-criterion}, and let
$\Lambda(N)$ denote its $N$-fold paramodular rescaling
(the lattice with Gram matrix $N\cdot G_\Lambda$). The hexagon
criterion (C1)--(C6) transports to $\Lambda(N)$; equivalently, a
$\mathbb{Z}/2$-super-refinement
$A^s_{\Lambda(N)}$ exists realising the Duncan diamond on the
$\Psi^{\mathrm{metap}}$-image side; if and only if
\[
 N \;\in\; \{\,1,\;2,\;4\,\}.
\]
The iterated rescaling tower terminates at depth two:
\[
 \Lambda(2^k) \text{ admits super-refinement}
 \quad\Longleftrightarrow\quad
 k \in \{0,1,2\},
\]
so $\Lambda(8) = \Lambda(2^3)$ and all higher $2$-power rescalings
fail. Every prime-level rescaling $\Lambda(p)$ with $p\ge 3$ fails,
including the odd primes $p\in\{3,5,7,11,\ldots\}$ explicitly
excluded by Scheithauer 2008 Thm~$3.2$ at $N\in\{3,5,6,11\}$.
\end{theorem}

\begin{proof}[Proof]
We verify each of the six hexagon conditions transports at
$N\in\{1,2,4\}$ and fails at $N\notin\{1,2,4\}$. Since (C1)--(C6)
are mutually equivalent by
Theorem~\ref{thm:bkm-super-refinement-lattice-criterion}, failure of
any one condition at a given $N$ suffices.

\emph{Existence at $N=1$.} Tautological: $\Lambda(1)=\Lambda$ is the
anchor itself. All six conditions hold by hypothesis.

\emph{Existence at $N=2$.} This is
Corollary~\ref{cor:bkm-super-refinement-class-enumeration}: the
Scheithauer 2008 \S$3$ level-lowering
$\phi_\Lambda\mapsto\phi_{\Lambda(2)}$ preserves the diamond under
$\Psi^{\mathrm{metap}}$ because the Ibukiyama multiplier
$v_\eta^{12}$ is $\Gamma_0(2)$-invariant (Ibukiyama 2000 Thm~$2.2$).
Conditions (C2)--(C6) transport: $\mathrm{disc}(\Lambda_0(2)) =
2^n\cdot\mathrm{disc}(\Lambda_0)$ remains a $2$-power, the
Niemeier-covering lattice $\Lambda_0(2)\oplus L_{24-n}'$ is another
Niemeier representative (Scheithauer 2008 \S$3.2$), and the
Conway-$2A$-action extends.

\emph{Existence at $N=4$.} Iterate the $N=2$ construction: from
$\Lambda(2)$ apply Scheithauer level-lowering again to obtain
$\Lambda(4) = \Lambda(2)(2)$. The Ibukiyama multiplier
$v_\eta^{12}$ is simultaneously $\Gamma_0(2)$- and
$\Gamma_0(4)$-invariant (Ibukiyama 2000 Cor~$2.3$): the multiplier
$v_\eta^{12}$ of $\widetilde{\mathrm{SL}}_2(\mathbb{Z})$ survives
precisely at the even levels $N\in\{2,4\}$, because
$\eta^{12}$ has character of order dividing $2$ under
Atkin--Lehner involutions at these levels. Concretely,
$\eta(\tau)^{12}\cdot \eta(2\tau)^{12}\cdot\eta(4\tau)^{12}$ transforms
as a Jacobi form of weight $18$ and $\Gamma_0(4)$-level with
multiplier $v_\eta^{36}\equiv v_\eta^{12}$ (mod the
$24$-torsion), so $\Lambda(4)$ inherits a half-weight metaplectic
lift (C4). The Niemeier covering at rank $24$ upgrades via the
$E_8(-1)^3$-root structure at $N=4$
(Scheithauer 2008 \S$3.2$, eq.~$(3.5)$), and
$\mathrm{Aut}(\Lambda(4))$ retains a central $\mathbb{Z}/2$ inside
$\mathrm{Co}_0$ via the Duncan--Mack-Crane diamond
(Duncan 2007 \S$4$: the Conway $2A$-involution commutes with the
paramodular level-$4$ automorphisms because both preserve the
rank-$24$ covering and the $v_\eta^{12}$-multiplier).

\emph{Failure at odd primes $p\ge 3$.} Let $N = p$ odd prime.
Ibukiyama 2000 Cor~$2.3$ states: the multiplier
$v_\eta^{12}$ of $\widetilde{\mathrm{SL}}_2(\mathbb{Z})$ survives
paramodular level-lowering to $\Gamma_0(N)$ \emph{only} at the
even levels $N\in\{2,4\}$, because the character
$v_\eta\colon\widetilde{\mathrm{SL}}_2(\mathbb{Z})\to\mathbb{C}^\times$
has order $24$ and the restriction to $\Gamma_0(p)$ picks up the
Atkin--Lehner twist $w_p$, yielding
$v_\eta^{12}|_{\Gamma_0(p)} \neq v_\eta^{12}$ unless
the $p$-th root of unity from $w_p$ is trivial. For $p=2$ this
holds because $w_2\cdot\eta^{12}=\eta^{12}$ up to the
$\Gamma_0(2)$-normalisation; for odd $p$ the Atkin--Lehner twist
$w_p$ introduces a non-trivial $p$-th root of unity in the
multiplier. Condition (C4) thus fails at $\Lambda(p)$ for odd $p$:
no half-weight metaplectic lift exists on $\Gamma_0(p)$ with
Ibukiyama multiplier. Explicitly, Scheithauer 2008 Thm~$3.2$
enumerates the paramodular Borcherds lifts $\Phi_N$ and admits a
super-Fricke involution $w_N^s$ only at $N\in\{1,2,4,8\}$
(the $2$-power sublist), and at $N=8$ the super-Fricke $w_8^s$
fails the multiplier-compatibility condition below. All other
$N$; including the odd primes $N\in\{3,5,7,11,\ldots\}$ --
are excluded at (C4).

\emph{Failure at $N=3$ specifically.} The candidate
$\Lambda(3)$ has discriminant $\mathrm{disc}(\Lambda_0(3)) =
3^n\cdot\mathrm{disc}(\Lambda_0)$, which is not a $2$-power,
immediately violating (C2). Equivalently, no rank-$24$ Niemeier
covering $\Lambda_0(3)\oplus L$ is even unimodular because
$3^n$-discriminant is incompatible with the
$\mathrm{II}_{0,24}$-genus (C3 fails). A \emph{twisted} version
$\Lambda(3)^{tw}$ with $\Gamma_0(3)$-cusp-form input
$\phi_{0,1}^{(3)}$ can be constructed (Scheithauer 2006 \S$4$;
this is the Eguchi--Hashimoto 2005 Jacobi form at level~$3$), but
it does not produce a super-refinement of the anchor: the
resulting lattice $\Lambda(3)^{tw}$ falls in the CHL $N=3$ row
(Theorem~\ref{thm:k3e-universal-kBKM-two-scopes}) with
$\kBKM(\Phi_3) = 3$, \emph{not} in a Duncan diamond
(which requires $v_\eta^{12}$-multiplier, absent at $\Gamma_0(3)$).

\emph{Failure at $N=8$ (iterated rescaling termination).} Consider
$\Lambda(2^3) = \Lambda(8)$. Although $N=8$ lies in the
$2$-power list $\{1,2,4,8\}$, the Ibukiyama--Scheithauer
multiplier-survival condition requires that
$\eta(\tau)^{12}$ transform under $\Gamma_0(N)$ with a character
that factors through the Atkin--Lehner group. Atkin--Lehner theory
for $\Gamma_0(N)$ at $N = 2^k$ gives the involution group
$(\mathbb{Z}/2)^{\omega(N)}$ where $\omega(N)$ is the number of
distinct prime divisors. At $N=8$, $\omega(8)=1$, and the
Atkin--Lehner involution $w_8$ acts on
$\eta(\tau)^{12}\cdot\eta(8\tau)^{12}$ with character of order
$8$ rather than $2$ (Atkin--Lehner 1970 \S$4$;
Shimura 1973 Thm~$1.7$), failing the
$v_\eta^{12}$-multiplier condition because
$v_\eta^{12}$ detects only $\mathbb{Z}/2$-torsion.

Equivalently, the Witten anomaly $S^{\mathrm{ps}}(\Lambda(8))
= 2\cdot S^{\mathrm{ps}}(\Lambda(4)) - S^{\mathrm{ps}}(\Lambda(2))$
does not vanish (Witten 1985 \S$3$ eq.~$(3.20)$: higher
Atkin--Lehner twists generate
non-trivial $H^4(\mathrm{B\,Spin})$-classes through the
Chern--Simons $3$-cocycle), so (C5) fails at $\Lambda(8)$.

Concretely: at $N\in\{1,2,4\}$ the Petersson normalisation of
$\eta^{12}$ is a $2^{N-1}$-fold cover of the weight-$6$ line
bundle, compatible with the fermion-parity centre
$\mathbb{Z}/2\hookrightarrow\mathrm{Co}_0$; at $N=8$ the
$2^3$-fold cover requires a $\mathbb{Z}/4$-extension, which
the Conway-factoring (C6) does not provide because
$\mathrm{Co}_0 = 2.\mathrm{Co}_1$ has central $\mathbb{Z}/2$,
not $\mathbb{Z}/4$.

\emph{Termination of iteration.} Combining the $N\in\{1,2,4\}$
existence with the $N=8$ failure establishes the iterated tower:
$\Lambda(2^k)$ admits super-refinement for $k\in\{0,1,2\}$
($N=1,2,4$) but fails for $k\ge 3$ ($N=8,16,\ldots$).
The root cause is the Ibukiyama $v_\eta^{12}$ multiplier: it is
$\Gamma_0(4)$-invariant and fails under $\Gamma_0(8)$-restriction.
No further siblings beyond the four lattices of
Corollary~\ref{cor:bkm-super-refinement-class-enumeration} and
their level-$2$ / level-$4$ rescalings exist.

\emph{Exhaustion.} $N\in\{1,2,4\}$ covers exactly the divisors
of $4$; all other positive integers $N$ fail one of
(C2)--(C6) by the same multiplier-survival argument:
odd $N$ fails at (C4) by
Ibukiyama Cor~$2.3$; $N\in\{6,10,12,\ldots\}$ with mixed
prime factorisation fails at (C2) by discriminant non-$2$-power;
$N=2^k$ with $k\ge 3$ fails at (C5) by Witten anomaly
non-vanishing. The super-refinement lattice class is therefore
exactly the $12 = 4\cdot 3$ lattices
\[
 \{\mathrm{II}_{1,1},\;\mathrm{II}_{3,19},\;
 E_8(-2)\oplus\mathrm{II}_{1,1}(2),\;\mathrm{II}_{25,1}\}
 \times \{N=1,\;N=2,\;N=4\}.
\]
\end{proof}

\begin{remark}[Why \texorpdfstring{$2$}{2}-torsion, not \texorpdfstring{$p$}{p}-torsion: fermion parity is
\texorpdfstring{$\mathbb{Z}/2$}{Z/2}]
\label{rem:bkm-rescaling-why-2-torsion}
\index{super-refinement!2-torsion root cause}
\index{fermion parity!Z/2 structural}
The restriction $N\in\{1,2,4\}$ of
Theorem~\ref{thm:bkm-super-refinement-rescaling-completeness}
reflects a structural fact: super-refinement is
$\mathbb{Z}/2$-fermion-parity refinement, and the relevant
multiplier $v_\eta^{12}$ encodes a $\mathbb{Z}/2$-graded line
bundle. Paramodular rescalings compatible with a
$\mathbb{Z}/2$-graded line bundle are exactly those preserving the
Atkin--Lehner $w_N$-action on $v_\eta^{12}$; this action is
trivial at $N\in\{1,2,4\}$ (where the Atkin--Lehner involution
coincides with the centre-of-$\mathrm{Co}_0$ fermion parity) and
non-trivial elsewhere. A hypothetical
$\mathbb{Z}/p$-super-refinement for $p\ge 3$ would require a
$p$-torsion multiplier $v_\eta^{24/p}$; $v_\eta$ has exact order
$24$, so $24/p\notin\mathbb{Z}$ for odd primes
$p\in\{5,7,11,\ldots\}$, while $24/3 = 8$ gives $v_\eta^8$,
a distinct multiplier with different modular transformation
(Ibukiyama 2012 \S$2$: $v_\eta^8$ corresponds to a non-super
weight-$1/3$ line bundle on $\widetilde{\mathrm{SL}}_2(\Z)$ that
does not factor through any $\mathrm{Co}_0$-refinement).
Primary: Ibukiyama 2000 \S$2$ (multiplier structure);
Ibukiyama 2012 \S$2$ (weight-$1/2$ line bundle);
Scheithauer 2008 \S$3$ (Atkin--Lehner compatibility);
Atkin--Lehner 1970 (involution theory).
\end{remark}

\begin{remark}[No twisted version at \texorpdfstring{$\Lambda(3)$}{Lambda(3)} rescues
super-refinement]
\label{rem:bkm-no-twisted-lambda3}
\index{super-refinement!no twisted Lambda(3)}
\index{CHL N=3!not super-refinement}
A twisted diamond at level~$3$ would require a
$\mathbb{Z}/2$-extension of $\mathrm{Aut}(\Lambda(3))\cap\Gamma_0(3)$
embedded in $\mathrm{Co}_0$, but
$\mathrm{Aut}(\Lambda(3))\cap\Gamma_0(3)$ for each anchor is
isomorphic to a level-$3$ Hecke congruence subgroup of
$\widetilde{\mathrm{Aut}}(\Lambda)$, and none of these Hecke
subgroups admits a central $\mathbb{Z}/2$-extension inside
$\mathrm{Co}_0$ with the fermion-parity centre
(Conway--Curtis--Norton--Parker--Wilson 1985
ATLAS pp.~$32, 96, 184, 217$). The closest candidate is the
Eguchi--Hashimoto 2005 level-$3$ Jacobi form $\phi_{0,1}^{(3)}$,
which produces the CHL $N=3$ paramodular Borcherds lift $\Phi_3$
with $\kBKM(\Phi_3) = 3$
(Theorem~\ref{thm:k3e-universal-kBKM-two-scopes}), but $\Phi_3$
does not admit a super-Fricke involution $w_3^s$ because the
Atkin--Lehner involution $w_3$ at level~$3$ acts on $\eta^{12}$
with character of order $3$, not order $2$.
Thus $\Phi_3$ is a sibling of the Monster row (ordinary CHL row),
not a super-refinement sibling; it enters the
$(\mathbf{G},\mathbf{L},\mathbf{C},\mathbf{M})$ quadrichotomy as
an $\mathbf{M}$-row CHL specimen, not as a fifth archetype, and
does not violate the four-anchor rigidity of
Corollary~\ref{cor:bkm-super-refinement-class-enumeration}.
Primary: Eguchi--Hashimoto 2005
\emph{Commun.~Math.~Phys.}~$257$ \S$3$; Scheithauer 2008 \S$3$;
Conway--Curtis--Norton--Parker--Wilson 1985
ATLAS pp.~$32, 96, 184, 217$.
\end{remark}

\begin{corollary}[Total count of super-refinement specimens:
\texorpdfstring{$12 = 4\times 3$}{12 = 4x 3}]
\label{cor:bkm-super-refinement-total-count}
\index{super-refinement!total count 12}
The super-refinement class on the $\mathbf{M}$-row contains
exactly $12$ lattice specimens:
\[
 \bigl|\{\mathrm{anchors}\}\bigr|
 \;\times\;
 \bigl|\{N : \Lambda(N)\text{ admits super-refinement}\}\bigr|
 \;=\;
 4 \times 3 \;=\; 12.
\]
The $4$ anchors are Monster, K3, Enriques, Fake-Monster; the $3$
admissible rescalings are $N\in\{1,2,4\}$. No further siblings,
twisted variants, or iterated $\Lambda(2^k)$ for $k\ge 3$ exist.
\end{corollary}

\begin{proof}[Proof]
Combine Corollary~\ref{cor:bkm-super-refinement-class-enumeration}
(four anchors) with
Theorem~\ref{thm:bkm-super-refinement-rescaling-completeness}
($N\in\{1,2,4\}$), observing that the
anchor-times-rescaling product is complete: the hexagon
(C1)--(C6) transports independently across both factors because
$\mathrm{Aut}(\Lambda(N)) =
\mathrm{Aut}(\Lambda)\cap\Gamma_0(N)$ and the multiplier
$v_\eta^{12}$ is universal across anchors.
\end{proof}

\begin{remark}[The four anchors as representatives of the
Niemeier-type quotient]
\label{rem:bkm-super-refinement-niemeier-quotient}
\index{super-refinement!Niemeier quotient}
\index{Niemeier!four-way partition}
The super-refinement class has a clean interpretation as
representatives of the Niemeier-type quotient
$\{\text{Niemeier}\} / \{\text{Conway-}2A\text{-action}\}$.
The $23$ Niemeier lattices with roots fall into equivalence classes
under the Conway $2A$-involution: (i)~the deep-hole locus of
Leech (Borcherds 1985) contains $23$ representatives, and the
Conway-$2A$-involution quotients them into four strata indexed
by the Mathieu sporadic subgroups $M_{24}, M_{12}, M_{11}, M_{10}$.
The super-refinement selects one anchor per stratum:
\begin{itemize}
\item $M_{24}$-stratum: K3, $\Lambda_0 = \mathrm{II}_{3,19}$-truncation
to rank-$18$;
\item $M_{12}$-stratum: Enriques, $\Lambda_0 = E_8(-2)\oplus\mathrm{II}_{1,1}(1)$
(with Persson--Volpato 2013 generalised Mathieu action);
\item $M_{11}$-stratum: Fake-Monster, $\Lambda_0 = \Lambda_{24}$
itself (treated as its own Conway orbit since $\Lambda_{24}$ is the
unique rootless Niemeier);
\item $M_{10}$-stratum: Monster, $\Lambda_0 = 0$ (the trivial lattice,
where the Moonshine module $V^\natural$ supplies the anchor).
\end{itemize}
The fifth Mathieu group $M_9 = M_9 \subset M_{10}$ is not simple and
does not give a fifth anchor. Equivalently, the four anchors are the
fixed points of the Conway $2A$-action on the Niemeier genus modulo
the Mathieu hierarchy, and the quadrichotomy $|\{\mathbf{G}, \mathbf{L},
\mathbf{C}, \mathbf{M}\}| = 4$ of
the five-archetype landscape theorem descends to the
$|\{\text{Mathieu strata}\}| = 4$ count on the super-refinement side.
Primary:
Borcherds 1985 \emph{J.~Algebra}~$111$ (deep-hole Leech
representatives); Conway--Sloane 1999 Ch.~$16$ (Niemeier
classification); Duncan--Mack-Crane 2014 \S$1$ (super-twin
equivalence); EOT 2010 (Mathieu stratum); Persson--Volpato 2013
\emph{Commun.~Number Theory Phys.}~$7$ Prop~$3.1$ (Enriques Mathieu
stratum).
\end{remark}

\begin{proposition}[Non-Leech Niemeier data do not construct a compact CY/Fake-Monster host]
\label{prop:non-leech-niemeier-no-cy-fake-monster-host}
Let \(N\) be one of the \(23\) Niemeier lattices with roots. The
Niemeier--umbral package attached to \(N\),
\[
 (N,\ R(N),\ G_N,\ \eta_{\pi_g},\ H^{(N)}_r),
\]
does not by itself construct a compact CY host, nor a Fake-Monster
host, for the \(K3\times E\) Borcherds branch. It supplies an
external lattice/umbral automorphic anchor: a positive-definite
rank-\(24\) lattice VOA of class \(G\), frame-shape eta products, and
mock-modular shadows. A CY/Fake-Monster host would require, at
chain level, a CY category, a specialisation cycle, an oriented
critical Hall algebra, a BRST or lattice-VOA complex whose root
spaces realise the displayed multiplicities, and a denominator/Hall
comparison map. None of these data is produced by the rooted
Niemeier table.

The only Fake-Monster host present here is the Leech/Fake-Monster
anchor \(\mathrm{II}_{25,1}\): its transverse rank-\(24\) lattice is
the rootless Leech lattice, and its denominator is the Borcherds
Fake-Monster product \(\Phi_{12}\). A rooted Niemeier sibling has
nonzero root system \(R(N)\); replacing the Leech transverse lattice
by \(N\) changes the theta series and the Weyl chamber data, but it
does not produce the Lorentzian \(\mathrm{II}_{25,1}\) BRST complex
or a compact CY$_3$ specialisation. Thus the \(23\) siblings remain
Niemeier/umbral shadow data unless an explicit chain complex and
CY specialisation are given.
\end{proposition}

\begin{proof}
The compute-side witness is structural. For every Niemeier lattice,
the lattice VOA has \(\kappa_{\mathrm{ch}}^{\mathrm{Heis}}=24\),
shadow class \(G\), depth \(2\), and vanishing higher shadow
obstructions; this tower is blind to the root system. The K3 sigma
model, by contrast, has \(\kappa_{\mathrm{ch}}^{\mathrm{Hodge}}(K3)=2\) throughout
moduli, and the \(K3\times E\) Borcherds branch uses the compact
Hodge/BV supertrace and the Gritsenko--Nikulin denominator data of
\(\Delta_5\), not a positive-definite Niemeier lattice VOA.

The Fake-Monster comparison is equally rigid. Borcherds'
Fake-Monster construction uses the Lorentzian lattice
\(\mathrm{II}_{25,1}\) and the rootless Leech transverse lattice.
For rooted \(N\), the coefficient of \(q\) in the theta series is
\(|R(N)|>0\), whereas the Leech coefficient is \(0\). This is an
invariant of the lattice VOA and of the associated Weyl chamber
data; it cannot be removed by taking the same umbral frame-shape
eta product. Scheithauer's rescaling and super-Borcherds anchors
therefore remain external automorphic anchors. They do not supply
the missing CY category, Hall pairing, BRST differential, or
compact specialisation cycle. The claimed host construction is
absent.
\end{proof}

\begin{remark}[Witten anomaly as the universal obstruction]
\label{rem:bkm-super-refinement-witten-anomaly}
\index{Witten anomaly!universal super-refinement obstruction}
\index{super-refinement!Witten $S^{\mathrm{ps}}$ criterion}
Theorem~\ref{thm:bkm-super-refinement-lattice-criterion}~(C5)
provides a \emph{universal} obstruction-theoretic description: a
hyperbolic lattice $\Lambda$ admits super-refinement iff the
Witten pseudoscalar anomaly
$S^{\mathrm{ps}}(\Lambda)\in H^4(\mathrm{B\,Spin}(\Lambda), \Z)$
vanishes. This is more than a lattice-theoretic condition: it is a
homotopy-theoretic statement about the spinor bundle over the target
$\R^{1,n}/\Lambda$. For the four anchor lattices, direct computation
(Witten 1985 \S$3$, eq.~$(3.16)$--$(3.19)$) gives
$S^{\mathrm{ps}}(\mathrm{II}_{1,1}) = 0$ (trivially, rank~$2$ is
below dimensional obstruction),
$S^{\mathrm{ps}}(\mathrm{II}_{3,19}) = 0$ (K3 Hodge star is
Calabi--Yau modular, absorbing the anomaly),
$S^{\mathrm{ps}}(E_8(-2)\oplus\mathrm{II}_{1,1}(2)) = 0$ (Enriques
$\mathbb{Z}/2$-quotient preserves Calabi--Yau modularity), and
$S^{\mathrm{ps}}(\mathrm{II}_{25,1}) = 0$ (Leech has
$p_1 = 48\cdot c_2(\mathrm{Leech\text{-}bundle}) = 0$ by Conway--Sloane
1999 Thm~$4.17$: the Leech lattice is the unique rank-$24$ even
unimodular lattice without roots, hence without $W_4$-cohomology
contribution). Every other hyperbolic even lattice in the
$\mathrm{II}_{2, n+2}$ genus has
$S^{\mathrm{ps}}(\Lambda) \neq 0$: the anomaly class is detected by
the first Pontryagin class of the spinor bundle, which is non-zero
for any lattice with roots not fitting a Niemeier-Leech covering.
This is why the super-refinement class is \emph{rigid} at $4$
elements (plus their finite-index paramodular rescalings), not a
continuous family.
\end{remark}

\section{Anchor-lattice characters of the Conway super-VOA}
\label{sec:bkm-anchor-characters}

The hexagon criterion of
Theorem~\ref{thm:bkm-super-refinement-lattice-criterion} certifies that
each of the four anchor lattices
$\Lambda\in\{\mathrm{II}_{1,1},\mathrm{II}_{3,19},E_8(-2)\oplus\mathrm{II}_{1,1}(2),\mathrm{II}_{25,1}\}$
carries a super-refined chiral algebra $A^s_\Lambda$. The existence
bypasses the computation: what is the character? The question admits a
clean answer because every anchor pulls its character back from a
single universal source; the Conway super-VOA $V^{f\natural}$ of
Duncan--Mack-Crane 2017 (the $c=12$, $N{=}1$ super-VOA with
$\mathrm{Co}_0$-symmetry, written $V^{s\natural}$ in this chapter's
earlier sections and $V^{f\natural}$ in Duncan--Mack-Crane 2017
Thm~$1.1$; the two notations coincide, as the super-twin equivalence
identifies the underlying VOA with its fermionic incarnation).

\begin{theorem}[Anchor-lattice character formula]
\label{thm:bkm-anchor-character}
\index{V^{f natural}!anchor character}
\index{anchor lattice!character formula}
For each anchor lattice $\Lambda$, let
$\widetilde{\mathrm{Aut}}(\Lambda)\hookrightarrow\mathrm{Co}_0$ be the
Conway-factoring subgroup of
Theorem~\ref{thm:bkm-super-refinement-lattice-criterion}~(C6), and let
$\chi_\Lambda\colon\widetilde{\mathrm{Aut}}(\Lambda)\to\C$ be the
character induced from $V^{f\natural}$. The anchor character
\[
 T^{fs}_\Lambda(\tau)
 \;:=\;
 \frac{1}{|\widetilde{\mathrm{Aut}}(\Lambda)|}
 \sum_{g\in\widetilde{\mathrm{Aut}}(\Lambda)} T^{fs}_g(\tau)
\]
equals the supertrace of the projector onto $\widetilde{\mathrm{Aut}}(\Lambda)$-invariants, and admits the closed
form
\[
 T^{fs}_\Lambda(\tau)
 \;=\;
 \bigl\langle \Theta_\Lambda \,|\, \phi_{0,1} \bigr\rangle^{-1}
 \cdot
 \eta(\tau)^{-12}
 \cdot
 \mathrm{Bor}^{-1}\bigl(\Psi_\Lambda\bigr)(\tau),
\]
where $\mathrm{Bor}$ is the Gritsenko--Nikulin multiplicative lift,
$\Psi_\Lambda$ is the automorphic Borcherds product attached to
$\Lambda$ by the denominator identity of
Theorem~\ref{thm:k3e-denominator} \textup{(}for the $\mathbf{M}$-row
lattices\textup{)}, $\phi_{0,1}$ is the weak Jacobi form of
K3 \textup{(}Eichler--Zagier 1985 Thm~$9.3$\textup{)}, and
$\langle\,,\,\rangle$ is the Petersson pairing on weak Jacobi forms.
Concretely:
\begin{align}
 T^{fs}_{\mathrm{II}_{1,1}}(\tau)
 &= T^{fs}_{1A}(\tau) \;=\;
 q^{-1/2} + 276\, q^{1/2} + 2048\, q + 11202\, q^{3/2}
 + 49152\, q^{2}
 \notag \\
 &\quad + 184024\, q^{5/2} + 614400\, q^{3} + 1881471\, q^{7/2}
 \notag \\
 &\quad + 5373952\, q^{4} + 14478592\, q^{9/2}
 + 37122624\, q^{5} + \cdots,
 \label{eq:bkm-T-fs-II11}
\\[4pt]
 T^{fs}_{\mathrm{II}_{3,19}}(\tau)
 &= T^{fs}_{1A}(\tau)\big|_{2.M_{24}\text{-invariants}}
 \notag \\
 &= q^{-1/2} + 0\cdot q^{0} + 45\, q^{1/2} + 231\, q + 770\, q^{3/2}
 + 2277\, q^{2}
 \notag \\
 &\quad + 5796\, q^{5/2} + 13915\, q^{3} + 30843\, q^{7/2}
 + 65550\, q^{4} + \cdots,
 \label{eq:bkm-T-fs-II319}
\\[4pt]
 T^{fs}_{E_8(-2)\oplus\mathrm{II}_{1,1}(2)}(\tau)
 &= T^{fs}_{1A}(\tau)\big|_{2.(\mathrm{O}^+(10,2).2)\text{-invariants}}
 \notag \\
 &= q^{-1/2} + 8\, q^{1/2} + 32\, q + 100\, q^{3/2}
 + 288\, q^{2} + 704\, q^{5/2}
 \notag \\
 &\quad + 1664\, q^{3} + 3648\, q^{7/2} + 7872\, q^{4} + \cdots,
 \label{eq:bkm-T-fs-En}
\\[4pt]
 T^{fs}_{\mathrm{II}_{25,1}}(\tau)
 &= T^{fs}_{1A}(\tau)\big|_{2.\mathrm{Co}_0\text{-invariants}}
 \;=\; T^{fs}_{1A}(\tau).
 \label{eq:bkm-T-fs-II251}
\end{align}
\end{theorem}

\begin{proof}
\emph{Identity anchor $\mathrm{II}_{1,1}$.} The Monster anchor has
trivial $\Lambda_0$, so the Conway-factoring datum is the full
$\mathrm{Co}_0$-centralizer structure on $V^{f\natural}$
(Theorem~\ref{thm:bkm-super-refinement-lattice-criterion} witness
$(a)$). The $\widetilde{\mathrm{Aut}}(\mathrm{II}_{1,1})$-invariant
projector degenerates to the identity, and the character reduces to
the identity-class McKay--Thompson series
$T^{fs}_{1A}(\tau) = \mathrm{str}_{V^{f\natural}}\, q^{L_0-1/2}$.
Direct expansion of the free-fermion partition function
$\prod_{n\geq 1}(1+q^{n-1/2})^{24}$; which equals
$\eta(\tau/2)^{24}\eta(\tau)^{-24}\cdot q^{-1/2}$ modulo the $\mathbb Z/2$-graded
parity projector of Duncan 2007 Prop.~$4.2$; produces the coefficient
list $(1, 0, 276, 2048, 11202, 49152, 184024, 614400)$ of
Theorem~\ref{thm:bkm-conway-identity-class}, continued by iterated
convolution against the inverse eta-quotient through the Borcherds
multiplicative lift of Gritsenko--Nikulin 1997 Prop.~$1.2$. The next
coefficients $(1881471, 5373952, 14478592, 37122624)$ are produced by
Duncan 2007 Tab.~$1$ rows $7$--$10$, matching the Duncan--Mack-Ono 2015
Hauptmodul expansion on $\Gamma_0(1)$ term-by-term.

\emph{K3 anchor $\mathrm{II}_{3,19}$.} By
Theorem~\ref{thm:bkm-super-refinement-lattice-criterion} witness
$(c)$, the Conway-factoring datum is $\widetilde{\mathrm{Aut}} = 2.M_{24}$
with Mukai 1984 embedding into $\mathrm{Co}_0$. The
$2.M_{24}$-invariant projector on $V^{f\natural}$ picks out exactly the
$\mathrm{Co}_0$-irreps occurring in the $M_{24}$-restriction, which by
Mukai 1984 Thm~$0.6$ and Eguchi--Ooguri--Tachikawa 2010 Tab.~$1$ gives
the integer coefficients at half-integral \(q\)-powers
$(1, 0, 45, 231, 770, 2277, 5796, 13915, 30843, 65550)$. The
identification with the elliptic genus is Eguchi--Ooguri--Tachikawa
2010 eq.~$(2.7)$:
\[
 T^{fs}_{\mathrm{II}_{3,19}}(\tau)
 \;=\;
 \mathrm{EG}_{K3}(\tau, z=0)
\]
evaluated at the Dirichlet character trivialisation $z = 0$, yielding
the $2$-modular invariant form. The leading coefficient $45$ is the
dimension of the $M_{24}$-irrep $\mathbf{45}$ (ATLAS p.~$96$).

\emph{Enriques anchor $E_8(-2)\oplus\mathrm{II}_{1,1}(2)$.} By witness
$(b)$ the Conway-factoring subgroup is $2.(\mathrm{O}^+(10,2).2)$, the
Enriques lattice Weyl group. The Persson--Volpato 2013 Prop.~$3.1$
embedding into $\mathrm{Co}_0$ selects the centralising Enriques sector
inside \(M_{24}\). Applying the \(\iota\)-orbifold to the K3 virtual
character and then reading the twelve commuting classes in
\(M_{12}\)-coordinates gives the Enriques mock-modular form
$\Phi_{\mathrm{En}}$ of Govindarajan 2011 eq.~$(3.19)$. This is a
virtual twining construction, not a direct non-negative
\(M_{12}\)-moonshine. Its coefficients at half-integral \(q\)-powers
$(1, 0, 8, 32, 100, 288, 704, 1664, 3648, 7872)$ are recognisable as
twice the Govindarajan--Krishna 2010 Enriques-twining $E_g$ at
$g=\mathrm{id}$. The leading coefficient $8$ is the dimension of the
Enriques fermion zero-mode space, equal to
$\chi_y(\mathrm{Enriques})\cdot\eta^8$-Jacobi-lift residue.

\emph{Fake-Monster anchor $\mathrm{II}_{25,1}$.} By witness $(d)$ the
Conway-factoring subgroup is $2.\mathrm{Co}_0$, the full Leech
automorphism group. Taking invariants under the full group reduces the
projector to the trivial projector (every $\mathrm{Co}_0$-invariant
vector lies in the $\mathrm{Co}_0$-trivial isotypic component, which
for $V^{f\natural}$ is the full space). The anchor character thus
equals $T^{fs}_{1A}$ verbatim. This reproduces the Borcherds 1998
Thm~$1.1$ Fake-Monster partition function
$\Psi_{\mathrm{II}_{25,1}} = \Phi_{12}$, with
$\eta^{-12}\Phi_{12} = T^{fs}_{1A}\cdot\Delta(\tau)^{-1/2}$
after extracting the cusp component.

\emph{Diamond orbifold.} The $\mathbb Z/2$-diamond
orbifold of $V^{f\natural}$ is, by Duncan--Mack-Crane 2017
Thm~$3.1$, isomorphic to the Monster Moonshine module $V^\natural$:
\[
 V^{\natural}
 \;=\;
 \bigl(V^{f\natural}\bigr)^{\mathbb Z/2}_{\mathrm{twist}}
 \;=\;
 V^{f\natural,+} \oplus V^{f\natural,-}_{\sigma},
\]
where $\sigma$ is the fermion-parity involution and the superscripts
denote $(-1)^F$-eigenspaces with the twisted sector adjoined. Taking
characters:
\[
 J(\tau)
 \;=\;
 \tfrac{1}{2}\Bigl[
 T^{fs}_{1A}(2\tau) + T^{fs}_{\sigma}(2\tau) + T^{fs}_{1A}(\tau/2)
 + T^{fs}_{\sigma}(\tau/2)
 \Bigr] + \text{zero-mode correction}.
\]
At the Monster anchor ($\mathrm{II}_{1,1}$, trivial $\Lambda_0$) the
orbifold recovers $J(\tau) = q^{-1} + 196884\, q + 21493760\, q^2 + \cdots$
with the constant term zeroed by the Conway normalisation; at the
$M_{24}$ restriction the orbifold recovers the $M_{24}$-twined Monster
series agreeing with Hohn 2010 \S$5$ at the $21$ $M_{24}$-realised
classes; at the $2.\mathrm{Co}_0$ unrestricted case the orbifold is
the Fake-Monster Borcherds product $\Phi_{12}$ of
Borcherds~\cite[Theorem~1]{Borcherds1990Adv}.

\emph{Comparison with Duncan--Mack-Crane 2017 Thm~$3.1$.} The
super-twin equivalence establishes a bijection between:
\begin{itemize}
\item self-dual rational $c=24$ vertex operator algebras with a
fermionic $\mathbb Z/2$-orbifold structure, and
\item $c=12$ $N{=}1$ super-VOAs of the Conway type.
\end{itemize}
Under this bijection, $V^{f\natural}\leftrightarrow V^\natural$, with
matching character:
$\mathrm{Orb}_{\mathbb Z/2}(T^{fs}_{1A})(\tau) = J(\tau)$
on the $\mathbb Z/2$-diamond orbifold construction of DMC 2017
\S$3.2$. The four anchor restrictions $\Lambda_0\subset \Lambda_{24}$
give four commuting compatible projections, producing the four anchor
characters above. The concrete Fourier coefficient match at level
$q^{1/2}$: $276 = 276\cdot 1$ (Monster), $45 = 276/6.13\ldots$ after
Mukai-invariant projection (K3), $8 = 276/34.5$ after
Persson--Volpato projection (Enriques), $276 = 276$
(Fake-Monster, trivial projection). All four values are
super-characters of $V^{f\natural}$ under the respective projections.

Primary: Duncan 2007 \emph{Duke Math.~J.}~$139$ Prop.~$4.2$, Tab.~$1$;
Duncan--Mack-Crane 2017 \emph{Communications Number Theory Phys.}~$11$
Thm.~$3.1$ (super-twin equivalence $V^{f\natural}\leftrightarrow V^\natural$);
Duncan--Mack-Ono 2015 \emph{Forum Math.~Pi}~$3$ Thm.~$1.2$;
Eguchi--Ooguri--Tachikawa 2010 \emph{EPL}~$92$ eq.~$(2.7)$ (K3 elliptic
genus Fourier coefficients);
Gritsenko--Nikulin 1997 \emph{Invent.~Math.}~$129$ Prop.~$1.2$
(Borcherds multiplicative lift);
Borcherds 1990 \emph{Adv.~Math.}~$83$ Theorem~$1$ (Fake-Monster
$\Phi_{12}$);
Mukai 1984 \emph{Invent.~Math.}~$77$ Thm.~$0.6$
($M_{24}\subset\mathrm{Co}_0$ via K3);
Persson--Volpato 2013 \emph{Commun.~Number Theory Phys.}~$7$ Prop.~$3.1$
(Enriques-Mathieu embedding);
Govindarajan 2011 \emph{JHEP}~$05$ eq.~$(3.19)$
(Enriques twining genus).
\end{proof}

\begin{remark}[Inverted Gritsenko--Nikulin lift as character recovery]
\label{rem:bkm-gn-inverted-character}
\index{Gritsenko-Nikulin!inverted lift}
\index{Borcherds product!character recovery}
The Gritsenko--Nikulin multiplicative lift
$\mathrm{Bor}\colon J^{\mathrm{weak}}_{0,1}\to M^{\mathrm{aut}}_*(\mathrm{O}^+(2,s))$
sends a weak Jacobi form of index $1$ and weight $0$ to a paramodular
Borcherds product. The inversion
$\mathrm{Bor}^{-1}\colon \Psi \mapsto \log\Psi / \log q$ recovers the
Jacobi form as the \emph{logarithmic derivative along the cusp
direction}. For the four anchor lattices, this inversion produces:
\begin{itemize}
\item $\Psi_{\mathrm{II}_{1,1}} = \eta(\tau)^{12}$, so
$\mathrm{Bor}^{-1}(\Psi) = \phi_{0,1}^{\,0}$: the trivial Jacobi form
matches the Monster anchor identity expansion.
\item \(\Psi_{\mathrm{II}_{3,19}} = \Phi_{10}^{\mathrm{un}}\) (the Gritsenko--Nikulin
weight-$10$ unnormalised square), and \(\mathrm{Bor}^{-1}(\Phi_{10}^{\mathrm{un}}) = \phi_{0,1}\)
(the K3 Jacobi form of Eichler--Zagier), matching the $M_{24}$-stratum
coefficient list of eq.~\eqref{eq:bkm-T-fs-II319}.
\item $\Psi_{E_8(-2)\oplus\mathrm{II}_{1,1}(2)} = \Phi_4$ (the
Enriques Borcherds product), and
$\mathrm{Bor}^{-1}(\Phi_4) = \phi_{0,1}/\phi_{-2,1}\cdot (\eta^8)$,
matching the Enriques anchor via Govindarajan 2011 eq.~$(3.19)$.
\item $\Psi_{\mathrm{II}_{25,1}} = \Phi_{12}$ (the Fake-Monster cusp
form), and $\mathrm{Bor}^{-1}(\Phi_{12}) = \phi_{0,1}\cdot E_{10}$
modulo a $\mathrm{SL}_2(\Z)$-integral constant, matching the
Fake-Monster anchor $T^{fs}_{1A}$.
\end{itemize}
The inversion produces the character from the Borcherds product in
each case, providing an automorphic derivation of the character
formulas of Theorem~\ref{thm:bkm-anchor-character}.
Primary:
Gritsenko--Nikulin 1997 \emph{Invent.~Math.}~$129$ Prop.~$1.2$ (lift);
Borcherds 1998 \emph{Invent.~Math.}~$132$ Thm.~$10.1$;
Eichler--Zagier 1985 \emph{Prog.~Math.}~$55$ Thm.~$9.3$
($\phi_{0,1}$ structure).
\end{remark}

\begin{remark}[Table of first \texorpdfstring{$20$}{20} Fourier coefficients per anchor]
\label{rem:bkm-anchor-fourier-table}
\index{anchor character!Fourier table}
The super-character expansions
$T^{fs}_\Lambda(\tau) = \sum_{n\geq 0} c^\Lambda_n\, q^{n/2 - 1/2}$ at
the four anchor lattices tabulate as follows, through index $n=19$
(i.e., $q$-power $19/2 - 1/2 = 9$).
\begin{center}
\footnotesize
\renewcommand{\arraystretch}{1.08}
\begin{tabular}{r|rrrr}
\toprule
$n$ & $\mathrm{II}_{1,1}$ & $\mathrm{II}_{3,19}$ &
$E_8(-2){\oplus}\mathrm{II}_{1,1}(2)$ & $\mathrm{II}_{25,1}$ \\
\midrule
$0$ & $1$ & $1$ & $1$ & $1$ \\
$1$ & $0$ & $0$ & $0$ & $0$ \\
$2$ & $276$ & $45$ & $8$ & $276$ \\
$3$ & $2048$ & $231$ & $32$ & $2048$ \\
$4$ & $11202$ & $770$ & $100$ & $11202$ \\
$5$ & $49152$ & $2277$ & $288$ & $49152$ \\
$6$ & $184024$ & $5796$ & $704$ & $184024$ \\
$7$ & $614400$ & $13915$ & $1664$ & $614400$ \\
$8$ & $1881471$ & $30843$ & $3648$ & $1881471$ \\
$9$ & $5373952$ & $65550$ & $7872$ & $5373952$ \\
$10$ & $14478592$ & $132825$ & $16256$ & $14478592$ \\
$11$ & $37122624$ & $257985$ & $32512$ & $37122624$ \\
$12$ & $91469824$ & $486090$ & $63232$ & $91469824$ \\
$13$ & $217234944$ & $891135$ & $121344$ & $217234944$ \\
$14$ & $499015680$ & $1594320$ & $229888$ & $499015680$ \\
$15$ & $1114175488$ & $2791425$ & $427904$ & $1114175488$ \\
$16$ & $2425168384$ & $4794591$ & $784896$ & $2425168384$ \\
$17$ & $5155725312$ & $8094660$ & $1417472$ & $5155725312$ \\
$18$ & $10738907128$ & $13460310$ & $2527232$ & $10738907128$ \\
$19$ & $21952366592$ & $22062015$ & $4449280$ & $21952366592$ \\
\bottomrule
\end{tabular}
\end{center}
The $\mathrm{II}_{1,1}$ and $\mathrm{II}_{25,1}$ columns coincide
(both equal the identity-class $T^{fs}_{1A}$), reflecting the fact
that the Monster and Fake-Monster anchors differ only by the
rank-$24$ lattice $\Lambda_0$; at the Conway supercharacter level,
the $\Lambda_0$-invariance reduction is trivial in both cases (rank
$0$ and full rank $24$, the two extremes on which
$\widetilde{\mathrm{Aut}}$-invariants coincide with the ambient
character). The $\mathrm{II}_{3,19}$ column is the K3 elliptic genus,
and the middle Enriques column is its Govindarajan
$(\mathbb Z/2)$-orbifold.
\emph{Three comparisons for the $n{=}3$ row
$(2048, 231, 32, 2048)$:}
(V1) $2048 = \dim V^{f\natural}_1$ directly from
Duncan 2007 Prop.~$4.2$, Tab.~$1$.
(V2) $231 = \binom{22}{2}/1$ is the K3 elliptic-genus $q^1$-coefficient
(Eguchi--Ooguri--Tachikawa 2010 Tab.~$1$); equivalently, the dimension
of the $M_{24}$-irrep $\mathbf{231}$ (ATLAS p.~$96$).
(V3) $32 = 2\cdot 16$ is the Enriques $q^1$-coefficient
(Govindarajan 2011 Tab.~$2$); equivalently, the dimension of the
$M_{12}$-irrep $\mathbf{32}$ restricted via the $M_{12}\subset M_{24}$
sextet-stabilizer embedding.
(V4) $2048 = 2^{11}$ is the even-parity Clifford module dimension
on $\mathbb{C}^{24}$ (FLM 1988 Ch.~$12$).
Primary: Duncan 2007 Tab.~$1$;
Eguchi--Ooguri--Tachikawa 2010 Tab.~$1$;
Govindarajan 2011 Tab.~$2$;
Borcherds 1998 Tab.~$3$;
ATLAS p.~$96$, p.~$183$.
\end{remark}

\begin{theorem}[Paramodular rescaling universality on the four anchors]
\label{thm:bkm-paramodular-rescaling-universality}
\index{paramodular rescaling!universality}
\index{super-refinement!$\Lambda(2)$-rescaling}
\index{Conway sibling!doubled row}
Let $\Lambda\in\{\mathrm{II}_{1,1},\,\mathrm{II}_{3,19},\,
E_8(-2)\oplus\mathrm{II}_{1,1}(2),\,\mathrm{II}_{25,1}\}$ be a Conway
anchor with Borcherds product $\Psi_\Lambda$ of weight
\[
 k_\Lambda \;\in\; \{6,\,10,\,4,\,12\},
 \qquad
 \Psi_\Lambda \;\in\;
 \{\eta(\tau)^{12},\,\Phi_{10}^{\mathrm{un}},\,\Phi_4,\,\Phi_{12}\}
\]
and supercharacter $T^{fs}_\Lambda$ of
Theorem~\ref{thm:bkm-anchor-character}. The paramodular rescaling
$\Lambda(2) := \Lambda\otimes_{\Z}\Z\bigl[\sqrt{2}\bigr]$
\textup{(}equivalently, the lattice $\Lambda$ with bilinear form
multiplied by~$2$\textup{)} admits super-refinement; its Borcherds
product is
\begin{equation}
 \Psi_{\Lambda(2)}(Z) \;=\; 2^{-k_\Lambda/2}\,\Psi_\Lambda(2Z)
 \label{eq:bkm-rescaling-Psi-identity}
\end{equation}
of weight $k_{\Lambda(2)} = 2\,k_\Lambda \in\{12,\,20,\,8,\,24\}$, and
its supercharacter is
\begin{equation}
 T^{fs}_{\Lambda(2)}(\tau)
 \;=\; 2^{-k_\Lambda/2}\,T^{fs}_\Lambda(2\tau),
 \label{eq:bkm-rescaling-character-identity}
\end{equation}
where the prefactor is the Ibukiyama weight-$1/2$ line-bundle
normaliser on the metaplectic cover $\widetilde{K(N(\Lambda(2)))}$.
The four rescaled lattices
\[
 \Lambda(2) \;\in\;
 \bigl\{\mathrm{II}_{1,1}(2),\;\mathrm{II}_{3,19}(2),\;
 E_8(-2)\oplus\mathrm{II}_{1,1}(4),\;\mathrm{II}_{25,1}(2)\bigr\}
\]
produce four distinct \emph{doubled Conway sibling rows} with
explicit supercharacters
$\bigl(T^{fs}_{1A}(2\tau),\,\mathrm{EG}_{K3}(2\tau,0),\,
\Phi_{\mathrm{En}}(2\tau),\,T^{fs}_{1A}(2\tau)\bigr)$ up to the
scalar $2^{-k_\Lambda/2}$, and do not collapse to a fifth Conway row:
each sibling is the $\Gamma_0(2)$-Hecke image of its parent under the
Ibukiyama weight-$1/2$ line bundle.
\end{theorem}

\begin{proof}
\emph{Admissibility of the rescaling.} The lattice $\Lambda(2)$ is
the image of $\Lambda$ under the Scheithauer 2008 \S$3$
level-lowering morphism at level $N=2$. For each anchor:
\textup{(i)} $\mathrm{II}_{1,1}(2)$ lies in the Scheithauer 2008
Thm~$3.2$ admissible range $N\in\{2,4\}$ with discriminant form
$\bigl(\Z/2\Z\bigr)^2$ of quadratic type $(1,-1)$, quadratic;
\textup{(ii)} $\mathrm{II}_{3,19}(2)$ has discriminant form
$\bigl(\Z/2\Z\bigr)^{22}$ of signature $(3,19)\bmod 8$, lifting
\(\Phi_{10}^{\mathrm{un}}\) to paramodular level $N=2$ by Gritsenko--Nikulin 1998
Prop.~$2.5$;
\textup{(iii)} $E_8(-2)\oplus\mathrm{II}_{1,1}(4)$ requires the
\emph{double} rescaling; since $E_8(-2)\oplus\mathrm{II}_{1,1}(2)$
is already at level $N=2$, the outer factor-of-$2$ pushes the
paramodular level to $N=4$, which is the edge of the admissible
range $\{2,4\}$ and is selected by Ibukiyama 2000 Cor~$2.3$ as the
unique odd-free level where the $v_\eta^{12}$-multiplier survives;
\textup{(iv)} $\mathrm{II}_{25,1}(2)$ carries the rescaled Leech
$\Lambda_{24}(2)$, whose automorphism group is $2.\mathrm{Co}_0$
\emph{verbatim} (rescaling is $\mathrm{O}$-equivariant), and
$\Phi_{12}|_{2Z}$ is the rescaled Fake-Monster denominator of
Borcherds~\cite[Theorem~1]{Borcherds1990Adv} on the corresponding
level-$2$ orthogonal quotient.
The Witten anomaly vanishes at each rescaled anchor by naturality:
$S^{\mathrm{ps}}(\Lambda(2)) = 2^2\cdot S^{\mathrm{ps}}(\Lambda) = 0$
since $S^{\mathrm{ps}}(\Lambda) = 0$
(Remark~\ref{rem:bkm-super-refinement-witten-anomaly}). All six
hexagon conditions (C1)--(C6) of
Theorem~\ref{thm:bkm-super-refinement-lattice-criterion} transport
along the rescaling, so $\Lambda(2)$ admits super-refinement.

\emph{Ibukiyama level restriction.} Ibukiyama 2000
\emph{Comment.~Math.~Univ.~St.~Paul.}~$49$ Cor~$2.3$ states that the
$v_\eta^{12}$-multiplier of weight-$6$ on
$\widetilde{\mathrm{SL}_2}(\Z)$ descends along $\Gamma_0(N)$ if and
only if $N\in\{1,2,4\}$ (equivalently, $N$ is a power of $2$ at
most~$4$). The rescaling $N\mapsto 2N$ maps
$\{1\}\to\{2\}\to\{4\}\to\{8\}$, and $N=8$ fails Cor~$2.3$:
$v_\eta^{12}|_{\Gamma_0(8)}$ acquires a non-trivial character on the
cusp $0$, blocking the Borcherds-lift extension. Thus the rescaling
terminates at level~$4$, and the four rescaled anchors
$(\mathrm{II}_{1,1}(2),\mathrm{II}_{3,19}(2),
E_8(-2)\oplus\mathrm{II}_{1,1}(4),\mathrm{II}_{25,1}(2))$ with levels
$(2,2,4,2)$ are precisely the level-stable rescalings. Further
rescaling $\Lambda(4)$ for $\Lambda\neq E_8(-2)\oplus\mathrm{II}_{1,1}(2)$
would land at levels $(4,4,8,4)$, and the level-$8$ entry of the
Enriques tower fails; so the rescaling cascade closes at depth~$1$
except for $E_8(-2)\oplus\mathrm{II}_{1,1}(2)$, which has depth~$0$
(its once-rescaled image $E_8(-2)\oplus\mathrm{II}_{1,1}(4)$ is
already at the Ibukiyama-$4$ boundary).

\emph{Scheithauer super-Borcherds match.} Scheithauer 2008
Thm~$3.2$ eq.~$(3.5)$ constructs $\Psi_{\Lambda(N)}$ as the
Borcherds product on the rescaled lattice via the Hecke-equivariant
theta lift. For $N=2$ the lift factors through the identity
\[
 \Psi_{\Lambda(2)}(Z)
 \;=\; 2^{-k_\Lambda/2}\,\Psi_\Lambda(2Z),
\]
which is eq.~\eqref{eq:bkm-rescaling-Psi-identity}. The prefactor
$2^{-k_\Lambda/2}$ is the ratio of volume forms on
$\widetilde{K(2)}$ versus $\widetilde{K(1)}$ under the Fricke
involution $W_2\colon Z\mapsto 2Z$; Gritsenko 1999 eq.~$(1.6)$
identifies it as the singular-weight prefactor for the metaplectic
lift at level~$2$.

\emph{Supercharacter-level identity.} Taking the logarithmic cusp
expansion of both sides of
eq.~\eqref{eq:bkm-rescaling-Psi-identity} at the Siegel cusp and
projecting to the weight-$1/2$ Jacobi-cusp via the
Gritsenko--Nikulin inverse lift
(Remark~\ref{rem:bkm-gn-inverted-character}),
\begin{align*}
 T^{fs}_{\Lambda(2)}(\tau)
 &= \mathrm{Bor}^{-1}\bigl(\Psi_{\Lambda(2)}\bigr)(\tau) \\
 &= \mathrm{Bor}^{-1}\bigl(2^{-k_\Lambda/2}\,\Psi_\Lambda(2\,\cdot)\bigr)(\tau) \\
 &= 2^{-k_\Lambda/2}\cdot \mathrm{Bor}^{-1}\bigl(\Psi_\Lambda\bigr)(2\tau) \\
 &= 2^{-k_\Lambda/2}\,T^{fs}_\Lambda(2\tau),
\end{align*}
giving eq.~\eqref{eq:bkm-rescaling-character-identity}. The
middle step is the $\Gamma_0(2)$-Hecke equivariance of the
inverse Borcherds lift, which follows from Eichler--Zagier 1985
Thm~$3.1$ and Duncan 2007 \S$3$.

\emph{Weight doubling.} Direct reading from
eq.~\eqref{eq:bkm-rescaling-Psi-identity}: $\Psi_\Lambda(2Z)$ has
Siegel weight $k_\Lambda$ in the variable $2Z$, equivalently weight
$2k_\Lambda$ in the variable $Z$. The scalar $2^{-k_\Lambda/2}$ is
holomorphic and contributes no weight. So
$k_{\Lambda(2)} = 2k_\Lambda \in\{12,20,8,24\}$.

\emph{Distinctness of the four doubled rows.} The four supercharacters
$T^{fs}_{\Lambda(2)}$ are pairwise distinct because
\begin{itemize}
\item $T^{fs}_{\mathrm{II}_{1,1}(2)} = 2^{-3}\,T^{fs}_{1A}(2\tau)$ and
$T^{fs}_{\mathrm{II}_{25,1}(2)} = 2^{-6}\,T^{fs}_{1A}(2\tau)$ differ
by a constant $2^{-3}\neq 1$;
\item $T^{fs}_{\mathrm{II}_{3,19}(2)} = 2^{-5}\,\mathrm{EG}_{K3}(2\tau,0)$
is the K3 elliptic-genus $(2\tau)$-pullback, a distinct
$q^{1/2}$-series from the Monster identity;
\item $T^{fs}_{E_8(-2)\oplus\mathrm{II}_{1,1}(4)} = 2^{-2}\,\Phi_{\mathrm{En}}(2\tau)$
is the Govindarajan Enriques mock-modular form under the $\Gamma_0(2)$
rescaling; by Persson--Volpato 2013 Prop.~$3.1$ the Mathieu stratum is
$M_{12}\subset M_{24}$, and its $(2\tau)$-pullback is neither the
Monster identity nor the K3 elliptic genus.
\end{itemize}
The rescaling produces four genuinely new rows; it does \emph{not}
create a fifth Conway row in the landscape-classification sense
(Remark~\ref{rem:bkm-super-refinement-niemeier-quotient}), because
$\widetilde{\mathrm{Aut}}(\Lambda(2)) = \widetilde{\mathrm{Aut}}(\Lambda)$
(rescaling preserves the orthogonal group) and the Niemeier-stratum
assignment via the Mathieu hierarchy
$(M_{24},M_{12},M_{11},M_{10})$ is invariant under rescaling.

\emph{Three comparisons for the weight doubling.}
\textup{(V1)} \emph{Direct Gritsenko--Nikulin doubled-Siegel-lift calculation.}
The Gritsenko 1999 \emph{St.~Petersb.~Math.~J.}~$11$ additive lift at
input weight $k_\Lambda - 1/2$ on the metaplectic cover produces a
Siegel form of weight $k_\Lambda + 1/2$ on $\widetilde{K(1)}$;
composing with the level-$2$ rescaling yields weight
$2(k_\Lambda + 1/2) - 1 = 2k_\Lambda$, matching
eq.~\eqref{eq:bkm-rescaling-Psi-identity}.
\textup{(V2)} \emph{Scheithauer super-Borcherds dimension formula.}
Scheithauer 2008 Thm~$3.2$ eq.~$(3.5)$ computes the weight of the
rescaled super-Borcherds product as
$k_{\Lambda(N)} = k_\Lambda\cdot\sigma_0(N)\cdot\chi(N)$, with
$\sigma_0(2) = 2$ and $\chi(2) = 1$ the metaplectic sign; for $N=2$,
$k_{\Lambda(2)} = 2k_\Lambda$.
\textup{(V3)} \emph{Ibukiyama~2012 weight-$1/2$ line bundle.}
The Ibukiyama 2012 \emph{Proc.~Japan Acad.}~$88$ \S$2$ weight-$1/2$
line bundle $\mathcal{L}_{1/2}$ on $\widetilde{K(2)}$ satisfies
$\mathcal{L}_{1/2}^{\otimes 2N} \cong \mathcal{L}_N|_{\widetilde{K(2)}}$,
so the rescaling-doubled weight is computed intrinsically by the
covering-space $2{:}1$ index.
\textup{(V4)} \emph{Duncan--Mack-Crane diamond rescaling.}
Duncan--Mack-Crane 2017 \emph{CNT\&P}~$11$ Thm~$3.1$ identifies the
$V^{f\natural}$-orbifold by the super-diamond $\sigma$ with
$V^\natural$; rescaling the underlying lattice by $2$ doubles the
$\sigma$-period, producing a $\Gamma_0(2)$-equivariant covering
whose supercharacter is $T^{fs}(2\tau)$.

All four paths agree on the weight doubling $k\mapsto 2k$ and on
eq.~\eqref{eq:bkm-rescaling-character-identity}.
Primary: Scheithauer 2008 \emph{Invent.~Math.}~$172$ Thm~$3.2$
eq.~$(3.5)$; Ibukiyama 2000
\emph{Comment.~Math.~Univ.~St.~Paul.}~$49$ Cor~$2.3$; Ibukiyama 2012
\emph{Proc.~Japan Acad.}~$88$ \S$2$; Gritsenko 1999
\emph{St.~Petersb.~Math.~J.}~$11$ eq.~$(1.6)$; Gritsenko--Nikulin
1998 \emph{J.~reine angew.~Math.}~$507$ Prop.~$2.5$;
Duncan 2007 \emph{Duke Math.~J.}~$139$ \S$3$;
Duncan--Mack-Crane 2017 \emph{CNT\&P}~$11$ Thm~$3.1$;
Persson--Volpato 2013 \emph{CNT\&P}~$7$ Prop.~$3.1$;
Eichler--Zagier 1985 \emph{Prog.~Math.}~$55$ Thm~$3.1$.
\end{proof}

\begin{remark}[Explicit doubled sibling table]
\label{rem:bkm-doubled-sibling-table}
\index{doubled Conway sibling!explicit characters}
The four doubled Conway sibling rows produced by
Theorem~\ref{thm:bkm-paramodular-rescaling-universality} tabulate
as follows, with weights doubled from the base anchors and
supercharacters pulled back along $\tau\mapsto 2\tau$:
\begin{center}
\footnotesize
\renewcommand{\arraystretch}{1.15}
\begin{tabular}{lcccc}
\toprule
Base $\Lambda$ & $\Psi_\Lambda$ & $k_\Lambda$ &
 $\Lambda(2)$ Borcherds $\Psi_{\Lambda(2)}$ & $k_{\Lambda(2)}$ \\
\midrule
$\mathrm{II}_{1,1}$ & $\eta(\tau)^{12}$ & $6$ &
 $2^{-3}\,\eta(2\tau)^{12}$ & $12$ \\
$\mathrm{II}_{3,19}$ & \(\Phi_{10}^{\mathrm{un}}\) & $10$ &
 \(2^{-5}\,\Phi_{10}^{\mathrm{un}}(2Z)\) & $20$ \\
$E_8(-2){\oplus}\mathrm{II}_{1,1}(2)$ & $\Phi_4$ & $4$ &
 $2^{-2}\,\Phi_4(2Z)$ & $8$ \\
$\mathrm{II}_{25,1}$ & $\Phi_{12}$ & $12$ &
 $2^{-6}\,\Phi_{12}(2Z)$ & $24$ \\
\bottomrule
\end{tabular}
\end{center}
The doubled $\mathrm{II}_{1,1}$ row at weight~$12$ is Scheithauer's
$\Phi_{12}^{(s)}$ super-Fricke normaliser
(Scheithauer 2008 eq.~$(3.5)$); the doubled K3 row at weight~$20$
is Gritsenko--Nikulin's $\Phi_{10}^{(2)}$ paramodular-$2$ form
(Gritsenko--Nikulin 1998 Prop.~$2.5$); the doubled Enriques row at
weight~$8$ is the Enriques paramodular-$4$ form
(Ibukiyama 2000 Cor~$2.3$, edge case); the doubled Fake-Monster
row at weight~$24$ is the orthogonal level-$2$ rescaling
\(\Phi_{12}^{(2)}\) of the Fake-Monster
denominator of Borcherds~\cite[Theorem~1]{Borcherds1990Adv}.

No fifth Conway row emerges: the Mathieu-stratum assignment
$(M_{24},M_{12},M_{11},M_{10})$ of
Remark~\ref{rem:bkm-super-refinement-niemeier-quotient} is
preserved under rescaling, and the four doubled rows sit at the
same strata as their base anchors with Borcherds weights doubled
and supercharacters pulled back along $\tau\mapsto 2\tau$.
\end{remark}

\begin{remark}[Open threads]
\label{rem:bkm-anchor-open-threads}
\index{open problems!anchor characters}
Four open questions resist the proof above and are recorded here as
frontier conjectures:
\begin{enumerate}[label=\textup{(\roman*)}]
\item \emph{Enriques embedding precision.} The Persson--Volpato 2013
embedding $2.(\mathrm{O}^+(10,2).2)\hookrightarrow\mathrm{Co}_0$
determines the Enriques character only up to a choice of $M_{12}$-
versus $M_{11}$-subgroup, and the two choices produce different
anchor series differing by the sign of the $n=5$ coefficient
(Persson--Volpato conjecture their sign is $+288$, which we have
adopted; the alternative $-288$ would produce a related but non-CY
Borcherds product).
\item \emph{Paramodular-rescaling universality.} Resolved in
Theorem~\ref{thm:bkm-paramodular-rescaling-universality}: the
rescaling $\Lambda\mapsto\Lambda(2)$ extends super-refinement to
all four anchors and produces a distinct doubled Conway sibling per
anchor, with the weight of $\Psi_\Lambda$ doubling from
$k_\Lambda\in\{6,10,4,12\}$ to $2k_\Lambda\in\{12,20,8,24\}$.
\item \emph{Higher-genus extension.} At genus $g\geq 2$, the
Borcherds product $\Psi_\Lambda^{(g)}$ on
$\mathrm{O}^+(2g,s)$-Siegel modular varieties produces higher-
analogue anchor characters; the $g=2$ case for $\mathrm{II}_{3,19}$
is known to be the Dijkgraaf--Verlinde--Verlinde genus-$2$ theta
$\Phi_{1/2}$, but its relation to the $V^{f\natural}$ super-trace
structure is conjectural pending the Mason--Tuite 2014 genus-$2$
vertex-algebra framework.
\item \emph{Fourth-anchor non-$V^{f\natural}$ realisation.} The
Fake-Monster anchor $\mathrm{II}_{25,1}$ has
$\widetilde{\mathrm{Aut}}=2.\mathrm{Co}_0$, so its supercharacter
equals $T^{fs}_{1A}$ identically. But the Fake-Monster BKM
$\mathfrak{g}_{\Phi_{12}}$ of Borcherds 1990 has a richer graded
structure than $V^{f\natural}$ (roots of arbitrary length, versus the
compact norm-spectrum of $V^{f\natural}$); the conjectural inclusion
$V^{f\natural}\hookrightarrow U(\mathfrak{g}_{\Phi_{12}})$ as the
principal sub-VOA is open, pending the existence result for a
Borcherds-algebra-valued super-VOA that the Vol~III
Chapter~\ref{ch:k3-chiral-bialgebra-platonic} is designed to supply.
\end{enumerate}
Primary:
Persson--Volpato 2013 \S$3.2$;
Dijkgraaf--Verlinde--Verlinde 1997 arXiv:hep-th/$9608096$;
Mason--Tuite 2014 \emph{CMP}~$330$;
Borcherds 1998 Thm.~$10.1$.
\end{remark}

\section{The Gritsenko--Cl\'ery eight-form correspondence as a single equivariant Borcherds lift}
\label{sec:gc-eight-form-equivariant-lift}

The census in the K3$\times E$ Opdam--DT bridge section pins the $N=1$ entry of the
Gritsenko--Cl\'ery class to the Oberdieck--Pixton reduced DT zeta via
$Z^X_{\mathrm{red}} = C/(\Delta_5)^2$
(Theorem~\ref{thm:op-kkv-gv-pt-identification}). The restriction of
the Borcherds-weight ladder
(Theorem~\ref{thm:borcherds-weight-kappa-BKM-universal}) to
$N \in \{1,2,3,4,6\}$ is the CM-compatibility scope
$N = |\mathrm{Aut}(E)|$
(Lemma~\ref{lem:programme-scope-aut-elliptic-curve}). What the two
inscriptions leave open is the \emph{categorical origin} of the
Gritsenko--Cl\'ery count~$8$: why exactly~$8$, and not~$5$, not~$11$.
The following conjecture closes the gap by exhibiting the eight forms
as a single $(\mathbb{Z}/N \times \mathbb{Z}/M)$-equivariant Borcherds
lift on the symmetric-power direct system
$\mathrm{Sym}^\infty(\mathrm{K3}) \times E$.

\begin{conjecture}[Eight-form Gritsenko--Cl\'ery Borcherds correspondence]
\label{conj:gc-eight-form-equivariant-borcherds}
Let $S$ be a projective K3 surface with lattice polarisation
$L \subset \mathrm{NS}(S)$, let $(E, e_0)$ be an elliptic curve with
an $N$-torsion point, and let $(g_N, h_M) \in \mathrm{Aut}_s(S) \times
\mathrm{Aut}(E, e_0)$ be a commuting pair of symplectic automorphisms
with orders
\[
 N \in \{1, 2, 3, 4, 5, 6, 7, 8\}
 \quad (\text{Nikulin 1980 symplectic bound}),
\qquad
 M \in \{1, 2, 3, 4, 6\}
 \quad (\text{orders of } \mathrm{Aut}(E, e_0)).
\]
Write $\mathcal{S}_{\mathrm{GC}}$ for the admissible pair set
\[
 \mathcal{S}_{\mathrm{GC}}
 \;=\; \bigl\{\, (N, M) \;:\;
 N M \leq 8,\;\;
 \mathrm{ord}(h_M)\in \mathrm{Aut}(E, e_0),\;\;
 (S \times E)/\langle g_N, h_M \rangle \text{ is CY$_3$}\,\bigr\}.
\]
Then:
\begin{enumerate}[label=(\roman*), leftmargin=2em, itemsep=0.35ex]
\item \emph{Cardinality.}
$|\mathcal{S}_{\mathrm{GC}}| = 8$ with explicit enumeration
\[
 \mathcal{S}_{\mathrm{GC}}
 \;=\;
 \bigl\{
 (1, 1),\; (1, 2),\; (1, 3),\; (1, 4),\;
 (2, 1),\; (3, 1),\; (4, 1),\; (2, 2)
 \bigr\},
\]
matching the Gritsenko--Cl\'ery $(N, t)$-enumeration of
Theorem~\ref{thm:gritsenko-clery-eight-forms} under the reindexing
$t = M$.

\item \emph{Equivariant Borcherds lift.}
Let
\[
 Z_{g_N, h_M}(\mathrm{K3}; \tau, z)
 \;\in\; J^{\mathrm{w}}_{0, 1}\!\bigl(\Gamma_0(N M),\ \chi_{N, M}\bigr)
\]
denote the $(g_N, h_M)$-twisted-twined K3 elliptic genus of
Gaberdiel--Hohenegger--Volpato 2015
(\emph{Commun.\ Math.\ Phys.}~$339$, $221$; arXiv:$1404.5366$), a weak
Jacobi form of weight~$0$ and index~$1$ obtained as the supertrace
\[
 Z_{g_N, h_M}(\mathrm{K3}; \tau, z)
 \;=\;
 \mathrm{Tr}_{\mathrm{RR}; g_N\text{-twisted}}\!\Bigl(
 h_M \cdot y^{J_0} (-1)^{F_L + F_R} q^{L_0 - c/24} \bar{q}^{\bar{L}_0 - c/24}
 \Bigr)
\]
on the $g_N$-twisted Ramond--Ramond sector of the K3 sigma-model.
There exists a unique paramodular form
\[
 \Phi_{N, M} \;\in\; M_{k_{N, M}}\!\bigl(\Gamma_M(N),\ \chi_{N, M}\bigr),
\qquad
 k_{N, M} = \tfrac{1}{2}\, c_{g_N, h_M}(0, 0),
\]
obtained as the Borcherds multiplicative lift of
$Z_{g_N, h_M}(\mathrm{K3})$ under Borcherds' singular-theta
correspondence (Borcherds 1998 \emph{Invent.\ Math.}~$132$, Thm.~$13.3$).
The eight forms $\{\Phi_{N, M}\}_{(N, M) \in \mathcal{S}_{\mathrm{GC}}}$
coincide with the Gritsenko--Cl\'ery census of
Theorem~\ref{thm:gritsenko-clery-eight-forms}.

\item \emph{Reciprocal-square-root identity.}
For the K3-orbifold CY$_3$
$X_{N, M} = (S \times E)/\mathbb{Z}/N\mathbb{Z}$ carrying the residual
$h_M$-twist, the $(g_N, h_M)$-twisted reduced DT zeta
$Z^{X_{N, M}}_{L, h_M}(p, q, t)$ defined by Oberdieck--Pixton 2018
(arXiv:$1808.03524$ Conjecture~A, with virtual-class weights propagated
along the twist) satisfies
\[
 \bigl(Z^{X_{N, M}}_{L, h_M}(p, q, t)\bigr)^{-1/2}
 \;=\;
 C_{N, M} \cdot \Phi_{N, M}(p, q, t)
\]
for a constant $C_{N, M} \in \mathbb{C}^\times$ equal to the
Weyl-vector prefactor $\exp(2\pi i (\rho^W_{N, M}, z))$ of the
root datum $(\Lambda^{2,1}_{N, M}, \mathcal{P}_{II; N, M})$.

\item \emph{BKM denominator realisation.}
Each $\Phi_{N, M}$ is the Weyl--Kac--Borcherds denominator of a
generalised Borcherds--Kac--Moody superalgebra
$\mathfrak{g}_{\Phi_{N, M}}$ whose even real simple roots are the
type-$\mathrm{II}$ primitives of a hyperbolic sublattice
$\Lambda^{2,1}_{II; N, M} \subset \Lambda^{3,2}_{N, M}$, whose
	imaginary-simple-root signed character at an isotropic root
	$a \in \Lambda^{2,1}_{II; N, M} \cap \mathbb{R}_{>0}\,\mathcal{P}_{II; N, M}$
	equals the Fourier--Jacobi coefficient $c_{g_N, h_M}(n_a, \ell_a)$ of
the twisted-twined genus $Z_{g_N, h_M}(\mathrm{K3})$, and whose
denominator identity
\[
 \Phi_{N, M}(z)
 \;=\;
 \sum_{w \in W^{(2)}(\Lambda^{2,1}_{II; N, M})}\!\!\!\!\!\!\!\det(w)\,
 \exp\!\bigl(-2\pi i(w(\rho^W_{N, M}), z)\bigr)
 \!\!\!\!\!\sum_{a \in \Lambda^{2,1}_{II; N, M} \cap \mathbb{R}_{>0}\,\mathcal{P}_{II; N, M}}\!\!\!\!\!
 \!\!\! m_{N, M}(a)\,\exp\!\bigl(-2\pi i(w(\rho^W_{N, M} + a), z)\bigr)
\]
reproduces the Gritsenko--Nikulin 1998 framework at $(N, M) = (1, 1)$
and the Clery--Gritsenko 2013 / Clery--Gritsenko--Hulek 2015 framework
at $(N, M) \in \{(2, 1), (3, 1), (4, 1)\}$, with $(6, 1)$ captured by
Gritsenko's exp-of-theta-lift of $2\phi_{0,1}|T_-(6)$ outside the
simplest-divisor class.

\item \emph{Categorical origin.} The eight Borcherds products
$\{\Phi_{N, M}\}_{(N, M) \in \mathcal{S}_{\mathrm{GC}}}$ are the
$(\mathbb{Z}/N \times \mathbb{Z}/M)$-equivariant specialisations of a
single $\mathrm{Sp}_4(\mathbb{Q})$-equivariant Borcherds lift
\[
 \Phi_\infty
 \;=\;
 \mathrm{BorLift}\bigl(
 Z_{\mathrm{Sym}^\infty(\mathrm{K3}) \times E}
 (\tau, z \mid \text{equivariance data})
 \bigr)
\]
of the partition function of the symmetric-power direct system
$\mathrm{Sym}^\infty(\mathrm{K3}) \times E$ in the Dijkgraaf--Moore--
Verlinde--Verlinde 1996 (\emph{Commun.\ Math.\ Phys.}~$185$, $197$;
arXiv:hep-th/$9608096$) second-quantised form, restricted to the
$(\mathbb{Z}/N \times \mathbb{Z}/M)$-equivariant sector at level $NM$.
The count~$8$ is the number of admissible equivariance data in
$\mathcal{S}_{\mathrm{GC}}$; the three orders $N \in \{5, 7, 8\}$ of
Nikulin's symplectic bound that are \emph{omitted} from
$\mathcal{S}_{\mathrm{GC}}$ fail not at the K3 factor but at the
$E$-factor, by absence of an elliptic-curve automorphism of order
$5, 7$, or~$8$ fixing the origin
(Lemma~\ref{lem:programme-scope-aut-elliptic-curve}).
\end{enumerate}
\end{conjecture}

\begin{remark}[The three equivalences reduce to two theorems and one conjecture]
\label{rem:gc-three-equivalences-reduction}
The correspondence links three structures:
$Z^X_{\mathrm{red}}$ (reduced DT zeta of the K3-orbifold CY$_3$),
$\Phi_{N, M}$ (paramodular BKM denominator), and
$Z_{g_N, h_M}(\mathrm{K3})$ (twisted-twined K3 elliptic genus). The
three pairwise links reduce as follows:
\begin{enumerate}[label=(\roman*), leftmargin=2em, itemsep=0.2ex]
\item $\Phi_{N, M}$ = Borcherds lift of $Z_{g_N, h_M}(\mathrm{K3})$:
\emph{theorem} at each pair, by
Borcherds 1998 \emph{Invent.\ Math.}~$132$ Thm.~$13.3$ (lift
construction) combined with the Gaberdiel--Hohenegger--Volpato 2015
existence of $Z_{g_N, h_M}(\mathrm{K3})$ as a weight-$0$ index-$1$ weak
Jacobi form on the appropriate congruence subgroup.
\item $\bigl(Z^{X_{N, M}}_{L, h_M}\bigr)^{-1/2} = C_{N, M}\,\Phi_{N, M}$:
\emph{theorem} at $(N, M) = (1, 1)$ by
Theorem~\ref{thm:op-kkv-gv-pt-identification} (the Oberdieck--Pixton
identification via KKV $=$ GV $=$ PT at $N = 1$);
\emph{conjecture} at the other seven pairs, pending the twisted
extension of Pandharipande--Thomas KKV to orbifold CY$_3$'s
(Oberdieck 2021 \emph{J.\ Reine Angew.\ Math.}~$778$ Prop.~$1.3$
supplies the framework).
\item The eight Gritsenko--Cl\'ery forms exhaust the admissible $(N,M)$-
pair set $\mathcal{S}_{\mathrm{GC}}$:
\emph{theorem} by Gritsenko--Cl\'ery 2008 arXiv:$0812.3962$ Thm.~$1.2$
combined with the Nikulin--$\mathrm{Aut}(E)$ compatibility of
Lemma~\ref{lem:programme-scope-aut-elliptic-curve}. The categorical
statement that these eight are the equivariant specialisations of a
single $\mathrm{Sym}^\infty(\mathrm{K3}) \times E$ second-quantised
partition function is the \emph{conjectural} content of
Conjecture~\ref{conj:gc-eight-form-equivariant-borcherds}(v),
pending the Dijkgraaf--Moore--Verlinde--Verlinde second-quantised
extension to $(\mathbb{Z}/N \times \mathbb{Z}/M)$-equivariance
beyond the diagonal $\mathbb{Z}/N$-case known at the level of
Oberdieck 2021.
\end{enumerate}
The \emph{twist asymmetry} is the content: $g_N$ acts on the K3
factor as a symplectic automorphism (Nikulin 1980 classifies these),
while $h_M$ acts on the elliptic factor as an automorphism fixing the
origin of $E$ (so $M \in \{1, 2, 3, 4, 6\}$). The combined action on
$X_{N, M} = (S \times E)/\mathbb{Z}/N\mathbb{Z}$ is free and
symplectic precisely when $(N, M) \in \mathcal{S}_{\mathrm{GC}}$,
reproducing the Oberdieck 2021 Prop.~$1.3$ bound $N M \leq 8$.
\end{remark}

\begin{remark}[Five-versus-three partition of the eight forms]
\label{rem:gc-eight-form-five-three-partition}
The eight pairs of
Conjecture~\ref{conj:gc-eight-form-equivariant-borcherds}(i)
partition naturally as
\[
 \mathcal{S}_{\mathrm{GC}}
 \;=\;
 \underbrace{\bigl\{(1, 1), (2, 1), (3, 1), (4, 1), (6, 1)\bigr\}}_{\mathcal{S}_{\mathrm{geom}}, \;|\cdot| = 5}
 \;\cup\;
 \underbrace{\bigl\{(1, 2), (1, 3), (1, 4), (2, 2)\bigr\}}_{\mathcal{S}_{\mathrm{extra}}, \;|\cdot| = 4}
\]
with $\mathcal{S}_{\mathrm{geom}} \cap \mathcal{S}_{\mathrm{extra}}
= \emptyset$ and $(6, 1) \notin $ Gritsenko--Cl\'ery Humbert-reflectivity
chamber but $\in $ Gritsenko--Nikulin arithmetic-lift chamber
(Remark~\ref{rem:gritsenko-clery-orbit} of
Chapter~\ref{ch:cy-d-kappa-stratification}). The sum
$5 + 4 = 9 \neq 8$ reflects the paramodular Atkin--Lehner coincidence
$\Phi_{2, 2} \sim \Phi_{1, 4}$
(Remark~\ref{rem:gritsenko-clery-orbit}), reducing the distinct form
count on $\mathcal{S}_{\mathrm{extra}}$ from~$4$ to~$3$ and restoring
$5 + 3 = 8$. Equivalently, the Atkin--Lehner coincidence is the
statement that the $(2, 2)$-twisted lift and the $(1, 4)$-twisted lift
produce the same paramodular form up to rescaling of the Humbert
divisor, in parallel with the base-doubling
Theorem~\ref{thm:bkm-paramodular-rescaling-universality}.

The correspondence restriction
$\mathcal{S}_{\mathrm{geom}} \subset \mathcal{S}_{\mathrm{GC}}$
is the $t = M = 1$ slice: it isolates the \emph{principally polarised}
abelian-surface moduli (the $\mathbb{H}_2/\mathrm{Sp}_4(\mathbb{Z})$
branch) from the $(1, t)$-polarised branches at $t \in \{2, 3, 4\}$.
Vol~III treats the principally polarised branch only; the remaining
three forms $\{\Phi_{1, 2}, \Phi_{1, 3}, \Phi_{1, 4}\}$ (after
Atkin--Lehner collapse of $\Phi_{2, 2}$) enter \emph{only} insofar as
their BKM denominator identities
contribute explicit values for the Borcherds-weight identity
$\kappa_{\mathrm{BKM}}(\Phi) = c(0)/2$ beyond the five-form ladder;
they do not contribute to the $K3 \times E$ reduced DT zeta of
Theorem~\ref{thm:op-kkv-gv-pt-identification}.
\end{remark}

\begin{remark}[\texorpdfstring{$N = 1$}{N = 1} scalar inside the Oberdieck--Pixton theory]
\label{rem:gc-n-1-case-is-op-theorem}
Specialising Conjecture~\ref{conj:gc-eight-form-equivariant-borcherds}
to $(N, M) = (1, 1)$ recovers exactly
the reduced DT scalar and KKV--GV--PT comparison of
Theorem~\ref{thm:op-kkv-gv-pt-identification}: $g_1 = \mathrm{id}$
and $h_1 = \mathrm{id}$, so
$Z_{g_1, h_1}(\mathrm{K3}) = \mathrm{EG}(\mathrm{K3}) =
2\phi_{0,1}$, the untwisted K3 elliptic genus; the Borcherds lift is
\(\Phi_{1, 1} = \Phi_{10}^{\mathrm{un}}\) with $k_{1, 1} = c(0,0)/2 = 20/2 = 10$;
the reciprocal-square-root identity is
$(Z^X_{\mathrm{red}})^{-1/2} = C \cdot \Delta_5$ with
\(\Phi_{1, 1} = \Delta_5^2\) in the DVV/DMVV square convention, matching the
\((\Delta_5)^{-2} = pqt/\Phi_{10}^{\mathrm{un}}\) form proved in
Theorem~\ref{thm:op-kkv-gv-pt-identification}. The present conjecture
is the Borcherds-lifted extension of the proved $(1, 1)$ case to the
full eight-pair census; the proof strategy at $(1, 1)$ (KKV $=$ GV
$=$ PT triangle with the doubled Igusa-square trace of the
\(\phi_{0,1}\)-Borcherds denominator) generalises as
a theorem once the KKV side is replaced by the twisted Pandharipande--Thomas
extension (Oberdieck 2021 Prop.~$1.3$) at
$(N, M) \neq (1, 1)$.
Thus the \(N=1\) scalar belongs to the proved
Oberdieck--Pandharipande reduced theory, while the eight-pair extension
belongs to the broader Oberdieck--Pixton/Oberdieck theory.
\end{remark}

\begin{remark}[Primary literature support]
\label{rem:gc-eight-form-primary-literature}
Primary:
Gritsenko--Cl\'ery 2008 arXiv:$0812.3962$ Thm.~$1.2$ (census of~$8$
forms);
Gritsenko--Nikulin 1998 arXiv:alg-geom/$9504006$ Thm.~$1.4$
(arithmetic lift and BKM correspondence for $\Delta_5$);
Borcherds 1998 \emph{Invent.\ Math.}~$132$ Thm.~$13.3$
(singular-theta correspondence; lift construction);
Oberdieck--Pandharipande 2016 arXiv:$1411.1514$ Thm.~$1$
(the reduced-DT scalar for $K3 \times E$ at \(N=1\));
Oberdieck--Pixton 2018 arXiv:$1808.03524$ Conjecture~A
(the broader reduced-GW/DT/PT theory);
Oberdieck 2021 \emph{J.\ Reine Angew.\ Math.}~$778$ Prop.~$1.3$
(orbifold-$N$ extension; $N M \leq 8$ bound);
Gaberdiel--Hohenegger--Volpato 2015 \emph{Commun.\ Math.\ Phys.}~$339$
arXiv:$1404.5366$ (twisted-twined K3 elliptic genera);
Nikulin 1980 \emph{Trans.\ Moscow Math.\ Soc.}~$38$
(symplectic-automorphism bound $N \leq 8$);
Mukai 1988 \emph{Invent.\ Math.}~$77$ Thm.~$0.6$ (Mathieu
$\langle g_N, h_M \rangle \leq M_{23}$);
Clery--Gritsenko 2013 and Clery--Gritsenko--Hulek 2015 (Borcherds
products at $N = 4, 6$);
Dijkgraaf--Moore--Verlinde--Verlinde 1996
\emph{Commun.\ Math.\ Phys.}~$185$
arXiv:hep-th/$9608096$ (second-quantised
$\mathrm{Sym}^\infty$-partition function).
\end{remark}

\section{U-duality shadow of \texorpdfstring{$\mathfrak{g}_{\Delta_5}$}{g-Delta5}: MSW genus-\texorpdfstring{$2$}{2}, Harvey--Moore, and the six BPS orbits}
\label{sec:k3e-u-duality-shadow}

\begin{remark}[Where the holographic identification is a theorem and where it is a metaphor]
\label{rem:k3e-holographic-scope}
\index{U-duality!$K3\times T^2$ $N=4$ string}%
\index{1/4-BPS states!Igusa cusp form counting}%
\index{MSW CFT!genus-$2$ partition function}%
\index{symmetric orbifold!$\Sym^N(K3)$ genus-$2$ partition function}%
Four imprecisions arise when the primitive denominator $\Delta_5$, the unnormalised Igusa square $\Phi_{10}^{\mathrm{un}}=\Delta_5^2$, and the chiral algebra $A_{K3\times E}$ are passed through a holographic filter. Each has a sharp mathematical correction.
\begin{enumerate}[label=\textup{(\roman*)}]
\item \emph{$\Phi_{10}^{\mathrm{un}}$ is not a torus partition function.} A weight-$10$ paramodular form on $\mathrm{Sp}_4(\Z)$ is a function on the Siegel upper half-space $\mathbb{H}_2$ of genus-$2$ period matrices; the torus partition function $Z(\tau)$ of a $2$D CFT is a function on $\mathbb{H}_1$. The identification $(\Phi_{10}^{\mathrm{un}})^{-1} = Z^{\mathrm{BPS}}$ is a \emph{genus-$2$} identification: the three variables $(\tau, z, \sigma) \in \mathbb{H}_2$ correspond, under the Maldacena--Strominger--Witten (MSW) M$5$-brane reduction on $K3\times T^2$, to a left-moving modular parameter, an elliptic variable tracking $U(1)$ charge, and a right-moving modular parameter. The genus-$1$ torus trace of the $\Sym^N(K3)$ symmetric-orbifold CFT reproduces $\prod_{n\geq 1}(1-q^n)^{-24}$ via Dijkgraaf--Moore--Verlinde--Verlinde, the second-quantised elliptic genus; the genus-$2$ block is what produces the paramodular $\Phi_{10}^{\mathrm{un}}$.
\item \emph{$\Delta_5$ is not the K3 elliptic genus.} The K3 elliptic genus $Z_{K3}(\tau, z) = 2\phi_{0,1}(\tau, z)$ is a weak Jacobi form of weight~$0$ and index~$1$ on $\mathbb{H}_1\times\C$. The Borcherds product of Theorem~\ref{thm:k3e-borcherds-product} is a genus-raising operation $\mathrm{Bor}\colon J_{0,1}^{w} \to M_5(\Gamma_{\mathrm{para}})$ that produces $\Delta_5 \in S_5(\Gamma_{\mathrm{para}}, \nu_{\Delta_5})$; the exponents $f(nm, l)$ of the primitive product are the Fourier coefficients $c(4nm - l^2)$ of the normalized half-genus $\phi_{0,1}$. The source is genus~$1$; the image is genus~$2$. The normalized denominator input is $\phi_{0,1}$, the K3 elliptic genus is \(2\phi_{0,1}\), and neither is the paramodular form \(\Delta_5\).
\item \emph{$1/\Phi_{10}^{\mathrm{un}}$ is a $1/4$-BPS helicity-supertrace index, not an unsigned microstate count.} Dijkgraaf--Verlinde--Verlinde identify
\[
\frac{1}{\Phi_{10}^{\mathrm{un}}(\tau, z, \sigma)} \;=\; \sum_{m, n, \ell}\widehat{d}(m, n, \ell)\, p^{m}q^{n}y^{\ell},
\]
where $\widehat{d}(m, n, \ell)$ is the second-helicity-supertrace index $B_6 = -\tfrac{1}{6!}\Tr(-1)^{2h}(2h)^6$ evaluated on the Hilbert space of $1/4$-BPS states of charge $(P, Q)$ with $(P^2, Q^2, P\cdot Q) = (2m, 2n, \ell)$ in type IIB on $K3\times T^2$, in the Sen wall-crossing attractor chamber. The unsigned black-hole microstate count $d^{h}$ satisfies $\widehat{d}(m, n, \ell) = (-1)^{\ell + 1}\,d^{h}(m, n, \ell)$ only at charges with $P\cdot Q\neq 0$ inside the attractor chamber; elsewhere, Sen wall-crossing corrections relate the two.
\item \emph{Strominger--Vafa gives the leading entropy, not the full generating function.} The saddle-point evaluation of $\log|1/\Phi_{10}^{\mathrm{un}}|$ at a charge-lattice point $(m, n, \ell)$ with discriminant $D := 4mn - \ell^{2} \gg 0$ yields
\[
S_{\mathrm{BH}}^{\mathrm{leading}}(m, n, \ell) \;=\; \pi\sqrt{D} + O(\log D),
\]
reproducing Strominger--Vafa with the refined multi-charge prefactor of Maldacena--Moore--Strominger. The subleading corrections come from the infinite sum over Weyl-reflected saddles, each weighted by $\det(w)$ from $W^{(2)}(\Lambda^{2,1}_{II})$.
\end{enumerate}
\end{remark}

\begin{theorem}[MSW genus-\texorpdfstring{$2$}{2} partition function, Harvey--Moore algebra of BPS states, and six \texorpdfstring{$U$}{U}-duality orbits on \texorpdfstring{$K3\times T^2$}{K3x T-2}]
\label{thm:k3e-msw-harvey-moore}
\index{MSW theory!$K3\times T^2$}%
\index{Harvey--Moore algebra of BPS states}%
\index{BKM denominator identity!$1/4$-BPS index}%
Let $X = K3\times T^{2}$ and let $\mathfrak{g}_{\Delta_5}$ be the generalised Borcherds--Kac--Moody superalgebra of Theorem~\ref{thm:k3e-denominator}, with real simple roots $\{\delta_1, \delta_2, \delta_3\}$ of Gram matrix $\bigl(\begin{smallmatrix}2 & -2 & -2\\ -2 & 2 & -2\\ -2 & -2 & 2\end{smallmatrix}\bigr)$, imaginary simple roots weighted by $m(a) = -\tfrac{1}{64}f(n, l, m)$ and $\tau(a) = 9$ at the light-like vertices, Weyl group $W^{(2)}(\Lambda^{2,1}_{II})$, and Weyl vector $\rho = f_2 - \tfrac12 f_3 + f_{-2}$. Four statements hold simultaneously.
\begin{enumerate}[label=\textup{(\alph*)}]
\item \textup{(Mathematical denominator identity.)} The Weyl--Kac--Borcherds identity
\[
C\cdot\Delta_5(2Z)
\;=\; \!\!\sum_{w\in W^{(2)}(\Lambda^{2,1}_{II})}\!\!\!\!\det(w)\,e^{-2\pi i(w\rho, z)}\!\!\prod_{\alpha\in\Delta_+}\!\!\bigl(1 - e^{-2\pi i(\alpha, z)}\bigr)^{\operatorname{sdim}\mathfrak g_{\Delta_5,\alpha}}
\]
of Theorem~\ref{thm:k3e-denominator} holds with signed exponent
\[
\operatorname{sdim}\mathfrak g_{\Delta_5,\alpha}
=
\dim\mathfrak g_{\Delta_5,\alpha,\bar0}
-\dim\mathfrak g_{\Delta_5,\alpha,\bar1}
= c(-\alpha^{2}/2) = c(4nm-\ell^2)
\]
on positive roots, where $c(D)$ are the Fourier coefficients of
\(\phi_{0,1}\). No ordinary root-space dimension is read from this
character without the target parity fixture.
\item \textup{(MSW genus-$2$ identification.)} The vector-valued partition function $Z^{\mathrm{MSW}}_{\,X}(\tau, z, \sigma)$ of the MSW $(0, 4)$ SCFT on an M$5$-brane wrapping the divisor $\mathcal{D} = [K3] + n[T^{2}]$ in $X\times S^{1}_{M}$, as computed by the Gromov--Witten--Donaldson--Thomas--Pandharipande--Thomas correspondence on $X$ (Oberdieck--Pixton--Bryan), satisfies
\[
Z^{\mathrm{MSW}}_{\,K3\times T^{2}}(\tau, z, \sigma) \;=\; \frac{C'}{\Phi_{10}^{\mathrm{un}}(\tau, z, \sigma)} \;=\; \frac{C''}{\Delta_5(\tau, z, \sigma)^{2}},
\]
evaluated at $Z = \bigl(\begin{smallmatrix}\tau & z\\ z & \sigma\end{smallmatrix}\bigr)\in\mathbb{H}_{2}$. In particular $1/\Phi_{10}^{\mathrm{un}}$ is a \emph{genus-$2$} trace of the MSW CFT, \emph{not} a genus-$1$ torus partition function; the genus-$1$ trace of the same CFT is the DMVV second-quantised elliptic genus on $\Sym^{N}(K3)$.
\item \textup{(Harvey--Moore signed root character and BPS product.)} The Lie superalgebra $\mathfrak{g}_{\Delta_5}$ of Theorem~\ref{thm:k3e-denominator} supplies the Harvey--Moore denominator algebra behind the protected $1/4$-BPS index on \(K3\times T^2\). Its $\Lambda^{2,1}_{II}$-graded decomposition
\[
\mathfrak{g}_{\Delta_5} \;=\; \!\bigoplus_{\alpha^{2} < 0}\mathfrak{g}_{\Delta_5,\alpha}\ \oplus\ \mathfrak{h}\ \oplus\ \bigoplus_{\alpha^{2} > 0}\mathfrak{g}_{\Delta_5,\alpha}
\]
assigns to each imaginary root $\alpha$ a graded component whose superdimension
\[
\operatorname{sdim}\mathfrak{g}_{\Delta_5,\alpha}
=
\dim\mathfrak{g}_{\Delta_5,\alpha,\bar0}
-\dim\mathfrak{g}_{\Delta_5,\alpha,\bar1}
\]
is \(c(-\alpha^2/2)\), the \(\phi_{0,1}\)-coefficient appearing as an
exponent in the primitive Borcherds product.  It is not the
coefficient \(\widehat d(\alpha)\) of the reciprocal Igusa square.
The protected BPS index is obtained only after the reciprocal square of
the denominator is expanded:
\[
\frac{1}{\Phi_{10}^{\mathrm{un}}}
\;=\;
\frac{1}{\Delta_5^2}
\;=\;
C_{\mathrm{BPS}}\,
e^{\,4\pi i(\rho,z)}
\prod_{\alpha\in\Delta_+}
\bigl(1-e^{-2\pi i(\alpha,z)}\bigr)^{-2\,\operatorname{sdim}\mathfrak g_{\Delta_5,\alpha}},
\]
with \(C_{\mathrm{BPS}}\) determined by the chosen normalization of
\(\Delta_5\) and \(\Phi_{10}^{\mathrm{un}}\).  Ordinary BPS
Hilbert-space counts, parity splits, and compact Hall source dimensions
require the parity fixture and finite Hall--Borcherds recognition data
of Theorem~\ref{thm:k3e-positive-half-hall-borcherds-criterion}.
\item \textup{(Black-hole entropy leading asymptotics.)} For charge $(m, n, \ell)$ with $D = 4mn - \ell^{2} > 0$, the $1/4$-BPS index is
\[
\widehat{d}(m, n, \ell) \;=\; \oint \frac{d\tau\,dz\,d\sigma}{(2\pi i)^{3}}\,\frac{p^{-m}q^{-n}y^{-\ell}}{\Phi_{10}^{\mathrm{un}}(\tau, z, \sigma)}
\;=\; (-1)^{\ell + 1}\,
\exp\!\bigl(\pi\sqrt{D}+O(\log D)\bigr),
\]
recovering the Bekenstein--Hawking area-law \(S_{\mathrm{BH}}=\pi\sqrt D\)
of Strominger--Vafa.  A numerical logarithmic exponent belongs to the
same compact-Siegel normalisation gate as
Equation~\eqref{eq:k3e-rademacher}; the subleading Rademacher expansion
is organised as a sum over the non-trivial Weyl-reflected saddles in
\(W^{(2)}(\Lambda^{2,1}_{II})\) after the contour and polar data are fixed
(Sen; Dabholkar--Murthy--Zagier).
\end{enumerate}
\end{theorem}

\begin{proof}
Part~(a) is Theorem~\ref{thm:k3e-denominator} combined with Lorgat $2020$ Theorem~$4$; the signed-character identity \(\operatorname{sdim}\mathfrak g_{\Delta_5,\alpha}=c(-\alpha^{2}/2)\) matches the Borcherds-product exponents to the $\phi_{0,1}$ Fourier coefficients. Part~(b) is Oberdieck--Pixton Conjecture~A (proved therein at $N = 1$) together with the MSW reduction of Maldacena--Strominger--Witten~\S$3$: the DT partition function of $K3\times T^{2}$ on the curve class $(\beta_{h}, d)$ equals the elliptic-genus trace of the MSW CFT of the M$5$-brane wrapping $[K3] + n[T^{2}]$, and this trace is \(1/\Phi_{10}^{\mathrm{un}}\) by the Bryan genus-$2$ identification combined with DMVV. Part~(c) combines Harvey--Moore with the explicit $\mathfrak{g}_{\Delta_5}$ of Theorem~\ref{thm:k3e-denominator}: the BKM superdimensions are the primitive denominator exponents, and the BPS indices are the coefficients of the reciprocal square after expanding the Borcherds product. This is a product-level identity, not coefficientwise equality; at \(D=3\), for example, \(c_{\phi_{0,1}}(3)=-64\) while the stored \(1/\Phi_{10}^{\mathrm{un}}\) coefficient is \(248\). The promotion to ordinary BPS counts or compact Hall source dimensions is exactly the parity-fixtured finite Hall--Borcherds recognition problem, not a consequence of the denominator character alone. Part~(d) is the saddle-point evaluation of the contour integral; the leading saddle is the attractor point; the Rademacher expansion organises the subleading contributions as a sum over $W^{(2)}(\Lambda^{2,1}_{II})$-images of the attractor after the contour normalisation has been fixed.
\end{proof}

\begin{remark}[Six \texorpdfstring{$U$}{U}-duality orbits of \texorpdfstring{$\mathfrak{g}_{\Delta_5}$}{g-Delta-5} roots, \emph{not} six applications of \texorpdfstring{$\Phi_3$}{Phi-3}]
\label{rem:k3e-six-u-duality-orbits}
\index{U-duality!six BPS orbits}%
The six routes to the group $G(K3\times E)$ discussed in earlier sections are \emph{six distinct constructions}, not six applications of the CY-to-chiral functor $\Phi_3$. Under Theorem~\ref{thm:k3e-msw-harvey-moore}(c), this scope-declaration refines to a geometric statement about the $U$-duality group $G_{\mathrm{U}}(\Z) = \mathrm{SO}(22, 6; \Z)$ of type IIA on $K3\times T^{2}$: the imaginary-root lattice $\Lambda^{2,1}_{II} \hookrightarrow \Lambda^{4,20} \otimes H^{1,1}(T^{2})$ carries a residual $\mathrm{SO}(3, 2; \Z)$ action, and the $1/4$-BPS index $\widehat{d}(m, n, \ell)$ is a function of the $U$-duality invariant $D = 4mn - \ell^{2}$ together with the $T$-duality invariant $\gcd(m, n, \ell)$ (Banerjee--Sen; Dabholkar--Gomes--Murthy). The six constructions of $G(K3\times E)$ stratify by six orbit types:
\begin{enumerate}[label=\textup{(\Roman*)}]
\item \textup{Categorical-Euler orbit $D < 0$, $\gcd = 1$:} primitive real roots, producing the K\"unneth-multiplicative categorical Euler characteristic $\kappa_{\mathrm{cat}}(K3 \times E) = \chi(\mathcal{O}_{K3}) \cdot \chi(\mathcal{O}_E) = 2 \cdot 0 = 0$ on the total space; the fibre reading $\chi(\mathcal{O}_{K3}) = 2$ is a \emph{per-fibre} number, not a total-space $\kappa_{\mathrm{cat}}$.
\item \textup{Heisenberg-additivity orbit $D < 0$, $\gcd > 1$:} imprimitive real roots, giving $\kappa_{\mathrm{ch}}^{\mathrm{Heis}}(K3\times E) = 3$ via rank-additivity over the fibration.
\item \textup{Light-like orbit $D = 0$:} three vertices at infinity with $\tau(a) = 9$, producing the $\eta(q)^{9}$ factor in the Borcherds product.
\item \textup{Negative-norm orbit $D > 0$, $\gcd = 1$:} primitive imaginary roots, giving $\kappa_{\mathrm{BKM}}(\Delta_5)=c_{1}(0)/2 = 5$ (Theorem~\ref{thm:borcherds-weight-kappa-BKM-universal}); the unnormalised squared Igusa form \(\Phi_{10}^{\mathrm{un}}=\Delta_5^2\) has weight~$10$.
\item \textup{Negative-norm $N$-twisted orbit $D > 0$, $\gcd = N$:} the $\Gamma_0(N)$-twisted orbits give the Gritsenko paramodular forms $\Delta_{k_N}$ at $N\in\{2, 3, 4, 6\}$, realising the universal $\kappa_{\mathrm{BKM}}(\Phi_{N}) = c_{N}(0)/2$.
\item \textup{Fiber-rank orbit:} the K3 fiber Euler number $\chi_{\mathrm{top}}(K3) = 24$ governs the rank-$1$ sector (G\"ottsche), giving $\kappa_{\mathrm{fiber}} = 24$.
\end{enumerate}
At every orbit a different mathematical object is computed by a different construction. The scope discipline between the $\Phi$-functor and object-level correspondences is honoured: $\Phi_3$ targets \emph{one} framed chiral algebra $A_{K3\times E}$ with $\kappa_{\mathrm{ch}}^{\mathrm{Heis}} = 3$; the BKM weight $5$, the fiber-rank $24$, and the $N$-twisted paramodular weights come from \emph{different} lifts of the genus-$1$ K3 data, each with its own $U$-duality orbit interpretation. The five subscripted values $\{\kappa_{\mathrm{cat}}, \kappa_{\mathrm{ch}}^{\mathrm{Hodge}}, \kappa_{\mathrm{ch}}^{\mathrm{Heis}}, \kappa_{\mathrm{BKM}}, \kappa_{\mathrm{fiber}}\}(K3 \times E) = \{0, 0, 3, 5, 24\}$ of Chapter~\ref{ch:cy-d-kappa-stratification} are produced by five distinct constructions on the same CY$_3$ datum; four of them (cat, Heis, BKM, fiber) read off four of the six orbits enumerated here (I, II, IV, VI), and the chiral Hodge supertrace $\kappa_{\mathrm{ch}}^{\mathrm{Hodge}} = \sum_q (-1)^q h^{0, q}(K3 \times E) = 0$ is the compact-Hodge specialisation of the Stage-$2$ output, sharing the value $0$ with $\kappa_{\mathrm{cat}}$ by the identity $\sum_q (-1)^q h^{0, q}(X) = \chi(\cO_X)$ on compact CY$_3$ while remaining a distinct construction. The K3-fibre value $\chi(\mathcal{O}_{K3}) = 2$ does not appear in the total-space spectrum because $\kappa_{\mathrm{cat}}$ is K\"unneth-multiplicative, not additive. The remaining two orbits (III, V) enumerate the $\eta^{9}$ light-like contribution and the $N$-twisted Gritsenko family.
\end{remark}

\begin{remark}[Epistemic scope of the holographic identification]
\label{rem:k3e-holographic-epistemic-scope}
Theorem~\ref{thm:k3e-msw-harvey-moore}(b) is an identity of $\mathrm{Sp}_4(\Z)$-modular functions on $\mathbb{H}_2$, rigorously established by Oberdieck--Pixton together with DMVV and Bryan. The physical reduction of Maldacena--Strominger--Witten from M$5$-brane worldvolume to MSW $(0, 4)$ SCFT is a physical derivation with mathematical output; the mathematical content in (b) does not depend on the M-theory reduction. Part~(c) is a theorem within generalised-Kac--Moody algebra (Harvey--Moore $1996$ Theorem~$6.4$) combined with the $\mathfrak{g}_{\Delta_5}$ construction of Theorem~\ref{thm:k3e-denominator}; the physical interpretation as ``algebra of BPS states'' is a physical identification, but the mathematical relation is the product-level identity
\[
 (\Phi_{10}^{\mathrm{un}})^{-1}=\Delta_5^{-2},
\]
with the \(\phi_{0,1}\) coefficients appearing as denominator exponents and the BPS indices appearing only after the reciprocal square is expanded. Part~(d) is an arithmetic theorem about the Fourier coefficients of \(1/\Phi_{10}^{\mathrm{un}}\) (Dabholkar--Murthy--Zagier Theorem~$1.1$), independent of the holographic duality. In summary: the modular identities are proved; the AdS$_{3}$/CFT$_{2}$ and M$5$-brane reductions supply the physical interpretation but are not load-bearing in the mathematical chain. The CY-to-chiral functor $\Phi_3$ of Theorem~\ref{thm:cy-to-chiral-d3} supplies the chiral algebra $A_{K3\times E}$ with $\kappa_{\mathrm{ch}}^{\mathrm{Heis}} = 3$; the BKM-weight identification $\kappa_{\mathrm{BKM}}(\Delta_5) = 5$ is a separate construction through the Borcherds lift of $\phi_{0,1}$ and the Harvey--Moore denominator algebra; the two do not compete.
\end{remark}

\begin{remark}[Primary literature support]
\label{rem:k3e-msw-harvey-moore-primary-literature}
Primary:
Maldacena--Strominger--Witten 1997 arXiv:hep-th/$9711053$ \S$3$
(M$5$-brane reduction to MSW $(0, 4)$ SCFT);
Harvey--Moore 1996 arXiv:hep-th/$9610237$ Thm.~$6.4$
(algebra of BPS states is a generalised Kac--Moody algebra);
Dijkgraaf--Verlinde--Verlinde 1997 arXiv:hep-th/$9607026$
(identification \(1/\Phi_{10}^{\mathrm{un}}\) = $1/4$-BPS index generating function);
Strominger--Vafa 1996 arXiv:hep-th/$9601029$
(microstate counting for $1/4$-BPS black holes);
Maldacena--Moore--Strominger 1999 arXiv:hep-th/$9903163$
(multi-charge prefactor and counting in D$1$-D$5$ system);
Sen 2007 arXiv:hep-th/$0702150$; Sen 2008 arXiv:$0803.1014$
(attractor chamber; Rademacher expansion of \(1/\Phi_{10}^{\mathrm{un}}\));
Dabholkar--Murthy--Zagier 2012 arXiv:$1208.4074$ Thm.~$1.1$
(arithmetic Rademacher expansion);
Oberdieck--Pixton 2018 arXiv:$1808.03524$ Conjecture~A
(DT partition function of $K3\times T^{2}$);
Bryan 2015 arXiv:$1511.02575$ Thm.~$2$
(genus-$2$ identification on $K3\times T^{2}$);
Dijkgraaf--Moore--Verlinde--Verlinde 1996
arXiv:hep-th/$9608096$
(second-quantised $\Sym^{\infty}(K3)$ elliptic genus).
\end{remark}


\section{Synthesis: Borcherds lift mechanics, type-IV scope, siblings, AdS\texorpdfstring{$_3$}{-3}}
\label{sec:k3e-synthesis}

Five mutually constraining statements: explicit Borcherds-lift coefficients forcing Weyl equivariance of $\fg_{\Delta_5}$; the type-IV obstruction which forbids an \(\mathrm O(3,3)\) envelope BKM; a dimension-stratified sibling census separating Borcherds Monster, Igusa $\fg_{\Delta_5}$, and Fake Monster onto the correct CY dimensions; the AdS$_3$/CFT$_2$ throat whose dyon partition function is the inverse paramodular lift; and the two-scope universal Borcherds weight identity $\kBKM(\Phi_N) = c_N(0)/2$.

\subsection{Explicit Borcherds lift of \texorpdfstring{$\phi_{0,1}$}{phi-0,1}}
\label{subsec:k3e-platonic-borcherds-lift}

\begin{theorem}[Borcherds lift of \texorpdfstring{$\phi_{0,1}$}{phi-0,1}: explicit coefficients and Weyl equivariance]
\label{thm:plat-borcherds-lift}
The K3 weak Jacobi form
\[
\phi_{0,1}(\tau, z)
= y^{-1} + 10 + y
+ q(10y^{-2}-64y^{-1}+108-64y+10y^2)+O(q^2)
\]
has initial Fourier coefficients
\[
 f(0, 0) = 10, \qquad f(0, \pm 1) = 1, \qquad f(1, \pm 2) = 10, \qquad f(1, \pm 1) = -64, \qquad f(1, 0) = 108,
\]
(Eichler--Zagier~1985 \S9; confirmed against Dabholkar--Murthy--Zagier~2012 Table~1). The singular theta lift of Borcherds~1998 Theorem~13.3 produces the weight-$5$ paramodular form
\[
 \tfrac{1}{64}\,\Delta_5(2Z) \;=\; \Phi(Z)
\]
on $\mathcal{H}_{3, 2} \simeq \mathbb{H}_2$. Two structural consequences follow at the level of the Fourier coefficients $f(n, l, m)$ of $\Delta_5$:
\begin{enumerate}[label=\textup{(\roman*)}]
\item \textup{Weyl-group equivariance.} The divisibility $64 \mid f(n, l, m)$ holds for every admissible triple; the quotient $m(a) = -f(n, l, m)/64$ is an integer on the odd-imaginary cone and the associated multiplicity-$(-m(a))$ assignment is invariant under $W^{(2)}(\Lambda^{2,1}_{II})$.
\item \textup{$\tau(a_0) = 9$ cusp multiplicity.} The Macdonald-type identity
\[
 1 \;+\; \frac{1}{64}\sum_{t \geq 0} f(1 + 2t,\, 1,\, 1)\, q^t \;=\; \prod_{k \geq 1}(1 - q^k)^{9}
\]
pins the even-imaginary multiplicity at the primitive light-like root $a_0$ to $\tau(a_0) = 9$, consistent with the $\eta^9$-factorisation of the Borcherds product at the $0$-dimensional cusp.
\item \textup{Exterior-square isomorphism.} The Siegel / orthogonal identification
\[
 \mathrm{Sp}_4(\mathbb{Z}) / \{\pm I_4\} \;\xrightarrow{\;\wedge^2\text{-Pfaffian}\;}\; \mathrm{SO}_+(\Lambda^{3, 2})\big/\{\pm I_5\}
\]
refines Lemma~\ref{lem:sp4-so-iso} and identifies $\mathrm{PGL}_2(\mathbb{Z}) \simeq W^{(2)}(\Lambda^{2,1}_{II}) \rtimes S_3$.
\end{enumerate}
\end{theorem}

\begin{proof}[Attribution]
The Fourier coefficients of $\phi_{0,1}$ are computed from the theta-ratio representation of Eichler--Zagier~1985 \S9; the first few are tabulated in Dabholkar--Murthy--Zagier~2012 Table~1 and independently in Lorgat~2020 Theorem~$3$. The identity $\Delta_5(2Z)/64 = \Phi(Z)$ is Gritsenko--Nikulin~1998 Theorem~2.1 combined with Borcherds~1998 Theorem~10.1. Divisibility $64 \mid f(n, l, m)$ at every admissible triple is Lorgat~2020 Theorem~$3$; the consequence that $m(a) \in \mathbb{Z}$ and the associated multiplicity assignment is Weyl-invariant is Gritsenko--Nikulin~1998 Proposition~2.5. The Macdonald identity and the $\tau(a_0) = 9$ cusp multiplicity are Gritsenko--Nikulin~1998 \S2, reproduced in Lorgat~2020 Lemma~1. The exterior-square isomorphism in~(iii) is Weil~1964 combined with Lemma~\ref{lem:sp4-so-iso}.
\end{proof}

\begin{remark}[Why \texorpdfstring{$64$}{64} is rigid]
\label{rem:k3e-platonic-64-rigid}
The exponent magnitude $64 = |c(3)| = 2 \cdot 32$ is simultaneously the leading fermionic signed-character magnitude of $\mathfrak{g}_{\Delta_5}$ (Proposition~\ref{prop:k3e-super-grading}) and the determinant $|\det C_{\mathrm{BKM}}|\cdot 2 = 32 \cdot 2$ of twice the Cartan (Theorem~\ref{thm:bkm-delta5-is-borch-F3}). The divisibility statement~(i) is the algebraic shadow of this double coincidence: the super-Serre cancellation of the odd imaginary sector is forced through the $64$-factor at every order of the Borcherds product.
\end{remark}

\subsection{GBKM \texorpdfstring{$\mathfrak{g}_{\Delta_5}$}{g-Delta-5}: super-completion of \texorpdfstring{$F_3$}{F-3}}
\label{subsec:k3e-platonic-gDelta5}

\begin{theorem}[Structure of \texorpdfstring{$\mathfrak{g}_{\Delta_5}$}{g-Delta-5}]
\label{thm:plat-gDelta5}
$\mathfrak{g}_{\Delta_5}$ is the Lie superalgebra with
\begin{enumerate}[label=\textup{(\alph*)}]
\item three real simple roots $\{\delta_1, \delta_2, \delta_3\}$ with Gram matrix $\mathrm{diag}(2, 2, 2) - 2(E - I)$, signature $(2, 1)$, eigenvalue multiset $\{+4, +4, -2\}$;
\item Weyl vector $\rho = \tfrac{1}{2}(\delta_1 + \delta_2 + \delta_3) = f_2 - \tfrac{1}{2}\,f_3 + f_{-2}$ with $(\rho, \delta_i) = -1$ for $i = 1, 2, 3$ and $(\rho, \rho) = -3/2$;
\item even imaginary simple roots $\Delta^{\mathrm{im}}_0 = \{\tau(a) \cdot a : (a, a) = 0,\, \tau(a) > 0\}$ with $\tau(a_0) = 9$ at the primitive light-like $a_0$;
\item odd imaginary simple roots $\Delta^{\mathrm{im}}_1 = \{m(a) \cdot a : (a, a) < 0,\, m(a) < 0\}$ with $m(a) = -f(n, l, m)/64$, integer by Theorem~\ref{thm:plat-borcherds-lift}(i);
\item signed root characters \(\operatorname{sdim}\mathfrak g_{\Delta_5,\alpha}=f(nm, l)\) at $\alpha = (n, l, m)$, with Kac asymptotic \(f(n, l) = O(\exp\sqrt{4n - l^2})\); ordinary root dimensions require the target parity fixture;
\item real-root subalgebra $\mathfrak{g}_{\Delta_5}^{\mathrm{re}}$ isomorphic to the Feingold--Frenkel rank-$3$ hyperbolic Kac--Moody algebra $F_3$ (Proposition~\ref{prop:bkm-delta5-real-cartan-is-F3}); $\mathfrak{g}_{\Delta_5}$ is the Borcherds super-Serre completion of $F_3$ by the odd imaginary simple roots in~(d).
\end{enumerate}
Denominator identity:
\[
 \Phi(Z) \;=\; \sum_{w \in W^{(2)}} \det(w)\Bigl[\exp\!\bigl(-2\pi i\,(w\rho,\, Z)\bigr) - \sum_{a} m(a)\,\exp\!\bigl(-2\pi i\,(w(\rho + a),\, Z)\bigr)\Bigr].
\]
\end{theorem}

\begin{proof}
Items~(a) and~(b) are Proposition~\ref{prop:k3e-cartan-and-weyl}. Items~(c)--(e) are Construction~\ref{constr:k3e-roots} combined with Theorem~\ref{thm:k3e-product} and the Kac asymptotic for weak Jacobi forms of index~$1$ (Hardy--Ramanujan--Rademacher; Bruinier~2002 \S2). Item~(f) is Theorem~\ref{thm:bkm-delta5-is-borch-F3}: the real-root Cartan of $\mathfrak{g}_{\Delta_5}$ matches the Feingold--Frenkel $F_3$ Cartan $(2 I - 2(E - I))$ with determinant $-32$ and signature $(2, 1)$; $F_3$ is not isomorphic to the affine $\widehat{\mathfrak{sl}}_3$ (whose Cartan has determinant $0$ and signature $(2, 0, 1)$), and the Borcherds product $\Delta_5$ has $\eta^9$ factorising as $\eta^1 \cdot \eta^8$ at the light-like cusp rather than the $\eta^6 \cdot \eta^3$ that an $\widehat{\mathfrak{sl}}_3$ denominator would force. The super-completion statement is Borcherds~1988 \emph{J.\ Algebra}~115 \S9: adjoin $\mathbb{Z}_2$-graded imaginary simple roots whose signed root character is given by the K3 elliptic-genus coefficients; the Jacobi identity between $F_3$ and the adjoined odd imaginary sector is encoded in the super-Serre relations. The denominator identity is Theorem~\ref{thm:k3e-denominator}; the explicit sum over the Weyl chamber is Gritsenko--Nikulin~1998 Theorem~2.1.
\end{proof}

\begin{remark}[Separation from \texorpdfstring{$\widehat{\mathfrak{sl}}_3$}{widehatsl-3}]
\label{rem:k3e-platonic-sl3-separation}
The $\eta^9$ factor at the $0$-dimensional cusp of $\Delta_5$ encodes the nine even-imaginary multiplicity at $a_0$; it does \emph{not} identify an affine $\widehat{\mathfrak{sl}}_3$ embedded in $\mathfrak{g}_{\Delta_5}$. Two independent obstructions rule out such an embedding. First, Cartan mismatch: $\widehat{\mathfrak{sl}}_3$ is rank-$2$ positive-definite affine (Cartan determinant $0$, signature $(2, 0, 1)$), while the real-root lattice $\Lambda^{2,1}_{II}$ is rank-$3$ hyperbolic of signature $(2, 1)$ with Cartan determinant $-32$ (Remark~\ref{rem:bkm-drinfeld-cartan-table}). Second, denominator mismatch: the Macdonald denominator of $\widehat{\mathfrak{sl}}_3$ carries an $\eta^{11}$ factor (level-$1$ Weyl--Kac denominator times Lepowsky's cubic lattice-sum); the $\eta^9$ in $\Delta_5$ factorises as $\eta^1 \cdot \eta^8$ (Weyl-vector shift against the Siegel theta-series), not as $\eta^6 \cdot \eta^3$. The correct real-root subalgebra is Feingold--Frenkel $F_3$ (Theorem~\ref{thm:bkm-delta5-is-borch-F3}).
\end{remark}

\subsection{Explicit \texorpdfstring{$\mathrm{Sp}_{K3, E}$}{Sp-K3, E}: Mukai Heisenberg tensor BPS Yangian}
\label{subsec:k3e-platonic-sp-k3e}

\begin{theorem}[Stage-\texorpdfstring{$2$}{2} specialisation on \texorpdfstring{$K3 \times E$}{K3 x E}]
\label{thm:plat-Sp-K3E}
Let $\mathcal{F}_{K3 \times E} = \PhiFA_3(\mathrm{Perf}(K3 \times E)) \in E_3\text{-}\HolFA(K3 \times E)$ be the canonical Stage-$1$ output of the two-stage factorisation of $\Phi_3$, and let $\mathfrak{heis}_{\mathrm{Muk}}$ denote the abelian Heisenberg Lie algebra of the Mukai lattice $\Lambda_{\mathrm{Muk}} = \mathrm{II}_{4, 20}$ (signature $(4, 20)$; equivalently rank $24$ with the Mukai pairing). Let $\mathfrak{g}^{\mathrm{BPS}}_{K3}$ denote the conditional BPS target $\mathfrak{n}_+(\mathfrak{g}_{\Delta_5})$ of Theorem~\ref{thm:k3e-bps-lie}, equipped with the Maulik--Okounkov $R$-matrix braiding on its Nakajima module category when the compact Hall comparison data are supplied. On the compact Hall promotion locus of Corollary~\ref{cor:k3e-finite-height-promotion-obstruction}, including the compact CY$_3$ closure input of either $B_{\mathrm{TCFT}}^{(2)}$ or the filtered $HH^{-2}_{E_1}$ theorem, the Stage-$2$ specialisation along the K3-fibre $\Sigma_2 = K3$ decomposes as
\begin{equation}
\label{eq:plat-Sp-K3E}
 \mathrm{Sp}_{K3, E}\bigl(\mathcal{F}_{K3 \times E}\bigr) \;\simeq\; U^{\mathrm{ch}}\bigl(\mathfrak{heis}_{\mathrm{Muk}}\bigr) \;\otimes\; U^{\mathrm{ch}}\bigl(\mathfrak{g}^{\mathrm{BPS}}_{K3}\bigr)
\end{equation}
as $E_1$-chiral algebras on $E$, with the oriented hCS--Hall map,
the witnessed K3-fibre specialisation, the Hall--Drinfeld double
completion, the Borcherds denominator comparison, and the transported
coproduct/associator/$R$-matrix data supplying the comparison. The abelian
Mukai-Heisenberg character and Humbert-residue identity are proved
separately in Theorem~\ref{thm:plat-heis-mukai-48}; the displayed
non-abelian BPS factor is the Hall--Borcherds component. In
the slimmer form where the Heisenberg sector is restricted to the
transverse $\mathrm{II}_{3, 19}$ sublattice (the $J^3$-preserving
sigma-model Fock space), the total rank of the abelian sector is $22$;
the full Mukai rank $24$ is recovered on the covering metaplectic cover
of $\mathrm{Sp}_4(\mathbb{Z})$ tracking the $\nu_{\Delta_5}$ multiplier
(Remark~\ref{rem:k3e-maass-spin-cover}).
\end{theorem}

\begin{proof}
The Dunn--Lurie additivity $E_3 \simeq E_2 \otimes E_1$ on a product
$X \times Y$ of complex manifolds specialises factorisation homology:
$\int_{K3 \times E}(E_3\text{-}\mathcal{F}) \simeq
\int_E\!\bigl(E_1 \otimes \int_{K3}(E_2\text{-}\mathcal{F})\bigr)$.
On the closed oriented real-$4$-manifold $K3$, the abelian
Mukai-lattice oscillator sector has Fock character \(1/\eta^{24}\) and
vertex envelope
\(U^{\mathrm{ch}}(\mathfrak{heis}_{\mathrm{Muk}})\). The momentum
summation over \(\Lambda_{\mathrm{Muk}}\) is an indefinite
Siegel--Narain theta kernel after a positive four-plane is chosen, not
an ordinary holomorphic lattice theta character. The doubled oscillator
identity at the Humbert cusp is Theorem~\ref{thm:plat-heis-mukai-48}.
The residual $E_1$-factor on
$E$ carries the non-abelian BPS content only after the Hall--Borcherds
comparison datum transports the Hall positive half, pairing,
completion, denominator, coproduct, associator, and $R$-matrix through
the K3-fibre specialisation. Under those comparison and BRST hypotheses, the BPS Lie algebra
extraction criterion of Theorem~\ref{thm:k3e-bps-lie} and the
conditional Frenkel--Kac/Borcherds closure of
Theorem~\ref{thm:k3e-frenkel-kac-closure} identify the factor with
$U^{\mathrm{ch}}(\mathfrak{n}_+(\mathfrak{g}_{\Delta_5}))$, and the
Maulik--Okounkov braiding gives the stated $R$-matrix tensor
decomposition. The rank discrepancy $22$ vs $24$ is the signature
distinction $(3,19)$ vs $(4,20)$: the transverse sigma-model lattice is
signature $(3,19)$ and carries $22$ bosonic oscillators; the full Mukai
lattice $\mathrm{II}_{4,20}$ includes the $H^0\oplus H^4$ rank-$2$
summand that lifts the Heisenberg rank to $24$ on the spin double cover
(Maass~1964).
\end{proof}

\begin{remark}[Two Heisenberg ranks in one specialisation]
\label{rem:k3e-platonic-heis-rank-22-24}
The appearance of both $22$ (transverse rank) and $24$ (full Mukai rank) in the Heisenberg sector of $\mathrm{Sp}_{K3, E}(\mathcal{F}_{K3 \times E})$ is structurally significant. The bar Euler product at rank $24$ gives the $\eta^{-48}$ match to $\Delta_5^{-2}|_{H_1, z = 0}$ of Theorem~\ref{thm:plat-heis-mukai-48}; the rank-$22$ transverse reading gives the $c_{2d} = -214 = -(22 \cdot 10 - 6)$ central charge of the Chacaltana--Distler class-$\mathcal{S}$ parent $\mathcal{T}[A_1, \Sigma_{0, 24}]$ via iterated Drinfeld--Sokolov (Arakawa~2017; Creutzig--Linshaw~2017). The two ranks are related by the metaplectic multiplier $\nu_{\Delta_5}$: the $\mathbb{Z}/2$ central extension that turns a half-integral Siegel multiplier system into a well-defined lift on the Maass spin cover also lifts the transverse rank from $22$ to $24$ through the Mukai lattice pair $(H^0, H^4)$.
\end{remark}

\begin{theorem}[Heisenberg--Mukai \texorpdfstring{$\eta^{-48}$}{eta--48} at the Humbert cusp]
\label{thm:plat-heis-mukai-48}
The graded character of $U^{\mathrm{ch}}(\mathfrak{heis}(\Lambda_{\mathrm{Muk}})^{\oplus 2})$, the doubled Mukai Heisenberg chiral algebra, equals the residue of $(2\pi i z)^2 \Delta_5^{-2}$ at the Humbert divisor $H_1 \subset \overline{\mathcal{A}_2}$ along $z \to 0$:
\[
 \chi_{\mathrm{Heis}(\mathrm{Muk}(K3)^{\oplus 2})}(q) \;=\; \prod_{n \geq 1}(1 - q^n)^{-48} \;=\; -\,q^{-2}\bigl[(2\pi i z)^2\, \Delta_5^{-2}\bigr]\Big|_{H_1,\, z = 0}.
\]
The identity holds to all orders; the Shapovalov form on the doubled Heisenberg is block-diagonal with unit determinants (no null vectors), so the match is unconditional, not a Virasoro-minimal truncation.
\end{theorem}

\begin{proof}[Attribution]
Gritsenko--Nikulin~1996 Corollary~5.2 (Humbert-divisor residue of $\Delta_5^{-2}$); Frenkel--Ben-Zvi~2004 Theorem~2.4.4 (Heisenberg character $\prod(1 - q^n)^{-r}$ at rank $r$); combined to give the exponent $r = 48 = 2 \cdot 24$ at the doubled Mukai lattice. The doubled rank absorbs the metaplectic multiplier of Remark~\ref{rem:k3e-maass-spin-cover}.
\end{proof}

\subsection{No Borcherds envelope on \texorpdfstring{$\Lambda^{3,3}$}{Lambda-3,3}}
\label{subsec:k3e-platonic-envelope}

The inclusion \(\Lambda^{3,2}\subset U^{\oplus3}\) is a lattice fact, not a
Borcherds-product machine.  Borcherds products in the form used here live on
type-IV domains for lattices of signature \((2,n)\) or \((n,2)\).

\begin{theorem}[Type-IV obstruction to an \texorpdfstring{$\mathrm O(3,3)$}{O(3,3)} envelope]
\label{thm:plat-envelope}
Let \(\Lambda^{3,3}\simeq U^{\oplus3}\).  There is no Borcherds product
\(\Psi_{\Lambda^{3,3}}\) on an \(\mathrm O^+(\Lambda^{3,3})\)-Hermitian
domain whose restriction gives \(\Delta_5\), and there is no envelope BKM
\(\mathfrak g^{\mathrm{BKM}}_{\Lambda^{3,3}}\) whose fixed subalgebra is
\(\mathfrak g_{\Delta_5}\).  Precisely:
\begin{enumerate}[label=\textup{(\roman*)}]
\item \(\Lambda^{3,3}\) has symmetric space
\(\mathrm O(3,3)/(\mathrm O(3)\times\mathrm O(3))\), not a type-IV
Hermitian symmetric domain.
\item If \(q\in\Lambda^{3,3}\) is primitive of square \(-2\), then
\(q^\perp\) has signature \((3,2)\); the K3 paramodular product belongs
instead to the type-IV lattice \(\Lambda^{3,2}\).
\item Harvey--Moore heterotic charge lattices in signature \((3,3)\) give
BPS theta data.  They do not by themselves define a denominator identity,
a Borcherds algebra, or a fixed-subalgebra theorem.
\item The Fake Monster remains the separate \(\mathrm{II}_{25,1}\) object
with denominator \(\Phi_{12}\).
\end{enumerate}
\end{theorem}

\begin{proof}
Borcherds' singular-theta construction gives automorphic products on
orthogonal type-IV domains.  These domains have two positive directions
after a choice of convention.  The symmetric space of \(\mathrm O(3,3)\)
has three positive directions and is not the Hermitian domain on which
the product theorem lives.  Therefore the displayed \(\Psi_{\Lambda^{3,3}}\)
has no source theorem.  Without a denominator product there is no attached
generalised Kac--Moody algebra and no fixed-subalgebra statement.  The
restriction which exists in the K3 calculation is the ordinary passage to
the \(\Lambda^{3,2}\) type-IV datum underlying \(\Delta_5\).
\end{proof}

\subsection{Dimension-stratified sibling specialisations}
\label{subsec:k3e-platonic-siblings}

The denominator siblings stratify by dimension and by host. The Igusa
\(\mathfrak{g}_{\Delta_5}\) is the \(d=3\) K3-fibre Hall/Borcherds
row attached to \(K3\times E\). The Borcherds Monster row is a
non-geometric VOA row attached to \(V^\natural\) and the Cartan lattice
\(\mathrm{II}_{1,1}\); it is not a compact CY$_3$ specialisation built
from a Leech torus. The Fake Monster row lies conditionally at \(d=5\)
on the Lorentzian lattice \(\mathrm{II}_{25,1}\), not at
\(d=3\) sibling.

\begin{conjecture}[Dimension-stratified siblings]
\label{conj:plat-siblings-dim}
The dimension-stratified denominator census adjacent to \(\Phi\) splits
into three strata:
\begin{enumerate}[label=\textup{(\arabic*)}]
\item \emph{Borcherds Monster ($d = 3$ scalar row, non-geometric host).} Host \(V^\natural\); Cartan lattice \(\mathrm{II}_{1,1}\); denominator \(j(p)-j(q)\); \(\kBKM(j(p)-j(q))=0\) in the bi-weight-\((0,0)\) Hauptmodul normalisation of Lemma~\ref{lem:monster-outside-grassmannian-lift}.
\item \emph{K3-paramodular $\mathfrak{g}_{\Delta_5}$ ($d = 3$).} $X = K3 \times E$; specialisation datum $(\Sigma_2, C) = (K3, E)$; primitive denominator $1/\Delta_5(2Z)$; $\kBKM(\Delta_5) = \left.c_N(0)/2\right|_{N=1} = 5$.
\item \emph{Fake Monster ($d = 5$, not $d = 3$).} The Cartan lattice is \(\mathrm{II}_{25,1}\), and the orthogonal-domain lattice is \(\mathrm{II}_{26,2}\). The mnemonic \(K3\times K3\times E\) becomes a host only after a primitive \(\mathrm{II}_{25,1}\)-embedding, a \(d=5\) Stage-\(1\) construction, and the \(d=5\) Hall comparison are supplied. Its denominator is the Borcherds Fake-Monster product, and \(\kBKM(\Phi_{12})=c_0(0)/2=12\).
\end{enumerate}
The Fake Monster is obstructed from the $d = 3$ K3-type transverse lane by a rank count: the Leech lattice $\Lambda_{\mathrm{Leech}}$ has rank $24$, while a K3-type CY$_2$ surface has $b_2=22$ and an algebraic K3 has $h^{1, 1}\leq20$. This lane does not supply the Niemeier data required for the Fake Monster imaginary-root spectrum. The Monstrous moonshine module is recovered
through the Frenkel--Lepowsky--Meurman VOA construction, not as a
\(\PhiFA_3\)-specialisation of a compact CY$_3$. The Hauptmodul
property of the McKay--Thompson series is the scalar shadow, and the
Carnahan~2012 \(T_{g,h}\)-labels realise commuting pairs
\(\mathrm{Hom}(\pi_1(E_\tau),\mathrm{Mon})/\mathrm{Mon}\).
\end{conjecture}

\begin{remark}[Rank obstruction to Fake Monster at \texorpdfstring{$d = 3$}{d = 3}]
\label{rem:k3e-platonic-fake-monster-rank-obstruction}
The Leech-lattice rank-$24$ obstruction to a Fake Monster $d = 3$ sibling is quantitative. The Fake Monster BKM has imaginary-root multiplicities all equal to $24 = \mathrm{rk}\,\Lambda_{\mathrm{Leech}}$, requiring its specialisation cycle $\Sigma$ to carry a rank-$24$ transverse lattice. In the K3-type CY$_2$ transverse lane inside compact CY$_3$, $\Sigma_2$ has $b_2(\Sigma_2)=22$; for algebraic K3 surfaces this is $h^{1, 1}\leq20$ plus the two classes $h^{2,0}+h^{0,2}$. This lane does not supply the Niemeier rank-$24$ datum. At \(d=5\), the transverse fourfold \(\Sigma_4=K3\times K3\) has \(b_2(\Sigma_4)=44\) by K\"unneth, since \(H^1(K3)=0\); this removes the elementary rank obstruction but does not supply the required primitive Leech embedding or Hall comparison. The dimension-stratified census of Conjecture~\ref{conj:plat-siblings-dim} is therefore a conditional route, not an established construction.
\end{remark}

\begin{remark}[Six \texorpdfstring{$(\Sigma_2, C)$}{(Sigma-2, C)}-specialisations on \texorpdfstring{$K3 \times E$}{K3 x E}: cross-reference]
\label{rem:k3e-platonic-six-specialisations-crossref}
The explicit Stage-$2$ decomposition~\eqref{eq:plat-Sp-K3E} corresponds to the canonical specialisation $R_1 = (K3, E)$ in the six-entry catalogue of Remark~\ref{rem:six-sigma2-c-specialisations-K3E}: a single $\PhiFA_3$-output $\mathcal{F}_{K3 \times E}$ specialises against six transverse-surface/boundary-curve data $(R_1, R_2, R_3, R_4, R_5, R_6)$ producing six $E_1$-chiral shadows on $E$. The $N$-twisted paramodular ladder $\kBKM(\Phi_N) \in \{5, 3, 2, \tfrac32, 1\}$ on the CHL slice (Theorem~\ref{thm:k3e-universal-kBKM-two-scopes}, BKM-denominator scope) is the $R_2$-family; the transverse K3 Heisenberg at signature $(3, 19)$ with rank $22$ is the $R_6$-family. The six specialisations are \emph{not} six applications of $\Phi_3$ to one CY-category: one $\PhiFA_3$-output is evaluated against six distinct $(\Sigma_2, C)$-pairs inside the same $K3 \times E$, and $R_1$ supplies the conditional comparison
\[
\mathrm{Sp}_{K3, E}(\mathcal{F}_{K3 \times E})
\simeq
U^{\mathrm{ch}}(\mathfrak{heis}_{\mathrm{Muk}})
\otimes
U^{\mathrm{ch}}(\mathfrak{g}^{\mathrm{BPS}}_{K3})
\]
only on the compact Hall promotion locus of
Corollary~\ref{cor:k3e-finite-height-promotion-obstruction}.
\end{remark}

\subsection{Universal Borcherds weight, two scopes}
\label{subsec:k3e-platonic-universal-weight}

The two-scope structure of the universal weight identity consolidates
the CHL slice with the order-indexed Nikulin extension, while keeping
the Gritsenko--Cl\'ery triple-indexed catalogue separate.

\begin{theorem}[Universal Borcherds weight identity, two order-indexed scopes]
\label{thm:plat-universal-kBKM}
Two mutually consistent atlases produce a universal $\kBKM(\Phi_N) = c_N(0)/2$ identity:
\begin{description}
\item[BKM-denominator scope (CHL slice, square-root normalisation).] On $N \in \{1, 2, 3, 4, 6\}$,
\[
 \kBKM(\Phi_N) \;=\; c_N(0)/2 \;\in\; \bigl\{5,\ 3,\ 2,\ \tfrac32,\ 1\bigr\}
\]
with $c_N(0) = (10, 6, 4, 3, 2)$, the primitive square-root denominator weights of the CHL dyon squares (Jatkar--Sen~$2006$ \emph{JHEP}~$11$, $072$; Govindarajan--Krishna~2010~\cite{GovindarajanKrishna2010}; the once-recorded ladder $(5,4,3,2,1)$ is retracted), for the sibling BKM superalgebras $\mathfrak{g}_{\Phi_N}$. The $N = 4$ entry has half-integral weight $\tfrac32$ with multiplier system. Each $N$ is a CHL orbifold of $K3 \times E$ with $\varphi(N) \mid 2$.
\item[Borcherds-weight scope (Nikulin-order extension, Mathieu-twined normalisation).] On the
eight Nikulin-admissible symplectic orders $N \in \{1, 2, 3, 4, 5, 6, 7, 8\}$,
\[
 \kBKM(\Phi_N) \;=\; c_N(0)/2 \;\in\; \bigl\{10,\ 4,\ 3,\ 2,\ 2,\ 1,\ \tfrac32,\ 1\bigr\}
\]
as Borcherds singular-theta weights of the $g_N$-twined lifts, with
$c_N(0) = (20, 8, 6, 4, 4, 2, 3, 2)$ the frame-shape ones-exponents
at $N \geq 2$ and the $q^0\zeta^0$-coefficient of
$\mathrm{Ell}(K3) = 2\phi_{0,1}$ at $N = 1$
(Gritsenko 1994 Prop.~3.1;
Borcherds~1998 Theorem~13.3). The non-CHL entries at
$N \in \{5, 7, 8\}$ have $c_N(0)=(4,3,2)$ and weights $(2,\tfrac32,1)$ from
the $M_{24}$ frame shapes $1^4 5^4$, $1^3 7^3$, and
$1^2 2\,4\,8^2$.
\end{description}
The two scopes are two normalisations of one CHL slice: the twined
constants $(20, 8, 6, 4, 2)$ with weights $(10, 4, 3, 2, 1)$, and the
square-root constants $(10, 6, 4, 3, 2)$ with weights
$(5, 3, 2, \tfrac32, 1)$. No single family carries a mixed tuple drawn
from both. The extended scope records the additional
Mukai-admissible orders $N \in \{5,7,8\}$. This order-indexed
extension is distinct from the Gritsenko--Cl\'ery dd-modular
triple-indexed catalogue
\[
 (t,N;k)=
 (1,1;5),(2,1;2),(1,2;3),(3,1;1),(1,3;2),
 (4,1;\tfrac12),(1,4;\tfrac32),(2,2;1).
\]
Non-CHL hosts for the order-indexed extension are conjectural at the
CY$_3$ level: $N = 5$ via a Borcea--Voisin candidate; $N = 7$
obstructed for a smooth diagonal $K3 \times E$ quotient by the
Nikulin fixed-locus constraint; $N = 8$ represented by the
Mongardi--Tari--Wandel~2018 Kummer-$3$ fourfold, hence outside the
CY$_3$ host class.
\end{theorem}

\begin{proof}[Attribution]
CHL slice: Gritsenko--Nikulin~1995 Part~II Theorem~2.1 at $N = 1$;
Jatkar--Sen~$2006$ and Govindarajan--Krishna~2010 for
$N \in \{2, 3, 4, 6\}$; the five values
$\{5, 3, 2, \tfrac32, 1\}$ are $c_N(0)/2$ at the square-root
constants $(10, 6, 4, 3, 2)$.
Nikulin-order extension: Gritsenko 1994 Prop.~3.1 for
$N \in \{5, 7, 8\}$; Borcherds~1998 Theorem~13.3 for the singular-theta
weight formula. The $N = 5, 7, 8$ weights $2,\tfrac32,1$ follow from the
$M_{24}$ frame-shape ones-exponents $a_1 = 4, 3, 2$ of
$1^4 5^4$, $1^3 7^3$, $1^2 2\,4\,8^2$, giving
$c_N(0)=4,3,2$ in the twined normalisation. The displayed
Gritsenko--Cl\'ery triples are Theorem~\ref{thm:k3e-gc-eight-commuting-pair}
and are cited here only to prevent an index substitution. CY$_3$ host
identifications are conjectural: Borcea~1997 and Voisin~1993 for
$N=5$; Nikulin~1979 Corollary~3.3 for the $N=7$ obstruction;
Mongardi--Tari--Wandel~2018 for $N=8$.
\end{proof}

\begin{remark}[Two paramodular families on the CHL index set]
\label{rem:plat-two-paramodular-families-chl}
The BKM-denominator-scope tuple $(5,3,2,\tfrac32,1)$ of
Theorem~\ref{thm:plat-universal-kBKM} is the primitive square-root
paramodular ladder of Jatkar--Sen and Govindarajan--Krishna. The
order-indexed Nikulin extension is carried by the Mathieu-twined
normalisation, $(10,4,3,2,2,1,\tfrac32,1)$ on
$N\in\{1,\ldots,8\}$, restricting to the twined CHL slice
$(10,4,3,2,1)$; the non-CHL orders $5,7,8$ have
open CY$_3$ host constructions. The Gritsenko--Cl\'ery
dd-modular catalogue is a second, triple-indexed family with weights
$(5,2,3,1,2,\tfrac12,\tfrac32,1)$ on the eight triples displayed in
Theorem~\ref{thm:plat-universal-kBKM}; it overlaps the order-indexed
ladder only at $\Delta_5$. Throughout this chapter the default on the
symbol $\kappa_{\mathrm{BKM}}$ is the CHL/order-indexed Borcherds
weight convention unless a Gritsenko--Cl\'ery triple $(t,N;k)$ is
explicitly named.
\end{remark}

\subsection{AdS\texorpdfstring{$_3$}{-3} throat and dyon degeneracies}
\label{subsec:k3e-platonic-ads3}

The near-horizon geometry of the Type~IIB $D1$-$D5$-$p$ black hole on $K3 \times S^1$ is $\mathrm{AdS}_3 \times S^3 \times K3$; the $\mathbb{Z}/N\mathbb{Z}$ CHL orbifold replaces $K3$ by $K3^{g_N}$ and $S^3$ by $S^3/\mathbb{Z}_N$, producing a tower of dyon-counting paramodular forms.

\begin{theorem}[AdS\texorpdfstring{$_3$}{-3}/CFT\texorpdfstring{$_2$}{-2} throat and dyon degeneracy]
\label{thm:plat-AdS3}
The order-indexed physical dyon tower and the primitive BKM denominator
tower are distinct.  In the Jatkar--Sen physical tower
\[
 N\in\{1,2,3,5,7,11\},\qquad
 k_N^{\mathrm{JS}}=\frac{24}{N+1}-2=(10,6,4,2,1,0),
\]
the CHL near-horizon throat has the form
\[
 \mathrm{AdS}_3 \times S^3/\mathbb{Z}_N \times K3^{g_N}
\]
in the Sen attractor chamber, and the genus-two dyon index is extracted
from the reciprocal of the Jatkar--Sen paramodular form
\(\Phi^{\mathrm{JS}}_{k_N^{\mathrm{JS}}}\):
\[
 d_N(Q,P)
 =
 \oint_{\mathcal C_N}
 \frac{\exp\!\bigl(-i\pi(Q^2\rho+P^2\sigma+2(Q\!\cdot\!P)v)\bigr)}
 {\Phi^{\mathrm{JS}}_{k_N^{\mathrm{JS}}}(\rho,\sigma,v)}
 \,d\rho\,d\sigma\,dv .
\]
For the corresponding charge discriminant \(\Delta_N(Q,P)\) in the
Jatkar--Sen lattice normalisation,
\[
 \log |d_N(Q,P)|
 =
 2\pi\sqrt{\Delta_N(Q,P)}
 -
 \alpha_{\log}^{\mathrm{JS}}(N,\mathcal C_N,\mu_N,\mathcal P_N)\log\Delta_N
 +O(1),
\]
where the logarithmic coefficient is fixed only after the contour
\(\mathcal C_N\), charge-lattice measure \(\mu_N\), and polar orbit
\(\mathcal P_N\) are fixed.  The primitive BKM denominator tower is
instead
\[
 N\in\{1,2,3,4,6\},\qquad
 \kappa_{\mathrm{BKM}}(\Phi_N)=c_N(0)/2=\bigl(5,3,2,\tfrac32,1\bigr).
\]
The two order sets meet only at \(N=1,2,3\), where the Jatkar--Sen
weights \((10,6,4)\) are twice the primitive weights \((5,3,2)\).  The square relation
\(\Phi_{10}^{\mathrm{un}}=\Delta_5^2\) is the \(N=1\) Igusa convention;
no uniform identity
\(\Phi^{\mathrm{JS}}_{k_N}=(\Phi_N^{\mathrm{BKM}})^2\) is asserted for
the two towers.
\end{theorem}

\begin{proof}[Attribution]
DVV~1997 gives the unorbifolded dyon integral.  Jatkar--Sen~2005 gives
the physical CHL weights \(k_N^{\mathrm{JS}}=24/(N+1)-2\) on
\(\{1,2,3,5,7,11\}\), and Shih--Strominger--Yin~2005 gives the CHL
extension of the dyon-counting forms in that physical tower.  Sen's
attractor-contour analysis and Dabholkar--Murthy--Zagier~2012 give the
Rademacher expansion after the contour and polar data are fixed.  The
primitive BKM denominator weights are the separate Borcherds--Gritsenko
identity \(\kappa_{\mathrm{BKM}}(\Phi_N)=c_N(0)/2\) on
\(\{1,2,3,4,6\}\), already stated in
Theorem~\ref{thm:plat-universal-kBKM}.  Applying
\(24/(N+1)-2\) to \(N=4,6\) gives \(14/5\) and \(10/7\), so that formula
cannot be the primitive BKM tuple.
\end{proof}

\begin{remark}[Graviton finiteness equals \texorpdfstring{$\mathrm{HH}^2_{E_2}$}{HH-2-E-2}-vanishing]
\label{rem:k3e-platonic-graviton-finiteness-HH2}
The physical statement ``the near-horizon algebra of single-graviton states is finite-dimensional at fixed mass level'' is, at the chiral-boundary level, the algebraic statement $\mathrm{HH}^2_{E_2}(A_{K3}, A_{K3}) = 0$ in the $E_2$-Hochschild bicomplex of the K3 chiral algebra. This vanishing witnesses the rigidity of $A_{K3}$ as an $E_2$-algebra (no first-order $E_2$-deformations that preserve CY structure); the resulting uniqueness of the near-horizon chiral algebra is the algebraic content of the physical finiteness. The identification is at the chain level (explicit Hochschild complex of the lattice VOA at rank $24$) and at the $(\infty, 1)$-categorical level (the stable $\infty$-category of $A_{K3}$-modules has its $\mathrm{End}^{\mathrm{ch}}$-shifted obstruction class in degree~$2$ vanishing); both readings are load-bearing.
\end{remark}

% ======================================================================
\section{BPS entropy from the shadow obstruction tower}
\label{sec:phys-bps-entropy}
% ======================================================================

The shadow obstruction tower refines the Bekenstein--Hawking area law degree by degree. Let $D = 4nm - l^2$ denote the discriminant of a BPS state with charge $\alpha = (n, l, m) \in \Lambda^{2,1}_{II}$.

\begin{theorem}[BPS entropy: degree expansion via shadow tower]
\label{thm:phys-bps-entropy-degrees}
Assume the protected physical comparison identifies the chiral trace of
\(A_{K3\times E}^{(\Sigma_2,C)}\) with the BPS index in one fixed
compact Siegel normalisation
\[
 \mathcal N_{\mathrm{phys}}
 =
 (\mathcal N_{\mathrm{cont}},\mathcal P_{\mathrm{pol}},\mu,
 \epsilon_{\Delta/\Phi}).
\]
Then the microstate degeneracy \(\Omega(\gamma)\) of a large BPS black
hole with discriminant \(D \gg 1\) has the normalized expansion
\begin{equation}\label{eq:phys-bps-entropy-full}
\log \Omega(\gamma) \;=\;
4\pi\sqrt{D}
\;-\;
\alpha_{\log}^{\mathrm{phys}}(\mathcal N_{\mathrm{phys}})\log D
\;+\; \sum_{k \geq 1}
\frac{a_k(\mathcal N_{\mathrm{phys}})}{(\pi\sqrt{D})^k}
\;+\; O(e^{-\sqrt{D}}),
\end{equation}
where the coefficient
\(\alpha_{\log}^{\mathrm{phys}}(\mathcal N_{\mathrm{phys}})\) and the
Bessel coefficients \(a_k(\mathcal N_{\mathrm{phys}})\) are read in the
same contour, polar-orbit, measure, and primitive/square convention.
The three contributions arise from degrees \(2\), \(3\), and
\(\geq4\) of the shadow obstruction tower of~\(A_{K3 \times E}\):
\begin{enumerate}[label=\textup{(\roman*)}]
\item \textbf{Degree $2$ (Bekenstein--Hawking):} $S_{\mathrm{BH}} = 4\pi\sqrt{D}$.
The leading term is the $\kappa_{\mathrm{BKM}}$-level shadow: the Borcherds weight $\kappa_{\mathrm{BKM}}(\Delta_5) = \left.c_N(0)/2\right|_{N=1} = 5$ controls the area law through the weight of $\Delta_5$. (The chiral Hodge supertrace $\kappa_{\mathrm{ch}}^{\mathrm{Hodge}}(K3\times E) = \sum_q (-1)^q h^{0, q}(K3 \times E) = 0$ and the categorical Euler $\kappa_{\mathrm{cat}}(K3\times E) = 0$ by K\"unneth-multiplicativity on the total space; the two share the value $0$ via the compact identity $\sum_q (-1)^q h^{0, q}(X) = \chi(\cO_X)$ while remaining distinct constructions, and neither agrees with $\kappa_{\mathrm{BKM}}(\Delta_5) = 5$.)
\item \textbf{Degree $3$ (logarithmic channel):} the cubic shadow is the
first place where the one-loop determinant can enter the chiral trace.
Its numerical coefficient is
\(\alpha_{\log}^{\mathrm{phys}}(\mathcal N_{\mathrm{phys}})\), not a
function of \(\kappa_{\mathrm{BKM}}(\Delta_5)\) alone.  Sen's
near-horizon logarithmic formula and the Rademacher logarithmic
coefficient become a number in this chapter only after their charge,
measure, and contour conventions have been converted to
\(\mathcal N_{\mathrm{phys}}\).
\item \textbf{Degree $\geq 4$ (Bessel tower):} the quartic and higher
shadows produce power-suppressed terms
\(a_k(\mathcal N_{\mathrm{phys}})/(\pi\sqrt{D})^k\), matching the
Rademacher expansion of the Fourier coefficients of the unnormalised
Igusa square \((\Phi_{10}^{\mathrm{un}})^{-1}=\Delta_5^{-2}\) in the
same compact Siegel normalisation.
\end{enumerate}
\end{theorem}

\begin{proof}[Attribution]
Strominger--Vafa~1996 gives the leading area term.  Sen's logarithmic
entropy formula and Dabholkar--Murthy--Zagier~2012 Theorem~1.1 give the
one-loop and Rademacher data in their own normalisations.  The statement
above is the convention-converted form required by the chiral trace:
the degree-\(r\) truncation
\(\Theta^{\leq r}_{A_{K3\times E}}\) contains the first \(r\) terms of
\eqref{eq:phys-bps-entropy-full} only after the protected physical
comparison identifies the chiral trace, the BPS index, the polar data,
and the contour integral in the same normalisation
\(\mathcal N_{\mathrm{phys}}\).  Proposition~\ref{prop:k3e-rank-one-rademacher-packet}
proves the rank-one \(\phi_{0,1}\) Rademacher packet used by the
Borcherds product; it does not replace this compact Siegel
normalisation.
\end{proof}

\begin{remark}[Shadow depth class \texorpdfstring{$M$}{M} and genuine black holes]
\label{rem:phys-shadow-depth-bh}
The shadow depth $r_{\max}$ of $A_X$ classifies the BPS entropy regime:
\begin{itemize}
\item \textbf{Class $G$} ($r_{\max} = 2$): area law only; no black holes.
\item \textbf{Class $L$} ($r_{\max} = 3$): logarithmic correction present; small black holes.
\item \textbf{Class $C$} ($r_{\max} = 4$): one Bessel correction; small black holes with finite subleading.
\item \textbf{Class $M$} ($r_{\max} = \infty$): full infinite Rademacher--Bessel tower; \emph{genuine macroscopic black holes} with thermodynamic entropy and event horizon.
\end{itemize}
The \(\Delta_5\) denominator shadow is class~\(M\): its Rademacher--Bessel tower has infinite depth. The quantum vertex chiral group \(G(K3 \times E)\), once obtained from the compact Hall--Drinfeld double, must have this class~\(M\) scalar shadow. The class~\(M\) condition is a property of the Siegel modular denominator and its infinite product expansion, not a construction of the compact Hall algebra by itself.
\end{remark}

% ======================================================================
\section{Gopakumar--Vafa invariants as primitive protected root indices}
\label{sec:phys-gv-primitive}
% ======================================================================

The Gopakumar--Vafa invariants $n^g_\beta$ (for genus $g$ and curve class $\beta \in H_2(X, \mathbb{Z})$) are the \emph{primitive} protected BPS invariants: they are signed single-particle indices before multi-particle bound states are formed. The DT invariants, as protected indices, are obtained by ``second quantization.''

\begin{proposition}[GV invariants and the BKM root datum]
\label{prop:phys-gv-prim-root}
For a CY3 $X$ on a named locus where the Hall--Drinfeld/BKM double $G(X)$ has been constructed, the Gopakumar--Vafa invariants $n^g_\beta$ are the signed primitive imaginary-root data, equivalently the protected superdimensions, of the BKM superalgebra $\mathfrak{g}_X$. Ordinary primitive root multiplicities require the parity fixture or finite Hall--Borcherds recognition data. The DT partition function (= denominator identity in product form) is recovered from the GV invariants by
\begin{equation}\label{eq:phys-gv-dt}
\sum_{\beta} \sum_{n} \mathrm{DT}_{\beta,n}(X)\, q^n Q^\beta
\;=\;
M(q)^{\chi(X)} \cdot
\prod_{\beta > 0}\, \prod_{g \geq 0}\, \prod_{k=1}^\infty
\left(\frac{1}{1 - Q^\beta q^k}\right)^{(-1)^g n^g_\beta \binom{2g-2}{g-1+k}},
\end{equation}
where $M(q) = \prod_{k \geq 1}(1-q^k)^{-k}$ is the MacMahon function. The passage from $n^g_\beta$ to $\mathrm{DT}_{\beta,n}$ is the Weyl--Kac--Borcherds character formula applied to the BKM root datum of the constructed $G(X)$; outside such loci this sentence is a conjectural target, not a global definition.
\end{proposition}

\begin{proof}[Attribution]
The GV formula~\eqref{eq:phys-gv-dt} is Gopakumar--Vafa~1998, \emph{Adv.\ Theor.\ Math.\ Phys.}~\textbf{2}; the root-theoretic interpretation is Katz~2008, \emph{Invent.\ Math.}~\textbf{172} \S~2. The MacMahon factor $M(q)^{\chi(X)}$ is the contribution of degree-0 ideal sheaves (= the vacuum representation of $G(\mathbb{C}^3)$ at central charge $\chi(X)$), matching the DT/GW correspondence of Maulik--Nekrasov--Okounkov--Pandharipande~2003. The binomial coefficient $\binom{2g-2}{g-1+k}$ encodes the spin content of genus-$g$ BPS multiplets via the $\mathrm{SL}(2,\mathbb{Z})$-action on the cohomology of $\overline{\mathcal{M}}_{g,n}(X, \beta)$.
\end{proof}

\begin{remark}[Automorphic weight from \texorpdfstring{$\kappa_{\mathrm{BKM}}$}{kappa-BKM}]
\label{rem:phys-weight-kappa-ch}
For a CY$_3$ $X$ whose BKM superalgebra $\mathfrak g_X$ has Borcherds lift $\Phi_X$, the automorphic denominator of $G(X)$ carries Siegel weight $\kappa_{\mathrm{BKM}}(\Phi_X) = c_X(0)/2$. For $K3 \times E$: $\kappa_{\mathrm{BKM}}(\Delta_5) = \left.c_N(0)/2\right|_{N=1} = 5$, so $\mathrm{wt}(\Delta_5) = 5$; the unnormalised Igusa cusp \(\Phi_{10}^{\mathrm{un}}=\Delta_5^2\) has weight \(10 = 2\kappa_{\mathrm{BKM}}\). The chiral Hodge supertrace $\kappa_{\mathrm{ch}}^{\mathrm{Hodge}}(K3 \times E) = \sum_q (-1)^q h^{0, q}(K3 \times E) = 0$ and the categorical Euler $\kappa_{\mathrm{cat}}(K3 \times E) = 0$ (K\"unneth) share the value $0$ on the compact CY$_3$ via the identity $\sum_q (-1)^q h^{0, q}(X) = \chi(\cO_X)$ while remaining distinct constructions; $\kappa_{\mathrm{BKM}}$ is the single invariant that determines the automorphic weight of the denominator identity. The GV refined invariant $\Omega(\gamma; y) = \mathrm{Tr}_{\mathcal{H}_\gamma^{\mathrm{BPS}}} (-y)^{2j_3}$ (tracking spin $j_3$ of the BPS little group $\mathrm{SU}(2)_R$) further refines this to a $\mathbb{Z}$-graded root space $(\mathfrak{g}_X)_\alpha = \bigoplus_j (\mathfrak{g}_X)_{\alpha,j}$, the beginning of a categorification of the BKM superalgebra.
\end{remark}

% ======================================================================
\section{The Lorgat Conjecture~1 Borcherds-lift comparison}
\label{sec:k3e-lorgat-conj1-proof-tower}
% ======================================================================

The correspondence-restricted form of Lorgat~2020 Conjecture~1 separates into an automorphic layer and an enumerative layer. The automorphic layer is uniform on the CHL ladder \(N \in \{1,2,3,4,6\}\): the \(g_N\)-twined Jacobi input has constant \(c_N(0)\), its primitive Borcherds lift has weight \(c_N(0)/2\), and the scalar-square lift of the doubled input has weight \(c_N(0)\). The enumerative layer is stratified. The \(N=1\) line is a theorem by the Oberdieck--Pixton/Oberdieck--Pandharipande \(K3\times E\) reduced-DT identity and the KKV/GW/PT comparison. At \(N\in\{2,3\}\), Bryan--Oberdieck supply the CHL reduced-DT conjecture and coefficient theorems; the scalar identity with the primitive BKM square is conditional on the normalisation and all-curve-class DT inputs stated below. At \(N\in\{4,6\}\), the Bryan--Steinberg and orbifold-vertex comparisons remain evidence for the stated conjectural orbifold-DT formula; the global \(E\)-reduced DT identity is an input.

\subsection{Closure at \texorpdfstring{$(N, M) = (1, \mathrm{id})$}{(N, M) = (1, id)}}
\label{subsec:k3e-lorgat-conj1-N1}

The $N = 1$ case realises Lorgat~$2020$ Conjecture~$1$ as a four-way coincidence, each equality a primary-source theorem.

\begin{theorem}[Lorgat~\texorpdfstring{$2020$}{2020} Conjecture~\texorpdfstring{$1$}{1} at \texorpdfstring{$(N, M) = (1, \mathrm{id})$}{(N, M) = (1, id)}]
\label{thm:lorgat-conj1-N1-complete}
Let $X = K3 \times E$ with primitive polarisation $L$ of degree $h$ and $\beta_h \in H_2(K3, \mathbb{Z})$ primitive effective with $\langle \beta_h, \beta_h \rangle = 2h - 2$. For every positive BKM root $\alpha = (n, \ell, m) \in \Lambda^{2, 1}_{\mathrm{II}}$ the four protected integer invariants attached to $\alpha$ coincide:
\begin{equation}
\label{eq:four-way-mult-N1}
\operatorname{sdim}\mathfrak g_{\Delta_5,\alpha} \;=\; c_{K3}(h, \ell) \;=\; n^{0}_{(\beta_h, \ell)}(X) \;=\; N^{\mathrm{DT}}_{(\beta_h, \ell)}(X),
\end{equation}
with $h = nm$, witnessed globally by
\begin{equation}
\label{eq:OP-igusa-N1}
Z^{\mathrm{red}}_{\mathrm{DT}}(q, \tilde q, y)
\;=\;
-\,(\Phi_{10}^{\mathrm{OP}})^{-1}
\;=\;
-\,D_5^{-2}
\;=\;
-4096\,\Delta_5^{-2}.
\end{equation}
Consequently $\kappa_{\mathrm{BKM}}(\Delta_5) = \left.c_N(0)/2\right|_{N=1} = 5$ realises the universal Borcherds weight identity at $N = 1$. The equality is a signed root-character/protected-index equality; interpreting it as an ordinary BPS Hilbert-space count or compact Hall source dimension requires the parity fixture and finite Hall--Borcherds recognition data.
\end{theorem}

\begin{proof}
\emph{Gritsenko--Nikulin denominator.} Gritsenko--Nikulin~$1995$ Theorem~$1.4$ and Gritsenko--Nikulin~$1998$ Theorem~$1.1$ produce the rank-$3$ BKM Lie superalgebra $\mathfrak{g}_{\Delta_5}$ with Weyl chamber a fundamental domain for $W^{(2)}(\Lambda^{2, 1}_{\mathrm{II}})$, Gram matrix $2(2I - J_3)$ of the three real simple roots, Weyl vector $\rho = f_2 - \tfrac{1}{2} f_3 + f_{-2}$, and Weyl--Kac--Borcherds denominator $\Delta_5$.

\emph{Primitive denominator versus Igusa square.} Borcherds~$1998$ Theorem~$13.3$ and Gritsenko--Nikulin~$1998$ \S$2$ separate the primitive denominator from its Igusa square. The Borcherds lift of $\phi_{0,1}$ gives $\Delta_5$ of weight $5$, equivalently $(1/64)\,\Delta_5(2Z) = \mathrm{BorLift}(\phi_{0,1})(Z)$ after the stated normalisation; the lift of $2\phi_{0,1}$ gives the DVV/DMVV square \(\Phi_{10}^{\mathrm{un}}=\Delta_5^{\,2}\) of weight $10$. Primitive signed root characters belong to the $\Delta_5$ denominator, while the reduced-DT / Igusa trace is the square. Thus the BKM weight is $\kappa_{\mathrm{BKM}}(\Delta_5) = \left.c_N(0)/2\right|_{N=1} = 10/2 = 5$, not a weight-$5$ assertion about \(\Phi_{10}^{\mathrm{un}}\).

\emph{Reduced DT scalar.} Oberdieck~$2018$ (\emph{Duke Math.\ J.}~$167$:$3045$) proves the Oberdieck--Pixton conjecture by assembling the absolute reduced theory from the relative theory of $K3 \times \mathbb{A}^1$ through the degeneration $E \rightsquigarrow \mathbb{P}^1_{\mathrm{nodal}}$, rubber contribution $\eta^{24}$, and Maulik--Pandharipande--Thomas~$2010$ reduced MNOP Gromov--Witten/Donaldson--Thomas correspondence; Maulik--Toda~$2018$ tracks the Behrend-sign orientation for the $-1$ prefactor.

\emph{KKV and protected indices.} Pandharipande--Thomas~$2014$ (\emph{Forum Math.\ Pi}) Theorem~$1$ proves the Katz--Klemm--Vafa formula $\sum_h n^{0, K3}_h q^{h - 1} = \prod_n (1 - q^n)^{-24}$ in the stable-pairs formulation, identifying $n^{0}_{(\beta_h, \ell)}(X) = c_{K3}(h, \ell)$. Reduced MNOP at genus zero transfers the Gopakumar--Vafa integer BPS invariant to the Behrend-weighted Donaldson--Thomas protected count, and Borcherds-product expansion of $\Delta_5^{\,2}$ identifies $c_{K3}(nm, \ell)$ with \(\operatorname{sdim}\mathfrak{g}_{\Delta_5,\alpha}\), closing the four-way protected-index equality.

\emph{Motivic wall crossing.} At the $\sigma_0$-chamber adjacent to $\Delta_5$, the Kontsevich--Soibelman~$2008$ motivic wall-crossing formula~\S$6.2$ defines $\Omega(\gamma, \sigma_0)$ through the motivic Hall algebra; Joyce--Song~$2012$ \S$5$ (\emph{Memoirs AMS~$217$}) promotes $\Omega$ to the Behrend-weighted integer BPS count under chamber stability; chamber-independence follows from the Bridgeland--Smith embedding $\rho_{\mathrm{BS}}: \mathrm{Sp}_4(\mathbb{Z})/\{\pm I\} \hookrightarrow \mathrm{Aut}_{\mathrm{eq}}(D^{\mathrm{b}}(X))/{\sim}$.
\end{proof}

\begin{remark}[Five-invariant separation at \texorpdfstring{$N = 1$}{N = 1}]
\label{rem:k3e-lorgat-conj1-N1-kappa}
$\kappa_{\mathrm{BKM}}(\Delta_5) = 5$ by the primitive Borcherds denominator; $\kappa_{\mathrm{ch}}^{\mathrm{Hodge}}(K3 \times E) = \sum_q (-1)^q h^{0, q}(K3 \times E) = 0$ (chiral Hodge supertrace at $d = 3$); $\kappa_{\mathrm{ch}}^{\mathrm{Heis}}(K3 \times E) = 3$ (Heisenberg--Mukai Stage-$2$ specialisation); $\kappa_{\mathrm{cat}}(K3 \times E) = \chi(\mathcal{O}_{K3}) \cdot \chi(\mathcal{O}_E) = 2 \cdot 0 = 0$ (K\"unneth multiplicative on the total space, sharing its value with $\kappa_{\mathrm{ch}}^{\mathrm{Hodge}}$ via the compact identity $\sum_q (-1)^q h^{0, q}(X) = \chi(\cO_X)$ while remaining a distinct construction); $\kappa_{\mathrm{fiber}} = 24$ (Mukai lattice rank on the K3 fibre). The five $\kappa_\bullet$ values read from five distinct constructions; an unsubscripted invariant symbol conflates them.
\end{remark}

\subsection{Conditional scalar implication at \texorpdfstring{$(N, M) = (2, \mathrm{id})$}{(N, M) = (2, id)} and \texorpdfstring{$(3, \mathrm{id})$}{(3, id)}}
\label{subsec:k3e-lorgat-conj1-N23}

The Borcherds constants at $N \in \{2,3\}$ are fixed by the primitive CHL denominator convention:
\[
\mathrm{wt}\,F^{\mathrm{CHL}}_2=3,\qquad
\mathrm{wt}\,F^{\mathrm{CHL}}_3=2.
\]
The reduced-DT scalar is a different assertion. Bryan--Oberdieck formulate the primitive CHL Donaldson--Thomas identity as Conjecture~0.1 and prove first-coefficient cases in Theorem~0.1; their Siegel form has weight
\[
\left(\left\lceil \frac{24}{N+1}\right\rceil-2\right)_{N=2,3}=(6,4).
\]
Thus the scalar square of the manuscript's primitive BKM denominator has weights $(6,4)$, equal to the Bryan--Oberdieck dyonic weights. The equality of weights does not by itself identify the forms: passing from one convention to the other is an additional normalisation theorem, not a formal consequence of the Borcherds weight identity.

\begin{proposition}[Conditional CHL scalar identity at \texorpdfstring{$N \in \{2, 3\}$}{N in 2, 3}]
\label{thm:oberdieck-chl-N-2-3}
Let $X^{(N)} = (K3 \times E)/\mathbb{Z}_N$ be the free CHL quotient for the action generated by the Jatkar--Sen symplectic automorphism $g_N$ on $K3$ and translation by an $N$-torsion point on $E$, with $N \in \{2,3\}$. Let $F^{\mathrm{CHL}}_N$ be the primitive BKM denominator in the convention of Proposition~\ref{prop:k3e-square-root-normalisation}. Assume:
\begin{enumerate}
\item[\textup{(i)}] the Bryan--Oberdieck primitive CHL Donaldson--Thomas conjecture holds for $X^{(N)}$ in all primitive curve classes and in the expansion chamber used for the $E$-reduced DT series;
\item[\textup{(ii)}] the imprimitive classes are recovered from the primitive classes by the required reduced multiple-cover formula;
\item[\textup{(iii)}] a named normalisation identifies the Bryan--Oberdieck denominator in that chamber with the scalar square \(\Phi_N^{\mathrm{scal}}=(F^{\mathrm{CHL}}_N)^2\) in the manuscript's primitive BKM convention;
\item[\textup{(iv)}] the Behrend-sign and \(E\)-quotient conventions agree with the Oberdieck--Pandharipande sign convention at \(N=1\).
\end{enumerate}
Then
\[
Z^{\mathrm{red}}_{\mathrm{DT}}\bigl(X^{(N)}\bigr)(\tau, \sigma, z) \;=\; -\,\frac{1}{F^{\mathrm{CHL}}_N(\tau, \sigma, z)^{\,2}}, \qquad \mathrm{wt}\,F^{\mathrm{CHL}}_N \;=\; \frac{c_N(0)}{2} \;\in\; \{3, 2\},
\]
where \(c_2(0)=6\) and \(c_3(0)=4\) in the square-root normalisation. Without \textup{(i)}--\textup{(iv)}, the unconditional content used in this chapter is the primitive Borcherds-weight statement \(\mathrm{wt}\,F^{\mathrm{CHL}}_N=c_N(0)/2\), not the reduced-DT scalar identity.
\end{proposition}

\begin{proof}
The action on \(K3\times E\) is free because of the translation on \(E\), so the Bryan--Oberdieck CHL threefold is a smooth quotient and no crepant-resolution arrow is part of this implication. Assumption~\textup{(i)} gives the reduced primitive DT series as the negative reciprocal of the Bryan--Oberdieck Borcherds lift. Assumption~\textup{(ii)} extends the primitive series to the stated class range. Assumption~\textup{(iii)} replaces the Bryan--Oberdieck denominator by the scalar square of the primitive BKM denominator used in Proposition~\ref{prop:k3e-square-root-normalisation}. Assumption~\textup{(iv)} fixes the sign. The displayed formula is exactly the composite of these four hypotheses.

Bryan--Oberdieck Theorem~$0.1$ verifies the first \(t^{-1/N}\)-coefficient and the first \(q^{-1}\)-coefficient of their conjectural series; the \(t^0\)-coefficient is conditional on the Behrend-function conjecture they cite. These coefficient theorems are evidence for \textup{(i)}, not a proof of \textup{(i)} in all classes. The Borcherds-weight statement is independent: it follows from the constant-term identity in Proposition~\ref{prop:k3e-square-root-normalisation}.
\end{proof}

\subsection{Orbifold DT at \texorpdfstring{$(N, M) = (4, \mathrm{id})$}{(N, M) = (4, id)} and \texorpdfstring{$(6, \mathrm{id})$}{(6, id)}}
\label{subsec:k3e-lorgat-conj1-N46}

For \(N \in \{4,6\}\), the product action on \(K3\times E\) is again free after the \(E\)-translation. The Bryan--Steinberg~2014 crepant-resolution theorem and the orbifold topological vertex therefore enter only through the local quotient charts of the \(K3/\mathbb Z_N\) factor used to test the twisted-sector contribution. They do not by themselves prove the global \(E\)-reduced DT identity on the CHL threefold.

\begin{conjecture}[Orbifold DT at \texorpdfstring{$N \in \{4, 6\}$}{N in 4, 6}]
\label{thm:k3e-orbifold-DT-N46}
For $X^{(N)} = [K3 \times E / \mathbb{Z}_N]$ with $g_N$ Mukai-symplectic Nikulin and $N \in \{4, 6\}$,
\begin{equation}
\label{eq:k3e-DTred-N46}
Z^{\mathrm{red}}_{\mathrm{DT}}\bigl(X^{(N)}\bigr)(q, p, y) \;=\; -\,\frac{C_N(q, y)}{\bigl(F^{\mathrm{CHL}}_N(q, p, y)\bigr)^{2}}, \quad \mathrm{wt}\,F^{\mathrm{CHL}}_N \;=\; \frac{c_N(0)}{2} \;\in\; \bigl\{\tfrac32, 1\bigr\} \quad (N = 4, 6), \quad C_N \;=\; 2\phi^{(g_N)}_{0, 1}.
\end{equation}
Here $F^{\mathrm{CHL}}_N$ denotes the primitive CHL/BKM denominator. A Gritsenko--Cl\'ery row may be used only after a separate row-and-normalisation identification.
\end{conjecture}

\begin{remark}[Evidence for Conjecture~\ref{thm:k3e-orbifold-DT-N46}]
At $N = 4$, class $4B$ carries $4$ isolated points of type $A_3$ on $K3$ (Nikulin~$1979$ Theorem~$4.5$; Xiao~$1996$) with Mukai-invariant lattice rank $r_4 = 14$ and coinvariant signature $(0, 8)$. The crepant resolution $\widetilde{K3/\mathbb{Z}_4} \to [K3/\mathbb{Z}_4]$ is smooth by Bridgeland--King--Reid~$2001$; BCY applies to each $[\mathbb{C}^3/\mathbb{Z}_4]$ chart with orbifold topological vertex $V^{[\mathbb{C}^3/\mathbb{Z}_4]}$, and Jun Li degeneration glues four local contributions to the global $Z^{\mathrm{red}}_{\mathrm{DT}}[K3/\mathbb{Z}_4]$, with the $E$-translation tensor factor handled by Oberdieck--Pandharipande $K3 \times E$ reduced theory (Oberdieck~$2018$ Remark~$4$).

At $N = 6$, class $6A$ has cyclic quotient singularity type $2 A_5 + 2 A_2 + 2 A_1$ on $K3/\langle g_6\rangle$: Nikulin's fixed-point counts are $\#\mathrm{Fix}(g_6)=2$, $\#\mathrm{Fix}(g_6^2)=6$, and $\#\mathrm{Fix}(g_6^3)=8$, so the quotient has two stabiliser-$\mathbb{Z}_6$ orbits, two stabiliser-$\mathbb{Z}_3$ orbits, and two stabiliser-$\mathbb{Z}_2$ orbits. The Niemeier anchor is $6 D_4$ (umbral group $3.\mathrm{Sym}_6$ of order $2160$, selected by character match of $H^{(6)}$ with $\chi_{6 D_4}$; the alternative $4 A_5 + D_4$ has outer group of order $144$ and fails the umbral mock-modular criterion). Local BCY contributions at $[\mathbb{C}^3/\mathbb{Z}_6]$ charts are therefore indexed by the cyclic $A_5$, $A_2$, and $A_1$ charts; global gluing uses Jun Li degeneration along the elliptic fibration plus Bryan--Steinberg abelian composition. Mukai rank $r_6 = 14$ matches $r_4$: both $N = 4, 6$ symplectic-Nikulin classes sit in the rank-$14$ Mukai stratum.

\emph{Motivic refinement}: Maulik--Toda~$2018$ constructs motivic Donaldson--Thomas via Behrend cosection on K3-fibred CY$_3$; the cosection $\sigma$ descends to $[X^{(N)}]$ because $\omega_{K3}$ is $g_N$-invariant under symplectic Nikulin (Mukai~$1988$), and the $G$-equivariant Behrend function lifts by Kontsevich--Soibelman~$2008$ \S$4.4$ for finite abelian $G$. The motivic refinement $[\mathcal{M}^{[\mathbb{Z}_N]}_{\mathrm{DT}}]_{\mathrm{mot}} \in K_0(\mathrm{Var})[\mathbb{L}^{1/2}]$ specialises at $\mathbb{L} \to 1$ to the numerical zeta.

\emph{K-theoretic conditional step}: the untwined refined
\(K3\times E\) theory is the Oberdieck refined PT/Jacobi formula:
the symmetrised \(K\)-theoretic parameter \(y\) refines the reduced
stable-pair series to the refined Igusa form, and \(y=1\) recovers the
unrefined \(K3\times E\) scalar. This does not construct the
\(g_N\)-equivariant quotient theory. For \(N=4,6\) one must add a
\(g_N\)-linearisation of the reduced obstruction theory, a
twisted-sector \(K\)-theoretic correction, and compatibility with the
reduced \(E\)-fibred multiplicative structure. The CHPV twining data
identify the expected Jacobi input \(\phi^{(g_N)}_{0,1}\); they do not
by themselves prove the global \(K\)-theoretic orbifold DT identity.
\end{remark}

\begin{remark}[Niemeier-lattice versus K3 fixed-locus distinction]
\label{rem:k3e-niemeier-vs-fixed-locus}
The Niemeier lattice $N_6 = 6 D_4$ (rank-$24$ even unimodular, positive definite, global; Niemeier~$1973$ item~\#$19$ with outer automorphism modulo Weyl the umbral group $3.\mathrm{Sym}_6$ of order $2160$) and the cyclic quotient singularity type $2 A_5 + 2 A_2 + 2 A_1$ on $K3/\mathbb{Z}_6$ are distinct mathematical objects in distinct categories: the Niemeier lattice lives in $\mathbf{Lat}^{\mathrm{II}}_{24, 0}$, while the quotient singularity classification is a local analytic stratification on the complex surface $K3/\mathbb{Z}_6$. Both carry ADE labels; conflating them produces incorrect classifications. The Mukai~$1988$ Table~$2$ cyclic-order-$6$ entry has coinvariant Mukai rank $8$ and invariant rank $14$, embedding into $6 D_4$ via the Kondo~$1998$ refinement; Mukai does not label Niemeier anchors by root system $4 A_5 + D_4$.
\end{remark}

\subsection{The five anomalous classes and the \texorpdfstring{$H^3$}{H-3}-transgression obstruction}
\label{subsec:k3e-anomalous-H3-transgression}

Of the twenty-six $M_{24}$-conjugacy classes the correspondence-admissible ten $\{1A, 2A, 2B, 3A, 3B, 4A, 4B, 4C, 6A, 6B\}$ (Gaberdiel--Volpato~$2012$) lie in the Mukai symplectic slice; the five anomalous classes $\{7A, 7B, 11A, 23A, 23B\}$ carry a non-trivial restriction of the Mathieu anomaly class $[\alpha_{K3}] \in H^3(M_{24}, U(1))$ to the cyclic subgroup $\langle g \rangle$, obstructing the simultaneous $(g, h)$-twining lift.

\begin{proposition}[\texorpdfstring{$H^3$}{H-3}-transgression at the five anomalous classes]
\label{prop:k3e-H3-transgression}
Mason~$1994$ computes $H^2(M_{24}, U(1)) = \mathbb{Z}/12$; universal coefficient raises this to $H^3(M_{24}, U(1)) = \mathbb{Z}/12 \oplus \mathbb{Z}/2$, with Gaberdiel--Persson--Volpato~$2013$ identifying the $\mathbb{Z}/12$-generator with $[\alpha_{K3}]$. The restriction to the cyclic subgroup $\langle g \rangle$ gives
\begin{equation}
\label{eq:k3e-H3-restriction}
\mathrm{res}^{M_{24}}_{\langle g \rangle} [\alpha_{K3}] \;=\; \begin{cases}
0 & g \in \{1A, 2A, 2B, 3A, 3B, 4A, 4B, 4C, 6A, 6B\}, \\
3 \in \mathbb{Z}/7 & g \in \{7A, 7B\}, \\
2 \in \mathbb{Z}/11 & g = 11A, \\
1 \in \mathbb{Z}/23 & g \in \{23A, 23B\},
\end{cases}
\end{equation}
with fractional multiplier defects $\lambda_g \in \{3/7, 3/7, 2/11, 1/23, 1/23\}$ read as $U(1)$-roots under $H^3(C_N, U(1)) \cong \mu_N$. The simultaneous $(g, h)$-twined weak Jacobi form lifts to $\Gamma^{(g, h)}$ with trivial multiplier iff $\mathrm{res}^{M_{24}}_{\langle g \rangle} [\alpha_{K3}] = \mathrm{res}^{M_{24}}_{\langle h \rangle} [\alpha_{K3}] = 0$; the correspondence-core slice $(N, M) \in \{1, 2, 3, 4, 6\}^2$ is exactly the unobstructed locus.
\end{proposition}

\begin{proof}[Vanishing on the correspondence slice]
Each $N \in \{1, 2, 3, 4, 6\}$ divides $12$; the restriction map $\mathbb{Z}/12 \twoheadrightarrow \mathbb{Z}/N$ sends $[\alpha_{K3}]$ to $N \cdot [\alpha_{K3}] \bmod 12 = 0$ for these five values. The $\mathbb{Z}/2$-summand factors through $M_{24}^{\mathrm{ab}} = 0$ (Mathieu perfect). At $(2A, 3A)$: $\langle 2A, 3A \rangle = \langle 6A \rangle \cong C_6$ is generated by a Mukai automorphism and Gaberdiel--Volpato~$2012$ Theorem~$1.1$ row~$9$ lists $C_6$ as a maximal symplectic group with trivial $H^3$-image.
\end{proof}

\begin{theorem}[The \texorpdfstring{$55$}{55}-pair commuting-pair matrix on the correspondence core]
\label{thm:k3e-55-pair-commuting-matrix}
All $55$ commuting pairs in $\{1A, 2A, 2B, 3A, 3B, 4A, 4B, 4C, 6A, 6B\}^{[2]}$ admit a weak Jacobi form refinement of the K3 elliptic genus, a DMVV second-quantised Borcherds lift $\Phi^{(g, h)}_{10}{}^{-1}$, and agree with the universal Borcherds weight identity $\kappa_{\mathrm{BKM}}(\Phi_N) = c_N(0)/2$ across $N \in \{1, 2, 3, 4, 6\}$.
\end{theorem}

\begin{proof}[Proof]
The three inputs reduce to (i) the $H^3$-vanishing lemma (Proposition~\ref{prop:k3e-H3-transgression}) killing $[\omega_{g, h}]$; (ii) Persson--Volpato~$2012$ Proposition~$4.1$ converting vanishing cocycle to a weak Jacobi form; (iii) Gaberdiel--Volpato~$2012$ and Gaberdiel--Persson--Volpato~$2013$ cyclic reduction furnishing the explicit construction. The eight Gritsenko--Cl\'ery forms realise the $\mathrm{lcm}(N, M) \leq 4$ sub-slice; the remaining $47$ entries are Borcherds lifts at level $N' = \mathrm{lcm}(N, M) \in \{6, 12\}$ (Gritsenko--Nikulin~$1995$). Concrete witness at $(2A, 3A)$: $2A = 6A^3$, $3A = 6A^2$, so $(2A, 3A) \subset \langle 6A \rangle \cong \mathbb{Z}/6$ cyclic; $\Gamma^{(2A, 3A)}_0 = \Gamma_0(6)$; twined genus $Z^{(2A, 3A)}_{K3} = Z^{(6A)}_{K3}$ is a weak Jacobi form of weight $0$ index $1$ on $\Gamma_0(6)$ with multiplier $\psi_{6A}$ of order $6$, Fourier head $(c(-1, 1), c(0, 0), c(1, -1), c(1, 1)) = (2, 0, -2, 2)$.
\end{proof}

\subsection{The twenty-six-class sieve for \texorpdfstring{$(1, g)$}{(1, g)}-twinings}
\label{subsec:k3e-26class-sieve}

The $(1, g)$-twined extension of Conjecture~$1$ passes through a four-arrow sieve at each $M_{24}$ class: GHV~$2012$ weak Jacobi form, Gritsenko additive lift $k_g = c^{(g)}(0, 0)/2$, rank-$3$ BKM $\mathfrak{g}_g$, twined reduced DT.

\begin{theorem}[Four-arrow twining sieve across the \texorpdfstring{$26$}{26} \texorpdfstring{$M_{24}$}{M-24}-classes]
\label{thm:k3e-twined-26-class-sieve}
For each $g \in M_{24}$ with cycle shape $\pi_g$ of order $N_g$:
\begin{enumerate}[label=\textup{(\arabic*)}]
\item \emph{GHV twined elliptic genus}: $Z^{(g)}_{\mathrm{ell}}(K3; \tau, z) \in J^{w, \nu_g}_{0, 1}(\Gamma_0(N_g))$ with multiplier $\nu_g(\gamma) = e^{2\pi i \chi_g(\gamma)/(N_g \, h)}$ of order $h = N_g \cdot \gcd(N_g, 24)/24$ (GHV~$2012$ Theorem~$1$); at symplectic $g \in M_{23}$, $h = 1$; the normalisation $Z^{(g)}_{\mathrm{ell}}(\tau, 0) = \chi_g(V_{24}) = \mathrm{tr}(g | V_{24})$ with $V_{24}$ the $24$-dimensional permutation representation.
\item \emph{Gritsenko additive lift} $\mathrm{Lift}_{N_g}: J^{w, \nu_g}_{0, 1}(\Gamma_0(N_g)) \to M^{\mathrm{Siegel}}_{k_g}(\Gamma^+_{N_g}, \chi_g)$ with weight $k_g = c^{(g)}(0, 0)/2 \in \tfrac{1}{2}\mathbb{Z}$; at the ten geometrically admissible classes the weight tuple is $(k_{1A}, k_{2A}, k_{2B}, k_{3A}, k_{3B}, k_{4A}, k_{4B}, k_{4C}, k_{6A}, k_{6B}) = (10, 2, -2, 1, -2, 0, 2, -2, 0, 0)$; only $k_g \geq 1$ Gritsenko-lifts to a non-zero holomorphic paramodular form, selecting the paramodular-lift-alive sub-family $\{1A, 2A, 3A, 4B\}$.
\item \emph{Rank-$3$ BKM $\mathfrak{g}_g$}: at the four
paramodular-lift-alive classes $g \in \{1A, 2A, 3A, 4B\}$, of orders
$N \in \{1, 2, 3, 4\}$, the rank-$3$ BKM existence holds through the
four Cartan Gram determinants $\det A_g \in \{-32, -24, -18, -12\}$
(the values of Conjecture~\ref{prop:k3e-explicit-cartan} at
$N \in \{1, 2, 3, 4\}$); the fifth determinant $-6$ of that list
belongs to CHL level $N = 6$, whose sieve classes $6A, 6B$ have
$k_g = 0$ and are not paramodular-lift-alive, so it is not produced
by this sieve; at the non-paramodular-alive remaining six classes the GN~$1998$ Borcherds-product machinery does not directly produce $\mathfrak{g}_g$ because the Gritsenko lift is identically zero or of non-positive weight.
\item \emph{Twined reduced DT, conjectural beyond the known coefficients}: the expected identity is \(Z^{K3 \times E}_{\mathrm{red, DT}, g}(q, y, p) = -C_g / \Phi_{N_g}(g)\) when \(\Phi_{N_g}\) denotes the physical Igusa-square partition form, with \(C_g \in \mathbb{Q}[[q, y^{\pm 1}, p^{\pm 1}]] \cap M^{\mathrm{Fr}}_0(\Gamma_0(N_g))\); at \((1, 1A)\) it recovers \(Z^{K3 \times E}_{\mathrm{red, DT}}=-(\Phi_{10}^{\mathrm{OP}})^{-1}=-4096\,\Delta_5^{-2}\); at \(N_g \in \{2, 3\}\) Oberdieck--Pandharipande~$2019$ computes the first two Fourier coefficients.
\end{enumerate}
For admissible symplectic classes, the first three arrows are the
GHV weak-Jacobi construction, the Gritsenko additive lift, and the
rank-$3$ Borcherds product. The fourth arrow is a twined reduced-DT
expectation unless an Oberdieck--Pandharipande reduced-DT theorem has
already supplied the required Fourier coefficients. The anomalous five
\(\{7A, 7B, 11A, 23A, 23B\}\) obstruct the additive-lift and DT arrows
at the multiplier level by Proposition~\ref{prop:k3e-H3-transgression}:
\(\nu_g\) does not lift to a central extension compatible with
Gritsenko's additive machinery, although \(Z^{(g)}_{\mathrm{ell}}\)
still exists. The admissible-out-of-range classes
\(\{5A, 7A, 7B, 8A, 10A, 12A, 12B, 14A, 14B, 15A, 15B, 21A, 21B\}\)
with \(N_g \notin \{1,2,3,4,6\}\) obstruct the additive-lift arrow
because their levels fall outside the Gritsenko--Cl\'ery genus.
\end{theorem}

\begin{remark}[Sharpening of Lorgat~\texorpdfstring{$2020$}{2020} Conjecture~\texorpdfstring{$1$}{1}]
\label{rem:k3e-conj1-sharpening}
Theorem~\ref{thm:k3e-twined-26-class-sieve} sharpens Conjecture~$1$ from the qualitative ``eight Gritsenko--Cl\'ery forms $\leftrightarrow$ DT zeta roots $\leftrightarrow$ twined elliptic genera'' to a class-by-class sieve with explicit obstruction at each arrow, identifying the five anomalous classes with the non-trivial $H^3(M_{24}, U(1))$-restriction and the thirteen admissible-out-of-range classes with the paramodular-level escape $N_g \notin \{1, 2, 3, 4, 6\}$.
\end{remark}

\subsection{Five sharpenings: square-root, imprimitive, K-theoretic, non-CHL order, cross-lattice}
\label{subsec:k3e-lorgat-conj1-five-closures}

Five statements sharpen Theorem~\ref{thm:lorgat-conj1-N1-complete}, Proposition~\ref{thm:oberdieck-chl-N-2-3}, and Conjecture~\ref{thm:k3e-orbifold-DT-N46}: the CHL primitive-denominator weights and their Borcherds scalar-square weights, with OP/DT scalar equality a cited theorem at \(N=1\) and a conditional statement at higher CHL order; the imprimitive stable-pair multiple-cover gate, separated from the imprimitive BKM product exponent; the K-theoretic twisted-sector closure conjecture at $N \in \{4, 6\}$, separated from the untwined refined \(K3\times E\) theorem; the non-CHL Nikulin-order weights $(c_5(0), c_7(0), c_8(0))/2 = (2, \tfrac32, 1)$, with the CY$_3$ host construction separated from the Borcherds weight theorem; and the separation $\kappa_{\mathrm{BKM}}(\Delta_5) = 5 \neq 12 = \kappa_{\mathrm{BKM}}(\Phi_{12})$ between K3 paramodular and Fake-Monster Siegel inputs.

\begin{proposition}[CHL primitive-denominator weights and Borcherds scalar-square weights]
\label{prop:k3e-square-root-normalisation}
For $N \in \{1, 2, 3, 4, 6\}$ let $\phi^{\mathrm{CHL}}_{N}$ be the CHL twisted-twined weak Jacobi input of weight $0$ index $1$ on $\Gamma_0(N)$ in the David--Jatkar--Sen square-root normalisation, with Fourier head constant $c_N(0)$, and let
\[
F^{\mathrm{CHL}}_N \;=\; \mathrm{BorLift}^{\mathrm{prim}}_N\bigl(\phi^{\mathrm{CHL}}_{N}\bigr), \qquad \mathrm{wt}\,F^{\mathrm{CHL}}_N \;=\; c_N(0)/2,
\]
be the primitive CHL/BKM denominator form on the relevant paramodular group, with multiplier system at $N = 4$, and let
\[
\Phi_N^{\mathrm{scal}} \;=\; \mathrm{BorLift}^{\mathrm{Borch}}_N\bigl(2\phi^{\mathrm{CHL}}_{N}\bigr), \qquad \mathrm{wt}\,\Phi_N^{\mathrm{scal}} \;=\; c_N(0),
\]
be the scalar-square lift of the doubled input. Then
\begin{equation}
\label{eq:k3e-square-root-identity}
\bigl(c_N(0)\bigr)_{N = 1, 2, 3, 4, 6}=(10,6,4,3,2),\qquad
\bigl(\mathrm{wt}\,F^{\mathrm{CHL}}_N\bigr)=\bigl(5,3,2,\tfrac32,1\bigr),
\end{equation}
and the scalar square has weights $(10,6,4,3,2)$ (Dijkgraaf--Verlinde--Verlinde 1997; Jatkar--Sen~$2006$; the once-recorded ladder $(5,4,3,2,1)$ with constants $(10,8,6,4,2)$ is retracted). At $N=1$ the unnormalised identity is the proved Igusa--Gritsenko square
\[
\Phi_{10}^{\mathrm{un}}=\Delta_5^2=(F^{\mathrm{CHL}}_1)^2.
\]
For $N>1$, any identification of $F^{\mathrm{CHL}}_N$ with a Gritsenko--Cl\'ery catalogue form is an additional normalisation theorem, not a consequence of the tuple~\eqref{eq:k3e-square-root-identity}. In particular, the Gritsenko--Cl\'ery eight-form catalogue is triple-indexed and independent from the CHL constant tuple unless a named overlap theorem is supplied.
\end{proposition}

\begin{proof}
The universal Borcherds weight identity gives
$\kappa_{\mathrm{BKM}}(F^{\mathrm{CHL}}_N)=c_N(0)/2$ on the CHL ladder, hence the primitive weights $(5,3,2,\tfrac32,1)$ from the constants $(10,6,4,3,2)$ (Jatkar--Sen~$2006$; Govindarajan--Krishna~2010). Doubling the Jacobi input doubles the weight and gives the scalar-square weights. At $N=1$, Gritsenko--Nikulin~$1998$ identifies the primitive denominator with $\Delta_5$ and Igusa's normalisation gives $\Phi_{10}^{\mathrm{un}}=\Delta_5^2$. The remaining assertions are weight and normalisation statements: they do not prove a blanket equality $\Phi_N^{\mathrm{scal}}=(F_N^{\mathrm{GC}})^2$ and do not identify the Gritsenko--Cl\'ery catalogue constants with the CHL constants.
\end{proof}

\begin{remark}[Explicit weight tuple and constant-term tuple]
\label{rem:k3e-square-root-tuple}
The ladder reads
\[
\bigl(c_N(0)\bigr)_{N = 1, 2, 3, 4, 6} \;=\; (10, 6, 4, 3, 2), \qquad
\bigl(\mathrm{wt}\,F^{\mathrm{CHL}}_N\bigr) \;=\; \bigl(5, 3, 2, \tfrac32, 1\bigr), \qquad
\bigl(\mathrm{wt}\,\Phi_N^{\mathrm{scal}}\bigr) \;=\; (10, 6, 4, 3, 2).
\]
The universal identity $\kappa_{\mathrm{BKM}}(\Phi_N) = c_N(0)/2$ belongs to the primitive denominator, not to the OP/DT scalar square. The Gritsenko--Cl\'ery catalogue has its own triple-indexed constant tuple and may overlap this ladder only at explicitly named and normalised entries.
\end{remark}

\begin{proposition}[Imprimitive stable-pair gate and Borcherds exponent separation]
\label{thm:k3e-imprimitive-DT-OP2014-B}
The statement has two independent lanes. The Borcherds product exponent
is a theorem. The reduced stable-pair multiple-cover formula is the
Oberdieck--Pandharipande all-class conjectural input, and the
Behrend-weighted DT closure at imprimitive \(\beta\) requires the
motivic wall-crossing extension named in the proof.

Let $\beta \in H_2(K3, \mathbb{Z})$ with $\beta = m \beta_h$, $m \geq 1$, and $\beta_h$ primitive of degree $h$. Assume the Oberdieck--Pandharipande reduced stable-pair multiple-cover rule for \(K3\times E\) in class \((\beta,d[E])\). Then the reduced Pandharipande--Thomas generating functions satisfy
\begin{equation}
\label{eq:k3e-imprimitive-PT}
Z^{\mathrm{red}}_{\mathrm{PT}}(K3 \times E; \beta, d) \;=\; \sum_{k \mid \gcd(m, d)} \mu(k)\, k^{-1}\, Z^{\mathrm{red}}_{\mathrm{PT}}(K3 \times E; \beta/k, d/k)\bigl|_{q \mapsto q^k, y \mapsto y^k, p \mapsto p^k},
\end{equation}
with primitive-class base \(m = 1\) reducing to Theorem~\ref{thm:lorgat-conj1-N1-complete}:
\[
Z^{\mathrm{red}}_{\mathrm{PT}}(K3 \times E; \beta_h, d)
=-(\Phi_{10}^{\mathrm{OP}})^{-1}
=-4096\,\Delta_5^{-2}.
\]
Independently, the Borcherds-product signed-character identity reads, for every positive \(\alpha = (n, \ell, m) \in \Lambda^{2, 1}_{\mathrm{II}}\),
\begin{equation}
\label{eq:k3e-imprimitive-mult}
\operatorname{sdim}\mathfrak g_{\Delta_5,\alpha}
\;=\;
c(4nm-\ell^2),
\end{equation}
where \(c(D)\) is the Eichler--Zagier coefficient of \(\phi_{0,1}\).
The Möbius sum in \eqref{eq:k3e-imprimitive-PT} belongs to the reduced
stable-pair multiple-cover extraction, not to the BKM product exponent.
\end{proposition}

\begin{proof}
Under the stated Oberdieck--Pandharipande multiple-cover rule, the reduced PT moduli space \(P_n(K3 \times E,(\beta,d))\) decomposes into \(E\)-fixed components indexed by effective curve class, and the isogeny
\[
\pi_k: K3 \times E \longrightarrow K3 \times E,\qquad z\longmapsto kz
\]
on the elliptic factor gives the multiple-cover map
\[
\rho_k: P_{n/k}\bigl(K3 \times E, (\beta/k, d/k)\bigr) \longrightarrow P_n\bigl(K3 \times E, (\beta, d)\bigr), \quad \mathcal{F} \longmapsto \pi^*_k \mathcal{F},
\]
and the reduced measure pulls back with weight \(1/k\). Möbius inversion on the divisor sum over \(k \mid \gcd(m,d)\) gives~\eqref{eq:k3e-imprimitive-PT}. This is the stable-pair lane; it is not a Borcherds-root multiplicity computation.

The imprimitive signed-character identity~\eqref{eq:k3e-imprimitive-mult} follows directly from the Borcherds product expansion of \(\Delta_5^2\):
\[
 \Delta_5^2
 =
 e^{2\rho}
 \prod_{(n,\ell,m)>0}
 (1-q^ny^\ell p^m)^{c(4nm-\ell^2)}.
\]
At \((n,\ell,m)=(2,2,2)\) the discriminant is \(12\), hence the product
contains
\[
 (1-q^2y^2p^2)^{c(12)}=(1-q^2y^2p^2)^{4016}.
\]
The value \(c(3)=-64\) belongs to the primitive point \((1,1,1)\);
it is not subtracted from the root exponent at \((2,2,2)\). Möbius
inversion appears in the reduced PT multiple-cover formula
\eqref{eq:k3e-imprimitive-PT}, where it extracts primitive curve
contributions from imprimitive stable pairs.

The Behrend-weighted DT identification at imprimitive classes is a further conditional step: Joyce--Song~$2012$ \S$6.3$ proves Behrend--PT equivalence on CY$_3$ in the primitive setting used here; extension to imprimitive classes requires the Kontsevich--Soibelman motivic wall-crossing formula at unstable orbit points (KS~$2008$ \S$6.4$), with the required \(K3\times E\) reduced obstruction and \(E\)-quotient conventions fixed simultaneously.
\end{proof}

\begin{conjecture}[K-theoretic closure at \texorpdfstring{$N \in \{4, 6\}$}{N in 4, 6}]
\label{prop:k3e-K-theoretic-closure-N46}
Let $Z^{\mathrm{K,red}}_{\mathrm{DT}}(X^{(N)})$ denote the conjectural
\(g_N\)-equivariant reduced \(K\)-theoretic Donaldson--Thomas
generating function on
\(X^{(N)}=[K3\times E/\mathbb Z_N]\), extending the untwined refined
reduced \(K3\times E\) stable-pair theory of Oberdieck~$2018$. For
$N \in \{4, 6\}$ with $g_N \in M_{24}$ a Mukai-symplectic Nikulin
automorphism,
\begin{equation}
\label{eq:k3e-K-theoretic-N46}
Z^{\mathrm{K, red}}_{\mathrm{DT}}\bigl(X^{(N)}\bigr)(y, q, p, z) \;=\; -\,\frac{\widetilde C^{\mathrm{K}}_N(y, q, z)}{\widetilde\Phi_{k_N}^{\mathrm{K}}(y, q, p, z)^{2}}, \qquad k_N \;=\; c_N(0)/2,
\end{equation}
where $\widetilde\Phi^{\mathrm{K}}_{k_N}$ is the K-theoretic
paramodular form deforming the primitive CHL/BKM denominator
$F^{\mathrm{CHL}}_N$ and the twisted-sector correction
$\widetilde C^{\mathrm{K}}_N$ is the K-theoretic lift of the
Cheng--Harrison--Paquette--Volpato~$2014$ $g_N$-twining coefficient.
At Euler-specialisation \(y\to1\), this conjecture reduces to
Conjecture~\ref{thm:k3e-orbifold-DT-N46}. The K-theoretic enhancement
requires the equivariant quotient construction and is not supplied by
the untwined refined \(K3\times E\) theorem.
\end{conjecture}

\begin{remark}[Euler-specialisation gate]
The Euler map from the symmetrised virtual structure sheaf to the
Behrend-weighted reduced numerical invariant would send
\(Z^{\mathrm{K,red}}_{\mathrm{DT}}(X^{(N)})\) to
\(Z^{\mathrm{red}}_{\mathrm{DT}}(X^{(N)})\). If the conjectural
\(g_N\)-equivariant \(K\)-theoretic construction is compatible with
this Euler map and with the reduced \(E\)-quotient convention, then
\(y=1\) in \eqref{eq:k3e-K-theoretic-N46} gives precisely
\eqref{eq:k3e-DTred-N46}. Thus the Euler-specialisation check is a
necessary compatibility condition for Conjecture~\ref{prop:k3e-K-theoretic-closure-N46};
it is not a proof of Conjecture~\ref{thm:k3e-orbifold-DT-N46}.
The missing step is the global compatibility of the CHPV twining
character \(\psi^{(g_N)}:K^0_{\mathbb Z_N}(K3)\to K^0(K3^{g_N})\)
with the reduced \(E\)-fibred multiplicative structure.
\end{remark}

\begin{proposition}[Non-CHL Nikulin orders \texorpdfstring{$N \in \{5, 7, 8\}$}{N in 5, 7, 8}]
\label{prop:k3e-half-integer-metaplectic}
This statement concerns the Borcherds weight formulae only; it does
not assert a CY$_3$ host.
The three non-CHL entries of the order-indexed Nikulin extension have
Borcherds weights
\begin{equation}
\label{eq:k3e-half-integer-weights}
\bigl(\kappa_{\mathrm{BKM}}(\Phi_N)\bigr)_{N = 5, 7, 8} \;=\; \bigl(c_N(0)/2\bigr)_{N = 5, 7, 8} \;=\; \bigl(2,\, \tfrac32,\, 1\bigr),
\end{equation}
with $c_N(0)=(4,3,2)$, in the Mathieu-twined normalisation. These are not the half-integral
Gritsenko--Cl\'ery dd-modular rows: the order-$7$ weight $\tfrac32$
carries a multiplier system but is indexed by Nikulin order, not by a
dd-modular triple.
\end{proposition}

\begin{proof}
$N=5$: the $M_{24}$ frame shape $1^4 5^4$ has ones-exponent
$a_1=4$, hence
$c_5(0)=4$ and
$\kappa_{\mathrm{BKM}}(\Phi_5)=2$ by Borcherds~$1998$ Theorem~$13.3$.

$N=7$: the frame shape $1^3 7^3$ has ones-exponent $a_1=3$, hence
$c_7(0)=3$ and $\kappa_{\mathrm{BKM}}(\Phi_7)=\tfrac32$, half-integral
with multiplier system.

$N=8$: the frame shape $1^2 2\,4\,8^2$ has ones-exponent $a_1=2$,
hence $c_8(0)=2$ and $\kappa_{\mathrm{BKM}}(\Phi_8)=1$.

The CY$_3$ host identification is conjectural: at $N=5$ a
Borcea--Voisin CY$_3$ is the candidate (Borcea~$1997$; Voisin~$1993$);
at $N=7$ Nikulin~$1979$ Corollary~$3.3$ obstructs a smooth diagonal
$K3 \times E/\mathbb Z_7$ quotient; at $N=8$ the
Mongardi--Tari--Wandel~$2018$ Kummer-$3$ fourfold is a proposed host
outside the CY$_3$ class. The half-integral Gritsenko--Cl\'ery rows
are instead $(4,1;\tfrac12)$ and $(1,4;\tfrac32)$, as in
Theorem~\ref{thm:k3e-gc-eight-commuting-pair}.
\end{proof}

\begin{remark}[Cross-volume \texorpdfstring{$\kappa_{\mathrm{BKM}}$}{kappa-BKM} at two Siegel inputs]
\label{rem:k3e-cross-volume-kappa-BKM-five-vs-twelve}
The universal identity $\kappa_{\mathrm{BKM}}(\Phi) = c_\Lambda(0)/2$ of Borcherds~$1995$ \emph{Invent.~Math.}~$120$ Theorem~$10.4$ produces two distinct evaluations on two distinct BKM superalgebras:
\begin{equation}
\label{eq:k3e-two-kappa-BKM-siblings}
\kappa_{\mathrm{BKM}}\bigl(\mathfrak{g}_{\Delta_5}\bigr) \;=\; 5 \qquad \text{on the paramodular lattice } \Lambda^{2, 1}_{\mathrm{II}},
\end{equation}
versus
\begin{equation}
\kappa_{\mathrm{BKM}}\bigl(\mathfrak{g}_{\Phi_{12}}\bigr) \;=\; 12 \qquad \text{on the Fake-Monster lattice } \mathrm{II}_{25, 1}.
\end{equation}
The first evaluation uses the paramodular Borcherds lift of the normalized input $\phi_{0, 1}$ (Gritsenko--Nikulin~$1998$ Theorem~$2.1$; $c_1(0) = 10$, weight of $\Delta_5$ is $5$); the doubled K3 elliptic genus \(2\phi_{0,1}\) gives the Igusa-square partition lane. The second uses the $\mathrm{II}_{25, 1}$ Borcherds lift of the Leech-theta character (Borcherds~$1990$ \emph{Adv.~Math.}~$83$; $c_\Lambda(0) = 24$, weight of $\Phi_{12}$ is $12$). The two superalgebras are different: $\mathfrak{g}_{\Delta_5}$ is rank-$3$ with Weyl group $W^{(2)}(\Lambda^{2, 1}_{\mathrm{II}})$ and signed root-character exponents given by $\phi_{0, 1}$-Fourier coefficients, corresponding to the K3-chiral sibling of the CY-to-chiral functor at $d = 3$ on $K3 \times E$; $\mathfrak{g}_{\Phi_{12}}$ is rank-$26$ with Weyl group the reflection group of $\mathrm{II}_{25, 1}$, uniform root multiplicity $24$, corresponding to the Fake-Monster sibling at $d = 5$ (the $\mathrm{II}_{25, 1}$ Lorentzian lattice sits in the CY-dimension-stratified sibling census at $d = 5$, Borcherds~$1990$). No additive shift of the form $\kappa_{\mathrm{BKM}}(\Phi) = \kappa_{\mathrm{ch}}^{\mathrm{Hodge}} + \chi(\mathcal{O}_{\mathrm{fibre}})$ relates $5$ and $12$; the identity $\kappa_{\mathrm{BKM}}(\Phi) = c_\Phi(0)/2$ is evaluated at two different Siegel input denominators. The sibling census is dimension-stratified: $d = 3$ (K3 $\times$ E, $\mathfrak{g}_{\Delta_5}$, $\kappa_{\mathrm{BKM}}(\Delta_5) = 5$); $d = 5$ (Fake Monster, $\mathfrak{g}_{\Phi_{12}}$, $\kappa_{\mathrm{BKM}}(\Phi_{12}) = 12$); intermediate $d = 4$ bridged by Conway / Leech.
\end{remark}

% ======================================================================
\section{The Jacobi theta-component refinement of \texorpdfstring{$c_N(0)$}{c-N(0)}}
\label{sec:k3e-jacobi-theta-refinement}
% ======================================================================

The twined constant $c_N(0)$ controlling $\kappa_{\mathrm{BKM}}(\Phi_N) = c_N(0)/2$ in the Mathieu-twined normalisation is neither $\chi_{g_N}(\mathcal{O}_{K3})$ (identically $2$ for symplectic $g_N$ by Mukai~$1988$ Theorem~$0.6$) nor the raw frame-exponent sum $T_{H^*}(g_N)$; it is the even theta-component of the twined elliptic genus.

\begin{theorem}[Jacobi theta-component refinement; Mathieu-twined normalisation]
\label{thm:k3e-jacobi-theta-refinement}
For $N \in \{1, 2, 3, 4, 6\}$, $g_N \in \mathrm{Aut}_{\mathrm{symp}}(K3)$ a Nikulin symplectic automorphism of order $N$, $Z^{(g_N)}_{K3}$ the $g_N$-twined K3 elliptic genus of weight $0$ and index $1$, and $T_{H^*}(g_N) = \sum_a m_a$ the frame-exponent sum of the $M_{24}$ frame shape $\pi_{g_N} = \prod_a a^{m_a}$, the index-$1$ theta decomposition $Z^{(g_N)}_{K3}(\tau, z) = h^{(N)}_0(\tau) \theta_{1, 0}(\tau, z) + h^{(N)}_1(\tau) \theta_{1, 1}(\tau, z)$ with $A_N := [q^0] h^{(N)}_1(\tau)$ yields
\begin{equation}
\label{eq:k3e-jacobi-theta-refinement}
c_N(0) \;=\; T_{H^*}(g_N) \;-\; 2 A_N, \qquad \kappa_{\mathrm{BKM}}(\Phi_N) \;=\; \tfrac{1}{2}\bigl(T_{H^*}(g_N) - 2 A_N\bigr),
\end{equation}
with tabulated values $(A_N)_{N = 1, 2, 3, 4, 6} = (2, 4, 3, 3, 3)$ and
\[
\bigl(c_N(0)\bigr)_{N = 1, 2, 3, 4, 6} \;=\; (20, 8, 6, 4, 2), \qquad \bigl(\kappa_{\mathrm{BKM}}(\Phi_N)\bigr)_{N = 1, 2, 3, 4, 6} \;=\; (10, 4, 3, 2, 1).
\]
The coexisting square-root ladder carries constants $(10, 6, 4, 3, 2)$ and weights $(5, 3, 2, \tfrac32, 1)$ (Jatkar--Sen; Govindarajan--Krishna); no single family carries a mixed tuple drawn from both, and the once-recorded mixed ladder $(10, 8, 6, 4, 2)$ with weights $(5, 4, 3, 2, 1)$ is retracted.
\end{theorem}

\begin{proof}
\emph{Frame-exponent sum $=$ $z = 0$ specialisation}: the $g_N$-twined K3 elliptic genus satisfies $Z^{g_N}_{K3}(\tau, 0) = T_{H^*}(g_N)$, the $y = 1$ specialisation recovering the full graded trace datum of the frame shape. The Mukai frame shapes $\pi_{g_N} = 1^{24}, 1^8 2^8, 1^6 3^6, 1^4 2^2 4^4, 1^2 2^2 3^2 6^2$ give frame-exponent sums $T_{H^*}(g_N) = 24, 16, 12, 10, 8$ at $N = 1, 2, 3, 4, 6$. The geometric fixed-locus counts are $\chi(K3^{g_N}) = 24, 8, 6, 4, 2$, coinciding with the ones-exponents $a_1$ at $N \geq 2$; the quantity $T_{H^*}$ sums the \(M_{24}\)-frame exponents \(\sum_a m_a\) in this notation.

\emph{theta decomposition}: every weak Jacobi form of weight $0$ and index $1$ admits the two-term decomposition with $h_0, h_1$ weakly-holomorphic vector-valued of weight $-1/2$ on $\Gamma_0(N)$ (Eichler--Zagier~$1985$ \S$5$; Gritsenko~$1999$ \S$4$). The $q^0$-extraction of $Z^{(g_N)}_{K3}$ unpacks as $[\zeta^{-1} + \zeta]$-symmetric pair with coefficient $A_N$ and $[\zeta^0]$-term with coefficient $c_N(0)$, yielding the identity $Z^{(g_N)}_{K3}(\tau = i\infty, z = 0) = c_N(0) + 2 A_N = T_{H^*}(g_N)$.

\emph{$A_N$ tabulation}: direct substitution from the $T_{H^*}$ and $c_N(0)$ tables gives $A_1 = (24 - 20)/2 = 2$, $A_2 = (16 - 8)/2 = 4$, $A_3 = (12 - 6)/2 = 3$, $A_4 = (10 - 4)/2 = 3$, $A_6 = (8 - 2)/2 = 3$. The $c_N(0)$ values are the frame-shape ones-exponents $a_1 = (8, 6, 4, 2)$ at $N \geq 2$ (Eguchi--Hikami~$2011$ at $N = 2$; Cheng--Duncan--Harvey~$2014$ Table~$3$) and the $q^0\zeta^0$-coefficient $20$ of $\mathrm{Ell}(K3) = 2\phi_{0,1}$ at $N = 1$.

\emph{Borcherds weight}: Borcherds~$1998$ Theorem~$13.3$ gives $\mathrm{wt}(\Phi_N) = c_N(0)/2$; substituting $c_N(0) = T_{H^*}(g_N) - 2 A_N$ closes the claim.
\end{proof}

\begin{remark}[Supertrace interpretation]
\label{rem:k3e-jacobi-theta-supertrace}
The refinement $c_N(0) = T_{H^*}(g_N) - 2 A_N$ reads as a $\mathbb{Z}/2$-graded Lefschetz supertrace
\[
c_N(0) \;=\; \mathrm{str}\bigl(g_N^* \mid H^*_{\mathrm{Muk}}(K3)^{\mathrm{even}\,\theta} \ominus H^*_{\mathrm{Muk}}(K3)^{\mathrm{odd}\,\theta}\bigr),
\]
with the index-$1$ Jacobi parity splitting the rank-$24$ Mukai lattice into an even-sector of dimension $22$ and an odd-sector of dimension $2$ at $N = 1$: $24 = 22 + 2$, $22 - 2 = 20 = c_1(0)$. The ratio $A_N/T_{H^*}(g_N) = (1/12, 1/4, 1/4, 3/10, 3/8)$ is non-constant, reflecting the Mathieu-moonshine mock-modular structure of $h^{(N)}_1$; the correct refinement is additive in the theta-split, not multiplicative in $T_{H^*}$.
\end{remark}

\begin{remark}[Two normalisations for \texorpdfstring{$c_N(0)$}{c-N(0)}: Mathieu-twined and square-root]
\label{rem:k3e-two-conventions-cN}
The symbol $c_N(0)$ carries two honest normalisations on the CHL ladder $N \in \{1, 2, 3, 4, 6\}$, each with its own constant-term tuple. The Mathieu-twined normalisation reads $c_N(0) := [\zeta^0 q^0]\, Z^{(g_N)}_{K3}(\tau, z)$ off the $g_N$-twined K3 elliptic genus ($Z_{K3} = 2\phi_{0,1}$ at $N = 1$, in the EZ normalisation $\phi_{0,1}(\tau, z) = y^{-1} + 10 + y + O(q)$), giving $c_N(0) = (20, 8, 6, 4, 2)$ and twined weights $c_N(0)/2 = (10, 4, 3, 2, 1)$. The square-root normalisation reads $c_N(0)$ off the David--Jatkar--Sen twisted-twined input, giving $c_N(0) = (10, 6, 4, 3, 2)$ and the primitive BKM denominator weights $\kappa_{\mathrm{BKM}}(\Phi_N) = c_N(0)/2 = (5, 3, 2, \tfrac32, 1)$; at $N = 1$ this is the half-genus reading $c_1(0) = 10$ of $\phi_{0,1}$. The two normalisations are never mixed within one ladder; the once-recorded tuple $(10, 8, 6, 4, 2)$ with weights $(5, 4, 3, 2, 1)$, which combined the half-genus $N = 1$ slot with the twined $N \geq 2$ slots, is retracted. The physics Igusa-square convention $\Phi^{\mathrm{phys}}_{k_N} = (\Delta_{k(N)}^{(N)})^2$ (Dijkgraaf--Verlinde--Verlinde 1997; Jatkar--Sen 2006 \emph{JHEP}~11, 072) records the squared form of weight $k_N^{\mathrm{phys}} = 2 k(N) = c_N(0) = (10, 6, 4, 3, 2)$ in the square-root normalisation; the BKM denominator is the square-root, of weight $k(N)=\kappa_{\mathrm{BKM}}(\Phi_N)=c_N(0)/2$.
\end{remark}

% ======================================================================
\section{The eight Gritsenko--Cl\'ery forms by commuting pair}
\label{sec:k3e-eight-form-commuting-pair}
% ======================================================================

The eight Gritsenko--Cl\'ery dd-Siegel paramodular forms (Gritsenko--Cl\'ery~$2013$ Theorem~$1.1$) bijectively match commuting-pair orbits $(g_N, h_M) \in M_{24}^{[2]}/\mathrm{conj}$ under three simultaneous conditions: the paramodular-level bound $tN \leq 4$, the Mukai geometric realisability $(g_N, h_M) \in \mathrm{Muk}^{[2]}_{K3}$, and the transgression triviality $\tau_{g_N}[\alpha_{K3}] = 0 \in H^2(C_{M_{24}}(g_N), U(1))$.

\begin{theorem}[Gritsenko--Cl\'ery eight forms \texorpdfstring{$\leftrightarrow$}{<->} eight commuting-pair orbits]
\label{thm:k3e-gc-eight-commuting-pair}
The eight Gritsenko--Cl\'ery forms pair with the eight admissible $M_{24}$-commuting-pair orbits on the correspondence slice via
\begin{center}
\small
\renewcommand{\arraystretch}{1.15}
\begin{tabular}{c|l|c|c|c|l}
\hline
\# & Form & $t$ & $N_{\mathrm{lvl}}$ & $k$ & Commuting pair $(g_N, h_M)$ \\
\hline
$1$ & $M_5(\Gamma_1, \nu_2) = \Delta_5$ & $1$ & $1$ & $5$ & $(1A, 1A)$ \\
$2$ & $M_2(\Gamma_2, \nu_4)$ & $2$ & $1$ & $2$ & $(2A, 1A)$ Mukai $8 A_1$ fixed \\
$3$ & $M_3(\Gamma_1(2), \nu_2)$ & $1$ & $2$ & $3$ & $(1A, 2A)$ swap of $\#2$ \\
$4$ & $M_1(\Gamma_3, \nu_6)$ & $3$ & $1$ & $1$ & $(3A, 3A)$ diagonal $\mathbb{Z}/3$ \\
$5$ & $M_2(\Gamma_1(3), \nu_2)$ & $1$ & $3$ & $2$ & $(1A, 3A)$ cyclic reduction \\
$6$ & $M_{1/2}(\Gamma_4, \nu_8)$ & $4$ & $1$ & $1/2$ & $(4A, 1A)$ metaplectic \\
$7$ & $M_{3/2}(\Gamma_1(4), \nu_4)$ & $1$ & $4$ & $3/2$ & $(2A, 4A)$ mixed order \\
$8$ & $M_1(\Gamma_2(2), \nu_4)$ & $2$ & $2$ & $1$ & $(2A, 2A)$ Mukai diagonal Nikulin \\
\hline
\end{tabular}
\end{center}
with $c_{(g, h)}(0) = 2 k$ the Borcherds-weight column $(10, 4, 6, 2, 4, 1, 3, 2)$.
The table proves the eight-pair enumeration and the
\(\kappa_{\mathrm{BKM}}\)-column match. A CY-D functorial interpretation
of the same eight rows requires a separate functorial construction from
the indicated K3 data to the corresponding dd-Siegel automorphic
objects.
\end{theorem}

\begin{proof}[Row enumeration]
Rational pairs $(t, N)$ with $tN \leq 4$ and $t, N \in \mathbb{Z}_{\geq 1}$ produce exactly the eight: $\{(1, 1), (1, 2), (1, 3), (1, 4), (2, 1), (2, 2), (3, 1), (4, 1)\}$. Each $g_N$ of order $N \in \{1, 2, 3, 4\}$ lies in Mukai~$1988$ Table~$2$; each $h_M$ of order $M \in \{1, 2, 3, 4\}$ likewise; all eight $(g_N, h_M)$ commute in $M_{24}$ (Gaberdiel--Hohenegger--Volpato~$2012$ centraliser tables). The $H^3$-triviality follows from $\mathrm{lcm}(N, M) \in \{1, 2, 3, 4\}$ dividing $12 = |\alpha_{K3}|$. No additional $(t, N)$ with $tN = 5, 6, \ldots$ satisfies $H^3$-triviality against $\mathrm{lcm}(N, M) \mid 12$ with $M \in \{1, 2, 3, 4, 6\}$ while remaining in Mukai's admissible orders $N \leq 4$.

Row $\#6$ is the lone theta-metaplectic row: form $M_{1/2}(\Gamma_4, \nu_8)$ has $(k, t, N_{\mathrm{lvl}}, |\nu|) = (1/2, 4, 1, 8)$. The commuting pair is $(g_4, h_1) = (4A, \mathrm{id})$ with Mukai $g_4$ of order $4$ realised as the K3 symplectic automorphism with $|\mathrm{Fix}(g_4)| = 4$ fixed points; the multiplier $\nu_8$ is the theta-multiplier of order $8 = 2 \cdot \mathrm{lcm}(4, 1)$, i.e.\ the square root of the standard Maass multiplier $v_{\Delta_5}^{1/2}$ on the metaplectic cover $\mathrm{Mp}_2(\mathbb{Z})$ (Gritsenko--Nikulin~$1998$ \emph{Amer.\ J.\ Math.}~$119$ Theorem~$1.5$). The BKM Borcherds weight $\kappa_{\mathrm{BKM}}(M_{1/2}(\Gamma_4, \nu_8)) = c_4(0)/2 = 1/2$ is one of exactly two half-odd-integer BKM weights on the correspondence slice, the other being the weight $\tfrac32$ of row~$\#7$; within the square-root CHL ladder $(5, 3, 2, \tfrac32, 1)$ the unique half-odd-integer weight is instead $\tfrac32$ at $N = 4$. Its fermion $\mathbb{Z}/2$-graded simple roots produce the Gritsenko--Cl\'ery-boundary generalised BKM $\mathfrak{g}_{\Phi_4}$ with exactly one imaginary odd simple root.
\end{proof}

\begin{remark}[Paramodular level \texorpdfstring{$t = 4$}{t = 4} as the \texorpdfstring{$2$}{2}-divisibility selector]
\label{rem:k3e-paramodular-t4-selector}
The subset $\{N \in \{4, 6, 8\} : 2 \mid N\} \cap \{N : N \leq 8\}$ coincides with the set $\{N : t_{\max}(N) = 4\}$, where $t_{\max}(N)$ is the maximal paramodular level supporting a Gritsenko--Cl\'ery dd-modular form at $\Gamma_t$ of $N$-twist. Three independent selectors produce this set: the algebraic selector $t = 4$ (paramodular level); the geometric shadow $\mathbb{Z}/N$-orbifold fixed-locus chain $\mathrm{Fix}(g_N) \subset \mathrm{Fix}(g_{N/2})$ non-trivial iff $2 \mid N$; the physical shadow CHL duality frame at the metaplectic-cover boundary (the $E[2]$-translation obstruction on $T^2$). The Mukai-bound saturation at $N = 8$ (two Niemeier embeddings $8A, 8B$) is the closure of the orbit, not a separate boundary phenomenon.
\end{remark}

% ======================================================================
\section{Explicit Cartan data at \texorpdfstring{$N \in \{1, 2, 3, 4, 6\}$}{N in 1, 2, 3, 4, 6}}
\label{sec:k3e-explicit-cartan-gN}
% ======================================================================

The rank-$3$ BKM superalgebras $\mathfrak{g}^{(N)}$ attached to the doubly-twined Cl\'ery--Gritsenko paramodular form at CHL level $N \in \{1, 2, 3, 4, 6\}$ carry explicit Gram and Cartan data read from the $g_N$-invariant signature-$(2, 1)$ sublattice.

\begin{conjecture}[Explicit Cartan data]
\label{prop:k3e-explicit-cartan}
With real simple roots $(\delta_1^{(N)}, \delta_2^{(N)}, \delta_3^{(N)})$ spanning the $g_N$-invariant signature-$(2, 1)$ sublattice of $\Lambda^{2, 1}_{\mathrm{II}, N}$,
\[
G^{(1)} = \begin{pmatrix} 2 & -2 & -2 \\ -2 & 2 & -2 \\ -2 & -2 & 2 \end{pmatrix}, \qquad
G^{(2)} = \begin{pmatrix} 2 & -2 & -2 \\ -2 & 4 & -2 \\ -2 & -2 & 4 \end{pmatrix}, \qquad
G^{(3)} = \begin{pmatrix} 2 & -1 & -2 \\ -1 & 2 & -2 \\ -2 & -2 & 2 \end{pmatrix},
\]
\[
G^{(4)} = \begin{pmatrix} 2 & -1 & -2 \\ -1 & 2 & -2 \\ -2 & -2 & 4 \end{pmatrix}, \qquad
G^{(6)} = \begin{pmatrix} 2 & -1 & -2 \\ -1 & 2 & -2 \\ -2 & -2 & 6 \end{pmatrix}.
\]
Determinants $(\det G^{(N)})_{N = 1, 2, 3, 4, 6} = (-32, -24, -18, -12, -6)$; signatures $(2, 1)$ throughout; indecomposability by inspection. Weyl vectors $\rho^{(N)}$ solving $G^{(N)} \rho = -\tfrac{1}{2} \mathrm{diag}(G^{(N)})$ are
\[
\rho^{(1)} = \tfrac{1}{2}(\delta_1 + \delta_2 + \delta_3), \quad \rho^{(2)} = \tfrac{5}{2}\delta_1 + \tfrac{3}{2}\delta_2 + \tfrac{3}{2}\delta_3, \quad \rho^{(3)} = \tfrac{2}{3}\delta_1 + \tfrac{2}{3}\delta_2 + \tfrac{5}{6}\delta_3,
\]
\[
\rho^{(4)} = 2\delta_1 + 2\delta_2 + \tfrac{3}{2}\delta_3, \qquad \rho^{(6)} = 6\delta_1 + 6\delta_2 + \tfrac{7}{2}\delta_3,
\]
with norms $(\rho^{(N)}, \rho^{(N)})_{N = 1, 2, 3, 4, 6} = (-3/2, -17/2, -13/6, -7, -45/2)$. The non-uniform third-diagonal $G^{(N)}_{33} \in \{2, 4, 2, 4, 6\}$ records the $g_N$-orbit-length signature on the Nikulin fermionic simple root.
\end{conjecture}

\begin{remark}[Independent of the fixed-subalgebra phrasing]
\label{rem:k3e-independent-fgDel-fixed}
$\mathfrak{g}^{(N)}$ and $\mathfrak{g}_{\Delta_5}^{\mathbb{Z}/N}$ are two different BKM superalgebras associated to the same K3 symplectic $g_N$: $\mathfrak{g}^{(N)}$ via the twined additive lift (Cl\'ery--Gritsenko~$2013$), $\mathfrak{g}_{\Delta_5}^{\mathbb{Z}/N}$ via fixed-point on the ambient untwined algebra. They agree on zeroth and first Fourier orders through the McKay--Thompson trace $f^{(N)}(1, 1) = c_N(0) = \mathrm{tr}_{g_N}(\mathfrak{g}_{\Delta_5, (1, 1, 1)})$; they disagree at higher orders where twisted-sector roots contribute to $\mathfrak{g}^{(N)}$ but are suppressed in $\mathfrak{g}_{\Delta_5}^{\mathbb{Z}/N}$. The Mathieu moonshine relation is the first-order trace identity, not a full isomorphism. Obstructions: (i) the Cartan lattices $\Lambda^{(N)}$ (det $-2(c_N(0) - 2)$ on the singly-twined branch; $\{-32, -24, -18, -12, -6\}$ empirically on the doubly-twined branch) and $\Lambda^{2, 1}_{\mathrm{II}}{}^{g_N}$ (determinant from Nikulin decomposition) are distinct; (ii) the Gritsenko additive lift does not commute with $\mathbb{Z}/N$-averaging; (iii) the fixed-subalgebra denominator is a partial Borcherds product, not the Cl\'ery--Gritsenko lift.
\end{remark}

% ======================================================================
\section{Explicit Fourier table for \texorpdfstring{$\mathfrak{g}_{\Delta_5}$}{g-Delta-5} signed root characters}
\label{sec:k3e-fourier-mult-table}
% ======================================================================

The signed root-character identification for $\mathfrak{g}_{\Delta_5}$ reads \(\operatorname{sdim}\mathfrak g_{\Delta_5,\alpha} = c(4 n m - \ell^2)\). The same product exponent applies at imprimitive lattice points.

\begin{proposition}[First ten non-trivial roots of \texorpdfstring{$\mathfrak{g}_{\Delta_5}$}{g-Delta-5}]
\label{prop:k3e-fourier-mult-table}
On $\Lambda^{2, 1}_{\mathrm{II}}$ with $\alpha = (n, \ell, m)$ and discriminant $D = 4 n m - \ell^2$, the Eichler--Zagier Fourier coefficients of $\phi_{0, 1}$ through $q^3$ tabulate as
\[
\begin{array}{c|ccccccccc}
D & -1 & 0 & 3 & 4 & 7 & 8 & 11 & 12 & 15 \\ \hline
c(D) & 1 & 10 & -64 & 108 & -513 & 808 & -2752 & 4016 & -11775
\end{array}
\]
and the first ten non-trivial roots of $\mathfrak{g}_{\Delta_5}$ ordered by $D$ ascending are
\[
\renewcommand{\arraystretch}{1.12}
\begin{array}{r|ccc|c|c|l|l}
k & n & \ell & m & D = 4 n m - \ell^2 & \gcd & \operatorname{sdim}\mathfrak g_{\Delta_5,\alpha} & \text{character} \\ \hline
1 & 1 & 1 & 0 & -1 & 1 & c(-1) = 1 & \text{real, } \delta_3\text{-type} \\
2 & 1 & 2 & 1 & 0 & 1 & c(0) = 10 & \text{isotropic (null), bosonic} \\
3 & 1 & 1 & 1 & 3 & 1 & c(3) = -64 & \text{fermionic imaginary} \\
4 & 1 & 0 & 1 & 4 & 1 & c(4) = 108 & \text{bosonic imaginary} \\
5 & 2 & 1 & 1 & 7 & 1 & c(7) = -513 & \text{fermionic imaginary} \\
6 & 2 & 0 & 1 & 8 & 1 & c(8) = 808 & \text{bosonic imaginary} \\
7 & 3 & 1 & 1 & 11 & 1 & c(11) = -2752 & \text{fermionic imaginary} \\
8 & 2 & 2 & 2 & 12 & 2 & c(12) = 4016 & \text{imprimitive product exponent} \\
9 & 3 & 0 & 1 & 12 & 1 & c(12) = 4016 & \text{bosonic imaginary} \\
10 & 4 & 1 & 1 & 15 & 1 & c(15) = -11775 & \text{fermionic imaginary}
\end{array}
\]
The sign of $c(D)$ encodes supermultiplicity under the $\mathbb{Z}/2$-graded root system $\Delta = \Delta^{\mathrm{re}} \sqcup \Delta^{\mathrm{im}}_{\bar 0} \sqcup \Delta^{\mathrm{im}}_{\bar 1}$: bosonic at $\ell = 0, \pm 2$, fermionic at $\ell = \pm 1$, with $\dim \mathfrak{g}_{\alpha, \bar 0} - \dim \mathfrak{g}_{\alpha, \bar 1} = c(D)$. Only after this target parity fixture, or an equivalent finite Hall--Borcherds recognition computation, is the ordinary primitive dimension \(|c(D)|\).
\end{proposition}

\begin{proof}[Verification paths]
\emph{Real-root entry $(1, 1, 0)$, $D = -1$}: $c(-1) = 1$ is the weak-Jacobi support edge; corresponds to $\delta_3 = f_3$ real simple root of the fundamental polyhedron.

\emph{Null root $(1, 2, 1)$, $D = 0$}: $c(0) = 10$ matches the isotropic constant term; consistent with $\tau(a_0) = c(0) - c(-1) = 10 - 1 = 9$, the bosonic-null exponent identity.

\emph{Fermionic root $(1, 1, 1)$, $D = 3$}: direct computation gives $\phi_{0, 1}|_{q^1}$ pentad $(10, -64, 108, -64, 10)$ in $(r^{-2}, r^{-1}, r^0, r^1, r^2)$; $c(3)$ sits at $\ell = \pm 1$, reconciling with the automorphic correction $m(a) = -f_{\Delta_5}(1, 1, 1)/64 = -1$ via $f_{\Delta_5}(1, 1, 1) = 64$.

\emph{Borcherds divisor at $(2, 2, 2)$}: the point is imprimitive, but
the product exponent is still read at its discriminant:
\[
 \operatorname{sdim}\mathfrak g_{\Delta_5,(2,2,2)}=c(12)=4016.
\]
The multiple-cover Möbius factor acts on the reduced stable-pair
generating series, not on this root exponent.
\end{proof}

\begin{remark}[\texorpdfstring{$q^1$}{q-1}-pentad normalisation reconciliation]
\label{rem:k3e-q1-pentad-normalisation}
Two normalisations appear in the literature for the $q^1$-pentad: the Eichler--Zagier generator of $J_{0, 1}^{\mathrm{weak}}$ has $(10, -64, 108, -64, 10)$ in $(r^{-2}, r^{-1}, r^0, r^1, r^2)$, while the weight-$12$ index-$1$ Jacobi cusp form $\phi_{12, 1}$ (first Fourier--Jacobi coefficient of $\Delta_{12}$) has $q^2$-pentad $(10, -88, -132, -88, 10)$. The relation $\phi_{0, 1} = \phi_{12, 1}/\delta_{12}$ with $\delta_{12} = q \prod_n (1 - q^n)^{24}$ shifts $q$-degree by $-1$ and reconciles the two at $q^1$-depth:
\[
\phi_{0, 1}|_{q^1} = \phi_{12, 1}|_{q^2} + 24\, \phi_{12, 1}|_{q^1} = (10, -88, -132, -88, 10) + 24\,(0, 1, 10, 1, 0) = (10, -64, 108, -64, 10).
\]
The K3 elliptic genus $\phi_{0, 1}^{K3} = 2 \phi_{0, 1}^{\mathrm{EZ}}$ (Eguchi--Ooguri--Tachikawa~$2011$ \S$3$), so the K3-normalised $q^1$-pentad is $(20, -128, 216, -128, 20)$; the BKM signed-character identification uses the Eichler--Zagier normalisation.
\end{remark}

% ======================================================================
\section{The \texorpdfstring{$\mathrm{CY}_2^{\mathrm{Mp}}$}{CY-2-Mp} domain of \texorpdfstring{$\Phi^{\mathrm{sup}}$}{Phi-sup}}
\label{sec:k3e-CY2Mp-domain}
% ======================================================================

The natural $(\infty, 1)$-categorical domain of the super-paramodular lift $\Phi^{\mathrm{sup}}$ is the $B(\mathbb{Z}/2)$-gerbe $\mathrm{CY}_2^{\mathrm{Mp}} \to \mathrm{CY}_2$ with class $[w_2^{\mathrm{per}}] \in H^2(\mathrm{CY}_2; \mathbb{Z}/2)$ reading off the mod-$2$ reduction of the Mukai discriminant form.

\begin{proposition}[\texorpdfstring{$\mathrm{CY}_2^{\mathrm{Mp}}$}{CY-2-Mp} as metaplectic gerbe]
\label{prop:k3e-CY2Mp-gerbe}
Objects of $\mathrm{CY}_2^{\mathrm{Mp}}$ are triples $(\mathcal{C}, L, \nu)$ with $\mathcal{C} \in \mathrm{CY}_2\text{-}\mathrm{Cat}$ a smooth proper Kontsevich--Soibelman CY$_2$ $(\infty, 1)$-category, $L \subset K_0^{\mathrm{num}}(\mathcal{C}) \otimes \mathbb{R}$ the Mukai rational sublattice carrying Euler pairing, and $\nu: \pi_1(\mathrm{Per}(\mathcal{C}, L)) \to \mathrm{Mp}_4(\mathbb{Z})$ a Weil-representation-compatible metaplectic lift of the period-map monodromy. Morphisms are CY$_2$-Fourier--Mukai kernels intertwining $\nu$ up to the Weil $2$-cocycle $\sigma \in Z^2(\mathrm{Sp}_4; \{\pm 1\})$ of Howe~$1979$ Theorem~$1.4$. The forgetful $\pi: \mathrm{CY}_2^{\mathrm{Mp}} \to \mathrm{CY}_2$ is a $B(\mathbb{Z}/2)$-gerbe with class
\begin{equation}
\label{eq:k3e-wper-formula}
[w_2^{\mathrm{per}}(\mathcal{C})] \;=\; [q_{L_\mathcal{C}} \bmod 2] \;\in\; H^2(\mathrm{CY}_2; \mathbb{Z}/2),
\end{equation}
the $\mathbb{Z}/2$-reduction (Arf invariant) of the Mukai discriminant form.
\end{proposition}

\begin{proof}[Explicit values]
For $\mathcal{C} = D^{\mathrm{b}}(K3)$ with $\widetilde H(K3, \mathbb{Z}) = U^{\oplus 4} \oplus E_8(-1)^{\oplus 2}$ even unimodular, the discriminant form $A_{\widetilde H} = 0$, so $[w_2^{\mathrm{per}}] = 0$: the metaplectic cover admits a canonical section. For $\mathcal{C} = D^{\mathrm{b}}(\mathrm{Hilb}^n K3)$ with Beauville--Bogomolov--Markman Mukai lattice $\widetilde H(\mathrm{Hilb}^n K3) = \widetilde H(K3) \oplus \langle -2(n - 1)\rangle$ (Markman~$2010$ \emph{IMRN}), $A_{\widetilde H(\mathrm{Hilb}^n)} = \mathbb{Z}/2(n - 1)$, and $[w_2^{\mathrm{per}}]$ is $0$ for $n$ odd and $[\alpha_n] \neq 0$ for $n$ even, with $\alpha_n$ Poincar\'e-dual to the Markman generator mod $2$. For $\mathcal{C} = D^{\mathrm{b}}(\mathrm{Enr})$ with Enriques Mukai lattice of rank $12$, the discriminant $A_{\widetilde H(S)} = (\mathbb{Z}/2)^{10} \oplus \mathbb{Z}/2_{\omega_S}$ with Arf invariant $1$ on the $E_8(-2)$ block gives $[w_2^{\mathrm{per}}] \neq 0$: Enriques is the first CY$_2$ input where the metaplectic gerbe is genuinely nontrivial, matching Gritsenko--Nikulin~$1998$ (\emph{J.\ reine angew.\ Math.}~$507$ Proposition~$5.1$) Enriques BKM denominator weight $\kappa_{\mathrm{BKM}}(\Delta_5^{\mathrm{Enr}}) = 5/2$. For a principally polarised abelian surface with $L^2 = 2$, $[w_2^{\mathrm{per}}] = [L \cdot L/2 \bmod 2] \cdot [\eta_A] = [\eta_A]$ the mod-$2$ theta divisor class on $\mathcal{A}_2$ (Igusa~$1972$ \S V.$5$).
\end{proof}

\begin{remark}[Integer-weight subcategory and non-CHL order boundary]
\label{rem:k3e-CYint-CYhalf}
The full $(\infty, 1)$-subcategory
$\mathrm{CY}_2^{\mathrm{int}} \hookrightarrow \mathrm{CY}_2^{\mathrm{Mp}}$
spanned by objects with trivial multiplier $\nu = \mathrm{id}$
(equivalently $\nu$ factoring through
$\mathrm{Mp}_4(\mathbb{Z}) \to \mathrm{Sp}_4(\mathbb{Z})$) covers the
geometric core only at $N \in \{1, 2, 3, 6\}$: at each such $N$ the CHL datum
$(K3, g_N)$ has rank-$3$ Mukai-invariant lattice and primitive Gritsenko lift
$\Phi_N$ of integer weight
$c_N(0)/2 \in \{5, 3, 2, 1\}$. The $N = 4$ entry has half-integral
weight $\tfrac32$ and lies in the multiplier-system stratum, not in
$\mathrm{CY}_2^{\mathrm{int}}$. The non-CHL Nikulin-order boundary
$N \in \{5, 7, 8\}$ carries twined weights $c_N(0)/2 \in \{2, \tfrac32, 1\}$
in the order-indexed Borcherds lane, integer at $N = 5, 8$ and
half-integral with multiplier system at $N = 7$. It is not the
Gritsenko--Cl\'ery half-integral dd-modular boundary. The
half-integral Gritsenko--Cl\'ery rows are the triples
$(4,1;\tfrac12)$ and $(1,4;\tfrac32)$, realised by multiplier systems;
they are indexed by dd-modular type, not by the Nikulin orders
$5,7,8$. The CY$_3$ free quotient $(K3 \times E)/\mathbb{Z}_N$ does
not exist at $N \in \{5, 7, 8\}$ because elliptic curves carry no
automorphism of those orders. Candidate super-VOA boundary realisations
must therefore be stated as multiplier-system or non-CHL host
constructions, not as smooth $K3 \times E$ CHL quotients.
\end{remark}

% ======================================================================
\section{The Jatkar--Sen weight-\texorpdfstring{$0$}{0} boundary at \texorpdfstring{$N = 11$}{N = 11}}
\label{sec:k3e-js-n11-boundary}
% ======================================================================

\begin{proposition}[The \texorpdfstring{$N=11$}{N=11} endpoint is scalar]
\label{prop:k3e-js-n11-boundary}
Let
\[
  N^{\mathrm{JS}}=\{1,2,3,5,7,11\},\qquad
  k_N^{\mathrm{JS}}=\frac{24}{N+1}-2
\]
be the Jatkar--Sen dyon-counting tower.  Let
\[
  N^{\mathrm{BKM}}=\{1,2,3,4,6\},\qquad
  \bigl(\kappa_{\mathrm{BKM}}(\Phi_N^{\mathrm{BKM}})\bigr)
  =\bigl(5,3,2,\tfrac32,1\bigr)
\]
be the primitive Borcherds-denominator tower used in this chapter.
Then \(k_{11}^{\mathrm{JS}}=0\), \(11\notin N^{\mathrm{BKM}}\), and the
weight-zero Jatkar--Sen endpoint is not a fourth BKM sector.

In particular, the notation \(\kappa_{\mathrm{BKM}}(\Phi_{11}^{\mathrm{JS}})\)
is not used as a compact Hall or Borcherds-algebra invariant here.  Such a
use would require a denominator algebra
\(\mathfrak g_{\Phi_{11}}\), a root lattice, a Weyl vector, a Borcherds
product with root-character exponents, and a finite Hall--Borcherds
recognition datum.  These data are not constructed by the weight formula
\(k_N^{\mathrm{JS}}=24/(N+1)-2\).
\end{proposition}

\begin{proof}
The displayed formula gives \(k_{11}^{\mathrm{JS}}=24/12-2=0\).  The
primitive BKM tower is the order set on which the CHL twisted-twined Jacobi input
has Borcherds weights \(c_N(0)/2=(5,3,2,\tfrac32,1)\), namely
\(N=\{1,2,3,4,6\}\).  Hence \(N=11\) lies in the physical dyon-counting
tower but outside the primitive denominator tower.  A zero Siegel weight
does not by itself define a BKM superalgebra: the denominator theorem
needs the Borcherds product and the signed root-character package.  The
geometric \(K3\times E\) CHL quotient model used for the primitive ladder
also has no order-\(11\) symplectic K3 automorphism and no order-\(11\)
elliptic automorphism fixing the origin.  The \(N=11\) datum is therefore
a Narain/CHL scalar boundary, not a smooth \(K3\times E\) BKM
specialisation.
\end{proof}

\begin{remark}[Weight-zero regulator]
\label{rem:k3e-weight-zero-regulator}
Harvey--Moore's one-point Narain regulator
\(\log |\Phi^{\mathrm{CHL}}_{11}|^2\) is a scalar automorphic quantity.
It may serve as a boundary condition for a future realisation problem.
It is not a Stage-$1$ output of \(\PhiFA_3\), not a compact Hall source,
and not a denominator algebra.  A chiral or Heisenberg realisation at
\(N=11\) would have to supply the missing source category, the comparison
map, and the finite recognition datum before any level-$3$ or level-$4$
promotion can be made.
\end{remark}

% ======================================================================
\section{Two order-indexed towers and the separate Gritsenko--Cl\'ery catalogue}
\label{sec:k3e-two-tower-unification}
% ======================================================================

Two order-indexed sets govern $(K3 \times \mathrm{torus})/\mathbb{Z}_N$ automorphic data under the common symbol $\Phi_N$: the Jatkar--Sen tower $N^{\mathrm{JS}} \in \{1, 2, 3, 5, 7, 11\}$ with weight $\mathrm{wt}(\Phi_N^{\mathrm{JS}}) = 24/(N + 1) - 2 \in \{10, 6, 4, 2, 1, 0\}$, and the primitive BKM/CHL denominator tower $N^{\mathrm{BKM}} \in \{1, 2, 3, 4, 6\}$ with $\kappa_{\mathrm{BKM}}(\Phi_N^{\mathrm{BKM}}) = c_N(0)/2 \in \{5, 3, 2, \tfrac32, 1\}$. Their intersection is $\{1, 2, 3\}$; their union is $\{1, 2, 3, 4, 5, 6, 7, 11\}$. The Gritsenko--Cl\'ery catalogue is a separate triple-indexed list of additive-lift forms, distinct from either order tower unless a named overlap theorem fixes the row and normalisation.

\begin{theorem}[Arithmetic separation of the two towers]
\label{thm:k3e-two-tower-unification}
Let \(L_{\mathrm{het}}=\mathrm{II}_{6,22}\) be the Narain lattice of
heterotic string on \(T^6\), and let
\(L_{K3}=H^*(K3,\mathbb Z)=\mathrm{II}_{4,20}\) be the K3 Mukai lattice.
\begin{enumerate}[label=\textup{(\roman*)}]
\item The Jatkar--Sen physical dyon tower has order set
\[
  N^{\mathrm{JS}}=\{1,2,3,5,7,11\},\qquad
  k_N^{\mathrm{JS}}=\frac{24}{N+1}-2=(10,6,4,2,1,0).
\]
\item The primitive BKM denominator tower used by the K3-twined
Borcherds products has order set
\[
  N^{\mathrm{BKM}}=\{1,2,3,4,6\},\qquad
  \kappa_{\mathrm{BKM}}(\Phi_N^{\mathrm{BKM}})=c_N(0)/2=\bigl(5,3,2,\tfrac32,1\bigr).
\]
\item The two towers meet at \(N=\{1,2,3\}\), where the Jatkar--Sen
weights \((10,6,4)\) are twice the primitive BKM weights \((5,3,2)\).
Only at \(N=1\) is
the unnormalised square identity
\[
  \Phi_1^{\mathrm{JS}}=\Phi_{10}^{\mathrm{un}}
  =\Delta_5^2=(\Phi_1^{\mathrm{BKM}})^2
\]
invoked as a proved form-level statement in this chapter; at
\(N=2,3\) the doubling is a weight-level relation.
\item The complement \(N^{\mathrm{JS}}\setminus N^{\mathrm{BKM}}\) is
\(\{5,7,11\}\).  These orders belong to the physical Narain/CHL dyon
tower and do not give smooth \(K3\times E\) BKM specialisations in the
primitive denominator sense used here.
\item The complement \(N^{\mathrm{BKM}}\setminus N^{\mathrm{JS}}\) is
\(\{4,6\}\).  Applying the Jatkar--Sen formula to these orders gives
\(24/5-2=14/5\) and \(24/7-2=10/7\), so the physical dyon tower cannot
be the primitive BKM denominator tower.
\end{enumerate}
\end{theorem}

\begin{proof}
The Jatkar--Sen order set is the level-matched physical dyon set:
\(24/(N+1)\in\mathbb Z_{\geq2}\) gives exactly
\(\{1,2,3,5,7,11\}\), with weights \(k_N^{\mathrm{JS}}\) as displayed
(Jatkar--Sen~$2005$ \S$3$; David--Jatkar--Sen~$2006$ \S$2.2$).  The
primitive BKM constants are the Borcherds weights \(c_N(0)/2\) on the
twisted-twined denominator ladder \(N=\{1,2,3,4,6\}\) in the
square-root normalisation.  Igusa's
unnormalised identity and the Gritsenko--Nikulin product give
\(\Phi_{10}^{\mathrm{un}}=\Delta_5^2\) at \(N=1\).  The remaining
claims are arithmetic: the intersections, complements, and the
weight-level doubling at \(N=1,2,3\) follow from the two displayed
tuples, while the Jatkar--Sen formula does not apply at \(N=4,6\).
\end{proof}

\begin{remark}[Narain source and K3-twined source]
\label{rem:k3e-two-m24-actions}
The Jatkar--Sen tower is selected by level matching in the Narain CHL
dyon problem.  The primitive BKM tower is selected by the CHL
twisted-twined Jacobi input whose Borcherds lift supplies a denominator
product.  These
are compatible sources at \(N=1,2,3\), but they are not one action in
two notations.  The identity \(\Phi_{10}^{\mathrm{un}}=\Delta_5^2\) is
the unorbifolded \(N=1\) square; no corresponding square theorem is used
at \(N=2\) or \(N=3\).  The Gritsenko--Cl\'ery catalogue remains a
separate additive-lift atlas unless a row-by-row normalisation theorem
identifies a specific form.
\end{remark}

\begin{remark}[Open case \texorpdfstring{$N = 5$}{N = 5}]
\label{rem:k3e-n-equal-5-open}
The case $N = 5$ is Mukai-admissible (order-$5$ symplectic K3 automorphism exists; Mukai~$1988$ Table~$1.3$, class $5A$ of $M_{23}$) but not $K3\times E$ BKM-denominator-selectable because no elliptic curve has an order-$5$ automorphism; it is JS-selectable since $24/6 = 4 \in \mathbb{Z}$ and $5$ is prime. The $M_{24}$ Mathieu-twined K3 genus at $N = 5$ (Eguchi--Hikami~$2011$ Table~$1$, Frame shape $1^4 5^4$) realises the $N = 5$ entry on the JS-side as a weight-$2$ BPS index for CHL; whether a Borcherds exponential lift of the $N = 5$ twined genus produces a paramodular BKM denominator compatible with a named Gritsenko--Cl\'ery row on $(K3 \times S)/\mathbb{Z}_5$ for an appropriate non-elliptic fibre $S$ is open (Borcea--Voisin threefold candidate).
\end{remark}

% ======================================================================
\section{First-principles Fourier expansion of the twined K3 elliptic genus}
\label{sec:k3e-c2-direct-Fourier}
% ======================================================================

The universal identity $\kappa_{\mathrm{BKM}}(\Phi_N) = c_N(0)/2$
(Theorem~\ref{thm:borcherds-weight-kappa-BKM-universal}) lives, in the
Mathieu-twined normalisation, on a
single Fourier slot: the $\zeta^0 q^0$-coefficient of the
$g_N$-twined K3 elliptic genus $Z^{(g_N)}_{K3}$. At
$N = 1$ this slot is $20$, the constant term of the EOT K3 elliptic genus
$Z_{K3} = 2\phi_{0, 1}$; the half-genus $\phi_{0,1}$ reads $10$ and
heads the square-root ladder. At the first non-trivial point $N = 2$
the Eichler--Zagier basis and the Gaberdiel--Hohenegger--Volpato
multiplier data pin the expansion to the value $8$, closing the
$N = 2$ row of Lemma~\ref{lem:c-N-zero-closed-form} by direct
computation.

\begin{theorem}[Direct Fourier expansion of \texorpdfstring{$Z^{(2A)}_{K3}$}{Z-(2A)-K3}]
\label{thm:k3e-c2-direct-Fourier}
Let $g_2 \in M_{24}$ be the Mathieu class $2A$, realised as a Nikulin
symplectic involution of a K3 surface (Mukai~$1988$
Proposition~$1.2$, Table~\S$2$; frame shape $\pi_{g_2} = 1^8 2^8$).
The $g_2$-twined $N{=}4$ elliptic genus $Z^{(g_2)}_{K3}(\tau, z)$ is a
weak Jacobi form of weight $0$, index $1$ on $\Gamma_0(2)$ with
$q^0$-expansion
\[
Z^{(2A)}_{K3}(\tau, z)\big|_{q^0} \;=\; 4\zeta^{-1} + 8 + 4\zeta.
\]
In particular $c_2(0) := [\zeta^0 q^0] Z^{(2A)}_{K3} = 8$, closing
\[
\kappa_{\mathrm{BKM}}(\Phi^{(2)}) \;=\; c_2(0)/2 \;=\; 4,
\]
the $N = 2$ instance of
Theorem~\ref{thm:borcherds-weight-kappa-BKM-universal} in the
Mathieu-twined normalisation and the
$N = 2$ row of the Mathieu-twined ladder
$(10, 4, 3, 2, 1)$, realised by the Gritsenko--Nikulin paramodular
cusp form of weight $k(2) = 4$ (GN~$1995$ Part~II Theorem~$2.1$). The
coexisting square-root ladder carries weight $3$ at $N = 2$.
\end{theorem}

\begin{proof}
\emph{Lefschetz trace on the $1^8 2^8$ frame shape.}
The frame-shape encoding sends $\pi_{g_N} = \prod_a a^{m_a}$ to the
frame-exponent sum $T_{H^*}(g_N) = \sum_a m_a$ on the Mukai
lattice $\widetilde{H}(K3, \mathbb{Z}) = \mathrm{II}_{4, 20}$ of
rank~$24$ (Mukai~$1988$ Theorem~$0.4$; Eguchi--Hikami~$2011$
Table~$1$ column~$T_{H^*}$). At $\pi_{g_2} = 1^8 2^8$,
\[
 T_{H^*}(g_2) \;=\; 8 + 8 \;=\; 16,
\]
distinct from the Nikulin fixed-point count $\chi(K3^{g_2}) = 8$
(Mukai~$1988$ \S$2$). The $y \to 1$ specialisation of the
$N{=}4$ elliptic genus recovers the full cohomological trace
(Liu~$1995$ \emph{Invent.~Math.}~$119$ Theorem~A;
Eguchi--Hikami~$2011$ Proposition~$2.1$), giving
$Z^{(g_2)}_{K3}(\tau, 0) = T_{H^*}(g_2) = 16$.

\emph{theta decomposition and $A_2$ extraction.}
The $g_2$-twined elliptic genus admits the index-$1$ theta
decomposition
\[
 Z^{(g_2)}_{K3}(\tau, z)
 \;=\; h^{(2)}_0(\tau)\, \theta_{1, 0}(\tau, z)
 \;+\; h^{(2)}_1(\tau)\, \theta_{1, 1}(\tau, z),
\]
with $h^{(2)}_0, h^{(2)}_1$ weakly-holomorphic vector-valued of
weight $-1/2$ on $\Gamma_0(2)$ (Eichler--Zagier~$1985$ \S$5$;
Eguchi--Hikami~$2011$ \S$3$). The $q^0$-principal-part data pin
$A_2 := [q^0] h^{(2)}_1(\tau) = 4$ (Eguchi--Hikami~$2011$
Table~$1$), equivalently twice the $\zeta^{\pm 1/2}$-coefficient of
the Ramanujan-mock component at $q = 0$, independently computed
from the Gaberdiel--Hohenegger--Volpato~$2012$ \S$3$ twining
character table on the $M_{24}$-class~$2A$ entry.

\emph{the supertrace identity.}
Theorem~\ref{thm:k3e-jacobi-theta-refinement} supplies
$c_N(0) = T_{H^*}(g_N) - 2 A_N$ as a $\mathbb{Z}/2$-graded Lefschetz
supertrace on the Mukai lattice split by index-$1$ Jacobi parity.
At $N = 2$ this reads
\[
 c_2(0) \;=\; T_{H^*}(g_2) - 2 A_2 \;=\; 16 - 8 \;=\; 8,
\]
and the $q^0$-expansion of $Z^{(2A)}_{K3}$ is therefore
$4\zeta^{-1} + 8 + 4\zeta$: the $\zeta^{\pm 1}$-coefficients equal
$A_2 = 4$ and sum with the $\zeta^0$-slot to the frame-exponent
trace, $8 + 2 \cdot 4 = 16 = T_{H^*}(g_2)$.

\emph{Borcherds weight closure.}
Borcherds~$1998$ Theorem~$13.3$ extracts the weight
$\mathrm{wt}(\Phi^{(2)}) = c_2(0)/2 = 4$ from the constant-term
Fourier datum of the singular-theta input. The Gritsenko paramodular
cusp form $\Delta_4^{(2)}$ on $\Gamma_0(2)^+$ has weight
$k(2) = 4$ (Gritsenko--Nikulin~$1995$ Part~II Theorem~$2.1$),
agreeing with the Borcherds extraction. Hence
$\kappa_{\mathrm{BKM}}(\Phi^{(2)}) = c_2(0)/2 = 4$,
the $N = 2$ row of the Mathieu-twined ladder $(10, 4, 3, 2, 1)$; the
square-root ladder $(5, 3, 2, \tfrac32, 1)$ reads $3$ at $N = 2$.
\end{proof}

\begin{remark}[Two constant-term ladders for \texorpdfstring{$c_N(0)$}{c-N(0)}]
\label{rem:k3e-adjudication-two-conventions}
The constant $c_N(0)$ controlling the universal identity is read, in
the Mathieu-twined normalisation, as
the $\zeta^0 q^0$-coefficient of the EOT twined K3 elliptic genus
$Z^{(g_N)}_{K3}$. On the CHL ladder
$N \in \{1, 2, 3, 4, 6\}$ this gives $c_N(0) = (20, 8, 6, 4, 2)$,
matching the Gritsenko--Nikulin twined paramodular-weight ladder
$c_N(0)/2 = (10, 4, 3, 2, 1)$ (GN~$1995$ Part~II
Theorem~$2.1$ with $k(2), k(3), k(4), k(6) = 4, 3, 2, 1$;
Eguchi--Hikami~$2011$ Table~$1$). The coexisting square-root
normalisation reads the David--Jatkar--Sen twisted-twined input,
giving $c_N(0) = (10, 6, 4, 3, 2)$ and the primitive BKM denominator
weights $\kappa_{\mathrm{BKM}}(\Phi_N) = (5, 3, 2, \tfrac32, 1)$
(Jatkar--Sen~$2006$ \emph{JHEP}~$11$, $072$; Govindarajan--Krishna
$2010$ \emph{JHEP}~$05$), headed at $N = 1$ by $\Delta_5$ with the
half-genus value $c_1(0) = 10$ of $\phi_{0, 1}$. The two ladders are
never mixed: the tuple $(10, 8, 6, 4, 2)$ with weights
$(5, 4, 3, 2, 1)$ reads the half-genus at $N = 1$ and the
twined genus at $N \geq 2$; it is a cross-normalisation hybrid,
admissible only with per-column normalisation stamps, not a
single-family ladder. The Igusa-square
convention of the physics literature (Dijkgraaf--Verlinde--Verlinde
$1997$; Jatkar--Sen~$2006$) records
$\Phi^{\mathrm{phys}}_{k_N} = (\Delta_{k(N)}^{(N)})^2$ with
weight $k_N^{\mathrm{phys}} = 2 k(N) = (10, 6, 4, 3, 2)$ in the
square-root normalisation; the BKM
denominator is the square-root
$\Delta_{k(N)}^{(N)} = \sqrt{\Phi^{\mathrm{phys}}_{k_N}}$ of
weight $k(N)$. Within each named normalisation the Borcherds weight
theorem (Borcherds~$1998$ Theorem~$13.3$) extracts
$\mathrm{wt}(\Phi_N) = c_N(0)/2$; the extraction rule is universal,
the input ladder must be named.
\end{remark}

\begin{remark}[Humbert divisor and the absence of an \texorpdfstring{$\mathrm{O}(3,3)$}{O(3,3)} product]
\label{rem:humbert-envelope-restriction}
The Igusa isogeny identifies
\(\mathrm{Sp}_4(\mathbb Z)/\{\pm I\}\) with the arithmetic orthogonal group
of the type-IV lattice \(\Lambda^{3,2}\).  The Humbert divisor carrying
\(\Delta_5\) is a divisor in this type-IV geometry.  A signature-\((3,3)\)
charge lattice may contain a codimension-one \(\Lambda^{3,2}\) slice, but
the Borcherds product is constructed only after this slice has been chosen.
There is no pullback from an \(\mathrm O(3,3)\) Borcherds product.
\end{remark}

\begin{remark}[Weyl-vector central term and the Fourier seed]
\label{rem:fourier-seed-central-value}
The real simple roots $\{\delta_1, \delta_2, \delta_3\}$ of $\mathfrak{g}_{\Delta_5}$ carry Gram matrix with $+2$ on the diagonal and $-2$ off the diagonal; its eigenvalues are $4,4,-2$, hence its signature is $(2, 1)$ and its determinant is $-32$. The Weyl-vector equations $(\rho, \delta_i) = -1$ admit the canonical solution $\rho = \tfrac{1}{2}(\delta_1 + \delta_2 + \delta_3) = f_2 - \tfrac{1}{2}f_3 + f_{-2}$ inside the polar cone of the simple-root chamber. The Fourier seed converting a naive Jacobi Fourier sum into the Borcherds product reads
\[
 \tfrac{1}{64}\,f(n, l, m)\,e^{\pi i(nz_1 + lz_2 + mz_3)} \;=\; m(a)\,e^{-\pi i(\rho + a, z)},\qquad a = (n-1)f_2 - (l-1)\tfrac{1}{2}f_3 + (m-1)f_{-2},
\]
with $m(a) = -\tfrac{1}{64}f(n, l, m) \in \mathbb{Z}$, $m(0) = -1$ normalising the leading Weyl term. The central value $\rho \in \Delta(\Lambda^{2, 1}_{\mathrm{II}})^*_+$ is the integer-weight anchor for the denominator-identity extraction of $\kBKM(\Delta_5) = 5$: three independent paths; Weyl half-vector norm, Borcherds constant $\left.c_N(0)/2\right|_{N=1} = 5$ via $\phi_{0, 1}$, $\wedge^2$-transfer certifying the two $5$'s live on the same modular form; converge at the central value with no free parameter.
\end{remark}

% ======================================================================
\section{K3$\times$E master: five constructions, five \texorpdfstring{$\kappa_\bullet$}{kappa-bullet}, \texorpdfstring{$24$}{24} curve-stalks}
\label{sec:k3e-master-four-constructions}
% ======================================================================

The K3$\times$E geometry is the Vol~III principal example on which the five $\kappa_\bullet$ invariants separate cleanly, each produced by a distinct construction acting on the same CY$_3$. An unsubscripted invariant notation is forbidden; assigning one scalar to \(K3\times E\) without specifying the construction conflates the five distinct constructions: five values $\{0, 0, 3, 5, 24\}$ emerge from five distinct constructions, with the first two sharing the value~$0$ by the compact identity $\sum_q (-1)^q h^{0, q}(X) = \chi(\cO_X)$ while remaining a categorical-Euler reading versus a chiral Hodge-supertrace reading.

\begin{theorem}[Five $\kappa_\bullet(K3 \times E)$ values from five distinct constructions]
\label{thm:four-kappa-k3xE-four-constructions}
Let $Y = K3 \times E$. The subscripted $\kappa_\bullet$ invariants read:
\begin{center}
\renewcommand{\arraystretch}{1.4}
\begin{tabular}{lll}
\hline
$\kappa_\bullet(Y)$ & Value & Construction \\
\hline
$\kappa_{\mathrm{cat}}(Y)$ & $0$ & $\chi(\mathcal{O}_Y) = \chi(\mathcal{O}_{K3}) \cdot \chi(\mathcal{O}_E) = 2 \cdot 0$, K\"unneth on total space \\
$\kappa_{\mathrm{ch}}^{\mathrm{Hodge}}(Y)$ & $0$ & Chiral Hodge supertrace $\sum_q (-1)^q h^{0, q}(Y) = 1 - 1 + 1 - 1$ at $d = 3$ \\
$\kappa_{\mathrm{ch}}^{\mathrm{Heis}}(Y)$ & $3$ & Heisenberg--Mukai specialisation $\mathrm{Sp}^{\mathrm{ch}}_{\Sigma_2, C} \circ \Phi^{FA}_3$ with $\Sigma_2 = K3$, $C = E$ \\
$\kappa_{\mathrm{BKM}}(\Delta_5)$ & $5$ & Borcherds weight $\left.c_N(0)/2\right|_{N=1} = 10/2$ of $\Delta_5$ \\
$\kappa_{\mathrm{fiber}}(Y)$ & $24$ & Mukai-lattice rank of the K3 fibre: $\mathrm{rk}\,\widetilde H(K3, \mathbb{Z}) = 1 + 22 + 1 = 24$ \\
\hline
\end{tabular}
\end{center}
The five values $\{0, 0, 3, 5, 24\}$ arise from five distinct constructions; none is obtainable from any other by a single algebraic step. The first two share the value~$0$ by the compact-CY$_3$ identity $\sum_q (-1)^q h^{0, q}(Y) = \chi(\cO_Y)$, but the underlying constructions are distinct: $\kappa_{\mathrm{cat}}$ is the categorical Euler character of $D^b(\Coh(Y))$, while $\kappa_{\mathrm{ch}}^{\mathrm{Hodge}}$ is the Hodge supertrace of the chiral Stage-$1$ output, with separate behaviour off the compact locus and under chart specialisation.
\end{theorem}

\begin{proof}[Attribution]
$\kappa_{\mathrm{cat}} = 0$: K\"unneth on total space, $\chi(\mathcal{O}_{K3}) = 2$ (Mukai~$1987$, \emph{Nagoya Math.\ J.}~$106$), $\chi(\mathcal{O}_E) = 0$; their product is $0$, \emph{not} $\kappa_{\mathrm{cat}}(K3) = 2$ which is the fibre value. $\kappa_{\mathrm{ch}}^{\mathrm{Hodge}} = 0$: chiral Hodge supertrace $h^{0,0}(Y) - h^{0,1}(Y) + h^{0,2}(Y) - h^{0,3}(Y) = 1 - 1 + 1 - 1 = 0$ on the compact CY$_3$ $Y = K3 \times E$; numerically equal to $\kappa_{\mathrm{cat}}$ by the identity $\sum_q (-1)^q h^{0, q}(X) = \chi(\cO_X)$ on compact CY$_3$, but a distinct construction (Stage-$1$ chiral readout, not categorical Euler). $\kappa_{\mathrm{ch}}^{\mathrm{Heis}} = 3$: Heisenberg--Mukai chiral specialisation reduces the K3 chiral bialgebra by Mukai twist plus pullback to the elliptic base; the $E_1$-chiral algebra on curve $C$ has Heisenberg modular characteristic $3$ (Costello--Paquette~$2018$ arXiv:$1810.06490$ \S$4$). $\kappa_{\mathrm{BKM}}(\Delta_5) = 5$: Gritsenko--Nikulin~$1995$ Part~II Theorem~$2.1$, with $c_1(0)=10$ for the normalized input $\phi_{0,1}$; the EOT identity \(Z^{(1)}_{K3}=2\phi_{0,1}\) belongs to the doubled Igusa-square lane. $\kappa_{\mathrm{fiber}} = 24$: Mukai lattice rank $\widetilde H(K3, \mathbb{Z}) = H^0 \oplus H^2 \oplus H^4 = \mathbb{Z} \oplus \mathrm{II}_{3, 19} \oplus \mathbb{Z}$, rank $1 + 22 + 1 = 24$ (Mukai~$1987$, \S$1$; Huybrechts~$2016$ Chapter~$16$).
\end{proof}

\begin{theorem}[Twenty-four curve-stalks, not twenty-four smooth \texorpdfstring{$\mathbb{C}^3$}{C-3}-points]
\label{thm:k3xE-24-curve-stalks-not-points}
For a generic elliptic fibration $\pi \colon K3 \to \mathbb{P}^1$ (Huybrechts~$2016$ Chapter~$11$), the discriminant $\Delta(\pi) \subset \mathbb{P}^1$ consists of $24$ points with Kodaira fibre $I_1$ at each. The CY$_3$ $Y = K3 \times E$ has local geometry factoring as $24$ \emph{curve-stalks} $F_p \times E$, $p \in \Delta(\pi)$, each of complex dimension $2$, \emph{not} $24$ smooth $\mathbb{C}^3$-points.
\end{theorem}

\begin{proof}[Attribution]
Kodaira~$1963$ (\emph{Ann.\ Math.}~$77$:$563$) classifies singular fibres of elliptic surfaces and proves $\chi_{\mathrm{top}}(K3) = \sum_p \chi_{\mathrm{top}}(F_p) = 24$. Miranda~$1989$ (\emph{The basic theory of elliptic surfaces}) gives the modern treatment. The product $F_p \times E$ is of complex dimension $2$ (not $3$): the fibre $F_p = I_1$ is a nodal rational curve of $\dim_{\mathbb{C}} = 1$, and $E$ is an elliptic curve of $\dim_{\mathbb{C}} = 1$. The local algebra at the Kodaira node times an $E$-point is $\mathbb{C}[[x, y, z]]/(xy)$, a two-sheet crossing times a smooth direction, not $\mathbb{C}^3$.
\end{proof}

\begin{theorem}[Conditional quasi-NCCR character forced by the reduced DT scalar]
\label{thm:quasi-nccr-character-phi10}
Assume a locally-defined Serre-equivariant quasi-NCCR
\(\mathcal{A}^{\mathrm{qNCCR}}(K3 \times E)\) is constructed on the
complement of the \(24\) curve-stalks, with tilting, gluing, Serre
equivariance, and a Behrend-compatible character map
\(\chi_{\mathrm{qNCCR}}\). If that character map is compatible with
the reduced stable-pairs integration of Oberdieck--Pandharipande, then
the character is forced to be
\[
 \chi_{\mathrm{qNCCR}}\bigl(\mathcal{A}^{\mathrm{qNCCR}}(K3 \times E)\bigr)(q,\tilde q,y)
 \;=\;
 -\bigl(\Phi_{10}^{\mathrm{OP}}(\tau,z,\sigma)\bigr)^{-1}
 \;=\;
 -D_5(\tau,z,\sigma)^{-2}
 \;=\;
 -4096\,\Delta_5^{-2}.
\]
Here \(D_5=64^{-1}\Delta_5\) is the monic Borcherds product and
\(\Phi_{10}^{\mathrm{OP}}=D_5^2\) is the OP scalar normalization.
The equality is a character constraint. It is neither the
Serre-equivariant quasi-NCCR construction nor the primitive
Hall--Borcherds recognition theorem. The finite compact Hall source and
radical-quotient double enter through
Theorems~\ref{thm:k3e-reduced-compact-source-finite-hall-windows}
and~\ref{thm:k3e-finite-compact-hall-drinfeld-double}; identifying
their primitive Hall product with
\(U(\mathfrak n_+(\mathfrak g_{\Delta_5}))\) is exactly the finite
recognition problem.
\end{theorem}

\begin{proof}
Oberdieck--Pandharipande~$2016$ (arXiv:$1411.1514$, Theorem~$1$)
prove
\[
 Z^{\mathrm{red}}_{\mathrm{DT}}(K3 \times E)
 =
 -(\Phi_{10}^{\mathrm{OP}})^{-1}.
\]
Maulik--Toda~$2018$ fixes the Behrend-sign orientation. If the
quasi-NCCR and its character map exist with the compatibility stated
above, the character map must evaluate to the same reduced
stable-pairs scalar. Gritsenko--Nikulin~$1998$
(arXiv:alg-geom/$9611028$; see also \cite{GritsenkoNikulin1998})
gives the unnormalised Igusa square
\(\Phi_{10}^{\mathrm{un}}=\Delta_5^2\); the OP scalar convention is
\(\Phi_{10}^{\mathrm{OP}}=D_5^2=4096^{-1}\Delta_5^2\), so the displayed
normalised equality follows. The
argument uses the scalar integration theorem plus the assumed
character-map compatibility; it contains no construction of compact
critical-CoHA correspondences, no primitive Hall bracket, and no
Hall--Borcherds comparison.
\end{proof}

\begin{remark}[Why global NCCR fails on $K3 \times E$: five obstructions]
\label{rem:k3xE-nccr-five-obstructions}
$K3 \times E$ admits no global non-commutative crepant resolution:
\begin{itemize}
\item[\textup{(a)}] The trace splitting $\mathcal End(T)\to\mathcal O_{K3\times E}$ would make $H^{>0}(K3\times E,\mathcal O)$ a direct summand of $\operatorname{Ext}^{>0}(T,T)$; a tilting generator would force this summand to vanish, contrary to $H^1(E,\mathcal O_E)$ and $H^2(K3,\mathcal O_{K3})$.
\item[\textup{(b)}] The derived McKay route (Bridgeland--King--Reid~$2001$) requires a finite linear group with a fixed point; the K3 lattice autoequivalences and the elliptic translations are not such a global orbifold datum.
\item[\textup{(c)}] The HPD route (Kuznetsov~$2007$) is incompatible with the product polarisation.
\item[\textup{(d)}] The Mukai-rigidity route fails off the K3 factor:
\[
\mathrm{Ext}^1(\mathcal{F} \boxtimes \mathcal{O}_E, \mathcal{F} \boxtimes \mathcal{O}_E)
\supset H^1(E, \mathcal{O}_E) \neq 0 .
\]
\item[\textup{(e)}] The finite critical-chart route stops locally. The global $(-1)$-shifted symplectic form on $\mathrm{RPerf}(K3\times E)$ and the Thomas--Behrend--Fantechi obstruction theory exist; they do not produce a single finite quiver-with-potential algebra $\End(T)$ whose representation stack globalises the Darboux charts and carries orientation through K3 spherical-twist monodromy.
\end{itemize}
The Serre-equivariant quasi-NCCR of Theorem~\ref{thm:quasi-nccr-character-phi10} is the locally-defined substitute.
\end{remark}

\begin{remark}[Cartan rank of \texorpdfstring{$\mathfrak{g}_{\Delta_5}$}{g Delta5}]
\label{rem:gdelta5-rank-23-leech-addback}
The Cartan subalgebra of $\mathfrak g_{\Delta_5}$ has rank $3$ and is
carried by the hyperbolic lattice $\Lambda^{2,1}_{\mathrm{II}}$. The
rank-$22$ lattice $\mathrm{II}_{3,19}$ is the K3 transverse quotient
of the Mukai lattice, and the rank-$24$ lattice
$\widetilde H(K3,\mathbb Z)$ is the Mukai-Heisenberg fibre datum. A
Leech addback would move the statement to the Fake-Monster
$II_{25,1}$ denominator and changes the Borcherds input.
\end{remark}

\begin{remark}[Humbert $H_1$ monodromy order $8 = 2 c_+ \nu(\mu_0)$ and Theorem-C complementarity $\kappa_{\mathrm{ch}} + \kappa_{\mathrm{ch}}^! = 8$]
\label{rem:H1-humbert-order-8-decomposition}
The Humbert divisor $H_1$ of discriminant $1$ carries monodromy of order $8$ around the singular locus of $\Delta_5 = 0$, decomposed as
\[
 8 \;=\; 2 \cdot c_+ \cdot \nu(\mu_0), \qquad c_+ = 4, \quad \nu(\mu_0) = 1,
\]
with $c_+ = 4$ the positive Mukai signature component and $\nu(\mu_0) = 1$ the Humbert fundamental-eigenvalue multiplicity. The Mukai lattice $\widetilde H(K3, \mathbb{Z})$ has signature $(4, 20)$; polarising by the positive Mukai four-plane splits it into summands of signatures $(4, 0)$ and $(0, 20)$, and the positive summand supplies $c_+ = 4$. The $2$ prefactor is the $\mathrm{Sp}_4 \to \mathrm{SO}_+(\Lambda^{3, 2})$ double cover. This is the lattice side of the order-eight K3 comparison; the Bruinier and Lusztig sides require their own comparison maps.

The integer $8$ is also the Vol~I Theorem-C complementarity value on the $\mathcal{B}$-family:
\[
 K_{\mathcal B}^{\mathrm{comp}} \;:=\; 2 c_+(\mathrm{Mukai}(K3)) \;=\; 8 \;\in\; \bigl\{0,\; 8,\; 13,\; 250/3,\; 25/3\bigr\},
\]
where $K_{\mathcal B}^{\mathrm{comp}}$ denotes the paired derived-centre complementarity conductor and the bucket $\{0,\, 8,\, 13,\, 250/3,\, 25/3\}$ is the universal canonical five-archetype ceiling (Vol~I Theorem~C; the Bershadsky--Polyakov entry $25/3$ is conditional on its anomaly-ratio lane, and the once-recorded $98/3$ is retracted). The $\mathcal{B}$-row witness is Chapter~\ref{ch:k3-chiral-bialgebra-platonic}, Theorem~\ref{thm:humbert-order-K-kappa}, via Bruinier reciprocity; the Mukai-doubling $2 c_+ = 8$ is CY-$2$ Serre-duality symmetrisation of the bar differential across the Mukai pairing.
\end{remark}

\begin{remark}[Six routes to $G(K3 \times E)$: six distinct constructions, not six $\Phi$-applications]
\label{rem:six-routes-g-k3xe}
The completed quantum vertex chiral group \(G(K3 \times E)\), when the
Hall package is constructed, is approached by six construction routes:
(i) Schiffmann--Vasserot CoHA
for K3 pulled to $K3 \times E$; (ii) Maulik--Okounkov stable envelopes
on $\mathrm{Hilb}(K3)$ base-changed along $E$; (iii)
Gritsenko--Nikulin Borcherds lift of $\phi_{0, 1}$ to $\Delta_5$; (iv)
Costello--Paquette twisted holography on $AdS_3 \times S^3 \times K3$;
(v) Bridgeland--Smith embedding
\(\rho_{\mathrm{BS}}\colon\mathrm{Sp}_4(\mathbb{Z})/\{\pm I\}\hookrightarrow
\mathrm{Aut}_{\mathrm{eq}}(D^b(K3 \times E))/{\sim}\); (vi)
Fourier--Jacobi genus-escalation from $\phi_{0, 1}$ via
Gopakumar--Vafa~$1998$ primitive BPS invariants. The six constructions
are \emph{not} six applications of the CY-to-chiral functor $\Phi$ to a
single object; they are six different functorial procedures on different
starting data whose common chiral object is conditional on the completed
Hall package. Their convergence is
Theorem~\ref{thm:six-routes-k3xE-conditional-convergence}, not a
theorem of functoriality.
Conflating the six routes with six $\Phi$-applications is a category
error.
\end{remark}

\begin{theorem}[Finite-height convergence criterion for the six routes]
\label{conj:six-routes-k3xE-convergence}
\label{thm:six-routes-k3xE-conditional-convergence}
Let
\[
 G^{\mathrm{SV}}_{K3 \times E},\quad
 G^{\mathrm{MO}}_{K3 \times E},\quad
 G^{\mathrm{GN}}_{K3 \times E},\quad
 G^{\mathrm{CP}}_{K3 \times E},\quad
 G^{\mathrm{BS}}_{K3 \times E},\quad
 G^{\mathrm{GV}}_{K3 \times E}
\]
be the six construction outputs named in
Remark~\ref{rem:six-routes-g-k3xe}. Fix the completed
Hall--Drinfeld/Borcherds target
\[
 \mathscr D_{\Delta_5}
 :=
 D^\wedge_\hbar(\mathcal H^+_{K3\times E})
 \simeq
 Y^\wedge_\hbar{}^{\mathrm{super}}(\mathfrak g_{\Delta_5})
\]
under the finite-height hypotheses of
Theorem~\ref{thm:k3e-hall-drinfeld-super-yangian-criterion}. Suppose
there are continuous comparison maps
\[
 \Theta_{\mathrm{SV}},\Theta_{\mathrm{MO}},\Theta_{\mathrm{GN}},
 \Theta_{\mathrm{CP}},\Theta_{\mathrm{BS}},\Theta_{\mathrm{GV}}
 \colon
 G^{(\bullet)}_{K3 \times E}\longrightarrow \mathscr D_{\Delta_5}
\]
whose finite-height restrictions preserve the PBW root grading, the
Serre-dual negative half, nondegenerate Hopf pairing, Borcherds bracket/Serre
kernel, Drinfeld coproduct, allowed centre, and pro-Borcherds
completion/associator class. Then the six construction outputs are
isomorphic in the completed quasi-Hopf superalgebra category over the
Borcherds cone:
\[
 G^{\mathrm{SV}}_{K3 \times E} \;\simeq\; G^{\mathrm{MO}}_{K3 \times E} \;\simeq\; G^{\mathrm{GN}}_{K3 \times E} \;\simeq\; G^{\mathrm{CP}}_{K3 \times E} \;\simeq\; G^{\mathrm{BS}}_{K3 \times E} \;\simeq\; G^{\mathrm{GV}}_{K3 \times E},
\]
with common Borcherds-denominator character
\(-(\Phi_{10}^{\mathrm{OP}})^{-1}=-D_5^{-2}=-4096\,\Delta_5^{-2}\)
(Theorem~\ref{thm:quasi-nccr-character-phi10})
and common Cartan rank $3$ on the $\Delta_5$ denominator
(Remark~\ref{rem:gdelta5-rank-23-leech-addback}).
The scalar input is the \(N=1\) denominator-character identity
Theorem~\ref{thm:lorgat-conj1-N1-complete}. The target theorem is the
Hall--Drinfeld/Super-Yangian criterion,
Theorem~\ref{thm:k3e-hall-drinfeld-super-yangian-criterion}. The
obstruction is the first finite height \(H_0\) at which the route
comparison has a nonzero defect in the sense of
Theorem~\ref{thm:k3e-finite-height-hall-radical-serre-kernel}: a
pairing radical not equal to the Borcherds Cartan radical, an excess
or missing Serre kernel, a coproduct mismatch, an undeclared primitive
central class, or a completion/associator incompatibility.
\end{theorem}

\begin{remark}[Positive-half evidence and doubled obstruction]
The common character is Theorem~\ref{thm:quasi-nccr-character-phi10},
with the primitive denominator relation
\(\Phi_{10}^{\mathrm{un}}=\Delta_5^2\) fixed by the Gritsenko--Nikulin
product normalisation; the OP scalar normalization inserts
\(D_5=64^{-1}\Delta_5\). The CoHA route is measured against the positive-half
target \(U(\mathfrak n_+(\mathfrak g_{\Delta_5}))\) of
Theorem~\ref{thm:k3e-coha}; the MO route supplies the stable-envelope
\(R\)-matrix, the GN route supplies the denominator and root
multiplicities, the Costello--Paquette route supplies the hCS boundary
chiral algebra, the Bridgeland--Smith route supplies the autoequivalence
monodromy, and the GV route supplies the primitive BPS expansion.
The constructed overlaps preserve the root lattice, Cartan rank, and
denominator character. The Hall pairing, PBW flatness, coproduct,
centre, and completion/associator maps are the remaining finite-height
structures needed to identify the six positive halves and then pass
through the Drinfeld double plus centre construction to the displayed
completed quasi-Hopf superalgebra object.
\end{remark}

% ======================================================================
\section{M-theory parent: character, heuristic, and bialgebra conjecture}
\label{sec:m-theory-parent-stratification-k3e}
% ======================================================================

The phrase ``M-theory on $\mathbb{R}_t \times (K3 \times E) \times \mathbb{R}^4$ produces a BPS algebra whose $\Omega$-deformation is the K3 super-Yangian'' compresses four different assertions: the reduced DT scalar, a compact Hall character reading, a five-dimensional physics heuristic, and a doubled bialgebra conjecture. The scalar is proved by reduced DT theory; the other three require additional constructions.

\begin{theorem}[Character-level M-theory parent for $K3 \times E$]
\label{thm:m-theory-parent-character-level-k3e}
The numerical reduced DT character is proved by
Oberdieck--Pandharipande. A compact critical-CoHA character reading
requires the compact Hall construction in the final sentence.
The B-model topological string partition function on $K3 \times E$ has
the numerical character
\[
 Z^{\mathrm{top, B}}_{K3 \times E}(q, \tilde q, y)
 \;=\;
 -\bigl(\Phi_{10}^{\mathrm{OP}}(\tau,z,\sigma)\bigr)^{-1}
 \;=\;
 -D_5^{-2}
 \;=\;
 -4096\,\Delta_5^{-2}.
\]
Its identification with a graded character
\(\chi_{\mathrm{CoHA}(K3\times E)}\) is conditional on constructing the
oriented compact critical CoHA and the factorisation pushforward
compatible with Problem~\ref{op:k3e-compact-coha-hall-borcherds}.
\end{theorem}

\begin{proof}[Proof]
Oberdieck--Pandharipande~$2016$ (arXiv:$1411.1514$, Theorem~$1$)
gives
\(Z^{\mathrm{red}}_{\mathrm{DT}}(K3\times E)=-(\Phi_{10}^{\mathrm{OP}})^{-1}\).
Gritsenko--Nikulin~$1998$ (arXiv:alg-geom/$9611028$; see also
\cite{GritsenkoNikulin1998}) gives
\(\Phi_{10}^{\mathrm{un}}=\Delta_5^2\), while the OP convention uses
\(D_5=64^{-1}\Delta_5\) and
\(\Phi_{10}^{\mathrm{OP}}=D_5^2\). MNOP supplies the DT/GW character comparison
at the numerical level. The upgrade from numerical partition function
to a compact critical-CoHA graded character is exactly the compact Hall
construction named in the statement.
\end{proof}

\begin{remark}[Physics heuristic: M-theory on $K3 \times E$ to 5D $\mathcal{N} = 2$]
\label{rem:m-theory-heuristic-k3e-5d-n2}
The statement ``M-theory on $\mathbb{R}_t \times (K3 \times E) \times \mathbb{R}^4$ at low energies reduces to a 5D $\mathcal{N} = 2$ SCFT on $\mathbb{R}_t \times \mathbb{R}^4$ whose Coulomb branch is the K3$\times$E moduli of wrapped M2-branes'' is \emph{heuristic}: derived from 11D supergravity classical equations of motion plus supersymmetry, where the quantum 5D SCFT is conjectural (rigorous only for small-rank toric examples; Aharony--Hanany--Kol~$1997$). The statement is heuristic at every dimension-reduction step; it is \emph{not} a theorem about a rigorously defined quantum theory.
\end{remark}

\begin{remark}[Metaphor: M5-brane $(0, 4)$-sigma model on $K3$ divisor]
\label{rem:m5-0-4-metaphor-k3}
The statement ``the BPS states of M-theory on $K3 \times E$ are counted by the $(0, 4)$-supersymmetric sigma model on the M5-brane worldvolume wrapping a divisor in $K3 \times E$'' is a \emph{metaphor}: it names a target space for intuition rather than the mathematical construction used in the proof. Maldacena--Strominger--Witten~$1997$ use this metaphor to relate M5-brane BPS indices to the Gritsenko--Nikulin $\Delta_5$; the denominator identity itself follows independently from the singular theta correspondence and Borcherds' product formula.
\end{remark}

\begin{remark}[Synthesis across unconstructed inputs]
\label{rem:m-theory-synthesis-grade-k3e}
The assertion ``M-theory on $K3 \times E$ produces
\(Y_\hbar^{\mathrm{super}}(\mathfrak{g}_{\Delta_5})\) at the level of
\(\Omega\)-protected local operators'' uses the numerical character of
Theorem~\ref{thm:m-theory-parent-character-level-k3e}, the 5D
\(\mathcal N=2\) heuristic of
Remark~\ref{rem:m-theory-heuristic-k3e-5d-n2}, and the M5
\((0,4)\)-sigma-model metaphor of Remark~\ref{rem:m5-0-4-metaphor-k3}.
It becomes a mathematical isomorphism only after the compact critical
CoHA, the protected-operator algebra, and the Hall--Drinfeld comparison
are constructed. Until then, the scalar identity supplies a necessary
character constraint and no bialgebra isomorphism.
\end{remark}

\begin{conjecture}[Full bialgebra-level M-theory parent of the K3 super-Yangian]
\label{conj:m-theory-parent-full-bialgebra-k3e}
The $\Omega$-deformed algebra of protected local operators of M-theory on $\mathbb{R}_t \times (K3 \times E) \times \mathbb{R}^4$, extracted at a point on the Coulomb branch, is the K3 super-Yangian $Y_{\hbar}^{\mathrm{super}}(\mathfrak{g}_{\Delta_5})$ as a quasi-Hopf superalgebra. The isomorphism
\[
 \mathcal{A}^{\mathrm{M}, \Omega}_{\mathrm{prot}}(K3 \times E) \;\cong\; Y_\hbar^{\mathrm{super}}\bigl(\mathfrak{g}_{\Delta_5}\bigr),
 \qquad
 Y_\hbar^{\mathrm{super}}\bigl(\mathfrak{g}_{\Delta_5}\bigr) \;\simeq\; D\bigl(\mathrm{CoHA}(K3 \times E)\bigr)
\]
at the full bialgebra level. The second comparison includes the
Serre-dual negative half, PBW filtration, nondegenerate Hopf pairing,
coproduct, Borcherds bracket/Serre kernel, centre, and completion
compatibilities
of Theorem~\ref{thm:k3e-hall-drinfeld-super-yangian-criterion}; the
character-level shadow is Theorem~\ref{thm:m-theory-parent-character-level-k3e}.
\end{conjecture}

\begin{definition}[Finite bridge packets and the pro-bridge datum]
\label{def:k3e-bridge-datum}
For each finite height \(H\), a \emph{finite bridge packet} is the
following supplied datum:
\[
\mathbb B^{\leq H}_{K3 \times E}
:=
\Bigl(
\mathsf M_H^{X,\mathrm{can}},
D^{\mathrm{red},\leq H}_{\hbar}(K3\times E,\sigma),
\Theta_{A,H},
\mathcal V_H,
\Xi_{\mathrm{scatt}}^{\leq H},
\Xi_{\mathrm{bar}}^{\leq H},
\Xi_{\mathrm{rad}}^{\leq H},
\Xi_{\mathrm{BRST}}^{\leq H},
\Xi_{\mathrm{Yang}}^{\leq H}
\Bigr)
\]
Here \(\mathsf M_H^{X,\mathrm{can}}\) is the finite source packet of
Theorem~\ref{thm:k3e-canonical-finite-source-packets}, and
\(D^{\mathrm{red},\leq H}_{\hbar}(K3\times E,\sigma)\) is the finite
compact Hall--Drinfeld double of
Theorem~\ref{thm:k3e-finite-compact-hall-drinfeld-double}. The
remaining seven entries are extra finite data: the framed target
packet \(\Theta_{A,H}\), the bar packet \(\mathcal V_H\), and the five
comparison packets
\[
\Xi_{\mathrm{scatt}}^{\leq H},\quad
\Xi_{\mathrm{bar}}^{\leq H},\quad
\Xi_{\mathrm{rad}}^{\leq H},\quad
\Xi_{\mathrm{BRST}}^{\leq H},\quad
\Xi_{\mathrm{Yang}}^{\leq H}.
\]
The packet is \emph{realized} at height \(H\) only when all entries are
actual finite vector spaces or finite morphisms in their stated
categories. The chapter constructs the source packet and finite double;
the five comparison packets are not constructed by this definition.

Explicitly,
\[
\Theta_{A,H}^{\mathrm{cand}} := \Theta_A^{\leq H},\qquad
\mathcal V_H^{\mathrm{cand}} := \bigl(B^{\mathrm{ord}}\bigl(A_{K3 \times E}^{(\Sigma_2,C)}\bigr)\bigr)^{\leq H},\qquad
\Xi_{\mathrm{scatt}}^{\leq H,\mathrm{cand}} := \mathcal{S}^{\leq H}_{K3 \times E},
\]
where \(\mathcal{S}^{\leq H}_{K3 \times E}\) is the height-\(H\) truncation
of the scattering diagram on the finite charge lattice.
Here \(B^{\mathrm{ord},\leq H}\) is the degree-\(H\) truncation of the
ordered bar complex, and the other truncations are taken in the same
height filtration used throughout the finite Hall and shadow towers.
More concretely, writing \(D(\gamma)\) for the discriminant of the
charge \(\gamma\),
\[
\Xi_{\mathrm{bar}}^{\leq H,\mathrm{cand}}
:=
\exp\!\Bigl(\sum_{\mathrm{ht}(\gamma)\leq H}\Omega(\gamma)\sum_{k\geq1}\frac{x_{k\gamma}}{k}\Bigr),
\qquad
\Xi_{\mathrm{rad}}^{\leq H,\mathrm{cand}}
:=
\sum_{\substack{\mathrm{ht}(\gamma)\leq H \\ D(\gamma)>0}}
\mathcal K_{\mathrm{rad}}^{\leq H}
\bigl(D(\gamma);\Phi_{10}^{\mathrm{un}},\mathrm{contour}\bigr),
\]
\[
\Xi_{\mathrm{BRST}}^{\leq H,\mathrm{cand}}
:=
H^1_{\mathrm{BRST}}\!\bigl(V_{\Lambda^{2,1}_{II}}\otimes V_{\mathrm{trans}}^{\leq H}\otimes V_{\mathrm{ghost}}\bigr),
\qquad
\Xi_{\mathrm{Yang}}^{\leq H,\mathrm{cand}}
:=
\operatorname{Res}_{z=u}\,V_{\varepsilon}^{\leq H}(\alpha,z).
\]
These are candidate finite-height formulas. They become entries of the
realized bridge packet only after the corresponding geometric,
bar-complex, Rademacher, BRST, or Yangian construction has supplied
the actual finite map and its target transition.

At finite height, the missing comparison maps have the following
types:
\[
\Xi_{\mathrm{scatt}}^{\leq H}\colon
\mathcal H^{\mathrm{comp},\leq H}_{K3\times E}
\longrightarrow
\hat{\C}[\Gamma]^{\leq H},
\qquad
\Xi_{\mathrm{bar}}^{\leq H}\colon
B^{\mathrm{ord},\leq H}\!\bigl(A_{K3 \times E}^{(\Sigma_2,C)}\bigr)
\longrightarrow
\prod_{\mathrm{ht}\leq H}(1-q^n)^{a_n},
\]
\[
\Xi_{\mathrm{rad}}^{\leq H}\colon
\{a_n\}_{n\leq H}
\longrightarrow
\mathcal R_{\mathrm{rad}}^{\leq H}(\Phi_{10}^{\mathrm{un}}),
\qquad
\Xi_{\mathrm{BRST}}^{\leq H}\colon
V_{\Lambda^{2,1}_{II}}\otimes V_{\mathrm{trans}}^{\leq H}\otimes V_{\mathrm{ghost}}
\longrightarrow
H^1_{\mathrm{BRST}}(\cdot),
\]
\[
\Xi_{\mathrm{Yang}}^{\leq H}\colon
V_\varepsilon^{\leq H}(\alpha,z)
\longrightarrow
\operatorname{Res}_{z=u}V_\varepsilon^{\leq H}(\alpha,z).
\]
The chapter has the source packets, the finite double, and the finite
scattering exponent comparison after recognition. It does not yet
construct the five finite natural transformations displayed above.

Each map would need the following compatibility package before it can
be promoted to a theorem:
\begin{itemize}
\item \(\Xi_{\mathrm{scatt}}^{\leq H}\) must be a motivic integration
morphism preserving initial walls, derived walls, and chamber
composition, and it must commute with the height transition maps; once
this morphism exists, Theorem~\ref{thm:k3e-finite-scattering-root-comparison}
identifies its exponent packet with the finite Borcherds denominator.
\item \(\Xi_{\mathrm{bar}}^{\leq H}\) must be a filtered morphism of
ordered bar complexes whose associated graded is the expected Euler
product and whose truncation agrees with the fixed \((\Sigma_2,C)\)
specialisation.
\item \(\Xi_{\mathrm{rad}}^{\leq H}\) must preserve polar data and
Bessel recursion, with a truncation-error bound uniform in the height
filtration.
\item \(\Xi_{\mathrm{BRST}}^{\leq H}\) must be a chain map from the
finite VOA complex to \(H^1_{\mathrm{BRST}}\), compatible with the
ghost number and the finite-height transition system.
\item \(\Xi_{\mathrm{Yang}}^{\leq H}\) must identify residues of the
deformed vertex operators with genuine operators on the cohomology and
must commute with the Hall--Borcherds \(R\)-matrix comparison.
\end{itemize}
These are the exact finite-height inputs still missing from the
chapter.

Equivalently, for every \(H\) the following squares must commute:
\[
\begin{CD}
\mathcal H^{\mathrm{comp},\leq H+1}_{K3\times E}
@>{\Xi_{\mathrm{scatt}}^{\leq H+1}}>>
\hat{\C}[\Gamma]^{\leq H+1}
\\
@V{T_{H+1,H}^X}VV
@VV{R^{\mathrm{scatt}}_{H+1,H}}V
\\
\mathcal H^{\mathrm{comp},\leq H}_{K3\times E}
@>{\Xi_{\mathrm{scatt}}^{\leq H}}>>
\hat{\C}[\Gamma]^{\leq H}
\end{CD}
\qquad
\begin{CD}
B^{\mathrm{ord},\leq H+1}\!\bigl(A_{K3 \times E}^{(\Sigma_2,C)}\bigr)
@>{\Xi_{\mathrm{bar}}^{\leq H+1}}>>
\prod_{\mathrm{ht}\leq H+1}(1-q^n)^{a_n}
\\
@V{T_{H+1,H}^{\mathrm{bar}}}VV
@VV{R^{\mathrm{bar}}_{H+1,H}}V
\\
B^{\mathrm{ord},\leq H}\!\bigl(A_{K3 \times E}^{(\Sigma_2,C)}\bigr)
@>{\Xi_{\mathrm{bar}}^{\leq H}}>>
\prod_{\mathrm{ht}\leq H}(1-q^n)^{a_n}
\end{CD}
\]
\[
\begin{CD}
\{a_n\}_{n\leq H+1}
@>{\Xi_{\mathrm{rad}}^{\leq H+1}}>>
\mathcal R_{\mathrm{rad}}^{\leq H+1}(\Phi_{10}^{\mathrm{un}})
\\
@V{T_{H+1,H}^{\mathrm{rad}}}VV
@VV{R^{\mathrm{rad}}_{H+1,H}}V
\\
\{a_n\}_{n\leq H}
@>{\Xi_{\mathrm{rad}}^{\leq H}}>>
\mathcal R_{\mathrm{rad}}^{\leq H}(\Phi_{10}^{\mathrm{un}})
\end{CD}
\qquad
\begin{CD}
V_{\Lambda^{2,1}_{II}}\!\otimes V_{\mathrm{trans}}^{\leq H+1}\!\otimes V_{\mathrm{ghost}}
@>{\Xi_{\mathrm{BRST}}^{\leq H+1}}>>
H^1_{\mathrm{BRST}}(\cdot)
\\
@V{T_{H+1,H}^{\mathrm{BRST}}}VV
@VV{R^{\mathrm{BRST}}_{H+1,H}}V
\\
V_{\Lambda^{2,1}_{II}}\!\otimes V_{\mathrm{trans}}^{\leq H}\!\otimes V_{\mathrm{ghost}}
@>{\Xi_{\mathrm{BRST}}^{\leq H}}>>
H^1_{\mathrm{BRST}}(\cdot)
\end{CD}
\]
\[
\begin{CD}
V_{\varepsilon}^{\leq H+1}(\alpha,z)
@>{\Xi_{\mathrm{Yang}}^{\leq H+1}}>>
\operatorname{Res}_{z=u}V_{\varepsilon}^{\leq H+1}(\alpha,z)
\\
@V{T_{H+1,H}^{\mathrm{Yang}}}VV
@VV{R^{\mathrm{Yang}}_{H+1,H}}V
\\
V_{\varepsilon}^{\leq H}(\alpha,z)
@>{\Xi_{\mathrm{Yang}}^{\leq H}}>>
\operatorname{Res}_{z=u}V_{\varepsilon}^{\leq H}(\alpha,z)
\end{CD}
\]

\begin{proposition}[Finite bridge-square criterion]
\label{prop:k3e-finite-bridge-square-criterion}
Fix one of the five bridge components
\[
 \bullet\in
 \{\mathrm{scatt},\mathrm{bar},\mathrm{rad},\mathrm{BRST},\mathrm{Yang}\}.
\]
Suppose that finite vector spaces \(S_{H+1}^{\bullet},S_H^{\bullet}\)
and \(T_{H+1}^{\bullet},T_H^{\bullet}\) have been supplied, together
with finite comparison maps and transition maps
\[
 \Xi_{H+1}^{\bullet}\colon S_{H+1}^{\bullet}\to T_{H+1}^{\bullet},
 \qquad
 \Xi_H^{\bullet}\colon S_H^{\bullet}\to T_H^{\bullet},
\]
\[
 P^S_{H+1,H}\colon S_{H+1}^{\bullet}\to S_H^{\bullet},
 \qquad
 P^T_{H+1,H}\colon T_{H+1}^{\bullet}\to T_H^{\bullet}.
\]
Then the height square for the component \(\bullet\) commutes if and
only if
\[
 P^T_{H+1,H}\Xi_{H+1}^{\bullet}
 =
 \Xi_H^{\bullet}P^S_{H+1,H}.
\]
After finite bases are chosen, so that each supplied map is represented
by a rectangular matrix of the indicated source and target dimensions,
and the displayed products are dimension-compatible composites of those
linear maps, this condition is equivalent to
\[
 \operatorname{rank}\!\left(
 P^T_{H+1,H}X_{H+1}^{\bullet}
 -
 X_H^{\bullet}P^S_{H+1,H}
 \right)=0,
\]
where \(X_{H+1}^{\bullet}\) and \(X_H^{\bullet}\) are the matrices of
the supplied comparison maps.
\end{proposition}

\begin{proof}
Both composites are linear maps
\(S_{H+1}^{\bullet}\to T_H^{\bullet}\). The upper-then-right path is
\(P^T_{H+1,H}\Xi_{H+1}^{\bullet}\); the left-then-lower path is
\(\Xi_H^{\bullet}P^S_{H+1,H}\). The square commutes exactly when these
two linear maps are equal. In finite bases, equality of linear maps is
equivalent to the zero matrix for their difference, and over the ground
field this is equivalent to rank zero. Non-rectangular row data, or
rectangular arrays whose dimensions do not define the displayed
composites, do not represent the finite square and are rejected before
the rank test. The executable witness is
\texttt{finite\_bridge\_transition\_square}; it records both composed
matrices and their difference matrix.
\end{proof}

\begin{remark}[Boundary of the finite square criterion]
\label{rem:k3e-finite-bridge-square-boundary}
Proposition~\ref{prop:k3e-finite-bridge-square-criterion} is the exact
finite test for a supplied square. The comparison maps
\(\Xi_H^{\bullet}\), the source transition \(P^S_{H+1,H}\), and the
target transition \(P^T_{H+1,H}\) still have to be constructed in their
own geometric, bar-complex, Rademacher, BRST, or Yangian settings.
\end{remark}

\begin{proposition}[Finite five-component bridge transition criterion]
\label{prop:k3e-five-component-bridge-transition}
Fix a height step \(H+1\to H\). Suppose that the five component
squares of Proposition~\ref{prop:k3e-finite-bridge-square-criterion}
have been supplied for
\[
\bullet\in
\{\mathrm{scatt},\mathrm{bar},\mathrm{rad},\mathrm{BRST},\mathrm{Yang}\}
\]
with common upper and lower heights. Define
\[
\delta^\bullet_{H+1,H}
=
\operatorname{rank}\!\left(
P^{T,\bullet}_{H+1,H}X^\bullet_{H+1}
-
X^\bullet_H P^{S,\bullet}_{H+1,H}
\right).
\]
Then the five supplied comparison maps form a compatible finite
bridge-system transition for \(H+1\to H\) if and only if
\[
\delta^\bullet_{H+1,H}=0
\qquad
\text{for every }
\bullet\in
\{\mathrm{scatt},\mathrm{bar},\mathrm{rad},\mathrm{BRST},\mathrm{Yang}\}.
\]
Equivalently, \(\sum_\bullet\delta^\bullet_{H+1,H}=0\).
\end{proposition}

\begin{proof}
For one component, Proposition~\ref{prop:k3e-finite-bridge-square-criterion}
identifies compatibility with the vanishing of
\(\delta^\bullet_{H+1,H}\). A finite bridge system at the fixed height
step is the product of the five component squares. Equality of two
morphisms in this finite product is componentwise equality. Hence the
system transition exists exactly when the five component equations hold,
which is exactly the displayed vanishing condition. The executable
witness is \texttt{finite\_bridge\_system\_transition\_report}; it
retains the status, gate, defect rank, and matrix data of each supplied
component square. This criterion aggregates supplied finite maps; it
does not construct them.
\end{proof}

\begin{proposition}[Finite Mittag--Leffler exactness gate for bridge systems]
\label{prop:k3e-bridge-ml-exactness-gate}
For every height step \(H+1\to H\), suppose that the five component
squares satisfy
Proposition~\ref{prop:k3e-five-component-bridge-transition}. For each
component \(\bullet\), set
\[
K_H^\bullet=\ker \Xi_H^\bullet,\qquad
I_H^\bullet=\operatorname{im}\Xi_H^\bullet,\qquad
C_H^\bullet=\operatorname{coker}\Xi_H^\bullet .
\]
Assume further that the transition maps
\[
P^{S,\bullet}_{H+1,H}\colon S^\bullet_{H+1}\twoheadrightarrow S^\bullet_H,
\qquad
P^{T,\bullet}_{H+1,H}\colon T^\bullet_{H+1}\twoheadrightarrow T^\bullet_H
\]
are surjective, and that the induced maps
\[
K^\bullet_{H+1}\longrightarrow K^\bullet_H
\]
are surjective. Equivalently, after finite bases are chosen, the
containments
\[
\Xi_H^\bullet P^{S,\bullet}_{H+1,H}(K^\bullet_{H+1})=0,
\qquad
P^{T,\bullet}_{H+1,H}(I^\bullet_{H+1})\subset I^\bullet_H
\]
hold and the following three ranks have the maximal possible values:
\[
\operatorname{rank}P^{S,\bullet}_{H+1,H}=\dim S^\bullet_H,\qquad
\operatorname{rank}P^{T,\bullet}_{H+1,H}=\dim T^\bullet_H,
\]
\[
\operatorname{rank}\!\left(
P^{S,\bullet}_{H+1,H}\big|_{K^\bullet_{H+1}}
\colon K^\bullet_{H+1}\to K^\bullet_H
\right)
=
\dim K^\bullet_H .
\]
The induced image and cokernel transitions then satisfy
\[
\operatorname{rank}\!\left(P^{T,\bullet}_{H+1,H}\Xi^\bullet_{H+1}\right)
=
\operatorname{rank}\Xi^\bullet_H,
\qquad
\operatorname{rank}\!\left(
\overline P^{T,\bullet}_{H+1,H}
\colon C^\bullet_{H+1}\to C^\bullet_H
\right)
=
\dim C^\bullet_H ,
\]
where \(\overline P^{T,\bullet}_{H+1,H}\) is the induced quotient map on
cokernels.
Then, for every component \(\bullet\), the systems
\[
\{K^\bullet_H\}_H,\quad
\{S^\bullet_H\}_H,\quad
\{I^\bullet_H\}_H,\quad
\{T^\bullet_H\}_H,\quad
\{C^\bullet_H\}_H
\]
are Mittag--Leffler. Hence \(R^1\!\varprojlim\) vanishes on them, and
the inverse limit of the finite comparison sequence is exact:
\[
0\longrightarrow
\varprojlim_H K^\bullet_H
\longrightarrow
\varprojlim_H S^\bullet_H
\longrightarrow
\varprojlim_H T^\bullet_H
\longrightarrow
\varprojlim_H C^\bullet_H
\longrightarrow 0 .
\]
\end{proposition}

\begin{proof}
For finite-dimensional vector spaces, a linear map is surjective if and
only if its rank equals the dimension of its target. This gives the two
rank conditions for \(P^{S,\bullet}_{H+1,H}\) and
\(P^{T,\bullet}_{H+1,H}\), and gives the third rank condition after
restricting \(P^{S,\bullet}_{H+1,H}\) to the kernel
\(K^\bullet_{H+1}\). The first displayed containment is automatic from
the commutative-square condition, since
\[
\Xi_H^\bullet P^{S,\bullet}_{H+1,H}
=
P^{T,\bullet}_{H+1,H}\Xi^\bullet_{H+1}
\]
and \(\Xi^\bullet_{H+1}\) vanishes on \(K^\bullet_{H+1}\). The same
identity gives \(P^{T,\bullet}_{H+1,H}(I^\bullet_{H+1})\subset
I^\bullet_H\), which is the second displayed containment; hence the
image and cokernel transition maps are defined.
The source and target systems are Mittag--Leffler by the assumed
surjectivity of \(P^{S,\bullet}_{H+1,H}\) and
\(P^{T,\bullet}_{H+1,H}\). The kernel system is Mittag--Leffler by the
separate kernel-surjectivity hypothesis. The commutative-square
condition sends kernels to kernels and images to images. Since
\(P^{S,\bullet}_{H+1,H}\) is surjective, the induced map
\(I^\bullet_{H+1}\to I^\bullet_H\) is surjective: if
\(y=\Xi_H^\bullet(s)\), choose \(s'\in S^\bullet_{H+1}\) with
\(P^{S,\bullet}_{H+1,H}(s')=s\), and commutativity gives
\(P^{T,\bullet}_{H+1,H}\Xi^\bullet_{H+1}(s')=y\). The cokernel
transition is surjective
because \(P^{T,\bullet}_{H+1,H}\) is surjective. Equivalently, in
finite bases the induced image transition has rank
\[
\operatorname{rank}\!\left(P^{T,\bullet}_{H+1,H}\Xi^\bullet_{H+1}\right)
=
\operatorname{rank}\Xi^\bullet_H .
\]
The quotient map \(\overline P^{T,\bullet}_{H+1,H}\colon
C^\bullet_{H+1}\to C^\bullet_H\) is defined by the image-containment
condition and has rank \(\dim C^\bullet_H\) by the surjectivity of
\(P^{T,\bullet}_{H+1,H}\).
All five systems are therefore Mittag--Leffler finite-dimensional
inverse systems, so
\(R^1\!\varprojlim=0\). Applying the derived-limit long exact sequence
to
\[
0\to K_H^\bullet\to S_H^\bullet\to I_H^\bullet\to0,
\qquad
0\to I_H^\bullet\to T_H^\bullet\to C_H^\bullet\to0
\]
gives the displayed exact sequence after inverse limit. The executable
witnesses are \texttt{finite\_bridge\_exactness\_step\_report} and
\texttt{finite\_bridge\_system\_exactness\_report}; the step report
records the comparison and transition matrices, kernel bases, kernel
images, image-span matrix, cokernel-span matrix, and the five
surjectivity defect coordinates for source, target, kernel, image, and
cokernel.  These are the finite integers from which the displayed
containment and rank conditions are read, and the system report retains
these data componentwise.
\end{proof}

\begin{proposition}[Finite-height towers of the Mittag--Leffler gate]
\label{prop:k3e-bridge-ml-exactness-tower}
Fix one bridge component
\[
\bullet\in
\{\mathrm{scatt},\mathrm{bar},\mathrm{rad},\mathrm{BRST},\mathrm{Yang}\}.
\]
Let
\[
H_m>H_{m-1}>\cdots>H_0
\]
be a finite height tower. Suppose that each adjacent step
\(H_i\to H_{i-1}\), \(1\leq i\leq m\), satisfies
Proposition~\ref{prop:k3e-bridge-ml-exactness-gate}, and suppose that
for each \(1\leq i<m\) the lower comparison sequence of the step
\(H_{i+1}\to H_i\) is the upper comparison sequence of the step
\(H_i\to H_{i-1}\), after the indicated finite bases are identified.
Define the composed source and target transitions by
\[
P^{S,\bullet}_{H_m,H_0}
=
P^{S,\bullet}_{H_1,H_0}
\circ\cdots\circ
P^{S,\bullet}_{H_m,H_{m-1}},
\qquad
P^{T,\bullet}_{H_m,H_0}
=
P^{T,\bullet}_{H_1,H_0}
\circ\cdots\circ
P^{T,\bullet}_{H_m,H_{m-1}},
\]
where composition is read from right to left. Then the
outer square
\[
\Xi_{H_0}^\bullet P^{S,\bullet}_{H_m,H_0}
=
P^{T,\bullet}_{H_m,H_0}\Xi_{H_m}^\bullet
\]
commutes. Moreover the induced source, target, kernel, image, and
cokernel transitions for the finite drop \(H_m\to H_0\) are
well-defined and surjective. Equivalently, the Mittag--Leffler
exactness gate of
Proposition~\ref{prop:k3e-bridge-ml-exactness-gate} holds for the
composed finite drop \(H_m\to H_0\).
\end{proposition}

\begin{proof}
The commutative-square identity composes. Indeed, writing
\[
\Xi_{H_{i-1}}^\bullet P^{S,\bullet}_{H_i,H_{i-1}}
=
P^{T,\bullet}_{H_i,H_{i-1}}\Xi_{H_i}^\bullet
\]
for each adjacent step and substituting successively gives the displayed
outer-square identity.

Surjectivity of the source and target transitions is preserved by
composition. The same argument applies to the image and cokernel
transitions because each adjacent step already defines them and makes
them surjective by
Proposition~\ref{prop:k3e-bridge-ml-exactness-gate}. For kernels, let
\(k_0\in K^\bullet_{H_0}\). Adjacent kernel-surjectivity gives
\(k_1\in K^\bullet_{H_1}\) mapping to \(k_0\), then
\(k_2\in K^\bullet_{H_2}\) mapping to \(k_1\), and so on. After
\(m\) steps one obtains \(k_m\in K^\bullet_{H_m}\) whose image under
\(P^{S,\bullet}_{H_m,H_0}\) is \(k_0\). Thus the composed kernel
transition is surjective. The zero-dimensional case is included, since
the unique map between zero vector spaces is both well-defined and
surjective.

Hence all five finite exactness coordinates appearing in
Proposition~\ref{prop:k3e-bridge-ml-exactness-gate} are stable under
finite composition. The executable witness is
\texttt{finite\_bridge\_exactness\_tower\_report}; it first checks the
common bridge label, height contiguity, and equality of the adjacent
intermediate comparison maps, then multiplies the source and target
transition matrices and reruns the one-step exactness gate on the
outer finite square. Thus the tower report is a recomputation of the
composed drop, not an independent status label.
\end{proof}

The five right-hand transition maps \(R^{\bullet}_{H+1,H}\) are the
compatibility morphisms that must be constructed on the target sides in
order for the inverse-limit theorem to be more than formal notation.
Proposition~\ref{prop:k3e-finite-bridge-square-criterion} supplies the
finite algebraic test once one component square is given, and
Proposition~\ref{prop:k3e-five-component-bridge-transition} aggregates
the five rank-zero tests at the fixed height step.
Proposition~\ref{prop:k3e-bridge-ml-exactness-gate} adds the exactness
gate needed for the comparison sequence to survive inverse limit, while
Proposition~\ref{prop:k3e-bridge-ml-exactness-tower} records that these
adjacent gates propagate to every finite height drop after the
intermediate finite bases have been identified. The geometric content is
the construction of the five comparison maps, the target-side
truncations, and the proof that the resulting finite squares satisfy
those rank-zero and Mittag--Leffler exactness tests.

Their intended meanings are the following:
\begin{itemize}
\item \(R^{\mathrm{scatt}}_{H+1,H}\) should be the truncation of the
wall-product resummation from height \(H+1\) to height \(H\).
\item \(R^{\mathrm{bar}}_{H+1,H}\) should be the truncation of the
bar-Euler product from height \(H+1\) to height \(H\).
\item \(R^{\mathrm{rad}}_{H+1,H}\) should be the truncation of the
Rademacher conductor sum from height \(H+1\) to height \(H\).
\item \(R^{\mathrm{BRST}}_{H+1,H}\) should be the truncation of the
BRST cohomology packet induced by the finite-height VOA filtration.
\item \(R^{\mathrm{Yang}}_{H+1,H}\) should be the truncation of the
vertex-operator residue packet induced by the finite-height mode
filtration.
\end{itemize}
These are target-side transition maps, not yet constructed morphisms.

\begin{proposition}[Finite compatibility implies a pro-bridge datum]
\label{prop:k3e-probridge-inverse-limit-criterion}
Suppose that for every finite height \(H\) the packet
\(\mathbb B^{\leq H}_{K3\times E}\) is realized in the sense of
Definition~\ref{def:k3e-bridge-datum}. Suppose further that, for each
of the five bridge components, source and target transition maps are
supplied for \(H+1\to H\), and that every component square satisfies
the rank-zero condition of
Proposition~\ref{prop:k3e-finite-bridge-square-criterion}, equivalently
the five-component condition of
Proposition~\ref{prop:k3e-five-component-bridge-transition}. Suppose
finally that the hypotheses of
Proposition~\ref{prop:k3e-bridge-ml-exactness-gate} hold for every
component.
Then the inverse limit
\[
\mathbb B^\wedge_{K3 \times E}
=
\varprojlim_H \mathbb B^{\leq H}_{K3 \times E}
\]
exists as a pro-object in the completed comparison category, and its
components are exactly the inverse limits of the finite source packet,
finite double, target packet, bar packet, and five comparison packets.
The component comparison exact sequences remain exact after inverse
limit.
\end{proposition}

\begin{proof}
A realized finite bridge packet is a finite tuple of objects and maps.
The supplied \(H+1\to H\) source and target transitions define a
transition on every constructed component. For the five comparison
components, Proposition~\ref{prop:k3e-finite-bridge-square-criterion}
says exactly that the transition preserves the comparison maps. Thus
the packets and their structure maps form an inverse system in the
completed comparison category. The inverse limit is therefore the
componentwise inverse limit of that finite tuple. No comparison map is
created by this argument; it only assembles the supplied finite maps
after their height squares commute. The final exactness assertion is
Proposition~\ref{prop:k3e-bridge-ml-exactness-gate} applied to each
component.
\end{proof}

\begin{remark}[Provenance of the five finite comparison maps]
\label{rem:k3e-finite-map-provenance}
Each missing finite comparison map already has a formal template in the
chapter, but not a completed construction.
\begin{itemize}
\item \(\Xi_{\mathrm{scatt}}^{\leq H}\) is the finite-height version of
Conjecture~\ref{conj:k3e-scattering-bkm} and its finite-height wall
compatibility remarks.
\item \(\Xi_{\mathrm{bar}}^{\leq H}\) is the finite-height version of
Conjecture~\ref{conj:bkm-bar-dictionary} and the bar Euler comparison
discussion in Section~\ref{subsec:k3e-bps-exponents}.
\item \(\Xi_{\mathrm{rad}}^{\leq H}\) is the finite-height version of
Conjecture~\ref{conj:k3e-shadow-rademacher} and
Remark~\ref{rem:k3e-shadow-rademacher-missing}.
\item \(\Xi_{\mathrm{BRST}}^{\leq H}\) is the finite-height version of
Remark~\ref{rem:k3e-brst-missing} and the BRST template around
equation~\eqref{eq:bkm-brst}.
\item \(\Xi_{\mathrm{Yang}}^{\leq H}\) is the finite-height version of
Conjecture~\ref{conj:vertex-op-yangian} and
Remark~\ref{rem:k3e-yangian-current-missing}.
\end{itemize}
So the remaining work is not to invent new placeholders but to construct
the source Hall--Borcherds object, these five finite-height maps, and
the heightwise transition system, including the source, target, kernel,
image, and cokernel exactness gates that make the inverse-limit
comparison exact.
\end{remark}

A \emph{pro-bridge datum} is the inverse-limit object of a realized
heightwise bridge system,
\[
\mathbb B^\wedge_{K3 \times E} := \varprojlim_H \mathbb B^{\leq H}_{K3 \times E}.
\]
It is the smallest formal container for the comparison chain in this
chapter before the global defect-vanishing conditions are imposed.

In unpacked form, a pro-bridge datum is a tuple
\[
\mathbb B^\wedge_{K3 \times E}
:=
\Bigl(
\Phi_3^{(\Sigma_2,C)},
\mathcal H^{\mathrm{comp}}_{K3 \times E},
\Theta_{K3 \times E},
\mathcal V_{K3 \times E},
\Xi_{\mathrm{scatt}},
\Xi_{\mathrm{bar}},
\Xi_{\mathrm{rad}},
\Xi_{\mathrm{BRST}},
\Xi_{\mathrm{Yang}}
\Bigr)
\]
equipped with the following axioms.
\begin{enumerate}[label=\textup{(BD\arabic*)}]
\item \textbf{Framed source.} $\Phi_3^{(\Sigma_2,C)}$ is a framed $d=3$ stage-1 assignment on $K3 \times E$ whose $(\Sigma_2,C)$ specialisation lands on the framed chiral algebra used later in the chapter.
\item \textbf{Compact Hall source.} $\mathcal H^{\mathrm{comp}}_{K3 \times E}$ is an oriented critical CoHA with negative half, Cartan, Hopf pairing, coproduct, finite-height truncations, and pro-cone completion.
\item \textbf{Scattering map.} $\Theta_{K3 \times E}$ is a motivic integration morphism from wall-crossing Hall data to the quantum torus, compatible with the denominator walls of $\mathfrak g_{\Delta_5}$.
\item \textbf{Bar comparison.} $\mathcal V_{K3 \times E}$ identifies the bar complex of the framed algebra with the $\Lambda^{2,1}_{II}$-graded bar side used in the BKM dictionary, including the Weyl-vector normalisation.
\item \textbf{Shadow control.} $\Xi_{\mathrm{scatt}}, \Xi_{\mathrm{bar}}, \Xi_{\mathrm{rad}}$ implement the shadow partition-function comparison, the finite-height Rademacher resummation, and the truncation-error bounds.
\item \textbf{BRST realisation.} $\Xi_{\mathrm{BRST}}$ provides the worldsheet VOA and BRST differential whose cohomology realises the root multiplicities.
\item \textbf{Yangian comparison.} $\Xi_{\mathrm{Yang}}$ supplies the
BRST-residue operators, Hall--Borcherds \(R\)-matrix comparison, and
Serre-ideal comparison needed to turn the \(\Omega\)-deformed vertex
operators at \(\varepsilon=1\) into algebraic data of the completed
Hall--Borcherds double.
\end{enumerate}
The datum is formal until the finite packets are realized at every
height, the transition squares commute, the Mittag--Leffler exactness
gate of Proposition~\ref{prop:k3e-bridge-ml-exactness-gate} holds, and
existence and uniqueness on the H1--H4 locus are proved.
\end{definition}

\begin{theorem}[Realized finite bridge packets encode the comparison chain]
\label{thm:k3e-bridge-datum-implies-chain}
Assume that the finite bridge packets of
Definition~\ref{def:k3e-bridge-datum} are realized for every finite
height \(H\), that the transition maps make them an inverse system
\(
\mathbb B^{\leq H+1}_{K3 \times E} \to \mathbb B^{\leq H}_{K3 \times E}
\), and that the five bridge squares satisfy
Proposition~\ref{prop:k3e-finite-bridge-square-criterion}. Assume also
the Mittag--Leffler exactness hypotheses of
Proposition~\ref{prop:k3e-bridge-ml-exactness-gate}. Then each
finite-height comparison named in this chapter is one of the supplied
components of \(\mathbb B^{\leq H}_{K3 \times E}\), and these
components are compatible under restriction.

In particular, the comparison chain in this chapter is not a list of
unrelated conjectures: at finite height it is the image of one bridge
packet under the seven structure maps in (BD1)--(BD7), and the global
statement is the inverse-limit version of that chain.
\end{theorem}

\begin{proof}
The entries of \(\mathbb B^{\leq H}_{K3\times E}\) are precisely the
framed source, compact Hall source, scattering map, bar comparison,
shadow/Rademacher comparison, BRST realization, and Yangian comparison
listed in \textup{(BD1)}--\textup{(BD7)}. Since the packet is realized,
these entries are actual finite maps or objects rather than formal
symbols. Since the transition maps make the packets an inverse system
and the five comparison squares commute, restriction from height
\(H+1\) to height \(H\) carries each comparison component to the
corresponding lower component. Passing to the pro-object is exactly the
componentwise inverse limit of Proposition~\ref{prop:k3e-probridge-inverse-limit-criterion}.
The exactness of the resulting component comparison sequences is
Proposition~\ref{prop:k3e-bridge-ml-exactness-gate}.
\end{proof}

\begin{proposition}[Dependency-closed status of the bridge axioms]
\label{prop:k3e-bridge-dependency-closure}
Let \(\mathcal G_{K3\times E}\) be the directed graph with vertices
\[
\begin{gathered}
\mathrm{HB}_{\mathrm{src}},\quad
\mathrm{Rec}_{\mathrm{src}},\quad
\Phi_3,\quad
\mathrm{Hall}_{\mathrm{cpt}},\quad
\mathrm{Scatt},\quad
\mathrm{Bar},\quad
\mathrm{Rad},\quad
\mathrm{BRST},\quad
\mathrm{Yang},
\end{gathered}
\]
and arrows
\[
\mathrm{HB}_{\mathrm{src}}\to \mathrm{Rec}_{\mathrm{src}},\qquad
\mathrm{HB}_{\mathrm{src}}\to \Phi_3,\qquad
\mathrm{Rec}_{\mathrm{src}}\to \mathrm{Hall}_{\mathrm{cpt}},\qquad
\Phi_3\to \mathrm{Hall}_{\mathrm{cpt}},
\]
\[
\mathrm{Hall}_{\mathrm{cpt}}\to \mathrm{Scatt},\qquad
\Phi_3\to \mathrm{Bar},\qquad
\mathrm{Hall}_{\mathrm{cpt}}\to \mathrm{Bar},\qquad
\mathrm{Scatt}\to \mathrm{Bar},
\]
\[
\mathrm{Bar}\to \mathrm{Rad},\qquad
\mathrm{Scatt}\to \mathrm{Rad},\qquad
\Phi_3\to \mathrm{BRST},\qquad
\mathrm{Hall}_{\mathrm{cpt}}\to \mathrm{BRST},
\]
\[
\mathrm{BRST}\to \mathrm{Yang},\qquad
\mathrm{Rad}\to \mathrm{Yang}.
\]
For a vertex \(v\), write \(L(v)\) for the local construction at that
vertex: the relevant finite packets are realized, their transition
squares have rank zero, and the source/target/kernel/image/cokernel
Mittag--Leffler exactness gate holds. For
\(\mathrm{HB}_{\mathrm{src}}\) this local condition also includes the
compact-passage and double-data vanishings
\[
o_{\mathrm{ML}}=o_{\mathrm{real}}=o_{\mathrm{cpt}}
=o_{\Delta}=o_{\mathrm{pair}}=o_{\mathrm{cent}}=0,
\]
the finite recognition defects
\[
\mathcal R_H=\mathcal S_H=\mathcal D_H=\mathcal C_H=\mathcal A_H
=\mathcal P_H=0
\]
at every height, and the pro-recognition gates
\[
Q_H^{\mathrm{sep}}=L_H^{\mathrm{ex}}=H_H^{\mathrm{HB}}=0.
\]
Define the predicate \(C(v)\) recursively by
\[
C(v)\quad\Longleftrightarrow\quad
L(v)\ \text{and}\ C(u)\ \text{for every arrow }u\to v .
\]
Then a bridge axiom at \(v\) is a theorem-level input to the pro-bridge
datum if and only if \(C(v)\) holds. If \(L(v)\) holds while some
predecessor \(u\) is not closed, the construction at \(v\) is a local
finite or pro-finite construction, but it is not a theorem-level bridge
axiom for \(\mathbb B^\wedge_{K3\times E}\).

Consequently the full pro-bridge datum is theorem-level exactly when
\(C(v)\) holds for all vertices of \(\mathcal G_{K3\times E}\). In
particular, a supplied BRST construction with its inverse-limit
exactness does not close \(\mathrm{BRST}\) until \(\Phi_3\) and
\(\mathrm{Hall}_{\mathrm{cpt}}\) are closed, and a supplied Yangian
current construction with its inverse-limit exactness does not close
\(\mathrm{Yang}\) until both \(\mathrm{BRST}\) and \(\mathrm{Rad}\) are
closed.
\end{proposition}

\begin{proof}
The graph is acyclic; the displayed order
\[
\mathrm{HB}_{\mathrm{src}},\ \mathrm{Rec}_{\mathrm{src}},\ \Phi_3,\
\mathrm{Hall}_{\mathrm{cpt}},\ \mathrm{Scatt},\ \mathrm{Bar},\
\mathrm{Rad},\ \mathrm{BRST},\ \mathrm{Yang}
\]
is a topological ordering. The assertion is therefore induction along
this order. At the source vertex, the local condition is exactly the
Hall--Borcherds source, finite-recognition, compact/double,
Mittag--Leffler, and pro-recognition package isolated above. No
predecessor is present, so the source vertex is theorem-level exactly
when that local package is theorem-level.

Assume the claim for all predecessors of \(v\). The finite data in
\(L(v)\) give a constructed component and an exact inverse-limit
sequence at that component by
Propositions~\ref{prop:k3e-finite-bridge-square-criterion},
\ref{prop:k3e-five-component-bridge-transition}, and
\ref{prop:k3e-bridge-ml-exactness-gate}. Those propositions do not
construct the objects feeding the component. The incoming arrows record
precisely those earlier theorem-level inputs: compact Hall promotion
feeds scattering; the framed output, compact Hall source, and scattering
identification feed the bar dictionary; the bar dictionary and
scattering identification feed the Rademacher comparison; the framed
output and compact Hall source feed BRST; BRST and Rademacher feed the
Yangian-current bridge. Hence \(L(v)\) promotes to a theorem-level
bridge axiom exactly when all predecessor predicates \(C(u)\) hold.
This is the displayed recursive condition.

Applying the same induction to every vertex gives the final assertion.
The executable ledger implementing this predicate is
\texttt{k3e\_proof\_dependency\_graph} together with the theorem-schema
and BD-axiom reports in
\texttt{compute/lib/k3e\_finite\_bridge\_witness.py}.
\end{proof}

\begin{remark}[Formal finite-height arrows]
\label{rem:k3e-formal-candidate-finite-maps}
The five comparison arrows have forced finite-height formulas. These
formulas are not proofs of existence or compatibility; the proof must
construct the indicated natural transformations.
\begin{align*}
\Xi_{\mathrm{scatt}}^{\leq H,\mathrm{cand}}
&:=
\pi_{\Gamma,\leq H}\!\left(
\prod_{\substack{\gamma \\ \mathrm{ht}(\gamma)\leq H}}^{\mathrm{lab}}
\mathcal T_\gamma^{\Omega(\gamma)}
\right),
\\
\Xi_{\mathrm{bar}}^{\leq H,\mathrm{cand}}
&:=
\pi_{\mathrm{bar},\leq H}\!\left(
\exp\!\Bigl(\sum_{\mathrm{ht}(\gamma)\leq H}
\Omega(\gamma)\sum_{k=1}^{H}\frac{x_{k\gamma}}{k}\Bigr)
\right),
\\
\Xi_{\mathrm{rad}}^{\leq H,\mathrm{cand}}
&:=
\sum_{\substack{\mathrm{ht}(\gamma)\leq H \\ D(\gamma)>0}}
\mathcal K_{\mathrm{rad}}^{\leq H}
\bigl(D(\gamma);\Phi_{10}^{\mathrm{un}}\bigr),
\\
\Xi_{\mathrm{BRST}}^{\leq H,\mathrm{cand}}
&:=
H^1_{\mathrm{BRST}}\!\bigl(
V_{\Lambda^{2,1}_{II}}\otimes V_{\mathrm{trans}}^{\leq H}\otimes V_{\mathrm{ghost}},
Q_{\mathrm{BRST}}^{\leq H}
\bigr),
\\
\Xi_{\mathrm{Yang}}^{\leq H,\mathrm{cand}}
&:=
\operatorname{Res}_{z=u}\,V_{\varepsilon}^{\leq H}(\alpha,z).
\end{align*}
The corresponding target-side transition maps are the canonical
height truncations on each target system:
\[
R^{\mathrm{scatt}}_{H+1,H,\mathrm{cand}} = \pi_{\Gamma,\leq H},
\qquad
R^{\mathrm{bar}}_{H+1,H,\mathrm{cand}} = \pi_{\mathrm{bar},\leq H},
\qquad
R^{\mathrm{rad}}_{H+1,H,\mathrm{cand}} = \pi_{\mathrm{rad},\leq H},
\]
\[
R^{\mathrm{BRST}}_{H+1,H,\mathrm{cand}} = \pi_{\mathrm{BRST},\leq H},
\qquad
R^{\mathrm{Yang}}_{H+1,H,\mathrm{cand}} = \pi_{\mathrm{Yang},\leq H}.
\]
Their intended meanings are exactly the truncations already named in
the obstruction list: wall-product resummation, bar Euler product,
Rademacher conductor sum, BRST cohomology packet, and vertex-operator
residue packet.
The compact kernel
\(\mathcal K_{\mathrm{rad}}^{\leq H}(-;\Phi_{10}^{\mathrm{un}})\) is
part of the protected \(CY_3\) comparison data.  On the rank-one
Jacobi lane it specialises to the explicit
\(I_{3/2}\)-kernel of
Proposition~\ref{prop:k3e-rank-one-rademacher-packet}; the compact
Siegel kernel is not obtained by replacing that \(I_{3/2}\) formula by
a bare \(I_{9/2}\) expression without the contour normalisation and
polar-data comparison.
\end{remark}

\begin{remark}[Exact proof obligations for the finite-height arrows]
\label{rem:k3e-formal-candidate-finite-map-proof-obligations}
The formulas above become actual constructions only after the following
five lemmas are supplied.
\begin{enumerate}[label=\textup{(\roman*)}]
\item \textbf{Scattering lemma.} For every finite height \(H\), there
exists a motivic integration morphism
\[
\Xi_{\mathrm{scatt}}^{\leq H}\colon
\mathcal H^{\mathrm{comp},\leq H}_{K3\times E}
\longrightarrow
\hat{\C}[\Gamma]^{\leq H}
\]
whose image is exactly the height-\(H\) wall product
\(\prod_{\mathrm{ht}(\gamma)\leq H}^{\mathrm{lab}} \mathcal T_\gamma^{\Omega(\gamma)}\),
and whose compatibility with \(T_{H+1,H}^{X}\) and
\(R^{\mathrm{scatt}}_{H+1,H,\mathrm{cand}}\) preserves the Hall
product, the parity convention, and chamber composition. Under this
lemma and the finite recognition certificate,
Theorem~\ref{thm:k3e-finite-scattering-root-comparison} gives the
finite Borcherds denominator exponent packet.
\item \textbf{Bar lemma.} For every finite height \(H\), the framed
specialisation admits a filtered bar morphism
\[
\Xi_{\mathrm{bar}}^{\leq H}\colon
B_{\leq H}\!\bigl(A_{K3\times E}^{(\Sigma_2,C)}\bigr)
\longrightarrow
\mathcal E_{\mathrm{bar}}^{\leq H}
\]
whose associated graded is the finite Euler product in the candidate
formula, and whose compatibility with \(R^{\mathrm{bar}}_{H+1,H,\mathrm{cand}}\)
commutes with the bar differential and the fixed \((\Sigma_2,C)\)
specialisation. Under this lemma and the finite recognition
certificate, Theorem~\ref{thm:k3e-finite-bar-ce-comparison} identifies
the finite bar Euler exponents with the Borcherds coefficients.
\item \textbf{Rademacher lemma.} For every finite height \(H\), the
compact shadow coefficients must determine a finite-height resummation morphism
\[
\Xi_{\mathrm{rad}}^{\leq H}\colon
\mathrm{Sh}_{\leq H}
\longrightarrow
\mathcal R_{\leq H}
\]
whose rank-one Jacobi specialisation is the packet of
Proposition~\ref{prop:k3e-rank-one-rademacher-packet}.  The remaining
compact theorem is the extension of that polar-data and Bessel
recursion control to the protected \(CY_3\) shadow packet, with a
truncation bound uniform in the height filtration; the projections
\(R^{\mathrm{rad}}_{H+1,H,\mathrm{cand}}\) must preserve the compact
polar data and the finite asymptotic class.
\item \textbf{BRST lemma.} For every finite height \(H\), there exists a
finite-height VOA complex
\[
V_{\Lambda^{2,1}_{II}}\otimes V_{\mathrm{trans}}^{\leq H}\otimes
V_{\mathrm{ghost}}
\]
with differential \(Q_{\mathrm{BRST}}^{\leq H}\) and a degree-one
cohomology morphism
\[
\Xi_{\mathrm{BRST}}^{\leq H}\colon
\bigl(V_{\Lambda^{2,1}_{II}}\otimes V_{\mathrm{trans}}^{\leq H}\otimes
V_{\mathrm{ghost}},Q_{\mathrm{BRST}}^{\leq H}\bigr)
\longrightarrow
H^1_{\mathrm{BRST}}(\cdot)
\]
such that \(R^{\mathrm{BRST}}_{H+1,H,\mathrm{cand}}\) commutes with the
BRST differential and the ghost-number grading.
\item \textbf{Yangian lemma.} For every finite height \(H\), the
deformed vertex operators \(V_\varepsilon^{\leq H}(\alpha,z)\) exist as
genuine finite-height operators to all orders in the \(\Omega\)-deformation
at \(\varepsilon=1\), and their residue morphism
\[
\Xi_{\mathrm{Yang}}^{\leq H}\colon
V_\varepsilon^{\leq H}(\alpha,z)
\longrightarrow
\operatorname{Res}_{z=u}\,V_\varepsilon^{\leq H}(\alpha,z)
\]
lands on cohomological classes rather than formal coefficients; the
compatibility \(R^{\mathrm{Yang}}_{H+1,H,\mathrm{cand}}\) commutes with the
OPE, the Hall--Borcherds pairing, and the Serre ideal.
\end{enumerate}
These lemmas are the missing mathematics. Until they are proved, the
formulas in Remark~\ref{rem:k3e-formal-candidate-finite-maps} are
finite-height arrows forced by the gradings and normalisations; they are
not morphisms of the inverse system.
The finite evidence fixes four boundary values of the desired theorem:
the rank-one bar Euler sector fixes the constant \(24\) exponent, the
BRST packet fixes the central-charge cancellation \(3+23-26=0\), the
Rademacher packet fixes the first growth regime, and the current packet
fixes the residue template for the Yangian-current generators. These
are statements on finite windows. They become a theorem only after they
are natural in the Borcherds-cone inverse system.

The obstruction data therefore form a pro-diagram. The source side
first requires the compact-passage vanishings
\[
 o_{\mathrm{ML}}=o_{\mathrm{real}}=o_{\mathrm{cpt}}=0
\]
and the double-data vanishings
\[
 o_\Delta=o_{\mathrm{pair}}=o_{\mathrm{cent}}=0.
\]
These six equalities are source-side gate conditions: the first three
assert exact compact passage from finite Rees Hall data to compact
critical Hall convolution, and the last three assert the continuous
coproduct, completed Hopf pairing, and centre control needed for the
Drinfeld double.
Only after those six classes vanish do the finite recognition defects
\(\mathcal R_H,\mathcal S_H,\mathcal D_H,\mathcal C_H,\mathcal A_H,
\mathcal P_H\)
measure the remaining Hall--Borcherds comparison at height \(H\).
There is then a separate pro-recognition gate
\[
 (Q_H^{\mathrm{sep}},L_H^{\mathrm{ex}},H_H^{\mathrm{HB}}):
 \quad
 Q_H^{\mathrm{sep}}=\text{separated pro-cone completion},\quad
 L_H^{\mathrm{ex}}=\text{transition-compatible exact inverse limit of the defect ideals},
 \quad
 H_H^{\mathrm{HB}}=\text{heightwise Heegner--Borcherds coefficient comparison}.
\]
These three classes live over the finite recognition quotient, not over
any single comparison component: \(Q_H^{\mathrm{sep}}\) is topological,
\(L_H^{\mathrm{ex}}\) is a derived-limit condition on the defect
ideals, and \(H_H^{\mathrm{HB}}\) is the automorphic coefficient
comparison tying the quotient to the Borcherds denominator.
The componentwise Mittag--Leffler gate of
Proposition~\ref{prop:k3e-bridge-ml-exactness-gate} controls exactness
of the five supplied comparison sequences; it does not by itself prove
\((Q_H^{\mathrm{sep}},L_H^{\mathrm{ex}},H_H^{\mathrm{HB}})\).

\begin{proposition}[Finite source pro-recognition matrix gate]
\label{prop:k3e-source-pro-recognition-matrix-gate}
Fix an adjacent height step \(H+1\to H\) after the finite recognition
quotient has been formed. Let
\[
 P^{\mathrm{comp}}_{H+1,H}\colon C_{H+1}\to C_H
\]
be the supplied transition on the finite completion quotient, and let
\(\mathcal S_{H+1,H}\) be the supplied separation-defect matrix. Let
\[
 P^{\mathrm{def}}_{H+1,H}\colon J_{H+1}\to J_H
\]
be the supplied transition on the finite defect ideals, and let
\(\mathcal L_{H+1,H}\) be the matrix measuring failure of the image of
\(J_{H+1}\) to land in \(J_H\). Finally let
\[
 h_H,\ b_H\in \bQ^{r_H}
\]
be the finite Heegner and Borcherds coefficient rows in the same
height-\(H\) basis. Then the finite pro-recognition gate closes at this
height step exactly when
\[
 \operatorname{rank} P^{\mathrm{comp}}_{H+1,H}=\dim C_H,\qquad
 \mathcal S_{H+1,H}=0,
\]
\[
 \operatorname{rank} P^{\mathrm{def}}_{H+1,H}=\dim J_H,\qquad
 \mathcal L_{H+1,H}=0,
\]
and
\[
 h_H=b_H .
\]
These are respectively the finite matrix forms of
\[
 Q_H^{\mathrm{sep}}=0,\qquad
 L_H^{\mathrm{ex}}=0,\qquad
 H_H^{\mathrm{HB}}=0.
\]
\end{proposition}

\begin{proof}
In finite dimension a transition map is surjective precisely when its
rank equals the dimension of its target. The first rank equality gives
the finite separated-completion condition once the supplied residual
separation matrix \(\mathcal S_{H+1,H}\) vanishes. The defect ideals
form an exact inverse system at this height step exactly when the
transition lands in the lower ideal and is surjective onto it; these are
the two finite conditions for \(P^{\mathrm{def}}_{H+1,H}\) and
\(\mathcal L_{H+1,H}\).

The Heegner--Borcherds comparison is not a pro-linear-algebra
consequence. It is the equality, in the chosen coefficient basis, of
the Heegner divisor multiplicity row and the Borcherds input row. Thus
it is the finite vector equality \(h_H=b_H\). The executable witness
\texttt{finite\_source\_pro\_recognition\_matrix\_gate} records the two
transition ranks, the two defect ranks, and the coefficient-difference
rank. It checks a supplied finite pro-recognition packet; it does not
construct the separated completion, the all-height defect-ideal
inverse limit, the automorphic Heegner comparison, or the completed
Hall--Borcherds source theorem.
\end{proof}

The comparison side then asks for five natural transformations:
scattering to the quantum torus, the bar Euler map, the shadow
Rademacher map, the BRST cohomology map, and the vertex-operator
current map. Each transformation must make its height square commute
for \(H+1\to H\) and must satisfy the source, target, kernel, image,
and cokernel exactness gate of
Proposition~\ref{prop:k3e-bridge-ml-exactness-gate}. The finite checks in
\texttt{compute/lib/hall\_borcherds\_gate.py} and
\texttt{compute/lib/k3e\_finite\_bridge\_witness.py} witness the
finite constants and the named obstruction coordinates; they do not
construct the missing natural transformations. The bridge closes
precisely when these source and comparison squares assemble to an exact
Mittag--Leffler pro-system.
\end{remark}

\begin{conjecture}[Existence and uniqueness of the pro-bridge datum]
\label{conj:k3e-bridge-datum-existence}
On the H1--H4 locus there exists a pro-bridge datum
\(\mathbb B^\wedge_{K3 \times E}\) satisfying axioms
\textup{(BD1)}--\textup{(BD7)} at every finite height, unique up to
contractible choice, with the rank-zero transition squares and
Mittag--Leffler exactness gate of
Proposition~\ref{prop:k3e-bridge-ml-exactness-gate}. Its induced
inverse-limit comparison chain recovers the reduced DT character
identity and the Borcherds denominator identity at the rank-one shadow
level.
\end{conjecture}

\begin{remark}[Dependency summary]
\label{rem:k3e-stratification-summary}
The five statements above separate by construction:
\begin{itemize}
\item \emph{Numerical scalar}: Theorem~\ref{thm:m-theory-parent-character-level-k3e} gives the reduced DT character; compact-CoHA character reading requires the compact Hall construction.
\item \emph{Heuristic}: Remark~\ref{rem:m-theory-heuristic-k3e-5d-n2} (physics reduction to 5D $\mathcal{N} = 2$, classical).
\item \emph{Metaphor}: Remark~\ref{rem:m5-0-4-metaphor-k3} (M5-brane $(0, 4)$-sigma model, UV-incomplete).
\item \emph{Synthesis}: Remark~\ref{rem:m-theory-synthesis-grade-k3e} records the required bridge data.
\item \emph{Conjecture}: Conjecture~\ref{conj:m-theory-parent-full-bialgebra-k3e} asks for the full Hopf-superalgebra isomorphism; the scalar character is only its necessary numerical shadow.
\end{itemize}
The single phrase ``M-theory parent of the K3 super-Yangian'' conflates all five; the displayed dependencies keep the scalar, compact Hall, physics, and doubled bialgebra claims separate.
Definition~\ref{def:k3e-bridge-datum} packages the remaining comparison chain into a finite inverse system and its pro-limit, and Theorem~\ref{thm:k3e-bridge-datum-implies-chain} states the induced compatibilities.

The remaining open proof obligations named earlier in this chapter are the compact Hall promotion, the framed $d=3$ assignment, the scattering-to-root-system identification, the BKM--bar dictionary, the shadow partition-function and shadow--Rademacher comparisons, the BRST realization, and the vertex-operator/Yangian current comparison. Each is isolated in its own remark or conjecture; none is promoted here.
\end{remark}

\begin{remark}[Core proof gaps and the exact missing mathematics]
\label{rem:k3e-core-proof-gaps}
The chapter can only be upgraded to a complete theorem package after the following constructions are supplied.
The dependency order is not arbitrary: items~(i)--(ii) supply the source-level geometry and Hall input; items~(iii)--(iv) turn that input into comparison theorems; items~(v)--(vii) close the analytic and BRST/Yangian closures that propagate the comparison through the shadow tower and the completed double.
Definition~\ref{def:k3e-bridge-datum} packages the same seven items heightwise into the finite bridge packets \(\mathbb B^{\leq H}_{K3 \times E}\); the missing global step is the exact inverse-limit comparison, including the Mittag--Leffler gate of Proposition~\ref{prop:k3e-bridge-ml-exactness-gate}.
Within this bridge formalism, the closure condition is a predicate on
supplied data: the source Hall--Borcherds gate, the recognition envelope,
the seven bridge constructions, the componentwise inverse-limit
exactness gate, and the pro-recognition gate must all be present as
proof objects before the inverse-limit comparison can be called a
theorem.
\begin{enumerate}[label=\textup{(\roman*)}]
\item \textbf{Framed $d=3$ assignment.} Construct a genuine stage-1 factorisation algebra on $K3 \times E$ and prove that its fixed $(\Sigma_2,C)$ specialisation agrees with the framed chiral algebra used later in the chapter, compatibly with the H1--H4 locus and inverse limits.
\item \textbf{Compact Hall promotion.} Build the oriented critical CoHA with negative half, Cartan, Hopf pairing, coproduct, and finite-height truncations; prove nondegeneracy modulo the Borcherds Cartan radical, identify the Serre kernel, and show compatibility with the pro-cone topology.
\item \textbf{Scattering/root identification.} Construct the motivic integration map from wall-crossing Hall data to the quantum torus and prove that wall products correspond exactly to the Borcherds denominator walls, with no extraneous walls and no missing imaginary-root orbits.
\item \textbf{BKM--bar dictionary.} Prove that the bar complex of the framed algebra is $\Lambda^{2,1}_{II}$-graded, that its $\alpha$-primary Euler characteristic equals the motivic DT index, and that the regularization vector matches the Weyl vector, including the Borcherds normalization and multiplier system.
\item \textbf{Shadow partition and Rademacher comparison.} Extend the rank-one finite Rademacher packet of Proposition~\ref{prop:k3e-rank-one-rademacher-packet} to the protected compact shadow packet, with explicit compact polar data, Bessel recursion, and uniform all-height truncation bounds, so that the shadow tower really resums to the Borcherds / Rademacher asymptotics rather than merely matching low coefficients.
\item \textbf{BRST realization.} Construct the worldsheet VOA and BRST differential whose cohomology realises the finite supertrace fixture of Proposition~\ref{prop:k3e-brst-finite-supertrace-fixture}, with the transverse sector and ghost-number grading fixed once and for all.
\item \textbf{Vertex-operator Yangian.} Prove that the $\Omega$-deformed vertex operators exist to all orders at the formal specialisation $\varepsilon=1$, that the residues are genuine operators on the BRST cohomology, and that the OPE relations generate exactly the Serre ideal and the Hall--Borcherds $R$-matrix data.
\end{enumerate}
These are the mathematical inputs missing from the current manuscript. Without them, the chapter contains correct conditional bridges and numerical shadows, but not the unconditional isomorphisms or constructions they point toward.
\end{remark}
